\documentclass[10pt]{article}
\usepackage[a4paper,margin=1.8cm]{geometry}

\usepackage{amsmath,amssymb,amsfonts,mathtools}
\usepackage{braket}
\usepackage{slashed}
\usepackage{cancel}
\usepackage[normalem]{ulem}
\usepackage{xcolor}
\usepackage{framed}
\usepackage{needspace}
\usepackage{enumitem}
\usepackage{longtable,booktabs,array}
\usepackage[hidelinks]{hyperref}
\allowdisplaybreaks
\sloppy
\emergencystretch=3em
\definecolor{notebar}{rgb}{0.75,0.15,0.15}
\definecolor{dkgreen}{rgb}{0.1,0.45,0.1}
\newenvironment{notebox}{\def\FrameCommand{{\color{notebar}\vrule width 1.5pt}\hspace{6pt}}\MakeFramed{\advance\hsize-\width\FrameRestore}\small}{\endMakeFramed}
\newcommand{\CGCket}{|{\rm CGC}\rangle}
\newcommand{\CGCbra}{\langle {\rm CGC}|}
\newcommand{\dbar}{d\hspace{-0.08em}\bar{}\hspace{0.1em}}
\def\wt#1{\widetilde{#1}}
\def\mf#1{\mathbb{#1}}
\def\v#1{\mathbf{#1}}
\def\gsim{ \,\, \vcenter{\hbox{$\buildrel{\displaystyle >}\over\sim$}} \,\,}
\def\lsim{ \,\, \vcenter{\hbox{$\buildrel{\displaystyle <}\over\sim$}} \,\,}
\newcommand{\A}[3]{\mf{A}^{#2}_{#3}(#1)}
\newcommand{\Bone}[3]{\mf{B}^{#2}_{1 #3}(#1)}
\newcommand{\Btwo}[3]{\mf{B}^{#2}_{2 #3}(#1)}
\newcommand{\Bthree}[3]{\mf{B}^{#2}_{3 #3}(#1)}
\newcommand{\Cone}[3]{\mf{C}^{#2}_{1 #3}(#1)}
\newcommand{\Ctwo}[3]{\mf{C}^{#2}_{2 #3}(#1)}
\newcommand{\Cthree}[3]{\mf{C}^{#2}_{3 #3}(#1)}
\newcommand{\Done}[3]{\mf{D}^{#2}_{1 #3}(#1)}
\newcommand{\Dtwo}[3]{\mf{D}^{#2}_{2 #3}(#1)}
\newcommand{\Dthree}[3]{\mf{D}^{#2}_{3 #3}(#1)}

\begin{document}
\title{Single inclusive gluon production at NLO:\\ record of all changes to the tex files and of the calculations behind them}
\author{}
\date{Compiled 20 November 2026}
\maketitle

\begin{abstract}
\noindent This document records every change that was made to the author's tex files while the calculation was checked, file by file and line by line, together with the reason and the numerical evidence that the checked copies state for each change, and it collects the derivations and verifications on which the substantive changes rest. Nothing here is new physics: every statement is taken from one of the files named in the text (checked copies, analysis documents, literature notes), and the file and line are given. Section~\ref{sec:howto} explains how to read the record and lists the global decisions; Section~\ref{sec:summary} is a summary table of the physics-relevant changes; Section~\ref{sec:files} lists every change of every file (old $\to$ new, rendered as formulas); Section~\ref{sec:calc} summarises the calculations and verifications; Section~\ref{sec:open} lists the open points and contains the sections on the problems P1--P6 that are being fixed now. A machine-readable list of all changes is in \texttt{changes\_list.csv} next to this file.
\end{abstract}

\tableofcontents
\clearpage

\setcounter{section}{-1}
\section{How to read this record}\label{sec:howto}

\subsection{Files}
\noindent Every file of the author that was checked has a \emph{checked copy}. The pairs (author's original $\to$ checked copy) are listed in Table~\ref{tab:pairs}. The originals are the author's uploads in \texttt{/root/.claude/uploads/0ab97fcf-94a2-5aec-8bad-586400481685/}; the checked copies are under \texttt{/home/user/correlations/nlo\_review/}. Except for \texttt{Limiting\_cases\_final.tex} (one line fewer than the original) and \texttt{Limiting\_cases\_final\_1\_checked.tex} (one line more), every checked copy has exactly the line numbering of the original, so a line number in this record refers to both files.

\begin{table}[h]
\centering\small
\begin{tabular}{@{}lll@{}}
\toprule
file & original (upload prefix) & checked copy (under \texttt{nlo\_review/})\\
\midrule
Calculation\_main & \texttt{99a3a399-Calculation\_main\_1.tex} & \texttt{calc\_main\_checked/Calculation\_main\_1\_checked.tex}\\
Group I & \texttt{0073ca94-Group\_I.tex} & \texttt{group\_I\_checked/Group\_I\_checked.tex}\\
Group II & \texttt{08bb31a9-Group\_II.tex} & \texttt{group\_II\_checked/Group\_II\_checked.tex}\\
Group II Row IV & \texttt{c62bee3b-Group\_II\_Row\_IV.tex} & \texttt{group\_II\_checked/Group\_II\_Row\_IV\_checked.tex}\\
Group III & \texttt{d98e045a-Group\_III.tex} & \texttt{group\_III\_checked/Group\_III\_checked.tex}\\
Four-$\rho$ & \texttt{43078eec-4rho.tex} & \texttt{four\_rho\_checked/4rho\_checked.tex}\\
Evolution, two $\rho$ & \texttt{294d9dfb-Evolution\_two\_rho.tex} & \texttt{evolution\_checked/Evolution\_two\_rho\_checked.tex}\\
Evolution, three $\rho$ & \texttt{6a6a9fd3-Evolution\_three\_rho.tex} & \texttt{evolution\_checked/Evolution\_three\_rho\_checked.tex}\\
Evolution (JIMWLK) terms & \texttt{4af0d15a-Evolution\_JIMWLK\_\_terms.tex} & \texttt{evolution\_checked/Evolution\_JIMWLK\_terms\_checked.tex}\\
BFKL & \texttt{9da5dc76-BFKL.tex} & \texttt{evolution\_checked/BFKL\_checked.tex}\\
Other terms (v1) & \texttt{99883d04-Other\_terms.tex} & \texttt{evolution\_checked/Other\_terms\_checked.tex}\\
Other terms (v2) & \texttt{ff9ece59-Other\_terms.tex} & \texttt{other\_terms\_v2/Other\_terms\_v2\_checked.tex}\\
Divergent terms & \texttt{371e60d2-Divergent\_terms.tex} & \texttt{limits\_checked/Divergent\_terms\_checked.tex}\\
scales & \texttt{8999d151-scales.tex} & \texttt{limits\_checked/scales\_checked.tex}\\
Limiting cases (flags) & \texttt{f164e8ee-Limiting\_cases.tex} & \texttt{limits\_checked/Limiting\_cases\_checked.tex}\\
Limiting cases (decisions) & \texttt{f164e8ee-Limiting\_cases.tex} & \texttt{limits\_checked/Limiting\_cases\_final.tex}\\
Limiting cases final 1 & \texttt{7daa274b-Limiting\_cases\_final\_1.tex} & \texttt{limits\_checked/Limiting\_cases\_final\_1\_checked.tex}\\
Final terms & \texttt{b2466302-Final\_terms.tex} & \texttt{final\_checked/Final\_terms\_checked.tex}\\
\bottomrule
\end{tabular}
\caption{The author's files and their checked copies, in the order of Section~\ref{sec:files}.}\label{tab:pairs}
\end{table}

\noindent Supporting analysis documents (derivations and numbers, no author's original): \texttt{limits\_checked/Other\_terms\_limits.tex}, \texttt{limits\_checked/Cancellation\_of\_flagged\_terms.tex}, \texttt{limits\_checked/Finite\_terms\_final.tex}; literature comparisons: \texttt{paper/notes/lit\_lo\_bfkl\_jimwlk.tex}, \texttt{lit\_dglap\_rc.tex}, \texttt{lit\_finite\_dilute.tex} and \texttt{paper/notes/README.md}; the brief for the open problems P1--P6: \texttt{paper/BRIEF\_fix.md}.

\subsection{Colour conventions of the checked copies}
\begin{itemize}
\item {\color{red}Red} \texttt{\{\textbackslash color\{red\}\dots\}}: corrected or added material, part of the expression.
\item Red text in square brackets \texttt{[\dots]}: a note of the check. It states the reason for the changes next to it and, usually, the numerical evidence. In this record such notes are reproduced in full in a box with a red bar.
\item {\color{orange}Orange} (in most files \texttt{\{\textbackslash color\{orange\}\textbackslash cancel\{\dots\}\}}, in \texttt{Divergent\_terms} and \texttt{scales} {\color{teal}teal} \texttt{\textbackslash cancel}): removed material. An orange note in square brackets (\texttt{[Removed: \dots]}, \texttt{[Removed here: \dots]}) records a removal and its reason.
\item {\color{magenta}Magenta} \fbox{FLAG $n$} or \fbox{OPEN}: an unresolved point (only in \texttt{Limiting\_cases\_checked} and in the analysis document \texttt{Finite\_terms\_final}).
\item {\color{dkgreen}Dark green} (only in \texttt{Limiting\_cases\_final}): the two rows that cancel each other (Group I Row IV, Group III Row I).
\item The author's own colours (blue, green, the author's red and orange highlighting in \texttt{Evolution\_three\_rho}, \texttt{BFKL}, \texttt{Divergent\_terms}, \texttt{scales}, \dots) were left unchanged. The extraction below separates them from the marks of the check (Section~\ref{sec:method}).
\end{itemize}

\subsection{How the list of changes was produced}\label{sec:method}
\noindent The per-file record of Section~\ref{sec:files} and the CSV file were produced by scripts (in the work directory \texttt{/tmp/claude-0/-home-user-correlations/0ab97fcf-94a2-5aec-8bad-586400481685/scratchpad/changes/}: \texttt{extract.py}, \texttt{extract2.py}, \texttt{render.py}), not by hand:
\begin{enumerate}
\item Every coloured group \texttt{\{\textbackslash color\{X\}\dots\}} with X $\in$ \{red, orange, teal, magenta, dark green\} of a checked copy is located (line, enclosing equation environment with its labels, section and subsection).
\item A group is attributed to the check only if it is \emph{not} in the original: the text of the group (with nested marks of the check undone) is searched in the original within three lines of the aligned line. The author's own red and orange groups are thereby excluded (e.g.\ 260 of them in \texttt{Evolution\_three\_rho}). The number of groups attributed to the check agrees, file by file, with the difference of the colour counts of the two files.
\item The original is reconstructed from the checked copy (red and magenta groups dropped, orange and teal groups and their \texttt{\textbackslash cancel} unwrapped) and compared token by token with the author's file. For each change the ``old'' text is taken \emph{from the author's file} (the text between the nearest unchanged neighbours), the ``new'' text is the checked copy with the marks accepted. Edits that the reconstruction does not explain are listed as \emph{unmarked edits}; almost all are layout (line breaks \texttt{\textbackslash notag\textbackslash\textbackslash\ \&} inserted in long lines, \texttt{\textbackslash par\textbackslash noindent} around inserted notes); the others are listed with their block.
\item Adjacent groups are merged into one change; each change is classified by rules on the old and new text: \emph{sign/factor correction}, \emph{phase convention}, \emph{index/Wilson-line orientation} (also argument and label relabellings), \emph{operator ordering} (also a term moved to another position), \emph{limits/range}, \emph{missing term added}, \emph{term removed/replaced}, \emph{typo/LaTeX} (also lists of integration variables, brackets, commas), \emph{note only}. The classification is automatic and is meant as a guide; for the changes of the summary table (Section~\ref{sec:summary}) it was checked by hand against the notes.
\item Formulas are rendered with a short piece of the surrounding expression (the original context on the left of the arrow, the checked one on the right); every rendered piece was compiled separately, and a piece that does not compile on its own (10 pieces, all of them misplaced-brace fixes, which are unbalanced by nature) is shown as source text.
\end{enumerate}

\subsection{Global decisions agreed with the author}\label{sec:decisions}
\noindent The following conventions and decisions apply to all files; each is stated in the checked copies at the place given.
\begin{enumerate}
\item \textbf{Phase} $e^{-ik\cdot(w-w')}$, with the conjugated bracket (primed coordinates, standing on the left) at $w'$ and the amplitude at $w$; correspondingly $e^{-ik\cdot(x-w')}$, $e^{-ik\cdot(y-y')}$, $e^{-ik\cdot(y-y')/\xi}$, and in Group III Rows III, V, VI $e^{-ik\cdot(w-y')}$ after the $w'$ integration. Reason: the definition of the cross section with the conjugated amplitude at $w'$ gives this sign; the original sign gives the result at $-\mathbf{k}$, and for non-commuting $\rho$ and a fixed target $k\to-k$ is not a symmetry (notes: Calculation\_main l.1537; Group I l.43, l.128; Group II l.54; Group II Row IV l.36). For the evolution lists the symmetrised two-$\rho$ parts do not depend on the sign, the three-$\rho$ and one-$\rho$ parts change sign together (BFKL l.39, checked with SU(2), SU(3) matrices).
\item \textbf{Plus prescription}: $\int dp^{+}\,[h(p^{+})]_{+}\,f(p^{+})\equiv\int dp^{+}\,h(p^{+})\,[f(p^{+})-f(0)]$, $f$ = everything that depends on $p^{+}$ (phase and denominators included) (Other\_terms\_v2 l.35; Limiting\_cases\_final l.36; Group III l.76). For DGLAP: $\int_{0}^{1}d\xi\,\big[\frac{1-\xi}{\xi}\big]_{+}F\equiv\int_{0}^{1}d\xi\,\frac{1-\xi}{\xi}\,[F(\xi)-F(0)]$ and $\int_{0}^{1}d\xi\,\frac{\xi}{(1-\xi)_{+}}F\equiv\int_{0}^{1}d\xi\,\frac{\xi F(\xi)-F(1)}{1-\xi}$ (Divergent\_terms l.86; Final\_terms l.229).
\item \textbf{No $1/\epsilon$ poles are written} in the final result (rule $\int d^{2}p/p^{2}=0$ for scaleless integrals; collinear poles absorbed); one $\overline{\rm MS}$ scale $\mu$, $\bar\mu_{\overline{MS}}=\mu_{\overline{MS}}$ (Divergent\_terms l.112, ``Decisions recorded''; Other\_terms l.35). The poles are kept in the intermediate files where they cannot be dropped by the scaleless rule (Group II Row IV l.143).
\item \textbf{Limits} $\Lambda\to0$ and $\vee\to\infty$ are taken in the finite terms (Limiting\_cases\_final l.36; \texttt{Finite\_terms\_final}, introduction).
\item \textbf{$\log(1-k^{+}/\vee)$ terms are dropped} (they vanish for $\vee\to\infty$) (Limiting\_cases\_final l.36, decision (5)). Caveat recorded: each is $\log(1-k^{+}/\vee)$ times the JIMWLK structure of its row, whose $z$ integral diverges at large $z$ for a non-neutral projectile; they vanish only with the large-distance regulator in place, and are best written with the JIMWLK term (Limiting\_cases\_checked l.35; \texttt{Cancellation\_of\_flagged\_terms}, Sec.~2.1).
\item \textbf{DGLAP} is written with the subtraction inside the splitting function, and the term $-\frac{1}{k^{+}}\times$(BFKL structure) of Group II Row IV is the left-over of the BFKL term and is counted among the finite terms (Limiting\_cases\_final l.36 decision (2) and l.153; Final\_terms l.229). The two readings of $\big[\frac{1-\xi}{\xi}\big]_{+}$ differ by exactly this term, $\int_{0}^{1}d\xi\,[\frac1\xi-\frac{1-\xi}{\xi}]=1$ (Limiting\_cases\_checked FLAG 5, l.183).
\item \textbf{The running-coupling term is the $\frac{\beta_{0}}{2}\delta(1-\xi)$ part of DGLAP} (added and subtracted): $N_{\rm DGLAP}+N_{\rm RC}$ does not depend on $\beta_{0}$; $g^{2}\to g^{2}(\mu)$ holds only together with the $\delta(1-\xi)$ term and the same $\mu$ (Divergent\_terms l.86 (v), l.112; Final\_terms l.229, l.239; Group II Row IV l.165).
\item \textbf{Group III Rows III, V, VI}: the large logarithm $\log\frac{\vee}{k^{+}}$ is removed completely, i.e.\ with all three terms of each row in the BFKL list, in the author's split $\big[\frac{k^{+}}{p^{+}(p^{+}+k^{+})}\big]_{+}+\frac{1}{p^{+}+k^{+}}$ with the subtracted $y'\to z$ terms (BFKL l.464, 493, 525, 551; Other\_terms l.212--236; Other\_terms\_v2 l.214--238; Limiting\_cases\_final l.36 decision (1)).
\item \textbf{Group I Row III} is written as one combined kernel $\mathcal{K}^{ij}$ (the separate equations $N_{1}$, $N_{2}$, $N_{3a}$, $N_{3b}$ are not integrable one by one), and the last term of $N_{3b}$ is replaced by the recoil term with $\mathcal{E}(s)=\int_{0}^{1}\frac{dt}{t}(e^{-ist}-1)$ (Limiting\_cases\_final l.36 decision (3), l.59, l.78; Limiting\_cases\_final\_1 l.107--108; Final\_terms l.312--313).
\item \textbf{Group I Row IV and Group III Row I cancel identically} and are not listed in the final terms (Limiting\_cases\_final l.36 decision (4), l.160--172; Final\_terms l.316, l.384).
\item \textbf{Operator ordering}: the three-$\rho$ part of $\mf{A}^{(3)}$ is $[\mf{A}^{(3)}_{2}]_{\rho\rho\rho}=\mf{A}^{(1)}_{2}\mf{N}_{2}$ ($\rho^{a}$ to the left of $\rho^{b}\rho^{b}$) (Calculation\_main l.555; Group I l.35). Consequence: the commutator term of Row I cancels, and the term $-2k^{i}\log\frac{\vee}{\Lambda}$ of Group I Row II is absent for the unbarred colour structure (Group I l.43, l.53).
\item \textbf{The two-$\rho$ section of the four-$\rho$ file is not used}: it was not confirmed (four-$\rho$ l.130); the two-$\rho$ remainder of the four-$\rho$ terms was obtained by an independent reordering of the corrected four-$\rho$ equation (BFKL l.101; JIMWLK l.284).
\end{enumerate}
\noindent Further conventions used throughout: $\rho$'s are non-commuting operators in the order written, $[\rho^{a}(x),\rho^{b}(y)]=if^{abc}\rho^{c}(x)\delta^{2}(x-y)$; c.c.\ implied for Groups I and III, not for Group II; energy denominators that vanish inside the integration region are principal values, so logarithms are $\log|\dots|$ (Calculation\_main l.62; Group I l.82); pure glue, $\beta_{0}=\frac{11N_{c}}{3}$.

\subsection{Chronology}
\noindent The checks were done in this order (file dates of the checked copies): Calculation\_main (wave functions, $\Omega$, cross-section assembly; 18 Nov); Group I, the four-$\rho$ file, Group II, Group II Row IV, Group III (18 Nov); the evolution lists (two- and three-$\rho$, JIMWLK, BFKL) and Other terms v1 (18 Nov); Divergent terms (18 Nov); the analysis of the flagged terms (\texttt{Cancellation\_of\_flagged\_terms}), scales, \texttt{Finite\_terms\_final} (18 Nov); Other terms v2, \texttt{Other\_terms\_limits} (19 Nov); Limiting cases (flags, then the author's decisions implemented), Limiting cases final 1 (19 Nov); Final terms (20 Nov). One open point of the evolution files (the form of the $\log\frac{k^{+}}{\Lambda}$ term of Group I Row III, $U^{bc}(z)$ versus $U^{bc}(w)$) was resolved later by the recoil identity (Section~\ref{sec:calc-recoil}); the notes of the earlier files still call it open.

\subsection{Reading the per-file record}
\noindent In Section~\ref{sec:files} each file starts with a header (paths, line counts, number of coloured groups of the check and of the author, number of records by type, the lines with unmarked layout edits). The changes follow in blocks; a block is one equation environment (with the notes that follow it) or one paragraph. Each change is shown as
\begin{center}
\textbf{l.$n$} \textsf{\small[type]} \emph{original (with a short context)} $\longrightarrow$ \emph{checked copy (same context)}\,,
\end{center}
or as \emph{added:} / \emph{removed:} with the context in grey. Long changes are shown as two paragraphs ``original'' / ``checked''. The notes of the check that follow a block are reproduced in full; they give the reason and the evidence for the changes of the block. Where equations were replaced as a whole (\texttt{Limiting\_cases\_final}), the author's equations and the new ones are printed in full before the blocks.
\clearpage
\section{Summary of the physics-relevant changes}\label{sec:summary}
\noindent In logical order (wave functions, cross-section rows, evolution lists, left-over terms, limits, final terms). ``l.'' = line of the checked copy (= line of the original, except in Limiting\_cases\_final and \_final\_1). Typo-level changes, repeated phase changes and pure notes are not listed here; they are all in Section~\ref{sec:files}. The evidence is quoted from the notes of the checked copies; ``Sec.~3.x'' points to the summary of the calculation.

{\small
\setlength{\LTleft}{0pt}\setlength{\LTright}{0pt}
\begin{longtable}{@{}p{0.4cm}p{2.7cm}p{5.0cm}p{4.0cm}p{4.0cm}@{}}
\toprule
\# & file, location & change & why & evidence\\
\midrule\endhead
\bottomrule\endfoot
1 & Calculation\_main l.420--432, 555--561, 1240, 1252, 1256 & $\rho^{a}(-\mathbf{k})$ moved to the left of $\rho^{b}(\mathbf{p})\rho^{b}(-\mathbf{p})$: $[\mf{A}^{(3)}_{2}]_{\rho\rho\rho}=\mf{A}^{(1)}_{2}\mf{N}_{2}$ & third-order perturbation theory with operator-valued matrix elements & explicit matrices vs exact $\Omega$ and exact diagonalisation: $3\times10^{-16}$; other ordering off at ${\cal O}(1)$ (Sec.~\ref{sec:calc-order})\\
2 & Calculation\_main l.1092--1101, 1146 & $\mf{N}^{(2)}\mf{A}^{(1)}=\mf{A}^{(1)}\mf{N}^{(2)}$ removed; $-[\mf{N}^{(2)},\mf{A}^{(1)}_{2}]$ kept & $[\rho^{a}(x)\rho^{a}(y),\rho^{b}(w)]\neq0$ & algebra (Sec.~\ref{sec:calc-comm})\\
3 & Calculation\_main l.1435, 1451, 1531 (CS1) & commutator $[\bar{\mf{N}}_{2},\bar{\mf{A}}^{(1)}]$ added & term of $\bar{\mf{A}}^{(3)\dagger}_{1}$ & with item 1 the $\mf{N}$-terms sum to $-\bar{\mf A}\mf N_2+U\bar{\mf N}_2\mf A+U\mf A\mf N_2-\bar{\mf N}_2\bar{\mf A}$, zero at $U=1$\\
4 & Calculation\_main l.1252 (A1coeff) & one-$\rho$ part: $-2\log\frac{\vee}{\Lambda}\to\log(\frac{\vee}{k^{+}}-1)-\log\frac{\vee}{\Lambda}$; constant $\frac{\pi^{2}}{3}-\frac{67}{9}+\log^{2}\frac{\Lambda}{k^{+}}-\log^{2}\frac{\vee}{k^{+}}$ removed; $\xi$ limits $\frac{k^{+}}{\vee}<\xi<\frac{k^{+}}{k^{+}+\Lambda}$ & relabelling $k^{+}\to\xi k^{+}$; constant was counted twice & ``checked numerically against Eqs.\ (combinedgrho), (combinedgrhorho) and the coordinate form'' (l.1253)\\
5 & Calculation\_main l.1260, 1268 & support of $\mf{B}_{2}$: $\Theta(k^{+}-\Lambda-p^{+})\to\Theta(\vee-k^{+}-p^{+})$ (single $\rho$), $\Theta(\vee-k^{+})\Theta(\vee-p^{+})$ ($\{\rho,\rho\}$) & parent gluon is a soft mode; two independent emissions & four-$\rho$ cancellation for commuting $\rho$ needs a common $\log\frac{\vee}{\Lambda}$ (Sec.~\ref{sec:calc-B2})\\
6 & Calculation\_main l.1276, 1289, 1302 (also 684, 706) & sign of the non-instantaneous $p^{+}>k^{+}$ term of $\mf{B}_{3}$: $(k^{+}-p^{+})\to(p^{+}-k^{+})$ & from (m2), (m3), cancellation of $1/(1-\xi)^{2}$, $a^{\dagger}a$ unitarity & $1/\Lambda$ terms of Group I Row III cancel (Sec.~\ref{sec:calc-B3}, \ref{sec:calc-scales})\\
7 & Calculation\_main l.1272, 1070 & labels of $\mf{B}_{2}$ ($ba\to ab$, $ji\to ij$, $z\leftrightarrow u$); arguments $p\leftrightarrow q$ in the $\Omega$ expansion & convention of Eq.~(Omegaonvacposzeta); normal ordering & notes l.1273, l.1076\\
8 & Calculation\_main l.1469--1490, 1531--1541 & $U^{bc}(v)\to U^{\dagger bc}(v)$, $U^{b'a}(w')\to U^{\dagger b'a}(w')$, transposed indices in CS1, CS2 & $\bar a^{\dagger f}=U^{fg}a^{\dagger g}$ & all sixteen products of CS2 present with the right weights (l.1542)\\
9 & Calculation\_main l.1530, 1539 and all cross-section and evolution files & phase $e^{-ik(w'-w)}\to e^{-ik(w-w')}$ (and $(w'-z)\to(z-w')$ etc.\ after integrations; $(y'-y)\to(y-y')$) & conjugated bracket at $w'$; $k\to-k$ not a symmetry for non-commuting $\rho$ & symmetrised two-$\rho$ parts unchanged, three- and one-$\rho$ parts change sign (SU(2), SU(3); BFKL l.39)\\
10 & Calculation\_main l.984--989 & ``$=0$'' of the normalisation lines removed; order $\rho^{a}(\mathbf{k})\rho^{a'}(-\mathbf{k}')$; primes of $a,i,k$ & only the sum of six terms vanishes; $\delta$ term zero only by $\int d^{2}p/p^{2}=0$ & notes l.985, 986\\
11 & Group I l.43--74 (Row II) & red terms in $N_{2}$, $N_{2a}$, $N_{2b}$: $-2k^{i}\log\frac{\vee}{\Lambda}$ removed for the unbarred structure; $R^{i}\to R^{i}+2k^{i}\log\frac{\vee}{k^{+}}$ & consequence of item 1: $-U[\mf N_2,\mf A^{(1)}]$ equals that term & same coefficient and sign, SU(2), SU(3) matrices; Other\_terms\_v2 l.49: residual $<10^{-14}$\\
12 & Group I l.80--84 (Row III) & window $|p^{+}-k^{+}|<\Lambda$ excluded in the instantaneous term; two logarithms added to $N_{1}$; logs are $\log|\cdot|$ & as in the wave functions; principal value & power terms: sum of eight factors identically zero (scales l.81)\\
13 & Group I l.139--262 (Row V) & every exponent reversed; range $\Lambda<p^{+}<k^{+}-\Lambda$; note on $\frac{2}{\epsilon}$ (loop-momentum continuation) & phase (item 9) & coefficient $2\log\frac{k^{+}}{\Lambda}-\frac{11}{6}$, ``UV''+``rest'' $=2b$ checked numerically (l.264)\\
14 & Group II Row IV l.93--180 & lower limit $0\to\frac{k^{+}}{\vee}$; $+\log(1-\frac{k^{+}}{\vee})$ added to $L$; $\frac{2}{\epsilon}$ kept (sentence ``Dropping all the $\frac1\epsilon$'' removed); $\beta_{0}$ added and subtracted & exact $\xi$ integrals; poles not scaleless (depend on $k$, target) & ``checked, also numerically with a test function'' (l.175)\\
15 & Group III l.51--52, 62--63 & Row II: garbled argument repaired, $p^{+}$ limits added; Row III: prefactor $-\frac{4}{(2\pi)^{3}}$ (minus sign added) & (CS2): $\mf{C}_{1}\mf{A}$ part has the opposite sign & ``checked numerically from the coefficient functions'' (Sec.~\ref{sec:calc-G3R3})\\
16 & four-$\rho$ l.67, l.130 & eight transposed index pairs of $U(z)$, $U(w')$; the two-$\rho$ section ``Not confirmed'' & written correctly in the third block & compared term by term with CS1, CS2 (SU(2), SU(3)); two-$\rho$ section does not vanish at $U=1$\\
17 & Evolution three-$\rho$ l.57, 67; 102--175; 291--313; 244 & Group I Row III: signs of the first two terms, $U^{bc}(z)\to U^{bc}(w)$; transposed $U$ indices (Group II Row I, Group III rows); kernels of the $\beta$, $\beta'$ headers $\mathcal{K}(x'-z)\to\mathcal{K}(w-z),\mathcal{K}(w'-z)$; a sign & $N_{2}$ of Group I has an overall minus; only then the two forms agree & lattice SU(2) check: symmetrised three-$\rho$ part vanishes identically (l.35, 253)\\
18 & JIMWLK l.668--706, 763; BFKL l.41--44, 579--594 & order $\rho(x)\rho(x')\to\rho(x')\rho(x)$ in the marked terms of the final lists; four missing brackets (BFKL) & equality with the grouped forms also in the contact terms & final JIMWLK $=-H_{\rm JIMWLK}$ on LO; BFKL $=$ one projectile-evolution step (JIMWLK l.735, BFKL l.633)\\
19 & Other terms (v1) l.45, 111, 143, 190--192 & added: Row II bracket $[\dots][\frac{11}{3}+\ell]$; Row III last term of $N_{3b}$ with $\Lambda$ kept; Row V $N_{2b\,\rm UV}$; Group II Row IV lower limit $\frac{k^{+}}{\vee}$ and four left-over terms ($-(1-\frac{k^{+}}{\vee})$, $\log(1-\frac{k^{+}}{\vee})$, DGLAP, RC) & what remains of the checked rows when exactly $\log\frac{\vee}{k^{+}}$ and $\log\frac{k^{+}}{\Lambda}$ are taken out & row-by-row comparison (l.35)\\
20 & Other terms (v2) l.42--217 & closing braces of the phases (factors $\frac{p^{+}}{k^{+}}$, $\frac{q^{+}}{k^{+}}$ back in the exponent); commas & misprints & residuals $\sim10^{-11},10^{-13},10^{-15}$ for $\frac{\Lambda}{k^{+}}=10^{-3},10^{-5},10^{-7}$ (l.35--238)\\
21 & Divergent terms l.90 and notes l.42--112 & $\log(1-\frac{k^{+}}{\vee})$ pole term added; pole bookkeeping; reading of the DGLAP plus prescriptions & completeness & pole sum $-\frac{\ell}{2}$ (Sec.~\ref{sec:calc-poles})\\
22 & scales l.80 & $\int_{w,w'}$ added & missing measure & power terms cancel identically (Sec.~\ref{sec:calc-scales})\\
23 & Limiting cases (flags) l.181 and FLAG 1--9 & $\xi$ lower limit $\frac{k^{+}}{\vee}\to0$; nine flags (not integrable as written / limit not taken) & --- & numbers in the flags (Sec.~\ref{sec:calc-recoil}, \ref{sec:calc-G3})\\
24 & Limiting cases final l.59--78 & Group I Row III: four equations replaced by the combined kernel $\mathcal{K}^{ij}$ + third term of $N_{3b}$ + recoil term with $\mathcal{E}$ & not integrable one by one; last term divergent at $z\to w$ & sum vs separate terms $2\times10^{-22}$--$7\times10^{-18}$ (relative $5\times10^{-16}$--$1.5\times10^{-12}$); recoil identity $3\times10^{-13}$\\
25 & Limiting cases final l.181--212 & Group III Rows III, V, VI replaced by the forms with all three BFKL terms removed, upper limit $\infty$ & the author's form still grows like $\log\frac{\vee}{k^{+}}$ & converges like $\frac{k^{+}}{\vee}$ (Sec.~\ref{sec:calc-G3})\\
26 & Limiting cases final l.81--85, 157--172; Final terms l.316, 384 & Group I Row IV and Group III Row I shown cancelling; not listed in the final terms & identical colour factors after renaming, opposite prefactors & sum $\sim10^{-23}$ vs rows $\sim10^{-8}$ (Sec.~\ref{sec:calc-cancelIV})\\
27 & Limiting cases final l.151 & DGLAP: subtraction in the splitting function; $-\frac{1}{k^{+}}$ remnant finite & author's decision & two readings differ by exactly the remnant (Sec.~\ref{sec:calc-dglap})\\
28 & Limiting cases final 1 l.107--108; Final terms l.312--313 & last term of $N_{3b}$ replaced by $-\frac{ig^{4}f^{abd}}{8\pi^{5}k^{+}}\int\mathcal{E}(k\cdot(z-w))K(x,w,z)K(x'-w')\cdot K(y-w)\,[\dots U^{bc}(w)\dots]$ & $\log^{2}$ divergence at $z\to w$ & recoil identity (Sec.~\ref{sec:calc-recoil})\\
29 & Final terms l.224 & DGLAP phase: $\frac{1}{\xi}$ moved into the exponent, $e^{-ik\cdot(y-y')\frac1\xi}$ & as in the checked Group II Row IV & character-by-character comparison with the verified files (l.35)\\
30 & fix notes (Sec.~\ref{sec:open}) & Group I Rows III, V and Group III Rows III, V, VI in integrable form (soft-recoil block $\mathcal{S}$); DGLAP in $\overline{\rm MS}$ form; singlet form & P1--P6 & see Sec.~\ref{sec:fix-soft}--\ref{sec:fix-singlet}\\
\end{longtable}
}

\clearpage
\section{Changes, file by file}\label{sec:files}
\noindent One subsection per file, in the order of Table~\ref{tab:pairs}. See Section~\ref{sec:method} for how the entries were obtained and Section~\ref{sec:howto} for the colours.
\needspace{15\baselineskip}\bigskip
\subsection{Calculation\_main}\label{sec:file-calc_main}
\noindent\textbf{Original:} \texttt{/root/.claude/uploads/0ab97fcf-94a2-5aec-8bad-586400481685/99a3a399-Calculation\_main\_1.tex}\\\textbf{Checked copy:} \texttt{/home/user/correlations/nlo\_review/calc\_main\_checked/Calculation\_main\_1\_checked.tex}\\
Lines: 1943 (original), 1943 (checked copy). Coloured groups of the check: 194 (194 top-level); author\textquotesingle s own coloured groups left unchanged: 0.\\
Records by type: sign/factor correction 9; phase convention 3; index/Wilson-line orientation 41; operator ordering 17; limits/range 8; missing term added 11; term removed/replaced 12; typo/LaTeX 14; note only 21.
\par\noindent Unmarked layout-only edits (line breaks \texttt{\textbackslash notag\textbackslash\textbackslash\ \&}, \texttt{\textbackslash par\textbackslash noindent} around inserted notes, spacing), checked-copy lines: 62, 421, 557, 561, 1076, 1100, 1101, 1146, 1147, 1243, 1248, 1249, 1252, 1253, 1256, 1260, 1261, 1269, 1273, 1277, 1290, 1303, 1358, 1430, 1435, 1448, 1491, 1531, 1537, 1540, 1542.
\par\noindent Unmarked edits that are not layout only (listed with their block below): 2.
\subsubsection*{Changes, block by block}
\needspace{3\baselineskip}\noindent\rule{\linewidth}{0.3pt}\par\noindent\textbf{Lines 62}\par
\begin{itemize}[leftmargin=2.2em,itemsep=1pt,topsep=2pt]
\item[\textbf{l.62}] \textsf{\small [note only]} \emph{added:} \textbf{Marking in this checked copy:} changes and additions are in red;  Line numbers of the source are those of the original file. Energy denominators that vanish inside the integration region are principal values throughout. $\vee=\Lambda e^{\delta\mathsf{Y}}$ in the notation of arXiv:1610.03453; $X=x-z$, $X^{\prime}=x-z^{\prime}$, $Y=y-z$.
\end{itemize}
\needspace{3\baselineskip}\noindent\rule{\linewidth}{0.3pt}\par\noindent\textbf{Lines 102} \quad \textsf{Review of the vacuum incoming state $\rangle$ Two outgoing gluons} \quad {\small (in \texttt{align}, lines 101--103)}\par
\begin{itemize}[leftmargin=2.2em,itemsep=1pt,topsep=2pt]
\item[\textbf{l.102}] \textsf{\small [typo/LaTeX]} $\displaystyle dp^{+}\, \frac{g^{2}p^{i}(k^{j}-p^{j})\,\{\rho^{a(-p),\rho^{b}(-k+p)}\}}{16\pi^{3}\,p^{2}(k-p)^{2}\, \sqrt{p^{+}(k^{+}-p^{+})}}\, \ket{g^{b}_{j}(k^{+}-p^{+}, k-p),g^{a}_{i}(p^{+},p)}$\ $\longrightarrow$\ $\displaystyle dp^{+}\, \frac{g^{2}p^{i}(k^{j}-p^{j})\,\{\rho^{a}(-p),\rho^{b}(-k+p)\}}{16\pi^{3}\,p^{2}(k-p)^{2}\, \sqrt{p^{+}(k^{+}-p^{+})}}\, \ket{g^{b}_{j}(k^{+}-p^{+}, k-p),g^{a}_{i}(p^{+},p)}$
\end{itemize}
\needspace{3\baselineskip}\noindent\rule{\linewidth}{0.3pt}\par\noindent\textbf{Lines 186} \quad \textsf{Review of the vacuum incoming state $\rangle$ One outgoing gluon} \quad {\small (in \texttt{align}, lines 185--187)}\par
\begin{itemize}[leftmargin=2.2em,itemsep=1pt,topsep=2pt]
\item[\textbf{l.186}] \textsf{\small [term removed/replaced]} \emph{removed:} $\displaystyle \xi(1-\xi)$\quad{\color{gray}\small(original context: }$\displaystyle 8\pi^{3/2}\,(2\pi)^{3}\, \xi(1-\xi)\,k^{4}(k^{2}\xi\,(1-\xi)+\tilde{p}${\color{gray}\small)}
\end{itemize}
\needspace{3\baselineskip}\noindent\rule{\linewidth}{0.3pt}\par\noindent\textbf{Lines 410}\par
\begin{itemize}[leftmargin=2.2em,itemsep=1pt,topsep=2pt]
\item[\textbf{l.410}] \textsf{\small [note only]} or $\Lambda$ is very small\ $\longrightarrow$\ ($\vee\gg\Lambda$; small $\Lambda$ alone does not justify dropping $1/\vee$, and the approximation is not needed because the exact forms cancel in $\left|\Psi_{g\rho\rho\rho}\right\rangle$)
\end{itemize}
\needspace{3\baselineskip}\noindent\rule{\linewidth}{0.3pt}\par\noindent\textbf{Lines 420--421} \quad {\small (in \texttt{align}, lines 419--421, label \texttt{degen})}\par
\begin{itemize}[leftmargin=2.2em,itemsep=1pt,topsep=2pt]
\item[\textbf{l.420}] \textsf{\small [operator ordering]} \emph{added:} $\displaystyle \langle g_{i}^{a}(k)|H_{g}|0\rangle$\quad{\color{gray}\small(in context: }$\displaystyle \langle g_{i}^{a}(k)|H_{g}|0\rangle|\langle g_{j}^{b}(p)|H_{g}|${\color{gray}\small)}
\item[\textbf{l.420}] \textsf{\small [operator ordering]} \emph{removed:} $\displaystyle \langle g_{i}^{a}(k)|H_{g}|0\rangle$\quad{\color{gray}\small(original context: }$\displaystyle (p)|H_{g}|0\rangle |^{2}\langle g_{i}^{a}(k)|H_{g}|0\rangle (2E_{g}(p)+E_{g}(k))${\color{gray}\small)}
\end{itemize}
\begin{notebox}\textsf{\small Note, l.421:} [Ordering: in third-order perturbation theory with operator-valued matrix elements the single-emission element $\left\langle g_{i}^{a}(k)\left|H_{g}\right|0\right\rangle\propto\rho^a(-\mathbf{k})$ stands to the left of the norm factor $\propto\rho^b(\mathbf{p})\rho^b(-\mathbf{p})$; see the note before Eq.~(\texttt{3rhosmom}).]\end{notebox}
\needspace{3\baselineskip}\noindent\rule{\linewidth}{0.3pt}\par\noindent\textbf{Lines 424} \quad {\small (in \texttt{align}, lines 423--425, label \texttt{virtual})}\par
\begin{itemize}[leftmargin=2.2em,itemsep=1pt,topsep=2pt]
\item[\textbf{l.424}] \textsf{\small [operator ordering]} \emph{added:} $\displaystyle \rho^{a}(-\mathbf{k})$\quad{\color{gray}\small(in context: }$\displaystyle g^{3}\mathbf{k}^{i}\rho^{a}(-\mathbf{k})\rho^{b}(\mathbf{p})\rho^{b}(-\mathbf{p})${\color{gray}\small)}
\item[\textbf{l.424}] \textsf{\small [operator ordering]} \emph{removed:} $\displaystyle \rho^{a}(-\mathbf{k})$\quad{\color{gray}\small(original context: }$\displaystyle \rho^{b}(\mathbf{p})\rho^{b}(-\mathbf{p})\rho^{a}(-\mathbf{k})${\color{gray}\small)}
\end{itemize}
\needspace{3\baselineskip}\noindent\rule{\linewidth}{0.3pt}\par\noindent\textbf{Lines 428--430} \quad {\small (in \texttt{align}, lines 427--429, label \texttt{zetainteg})}\par
\begin{itemize}[leftmargin=2.2em,itemsep=1pt,topsep=2pt]
\item[\textbf{l.428}] \textsf{\small [operator ordering]} \emph{added:} $\displaystyle \rho^{a}(-\mathbf{k})$\quad{\color{gray}\small(in context: }$\displaystyle g^{3}\mathbf{k}^{i}\rho^{a}(-\mathbf{k})\rho^{b}(\mathbf{p})\rho^{b}(-\mathbf{p})${\color{gray}\small)}
\item[\textbf{l.428}] \textsf{\small [operator ordering]} \emph{removed:} $\displaystyle \rho^{a}(-\mathbf{k})$\quad{\color{gray}\small(original context: }$\displaystyle \rho^{b}(\mathbf{p})\rho^{b}(-\mathbf{p})\rho^{a}(-\mathbf{k})${\color{gray}\small)}
\item[\textbf{l.430}] \textsf{\small [typo/LaTeX]} or $\Lambda$ is very small\ $\longrightarrow$\ ($\vee\gg\Lambda$)
\end{itemize}
\needspace{3\baselineskip}\noindent\rule{\linewidth}{0.3pt}\par\noindent\textbf{Lines 432} \quad {\small (in \texttt{align}, lines 431--433)}\par
\begin{itemize}[leftmargin=2.2em,itemsep=1pt,topsep=2pt]
\item[\textbf{l.432}] \textsf{\small [operator ordering]} \emph{added:} $\displaystyle \rho^{a}(-\mathbf{k})$\quad{\color{gray}\small(in context: }$\displaystyle g^{3}\mathbf{k}^{i}\rho^{a}(-\mathbf{k})\rho^{b}(\mathbf{p})\rho^{b}(-\mathbf{p})${\color{gray}\small)}
\item[\textbf{l.432}] \textsf{\small [operator ordering]} \emph{removed:} $\displaystyle \rho^{a}(-\mathbf{k})$\quad{\color{gray}\small(original context: }$\displaystyle \rho^{b}(\mathbf{p})\rho^{b}(-\mathbf{p})\rho^{a}(-\mathbf{k})${\color{gray}\small)}
\end{itemize}
\needspace{3\baselineskip}\noindent\rule{\linewidth}{0.3pt}\par\noindent\textbf{Lines 555} \quad \textsf{Review of the vacuum incoming state $\rangle$ Final results}\par
\begin{itemize}[leftmargin=2.2em,itemsep=1pt,topsep=2pt]
\item[\textbf{l.555}] \textsf{\small [typo/LaTeX]} (\texttt{threerhoafter})\ $\longrightarrow$\ the exact form of $\left|\Psi_{g\rho\rho\rho}^{6\,c}\right\rangle$ (the equation just above (\texttt{threerhoafter}), with $\left[\frac{1}{\vee}-\frac{1}{\Lambda}\right]$)
\end{itemize}
\begin{notebox}\textsf{\small Note, l.555:} ; the $1/\Lambda$ and $1/\vee$ terms cancel exactly. Operator ordering: the surviving $\log(\vee/\Lambda)$ term comes from $\left|\Psi^{7}_{g\rho\rho\rho}\right\rangle$ and has $\rho^a(-\mathbf{k})$ to the \emph{left} of $\rho^b(\mathbf{p})\rho^b(-\mathbf{p})$, i.e. $[\mf{A}^{(3)}_2]_{\rho\rho\rho}=\mf{A}^{(1)}_2\,\mf{N}_2$ (in coordinate space $\rho^a(\mathbf{w})\rho^b(\mathbf{x})\rho^b(\mathbf{y})$). This follows from third-order perturbation theory with operator-valued matrix elements (wave operator with canonical normalisation) and was verified with explicit matrices against an exact construction of $\Omega$ and against exact diagonalisation (agreement to $3\times10^{-16}$; the ordering $\rho^b\rho^b\rho^a$ deviates at order one). It differs from the ordering written in the paper, Eqs.~(3.81)--(3.83) and (3.102)--(3.103). The power terms still cancel between $\left|\Psi^{6\,c}_{g\rho\rho\rho}\right\rangle$ (ordering $\rho^b\rho^b\rho^a$) and $\left|\Psi^{7}_{g\rho\rho\rho}\right\rangle$ (ordering $\rho^a\rho^b\rho^b$), because $\int d^{2}\mathbf{p}\,\rho^b(\mathbf{p})\rho^b(-\mathbf{p})$ (no $1/\mathbf{p}^2$ weight) is local in coordinate space and commutes with $\rho^a$. The two-$\rho$ part, Eq.~(\texttt{combinedgrhorho}), including its $-2k^i\log(\vee/\Lambda)$ term, is correct as it stands together with this ordering\end{notebox}
\needspace{3\baselineskip}\noindent\rule{\linewidth}{0.3pt}\par\noindent\textbf{Lines 557} \quad {\small (in \texttt{align}, lines 556--558, label \texttt{3rhosmom})}\par
\begin{itemize}[leftmargin=2.2em,itemsep=1pt,topsep=2pt]
\item[\textbf{l.557}] \textsf{\small [operator ordering]} \emph{added:} $\displaystyle \rho^{a}(-\mathbf{k})$\quad{\color{gray}\small(in context: }$\displaystyle g^{3}\,\rho^{a}(-\mathbf{k})\,\rho^{b}(\mathbf{p})\,\rho^{b}(-\mathbf{p})${\color{gray}\small)}
\item[\textbf{l.557}] \textsf{\small [operator ordering]} \emph{removed:} $\displaystyle \,\rho^{a}(-\mathbf{k})$\quad{\color{gray}\small(original context: }$\displaystyle \,\rho^{b}(\mathbf{p})\,\rho^{b}(-\mathbf{p})\,\rho^{a}(-\mathbf{k})\, \mathbf{k}^{i}${\color{gray}\small)}
\end{itemize}
\needspace{3\baselineskip}\noindent\rule{\linewidth}{0.3pt}\par\noindent\textbf{Lines 561} \quad {\small (in \texttt{align}, lines 560--562, label \texttt{3rhospos})}\par
\begin{itemize}[leftmargin=2.2em,itemsep=1pt,topsep=2pt]
\item[\textbf{l.561}] \textsf{\small [operator ordering]} \emph{added:} $\displaystyle \rho^{a}(\mathbf{w})$\quad{\color{gray}\small(in context: }$\displaystyle ig^{3}\,\rho^{a}(\mathbf{w})\,\rho^{b}(\mathbf{x})\,\rho^{b}(\mathbf{y})${\color{gray}\small)}
\item[\textbf{l.561}] \textsf{\small [operator ordering]} \emph{removed:} $\displaystyle \,\rho^{a}(\mathbf{w})$\quad{\color{gray}\small(original context: }$\displaystyle \,\rho^{b}(\mathbf{x})\,\rho^{b}(\mathbf{y})\,\rho^{a}(\mathbf{w})\, (x-u)\cdot (y-u)\,(w-z)^{i}${\color{gray}\small)}
\end{itemize}
\needspace{3\baselineskip}\noindent\rule{\linewidth}{0.3pt}\par\noindent\textbf{Lines 598} \quad \textsf{Calculation of the single gluon incoming state $\rangle$ No outgoing gluon}\par
\begin{itemize}[leftmargin=2.2em,itemsep=1pt,topsep=2pt]
\item[\textbf{l.598}] \textsf{\small [typo/LaTeX]} he\ $\longrightarrow$\ the
\end{itemize}
\needspace{3\baselineskip}\noindent\rule{\linewidth}{0.3pt}\par\noindent\textbf{Lines 684} \quad \textsf{Calculation of the single gluon incoming state $\rangle$ One outgoing gluon} \quad {\small (in \texttt{align}, lines 683--685, label \texttt{onegrho12})}\par
\begin{itemize}[leftmargin=2.2em,itemsep=1pt,topsep=2pt]
\item[\textbf{l.684}] \textsf{\small [sign/factor correction]} $\displaystyle -\frac{ig^{2}f^{abc}}{8\pi^{3}}\,\int_{k^{+}+\Lambda}^{\vee} dp^{+}\,\int$\ $\longrightarrow$\ $\displaystyle +\frac{ig^{2}f^{abc}}{8\pi^{3}}\,\int_{k^{+}+\Lambda}^{\vee} dp^{+}\,\int$
\end{itemize}
\needspace{3\baselineskip}\noindent\rule{\linewidth}{0.3pt}\par\noindent\textbf{Lines 706--720} \quad {\small (in \texttt{align}, lines 700--719, label \texttt{Eq:onegrhoX1})}\par
\begin{itemize}[leftmargin=2.2em,itemsep=1pt,topsep=2pt]
\item[\textbf{l.706}] \textsf{\small [sign/factor correction]} \emph{added:} $\displaystyle -$\quad{\color{gray}\small(in context: }$\displaystyle -\frac{ig^2 f^{abc}} {4 \pi^3}${\color{gray}\small)}
\item[\textbf{l.720}] \textsf{\small [typo/LaTeX]} two\ $\longrightarrow$\ to
\end{itemize}
\needspace{3\baselineskip}\noindent\rule{\linewidth}{0.3pt}\par\noindent\textbf{Lines 732} \quad {\small (in \texttt{align}, lines 731--733, label \texttt{onegrho22})}\par
\begin{itemize}[leftmargin=2.2em,itemsep=1pt,topsep=2pt]
\item[\textbf{l.732}] \textsf{\small [sign/factor correction]} \emph{removed:} $\displaystyle +$\quad{\color{gray}\small(original context: }$\displaystyle \frac{k^{+}+p^{+}}{(k^{+}-p^{+})}+(k-p)^{l}${\color{gray}\small)}
\end{itemize}
\needspace{3\baselineskip}\noindent\rule{\linewidth}{0.3pt}\par\noindent\textbf{Lines 913} \quad \textsf{Calculation of the single gluon incoming state $\rangle$ Two outgoing gluons} \quad {\small (in \texttt{align}, lines 912--917, label \texttt{gindeppos})}\par
\begin{itemize}[leftmargin=2.2em,itemsep=1pt,topsep=2pt]
\item[\textbf{l.913}] \textsf{\small [missing term added]} \emph{added:} $\displaystyle gg$\quad{\color{gray}\small(in context: }$\displaystyle ,z^{\prime}))_{gg}${\color{gray}\small)}
\end{itemize}
\noindent{\small\textsf{Unmarked edit, l.913 (inserted):} \texttt{\textbackslash{}\allowbreak{}left(}}\par
\noindent{\small\textsf{Unmarked edit, l.913 (inserted):} \texttt{\textbackslash{}\allowbreak{}right)\allowbreak{}\_\{\}\allowbreak{}}}\par
\needspace{3\baselineskip}\noindent\rule{\linewidth}{0.3pt}\par\noindent\textbf{Lines 953} \quad {\small (in \texttt{align}, lines 949--954)}\par
\begin{itemize}[leftmargin=2.2em,itemsep=1pt,topsep=2pt]
\item[\textbf{l.953}] \textsf{\small [term removed/replaced]} $\displaystyle \xi k^+,\v{z}$\ $\longrightarrow$\ $\displaystyle p^+,\v{z}$
\end{itemize}
\needspace{3\baselineskip}\noindent\rule{\linewidth}{0.3pt}\par\noindent\textbf{Lines 984--986} \quad \textsf{Calculation of the single gluon incoming state $\rangle$ Normalization}\par
\begin{itemize}[leftmargin=2.2em,itemsep=1pt,topsep=2pt]
\item[\textbf{l.984}] \textsf{\small [term removed/replaced]} \emph{removed:} $\displaystyle = 0$
\item[\textbf{l.985}] \textsf{\small [sign/factor correction]} \emph{added:} $\displaystyle +$\quad{\color{gray}\small(in context: }$\displaystyle $$+\bra{(\Psi^{a'}_{i'}(k'))_{gg\rho} }\ket{ (\Psi^{a}_{i}(k) )_{gg} }${\color{gray}\small)}
\item[\textbf{l.985}] \textsf{\small [term removed/replaced]} $\displaystyle $$\bra{(\Psi^{a'}_{i'}(k'))_{gg\rho} }\ket{ (\Psi^{a}_{i}(k) )_{gg} } + \bra{(\Psi^{a}_{i}(k))_{gg} }\ket{(\Psi^{a'}_{i'}(k'))_{gg\rho} }$\ $\longrightarrow$\ $\displaystyle $$+\bra{(\Psi^{a'}_{i'}(k'))_{gg\rho} }\ket{ (\Psi^{a}_{i}(k) )_{gg} } + \bra{(\Psi^{a'}_{i'}(k'))_{gg} }\ket{(\Psi^{a}_{i}(k))_{gg\rho} }$
\item[\textbf{l.985}] \textsf{\small [typo/LaTeX]} \emph{added:} (only the sum of the six terms of these two lines vanishes; the two lines do not vanish separately. $\mf{N}$ is not affected.)
\item[\textbf{l.986}] \textsf{\small [missing term added]} \emph{added:} $\displaystyle \delta^{aa'}\delta_{ii'}\delta^{3}(k-k')\,$\quad{\color{gray}\small(in context: }$\displaystyle -\delta^{aa'}\delta_{ii'}\delta^{3}(k-k')\,\frac{g^2N_c}{8\pi^3} ( 2 \ln${\color{gray}\small)}
\item[\textbf{l.986}] \textsf{\small [typo/LaTeX]} \emph{added:} (zero only by the scaleless-integral rule $\int d^{2}p/p^{2}=0$)
\end{itemize}
\needspace{3\baselineskip}\noindent\rule{\linewidth}{0.3pt}\par\noindent\textbf{Lines 989} \quad {\small (in \texttt{align}, lines 988--993)}\par
\begin{itemize}[leftmargin=2.2em,itemsep=1pt,topsep=2pt]
\item[\textbf{l.989}] \textsf{\small [index/Wilson-line orientation]} $\displaystyle \bra{g^{a}_{}()}\ket{(\Psi^{a}_{i}(k))^{2}_{g\rho\rho} }$\ $\longrightarrow$\ $\displaystyle \bra{g^{a'}_{i'}(k')}\ket{(\Psi^{a}_{i}(k))^{2}_{g\rho\rho} }$
\item[\textbf{l.989}] \textsf{\small [index/Wilson-line orientation]} $\displaystyle i$\ $\longrightarrow$\ $\displaystyle i'$
\item[\textbf{l.989}] \textsf{\small [index/Wilson-line orientation]} $\displaystyle k$\ $\longrightarrow$\ $\displaystyle k'$
\item[\textbf{l.989}] \textsf{\small [operator ordering]} $\displaystyle \frac{g^2}{4 \pi^3 \sqrt{k^+ k^{'+}}} \frac{k^{i} k^{'i'}}{\v{k}^2\v{k}^{'2'}} \rho^{a'} (-\v{k}') \rho^a (\v{k})$\ $\longrightarrow$\ $\displaystyle \frac{g^2}{4 \pi^3 \sqrt{k^+ k^{'+}}} \frac{k^{i} k^{'i'}}{\v{k}^2\v{k}^{'2'}} \rho^a (\v{k})\rho^{a'} (-\v{k}')$
\item[\textbf{l.989}] \textsf{\small [index/Wilson-line orientation]} $\displaystyle \Psi^{a}_{}($\ $\longrightarrow$\ $\displaystyle \Psi^{a'}_{i'}($
\item[\textbf{l.989}] \textsf{\small [index/Wilson-line orientation]} $\displaystyle i$\ $\longrightarrow$\ $\displaystyle i'$
\item[\textbf{l.989}] \textsf{\small [index/Wilson-line orientation]} $\displaystyle k$\ $\longrightarrow$\ $\displaystyle k'$
\end{itemize}
\needspace{3\baselineskip}\noindent\rule{\linewidth}{0.3pt}\par\noindent\textbf{Lines 1070--1076} \quad \textsf{Parametrization of the Evolution operator $(\Omega)$} \quad {\small (in \texttt{align}, lines 1068--1076)}\par
\begin{itemize}[leftmargin=2.2em,itemsep=1pt,topsep=2pt]
\item[\textbf{l.1070}] \textsf{\small [index/Wilson-line orientation]} $\displaystyle \mf{A}^{\dagger c}_{2k}(p)\mf{A}^{b}_{2j}$\ $\longrightarrow$\ $\displaystyle \mf{A}^{\dagger c}_{2k}(q)\mf{A}^{b}_{2j}$
\item[\textbf{l.1070}] \textsf{\small [index/Wilson-line orientation]} $\displaystyle \mf{A}^{b}_{2j}(q)$\ $\longrightarrow$\ $\displaystyle \mf{A}^{b}_{2j}(p)$
\end{itemize}
\begin{notebox}\textsf{\small Note, l.1076:} [Arguments fixed: $(\mf{A}_2^\dagger a)(\mf{A}_2 a^\dagger)$ normal-orders to $\mf{A}^{\dagger c}_{2k}(q)\mf{A}^{b}_{2j}(p)\,a^{\dagger b}_j(p)a^c_k(q)$; the $\rho$-independent term $2\int\mf{C}_1^\dagger\mf{C}_1$ of the $a^\dagger a$ bracket is not written.]\end{notebox}
\needspace{3\baselineskip}\noindent\rule{\linewidth}{0.3pt}\par\noindent\textbf{Lines 1092}\par
\begin{itemize}[leftmargin=2.2em,itemsep=1pt,topsep=2pt]
\item[\textbf{l.1092}] \textsf{\small [typo/LaTeX]} \emph{removed:} The fact that $A^{(1)}$ can at most be linear in $\rho$ at $g$ order we find
\end{itemize}
\begin{notebox}\textsf{\small Note, l.1092:} [Removed: $\mf{N}^{(2)}\propto\int\rho^a\rho^a/\v{p}^2$ does not commute with a single $\rho^b$, since $[\rho^a(x)\rho^a(y),\rho^b(w)]\neq 0$ for $x\neq y$, so $[\mf{N}^{(2)},\mf{A}^{(1)}]\propto g^3 f^{abc}\{\rho^a,\rho^c\}\neq 0$.]\end{notebox}
\needspace{3\baselineskip}\noindent\rule{\linewidth}{0.3pt}\par\noindent\textbf{Lines 1094--1098} \quad {\small (in \texttt{align}, lines 1093--1096)}\par
\begin{itemize}[leftmargin=2.2em,itemsep=1pt,topsep=2pt]
\item[\textbf{l.1094}] \textsf{\small [term removed/replaced]} \emph{removed:} $\displaystyle \mf {N}^{(2)}\mf{A}^{(1)c}_{1k}(q)$
\item[\textbf{l.1094}] \textsf{\small [term removed/replaced]} \emph{removed:} $\displaystyle = \mf{A}^{(1) c}_{1k}(q) \mf {N}^{(2)}$
\item[\textbf{l.1095}] \textsf{\small [term removed/replaced]} \emph{removed:} $\displaystyle \mf {N}^{(2)}\mf{A}^{(1)c}_{2k}(q)$
\item[\textbf{l.1095}] \textsf{\small [term removed/replaced]} \emph{removed:} $\displaystyle = \mf{A}^{(1) c}_{2k}(q) \mf {N}^{(2)}$
\item[\textbf{l.1097}] \textsf{\small [typo/LaTeX]} and thus 
the relation for $\mf{A}^{(3)\, \dagger\, b}_{1j}(p)$ simplifies to:\ $\longrightarrow$\ Keeping the commutator, the relation for $\mf{A}^{(3)\, \dagger\, b}_{1j}(p)$ reads:
\end{itemize}
\needspace{3\baselineskip}\noindent\rule{\linewidth}{0.3pt}\par\noindent\textbf{Lines 1100--1101} \quad {\small (in \texttt{align}, lines 1099--1101)}\par
\begin{itemize}[leftmargin=2.2em,itemsep=1pt,topsep=2pt]
\item[\textbf{l.1100}] \textsf{\small [missing term added]} \emph{added:} $\displaystyle -\big[\mf{N}^{(2)},\,\mf{A}^{(1)\, b}_{2j}(p)\big]$\quad{\color{gray}\small(in context: }$\displaystyle \mf{A}^{(3)\, b}_{2j}(p) -\big[\mf{N}^{(2)},\,\mf{A}^{(1)\, b}_{2j}(p)\big]${\color{gray}\small)}
\end{itemize}
\begin{notebox}\textsf{\small Note, l.1101:} [With the three-$\rho$ part of $\mf{A}^{(3)}_2$ ordered as $[\mf{A}^{(3)}_2]_{\rho\rho\rho}=\mf{A}^{(1)}_2\mf{N}^{(2)}$ (see the note before Eq.~(\texttt{3rhosmom})), the first two terms give $-[\mf{A}^{(3)\,b}_{2j}]_{\rho\rho\rho}-\big[\mf{N}^{(2)},\mf{A}^{(1)\,b}_{2j}\big]=-\mf{N}^{(2)}\mf{A}^{(1)\,b}_{2j}$.]\end{notebox}
\needspace{3\baselineskip}\noindent\rule{\linewidth}{0.3pt}\par\noindent\textbf{Lines 1146--1147} \quad {\small (in \texttt{align}, lines 1145--1147)}\par
\begin{itemize}[leftmargin=2.2em,itemsep=1pt,topsep=2pt]
\item[\textbf{l.1146}] \textsf{\small [missing term added]} \emph{added:} $\displaystyle -\big[\mf{N}^{(2)},\,\mf{A}^{(1)\, b}_{2j}(p^{+},-x)\big]$\quad{\color{gray}\small(in context: }$\displaystyle \mf{A}^{(3)\, b}_{2j}(p^{+}, -x) -\big[\mf{N}^{(2)},\,\mf{A}^{(1)\, b}_{2j}(p^{+},-x)\big]${\color{gray}\small)}
\end{itemize}
\begin{notebox}\textsf{\small Note, l.1147:} [As in momentum space: $-[\mf{A}^{(3)}_2]_{\rho\rho\rho}-[\mf{N}^{(2)},\mf{A}^{(1)}_2]=-\mf{N}^{(2)}\mf{A}^{(1)}_2$ for $[\mf{A}^{(3)}_2]_{\rho\rho\rho}=\mf{A}^{(1)}_2\mf{N}^{(2)}$.]\end{notebox}
\needspace{3\baselineskip}\noindent\rule{\linewidth}{0.3pt}\par\noindent\textbf{Lines 1240}\par
\begin{itemize}[leftmargin=2.2em,itemsep=1pt,topsep=2pt]
\item[\textbf{l.1240}] \textsf{\small [typo/LaTeX]} \emph{added:} (i.e. $\mf{N}_0\equiv\mf{N}^{(0)}$, $\mf{N}_2\equiv\mf{N}^{(2)}$)
\item[\textbf{l.1240}] \textsf{\small [typo/LaTeX]} $\mf{A}^{(3)}_{2}$\ $\longrightarrow$\ $\mf{A}^{(1)}_{2}$
\end{itemize}
\needspace{3\baselineskip}\noindent\rule{\linewidth}{0.3pt}\par\noindent\textbf{Lines 1243} \quad {\small (in \texttt{align}, lines 1241--1244, label \texttt{A1coeffM})}\par
\begin{itemize}[leftmargin=2.2em,itemsep=1pt,topsep=2pt]
\item[\textbf{l.1243}] \textsf{\small [missing term added]} \emph{added:} $\displaystyle (1)$
\item[\textbf{l.1243}] \textsf{\small [missing term added]} \emph{added:} $\displaystyle (1)$
\end{itemize}
\needspace{3\baselineskip}\noindent\rule{\linewidth}{0.3pt}\par\noindent\textbf{Lines 1248--1249} \quad {\small (in \texttt{align}, lines 1246--1249, label \texttt{A1coeff})}\par
\begin{itemize}[leftmargin=2.2em,itemsep=1pt,topsep=2pt]
\item[\textbf{l.1248}] \textsf{\small [missing term added]} \emph{added:} $\displaystyle (1)$
\item[\textbf{l.1248}] \textsf{\small [sign/factor correction]} \emph{added:} $\displaystyle -$\quad{\color{gray}\small(in context: }$\displaystyle k^{+},-\mathbf{z}${\color{gray}\small)}
\item[\textbf{l.1248}] \textsf{\small [missing term added]} \emph{added:} $\displaystyle (1)$
\end{itemize}
\begin{notebox}\textsf{\small Note, l.1249:} [Argument fixed: with $f(k)=\int d^2x\,e^{ikx}f(x)$ the right-hand side is $-\mf{A}^{(1)}_2(k^+,-\mathbf{z})=\mf{A}^{(1)\dagger}_1(k^+,\mathbf{z})$.]\end{notebox}
\needspace{3\baselineskip}\noindent\rule{\linewidth}{0.3pt}\par\noindent\textbf{Lines 1252--1253} \quad {\small (in \texttt{align}, lines 1251--1253)}\par
\begin{itemize}[leftmargin=2.2em,itemsep=1pt,topsep=2pt]
\item[\textbf{l.1252}] \textsf{\small [term removed/replaced]} $\displaystyle \log{(\frac{\Lambda}{k^{+}})}-2\,\log{(\frac{\vee}{\Lambda})}]$\ $\longrightarrow$\ $\displaystyle \log{(\frac{\Lambda}{k^{+}})}+\log{(\frac{\vee}{k^{+}}-1)}-\log{(\frac{\vee}{\Lambda})}]$
\item[\textbf{l.1252}] \textsf{\small [term removed/replaced]} \emph{removed:} $\displaystyle + \frac{\pi^{2}}{3} - \frac{67}{9} +\log^{2}(\frac{\Lambda}{k^{+}})-\log^{2}(\frac{\vee}{k^{+}})$
\item[\textbf{l.1252}] \textsf{\small [limits/range]} $\displaystyle \frac{\Lambda}{k^{+}}$\ $\longrightarrow$\ $\displaystyle \frac{k^{+}}{\vee}$
\item[\textbf{l.1252}] \textsf{\small [limits/range]} $\displaystyle 1-\frac{\Lambda}{k^{+}}$\ $\longrightarrow$\ $\displaystyle \frac{k^{+}}{k^{+}+\Lambda}$
\item[\textbf{l.1252}] \textsf{\small [operator ordering]} \emph{added:} $\displaystyle \rho^{a}(-\mathbf{k})$\quad{\color{gray}\small(in context: }$\displaystyle g^{3}\,\rho^{a}(-\mathbf{k})\,\rho^{b}(\mathbf{p})\,\rho^{b}(-\mathbf{p})${\color{gray}\small)}
\item[\textbf{l.1252}] \textsf{\small [operator ordering]} \emph{removed:} $\displaystyle \,\rho^{a}(-\mathbf{k})$\quad{\color{gray}\small(original context: }$\displaystyle \,\rho^{b}(\mathbf{p})\,\rho^{b}(-\mathbf{p})\,\rho^{a}(-\mathbf{k})\, \mathbf{k}^{i}${\color{gray}\small)}
\end{itemize}
\begin{notebox}\textsf{\small Note, l.1253:} [One-$\rho$ part: the $\xi$-integral of Eq.~(\texttt{combinedgrho}), after relabelling $k^+\to\xi k^+$ (limits $k^+/\vee<\xi<k^+/(k^++\Lambda)$), adds $\log(\frac{\vee}{k^+}-1)+\log\frac{\vee}{\Lambda}$ to the factor of the pole bracket and turns the constant of the first line into the one of the second line, which was counted in addition. Two-$\rho$ part: the same relabelling fixes the $\xi$ limits. Three-$\rho$ part: $\rho^a$ moved to the left of $\rho^b\rho^b$, so that it equals $\mf{A}^{(1)}_2\mf{N}_2$ (note before Eq.~(\texttt{3rhosmom})); the two-$\rho$ part with its $-2k^i\log(\vee/\Lambda)$ term is unchanged. Checked numerically against Eqs.~(\texttt{combinedgrho}), (\texttt{combinedgrhorho}) and the coordinate form below.]\end{notebox}
\needspace{3\baselineskip}\noindent\rule{\linewidth}{0.3pt}\par\noindent\textbf{Lines 1256} \quad {\small (in \texttt{align}, lines 1255--1257)}\par
\begin{itemize}[leftmargin=2.2em,itemsep=1pt,topsep=2pt]
\item[\textbf{l.1256}] \textsf{\small [operator ordering]} \emph{added:} $\displaystyle \rho^{a}(\mathbf{w})$\quad{\color{gray}\small(in context: }$\displaystyle ig^{3}\,\rho^{a}(\mathbf{w})\,\rho^{b}(\mathbf{x})\,\rho^{b}(\mathbf{y})${\color{gray}\small)}
\item[\textbf{l.1256}] \textsf{\small [operator ordering]} \emph{removed:} $\displaystyle \,\rho^{a}(\mathbf{w})$\quad{\color{gray}\small(original context: }$\displaystyle \,\rho^{b}(\mathbf{x})\,\rho^{b}(\mathbf{y})\,\rho^{a}(\mathbf{w})\, (x-u)\cdot (y-u)\,(w-z)^{i}${\color{gray}\small)}
\item[\textbf{l.1256}] \textsf{\small [term removed/replaced]} $\displaystyle V$\ $\longrightarrow$\ $\displaystyle \vee$
\end{itemize}
\needspace{3\baselineskip}\noindent\rule{\linewidth}{0.3pt}\par\noindent\textbf{Lines 1260--1261} \quad {\small (in \texttt{align}, lines 1259--1261)}\par
\begin{itemize}[leftmargin=2.2em,itemsep=1pt,topsep=2pt]
\item[\textbf{l.1260}] \textsf{\small [limits/range]} $\displaystyle \frac{1}{k^{2}\,p^{+}}]\,\Theta(k^{+}-\Lambda-p^{+})$\ $\longrightarrow$\ $\displaystyle \frac{1}{k^{2}\,p^{+}}]\,\Theta(\vee-k^{+}-p^{+})$
\item[\textbf{l.1260}] \textsf{\small [limits/range]} $\displaystyle \frac{g^{2}p^{j}k^{i}\,\{\rho^{b}(-p),\rho^{a}(-k)\}}{2\,p^{2}k^{2}\, \sqrt{p^{+}k^{+}}}\,\Theta(k^{+}-\Lambda-p^{+})$\ $\longrightarrow$\ $\displaystyle \frac{g^{2}p^{j}k^{i}\,\{\rho^{b}(-p),\rho^{a}(-k)\}}{2\,p^{2}k^{2}\, \sqrt{p^{+}k^{+}}}\,\Theta(\vee-k^{+})\Theta(\vee-p^{+})$
\item[\textbf{l.1260}] \textsf{\small [limits/range]} $\displaystyle \,\Theta(k^{+}-\Lambda-p^{+})$\ $\longrightarrow$\ $\displaystyle \Theta(\vee-k^{+}-p^{+})$
\item[\textbf{l.1260}] \textsf{\small [limits/range]} $\displaystyle \,\frac{ig^{2}f^{abc}\rho^{c}(-p-k)\,(p^{+}-k^{+})\delta_{ji}}{2\,\sqrt{p^{+}k^{+}}\,(p^{+}+k^{+})^{2}\,(\frac{k^{2}}{k^{+}}+\frac{p^{2}}{p^{+}})}\,\Theta(k^{+}-\Lambda-p^{+})$\ $\longrightarrow$\ $\displaystyle \,\frac{ig^{2}f^{abc}\rho^{c}(-p-k)\,(p^{+}-k^{+})\delta_{ji}}{2\,\sqrt{p^{+}k^{+}}\,(p^{+}+k^{+})^{2}\,(\frac{k^{2}}{k^{+}}+\frac{p^{2}}{p^{+}})}\,\Theta(\vee-k^{+}-p^{+})$
\end{itemize}
\begin{notebox}\textsf{\small Note, l.1261:} [$\Theta$ replaced; the support follows from the kinematics ($k^+,p^+>\Lambda$ throughout). Third term (triple vertex): the parent gluon carries $k^++p^+$ and is itself a soft mode, hence $\Theta(\vee-k^+-p^+)$; the instantaneous fourth term has the same support in Eq.~(\texttt{ggthree}) (a convention; no particle carries $k^++p^+$ there). Second term ($\{\rho,\rho\}$): the two gluons are emitted independently by the valence charges, so $\Lambda<k^+<\vee$ and $\Lambda<p^+<\vee$ with no condition on the sum, as for the same intermediate state in Eq.~(\texttt{3rho}); with $k^++p^+<\vee$ the four-$\rho$ terms of the cross section would not cancel for commuting $\rho$ (coherent state), see the note after Eq.~(\texttt{eq:CS1}). The first term is the commutator part of the same diagram, so by the same kinematics its support is also the full square; it is left unchanged here for the author to decide. $\Theta(k^+-\Lambda-p^+)$ is not symmetric under exchange of the two gluons and removes $p^+>k^+$.]\end{notebox}
\needspace{3\baselineskip}\noindent\rule{\linewidth}{0.3pt}\par\noindent\textbf{Lines 1268--1269} \quad {\small (in \texttt{align}, lines 1267--1269)}\par
\begin{itemize}[leftmargin=2.2em,itemsep=1pt,topsep=2pt]
\item[\textbf{l.1268}] \textsf{\small [limits/range]} $\displaystyle \Theta(k^{+}-\Lambda-p^{+})$\ $\longrightarrow$\ $\displaystyle \Theta(\vee-k^{+}-p^{+})$
\item[\textbf{l.1268}] \textsf{\small [limits/range]} $\displaystyle \Theta(k^{+}-\Lambda-p^{+})$\ $\longrightarrow$\ $\displaystyle \Theta(\vee-k^{+})\Theta(\vee-p^{+})$
\end{itemize}
\begin{notebox}\textsf{\small Note, l.1269:} [$\Theta$ replaced as in the momentum-space expression above: $k^+,p^+>\Lambda$, with $k^++p^+<\vee$ in the single-$\rho$ term (parent gluon of the triple vertex; for $k^++p^+>\vee$ only the commutator part of the two-emission diagram could remain, a point for the author to decide) and $k^+<\vee$, $p^+<\vee$ separately in the $\{\rho,\rho\}$ term (two independent emissions).]\end{notebox}
\needspace{3\baselineskip}\noindent\rule{\linewidth}{0.3pt}\par\noindent\textbf{Lines 1272--1273} \quad {\small (in \texttt{align}, lines 1271--1273)}\par
\begin{itemize}[leftmargin=2.2em,itemsep=1pt,topsep=2pt]
\item[\textbf{l.1272}] \textsf{\small [index/Wilson-line orientation]} $\displaystyle ba$\ $\longrightarrow$\ $\displaystyle ab$
\item[\textbf{l.1272}] \textsf{\small [index/Wilson-line orientation]} $\displaystyle ji$\ $\longrightarrow$\ $\displaystyle ij$
\item[\textbf{l.1272}] \textsf{\small [index/Wilson-line orientation]} $\displaystyle y-z$\ $\longrightarrow$\ $\displaystyle y-u$
\item[\textbf{l.1272}] \textsf{\small [index/Wilson-line orientation]} $\displaystyle x-u$\ $\longrightarrow$\ $\displaystyle x-z$
\item[\textbf{l.1272}] \textsf{\small [index/Wilson-line orientation]} $\displaystyle y-z$\ $\longrightarrow$\ $\displaystyle y-u$
\item[\textbf{l.1272}] \textsf{\small [index/Wilson-line orientation]} $\displaystyle x-u$\ $\longrightarrow$\ $\displaystyle x-z$
\end{itemize}
\begin{notebox}\textsf{\small Note, l.1273:} [Labels fixed: with the convention of Eq.~(\texttt{Omegaonvacposzeta}) (second index pair $\leftrightarrow$ first argument) the single-$\rho$ term as written is $\mf{B}^{ab}_{2ij}(\xi k^+,-\v{z};(1-\xi)k^+,-\v{u})$; the two-$\rho$ term is relabelled to match.]\end{notebox}
\needspace{3\baselineskip}\noindent\rule{\linewidth}{0.3pt}\par\noindent\textbf{Lines 1276--1277} \quad {\small (in \texttt{align}, lines 1275--1277)}\par
\begin{itemize}[leftmargin=2.2em,itemsep=1pt,topsep=2pt]
\item[\textbf{l.1276}] \textsf{\small [term removed/replaced]} $\displaystyle [\frac{(k^{+}-p^{+})}{(k-p)^{2}}\,\Theta(p^{+}-k^{+}-\Lambda)$\ $\longrightarrow$\ $\displaystyle [\frac{(p^{+}-k^{+})}{(k-p)^{2}}\,\Theta(p^{+}-k^{+}-\Lambda)$
\end{itemize}
\begin{notebox}\textsf{\small Note, l.1277:} [Sign of the non-instantaneous $p^+>k^+$ term fixed (Eqs.~(\texttt{onegrho12}) and (\texttt{Eq:onegrhoX1}) carry the wrong overall sign): it follows from (\texttt{m2}), (\texttt{m3}), from the cancellation of the $1/(1-\xi)^2$ singularity, and from the $a^\dagger a$ unitarity constraint.]\end{notebox}
\needspace{3\baselineskip}\noindent\rule{\linewidth}{0.3pt}\par\noindent\textbf{Lines 1289--1290} \quad {\small (in \texttt{align}, lines 1279--1290)}\par
\begin{itemize}[leftmargin=2.2em,itemsep=1pt,topsep=2pt]
\item[\textbf{l.1289}] \textsf{\small [sign/factor correction]} \emph{added:} $\displaystyle -$\quad{\color{gray}\small(in context: }$\displaystyle [-\frac{1}{((x-z) - (x-u))^2}\Theta(p^{+} - k^{+} - \Lambda )${\color{gray}\small)}
\end{itemize}
\begin{notebox}\textsf{\small Note, l.1290:} [Sign of the non-instantaneous $p^+>k^+$ term fixed (Eqs.~(\texttt{onegrho12}) and (\texttt{Eq:onegrhoX1}) carry the wrong overall sign): it follows from (\texttt{m2}), (\texttt{m3}), from the cancellation of the $1/(1-\xi)^2$ singularity, and from the $a^\dagger a$ unitarity constraint.]\end{notebox}
\needspace{3\baselineskip}\noindent\rule{\linewidth}{0.3pt}\par\noindent\textbf{Lines 1302--1303} \quad {\small (in \texttt{align}, lines 1292--1303)}\par
\begin{itemize}[leftmargin=2.2em,itemsep=1pt,topsep=2pt]
\item[\textbf{l.1302}] \textsf{\small [sign/factor correction]} \emph{added:} $\displaystyle -$\quad{\color{gray}\small(in context: }$\displaystyle [-\frac{1}{((x-z) - (x-u))^2}\,\Theta(\xi - 1 - \frac{\Lambda}{k^{+}})${\color{gray}\small)}
\end{itemize}
\begin{notebox}\textsf{\small Note, l.1303:} [Sign of the non-instantaneous $p^+>k^+$ term fixed (Eqs.~(\texttt{onegrho12}) and (\texttt{Eq:onegrhoX1}) carry the wrong overall sign): it follows from (\texttt{m2}), (\texttt{m3}), from the cancellation of the $1/(1-\xi)^2$ singularity, and from the $a^\dagger a$ unitarity constraint.]\end{notebox}
\needspace{3\baselineskip}\noindent\rule{\linewidth}{0.3pt}\par\noindent\textbf{Lines 1358} \quad \textsf{Production cross section}\par
\begin{notebox}\textsf{\small Note, l.1358:} [Note: the two-gluon term needs the factor $\frac12$ (present in the next equation) and the vacuum projector $\ket{0}\bra{0}$ is not written; what is used below is the sum over states inserted between $\bar a^{\dagger}$ and $\bar a$.]\end{notebox}
\needspace{3\baselineskip}\noindent\rule{\linewidth}{0.3pt}\par\noindent\textbf{Lines 1386}\par
\begin{itemize}[leftmargin=2.2em,itemsep=1pt,topsep=2pt]
\item[\textbf{l.1386}] \textsf{\small [typo/LaTeX]} $\mf{A}^\dagger\,$, and $\mf{B}^\dagger\,$\ $\longrightarrow$\ $\mf{A}_2$ and $\mf{B}_2$ (the creation terms)
\end{itemize}
\needspace{3\baselineskip}\noindent\rule{\linewidth}{0.3pt}\par\noindent\textbf{Lines 1430} \quad {\small (in \texttt{align}, lines 1429--1431)}\par
\begin{itemize}[leftmargin=2.2em,itemsep=1pt,topsep=2pt]
\item[\textbf{l.1430}] \textsf{\small [typo/LaTeX]} \emph{added:} $\displaystyle j$
\item[\textbf{l.1430}] \textsf{\small [typo/LaTeX]} \emph{added:} $\displaystyle j$
\end{itemize}
\needspace{3\baselineskip}\noindent\rule{\linewidth}{0.3pt}\par\noindent\textbf{Lines 1435--1448} \quad {\small (in \texttt{align}, lines 1433--1448)}\par
\begin{itemize}[leftmargin=2.2em,itemsep=1pt,topsep=2pt]
\item[\textbf{l.1435}] \textsf{\small [missing term added]} \emph{added:} $\displaystyle +\big[\bar{\mf{N}}_{2},\bar{\mf{A}}^{\dagger\,(1)\,a}_{1i}(k^{+},w^{\prime})\big]$\quad{\color{gray}\small(in context: }$\displaystyle \mf{A}^{\dagger\, (1)\, b}_{1i}(k^{+},w^{\prime})+\big[\bar{\mf{N}}_{2},\bar{\mf{A}}^{\dagger\,(1)\,a}_{1i}(k^{+},w^{\prime})\big]${\color{gray}\small)}
\end{itemize}
\begin{notebox}\textsf{\small Note, l.1448:} [Added: the term $[\bar{\mf{N}}_2,\bar{\mf{A}}^{(1)\dagger}_1]$ of $\bar{\mf{A}}^{(3)\dagger}_1$. The formula holds with the three-$\rho$ part of $\mf{A}^{(3)}_2$ ordered as $\mf{A}^{(1)}_2\mf{N}_2$ (note before Eq.~(\texttt{3rhosmom})). In the sum of the $\mf{N}$-type terms the commutator then cancels: in the form of the next equation, $-\bar{\mf{A}}\mf{N}_2+U\bar{\mf{N}}_2\mf{A}-[\bar{\mf{N}}_2,\bar{\mf{A}}]+U[\mf{A}^{(3)}]_{\rho\rho\rho}-[\bar{\mf{A}}^{(3)}]_{\rho\rho\rho}=-\bar{\mf{A}}\mf{N}_2+U\bar{\mf{N}}_2\mf{A}+U\mf{A}\mf{N}_2-\bar{\mf{N}}_2\bar{\mf{A}}$ with $\mf{A}\equiv\mf{A}^{(1)}_2$, which vanishes for $U=1$, as $\bar\Omega^\dagger\Omega=1$ requires.]\end{notebox}
\needspace{3\baselineskip}\noindent\rule{\linewidth}{0.3pt}\par\noindent\textbf{Lines 1451} \quad {\small (in \texttt{align}, lines 1449--1464)}\par
\begin{itemize}[leftmargin=2.2em,itemsep=1pt,topsep=2pt]
\item[\textbf{l.1451}] \textsf{\small [missing term added]} \emph{added:} $\displaystyle -\big[\bar{\mf{N}}_{2},\bar{\mf{A}}^{(1)\,a}_{2i}(k^{+},-w^{\prime})\big]$\quad{\color{gray}\small(in context: }$\displaystyle (k^{+},-w^{\prime})-\big[\bar{\mf{N}}_{2},\bar{\mf{A}}^{(1)\,a}_{2i}(k^{+},-w^{\prime})\big]${\color{gray}\small)}
\end{itemize}
\needspace{3\baselineskip}\noindent\rule{\linewidth}{0.3pt}\par\noindent\textbf{Lines 1469} \quad {\small (in \texttt{align}, lines 1468--1470, label \texttt{eq: A1a})}\par
\begin{itemize}[leftmargin=2.2em,itemsep=1pt,topsep=2pt]
\item[\textbf{l.1469}] \textsf{\small [index/Wilson-line orientation]} $\displaystyle \frac{2}{\sqrt{2\pi}}\int_{w^{\prime}}e^{-i k w^\prime} U^{ bc}(v)\Bar{\mf{B}}^{\dagger ca}_{1ji}(k^{+},w^{\prime};$\ $\longrightarrow$\ $\displaystyle \frac{2}{\sqrt{2\pi}}\int_{w^{\prime}}e^{-i k w^\prime} U^{\dagger\, bc}(v)\Bar{\mf{B}}^{\dagger ca}_{1ji}(k^{+},w^{\prime};$
\end{itemize}
\needspace{3\baselineskip}\noindent\rule{\linewidth}{0.3pt}\par\noindent\textbf{Lines 1478} \quad {\small (in \texttt{align}, lines 1477--1479)}\par
\begin{itemize}[leftmargin=2.2em,itemsep=1pt,topsep=2pt]
\item[\textbf{l.1478}] \textsf{\small [index/Wilson-line orientation]} $\displaystyle (k^{+},-w^{\prime};q^{+},-v)+U^{ bc}(v)\Bar{\mf{B}}^{\dagger ca}_{1ji}(k^{+},w^{\prime};$\ $\longrightarrow$\ $\displaystyle (k^{+},-w^{\prime};q^{+},-v)+U^{\dagger\, bc}(v)\Bar{\mf{B}}^{\dagger ca}_{1ji}(k^{+},w^{\prime};$
\end{itemize}
\needspace{3\baselineskip}\noindent\rule{\linewidth}{0.3pt}\par\noindent\textbf{Lines 1486} \quad {\small (in \texttt{align}, lines 1485--1487)}\par
\begin{itemize}[leftmargin=2.2em,itemsep=1pt,topsep=2pt]
\item[\textbf{l.1486}] \textsf{\small [index/Wilson-line orientation]} $\displaystyle ,z-w^{\prime}; q^{+},w^{\prime}-v)U^{bd}(v)U^{ce}(z) \mf{A}^{\dagger\,(1)\,e}_{1k}(q^{+}$\ $\longrightarrow$\ $\displaystyle ,z-w^{\prime}; q^{+},w^{\prime}-v)U^{\dagger\, bd}(v)U^{ce}(z) \mf{A}^{\dagger\,(1)\,e}_{1k}(q^{+}$
\end{itemize}
\needspace{3\baselineskip}\noindent\rule{\linewidth}{0.3pt}\par\noindent\textbf{Lines 1490--1491} \quad {\small (in \texttt{align}, lines 1489--1491)}\par
\begin{itemize}[leftmargin=2.2em,itemsep=1pt,topsep=2pt]
\item[\textbf{l.1490}] \textsf{\small [index/Wilson-line orientation]} $\displaystyle (q^{+},v)U^{\dagger db}(v) \mf{A}^{\dagger\,(1)\, c}$\ $\longrightarrow$\ $\displaystyle (q^{+},v)U^{\dagger bd}(v) \mf{A}^{\dagger\,(1)\, c}$
\end{itemize}
\begin{notebox}\textsf{\small Note, l.1491:} [Orientation of $U(v)$ fixed in the four places marked above: $\bar a^{\dagger f}=U^{fg}a^{\dagger g}$ gives $U^{cb}(v)=U^{\dagger\, bc}(v)$, as in the next equation.]\end{notebox}
\needspace{3\baselineskip}\noindent\rule{\linewidth}{0.3pt}\par\noindent\textbf{Lines 1530--1537} \quad \textsf{Production cross section $\rangle$ Production cross section at the Next to leading order} \quad {\small (in \texttt{align}, lines 1528--1537, label \texttt{eq:CS1})}\par
\begin{itemize}[leftmargin=2.2em,itemsep=1pt,topsep=2pt]
\item[\textbf{l.1530}] \textsf{\small [phase convention]} $\displaystyle \int_{w, w^\prime} e^{-ik(w^\prime-w)}$\ $\longrightarrow$\ $\displaystyle \int_{w, w^\prime} e^{-ik(w-w^\prime)}$
\item[\textbf{l.1531}] \textsf{\small [index/Wilson-line orientation]} $\displaystyle (k^{+},-w^{\prime})U^{b^{\prime}a}(w^{\prime})$\ $\longrightarrow$\ $\displaystyle (k^{+},-w^{\prime})U^{\dagger\, b^{\prime}a}(w^{\prime})$
\item[\textbf{l.1531}] \textsf{\small [missing term added]} \emph{added:} $\displaystyle -\big[\bar{\mf{N}}_{2},\bar{\mf{A}}^{(1)\,a}_{i}(k^{+},-w)\big]$\quad{\color{gray}\small(in context: }$\displaystyle \mf{A}^{(1)\, b}_{i}(k^{+},-w)-\big[\bar{\mf{N}}_{2},\bar{\mf{A}}^{(1)\,a}_{i}(k^{+},-w)\big]${\color{gray}\small)}
\item[\textbf{l.1531}] \textsf{\small [index/Wilson-line orientation]} $\displaystyle (k^{+},-w^{\prime})U^{b^{\prime}a}(w^{\prime})$\ $\longrightarrow$\ $\displaystyle (k^{+},-w^{\prime})U^{\dagger\, b^{\prime}a}(w^{\prime})$
\item[\textbf{l.1531}] \textsf{\small [index/Wilson-line orientation]} $\displaystyle (k^{+},-w^{\prime})U^{b^{\prime}a}(w^{\prime})$\ $\longrightarrow$\ $\displaystyle (k^{+},-w^{\prime})U^{\dagger\, b^{\prime}a}(w^{\prime})$
\item[\textbf{l.1532}] \textsf{\small [index/Wilson-line orientation]} $\displaystyle (k^{+},-w^{\prime})U^{b^{\prime}a}(w^{\prime})$\ $\longrightarrow$\ $\displaystyle (k^{+},-w^{\prime})U^{\dagger\, b^{\prime}a}(w^{\prime})$
\item[\textbf{l.1532}] \textsf{\small [index/Wilson-line orientation]} $\displaystyle [U^{ac}(w)U^{bd}(x$\ $\longrightarrow$\ $\displaystyle [U^{ad}(w)U^{bc}(x$
\item[\textbf{l.1532}] \textsf{\small [index/Wilson-line orientation]} $\displaystyle [U^{ac}(w)U^{bd}(x) \bar{\mf{A}}^{\dagger\, (1)\, b}$\ $\longrightarrow$\ $\displaystyle [U^{ad}(w)U^{bc}(x) \bar{\mf{A}}^{\dagger\, (1)\, b}$
\item[\textbf{l.1532}] \textsf{\small [index/Wilson-line orientation]} $\displaystyle (k^{+},-w^{\prime})U^{b^{\prime}a}(w^{\prime})$\ $\longrightarrow$\ $\displaystyle (k^{+},-w^{\prime})U^{\dagger\, b^{\prime}a}(w^{\prime})$
\end{itemize}
\begin{notebox}\textsf{\small Note, l.1537:} [Changes: (i) with the conjugated bracket at $w^\prime$ and the other at $w$, as written, the phase is $e^{-ik(w-w^\prime)}$ (definition and LO formula have ${\cal A}$ at $w^\prime$); (ii) $U^{b^\prime a}\to U^{\dagger\, b^\prime a}=U^{ab^\prime}$; (iii) commutator term added in the first line; with the three-$\rho$ part of $\mf{A}^{(3)}$ ordered as $\mf{A}^{(1)}\mf{N}_2$ it cancels in the sum of the $\mf{N}$-type terms of the first two rows, whose total is $-\bar{\mf{A}}\mf{N}_2+U\bar{\mf{N}}_2\mf{A}+U\mf{A}\mf{N}_2-\bar{\mf{N}}_2\bar{\mf{A}}$ (zero at $U=1$), so no three-$\rho$ term arises from the commutator; (iv) colour indices of the $UU\bar{\mf{A}}^\dagger\mf{B}_2$ term as in the amplitude. Ranges of $p^+$: $\Lambda<p^+<\vee$ with $|p^+-k^+|>\Lambda$ in the single-$\rho$ parts for the $\bar{\mf{B}}_3^\dagger$ terms, for the $\bar{\mf{A}}^\dagger\mf{B}_2$ terms $\Lambda<p^+<\vee-k^+$ in the single-$\rho$ part of $\mf{B}_2$ and $\Lambda<p^+<\vee$ in its $\{\rho,\rho\}$ part (for commuting $\rho$ the four-$\rho$ terms of the first four rows and of Eq.~(\texttt{eq:CS2}) cancel only if all of them carry the same $\log(\vee/\Lambda)$), $\Lambda<p^+<k^+-\Lambda$ for the $\bar{\mf{C}}_1^\dagger\mf{B}_2$ terms.]\end{notebox}
\needspace{3\baselineskip}\noindent\rule{\linewidth}{0.3pt}\par\noindent\textbf{Lines 1539--1542} \quad {\small (in \texttt{align}, lines 1538--1542, label \texttt{eq:CS2})}\par
\begin{itemize}[leftmargin=2.2em,itemsep=1pt,topsep=2pt]
\item[\textbf{l.1539}] \textsf{\small [phase convention]} $\displaystyle \int_{z,w,w^{\prime}} e^{-ik(w^{\prime}-w)}$\ $\longrightarrow$\ $\displaystyle \int_{z,w,w^{\prime}} e^{-ik(w-w^{\prime})}$
\item[\textbf{l.1540}] \textsf{\small [index/Wilson-line orientation]} $\displaystyle (k^{+},-w^{\prime};p^{+},-z)U^{ c^{\prime}a}(w^{\prime}) -\Bar{\mf{B}}^{\dagger\, c^{\prime}a}_{2ji}(k^{+},-w^{\prime}$\ $\longrightarrow$\ $\displaystyle (k^{+},-w^{\prime};p^{+},-z)U^{\dagger\, c^{\prime}a}(w^{\prime}) -\Bar{\mf{B}}^{\dagger\, c^{\prime}a}_{2ji}(k^{+},-w^{\prime}$
\item[\textbf{l.1540}] \textsf{\small [index/Wilson-line orientation]} $\displaystyle ; p^{+},-z)U^{bc^{\prime}}(z)$\ $\longrightarrow$\ $\displaystyle ; p^{+},-z)U^{c^{\prime}b}(z)$
\item[\textbf{l.1540}] \textsf{\small [index/Wilson-line orientation]} $\displaystyle ,- w^{\prime})U^{bd^{\prime}}(z)\Bar{\mf{A}}^{\dagger\, (1)\, d^{\prime}}$\ $\longrightarrow$\ $\displaystyle ,- w^{\prime})U^{d^{\prime}b}(z)\Bar{\mf{A}}^{\dagger\, (1)\, d^{\prime}}$
\item[\textbf{l.1540}] \textsf{\small [index/Wilson-line orientation]} $\displaystyle (z)\Bar{\mf{A}}^{\dagger\, (1)\, d^{\prime}}_{j}(p^{+},-z)U^{c^{\prime}a}(w^{\prime})$\ $\longrightarrow$\ $\displaystyle (z)\Bar{\mf{A}}^{\dagger\, (1)\, d^{\prime}}_{j}(p^{+},-z)U^{\dagger\, c^{\prime}a}(w^{\prime})$
\item[\textbf{l.1540}] \textsf{\small [index/Wilson-line orientation]} $\displaystyle [U^{bc}(z)\Bar{\mf{A}}^{(1)\, c}_{j}(p^{+},-z)\Bar{\mf{A}}$\ $\longrightarrow$\ $\displaystyle [U^{\dagger\, bc}(z)\Bar{\mf{A}}^{(1)\, c}_{j}(p^{+},-z)\Bar{\mf{A}}$
\item[\textbf{l.1540}] \textsf{\small [index/Wilson-line orientation]} $\displaystyle ,- w^{\prime})U^{bc^{\prime}}(z)$\ $\longrightarrow$\ $\displaystyle ,- w^{\prime})U^{c^{\prime}b}(z)$
\item[\textbf{l.1540}] \textsf{\small [index/Wilson-line orientation]} $\displaystyle [U^{bc}(z)\Bar{\mf{A}}^{(1)\, a}_{i}(k^{+},- w)\Bar{\mf{A}}$\ $\longrightarrow$\ $\displaystyle [U^{\dagger\, bc}(z)\Bar{\mf{A}}^{(1)\, a}_{i}(k^{+},- w)\Bar{\mf{A}}$
\item[\textbf{l.1540}] \textsf{\small [index/Wilson-line orientation]} $\displaystyle +k^{+},y-w; p^{+},w-z)U^{bd}(z)U^{ce}(y) \mf{A}^{(1)\,e}_{k}(p^{+}$\ $\longrightarrow$\ $\displaystyle +k^{+},y-w; p^{+},w-z)U^{\dagger\, bd}(z)U^{ce}(y) \mf{A}^{(1)\,e}_{k}(p^{+}$
\item[\textbf{l.1540}] \textsf{\small [index/Wilson-line orientation]} $\displaystyle -U^{bd}(z) \Bar{\mf{C}}^{dac}_{1jik}(p^{+}+k^{+},$\ $\longrightarrow$\ $\displaystyle -U^{\dagger\, bd}(z) \Bar{\mf{C}}^{dac}_{1jik}(p^{+}+k^{+},$
\item[\textbf{l.1540}] \textsf{\small [phase convention]} $\displaystyle \int_{z,w,w^{\prime}} e^{-ik(w^{\prime}-w)}$\ $\longrightarrow$\ $\displaystyle \int_{z,w,w^{\prime}} e^{-ik(w-w^{\prime})}$
\item[\textbf{l.1541}] \textsf{\small [index/Wilson-line orientation]} $\displaystyle ,- w^{\prime})U^{bd^{\prime}}(z)\Bar{\mf{A}}^{\dagger\, (1)\, d^{\prime}}$\ $\longrightarrow$\ $\displaystyle ,- w^{\prime})U^{d^{\prime}b}(z)\Bar{\mf{A}}^{\dagger\, (1)\, d^{\prime}}$
\item[\textbf{l.1541}] \textsf{\small [index/Wilson-line orientation]} $\displaystyle (z)\Bar{\mf{A}}^{\dagger\, (1)\, d^{\prime}}_{j}(p^{+},-z)U^{c^{\prime}a}(w^{\prime})$\ $\longrightarrow$\ $\displaystyle (z)\Bar{\mf{A}}^{\dagger\, (1)\, d^{\prime}}_{j}(p^{+},-z)U^{\dagger\, c^{\prime}a}(w^{\prime})$
\item[\textbf{l.1541}] \textsf{\small [index/Wilson-line orientation]} $\displaystyle ,- w^{\prime})U^{bc^{\prime}}(z)$\ $\longrightarrow$\ $\displaystyle ,- w^{\prime})U^{c^{\prime}b}(z)$
\item[\textbf{l.1541}] \textsf{\small [index/Wilson-line orientation]} $\displaystyle ,- w^{\prime})U^{bc^{\prime}}(z)$\ $\longrightarrow$\ $\displaystyle ,- w^{\prime})U^{c^{\prime}b}(z)$
\item[\textbf{l.1541}] \textsf{\small [index/Wilson-line orientation]} $\displaystyle [U^{bc}(z)\Bar{\mf{A}}^{(1)\, c}_{j}(p^{+},-z)\Bar{\mf{A}}$\ $\longrightarrow$\ $\displaystyle [U^{\dagger\, bc}(z)\Bar{\mf{A}}^{(1)\, c}_{j}(p^{+},-z)\Bar{\mf{A}}$
\item[\textbf{l.1541}] \textsf{\small [index/Wilson-line orientation]} $\displaystyle [U^{bc}(z)\Bar{\mf{A}}^{(1)\, c}_{j}(p^{+},-z)\Bar{\mf{A}}$\ $\longrightarrow$\ $\displaystyle [U^{\dagger\, bc}(z)\Bar{\mf{A}}^{(1)\, c}_{j}(p^{+},-z)\Bar{\mf{A}}$
\item[\textbf{l.1541}] \textsf{\small [index/Wilson-line orientation]} $\displaystyle [U^{bc}(z)\Bar{\mf{A}}^{(1)\, a}_{i}(k^{+},- w)\Bar{\mf{A}}$\ $\longrightarrow$\ $\displaystyle [U^{\dagger\, bc}(z)\Bar{\mf{A}}^{(1)\, a}_{i}(k^{+},- w)\Bar{\mf{A}}$
\end{itemize}
\begin{notebox}\textsf{\small Note, l.1542:} [Changes: phase as in the previous equation; orientation of the $U$ matrices made consistent with the amplitude ${\cal A}^{ab}_{ij}$ above and its conjugate ($U$ real: $(U^{\dagger\, bc})^{*}=U^{cb}$). All sixteen products are present with the right weights. Ranges of $p^+$: $p^+>\Lambda$, and $p^++k^+<\vee$ in every term containing $\mf{C}_1$ or the single-$\rho$ part of $\mf{B}_2$ ($p^+<\vee$ otherwise, including the $\{\rho,\rho\}$ part of $\mf{B}_2$).]\end{notebox}
\needspace{3\baselineskip}\noindent\rule{\linewidth}{0.3pt}\par\noindent\textbf{Lines 1636} \quad \textsf{The CGC Hamiltonian: Eikonal Approximation for QCD Hamiltonian} \quad {\small (in \texttt{equation}, lines 1633--1637)}\par
\begin{itemize}[leftmargin=2.2em,itemsep=1pt,topsep=2pt]
\item[\textbf{l.1636}] \textsf{\small [sign/factor correction]} \emph{added:} $\displaystyle -$\quad{\color{gray}\small(in context: }$\displaystyle -\mathbf{p}^{i}${\color{gray}\small)}
\item[\textbf{l.1636}] \textsf{\small [sign/factor correction]} $\displaystyle \mathbf{p}^{i}+\mathbf{q}^{i}$\ $\longrightarrow$\ $\displaystyle -\mathbf{p}^{i}-\mathbf{q}^{i}$
\end{itemize}
\needspace{15\baselineskip}\bigskip
\subsection{Group I}\label{sec:file-group_I}
\noindent\textbf{Original:} \texttt{/root/.claude/uploads/0ab97fcf-94a2-5aec-8bad-586400481685/0073ca94-Group\_I.tex}\\\textbf{Checked copy:} \texttt{/home/user/correlations/nlo\_review/group\_I\_checked/Group\_I\_checked.tex}\\
Lines: 268 (original), 268 (checked copy). Coloured groups of the check: 219 (219 top-level); author\textquotesingle s own coloured groups left unchanged: 10.\\
Records by type: sign/factor correction 1; phase convention 79; index/Wilson-line orientation 4; limits/range 1; missing term added 7; typo/LaTeX 30; note only 16.
\par\noindent Unmarked layout-only edits (line breaks \texttt{\textbackslash notag\textbackslash\textbackslash\ \&}, \texttt{\textbackslash par\textbackslash noindent} around inserted notes, spacing), checked-copy lines: 35, 43, 52, 53, 59, 60, 66, 67, 74, 80, 84, 101, 102, 128, 136, 140, 160, 164, 168, 173, 177, 180, 184, 214, 224, 233, 234, 245, 251, 260, 261, 264.
\subsubsection*{Changes, block by block}
\needspace{3\baselineskip}\noindent\rule{\linewidth}{0.3pt}\par\noindent\textbf{Lines 35}\par
\begin{itemize}[leftmargin=2.2em,itemsep=1pt,topsep=2pt]
\item[\textbf{l.35}] \textsf{\small [typo/LaTeX]} \emph{added:} \textbf{Marking in this checked copy:} changes and additions are in red;
\end{itemize}
\begin{notebox}\textsf{\small Note, l.35:} Line numbers of the source are those of the original file. The phase is $e^{-ik\cdot(w-w^{\prime})}$ throughout (conjugated LO bracket at $w^{\prime}$, order-$g^{3}$ bracket at $w$); c.c. is implied; four-$\rho$ terms are collected in the separate four-$\rho$ file. Operator ordering: the three-$\rho$ term of $\mf{A}^{(3)}$ is $\mf{A}^{(1)}\mf{N}_{2}$ ($\rho^{a}$ to the left; derived in third-order perturbation theory and verified with explicit matrices, see the main file); its consequences are in Rows I and II.\end{notebox}
\needspace{3\baselineskip}\noindent\rule{\linewidth}{0.3pt}\par\noindent\textbf{Lines 38} \quad \textsf{Group I (with c.c.) $\rangle$ Row I { $\checkmark$]}}\par
\begin{notebox}\textsf{\small Note, l.38:} [Confirmed, together with Row II. The $\mf{N}$-terms of the order-$g^{3}$ amplitude at $w$ (first row of the corrected Eq.~(CS1) of the main file) are $-\bar{\mf{A}}^{(1)\,a}_{i}(k^{+},-w)\,\mf{N}_{2}+U^{ab}(w)\,\bar{\mf{N}}_{2}\,\mf{A}^{(1)\,b}_{i}(k^{+},-w)-\big[\bar{\mf{N}}_{2},\bar{\mf{A}}^{(1)\,a}_{i}(k^{+},-w)\big]$ (bars: $\rho^{a}(x)\to U^{ab}(x)\rho^{b}(x)$). The commutator is a $\{\rho,\rho\}$ term, but it cancels against the correctly ordered three-$\rho$ term $-\bar{\mf{A}}^{(1)}\bar{\mf{N}}_{2}$ of $-\bar{\mf{A}}^{(3)}$ in Row II: $-\bar{\mf{A}}^{(1)}\bar{\mf{N}}_{2}-[\bar{\mf{N}}_{2},\bar{\mf{A}}^{(1)}]=-\bar{\mf{N}}_{2}\bar{\mf{A}}^{(1)}$. Rows I and II together therefore contain the four-$\rho$ terms $\big[{\rm LO}\big]^{\dagger}\times\big[-\bar{\mf{A}}\,\mf{N}_{2}+U\,\bar{\mf{N}}_{2}\,\mf{A}+U\,\mf{A}\,\mf{N}_{2}-\bar{\mf{N}}_{2}\,\bar{\mf{A}}\big]$ ($\mf{A}\equiv\mf{A}^{(1)}$; the bracket vanishes for $U=1$) and no three-$\rho$ term from the commutator. The four-$\rho$ terms are collected in the separate four-$\rho$ file.]\end{notebox}
\needspace{3\baselineskip}\noindent\rule{\linewidth}{0.3pt}\par\noindent\textbf{Lines 41--43} \quad \textsf{Group I (with c.c.) $\rangle$ Row II { $\checkmark$]}} \quad {\small (in \texttt{align}, lines 40--42)}\par
\begin{itemize}[leftmargin=2.2em,itemsep=1pt,topsep=2pt]
\item[\textbf{l.41}] \textsf{\small [phase convention]} $\displaystyle \int_{w, w^\prime} e^{-ik(w^\prime-w)}\,$\ $\longrightarrow$\ $\displaystyle \int_{w, w^\prime} e^{-ik(w-w^\prime)}\,$
\item[\textbf{l.41}] \textsf{\small [index/Wilson-line orientation]} $\displaystyle (k^{+},-w^{\prime})U^{b^{\prime}a}(w^{\prime})$\ $\longrightarrow$\ $\displaystyle (k^{+},-w^{\prime})U^{\dagger\, b^{\prime}a}(w^{\prime})$
\end{itemize}
\begin{notebox}\textsf{\small Note, l.43:} [Phase: the conjugated LO bracket is at $w'$ and the $g^{3}$ bracket at $w$, so the definition gives $e^{-ik\cdot(w-w')}$ (as in the LO formula); the original phase gives the result at $-\mathbf{k}$. Corrected here and in all equations of this row. The last, three-$\rho$ term of $\mf{A}^{(3)}$ in coordinate space is $\mf{A}^{(1)}\mf{N}_{2}$ ($\rho^{a}$ to the left of $\rho^{b}\rho^{b}$; main file) and is not contained in the two parts below. Barred part: with the commutator of Row I, $-\bar{\mf{A}}^{(1)}\bar{\mf{N}}_{2}-[\bar{\mf{N}}_{2},\bar{\mf{A}}^{(1)}]=-\bar{\mf{N}}_{2}\bar{\mf{A}}^{(1)}$. Unbarred part: $U\mf{A}^{(1)}\mf{N}_{2}=U\mf{N}_{2}\mf{A}^{(1)}-U[\mf{N}_{2},\mf{A}^{(1)}]$. In the ordering of the four-$\rho$ file ($U\mf{N}_{2}\mf{A}^{(1)}-\bar{\mf{N}}_{2}\bar{\mf{A}}^{(1)}$) the four-$\rho$ term is]\end{notebox}
\needspace{3\baselineskip}\noindent\rule{\linewidth}{0.3pt}\par\noindent\textbf{Lines 43} \quad {\small (in \texttt{align}, lines 43--43)}\par
\begin{itemize}[leftmargin=2.2em,itemsep=1pt,topsep=2pt]
\item[\textbf{l.43}] \textsf{\small [missing term added]} \emph{added:} $\displaystyle \frac{d^{3}N^{\rm II}_{4\rho}}{d^{3}k}$
\item[\textbf{l.43}] \textsf{\small [phase convention]} \emph{added:} $\displaystyle = \frac{1}{(2\pi)^{3}}\,\frac{g^{4}}{16\pi^{5}}\,\frac{1}{k^{+}}\,\log\frac{\vee}{\Lambda}\int_{w,w',x',x,u,v,z}e^{-ik\cdot(w-w')}\,\frac{(x'-w')\cdot(x-w)}{(x'-w')^{2}(x-w)^{2}}\,\frac{(u-z)\cdot(v-z)}{(u-z)^{2}(v-z)^{2}}$
\item[\textbf{l.43}] \textsf{\small [missing term added]} \emph{added:} $\displaystyle \quad\times\big[U^{ab'}(x')-U^{ab'}(w')\big]\rho^{b'}(x')\Big[U^{ab}(w)\,\rho^{c}(u)\rho^{c}(v)\,\rho^{b}(x)-U^{cd}(u)U^{ce}(v)\,U^{ab}(x)\,\rho^{d}(u)\rho^{e}(v)\,\rho^{b}(x)\Big]$
\item[\textbf{l.43}] \textsf{\small [typo/LaTeX]} it\ $\longrightarrow$\ the one- and two-$\rho$ terms of $\mf{A}^{(3)}$
\end{itemize}
\begin{notebox}\textsf{\small Note, l.43:} [The remainder $-U[\mf{N}_{2},\mf{A}^{(1)}]$ is a $\{\rho,\rho\}$ term. Since $[\mf{N}_{2},\mf{A}^{(1)\,a}_{i}(k^{+},-w)]=-\frac{g^{3}}{8\sqrt{2}\pi^{4}\sqrt{k^{+}}}\log\frac{\vee}{\Lambda}\,f^{abc}\int_{x,y,z}\frac{(x-z)\cdot(y-z)}{(x-z)^{2}(y-z)^{2}}\frac{(y-w)^{i}}{(y-w)^{2}}\{\rho^{b}(x),\rho^{c}(y)\}$ is exactly the term $-2k^{i}\log\frac{\vee}{\Lambda}$ of the two-$\rho$ part of $\mf{A}^{(3)}$ (checked: same coefficient and sign, and with explicit SU(2), SU(3) matrices), it removes that term for the unbarred colour structure $f^{hbc}U^{ah}(w)\{\rho^{b}(x),\rho^{c}(y)\}$; for the barred structure the term stays. This is implemented by the red terms added to $N_{2}$, $N_{2a}$, $N_{2b}$ below. Equivalent alternative: no red terms in $N_{2}$, $N_{2a}$, $N_{2b}$ (the bracket of $N_{2}$ as written for both colour structures), and the unbarred four-$\rho$ term written as $U\mf{A}^{(1)}\mf{N}_{2}$, i.e. $U^{ab}(w)\,\rho^{b}(x)\,\rho^{c}(u)\rho^{c}(v)$ here and in Eq.~(4rho) of the four-$\rho$ file.]\end{notebox}
\begin{notebox}\textsf{\small Note, l.43:} [In $N_{2}$, $N_{2b}$ the integration variable $k$ in $d^{2}p\,d^{2}k$ is not the external momentum of the phase.]\end{notebox}
\needspace{3\baselineskip}\noindent\rule{\linewidth}{0.3pt}\par\noindent\textbf{Lines 45} \quad {\small (in \texttt{align}, lines 44--46)}\par
\begin{itemize}[leftmargin=2.2em,itemsep=1pt,topsep=2pt]
\item[\textbf{l.45}] \textsf{\small [phase convention]} $\displaystyle \, \int_{w, w',x,x',} e^{-ik(w^\prime-w)}\, \frac{(x'-w')\cdot (x-w)}{(x'-w')^{2}\, (x-w)^{2}}\,$\ $\longrightarrow$\ $\displaystyle \, \int_{w, w',x,x',} e^{-ik(w-w^\prime)}\, \frac{(x'-w')\cdot (x-w)}{(x'-w')^{2}\, (x-w)^{2}}\,$
\end{itemize}
\needspace{3\baselineskip}\noindent\rule{\linewidth}{0.3pt}\par\noindent\textbf{Lines 49} \quad {\small (in \texttt{align}, lines 48--50)}\par
\begin{itemize}[leftmargin=2.2em,itemsep=1pt,topsep=2pt]
\item[\textbf{l.49}] \textsf{\small [phase convention]} $\displaystyle \, \int_{w, w',x,x',} e^{-ik(w^\prime-w)}\, \frac{(x'-w')\cdot (x-w)}{(x'-w')^{2}\, (x-w)^{2}}\,$\ $\longrightarrow$\ $\displaystyle \, \int_{w, w',x,x',} e^{-ik(w-w^\prime)}\, \frac{(x'-w')\cdot (x-w)}{(x'-w')^{2}\, (x-w)^{2}}\,$
\end{itemize}
\needspace{3\baselineskip}\noindent\rule{\linewidth}{0.3pt}\par\noindent\textbf{Lines 52--53} \quad {\small (in \texttt{align}, lines 51--53)}\par
\begin{itemize}[leftmargin=2.2em,itemsep=1pt,topsep=2pt]
\item[\textbf{l.52}] \textsf{\small [phase convention]} $\displaystyle \,\int_{w,w',x'}\, e^{-ik\cdot (w'-w)}\,\int d^{2}p\,d$\ $\longrightarrow$\ $\displaystyle \,\int_{w,w',x'}\, e^{-ik\cdot (w-w')}\,\int d^{2}p\,d$
\item[\textbf{l.52}] \textsf{\small [phase convention]} \emph{added:} $\displaystyle -\,\frac{1}{(2\pi)^{3}}\,\frac{g^{4}}{64\pi^{6}}\, \frac{1}{k^{+}}\,\int_{w,w^{\prime},x^{\prime}}\, e^{-ik\cdot (w-w^{\prime})}\,\int d^{2}p\,d^{2}k\,\int_{x,y}\, e^{ik\cdot (w-y)} e^{ip\cdot (y-x)}\, \frac{(x^{\prime}-w^{\prime})^{i}}{(x^{\prime}-w^{\prime})^{2}\,\mathbf{p}^{2}\mathbf{k}^{2} }\,\Bigg[+2\,k^{i}\,\log{(\frac{\vee}{\Lambda})}\Bigg]\times$
\item[\textbf{l.52}] \textsf{\small [missing term added]} \emph{added:} $\displaystyle \Bigg[U^{ab^{\prime}}(x^{\prime})\,\rho^{b^{\prime}}(x^{\prime}) - U^{ab^{\prime}}(w^{\prime})\,\rho^{b^{\prime}}(x^{\prime})\Bigg]\times\Bigg[f^{hbc}\,U^{ah}(w)\,\{ \rho^{b}(\mathbf{x}),\,\rho^{c}(\mathbf{y})\}\Bigg]$
\end{itemize}
\begin{notebox}\textsf{\small Note, l.53:} [Red term: it cancels the term $-2k^{i}\log\frac{\vee}{\Lambda}$ of the bracket for the unbarred colour structure $f^{hbc}U^{ah}(w)\{\rho^{b}(x),\rho^{c}(y)\}$, so that this term multiplies only $-f^{abc}U^{bd}(x)U^{ce}(y)\{\rho^{d}(x),\rho^{e}(y)\}$. Reason: the three-$\rho$ term of $\mf{A}^{(3)}$ is $\mf{A}^{(1)}\mf{N}_{2}$, while the four-$\rho$ file has the unbarred term as $U\mf{N}_{2}\mf{A}^{(1)}$; the difference $-U[\mf{N}_{2},\mf{A}^{(1)}]$ is this red term (note above). All other terms of the bracket are common to both structures.]\end{notebox}
\needspace{3\baselineskip}\noindent\rule{\linewidth}{0.3pt}\par\noindent\textbf{Lines 58--60} \quad {\small (in \texttt{align}, lines 57--60)}\par
\begin{itemize}[leftmargin=2.2em,itemsep=1pt,topsep=2pt]
\item[\textbf{l.58}] \textsf{\small [phase convention]} $\displaystyle \,{\color{blue}\log\frac{k^{+}}{\Lambda}}\int_{w,w'} e^{-ik\cdot(w'-w)}\int_{x,y,x'}\frac{(x'-w')^{i}}{(x'-w')^{2}}$\ $\longrightarrow$\ $\displaystyle \,{\color{blue}\log\frac{k^{+}}{\Lambda}}\int_{w,w'} e^{-ik\cdot(w-w')}\int_{x,y,x'}\frac{(x'-w')^{i}}{(x'-w')^{2}}$
\item[\textbf{l.59}] \textsf{\small [phase convention]} \emph{added:} $\displaystyle +\frac{1}{(2\pi)^{3}}\,\frac{i\,g^{4}}{16\pi^{5}}\,\frac{1}{k^{+}}\,\log\frac{k^{+}}{\Lambda}\int_{w,w^{\prime}} e^{-ik\cdot(w-w^{\prime})}\int_{x,y,x^{\prime},z}\frac{(x^{\prime}-w^{\prime})^{i}}{(x^{\prime}-w^{\prime})^{2}}\,\frac{(y-w)^{i}}{(y-w)^{2}}\,\frac{(x-z)\cdot(y-z)}{(x-z)^{2}\,(y-z)^{2}}$
\item[\textbf{l.59}] \textsf{\small [missing term added]} \emph{added:} $\displaystyle \times\Big[U^{ab^{\prime}}(x^{\prime})\rho^{b^{\prime}}(x^{\prime})-U^{ab^{\prime}}(w^{\prime})\rho^{b^{\prime}}(x^{\prime})\Big]\Big[f^{hbc}U^{ah}(w)\{\rho^{b}(x),\rho^{c}(y)\}\Big]$
\end{itemize}
\begin{notebox}\textsf{\small Note, l.60:} [Red term: the $\log\frac{k^{+}}{\Lambda}$ part of the red term of $N_{2}$, with $\int d^{2}k\,e^{ik\cdot(w-y)}\frac{k^{i}}{k^{2}}=-2\pi i\frac{(y-w)^{i}}{(y-w)^{2}}$ and $\int d^{2}p\,\frac{e^{ip\cdot(y-x)}}{p^{2}}=\int_{z}\frac{(x-z)\cdot(y-z)}{(x-z)^{2}(y-z)^{2}}$. The $z$ integral is logarithmically divergent at large $z$ and is left unevaluated; with the rule $\int d^{2}p/p^{2}=0$ for its constant it is $-\pi\log(x-y)^{2}$ (checked numerically: the difference of the integral at two separations $r$, $r_{0}$ is $-\pi\log(r^{2}/r_{0}^{2})$). In the bracket of the first two lines the piece coming from $-2k^{i}$ is, by the antisymmetry of the colour structure under $x\leftrightarrow y$, the $\log(x-y)^{2}$ part; hence for the unbarred structure the kernel reduces to $\frac{(x-w)^{i}}{(x-w)^{2}}\log(w-y)^{2}-\frac{(y-w)^{i}}{(y-w)^{2}}\log(w-x)^{2}$, and for the barred structure it is as written. The $\log\frac{k^{+}}{\Lambda}$ term removed in this way (unbarred structure with the $\log(x-y)^{2}$ kernel) is the one that the JIMWLK and three-$\rho$ lists do not contain.]\end{notebox}
\needspace{3\baselineskip}\noindent\rule{\linewidth}{0.3pt}\par\noindent\textbf{Lines 63--67} \quad {\small (in \texttt{align}, lines 62--67)}\par
\begin{itemize}[leftmargin=2.2em,itemsep=1pt,topsep=2pt]
\item[\textbf{l.63}] \textsf{\small [typo/LaTeX]} \emph{added:} $\displaystyle ,x'$\quad{\color{gray}\small(in context: }$\displaystyle \,\frac{1}{k^{+}}\int_{w,w',x'} e^{-ik\cdot(w-w')}${\color{gray}\small)}
\item[\textbf{l.63}] \textsf{\small [phase convention]} $\displaystyle \int_{w,w'} e^{-ik\cdot(w'-w)}$\ $\longrightarrow$\ $\displaystyle \int_{w,w',x'} e^{-ik\cdot(w-w')}$
\item[\textbf{l.66}] \textsf{\small [phase convention]} \emph{added:} $\displaystyle +\frac{1}{(2\pi)^{3}}\,\frac{i\,g^{4}}{16\pi^{5}}\,\frac{1}{k^{+}}\,\log\frac{\vee}{k^{+}}\int_{w,w^{\prime}} e^{-ik\cdot(w-w^{\prime})}\int_{x,y,x^{\prime},z}\frac{(x^{\prime}-w^{\prime})^{i}}{(x^{\prime}-w^{\prime})^{2}}\,\frac{(y-w)^{i}}{(y-w)^{2}}\,\frac{(x-z)\cdot(y-z)}{(x-z)^{2}\,(y-z)^{2}}$
\item[\textbf{l.66}] \textsf{\small [missing term added]} \emph{added:} $\displaystyle \times\Big[U^{ab^{\prime}}(x^{\prime})\rho^{b^{\prime}}(x^{\prime})-U^{ab^{\prime}}(w^{\prime})\rho^{b^{\prime}}(x^{\prime})\Big]\Big[f^{hbc}U^{ah}(w)\{\rho^{b}(x),\rho^{c}(y)\}\Big]$
\end{itemize}
\begin{notebox}\textsf{\small Note, l.67:} [Red term: the $\log\frac{\vee}{k^{+}}$ part of the red term of $N_{2}$; in momentum space it is $R^{i}\to R^{i}+2k^{i}\log\frac{\vee}{k^{+}}$ for the unbarred structure. This is the term which the BFKL file has as the Row II contribution: with $\int_{z}\frac{(a-z)\cdot(b-z)}{(a-z)^{2}(b-z)^{2}}=-\pi\log(a-b)^{2}$ (constant dropped) and the antisymmetry under $x\leftrightarrow y$, both are $-\frac{1}{(2\pi)^{3}}\frac{i\,g^{4}}{16\pi^{4}}\frac{1}{k^{+}}\log\frac{\vee}{k^{+}}\,\frac{(y-w)^{i}}{(y-w)^{2}}\log(x-y)^{2}$ times the same factors.]\end{notebox}
\needspace{3\baselineskip}\noindent\rule{\linewidth}{0.3pt}\par\noindent\textbf{Lines 74}\par
\begin{notebox}\textsf{\small Note, l.74:} [$R^{i}$ as written holds for the barred colour structure; its first term contains the $-2k^{i}\log\frac{\vee}{k^{+}}$ from $-2k^{i}\log\frac{\vee}{\Lambda}$. For the unbarred structure $R^{i}\to R^{i}+2k^{i}\log\frac{\vee}{k^{+}}$, i.e. the first term is $-2k^{i}\log\Big(\frac{k^{+}p^{2}}{\vee k^{2}+k^{+}p^{2}}\Big)$; this is the red term of $N_{2b}$.]\end{notebox}
\needspace{3\baselineskip}\noindent\rule{\linewidth}{0.3pt}\par\noindent\textbf{Lines 77--78} \quad \textsf{Group I (with c.c.) $\rangle$ Row III { $\checkmark$]}} \quad {\small (in \texttt{align}, lines 76--78)}\par
\begin{itemize}[leftmargin=2.2em,itemsep=1pt,topsep=2pt]
\item[\textbf{l.77}] \textsf{\small [phase convention]} $\displaystyle \int_{w, w^\prime} e^{-ik(w^\prime-w)}\, \int\,\frac{dp^{+}}{2\pi}$\ $\longrightarrow$\ $\displaystyle \int_{w, w^\prime} e^{-ik(w-w^\prime)}\, \int\,\frac{dp^{+}}{2\pi}$
\item[\textbf{l.77}] \textsf{\small [index/Wilson-line orientation]} $\displaystyle (k^{+},-w^{\prime})U^{b^{\prime}a}(w^{\prime})$\ $\longrightarrow$\ $\displaystyle (k^{+},-w^{\prime})U^{\dagger\, b^{\prime}a}(w^{\prime})$
\end{itemize}
\begin{notebox}\textsf{\small Note, l.78:} [Phase: with the conjugated LO bracket at $w'$ and the $g^{3}$ bracket at $w$ the definition gives $e^{-ik(w-w')}$, as in the corrected CS1; without averaging the expression is not even in $k$, so this is $k\to-k$ of what was written. Colour: $U^{b'a}(w')\to U^{\dagger\,b'a}(w')=U^{ab'}(w')$, which is what is used from the next equation on.]\end{notebox}
\needspace{3\baselineskip}\noindent\rule{\linewidth}{0.3pt}\par\noindent\textbf{Lines 80--82} \quad {\small (in \texttt{align}, lines 79--81)}\par
\begin{itemize}[leftmargin=2.2em,itemsep=1pt,topsep=2pt]
\item[\textbf{l.80}] \textsf{\small [phase convention]} $\displaystyle \int_{w, w^\prime} e^{-ik(w^\prime-w)}\,$\ $\longrightarrow$\ $\displaystyle \int_{w, w^\prime} e^{-ik(w-w^\prime)}\,$
\item[\textbf{l.80}] \textsf{\small [limits/range]} $\displaystyle \, \frac{ig^{4}\,f^{abd}}{8\pi^{4}}\, \int_{\Lambda}^{\vee}\, \frac{dp^{+}}{2\pi}\,$\ $\longrightarrow$\ $\displaystyle \, \frac{ig^{4}\,f^{abd}}{8\pi^{4}}\, \Big(\int_{\Lambda}^{k^{+}-\Lambda}+\int_{k^{+}+\Lambda}^{\vee}\Big)\, \frac{dp^{+}}{2\pi}\,$
\item[\textbf{l.80}] \textsf{\small [typo/LaTeX]} \emph{added:} $\displaystyle \int_{x,y,z,x'}$\quad{\color{gray}\small(in context: }$\displaystyle \, \frac{ig^{4}\,f^{abd}}{8\pi^{4}}\, \Big(\int_{\Lambda}^{k^{+}-\Lambda}+\int_{k^{+}+\Lambda}^{\vee}\Big)\, \frac{dp^{+}}{2\pi}\, \int_{x,y,z,x'}\,\frac{(x'-w')\cdot (x-w)\, (x-z)\cdot (y-z)}{(x'-w')^{2}\, (x-w)^{2}\, (x-z)^{2}\, (y-z)^{2}}\, \frac{[k^{+}\,(x-w)^{2}+p^{+}\,(x-z)^{2}]}{[p^{+}\,(x-z)^{2}-k^{+}\,(x-w)^{2}]}\, \frac{1}{p^{+}\,k^{+}}${\color{gray}\small)}
\end{itemize}
\begin{notebox}\textsf{\small Note, l.82:} [Checked: the equation above is what the corrected $\mf{B}_{3}$ gives; the sign of its $p^{+}>k^{+}$ non-instantaneous term (fourth term) already agrees with the corrected sign, so the $1/\Lambda$ terms cancel between $N_{1}$ and $N_{2}$, $N_{3a}$. The poles at $p^{+}(x-z)^{2}=k^{+}(x-w)^{2}$ are principal values, so every logarithm below is $\log|\dots|$. In the third term the window $|p^{+}-k^{+}|<\Lambda$ is now excluded, as in the wave functions; this adds the two red logarithms in the last two terms of $N_{1}$. The four-$\rho$ term from the $\{\rho,\rho\}$ piece of $\mf{B}_{3}$ is collected in the separate four-$\rho$ file.]\end{notebox}
\needspace{3\baselineskip}\noindent\rule{\linewidth}{0.3pt}\par\noindent\textbf{Lines 84} \quad {\small (in \texttt{align}, lines 83--85)}\par
\begin{itemize}[leftmargin=2.2em,itemsep=1pt,topsep=2pt]
\item[\textbf{l.84}] \textsf{\small [phase convention]} $\displaystyle \int_{w, w^\prime} e^{-ik(w^\prime-w)}\,$\ $\longrightarrow$\ $\displaystyle \int_{w, w^\prime} e^{-ik(w-w^\prime)}\,$
\item[\textbf{l.84}] \textsf{\small [missing term added]} \emph{added:} $\displaystyle -\log\frac{k^{+}+\Lambda}{k^{+}-\Lambda}$\quad{\color{gray}\small(in context: }$\displaystyle \log\frac{k^{+}}{\Lambda}-\log\frac{k^{+}+\Lambda}{k^{+}-\Lambda}\Big]${\color{gray}\small)}
\item[\textbf{l.84}] \textsf{\small [missing term added]} \emph{added:} $\displaystyle -\log\frac{(k^{+}+\Lambda)(x-z)^{2}-k^{+}(x-w)^{2}}{(k^{+}-\Lambda)(x-z)^{2}-k^{+}(x-w)^{2}}$\quad{\color{gray}\small(in context: }$\displaystyle \qquad-\log\frac{(k^{+}+\Lambda)(x-z)^{2}-k^{+}(x-w)^{2}}{(k^{+}-\Lambda)(x-z)^{2}-k^{+}(x-w)^{2}}\Big]${\color{gray}\small)}
\end{itemize}
\needspace{3\baselineskip}\noindent\rule{\linewidth}{0.3pt}\par\noindent\textbf{Lines 88} \quad {\small (in \texttt{align}, lines 87--89)}\par
\begin{itemize}[leftmargin=2.2em,itemsep=1pt,topsep=2pt]
\item[\textbf{l.88}] \textsf{\small [phase convention]} $\displaystyle \int_{w, w^\prime} e^{-ik(w^\prime-w)}\,$\ $\longrightarrow$\ $\displaystyle \int_{w, w^\prime} e^{-ik(w-w^\prime)}\,$
\end{itemize}
\needspace{3\baselineskip}\noindent\rule{\linewidth}{0.3pt}\par\noindent\textbf{Lines 92--108} \quad {\small (in \texttt{align}, lines 91--110)}\par
\begin{itemize}[leftmargin=2.2em,itemsep=1pt,topsep=2pt]
\item[\textbf{l.92}] \textsf{\small [phase convention]} $\displaystyle \int_{w,w'}e^{-ik(w'-w)}$\ $\longrightarrow$\ $\displaystyle \int_{w,w'}e^{-ik(w-w')}$
\item[\textbf{l.97}] \textsf{\small [typo/LaTeX]} \emph{added:} $\displaystyle \int_{x,y,z,x'}$\quad{\color{gray}\small(in context: }$\displaystyle \frac{ig^{4}f^{abd}}{8\pi^{5}}\,\frac{1}{k^{+}}\,\int_{x,y,z,x'}\log{[\frac{(x-w)^{2}\,k^{+}-(x-z)^{2}\,(k^{+}-\Lambda)}{(x-w)^{2}\,k^{+}-(x-z)^{2}\,\Lambda}]}\,${\color{gray}\small)}
\item[\textbf{l.97}] \textsf{\small [typo/LaTeX]} \emph{removed:} $\displaystyle \int_{x,y,z,x'}$\quad{\color{gray}\small(original context: }$\displaystyle \frac{ig^{4}f^{abd}}{8\pi^{5}}\,\frac{1}{k^{+}}\,\log{[\frac{(x-w)^{2}\,k^{+}-(x-z)^{2}\,(k^{+}-\Lambda)}{(x-w)^{2}\,k^{+}-(x-z)^{2}\,\Lambda}]}\,\int_{x,y,z,x'}${\color{gray}\small)}
\item[\textbf{l.98}] \textsf{\small [typo/LaTeX]} \emph{added:} $\displaystyle \int_{x,y,z,x'}$\quad{\color{gray}\small(in context: }$\displaystyle \frac{ig^{4}f^{abd}}{8\pi^{5}}\,\frac{1}{k^{+}}\,\int_{x,y,z,x'}\log{[\frac{(x-w)^{2}\,k^{+}-(x-z)^{2}\,(k^{+}-\Lambda)}{(x-w)^{2}\,k^{+}-(x-z)^{2}\,\Lambda}]}\,${\color{gray}\small)}
\item[\textbf{l.98}] \textsf{\small [typo/LaTeX]} \emph{removed:} $\displaystyle \int_{x,y,z,x'}$\quad{\color{gray}\small(original context: }$\displaystyle \frac{ig^{4}f^{abd}}{8\pi^{5}}\,\frac{1}{k^{+}}\,\log{[\frac{(x-w)^{2}\,k^{+}-(x-z)^{2}\,(k^{+}-\Lambda)}{(x-w)^{2}\,k^{+}-(x-z)^{2}\,\Lambda}]}\,\int_{x,y,z,x'}${\color{gray}\small)}
\item[\textbf{l.100}] \textsf{\small [typo/LaTeX]} \emph{added:} $\displaystyle \int_{x,y,z,x'}$\quad{\color{gray}\small(in context: }$\displaystyle \frac{ig^{4}f^{abd}}{8\pi^{5}}\,\frac{1}{k^{+}}\,\int_{x,y,z,x'}\log{[\frac{(x-w)^{2}\,k^{+}-(x-z)^{2}\,(k^{+}-\Lambda)}{(x-w)^{2}\,k^{+}-(x-z)^{2}\,\Lambda}]}\,${\color{gray}\small)}
\item[\textbf{l.100}] \textsf{\small [typo/LaTeX]} \emph{removed:} $\displaystyle \int_{x,y,z,x'}$\quad{\color{gray}\small(original context: }$\displaystyle \frac{ig^{4}f^{abd}}{8\pi^{5}}\,\frac{1}{k^{+}}\,\log{[\frac{(x-w)^{2}\,k^{+}-(x-z)^{2}\,(k^{+}-\Lambda)}{(x-w)^{2}\,k^{+}-(x-z)^{2}\,\Lambda}]}\,\int_{x,y,z,x'}${\color{gray}\small)}
\item[\textbf{l.101}] \textsf{\small [typo/LaTeX]} \emph{added:} $\displaystyle \int_{x,y,z,x'}$\quad{\color{gray}\small(in context: }$\displaystyle \frac{ig^{4}f^{abd}}{8\pi^{5}}\,\frac{1}{k^{+}}\,\int_{x,y,z,x'}\log{[\frac{(x-w)^{2}\,k^{+}-(x-z)^{2}\,(k^{+}-\Lambda)}{(x-w)^{2}\,k^{+}-(x-z)^{2}\,\Lambda}]}\,${\color{gray}\small)}
\item[\textbf{l.101}] \textsf{\small [typo/LaTeX]} \emph{removed:} $\displaystyle \int_{x,y,z,x'}$\quad{\color{gray}\small(original context: }$\displaystyle \frac{ig^{4}f^{abd}}{8\pi^{5}}\,\frac{1}{k^{+}}\,\log{[\frac{(x-w)^{2}\,k^{+}-(x-z)^{2}\,(k^{+}-\Lambda)}{(x-w)^{2}\,k^{+}-(x-z)^{2}\,\Lambda}]}\,\int_{x,y,z,x'}${\color{gray}\small)}
\item[\textbf{l.102}] \textsf{\small [typo/LaTeX]} \emph{added:} $\displaystyle \int_{x,y,z,x'}$\quad{\color{gray}\small(in context: }$\displaystyle \frac{ig^{4}f^{abd}}{8\pi^{5}}\,\frac{1}{k^{+}}\,\int_{x,y,z,x'}\log{[\frac{(x-w)^{2}\,k^{+}-(x-z)^{2}\,(k^{+}-\Lambda)}{(x-w)^{2}\,k^{+}-(x-z)^{2}\,\Lambda}]}\,${\color{gray}\small)}
\item[\textbf{l.102}] \textsf{\small [typo/LaTeX]} \emph{removed:} $\displaystyle \int_{x,y,z,x'}$\quad{\color{gray}\small(original context: }$\displaystyle \frac{ig^{4}f^{abd}}{8\pi^{5}}\,\frac{1}{k^{+}}\,\log{[\frac{(x-w)^{2}\,k^{+}-(x-z)^{2}\,(k^{+}-\Lambda)}{(x-w)^{2}\,k^{+}-(x-z)^{2}\,\Lambda}]}\,\int_{x,y,z,x'}${\color{gray}\small)}
\item[\textbf{l.104}] \textsf{\small [typo/LaTeX]} \emph{added:} $\displaystyle \int_{x,y,z,x'}$\quad{\color{gray}\small(in context: }$\displaystyle \frac{ig^{4}f^{abd}}{8\pi^{5}}\,\frac{1}{k^{+}}\,\int_{x,y,z,x'}\log{[\frac{(x-w)^{2}\,k^{+}-(x-z)^{2}\,(k^{+}-\Lambda)}{(x-w)^{2}\,k^{+}-(x-z)^{2}\,\Lambda}]}\,${\color{gray}\small)}
\item[\textbf{l.104}] \textsf{\small [typo/LaTeX]} \emph{removed:} $\displaystyle \int_{x,y,z,x'}$\quad{\color{gray}\small(original context: }$\displaystyle \frac{ig^{4}f^{abd}}{8\pi^{5}}\,\frac{1}{k^{+}}\,\log{[\frac{(x-w)^{2}\,k^{+}-(x-z)^{2}\,(k^{+}-\Lambda)}{(x-w)^{2}\,k^{+}-(x-z)^{2}\,\Lambda}]}\,\int_{x,y,z,x'}${\color{gray}\small)}
\item[\textbf{l.105}] \textsf{\small [typo/LaTeX]} \emph{added:} $\displaystyle \int_{x,y,z,x'}$\quad{\color{gray}\small(in context: }$\displaystyle \frac{ig^{4}f^{abd}}{8\pi^{5}}\,\frac{1}{k^{+}}\,\int_{x,y,z,x'}\mathcal{A}\,${\color{gray}\small)}
\item[\textbf{l.105}] \textsf{\small [typo/LaTeX]} \emph{removed:} $\displaystyle \int_{x,y,z,x'}$\quad{\color{gray}\small(original context: }$\displaystyle \frac{ig^{4}f^{abd}}{8\pi^{5}}\,\frac{1}{k^{+}}\,\mathcal{A}\,\int_{x,y,z,x'}${\color{gray}\small)}
\item[\textbf{l.107}] \textsf{\small [typo/LaTeX]} \emph{added:} $\displaystyle \int_{x,y,z,x'}$\quad{\color{gray}\small(in context: }$\displaystyle \frac{ig^{4}f^{abd}}{8\pi^{5}}\,\frac{1}{k^{+}}\,\int_{x,y,z,x'}\mathcal{A}\,${\color{gray}\small)}
\item[\textbf{l.107}] \textsf{\small [typo/LaTeX]} \emph{removed:} $\displaystyle \int_{x,y,z,x'}$\quad{\color{gray}\small(original context: }$\displaystyle \frac{ig^{4}f^{abd}}{8\pi^{5}}\,\frac{1}{k^{+}}\,\mathcal{A}\,\int_{x,y,z,x'}${\color{gray}\small)}
\item[\textbf{l.108}] \textsf{\small [typo/LaTeX]} \emph{added:} $\displaystyle \int_{x,y,z,x'}$\quad{\color{gray}\small(in context: }$\displaystyle \frac{ig^{4}f^{abd}}{8\pi^{5}}\frac{1}{k^{+}}\,\int_{x,y,z,x'}\mathcal{A}\,${\color{gray}\small)}
\item[\textbf{l.108}] \textsf{\small [typo/LaTeX]} \emph{removed:} $\displaystyle \int_{x,y,z,x'}$\quad{\color{gray}\small(original context: }$\displaystyle \frac{ig^{4}f^{abd}}{8\pi^{5}}\frac{1}{k^{+}}\,\mathcal{A}\,\int_{x,y,z,x'}${\color{gray}\small)}
\end{itemize}
\needspace{3\baselineskip}\noindent\rule{\linewidth}{0.3pt}\par\noindent\textbf{Lines 113--118} \quad {\small (in \texttt{align}, lines 112--120)}\par
\begin{itemize}[leftmargin=2.2em,itemsep=1pt,topsep=2pt]
\item[\textbf{l.113}] \textsf{\small [phase convention]} $\displaystyle \int_{w,w'}e^{-ik(w'-w)}$\ $\longrightarrow$\ $\displaystyle \int_{w,w'}e^{-ik(w-w')}$
\item[\textbf{l.113}] \textsf{\small [typo/LaTeX]} \emph{added:} $\displaystyle \int_{x,y,z,x'}$\quad{\color{gray}\small(in context: }$\displaystyle \frac{ig^{4}f^{abd}}{8\pi^{5}}\,\frac{1}{k^{+}}\,\int_{x,y,z,x'}\mathcal{A}\,${\color{gray}\small)}
\item[\textbf{l.113}] \textsf{\small [typo/LaTeX]} \emph{removed:} $\displaystyle \int_{x,y,z,x'}$\quad{\color{gray}\small(original context: }$\displaystyle \frac{ig^{4}f^{abd}}{8\pi^{5}}\,\frac{1}{k^{+}}\,\mathcal{A}\,\int_{x,y,z,x'}${\color{gray}\small)}
\item[\textbf{l.114}] \textsf{\small [typo/LaTeX]} \emph{added:} $\displaystyle \int_{x,y,z,x'}$\quad{\color{gray}\small(in context: }$\displaystyle \frac{ig^{4}f^{abd}}{16\pi^{5}}\,\frac{1}{k^{+}}\,\int_{x,y,z,x'}\log{[\frac{[k^{+}\,(x-w)-(k^{+}-\Lambda)\,(x-z)]^{2}}{[k^{+}\,(x-w)-\Lambda\,(x-z)]^{2}}]}\,${\color{gray}\small)}
\item[\textbf{l.114}] \textsf{\small [typo/LaTeX]} \emph{removed:} $\displaystyle \int_{x,y,z,x'}$\quad{\color{gray}\small(original context: }$\displaystyle \frac{ig^{4}f^{abd}}{16\pi^{5}}\,\frac{1}{k^{+}}\,\log{[\frac{[k^{+}\,(x-w)-(k^{+}-\Lambda)\,(x-z)]^{2}}{[k^{+}\,(x-w)-\Lambda\,(x-z)]^{2}}]}\,\int_{x,y,z,x'}${\color{gray}\small)}
\item[\textbf{l.116}] \textsf{\small [typo/LaTeX]} \emph{added:} $\displaystyle \int_{x,y,z,x'}$\quad{\color{gray}\small(in context: }$\displaystyle \frac{ig^{4}f^{abd}}{16\pi^{5}}\,\frac{1}{k^{+}}\,\int_{x,y,z,x'}\log{[\frac{[k^{+}\,(x-w)-(k^{+}-\Lambda)\,(x-z)]^{2}}{[k^{+}\,(x-w)-\Lambda\,(x-z)]^{2}}]}\,${\color{gray}\small)}
\item[\textbf{l.116}] \textsf{\small [typo/LaTeX]} \emph{removed:} $\displaystyle \int_{x,y,z,x'}$\quad{\color{gray}\small(original context: }$\displaystyle \frac{ig^{4}f^{abd}}{16\pi^{5}}\,\frac{1}{k^{+}}\,\log{[\frac{[k^{+}\,(x-w)-(k^{+}-\Lambda)\,(x-z)]^{2}}{[k^{+}\,(x-w)-\Lambda\,(x-z)]^{2}}]}\,\int_{x,y,z,x'}${\color{gray}\small)}
\item[\textbf{l.118}] \textsf{\small [typo/LaTeX]} \emph{added:} $\displaystyle \int_{x,y,z,x'}$\quad{\color{gray}\small(in context: }$\displaystyle \frac{ig^{4}f^{abd}}{16\pi^{5}}\,\frac{1}{k^{+}}\,\int_{x,y,z,x'}\log{[\frac{[k^{+}\,(x-w)-(k^{+}-\Lambda)\,(x-z)]^{2}}{[k^{+}\,(x-w)-\Lambda\,(x-z)]^{2}}]}\,${\color{gray}\small)}
\item[\textbf{l.118}] \textsf{\small [typo/LaTeX]} \emph{removed:} $\displaystyle \int_{x,y,z,x'}$\quad{\color{gray}\small(original context: }$\displaystyle \frac{ig^{4}f^{abd}}{16\pi^{5}}\,\frac{1}{k^{+}}\,\log{[\frac{[k^{+}\,(x-w)-(k^{+}-\Lambda)\,(x-z)]^{2}}{[k^{+}\,(x-w)-\Lambda\,(x-z)]^{2}}]}\,\int_{x,y,z,x'}${\color{gray}\small)}
\end{itemize}
\needspace{3\baselineskip}\noindent\rule{\linewidth}{0.3pt}\par\noindent\textbf{Lines 127--128} \quad \textsf{Group I (with c.c.) $\rangle$ Row IV { $\checkmark$]}} \quad {\small (in \texttt{align}, lines 126--128)}\par
\begin{itemize}[leftmargin=2.2em,itemsep=1pt,topsep=2pt]
\item[\textbf{l.127}] \textsf{\small [phase convention]} $\displaystyle \int_{w, w^\prime} e^{-ik(w^\prime-w)}\, \int_{\Lambda}^{\vee - k^{+}}$\ $\longrightarrow$\ $\displaystyle \int_{w, w^\prime} e^{-ik(w-w^\prime)}\, \int_{\Lambda}^{\vee - k^{+}}$
\item[\textbf{l.127}] \textsf{\small [index/Wilson-line orientation]} $\displaystyle (k^{+},-w^{\prime})U^{b^{\prime}a}(w^{\prime})$\ $\longrightarrow$\ $\displaystyle (k^{+},-w^{\prime})U^{\dagger\, b^{\prime}a}(w^{\prime})$
\end{itemize}
\begin{notebox}\textsf{\small Note, l.128:} [Phase: with the conjugated LO bracket at $w^\prime$ and the $g^3$ bracket at $w$ the definition gives $e^{-ik(w-w^\prime)}$ (the original phase is the result at $-k$, not equal for a given $\rho$, $U$); $U^{b^\prime a}\to U^{\dagger\, b^\prime a}=U^{ab^\prime}$, as already used in the next equation. The colour indices $U^{ac}(w)U^{bd}(z)\mf{B}_2^{dc}$ and the range $\Lambda<p^+<\vee-k^+$ agree with the corrected main file. The next two equations contain only the single-$\rho$ part of $\mf{B}_2$; see the note after the last equation.]\end{notebox}
\needspace{3\baselineskip}\noindent\rule{\linewidth}{0.3pt}\par\noindent\textbf{Lines 130} \quad {\small (in \texttt{align}, lines 129--131)}\par
\begin{itemize}[leftmargin=2.2em,itemsep=1pt,topsep=2pt]
\item[\textbf{l.130}] \textsf{\small [phase convention]} $\displaystyle (w^\prime-w)$\ $\longrightarrow$\ $\displaystyle (w-w^\prime)$
\end{itemize}
\needspace{3\baselineskip}\noindent\rule{\linewidth}{0.3pt}\par\noindent\textbf{Lines 135--136} \quad {\small (in \texttt{align}, lines 133--136)}\par
\begin{itemize}[leftmargin=2.2em,itemsep=1pt,topsep=2pt]
\item[\textbf{l.135}] \textsf{\small [phase convention]} $\displaystyle \int_{x,y,z,w,x',w'}\,e^{-ik\cdot (w^\prime-w)}\,$\ $\longrightarrow$\ $\displaystyle \int_{x,y,z,w,x',w'}\,e^{-ik\cdot (w-w^\prime)}\,$
\end{itemize}
\begin{notebox}\textsf{\small Note, l.136:} [Checked: prefactor, colour structure with the order of the $\rho$'s, and the $p^+$ integral over $\Lambda<p^+<\vee-k^+$ (numerically); no $1/\Lambda$ term arises. Not included above: the $\{\rho,\rho\}$ part of $\mf{B}_2$, which gives a term with four $\rho$'s; it is collected in the separate four-$\rho$ file. For that part the two gluons are emitted independently, so its range is $\Lambda<p^+<\vee$ and it carries $\log\frac{\vee}{\Lambda}$.]\end{notebox}
\needspace{3\baselineskip}\noindent\rule{\linewidth}{0.3pt}\par\noindent\textbf{Lines 139--140} \quad \textsf{Group I (with c.c.) $\rangle$ Row V { $\checkmark$]}} \quad {\small (in \texttt{align}, lines 138--140)}\par
\begin{itemize}[leftmargin=2.2em,itemsep=1pt,topsep=2pt]
\item[\textbf{l.139}] \textsf{\small [phase convention]} $\displaystyle \int_{w, w^\prime} e^{-ik(w^\prime-w)}\, \int\, \frac{dp^{+}}{2\pi}$\ $\longrightarrow$\ $\displaystyle \int_{w, w^\prime} e^{-ik(w-w^\prime)}\, \int\, \frac{dp^{+}}{2\pi}$
\item[\textbf{l.139}] \textsf{\small [index/Wilson-line orientation]} $\displaystyle (k^{+},-w^{\prime})U^{b^{\prime}a}(w^{\prime})$\ $\longrightarrow$\ $\displaystyle (k^{+},-w^{\prime})U^{\dagger\, b^{\prime}a}(w^{\prime})$
\end{itemize}
\begin{notebox}\textsf{\small Note, l.140:} [Phase: with the conjugated LO bracket at $w^\prime$ and the $g^{3}$ bracket at $w$ the definition gives $e^{-ik\cdot(w-w^\prime)}$; for given $\rho$, $U$ the original phase is the result at $-k$. After the $w$ integration this reverses every exponent of Row V, as marked below. $U^{b^\prime a}\to U^{\dagger\, b^\prime a}=U^{ab^\prime}$, as already used below. Range: $\Lambda<p^{+}<k^{+}-\Lambda$.]\end{notebox}
\needspace{3\baselineskip}\noindent\rule{\linewidth}{0.3pt}\par\noindent\textbf{Lines 144} \quad {\small (in \texttt{align}, lines 143--145)}\par
\begin{itemize}[leftmargin=2.2em,itemsep=1pt,topsep=2pt]
\item[\textbf{l.144}] \textsf{\small [phase convention]} $\displaystyle \int_{w, w^\prime} e^{-ik(w^\prime-w)}\,$\ $\longrightarrow$\ $\displaystyle \int_{w, w^\prime} e^{-ik(w-w^\prime)}\,$
\end{itemize}
\needspace{3\baselineskip}\noindent\rule{\linewidth}{0.3pt}\par\noindent\textbf{Lines 148} \quad {\small (in \texttt{align}, lines 147--149)}\par
\begin{itemize}[leftmargin=2.2em,itemsep=1pt,topsep=2pt]
\item[\textbf{l.148}] \textsf{\small [phase convention]} $\displaystyle (w^\prime-z)$\ $\longrightarrow$\ $\displaystyle (z-w^\prime)$
\item[\textbf{l.148}] \textsf{\small [phase convention]} $\displaystyle \,e^{-ik\cdot (w^\prime-z)}\,e^{-ik\cdot (z-x)\frac{p^{+}}{k^{+}}}\,\frac{(x'-w')^{i}\,(z-x)^{m}\,(y-z)^{k}\,(y'-x)^{j}}{(x'-w')^{2}\,(z-x)^{2}\,(y-z)^{2}\,(y'-x)^{2}}\,$\ $\longrightarrow$\ $\displaystyle \,e^{-ik\cdot (z-w^\prime)}\,e^{-ik\cdot (x-z)\frac{p^{+}}{k^{+}}}\,\frac{(x'-w')^{i}\,(z-x)^{m}\,(y-z)^{k}\,(y'-x)^{j}}{(x'-w')^{2}\,(z-x)^{2}\,(y-z)^{2}\,(y'-x)^{2}}\,$
\end{itemize}
\needspace{3\baselineskip}\noindent\rule{\linewidth}{0.3pt}\par\noindent\textbf{Lines 152} \quad {\small (in \texttt{align}, lines 151--153)}\par
\begin{itemize}[leftmargin=2.2em,itemsep=1pt,topsep=2pt]
\item[\textbf{l.152}] \textsf{\small [phase convention]} $\displaystyle (w^\prime-z)$\ $\longrightarrow$\ $\displaystyle (z-w^\prime)$
\item[\textbf{l.152}] \textsf{\small [phase convention]} $\displaystyle \,e^{-ik\cdot (w^\prime-z)}\,e^{-ik\cdot (z-x)\frac{p^{+}}{k^{+}}}\,\frac{(x'-w')^{i}\,(z-x)^{m}\,(y-z)^{k}\,(y'-x)^{j}}{(x'-w')^{2}\,(z-x)^{2}\,(y-z)^{2}\,(y'-x)^{2}}\,$\ $\longrightarrow$\ $\displaystyle \,e^{-ik\cdot (z-w^\prime)}\,e^{-ik\cdot (x-z)\frac{p^{+}}{k^{+}}}\,\frac{(x'-w')^{i}\,(z-x)^{m}\,(y-z)^{k}\,(y'-x)^{j}}{(x'-w')^{2}\,(z-x)^{2}\,(y-z)^{2}\,(y'-x)^{2}}\,$
\item[\textbf{l.152}] \textsf{\small [phase convention]} $\displaystyle (w^\prime-x)$\ $\longrightarrow$\ $\displaystyle (x-w^\prime)$
\item[\textbf{l.152}] \textsf{\small [phase convention]} $\displaystyle (x-z)$\ $\longrightarrow$\ $\displaystyle (z-x)$
\item[\textbf{l.152}] \textsf{\small [phase convention]} $\displaystyle (w^\prime-z)$\ $\longrightarrow$\ $\displaystyle (z-w^\prime)$
\item[\textbf{l.152}] \textsf{\small [phase convention]} $\displaystyle (w^\prime-x)$\ $\longrightarrow$\ $\displaystyle (x-w^\prime)$
\end{itemize}
\needspace{3\baselineskip}\noindent\rule{\linewidth}{0.3pt}\par\noindent\textbf{Lines 156} \quad {\small (in \texttt{align}, lines 155--157)}\par
\begin{itemize}[leftmargin=2.2em,itemsep=1pt,topsep=2pt]
\item[\textbf{l.156}] \textsf{\small [phase convention]} $\displaystyle \int_{w, w^\prime} e^{-ik\cdot(w^\prime-w)}\,$\ $\longrightarrow$\ $\displaystyle \int_{w, w^\prime} e^{-ik\cdot(w-w^\prime)}\,$
\end{itemize}
\needspace{3\baselineskip}\noindent\rule{\linewidth}{0.3pt}\par\noindent\textbf{Lines 160} \quad {\small (in \texttt{align}, lines 159--161)}\par
\begin{itemize}[leftmargin=2.2em,itemsep=1pt,topsep=2pt]
\item[\textbf{l.160}] \textsf{\small [phase convention]} $\displaystyle \,\int_{ w^\prime} e^{-ik\cdot(w'-z)}\,e^{-ik\cdot (z-x)\,\frac{p^{+}}{k^{+}}}\, \int_{x,y,z,x'}$\ $\longrightarrow$\ $\displaystyle \,\int_{ w^\prime} e^{-ik\cdot(z-w')}\,e^{-ik\cdot (x-z)\,\frac{p^{+}}{k^{+}}}\, \int_{x,y,z,x'}$
\item[\textbf{l.160}] \textsf{\small [phase convention]} $\displaystyle e^{-ik\cdot(w'-z)}\,e^{-ik\cdot (z-x)\,\frac{p^{+}}{k^{+}}}\, \int_{x,y,z,x'}\,$\ $\longrightarrow$\ $\displaystyle e^{-ik\cdot(z-w')}\,e^{-ik\cdot (x-z)\,\frac{p^{+}}{k^{+}}}\, \int_{x,y,z,x'}\,$
\item[\textbf{l.160}] \textsf{\small [sign/factor correction]} \emph{removed:} $\displaystyle \Bigg\}$\quad{\color{gray}\small(original context: }$\displaystyle ]\,\rho^{b'}(x')\,\rho^{f}(y)\Bigg\}${\color{gray}\small)}
\end{itemize}
\needspace{3\baselineskip}\noindent\rule{\linewidth}{0.3pt}\par\noindent\textbf{Lines 164} \quad {\small (in \texttt{align}, lines 163--165)}\par
\begin{itemize}[leftmargin=2.2em,itemsep=1pt,topsep=2pt]
\item[\textbf{l.164}] \textsf{\small [phase convention]} $\displaystyle \,\int_{x,y,z,x',w'}\,e^{-ik\cdot(w'-z)}\,e^{-ik\cdot (z-x)\,\frac{p^{+}}{k^{+}}}\,$\ $\longrightarrow$\ $\displaystyle \,\int_{x,y,z,x',w'}\,e^{-ik\cdot(z-w')}\,e^{-ik\cdot (x-z)\,\frac{p^{+}}{k^{+}}}\,$
\item[\textbf{l.164}] \textsf{\small [phase convention]} $\displaystyle \,e^{-ik\cdot(w'-z)}\,e^{-ik\cdot (z-x)\,\frac{p^{+}}{k^{+}}}\,$\ $\longrightarrow$\ $\displaystyle \,e^{-ik\cdot(z-w')}\,e^{-ik\cdot (x-z)\,\frac{p^{+}}{k^{+}}}\,$
\end{itemize}
\needspace{3\baselineskip}\noindent\rule{\linewidth}{0.3pt}\par\noindent\textbf{Lines 168} \quad {\small (in \texttt{align}, lines 167--169)}\par
\begin{itemize}[leftmargin=2.2em,itemsep=1pt,topsep=2pt]
\item[\textbf{l.168}] \textsf{\small [phase convention]} $\displaystyle \, \int_{x,y,z,x',w'}\,e^{-ik\cdot(w'-z)}\,e^{-ik\cdot (z-x)\,\frac{p^{+}}{k^{+}}}\,$\ $\longrightarrow$\ $\displaystyle \, \int_{x,y,z,x',w'}\,e^{-ik\cdot(z-w')}\,e^{-ik\cdot (x-z)\,\frac{p^{+}}{k^{+}}}\,$
\item[\textbf{l.168}] \textsf{\small [phase convention]} $\displaystyle \,e^{-ik\cdot(w'-z)}\,e^{-ik\cdot (z-x)\,\frac{p^{+}}{k^{+}}}\,$\ $\longrightarrow$\ $\displaystyle \,e^{-ik\cdot(z-w')}\,e^{-ik\cdot (x-z)\,\frac{p^{+}}{k^{+}}}\,$
\end{itemize}
\needspace{3\baselineskip}\noindent\rule{\linewidth}{0.3pt}\par\noindent\textbf{Lines 173} \quad {\small (in \texttt{align}, lines 172--174)}\par
\begin{itemize}[leftmargin=2.2em,itemsep=1pt,topsep=2pt]
\item[\textbf{l.173}] \textsf{\small [phase convention]} $\displaystyle \, \int_{x,y,z,x',w'}\,e^{-ik\cdot(w'-z)}\,e^{-ik\cdot (z-x)\,\frac{p^{+}}{k^{+}}}\,$\ $\longrightarrow$\ $\displaystyle \, \int_{x,y,z,x',w'}\,e^{-ik\cdot(z-w')}\,e^{-ik\cdot (x-z)\,\frac{p^{+}}{k^{+}}}\,$
\item[\textbf{l.173}] \textsf{\small [phase convention]} $\displaystyle \,e^{-ik\cdot(w'-z)}\,e^{-ik\cdot (z-x)\,\frac{p^{+}}{k^{+}}}\,$\ $\longrightarrow$\ $\displaystyle \,e^{-ik\cdot(z-w')}\,e^{-ik\cdot (x-z)\,\frac{p^{+}}{k^{+}}}\,$
\end{itemize}
\needspace{3\baselineskip}\noindent\rule{\linewidth}{0.3pt}\par\noindent\textbf{Lines 177} \quad {\small (in \texttt{align}, lines 176--178)}\par
\begin{itemize}[leftmargin=2.2em,itemsep=1pt,topsep=2pt]
\item[\textbf{l.177}] \textsf{\small [phase convention]} $\displaystyle \, \int_{x,y,z,x',w'}\,e^{-ik\cdot(w'-z)}\,e^{-ik\cdot (z-x)\,\frac{p^{+}}{k^{+}}}$\ $\longrightarrow$\ $\displaystyle \, \int_{x,y,z,x',w'}\,e^{-ik\cdot(z-w')}\,e^{-ik\cdot (x-z)\,\frac{p^{+}}{k^{+}}}$
\item[\textbf{l.177}] \textsf{\small [phase convention]} $\displaystyle \,e^{-ik\cdot(w'-z)}\,e^{-ik\cdot (z-x)\,\frac{p^{+}}{k^{+}}}$\ $\longrightarrow$\ $\displaystyle \,e^{-ik\cdot(z-w')}\,e^{-ik\cdot (x-z)\,\frac{p^{+}}{k^{+}}}$
\end{itemize}
\needspace{3\baselineskip}\noindent\rule{\linewidth}{0.3pt}\par\noindent\textbf{Lines 180} \quad {\small (in \texttt{align}, lines 179--181)}\par
\begin{itemize}[leftmargin=2.2em,itemsep=1pt,topsep=2pt]
\item[\textbf{l.180}] \textsf{\small [phase convention]} $\displaystyle \, \int_{x,y,z,x',w'}\,e^{-ik\cdot(w'-z)}\,e^{-ik\cdot (z-x)\,\frac{p^{+}}{k^{+}}}\,$\ $\longrightarrow$\ $\displaystyle \, \int_{x,y,z,x',w'}\,e^{-ik\cdot(z-w')}\,e^{-ik\cdot (x-z)\,\frac{p^{+}}{k^{+}}}\,$
\item[\textbf{l.180}] \textsf{\small [phase convention]} $\displaystyle \,e^{-ik\cdot(w'-z)}\,e^{-ik\cdot (z-x)\,\frac{p^{+}}{k^{+}}}\,$\ $\longrightarrow$\ $\displaystyle \,e^{-ik\cdot(z-w')}\,e^{-ik\cdot (x-z)\,\frac{p^{+}}{k^{+}}}\,$
\end{itemize}
\needspace{3\baselineskip}\noindent\rule{\linewidth}{0.3pt}\par\noindent\textbf{Lines 184} \quad {\small (in \texttt{align}, lines 183--185)}\par
\begin{itemize}[leftmargin=2.2em,itemsep=1pt,topsep=2pt]
\item[\textbf{l.184}] \textsf{\small [phase convention]} $\displaystyle \,\int_{x,y,z,x',w'}\, e^{-ik\cdot(w'-z)}\,e^{-ik\cdot (z-x)\,\frac{p^{+}}{k^{+}}}\,$\ $\longrightarrow$\ $\displaystyle \,\int_{x,y,z,x',w'}\, e^{-ik\cdot(z-w')}\,e^{-ik\cdot (x-z)\,\frac{p^{+}}{k^{+}}}\,$
\item[\textbf{l.184}] \textsf{\small [phase convention]} $\displaystyle \, e^{-ik\cdot(w'-z)}\,e^{-ik\cdot (z-x)\,\frac{p^{+}}{k^{+}}}\,$\ $\longrightarrow$\ $\displaystyle \, e^{-ik\cdot(z-w')}\,e^{-ik\cdot (x-z)\,\frac{p^{+}}{k^{+}}}\,$
\item[\textbf{l.184}] \textsf{\small [phase convention]} $\displaystyle \, \int_{x,y,z,x',w'}\,e^{-ik\cdot(w'-x)}\,e^{-ik\cdot (x-z)\,\frac{q^{+}}{k^{+}}}\,$\ $\longrightarrow$\ $\displaystyle \, \int_{x,y,z,x',w'}\,e^{-ik\cdot(x-w')}\,e^{-ik\cdot (z-x)\,\frac{q^{+}}{k^{+}}}\,$
\item[\textbf{l.184}] \textsf{\small [phase convention]} $\displaystyle \,e^{-ik\cdot(w'-x)}\,e^{-ik\cdot (x-z)\,\frac{q^{+}}{k^{+}}}\,$\ $\longrightarrow$\ $\displaystyle \,e^{-ik\cdot(x-w')}\,e^{-ik\cdot (z-x)\,\frac{q^{+}}{k^{+}}}\,$
\item[\textbf{l.184}] \textsf{\small [phase convention]} $\displaystyle \Big)\Big] \int_{x,y,z,x',w'}\,e^{-ik\cdot(w'-z)}\, \frac{(x'-w')^{i}\,(z-x)^{m}}{(x'-w')^{2}\,(z-x)^{2}\, (y-z)^{2}}\,$\ $\longrightarrow$\ $\displaystyle \Big)\Big] \int_{x,y,z,x',w'}\,e^{-ik\cdot(z-w')}\, \frac{(x'-w')^{i}\,(z-x)^{m}}{(x'-w')^{2}\,(z-x)^{2}\, (y-z)^{2}}\,$
\item[\textbf{l.184}] \textsf{\small [phase convention]} $\displaystyle )\Big]\, \int_{x,y,z,x',w'}\,e^{-ik\cdot(w'-x)}\,\frac{(x'-w')^{i}\,(z-x)^{m}}{(x'-w')^{2}\,(z-x)^{2}\,(y-x)^{2}}$\ $\longrightarrow$\ $\displaystyle )\Big]\, \int_{x,y,z,x',w'}\,e^{-ik\cdot(x-w')}\,\frac{(x'-w')^{i}\,(z-x)^{m}}{(x'-w')^{2}\,(z-x)^{2}\,(y-x)^{2}}$
\end{itemize}
\needspace{3\baselineskip}\noindent\rule{\linewidth}{0.3pt}\par\noindent\textbf{Lines 192} \quad {\small (in \texttt{align}, lines 191--194)}\par
\begin{itemize}[leftmargin=2.2em,itemsep=1pt,topsep=2pt]
\item[\textbf{l.192}] \textsf{\small [phase convention]} $\displaystyle \int_{x,y,z,x',w'}e^{-ik\cdot(w'-x)}\,\frac{1}{k^{+}}[$\ $\longrightarrow$\ $\displaystyle \int_{x,y,z,x',w'}e^{-ik\cdot(x-w')}\,\frac{1}{k^{+}}[$
\end{itemize}
\needspace{3\baselineskip}\noindent\rule{\linewidth}{0.3pt}\par\noindent\textbf{Lines 199}\par
\begin{notebox}\textsf{\small Note, l.199:} [Note: this value is the continuation in the loop momentum, $-(2\pi)^{2}\mu^{\epsilon}\int\frac{d^{d}p}{(2\pi)^{d}}$ $e^{ip\cdot(y-x)}/p^{2}$ with $d=2-\epsilon$, which is the convention of the main file (same combination as in $\mf{A}^{(3)}_{2}$); the $z$ integral of the two-dimensional integrand taken literally in $d=2-\epsilon$ gives the same expression with $-\frac{2}{\epsilon}$ in place of $\frac{2}{\epsilon}$. The loop-momentum continuation is the intended one; for the $U(x)$ terms the split ``UV in dimensional regularisation $+$ rest in $d=2$'' is the definition of the scheme.]\end{notebox}
\needspace{3\baselineskip}\noindent\rule{\linewidth}{0.3pt}\par\noindent\textbf{Lines 205} \quad {\small (in \texttt{align}, lines 204--207)}\par
\begin{itemize}[leftmargin=2.2em,itemsep=1pt,topsep=2pt]
\item[\textbf{l.205}] \textsf{\small [phase convention]} $\displaystyle ]\int_{x,y,x',w'}e^{-ik\cdot(w'-x)}$\ $\longrightarrow$\ $\displaystyle ]\int_{x,y,x',w'}e^{-ik\cdot(x-w')}$
\end{itemize}
\needspace{3\baselineskip}\noindent\rule{\linewidth}{0.3pt}\par\noindent\textbf{Lines 210} \quad {\small (in \texttt{align}, lines 209--211)}\par
\begin{itemize}[leftmargin=2.2em,itemsep=1pt,topsep=2pt]
\item[\textbf{l.210}] \textsf{\small [phase convention]} $\displaystyle ]\int_{x,y,x',w'}e^{-ik\cdot(w'-x)}\,[\frac{2}{\epsilon}$\ $\longrightarrow$\ $\displaystyle ]\int_{x,y,x',w'}e^{-ik\cdot(x-w')}\,[\frac{2}{\epsilon}$
\end{itemize}
\needspace{3\baselineskip}\noindent\rule{\linewidth}{0.3pt}\par\noindent\textbf{Lines 214--219} \quad {\small (in \texttt{align}, lines 213--221)}\par
\begin{itemize}[leftmargin=2.2em,itemsep=1pt,topsep=2pt]
\item[\textbf{l.214}] \textsf{\small [phase convention]} $\displaystyle \{e^{-ik\cdot(w'-z)}\,e^{-ik\cdot(z-x)\,\frac{p^{+}}{k^{+}}}\,$\ $\longrightarrow$\ $\displaystyle \{e^{-ik\cdot(z-w')}\,e^{-ik\cdot(x-z)\,\frac{p^{+}}{k^{+}}}\,$
\item[\textbf{l.214}] \textsf{\small [phase convention]} $\displaystyle \{e^{-ik\cdot(w'-z)}\,e^{-ik\cdot(z-x)\,\frac{p^{+}}{k^{+}}}\,$\ $\longrightarrow$\ $\displaystyle \{e^{-ik\cdot(z-w')}\,e^{-ik\cdot(x-z)\,\frac{p^{+}}{k^{+}}}\,$
\item[\textbf{l.219}] \textsf{\small [phase convention]} $\displaystyle -e^{-ik\cdot(w'-x)}\,\frac{1}{k^{+}}[$\ $\longrightarrow$\ $\displaystyle -e^{-ik\cdot(x-w')}\,\frac{1}{k^{+}}[$
\end{itemize}
\needspace{3\baselineskip}\noindent\rule{\linewidth}{0.3pt}\par\noindent\textbf{Lines 224--235} \quad {\small (in \texttt{align}, lines 223--237)}\par
\begin{itemize}[leftmargin=2.2em,itemsep=1pt,topsep=2pt]
\item[\textbf{l.224}] \textsf{\small [phase convention]} $\displaystyle \{e^{-ik\cdot(w'-z)}\,e^{-ik\cdot(z-x)\frac{p^{+}}{k^{+}}}\,$\ $\longrightarrow$\ $\displaystyle \{e^{-ik\cdot(z-w')}\,e^{-ik\cdot(x-z)\frac{p^{+}}{k^{+}}}\,$
\item[\textbf{l.224}] \textsf{\small [phase convention]} $\displaystyle \{e^{-ik\cdot(w'-z)}\,e^{-ik\cdot(z-x)\frac{p^{+}}{k^{+}}}\,$\ $\longrightarrow$\ $\displaystyle \{e^{-ik\cdot(z-w')}\,e^{-ik\cdot(x-z)\frac{p^{+}}{k^{+}}}\,$
\item[\textbf{l.232}] \textsf{\small [phase convention]} $\displaystyle -e^{-ik\cdot(w'-x)}\,\frac{1}{k^{+}}[$\ $\longrightarrow$\ $\displaystyle -e^{-ik\cdot(x-w')}\,\frac{1}{k^{+}}[$
\item[\textbf{l.233}] \textsf{\small [phase convention]} $\displaystyle \,{\color{blue}\log\frac{k^{+}}{\Lambda}}\int_{x,y,z,x',w'}\,e^{-ik\cdot(w'-z)}\,$\ $\longrightarrow$\ $\displaystyle \,{\color{blue}\log\frac{k^{+}}{\Lambda}}\int_{x,y,z,x',w'}\,e^{-ik\cdot(z-w')}\,$
\item[\textbf{l.234}] \textsf{\small [phase convention]} $\displaystyle \,{\color{blue}\log\frac{k^{+}}{\Lambda}}\int_{x,y,z,x',w'}\,e^{-ik\cdot(w'-x)}\,$\ $\longrightarrow$\ $\displaystyle \,{\color{blue}\log\frac{k^{+}}{\Lambda}}\int_{x,y,z,x',w'}\,e^{-ik\cdot(x-w')}\,$
\item[\textbf{l.235}] \textsf{\small [phase convention]} $\displaystyle (w'-x)$\ $\longrightarrow$\ $\displaystyle (x-w')$
\end{itemize}
\needspace{3\baselineskip}\noindent\rule{\linewidth}{0.3pt}\par\noindent\textbf{Lines 241} \quad {\small (in \texttt{align}, lines 240--242)}\par
\begin{itemize}[leftmargin=2.2em,itemsep=1pt,topsep=2pt]
\item[\textbf{l.241}] \textsf{\small [phase convention]} $\displaystyle (w^\prime-z)$\ $\longrightarrow$\ $\displaystyle (z-w^\prime)$
\item[\textbf{l.241}] \textsf{\small [phase convention]} $\displaystyle \,e^{-ik\cdot (w^\prime-z)}\,e^{-ik\cdot (z-x)\frac{p^{+}}{k^{+}}}\,\frac{(x'-w')^{i}\,(z-x)^{m}\,(y-z)^{k}\,(y'-x)^{j}}{(x'-w')^{2}\,(z-x)^{2}\,(y-z)^{2}\,(y'-x)^{2}}\,$\ $\longrightarrow$\ $\displaystyle \,e^{-ik\cdot (z-w^\prime)}\,e^{-ik\cdot (x-z)\frac{p^{+}}{k^{+}}}\,\frac{(x'-w')^{i}\,(z-x)^{m}\,(y-z)^{k}\,(y'-x)^{j}}{(x'-w')^{2}\,(z-x)^{2}\,(y-z)^{2}\,(y'-x)^{2}}\,$
\item[\textbf{l.241}] \textsf{\small [phase convention]} $\displaystyle (w^\prime-x)$\ $\longrightarrow$\ $\displaystyle (x-w^\prime)$
\item[\textbf{l.241}] \textsf{\small [phase convention]} $\displaystyle (x-z)$\ $\longrightarrow$\ $\displaystyle (z-x)$
\item[\textbf{l.241}] \textsf{\small [phase convention]} $\displaystyle (w^\prime-z)$\ $\longrightarrow$\ $\displaystyle (z-w^\prime)$
\item[\textbf{l.241}] \textsf{\small [phase convention]} $\displaystyle (w^\prime-x)$\ $\longrightarrow$\ $\displaystyle (x-w^\prime)$
\end{itemize}
\needspace{3\baselineskip}\noindent\rule{\linewidth}{0.3pt}\par\noindent\textbf{Lines 245} \quad {\small (in \texttt{align}, lines 244--246)}\par
\begin{itemize}[leftmargin=2.2em,itemsep=1pt,topsep=2pt]
\item[\textbf{l.245}] \textsf{\small [phase convention]} $\displaystyle \,\int_{x,y,z,x',w'}\, e^{-ik\cdot(w'-z)}\,e^{-ik\cdot (z-x)\,\frac{p^{+}}{k^{+}}}\,$\ $\longrightarrow$\ $\displaystyle \,\int_{x,y,z,x',w'}\, e^{-ik\cdot(z-w')}\,e^{-ik\cdot (x-z)\,\frac{p^{+}}{k^{+}}}\,$
\item[\textbf{l.245}] \textsf{\small [phase convention]} $\displaystyle \, e^{-ik\cdot(w'-z)}\,e^{-ik\cdot (z-x)\,\frac{p^{+}}{k^{+}}}\,$\ $\longrightarrow$\ $\displaystyle \, e^{-ik\cdot(z-w')}\,e^{-ik\cdot (x-z)\,\frac{p^{+}}{k^{+}}}\,$
\item[\textbf{l.245}] \textsf{\small [phase convention]} $\displaystyle \, \int_{x,y,z,x',w'}\,e^{-ik\cdot(w'-x)}\,e^{-ik\cdot (x-z)\,\frac{q^{+}}{k^{+}}}\,$\ $\longrightarrow$\ $\displaystyle \, \int_{x,y,z,x',w'}\,e^{-ik\cdot(x-w')}\,e^{-ik\cdot (z-x)\,\frac{q^{+}}{k^{+}}}\,$
\item[\textbf{l.245}] \textsf{\small [phase convention]} $\displaystyle \,e^{-ik\cdot(w'-x)}\,e^{-ik\cdot (x-z)\,\frac{q^{+}}{k^{+}}}\,$\ $\longrightarrow$\ $\displaystyle \,e^{-ik\cdot(x-w')}\,e^{-ik\cdot (z-x)\,\frac{q^{+}}{k^{+}}}\,$
\item[\textbf{l.245}] \textsf{\small [phase convention]} $\displaystyle \Big)\Big] \int_{x,y,z,x',w'}\,e^{-ik\cdot(w'-z)}\, \frac{(x'-w')^{i}\,(z-x)^{m}}{(x'-w')^{2}\,(z-x)^{2}\, (y-z)^{2}}\,$\ $\longrightarrow$\ $\displaystyle \Big)\Big] \int_{x,y,z,x',w'}\,e^{-ik\cdot(z-w')}\, \frac{(x'-w')^{i}\,(z-x)^{m}}{(x'-w')^{2}\,(z-x)^{2}\, (y-z)^{2}}\,$
\item[\textbf{l.245}] \textsf{\small [phase convention]} $\displaystyle )\Big]\, \int_{x,y,z,x',w'}\,e^{-ik\cdot(w'-x)}\,\frac{(x'-w')^{i}\,(z-x)^{m}}{(x'-w')^{2}\,(z-x)^{2}\,(y-x)^{2}}$\ $\longrightarrow$\ $\displaystyle )\Big]\, \int_{x,y,z,x',w'}\,e^{-ik\cdot(x-w')}\,\frac{(x'-w')^{i}\,(z-x)^{m}}{(x'-w')^{2}\,(z-x)^{2}\,(y-x)^{2}}$
\end{itemize}
\needspace{3\baselineskip}\noindent\rule{\linewidth}{0.3pt}\par\noindent\textbf{Lines 248} \quad {\small (in \texttt{align}, lines 247--249)}\par
\begin{itemize}[leftmargin=2.2em,itemsep=1pt,topsep=2pt]
\item[\textbf{l.248}] \textsf{\small [phase convention]} $\displaystyle ]\int_{x,y,x',w'}e^{-ik\cdot(w'-x)}\,[\frac{2}{\epsilon}$\ $\longrightarrow$\ $\displaystyle ]\int_{x,y,x',w'}e^{-ik\cdot(x-w')}\,[\frac{2}{\epsilon}$
\end{itemize}
\needspace{3\baselineskip}\noindent\rule{\linewidth}{0.3pt}\par\noindent\textbf{Lines 251--264} \quad {\small (in \texttt{align}, lines 250--264)}\par
\begin{itemize}[leftmargin=2.2em,itemsep=1pt,topsep=2pt]
\item[\textbf{l.251}] \textsf{\small [phase convention]} $\displaystyle \{e^{-ik\cdot(w'-z)}\,e^{-ik\cdot(z-x)\frac{p^{+}}{k^{+}}}\,$\ $\longrightarrow$\ $\displaystyle \{e^{-ik\cdot(z-w')}\,e^{-ik\cdot(x-z)\frac{p^{+}}{k^{+}}}\,$
\item[\textbf{l.251}] \textsf{\small [phase convention]} $\displaystyle \{e^{-ik\cdot(w'-z)}\,e^{-ik\cdot(z-x)\frac{p^{+}}{k^{+}}}\,$\ $\longrightarrow$\ $\displaystyle \{e^{-ik\cdot(z-w')}\,e^{-ik\cdot(x-z)\frac{p^{+}}{k^{+}}}\,$
\item[\textbf{l.259}] \textsf{\small [phase convention]} $\displaystyle -e^{-ik\cdot(w'-x)}\,\frac{1}{k^{+}}[$\ $\longrightarrow$\ $\displaystyle -e^{-ik\cdot(x-w')}\,\frac{1}{k^{+}}[$
\item[\textbf{l.260}] \textsf{\small [phase convention]} $\displaystyle \,{\color{blue}\log\frac{k^{+}}{\Lambda}}\int_{x,y,z,x',w'}\,e^{-ik\cdot(w'-z)}\,$\ $\longrightarrow$\ $\displaystyle \,{\color{blue}\log\frac{k^{+}}{\Lambda}}\int_{x,y,z,x',w'}\,e^{-ik\cdot(z-w')}\,$
\item[\textbf{l.261}] \textsf{\small [phase convention]} $\displaystyle \,{\color{blue}\log\frac{k^{+}}{\Lambda}}\int_{x,y,z,x',w'}\,e^{-ik\cdot(w'-x)}\,$\ $\longrightarrow$\ $\displaystyle \,{\color{blue}\log\frac{k^{+}}{\Lambda}}\int_{x,y,z,x',w'}\,e^{-ik\cdot(x-w')}\,$
\item[\textbf{l.262}] \textsf{\small [phase convention]} $\displaystyle (w'-x)$\ $\longrightarrow$\ $\displaystyle (x-w')$
\end{itemize}
\begin{notebox}\textsf{\small Note, l.264:} [Checked for Row V: starting expressions, contractions, colour factors and order of the $\rho$'s; the separation of the $p^{+}$ integral (plus prescriptions and logarithms reproduce the integral over $\Lambda<p^{+}<k^{+}-\Lambda$ up to ${\cal O}(\Lambda)$, numerically); the coefficient $2\log\frac{k^{+}}{\Lambda}-\frac{11}{6}$; ``UV''+``rest''$\,=2b$. In $2a$ and in ``rest'' $d_{\perp}=2$, and the $z$ integral of ``rest'' at $z\to x$ is understood after the angular average. $\int^{0}_{k^{+}}dq^{+}[k^{+}/q^{+}]_{+}f$ stands for $-\int_{0}^{k^{+}}dq^{+}\frac{k^{+}}{q^{+}}\,(f(q^{+})-f(0))$. The remaining integrals of $2a$ and ``rest'' are not evaluated here.]\end{notebox}
\needspace{15\baselineskip}\bigskip
\subsection{Group II}\label{sec:file-group_II}
\noindent\textbf{Original:} \texttt{/root/.claude/uploads/0ab97fcf-94a2-5aec-8bad-586400481685/08bb31a9-Group\_II.tex}\\\textbf{Checked copy:} \texttt{/home/user/correlations/nlo\_review/group\_II\_checked/Group\_II\_checked.tex}\\
Lines: 277 (original), 277 (checked copy). Coloured groups of the check: 60 (60 top-level); author\textquotesingle s own coloured groups left unchanged: 6.\\
Records by type: phase convention 24; missing term added 1; typo/LaTeX 1; note only 8.
\par\noindent Unmarked layout-only edits (line breaks \texttt{\textbackslash notag\textbackslash\textbackslash\ \&}, \texttt{\textbackslash par\textbackslash noindent} around inserted notes, spacing), checked-copy lines: 35, 40, 53, 54, 57, 61, 85, 86, 92.
\subsubsection*{Changes, block by block}
\needspace{3\baselineskip}\noindent\rule{\linewidth}{0.3pt}\par\noindent\textbf{Lines 35}\par
\begin{itemize}[leftmargin=2.2em,itemsep=1pt,topsep=2pt]
\item[\textbf{l.35}] \textsf{\small [note only]} \emph{added:} \textbf{Marking in this checked copy:} changes and additions are in red;  Line numbers of the source are those of the original file. The phase is $e^{-ik\cdot(w-w^{\prime})}$ throughout (conjugated bracket at $w^{\prime}$, unconjugated bracket at $w$); no c.c. is to be added; four-$\rho$ terms are collected in the separate four-$\rho$ file.
\end{itemize}
\needspace{3\baselineskip}\noindent\rule{\linewidth}{0.3pt}\par\noindent\textbf{Lines 39--40} \quad \textsf{Group II $\rangle$ Row I { $\checkmark$]}} \quad {\small (in \texttt{align}, lines 38--40)}\par
\begin{itemize}[leftmargin=2.2em,itemsep=1pt,topsep=2pt]
\item[\textbf{l.39}] \textsf{\small [phase convention]} $\displaystyle \int_{z,w,w^{\prime}} e^{-ik(w^{\prime}-w)}\,$\ $\longrightarrow$\ $\displaystyle \int_{z,w,w^{\prime}} e^{-ik(w-w^{\prime})}\,$
\end{itemize}
\begin{notebox}\textsf{\small Note, l.40:} [Phase: the conjugated bracket stands at $w^\prime$ and the amplitude at $w$, so the definition gives $e^{-ik(w-w^\prime)}$; the row is Hermitian by itself (no c.c.), but the direction matters ($k\to-k$) because the $\rho$'s do not commute and nothing is averaged. Range: $\Lambda<p^+<\vee-k^+$ holds for every product containing the single-$\rho$ term of $\mf{B}_2$ (first and second part below); the product of the two $\{\rho,\rho\}$ terms has $\Lambda<p^+<\vee$, contains four $\rho$'s and is kept in the separate four-$\rho$ file.]\end{notebox}
\needspace{3\baselineskip}\noindent\rule{\linewidth}{0.3pt}\par\noindent\textbf{Lines 43}\par
\begin{itemize}[leftmargin=2.2em,itemsep=1pt,topsep=2pt]
\item[\textbf{l.43}] \textsf{\small [typo/LaTeX]} \emph{removed:} We write it in two parts:
\end{itemize}
\begin{notebox}\textsf{\small Note, l.43:} [Repeated sentence. First part $=$ interference of the $\{\rho,\rho\}$ term of $\mf{B}_2$ with its single-$\rho$ term, in both orders: the second term is the Hermitian conjugate of the first, so nothing is to be added.]\end{notebox}
\needspace{3\baselineskip}\noindent\rule{\linewidth}{0.3pt}\par\noindent\textbf{Lines 45} \quad {\small (in \texttt{align}, lines 44--46)}\par
\begin{itemize}[leftmargin=2.2em,itemsep=1pt,topsep=2pt]
\item[\textbf{l.45}] \textsf{\small [phase convention]} $\displaystyle (w^{\prime}-w)$\ $\longrightarrow$\ $\displaystyle (w-w^{\prime})$
\item[\textbf{l.45}] \textsf{\small [phase convention]} $\displaystyle (w^{\prime}-w)$\ $\longrightarrow$\ $\displaystyle (w-w^{\prime})$
\end{itemize}
\needspace{3\baselineskip}\noindent\rule{\linewidth}{0.3pt}\par\noindent\textbf{Lines 49} \quad {\small (in \texttt{align}, lines 48--50)}\par
\begin{itemize}[leftmargin=2.2em,itemsep=1pt,topsep=2pt]
\item[\textbf{l.49}] \textsf{\small [phase convention]} $\displaystyle \,\int_{x,y,z,w,x',w^{\prime}} e^{-ik(w^{\prime}-w)}$\ $\longrightarrow$\ $\displaystyle \,\int_{x,y,z,w,x',w^{\prime}} e^{-ik(w-w^{\prime})}$
\item[\textbf{l.49}] \textsf{\small [phase convention]} $\displaystyle \,\int_{x,y,z,w,x',w^{\prime}} e^{-ik(w^{\prime}-w)}$\ $\longrightarrow$\ $\displaystyle \,\int_{x,y,z,w,x',w^{\prime}} e^{-ik(w-w^{\prime})}$
\end{itemize}
\needspace{3\baselineskip}\noindent\rule{\linewidth}{0.3pt}\par\noindent\textbf{Lines 53--54} \quad {\small (in \texttt{align}, lines 52--54)}\par
\begin{itemize}[leftmargin=2.2em,itemsep=1pt,topsep=2pt]
\item[\textbf{l.53}] \textsf{\small [phase convention]} $\displaystyle (w^{\prime}-w)$\ $\longrightarrow$\ $\displaystyle (w-w^{\prime})$
\end{itemize}
\begin{notebox}\textsf{\small Note, l.54:} [Checked against Eq.~(CS2) of the main file with the single-$\rho$ part of $\mf{B}_{2}$: prefactor, range $\Lambda<p^{+}<\vee-k^{+}$, tensors and colour factors agree. Phase: with the conjugated bracket (at $w^{\prime}$, $x^{\prime}$) standing left and the other at $w$, the definition gives $e^{-ik(w-w^{\prime})}$; this term is hermitian by itself, but the two phases differ by $k\to-k$, which is not a symmetry for a fixed target, so the exponent is changed here and below.]\end{notebox}
\needspace{3\baselineskip}\noindent\rule{\linewidth}{0.3pt}\par\noindent\textbf{Lines 57} \quad {\small (in \texttt{align}, lines 56--58)}\par
\begin{itemize}[leftmargin=2.2em,itemsep=1pt,topsep=2pt]
\item[\textbf{l.57}] \textsf{\small [phase convention]} $\displaystyle (w^{\prime}-w)$\ $\longrightarrow$\ $\displaystyle (w-w^{\prime})$
\end{itemize}
\needspace{3\baselineskip}\noindent\rule{\linewidth}{0.3pt}\par\noindent\textbf{Lines 61} \quad {\small (in \texttt{align}, lines 60--62)}\par
\begin{itemize}[leftmargin=2.2em,itemsep=1pt,topsep=2pt]
\item[\textbf{l.61}] \textsf{\small [phase convention]} $\displaystyle (w^{\prime}-w)$\ $\longrightarrow$\ $\displaystyle (w-w^{\prime})$
\end{itemize}
\needspace{3\baselineskip}\noindent\rule{\linewidth}{0.3pt}\par\noindent\textbf{Lines 65} \quad {\small (in \texttt{align}, lines 64--66)}\par
\begin{itemize}[leftmargin=2.2em,itemsep=1pt,topsep=2pt]
\item[\textbf{l.65}] \textsf{\small [phase convention]} $\displaystyle \,\int_{z,w,w^{\prime},x,x'}\, e^{-ik(w^{\prime}-w)}\,$\ $\longrightarrow$\ $\displaystyle \,\int_{z,w,w^{\prime},x,x'}\, e^{-ik(w-w^{\prime})}\,$
\end{itemize}
\needspace{3\baselineskip}\noindent\rule{\linewidth}{0.3pt}\par\noindent\textbf{Lines 69} \quad {\small (in \texttt{align}, lines 68--70)}\par
\begin{itemize}[leftmargin=2.2em,itemsep=1pt,topsep=2pt]
\item[\textbf{l.69}] \textsf{\small [phase convention]} $\displaystyle )\Big]\,\int_{z,w,w^{\prime},x,x'}\, e^{-ik(w^{\prime}-w)}\,$\ $\longrightarrow$\ $\displaystyle )\Big]\,\int_{z,w,w^{\prime},x,x'}\, e^{-ik(w-w^{\prime})}\,$
\item[\textbf{l.69}] \textsf{\small [phase convention]} $\displaystyle \,\int_{z,w,w^{\prime},x,x'}\, e^{-ik(w^{\prime}-w)}\, \frac{(x-z)^{2}}{[(x'-w')^{2}\,(x-z)^{2}\,(x-w)^{2} - (x-w)^{4}\, (x'-z)^{2}]}\,$\ $\longrightarrow$\ $\displaystyle \,\int_{z,w,w^{\prime},x,x'}\, e^{-ik(w-w^{\prime})}\, \frac{(x-z)^{2}}{[(x'-w')^{2}\,(x-z)^{2}\,(x-w)^{2} - (x-w)^{4}\, (x'-z)^{2}]}\,$
\item[\textbf{l.69}] \textsf{\small [phase convention]} $\displaystyle \,\int_{z,w,w^{\prime},x,x'}\, e^{-ik(w^{\prime}-w)}\, \frac{(x'-z)^{2}}{[(x'-w')^{4}\,(x-z)^{2}\, - (x-w)^{2}\, (x'-z)^{2}\,(x'-w')^{2}]}\,$\ $\longrightarrow$\ $\displaystyle \,\int_{z,w,w^{\prime},x,x'}\, e^{-ik(w-w^{\prime})}\, \frac{(x'-z)^{2}}{[(x'-w')^{4}\,(x-z)^{2}\, - (x-w)^{2}\, (x'-z)^{2}\,(x'-w')^{2}]}\,$
\end{itemize}
\needspace{3\baselineskip}\noindent\rule{\linewidth}{0.3pt}\par\noindent\textbf{Lines 73} \quad {\small (in \texttt{align}, lines 72--74)}\par
\begin{itemize}[leftmargin=2.2em,itemsep=1pt,topsep=2pt]
\item[\textbf{l.73}] \textsf{\small [phase convention]} $\displaystyle (w^{\prime}-w)$\ $\longrightarrow$\ $\displaystyle (w-w^{\prime})$
\item[\textbf{l.73}] \textsf{\small [phase convention]} $\displaystyle (w^{\prime}-w)$\ $\longrightarrow$\ $\displaystyle (w-w^{\prime})$
\end{itemize}
\needspace{3\baselineskip}\noindent\rule{\linewidth}{0.3pt}\par\noindent\textbf{Lines 77} \quad {\small (in \texttt{align}, lines 76--78)}\par
\begin{itemize}[leftmargin=2.2em,itemsep=1pt,topsep=2pt]
\item[\textbf{l.77}] \textsf{\small [phase convention]} $\displaystyle (w^{\prime}-w)$\ $\longrightarrow$\ $\displaystyle (w-w^{\prime})$
\item[\textbf{l.77}] \textsf{\small [phase convention]} $\displaystyle (w^{\prime}-w)$\ $\longrightarrow$\ $\displaystyle (w-w^{\prime})$
\item[\textbf{l.77}] \textsf{\small [phase convention]} $\displaystyle (w^{\prime}-w)$\ $\longrightarrow$\ $\displaystyle (w-w^{\prime})$
\end{itemize}
\needspace{3\baselineskip}\noindent\rule{\linewidth}{0.3pt}\par\noindent\textbf{Lines 81} \quad {\small (in \texttt{align}, lines 80--82)}\par
\begin{itemize}[leftmargin=2.2em,itemsep=1pt,topsep=2pt]
\item[\textbf{l.81}] \textsf{\small [phase convention]} $\displaystyle (w^{\prime}-w)$\ $\longrightarrow$\ $\displaystyle (w-w^{\prime})$
\item[\textbf{l.81}] \textsf{\small [phase convention]} $\displaystyle (w^{\prime}-w)$\ $\longrightarrow$\ $\displaystyle (w-w^{\prime})$
\item[\textbf{l.81}] \textsf{\small [phase convention]} $\displaystyle (w^{\prime}-w)$\ $\longrightarrow$\ $\displaystyle (w-w^{\prime})$
\end{itemize}
\needspace{3\baselineskip}\noindent\rule{\linewidth}{0.3pt}\par\noindent\textbf{Lines 85--88} \quad {\small (in \texttt{align}, lines 84--86)}\par
\begin{itemize}[leftmargin=2.2em,itemsep=1pt,topsep=2pt]
\item[\textbf{l.85}] \textsf{\small [phase convention]} $\displaystyle \,\int_{z,w,w^{\prime},x,x'} e^{-ik(w^{\prime}-w)}\,$\ $\longrightarrow$\ $\displaystyle \,\int_{z,w,w^{\prime},x,x'} e^{-ik(w-w^{\prime})}\,$
\item[\textbf{l.85}] \textsf{\small [phase convention]} $\displaystyle \,\int_{z,w,w^{\prime},x,x'} e^{-ik(w^{\prime}-w)}\,\frac{(x-z)^{2}\, (x-w)^{2}}{(z-w)^{2}\, (z-w')^{2}\,[(x-z)^{2} - (x-w)^{2}]}\,$\ $\longrightarrow$\ $\displaystyle \,\int_{z,w,w^{\prime},x,x'} e^{-ik(w-w^{\prime})}\,\frac{(x-z)^{2}\, (x-w)^{2}}{(z-w)^{2}\, (z-w')^{2}\,[(x-z)^{2} - (x-w)^{2}]}\,$
\item[\textbf{l.85}] \textsf{\small [phase convention]} $\displaystyle ]\,\int_{z,w,w^{\prime},x,x'} e^{-ik(w^{\prime}-w)}\,\frac{1}{(z-w)^{2}\, (z-w')^{2}}$\ $\longrightarrow$\ $\displaystyle ]\,\int_{z,w,w^{\prime},x,x'} e^{-ik(w-w^{\prime})}\,\frac{1}{(z-w)^{2}\, (z-w')^{2}}$
\item[\textbf{l.85}] \textsf{\small [missing term added]} \emph{added:} $\displaystyle U$\quad{\color{gray}\small(in context: }$\displaystyle f^{ac'd'}\,f^{cbd}\,U^{ac}(w)\,U^{c'b}(z)\,U^{d'e'}(x')\,\rho^{e'}(x${\color{gray}\small)}
\item[\textbf{l.85}] \textsf{\small [phase convention]} $\displaystyle \,\int_{z,w,w^{\prime},x,x'} e^{-ik(w^{\prime}-w)}\,\frac{(x'-w')^{2}\, (x'-z)^{2}}{(z-w)^{2}\, (z-w')^{2}}$\ $\longrightarrow$\ $\displaystyle \,\int_{z,w,w^{\prime},x,x'} e^{-ik(w-w^{\prime})}\,\frac{(x'-w')^{2}\, (x'-z)^{2}}{(z-w)^{2}\, (z-w')^{2}}$
\end{itemize}
\begin{notebox}\textsf{\small Note, l.86:} [Checked: $2a+\dots+2f$ equals the $p^{+}$ integral of the previous expression exactly (no expansion in $\Lambda/k^{+}$ or $k^{+}/\vee$). The separate terms inside each of $2a$--$2f$ are singular at $(x^{\prime}-w^{\prime})^{2}(x-z)^{2}=(x-w)^{2}(x^{\prime}-z)^{2}$, and those of $2d$, $2e$, $2f$ also at $(x-z)^{2}=(x-w)^{2}$ and $(x^{\prime}-z)^{2}=(x^{\prime}-w^{\prime})^{2}$; these poles cancel within each of the six pieces, so the terms of one piece must be kept under one integral.]\end{notebox}
\begin{notebox}\textsf{\small Note, l.88:} [Checked: this row is $\frac{4}{(2\pi)^3}\cdot\frac14$ times a product of four coefficients $\mf{A}^{(1)}$, each linear in $\rho$, so every term has exactly four $\rho$'s before any reordering; $\Lambda<p^+<\vee$.]\end{notebox}
\needspace{3\baselineskip}\noindent\rule{\linewidth}{0.3pt}\par\noindent\textbf{Lines 90} \quad \textsf{Group II $\rangle$ Row III { $\checkmark$]}}\par
\begin{notebox}\textsf{\small Note, l.90:} [Checked: this row is $\frac{4}{(2\pi)^3}\cdot\frac14$ times a product of four coefficients $\mf{A}^{(1)}$, each linear in $\rho$, so every term has exactly four $\rho$'s before any reordering; $\Lambda<p^+<\vee$.]\end{notebox}
\needspace{3\baselineskip}\noindent\rule{\linewidth}{0.3pt}\par\noindent\textbf{Lines 92} \quad \textsf{Group II $\rangle$ Row IV { $\checkmark$]}}\par
\begin{notebox}\textsf{\small Note, l.92:} [This section is an older version of Row IV. The checked version is the separate file Group\_II\_Row\_IV\_checked; the same corrections apply here: phase $e^{-ik\cdot(w-w')}$, $\int_{k^+/\vee}^{1}d\xi$ with $\log\frac{k^+}{\Lambda}+\log(1-\frac{k^+}{\vee})$ in place of $\int_0^1 d\xi$ with $\log\frac{k^+}{\Lambda}$, the $\frac{2}{\epsilon}$ poles kept, and $\mathcal{P}_{gg}$ with its factor $2N_c$. It has not been re-marked line by line.]\end{notebox}
\needspace{15\baselineskip}\bigskip
\subsection{Group II Row IV}\label{sec:file-group_II_row_IV}
\noindent\textbf{Original:} \texttt{/root/.claude/uploads/0ab97fcf-94a2-5aec-8bad-586400481685/c62bee3b-Group\_II\_Row\_IV.tex}\\\textbf{Checked copy:} \texttt{/home/user/correlations/nlo\_review/group\_II\_checked/Group\_II\_Row\_IV\_checked.tex}\\
Lines: 204 (original), 204 (checked copy). Coloured groups of the check: 153 (153 top-level); author\textquotesingle s own coloured groups left unchanged: 0.\\
Records by type: phase convention 39; limits/range 17; missing term added 11; typo/LaTeX 18; note only 10.
\par\noindent Unmarked layout-only edits (line breaks \texttt{\textbackslash notag\textbackslash\textbackslash\ \&}, \texttt{\textbackslash par\textbackslash noindent} around inserted notes, spacing), checked-copy lines: 35, 38, 47, 49, 53, 55, 122, 138, 146, 150.
\subsubsection*{Changes, block by block}
\needspace{3\baselineskip}\noindent\rule{\linewidth}{0.3pt}\par\noindent\textbf{Lines 35--36}\par
\begin{itemize}[leftmargin=2.2em,itemsep=1pt,topsep=2pt]
\item[\textbf{l.35}] \textsf{\small [note only]} \emph{added:} \textbf{Marking in this checked copy:} changes and additions are in red;  Line numbers of the source are those of the original file. The phase is $e^{-ik\cdot(w-w^{\prime})}$ throughout (conjugated bracket at $w^{\prime}$, unconjugated bracket at $w$); no c.c. is to be added; four-$\rho$ terms are collected in the separate four-$\rho$ file.
\end{itemize}
\begin{notebox}\textsf{\small Note, l.36:} [Checked against the corrected Eq.~(CS2) of the main file: factor $\frac{4}{(2\pi)^{3}}\int\frac{dp^{+}}{2\pi}$, $U^{db}(z)=U^{\dagger\,bd}(z)$, range $\Lambda<p^{+}<\vee-k^{+}$ (the parent gluon is a soft mode). Phase: the conjugated bracket is at $w'$ and stands left, so the definition gives $e^{-ik\cdot(w-w')}$. For this diagonal term the two phases differ only by the commutator $[\rho^{e'}(x'),\rho^{e}(x)]$ (a single-$\rho$ term, odd in $k$); the $\{\rho,\rho\}$ part is even in $k$. The phase is corrected in all equations below; the result is Hermitian. No four-$\rho$ term occurs in this row.]\end{notebox}
\needspace{3\baselineskip}\noindent\rule{\linewidth}{0.3pt}\par\noindent\textbf{Lines 38} \quad {\small (in \texttt{align}, lines 37--39)}\par
\begin{itemize}[leftmargin=2.2em,itemsep=1pt,topsep=2pt]
\item[\textbf{l.38}] \textsf{\small [phase convention]} $\displaystyle \int_{z,w,w^{\prime}} e^{-ik(w^{\prime}-w)}\,$\ $\longrightarrow$\ $\displaystyle \int_{z,w,w^{\prime}} e^{-ik(w-w^{\prime})}\,$
\item[\textbf{l.38}] \textsf{\small [typo/LaTeX]} \emph{added:} $\displaystyle ;$\quad{\color{gray}\small(in context: }$\displaystyle w^{\prime};\, p^{+},w^{\prime}${\color{gray}\small)}
\item[\textbf{l.38}] \textsf{\small [typo/LaTeX]} \emph{added:} $\displaystyle ;$\quad{\color{gray}\small(in context: }$\displaystyle w^{\prime};\, p^{+},w^{\prime}${\color{gray}\small)}
\item[\textbf{l.38}] \textsf{\small [typo/LaTeX]} \emph{added:} $\displaystyle ;$\quad{\color{gray}\small(in context: }$\displaystyle k^{+},y-w;\, p^{+},w${\color{gray}\small)}
\item[\textbf{l.38}] \textsf{\small [typo/LaTeX]} \emph{added:} $\displaystyle ;$\quad{\color{gray}\small(in context: }$\displaystyle k^{+},y-w;\, p^{+},w${\color{gray}\small)}
\end{itemize}
\needspace{3\baselineskip}\noindent\rule{\linewidth}{0.3pt}\par\noindent\textbf{Lines 41} \quad {\small (in \texttt{align}, lines 40--44)}\par
\begin{itemize}[leftmargin=2.2em,itemsep=1pt,topsep=2pt]
\item[\textbf{l.41}] \textsf{\small [phase convention]} $\displaystyle (w^{\prime}-w)$\ $\longrightarrow$\ $\displaystyle (w-w^{\prime})$
\end{itemize}
\needspace{3\baselineskip}\noindent\rule{\linewidth}{0.3pt}\par\noindent\textbf{Lines 47--49} \quad {\small (in \texttt{align}, lines 46--50)}\par
\begin{itemize}[leftmargin=2.2em,itemsep=1pt,topsep=2pt]
\item[\textbf{l.47}] \textsf{\small [typo/LaTeX]} \emph{removed:} $\displaystyle ,$\quad{\color{gray}\small(original context: }$\displaystyle \frac{1}{(2\pi)^{3}},\,\frac{g^{4}}{4\pi^{4}}${\color{gray}\small)}
\item[\textbf{l.47}] \textsf{\small [typo/LaTeX]} \emph{removed:} $\displaystyle ,$\quad{\color{gray}\small(original context: }$\displaystyle \frac{1}{(2\pi)^{3}},\,\frac{g^{4}}{4\pi^{4}},\,\frac{1}{k^{+}}\int_{\frac{k^{+}}{\vee}}^{\frac{k^{+}}{k^{+}+\Lambda}}\,\frac{d\xi}{2\pi}${\color{gray}\small)}
\item[\textbf{l.47}] \textsf{\small [phase convention]} $\displaystyle (w’-w)$\ $\longrightarrow$\ $\displaystyle (w-w')$
\item[\textbf{l.49}] \textsf{\small [typo/LaTeX]} \emph{removed:} $\displaystyle ,$\quad{\color{gray}\small(original context: }$\displaystyle \frac{1}{1-\xi},\,\delta_{jm}\delta_{ik}${\color{gray}\small)}
\item[\textbf{l.49}] \textsf{\small [typo/LaTeX]} \emph{removed:} $\displaystyle ,$\quad{\color{gray}\small(original context: }$\displaystyle \frac{1}{\xi},\,\delta_{im}\delta_{kj}\Big]${\color{gray}\small)}
\end{itemize}
\needspace{3\baselineskip}\noindent\rule{\linewidth}{0.3pt}\par\noindent\textbf{Lines 53--55} \quad {\small (in \texttt{align}, lines 52--56)}\par
\begin{itemize}[leftmargin=2.2em,itemsep=1pt,topsep=2pt]
\item[\textbf{l.53}] \textsf{\small [typo/LaTeX]} \emph{removed:} $\displaystyle ,$\quad{\color{gray}\small(original context: }$\displaystyle \frac{1}{(2\pi)^{3}},\,\frac{g^{4}\,N_{c}}{4\pi^{4}}${\color{gray}\small)}
\item[\textbf{l.53}] \textsf{\small [typo/LaTeX]} \emph{removed:} $\displaystyle ,$\quad{\color{gray}\small(original context: }$\displaystyle \frac{1}{(2\pi)^{3}},\,\frac{g^{4}\,N_{c}}{4\pi^{4}},\,\frac{1}{k^{+}}\int_{\frac{k^{+}}{\vee}}^{\frac{k^{+}}{k^{+}+\Lambda}}\,\frac{d\xi}{2\pi}${\color{gray}\small)}
\item[\textbf{l.53}] \textsf{\small [phase convention]} $\displaystyle (w’-w)$\ $\longrightarrow$\ $\displaystyle (w-w')$
\item[\textbf{l.55}] \textsf{\small [typo/LaTeX]} \emph{removed:} $\displaystyle ,$\quad{\color{gray}\small(original context: }$\displaystyle \frac{1}{1-\xi},\,\delta_{jm}\delta_{ik}${\color{gray}\small)}
\item[\textbf{l.55}] \textsf{\small [typo/LaTeX]} \emph{removed:} $\displaystyle ,$\quad{\color{gray}\small(original context: }$\displaystyle \frac{1}{\xi},\,\delta_{im}\delta_{kj}\Big]${\color{gray}\small)}
\end{itemize}
\needspace{3\baselineskip}\noindent\rule{\linewidth}{0.3pt}\par\noindent\textbf{Lines 59} \quad {\small (in \texttt{align}, lines 58--60)}\par
\begin{itemize}[leftmargin=2.2em,itemsep=1pt,topsep=2pt]
\item[\textbf{l.59}] \textsf{\small [phase convention]} $\displaystyle \int_{z,w,w',x'}e^{-ik(w'-w)}\int_{x,y,y'}\frac{(w-z)^{m}(x-y)^{k}(w'-z)^{m'}(x'-y')^{k'}}{(w-z)^{2}(x-y)^{2}(w'-z)^{2}(x'-y')^{2}}$\ $\longrightarrow$\ $\displaystyle \int_{z,w,w',x'}e^{-ik(w-w')}\int_{x,y,y'}\frac{(w-z)^{m}(x-y)^{k}(w'-z)^{m'}(x'-y')^{k'}}{(w-z)^{2}(x-y)^{2}(w'-z)^{2}(x'-y')^{2}}$
\end{itemize}
\needspace{3\baselineskip}\noindent\rule{\linewidth}{0.3pt}\par\noindent\textbf{Lines 63} \quad {\small (in \texttt{align}, lines 62--66)}\par
\begin{itemize}[leftmargin=2.2em,itemsep=1pt,topsep=2pt]
\item[\textbf{l.63}] \textsf{\small [phase convention]} $\displaystyle \int_{z,w,w',x,x'}e^{-ik(w'-w)}\,\frac{(w-z)^{m}(x-y)^{k}(w'-z)^{m'}(x'-y')^{k'}}{(w-z)^{2}(x-y)^{2}(w'-z)^{2}(x'-y')^{2}}$\ $\longrightarrow$\ $\displaystyle \int_{z,w,w',x,x'}e^{-ik(w-w')}\,\frac{(w-z)^{m}(x-y)^{k}(w'-z)^{m'}(x'-y')^{k'}}{(w-z)^{2}(x-y)^{2}(w'-z)^{2}(x'-y')^{2}}$
\end{itemize}
\needspace{3\baselineskip}\noindent\rule{\linewidth}{0.3pt}\par\noindent\textbf{Lines 69--73} \quad {\small (in \texttt{align}, lines 68--72)}\par
\begin{itemize}[leftmargin=2.2em,itemsep=1pt,topsep=2pt]
\item[\textbf{l.69}] \textsf{\small [phase convention]} $\displaystyle -$\ $\longrightarrow$\ $\displaystyle +$
\end{itemize}
\begin{notebox}\textsf{\small Note, l.73:} [Checked: this is the continuation in the momentum conjugate to $r$, $(2\pi)^{2}\mu^{\epsilon}\int\frac{d^{d}q}{(2\pi)^{d}}\frac{q^{m}q^{m'}}{q^{4}}e^{iq\cdot u}$, $d=2-\epsilon$; the divergence is at $q\to0$ (large $r$, collinear), so the integral exists for $\epsilon<0$ and the pole term $-\pi\delta^{mm'}/\epsilon$ is positive there.]\end{notebox}
\needspace{3\baselineskip}\noindent\rule{\linewidth}{0.3pt}\par\noindent\textbf{Lines 81--86} \quad {\small (in \texttt{align}, lines 80--85)}\par
\begin{itemize}[leftmargin=2.2em,itemsep=1pt,topsep=2pt]
\item[\textbf{l.81}] \textsf{\small [phase convention]} $\displaystyle \int_{u,Y,x,x'}e^{-ik\cdot u}\,$\ $\longrightarrow$\ $\displaystyle \int_{u,Y,x,x'}e^{+ik\cdot u}\,$
\end{itemize}
\begin{notebox}\textsf{\small Note, l.86:} [Checked. The result holds with $d_{\perp}=\delta_{mm}=2-\epsilon$ kept in the contraction: the term $-\frac12\frac{(x-y)\cdot(x'-y')}{\dots}$ in the last bracket comes from $-\epsilon\times\frac{1}{\epsilon}$ together with $\frac{u^{m}u^{m'}}{u^{2}}\delta_{mm'}\delta_{kk'}$. With $d_{\perp}=2$ the bracket $[\frac{2}{\epsilon}+\dots]$ would be replaced by $[\frac{2}{\epsilon}+\dots+1]$.]\end{notebox}
\needspace{3\baselineskip}\noindent\rule{\linewidth}{0.3pt}\par\noindent\textbf{Lines 88--93} \quad {\small (in \texttt{align}, lines 87--92)}\par
\begin{itemize}[leftmargin=2.2em,itemsep=1pt,topsep=2pt]
\item[\textbf{l.88}] \textsf{\small [phase convention]} $\displaystyle d\xi\int_{u,Y,x,x'}e^{-ik\cdot u}\,$\ $\longrightarrow$\ $\displaystyle d\xi\int_{u,Y,x,x'}e^{+ik\cdot u}\,$
\end{itemize}
\begin{notebox}\textsf{\small Note, l.93:} [$\int_{k^{+}/\vee}^{k^{+}/(k^{+}+\Lambda)}\frac{d\xi}{1-\xi}=\log\frac{k^{+}}{\Lambda}+\log(1-\frac{k^{+}}{\vee})+O(\frac{\Lambda}{k^{+}})$, $\int\frac{d\xi}{\xi}=\log\frac{\vee}{k^{+}}+O(\frac{\Lambda}{k^{+}})$. The subtracted integrals keep the lower limit $\frac{k^{+}}{\vee}$ (only the upper limit can be set to 1, at the cost of $O(\Lambda/k^{+})$); with $\int_{0}^{1}$ the equations of this section would differ from the starting expression by terms of order $k^{+}/\vee$. Corrected from here to the first form of the DGLAP term.]\end{notebox}
\needspace{3\baselineskip}\noindent\rule{\linewidth}{0.3pt}\par\noindent\textbf{Lines 95--101} \quad {\small (in \texttt{align}, lines 94--103)}\par
\begin{itemize}[leftmargin=2.2em,itemsep=1pt,topsep=2pt]
\item[\textbf{l.95}] \textsf{\small [phase convention]} $\displaystyle \int_{u,Y,x,x'}e^{-ik\cdot u}\Big[\frac{2}{\epsilon}$\ $\longrightarrow$\ $\displaystyle \int_{u,Y,x,x'}e^{+ik\cdot u}\Big[\frac{2}{\epsilon}$
\item[\textbf{l.96}] \textsf{\small [typo/LaTeX]} \emph{added:} $\displaystyle \Big[$\quad{\color{gray}\small(in context: }$\displaystyle \{\Big[\log\frac{k^{+}}{\Lambda}${\color{gray}\small)}
\item[\textbf{l.96}] \textsf{\small [missing term added]} \emph{added:} $\displaystyle +\log\Big(1-\frac{k^{+}}{\vee}\Big)\Big]$\quad{\color{gray}\small(in context: }$\displaystyle \{\Big[\log\frac{k^{+}}{\Lambda}+\log\Big(1-\frac{k^{+}}{\vee}\Big)\Big]\,\frac{(x-y)\cdot(x'-y')}{(x-y)^{2}(x'-y')^{2}}[U^{ce'}(y')${\color{gray}\small)}
\item[\textbf{l.98}] \textsf{\small [limits/range]} $\displaystyle +\int_{0}^{1}d\xi\,$\ $\longrightarrow$\ $\displaystyle +\int_{\frac{k^{+}}{\vee}}^{1}d\xi\,$
\item[\textbf{l.101}] \textsf{\small [limits/range]} $\displaystyle \,\frac{1}{k^{+}}\int_{0}^{1}d\xi\int_{u,Y,x,x'}$\ $\longrightarrow$\ $\displaystyle \,\frac{1}{k^{+}}\int_{\frac{k^{+}}{\vee}}^{1}d\xi\int_{u,Y,x,x'}$
\item[\textbf{l.101}] \textsf{\small [phase convention]} $\displaystyle d\xi\int_{u,Y,x,x'}e^{-ik\cdot u}\,2\,\xi(1-\xi)$\ $\longrightarrow$\ $\displaystyle d\xi\int_{u,Y,x,x'}e^{+ik\cdot u}\,2\,\xi(1-\xi)$
\end{itemize}
\needspace{3\baselineskip}\noindent\rule{\linewidth}{0.3pt}\par\noindent\textbf{Lines 106--113} \quad {\small (in \texttt{align}, lines 105--115)}\par
\begin{itemize}[leftmargin=2.2em,itemsep=1pt,topsep=2pt]
\item[\textbf{l.106}] \textsf{\small [typo/LaTeX]} \emph{added:} $\displaystyle \Big[$\quad{\color{gray}\small(in context: }$\displaystyle \{\Big[\log\frac{k^{+}}{\Lambda}${\color{gray}\small)}
\item[\textbf{l.106}] \textsf{\small [missing term added]} \emph{added:} $\displaystyle +\log\Big(1-\frac{k^{+}}{\vee}\Big)\Big]$\quad{\color{gray}\small(in context: }$\displaystyle \{\Big[\log\frac{k^{+}}{\Lambda}+\log\Big(1-\frac{k^{+}}{\vee}\Big)\Big]\int_{u,Y,x,x'}e^{+ik\cdot u}\Big[\frac{2}{\epsilon}${\color{gray}\small)}
\item[\textbf{l.106}] \textsf{\small [phase convention]} $\displaystyle \int_{u,Y,x,x'}e^{-ik\cdot u}\Big[\frac{2}{\epsilon}$\ $\longrightarrow$\ $\displaystyle \int_{u,Y,x,x'}e^{+ik\cdot u}\Big[\frac{2}{\epsilon}$
\item[\textbf{l.108}] \textsf{\small [limits/range]} $\displaystyle +\int_{0}^{1}d\xi\,$\ $\longrightarrow$\ $\displaystyle +\int_{\frac{k^{+}}{\vee}}^{1}d\xi\,$
\item[\textbf{l.108}] \textsf{\small [phase convention]} $\displaystyle [\int_{u,Y,x,x'}e^{-ik\cdot u}\Big[\frac{2}{\epsilon}$\ $\longrightarrow$\ $\displaystyle [\int_{u,Y,x,x'}e^{+ik\cdot u}\Big[\frac{2}{\epsilon}$
\item[\textbf{l.113}] \textsf{\small [limits/range]} $\displaystyle \,\frac{1}{k^{+}}\int_{0}^{1}d\xi\int_{u,Y,x,x'}$\ $\longrightarrow$\ $\displaystyle \,\frac{1}{k^{+}}\int_{\frac{k^{+}}{\vee}}^{1}d\xi\int_{u,Y,x,x'}$
\item[\textbf{l.113}] \textsf{\small [phase convention]} $\displaystyle d\xi\int_{u,Y,x,x'}e^{-ik\cdot u}\,2\,\xi(1-\xi)$\ $\longrightarrow$\ $\displaystyle d\xi\int_{u,Y,x,x'}e^{+ik\cdot u}\,2\,\xi(1-\xi)$
\end{itemize}
\needspace{3\baselineskip}\noindent\rule{\linewidth}{0.3pt}\par\noindent\textbf{Lines 118--125} \quad {\small (in \texttt{align}, lines 117--127)}\par
\begin{itemize}[leftmargin=2.2em,itemsep=1pt,topsep=2pt]
\item[\textbf{l.118}] \textsf{\small [typo/LaTeX]} \emph{added:} $\displaystyle \Big[$\quad{\color{gray}\small(in context: }$\displaystyle \{\Big[\log\frac{k^{+}}{\Lambda}${\color{gray}\small)}
\item[\textbf{l.118}] \textsf{\small [missing term added]} \emph{added:} $\displaystyle +\log\Big(1-\frac{k^{+}}{\vee}\Big)\Big]$\quad{\color{gray}\small(in context: }$\displaystyle \{\Big[\log\frac{k^{+}}{\Lambda}+\log\Big(1-\frac{k^{+}}{\vee}\Big)\Big]\int_{x,y,x',y'}e^{-ik\cdot(y-y')}\Big[\frac{2}{\epsilon}${\color{gray}\small)}
\item[\textbf{l.118}] \textsf{\small [phase convention]} $\displaystyle \int_{x,y,x',y'}e^{-ik\cdot(y'-y)}\Big[\frac{2}{\epsilon}$\ $\longrightarrow$\ $\displaystyle \int_{x,y,x',y'}e^{-ik\cdot(y-y')}\Big[\frac{2}{\epsilon}$
\item[\textbf{l.120}] \textsf{\small [limits/range]} $\displaystyle +\int_{0}^{1}d\xi\,$\ $\longrightarrow$\ $\displaystyle +\int_{\frac{k^{+}}{\vee}}^{1}d\xi\,$
\item[\textbf{l.120}] \textsf{\small [phase convention]} $\displaystyle \Big)\int_{x,y,x',y'}e^{-ik\cdot(y'-y)\frac{1}{\xi}}\Big[\frac{2}{\epsilon}$\ $\longrightarrow$\ $\displaystyle \Big)\int_{x,y,x',y'}e^{-ik\cdot(y-y')\frac{1}{\xi}}\Big[\frac{2}{\epsilon}$
\item[\textbf{l.122}] \textsf{\small [phase convention]} $\displaystyle \int_{x,y,x',y'}e^{-ik\cdot(y'-y)}\Big[\frac{2}{\epsilon}$\ $\longrightarrow$\ $\displaystyle \int_{x,y,x',y'}e^{-ik\cdot(y-y')}\Big[\frac{2}{\epsilon}$
\item[\textbf{l.125}] \textsf{\small [limits/range]} $\displaystyle 0$\ $\longrightarrow$\ $\displaystyle \frac{k^{+}}{\vee}$
\item[\textbf{l.125}] \textsf{\small [phase convention]} $\displaystyle \int_{x,y,x',y'}e^{-ik\cdot(y'-y)\frac{1}{\xi}}$\ $\longrightarrow$\ $\displaystyle \int_{x,y,x',y'}e^{-ik\cdot(y-y')\frac{1}{\xi}}$
\end{itemize}
\needspace{3\baselineskip}\noindent\rule{\linewidth}{0.3pt}\par\noindent\textbf{Lines 130} \quad {\small (in \texttt{align}, lines 129--132)}\par
\begin{itemize}[leftmargin=2.2em,itemsep=1pt,topsep=2pt]
\item[\textbf{l.130}] \textsf{\small [typo/LaTeX]} \emph{added:} $\displaystyle \Big[$\quad{\color{gray}\small(in context: }$\displaystyle \frac{1}{(2\pi)^{3}}\,\frac{g^{4}\,N_{c}}{8\pi^{4}}\,\frac{1}{k^{+}}\,\Big[\log\frac{k^{+}}{\Lambda}${\color{gray}\small)}
\item[\textbf{l.130}] \textsf{\small [missing term added]} \emph{added:} $\displaystyle +\log\Big(1-\frac{k^{+}}{\vee}\Big)\Big]$\quad{\color{gray}\small(in context: }$\displaystyle \,\frac{1}{k^{+}}\,\Big[\log\frac{k^{+}}{\Lambda}+\log\Big(1-\frac{k^{+}}{\vee}\Big)\Big]\int_{x,y,x',y'}e^{-ik\cdot(y-y')}\Big[\frac{2}{\epsilon}${\color{gray}\small)}
\item[\textbf{l.130}] \textsf{\small [phase convention]} $\displaystyle \int_{x,y,x',y'}e^{-ik\cdot(y'-y)}\Big[\frac{2}{\epsilon}$\ $\longrightarrow$\ $\displaystyle \int_{x,y,x',y'}e^{-ik\cdot(y-y')}\Big[\frac{2}{\epsilon}$
\end{itemize}
\needspace{3\baselineskip}\noindent\rule{\linewidth}{0.3pt}\par\noindent\textbf{Lines 137--143} \quad {\small (in \texttt{align}, lines 136--142)}\par
\begin{itemize}[leftmargin=2.2em,itemsep=1pt,topsep=2pt]
\item[\textbf{l.137}] \textsf{\small [limits/range]} $\displaystyle \{\int_{0}^{1}d\xi\,$\ $\longrightarrow$\ $\displaystyle \{\int_{\frac{k^{+}}{\vee}}^{1}d\xi\,$
\item[\textbf{l.137}] \textsf{\small [phase convention]} $\displaystyle \Big)\int_{x,y,x',y'}e^{-ik\cdot(y'-y)\frac{1}{\xi}}\Big[\frac{2}{\epsilon}$\ $\longrightarrow$\ $\displaystyle \Big)\int_{x,y,x',y'}e^{-ik\cdot(y-y')\frac{1}{\xi}}\Big[\frac{2}{\epsilon}$
\item[\textbf{l.138}] \textsf{\small [phase convention]} $\displaystyle \int_{x,y,x',y'}e^{-ik\cdot(y'-y)}\Big[\frac{2}{\epsilon}$\ $\longrightarrow$\ $\displaystyle \int_{x,y,x',y'}e^{-ik\cdot(y-y')}\Big[\frac{2}{\epsilon}$
\item[\textbf{l.140}] \textsf{\small [limits/range]} $\displaystyle 0$\ $\longrightarrow$\ $\displaystyle \frac{k^{+}}{\vee}$
\item[\textbf{l.140}] \textsf{\small [phase convention]} $\displaystyle \int_{x,y,x',y'}e^{-ik\cdot(y'-y)\frac{1}{\xi}}$\ $\longrightarrow$\ $\displaystyle \int_{x,y,x',y'}e^{-ik\cdot(y-y')\frac{1}{\xi}}$
\item[\textbf{l.143}] \textsf{\small [typo/LaTeX]} \emph{removed:} Dropping all the $\frac{1}{\epsilon}$, we get
\end{itemize}
\begin{notebox}\textsf{\small Note, l.143:} [The poles cannot be dropped by the rule for scaleless integrals: their coefficient (the $x,y,x',y'$ integrals with the phases $e^{-ik\cdot(y-y')/\xi}$ and $e^{-ik\cdot(y-y')}$) depends on $k$ and on the target. They are kept here, as in the final list below (they are the collinear poles of the splitting of the parent gluon). They are removed only when the final terms are written. The contraction uses the $d$-dimensional polarisation sum, $d_{\perp}=2-\epsilon$, and the same $\epsilon$ as in the UV poles of the virtual terms. In the next equation the two $\xi\to1$ terms are combined.]\end{notebox}
\needspace{3\baselineskip}\noindent\rule{\linewidth}{0.3pt}\par\noindent\textbf{Lines 145--150} \quad {\small (in \texttt{align}, lines 144--150)}\par
\begin{itemize}[leftmargin=2.2em,itemsep=1pt,topsep=2pt]
\item[\textbf{l.145}] \textsf{\small [limits/range]} $\displaystyle \{\int_{0}^{1}d\xi\,$\ $\longrightarrow$\ $\displaystyle \{\int_{\frac{k^{+}}{\vee}}^{1}d\xi\,$
\item[\textbf{l.145}] \textsf{\small [phase convention]} $\displaystyle \Big)\int_{x,y,x',y'}e^{-ik\cdot(y'-y)\frac{1}{\xi}}\Big[$\ $\longrightarrow$\ $\displaystyle \Big)\int_{x,y,x',y'}e^{-ik\cdot(y-y')\frac{1}{\xi}}\Big[$
\item[\textbf{l.145}] \textsf{\small [missing term added]} \emph{added:} $\displaystyle \frac{2}{\epsilon}+$\quad{\color{gray}\small(in context: }$\displaystyle +\frac{1-\xi}{\xi}\Big)\int_{x,y,x',y'}e^{-ik\cdot(y-y')\frac{1}{\xi}}\Big[\frac{2}{\epsilon}+2\gamma_{E}${\color{gray}\small)}
\item[\textbf{l.146}] \textsf{\small [phase convention]} $\displaystyle \int_{x,y,x',y'}e^{-ik\cdot(y'-y)}\Big[$\ $\longrightarrow$\ $\displaystyle \int_{x,y,x',y'}e^{-ik\cdot(y-y')}\Big[$
\item[\textbf{l.146}] \textsf{\small [missing term added]} \emph{added:} $\displaystyle \frac{2}{\epsilon}+$\quad{\color{gray}\small(in context: }$\displaystyle \frac{1}{1-\xi}\int_{x,y,x',y'}e^{-ik\cdot(y-y')}\Big[\frac{2}{\epsilon}+2\gamma_{E}${\color{gray}\small)}
\item[\textbf{l.148}] \textsf{\small [limits/range]} $\displaystyle 0$\ $\longrightarrow$\ $\displaystyle \frac{k^{+}}{\vee}$
\item[\textbf{l.148}] \textsf{\small [phase convention]} $\displaystyle \int_{x,y,x',y'}e^{-ik\cdot(y'-y)\frac{1}{\xi}}$\ $\longrightarrow$\ $\displaystyle \int_{x,y,x',y'}e^{-ik\cdot(y-y')\frac{1}{\xi}}$
\end{itemize}
\begin{notebox}\textsf{\small Note, l.150:} [Here and below $\int_{k^{+}/\vee}^{1}d\xi\,\frac{\xi}{(1-\xi)_{+}}f(\xi)\equiv\int_{k^{+}/\vee}^{1}d\xi\,\frac{\xi f(\xi)-f(1)}{1-\xi}$, as defined at the end.]\end{notebox}
\needspace{3\baselineskip}\noindent\rule{\linewidth}{0.3pt}\par\noindent\textbf{Lines 152--154} \quad {\small (in \texttt{align}, lines 151--156)}\par
\begin{itemize}[leftmargin=2.2em,itemsep=1pt,topsep=2pt]
\item[\textbf{l.152}] \textsf{\small [limits/range]} $\displaystyle \{\int_{0}^{1}d\xi\,$\ $\longrightarrow$\ $\displaystyle \{\int_{\frac{k^{+}}{\vee}}^{1}d\xi\,$
\item[\textbf{l.152}] \textsf{\small [phase convention]} $\displaystyle \Big)\int_{x,y,x',y'}e^{-ik\cdot(y'-y)\frac{1}{\xi}}\Big[$\ $\longrightarrow$\ $\displaystyle \Big)\int_{x,y,x',y'}e^{-ik\cdot(y-y')\frac{1}{\xi}}\Big[$
\item[\textbf{l.152}] \textsf{\small [missing term added]} \emph{added:} $\displaystyle \frac{2}{\epsilon}+$\quad{\color{gray}\small(in context: }$\displaystyle +\frac{1-\xi}{\xi}\Big)\int_{x,y,x',y'}e^{-ik\cdot(y-y')\frac{1}{\xi}}\Big[\frac{2}{\epsilon}+2\gamma_{E}${\color{gray}\small)}
\item[\textbf{l.154}] \textsf{\small [limits/range]} $\displaystyle 0$\ $\longrightarrow$\ $\displaystyle \frac{k^{+}}{\vee}$
\item[\textbf{l.154}] \textsf{\small [phase convention]} $\displaystyle \int_{x,y,x',y'}e^{-ik\cdot(y'-y)\frac{1}{\xi}}$\ $\longrightarrow$\ $\displaystyle \int_{x,y,x',y'}e^{-ik\cdot(y-y')\frac{1}{\xi}}$
\end{itemize}
\needspace{3\baselineskip}\noindent\rule{\linewidth}{0.3pt}\par\noindent\textbf{Lines 159--165} \quad {\small (in \texttt{align}, lines 158--164)}\par
\begin{itemize}[leftmargin=2.2em,itemsep=1pt,topsep=2pt]
\item[\textbf{l.159}] \textsf{\small [limits/range]} $\displaystyle \frac{1}{k^{+}}\int_{0}^{1}d\xi\,$\ $\longrightarrow$\ $\displaystyle \frac{1}{k^{+}}\int_{\frac{k^{+}}{\vee}}^{1}d\xi\,$
\item[\textbf{l.159}] \textsf{\small [phase convention]} $\displaystyle \int_{x,y,x',y'}e^{-ik\cdot(y'-y)\frac{1}{\xi}}\Big[$\ $\longrightarrow$\ $\displaystyle \int_{x,y,x',y'}e^{-ik\cdot(y-y')\frac{1}{\xi}}\Big[$
\item[\textbf{l.159}] \textsf{\small [missing term added]} \emph{added:} $\displaystyle \frac{2}{\epsilon}+$\quad{\color{gray}\small(in context: }$\displaystyle \{\frac{\mathcal P_{gg\,+}(\xi)}{\xi^{2}}\int_{x,y,x',y'}e^{-ik\cdot(y-y')\frac{1}{\xi}}\Big[\frac{2}{\epsilon}+2\gamma_{E}${\color{gray}\small)}
\item[\textbf{l.161}] \textsf{\small [typo/LaTeX]} \emph{removed:} $\displaystyle \bigg]$\quad{\color{gray}\small(original context: }$\displaystyle (x)]\rho^{e'}(x')\rho^{e}(x)\bigg]${\color{gray}\small)}
\item[\textbf{l.162}] \textsf{\small [limits/range]} $\displaystyle 0$\ $\longrightarrow$\ $\displaystyle \frac{k^{+}}{\vee}$
\item[\textbf{l.162}] \textsf{\small [phase convention]} $\displaystyle \int_{x,y,x',y'}e^{-ik\cdot(y'-y)\frac{1}{\xi}}$\ $\longrightarrow$\ $\displaystyle \int_{x,y,x',y'}e^{-ik\cdot(y-y')\frac{1}{\xi}}$
\end{itemize}
\begin{notebox}\textsf{\small Note, l.165:} [The $\beta_{0}$ term is added (first term) and subtracted (third term), with $\int^{1}d\xi\,\delta(1-\xi)=1$; the sum does not depend on $\beta_{0}$. This real-emission term alone does not produce $\beta_{0}$.]\end{notebox}
\needspace{3\baselineskip}\noindent\rule{\linewidth}{0.3pt}\par\noindent\textbf{Lines 167--175} \quad {\small (in \texttt{align}, lines 166--174)}\par
\begin{itemize}[leftmargin=2.2em,itemsep=1pt,topsep=2pt]
\item[\textbf{l.167}] \textsf{\small [limits/range]} $\displaystyle \frac{1}{k^{+}}\int_{0}^{1}d\xi\,$\ $\longrightarrow$\ $\displaystyle \frac{1}{k^{+}}\int_{\frac{k^{+}}{\vee}}^{1}d\xi\,$
\item[\textbf{l.167}] \textsf{\small [phase convention]} $\displaystyle \int_{x,y,x',y'}e^{-ik\cdot(y'-y)\frac{1}{\xi}}\Big[$\ $\longrightarrow$\ $\displaystyle \int_{x,y,x',y'}e^{-ik\cdot(y-y')\frac{1}{\xi}}\Big[$
\item[\textbf{l.167}] \textsf{\small [missing term added]} \emph{added:} $\displaystyle \frac{2}{\epsilon}+$\quad{\color{gray}\small(in context: }$\displaystyle \{\frac{\mathcal P_{gg}(\xi)}{\xi^{2}}\int_{x,y,x',y'}e^{-ik\cdot(y-y')\frac{1}{\xi}}\Big[\frac{2}{\epsilon}+2\gamma_{E}${\color{gray}\small)}
\item[\textbf{l.170}] \textsf{\small [phase convention]} $\displaystyle \int_{x,y,x',y'}e^{-ik\cdot(y'-y)}\Big[$\ $\longrightarrow$\ $\displaystyle \int_{x,y,x',y'}e^{-ik\cdot(y-y')}\Big[$
\item[\textbf{l.170}] \textsf{\small [missing term added]} \emph{added:} $\displaystyle \frac{2}{\epsilon}+$\quad{\color{gray}\small(in context: }$\displaystyle \frac{1}{k^{+}}\int_{x,y,x',y'}e^{-ik\cdot(y-y')}\Big[\frac{2}{\epsilon}+2\gamma_{E}${\color{gray}\small)}
\item[\textbf{l.172}] \textsf{\small [limits/range]} $\displaystyle 0$\ $\longrightarrow$\ $\displaystyle \frac{k^{+}}{\vee}$
\item[\textbf{l.172}] \textsf{\small [phase convention]} $\displaystyle \int_{x,y,x',y'}e^{-ik\cdot(y'-y)\frac{1}{\xi}}$\ $\longrightarrow$\ $\displaystyle \int_{x,y,x',y'}e^{-ik\cdot(y-y')\frac{1}{\xi}}$
\end{itemize}
\begin{notebox}\textsf{\small Note, l.175:} [Checked, also numerically with a test function: the sum of the five terms below equals the expression before the extraction of the end-point logarithms up to $O(\Lambda/k^{+})$, with the coefficients $\log\frac{k^{+}}{\Lambda}+\log(1-\frac{k^{+}}{\vee})$ and $\log\frac{\vee}{k^{+}}-(1-\frac{k^{+}}{\vee})$ and the lower limit $\frac{k^{+}}{\vee}$.]\end{notebox}
\needspace{3\baselineskip}\noindent\rule{\linewidth}{0.3pt}\par\noindent\textbf{Lines 180--184} \quad {\small (in \texttt{align}, lines 179--183)}\par
\begin{itemize}[leftmargin=2.2em,itemsep=1pt,topsep=2pt]
\item[\textbf{l.180}] \textsf{\small [limits/range]} $\displaystyle \frac{1}{k^{+}}\int_{0}^{1}d\xi\,$\ $\longrightarrow$\ $\displaystyle \frac{1}{k^{+}}\int_{\frac{k^{+}}{\vee}}^{1}d\xi\,$
\item[\textbf{l.180}] \textsf{\small [phase convention]} $\displaystyle \int_{x,y,x',y'}e^{-ik\cdot(y'-y)\frac{1}{\xi}}\Big[$\ $\longrightarrow$\ $\displaystyle \int_{x,y,x',y'}e^{-ik\cdot(y-y')\frac{1}{\xi}}\Big[$
\item[\textbf{l.180}] \textsf{\small [missing term added]} \emph{added:} $\displaystyle \frac{2}{\epsilon}+$\quad{\color{gray}\small(in context: }$\displaystyle \{\frac{\mathcal P_{gg}(\xi)}{\xi^{2}}\int_{x,y,x',y'}e^{-ik\cdot(y-y')\frac{1}{\xi}}\Big[\frac{2}{\epsilon}+2\gamma_{E}${\color{gray}\small)}
\end{itemize}
\begin{notebox}\textsf{\small Note, l.184:} [$\widetilde{\mathcal{P}}_{gg}$ is not used below. The change from the previous form is $\frac{1}{\xi}=\frac{1-\xi}{\xi}+1$ in the last term; $\int_{k^{+}/\vee}^{1}d\xi=1-\frac{k^{+}}{\vee}$ is moved to the BFKL term.]\end{notebox}
\needspace{3\baselineskip}\noindent\rule{\linewidth}{0.3pt}\par\noindent\textbf{Lines 187--201} \quad {\small (in \texttt{align}, lines 186--203)}\par
\begin{itemize}[leftmargin=2.2em,itemsep=1pt,topsep=2pt]
\item[\textbf{l.187}] \textsf{\small [phase convention]} $\displaystyle \int_{x,y,x',y'}e^{-ik\cdot(y'-y)\frac{1}{\xi}}\Big[\frac{2}{\epsilon}$\ $\longrightarrow$\ $\displaystyle \int_{x,y,x',y'}e^{-ik\cdot(y-y')\frac{1}{\xi}}\Big[\frac{2}{\epsilon}$
\item[\textbf{l.192}] \textsf{\small [phase convention]} $\displaystyle \int_{x,y,x',y'}e^{-ik\cdot(y'-y)}\Big[\frac{2}{\epsilon}$\ $\longrightarrow$\ $\displaystyle \int_{x,y,x',y'}e^{-ik\cdot(y-y')}\Big[\frac{2}{\epsilon}$
\item[\textbf{l.194}] \textsf{\small [phase convention]} $\displaystyle \Big)\Big]\int_{x,z,w,w',x'}e^{-ik(w'-w)}\,\frac{(w-z)\cdot (w'-z)\,(x-w)\cdot (x'-w')}{(w-z)^{2}(x-w)^{2}(w'-z)^{2}(x'-w')^{2}}$\ $\longrightarrow$\ $\displaystyle \Big)\Big]\int_{x,z,w,w',x'}e^{-ik(w-w')}\,\frac{(w-z)\cdot (w'-z)\,(x-w)\cdot (x'-w')}{(w-z)^{2}(x-w)^{2}(w'-z)^{2}(x'-w')^{2}}$
\item[\textbf{l.196}] \textsf{\small [phase convention]} $\displaystyle \Big)\Big]\int_{x,y,x',y'}e^{-ik\cdot(y'-y)}\Big[\frac{2}{\epsilon}$\ $\longrightarrow$\ $\displaystyle \Big)\Big]\int_{x,y,x',y'}e^{-ik\cdot(y-y')}\Big[\frac{2}{\epsilon}$
\item[\textbf{l.198}] \textsf{\small [phase convention]} $\displaystyle \Big)\Big]\int_{x,z,w,w',x'}e^{-ik(w'-w)}\,\frac{(w-z)\cdot (w'-z)\,(x-z)\cdot (x'-z)}{(w-z)^{2}(x-z)^{2}(w'-z)^{2}(x'-z)^{2}}$\ $\longrightarrow$\ $\displaystyle \Big)\Big]\int_{x,z,w,w',x'}e^{-ik(w-w')}\,\frac{(w-z)\cdot (w'-z)\,(x-z)\cdot (x'-z)}{(w-z)^{2}(x-z)^{2}(w'-z)^{2}(x'-z)^{2}}$
\item[\textbf{l.201}] \textsf{\small [phase convention]} $\displaystyle \int_{x,y,x',y'}e^{-ik\cdot(y'-y)\frac{1}{\xi}}$\ $\longrightarrow$\ $\displaystyle \int_{x,y,x',y'}e^{-ik\cdot(y-y')\frac{1}{\xi}}$
\end{itemize}
\needspace{15\baselineskip}\bigskip
\subsection{Group III}\label{sec:file-group_III}
\noindent\textbf{Original:} \texttt{/root/.claude/uploads/0ab97fcf-94a2-5aec-8bad-586400481685/d98e045a-Group\_III.tex}\\\textbf{Checked copy:} \texttt{/home/user/correlations/nlo\_review/group\_III\_checked/Group\_III\_checked.tex}\\
Lines: 125 (original), 125 (checked copy). Coloured groups of the check: 95 (95 top-level); author\textquotesingle s own coloured groups left unchanged: 8.\\
Records by type: sign/factor correction 1; phase convention 33; limits/range 1; term removed/replaced 1; typo/LaTeX 8; note only 13.
\par\noindent Unmarked layout-only edits (line breaks \texttt{\textbackslash notag\textbackslash\textbackslash\ \&}, \texttt{\textbackslash par\textbackslash noindent} around inserted notes, spacing), checked-copy lines: 35, 41, 51, 52, 59, 63, 76, 81, 92, 97, 98, 100, 104, 105, 108, 112, 113, 115, 119, 123, 124.
\subsubsection*{Changes, block by block}
\needspace{3\baselineskip}\noindent\rule{\linewidth}{0.3pt}\par\noindent\textbf{Lines 35}\par
\begin{itemize}[leftmargin=2.2em,itemsep=1pt,topsep=2pt]
\item[\textbf{l.35}] \textsf{\small [note only]} \emph{added:} \textbf{Marking in this checked copy:} changes and additions are in red;  Line numbers of the source are those of the original file. The phase is $e^{-ik\cdot(w-w^{\prime})}$ throughout (conjugated bracket at $w^{\prime}$, unconjugated bracket at $w$); c.c. is implied; four-$\rho$ terms are collected in the separate four-$\rho$ file.
\end{itemize}
\needspace{3\baselineskip}\noindent\rule{\linewidth}{0.3pt}\par\noindent\textbf{Lines 39--41} \quad \textsf{Group III (with c.c.) $\rangle$ Row I { $\checkmark$]}} \quad {\small (in \texttt{align}, lines 38--41)}\par
\begin{itemize}[leftmargin=2.2em,itemsep=1pt,topsep=2pt]
\item[\textbf{l.39}] \textsf{\small [phase convention]} $\displaystyle \int_{z,w,w^{\prime}} e^{-ik(w^{\prime}-w)}\,$\ $\longrightarrow$\ $\displaystyle \int_{z,w,w^{\prime}} e^{-ik(w-w^{\prime})}\,$
\end{itemize}
\begin{notebox}\textsf{\small Note, l.41:} [Phase: with the conjugated bracket at $w^{\prime}$ (left) and the $\mf{B}_{2}$ bracket at $w$, as written, the definition gives $e^{-ik(w-w^{\prime})}$ (as in the corrected main file); changed here and below. Otherwise this is the cross term of the main file with weight $4\times\frac{1}{2}=2$, and the orientations of $U(z)$, $U(w^{\prime})$ agree with the corrected expression there. Only the single-$\rho$ part of $\mf{B}_{2}$ is kept: it has support $p^{+}+k^{+}<\vee$, hence $\Lambda<p^{+}<\vee-k^{+}$ although the two $\bar{\mf{A}}^{\dagger}$ have the full range. The $\{\rho,\rho\}$ part of $\mf{B}_{2}$ (range $\Lambda<p^{+}<\vee$) gives four-$\rho$ terms, collected separately.]\end{notebox}
\needspace{3\baselineskip}\noindent\rule{\linewidth}{0.3pt}\par\noindent\textbf{Lines 43} \quad {\small (in \texttt{align}, lines 42--44)}\par
\begin{itemize}[leftmargin=2.2em,itemsep=1pt,topsep=2pt]
\item[\textbf{l.43}] \textsf{\small [phase convention]} $\displaystyle (w^{\prime}-w)$\ $\longrightarrow$\ $\displaystyle (w-w^{\prime})$
\end{itemize}
\needspace{3\baselineskip}\noindent\rule{\linewidth}{0.3pt}\par\noindent\textbf{Lines 47} \quad {\small (in \texttt{align}, lines 46--48)}\par
\begin{itemize}[leftmargin=2.2em,itemsep=1pt,topsep=2pt]
\item[\textbf{l.47}] \textsf{\small [phase convention]} $\displaystyle \, \int_{x,y,x',z,w,w^{\prime}} e^{-ik(w^{\prime}-w)}\,$\ $\longrightarrow$\ $\displaystyle \, \int_{x,y,x',z,w,w^{\prime}} e^{-ik(w-w^{\prime})}\,$
\end{itemize}
\needspace{3\baselineskip}\noindent\rule{\linewidth}{0.3pt}\par\noindent\textbf{Lines 51--52} \quad \textsf{Group III (with c.c.) $\rangle$ Row II { $\checkmark$]}} \quad {\small (in \texttt{align}, lines 50--52)}\par
\begin{itemize}[leftmargin=2.2em,itemsep=1pt,topsep=2pt]
\item[\textbf{l.51}] \textsf{\small [limits/range]} \emph{added:} $\displaystyle _{\Lambda}^{\vee-k^{+}}$\quad{\color{gray}\small(in context: }$\displaystyle \frac{2}{(2\pi)^{3}} \int_{\Lambda}^{\vee-k^{+}}\, \frac{dp^{+}}{2\pi}\int_{z,w,w^{\prime}} e^{-ik(w-w^{\prime})}\,${\color{gray}\small)}
\item[\textbf{l.51}] \textsf{\small [phase convention]} $\displaystyle \int_{z,w,w^{\prime}} e^{-ik(w^{\prime}-w)}\,$\ $\longrightarrow$\ $\displaystyle \int_{z,w,w^{\prime}} e^{-ik(w-w^{\prime})}\,$
\item[\textbf{l.51}] \textsf{\small [term removed/replaced]} $\displaystyle k^{+},- U(w^{\prime})^{c^{\prime}b}(z)$\ $\longrightarrow$\ $\displaystyle k^{+},-w^{\prime})\,U^{c^{\prime}b}(z)$
\end{itemize}
\begin{notebox}\textsf{\small Note, l.52:} [Garbled argument repaired: the first term is $\bar{\mf{A}}^{\dagger\,(1)\,c^{\prime}}_{j}(p^{+},-z)\bar{\mf{A}}^{\dagger\,(1)\,a}_{i}(k^{+},-w^{\prime})\,U^{c^{\prime}b}(z)$, as in the main file; limits of $p^{+}$ added (single-$\rho$ part of $\mf{B}_{2}$, as in Row I); phase as in Row I.]\end{notebox}
\needspace{3\baselineskip}\noindent\rule{\linewidth}{0.3pt}\par\noindent\textbf{Lines 54} \quad {\small (in \texttt{align}, lines 53--55)}\par
\begin{itemize}[leftmargin=2.2em,itemsep=1pt,topsep=2pt]
\item[\textbf{l.54}] \textsf{\small [phase convention]} $\displaystyle (w^{\prime}-w)$\ $\longrightarrow$\ $\displaystyle (w-w^{\prime})$
\end{itemize}
\needspace{3\baselineskip}\noindent\rule{\linewidth}{0.3pt}\par\noindent\textbf{Lines 58--59} \quad {\small (in \texttt{align}, lines 57--59)}\par
\begin{itemize}[leftmargin=2.2em,itemsep=1pt,topsep=2pt]
\item[\textbf{l.58}] \textsf{\small [phase convention]} $\displaystyle \, \int_{x,y,x',z,w,w^{\prime}} e^{-ik(w^{\prime}-w)}\,$\ $\longrightarrow$\ $\displaystyle \, \int_{x,y,x',z,w,w^{\prime}} e^{-ik(w-w^{\prime})}\,$
\end{itemize}
\begin{notebox}\textsf{\small Note, l.59:} [Rows I and II have the same kernel and the same last bracket and differ only in the order of the two $\rho$'s of the first bracket. In the sum the first bracket is $\big[U^{ab'}(x')-\frac{1}{2}U^{ab'}(w')\big]\big(U^{\dagger}(z)U(y)\big)^{bd'}\{\rho^{b'}(x'),\rho^{d'}(y)\}-\frac{1}{2}U^{ab'}(x')\{\rho^{b'}(x'),\rho^{b}(y)\}-\frac{i}{2}\,\delta^{2}(x'-y)\,f^{b'd'e}\Big[U^{ab'}(w')\big(U^{\dagger}(z)U(x')\big)^{bd'}-U^{ab'}(x')\,\delta^{bd'}\Big]\rho^{e}(x')$; the last term (from the commutator) has one $\rho$ less.]\end{notebox}
\needspace{3\baselineskip}\noindent\rule{\linewidth}{0.3pt}\par\noindent\textbf{Lines 62--63} \quad \textsf{Group III (with c.c.) $\rangle$ Row III { $\checkmark$]}} \quad {\small (in \texttt{align}, lines 61--63)}\par
\begin{itemize}[leftmargin=2.2em,itemsep=1pt,topsep=2pt]
\item[\textbf{l.62}] \textsf{\small [sign/factor correction]} \emph{added:} $\displaystyle -$\quad{\color{gray}\small(in context: }$\displaystyle -\frac{4}{(2\pi)^{3}} \int_{\Lambda}^{\vee -k^{+}}\, \frac{dp^{+}}{2\pi}\int${\color{gray}\small)}
\item[\textbf{l.62}] \textsf{\small [phase convention]} $\displaystyle \int_{z,w,w^{\prime}} e^{-ik(w^{\prime}-w)}\,\int_{y^{\prime}}$\ $\longrightarrow$\ $\displaystyle \int_{z,w,w^{\prime}} e^{-ik(w-w^{\prime})}\,\int_{y^{\prime}}$
\end{itemize}
\begin{notebox}\textsf{\small Note, l.63:} [Sign: in Eq.~(CS2) of the main file this cross term enters as $-\int_{y^{\prime}}[\dots]\times[\dots]$ (the $\mf{C}_1\mf{A}$ part of the amplitude has the opposite sign to the $\mf{B}_2$ part), so the prefactor is $-4/(2\pi)^3$; the expressions below already correspond to this sign (checked numerically from the coefficient functions). Phase: with the conjugated bracket at $w^{\prime}$ and the other at $w$ it is $e^{-ik\cdot(w-w^{\prime})}$, hence after the $w^{\prime}$ integration $e^{-ik\cdot(w-y^{\prime})}\,e^{-ik\cdot\frac{p^{+}}{k^{+}}(z-y^{\prime})}$. The range $\Lambda<p^{+}<\vee-k^{+}$ is correct for both parts, since the parent gluon of $\mf{C}_1$ carries $p^{+}+k^{+}<\vee$.]\end{notebox}
\needspace{3\baselineskip}\noindent\rule{\linewidth}{0.3pt}\par\noindent\textbf{Lines 67} \quad {\small (in \texttt{align}, lines 66--68)}\par
\begin{itemize}[leftmargin=2.2em,itemsep=1pt,topsep=2pt]
\item[\textbf{l.67}] \textsf{\small [phase convention]} $\displaystyle \, \int_{w,w',x',y',z,y,x}\, e^{-ik\cdot (w'-w)}\, \frac{(w'-z)^{m}\,(x'-y')^{k'}\,(y-w)^{i}\,(x-z)^{j}}{(w'-z)^{2}\,(x'-y')^{2}\,(y-w)^{2}\,(x-z)^{2}}$\ $\longrightarrow$\ $\displaystyle \, \int_{w,w',x',y',z,y,x}\, e^{-ik\cdot (w-w')}\, \frac{(w'-z)^{m}\,(x'-y')^{k'}\,(y-w)^{i}\,(x-z)^{j}}{(w'-z)^{2}\,(x'-y')^{2}\,(y-w)^{2}\,(x-z)^{2}}$
\end{itemize}
\needspace{3\baselineskip}\noindent\rule{\linewidth}{0.3pt}\par\noindent\textbf{Lines 71} \quad {\small (in \texttt{align}, lines 70--72)}\par
\begin{itemize}[leftmargin=2.2em,itemsep=1pt,topsep=2pt]
\item[\textbf{l.71}] \textsf{\small [phase convention]} $\displaystyle (y'-w)$\ $\longrightarrow$\ $\displaystyle (w-y')$
\item[\textbf{l.71}] \textsf{\small [phase convention]} $\displaystyle \, e^{-ik\cdot (y'-w)}\, e^{-ik\cdot \frac{p^{+}}{k^{+}}\,(y'-z)}\,\frac{(y'-z)^{m}\,(x'-y')^{k'}\,(y-w)^{i}\,(x-z)^{j}}{(y'-z)^{2}\,(x'-y')^{2}\,(y-w)^{2}\,(x-z)^{2}}$\ $\longrightarrow$\ $\displaystyle \, e^{-ik\cdot (w-y')}\, e^{-ik\cdot \frac{p^{+}}{k^{+}}\,(z-y')}\,\frac{(y'-z)^{m}\,(x'-y')^{k'}\,(y-w)^{i}\,(x-z)^{j}}{(y'-z)^{2}\,(x'-y')^{2}\,(y-w)^{2}\,(x-z)^{2}}$
\end{itemize}
\needspace{3\baselineskip}\noindent\rule{\linewidth}{0.3pt}\par\noindent\textbf{Lines 75--76} \quad {\small (in \texttt{align}, lines 74--76)}\par
\begin{itemize}[leftmargin=2.2em,itemsep=1pt,topsep=2pt]
\item[\textbf{l.75}] \textsf{\small [phase convention]} $\displaystyle (y'-w)$\ $\longrightarrow$\ $\displaystyle (w-y')$
\item[\textbf{l.75}] \textsf{\small [phase convention]} $\displaystyle \, e^{-ik\cdot (y'-w)}\, e^{-ik\cdot \frac{p^{+}}{k^{+}}\,(y'-z)}\,\frac{(y'-z)^{m}\,(x'-y')^{k'}\,(y-w)^{i}\,(x-z)^{j}}{(y'-z)^{2}\,(x'-y')^{2}\,(y-w)^{2}\,(x-z)^{2}}$\ $\longrightarrow$\ $\displaystyle \, e^{-ik\cdot (w-y')}\, e^{-ik\cdot \frac{p^{+}}{k^{+}}\,(z-y')}\,\frac{(y'-z)^{m}\,(x'-y')^{k'}\,(y-w)^{i}\,(x-z)^{j}}{(y'-z)^{2}\,(x'-y')^{2}\,(y-w)^{2}\,(x-z)^{2}}$
\item[\textbf{l.75}] \textsf{\small [phase convention]} $\displaystyle )\Big]\, \int_{w,x',y',z,y,x}\, e^{-ik\cdot (y'-w)}\,\frac{(y'-z)\cdot (x-z)\,(x'-y')\cdot (y-w)}{(y'-z)^{2}\,(x'-y')^{2}\,(y-w)^{2}\,(x-z)^{2}}$\ $\longrightarrow$\ $\displaystyle )\Big]\, \int_{w,x',y',z,y,x}\, e^{-ik\cdot (w-y')}\,\frac{(y'-z)\cdot (x-z)\,(x'-y')\cdot (y-w)}{(y'-z)^{2}\,(x'-y')^{2}\,(y-w)^{2}\,(x-z)^{2}}$
\end{itemize}
\begin{notebox}\textsf{\small Note, l.76:} [Checked: prefactors, $\delta$-function, tensor bracket, colour structure and the explicit logarithmic term, which is the $p^{+}\to 0$ limit of the $1/p^{+}$ term; with $\int_{0}^{\vee-k^{+}}dp^{+}\,[1/p^{+}]_{+}f=\int_{0}^{\vee-k^{+}}dp^{+}\,(f(p^{+})-f(0))/p^{+}$ the result equals the cut-off integral up to terms of order $\Lambda$ at fixed transverse coordinates.]\end{notebox}
\needspace{3\baselineskip}\noindent\rule{\linewidth}{0.3pt}\par\noindent\textbf{Lines 79--81} \quad {\small (in \texttt{align}, lines 78--81)}\par
\begin{itemize}[leftmargin=2.2em,itemsep=1pt,topsep=2pt]
\item[\textbf{l.79}] \textsf{\small [typo/LaTeX]} \emph{added:} $\displaystyle x',$
\item[\textbf{l.79}] \textsf{\small [typo/LaTeX]} \emph{removed:} $\displaystyle y,$
\item[\textbf{l.79}] \textsf{\small [phase convention]} $\displaystyle (w'-w)$\ $\longrightarrow$\ $\displaystyle (w-w')$
\end{itemize}
\begin{notebox}\textsf{\small Note, l.81:} [Only the phase and the list of integration variables are changed. The transverse integrals of this part are locally integrable in two dimensions (the integrand behaves as $1/|z-w|$, $1/|x-w|$, $1/|x-z|$, $1/|y^{\prime}-z|$ at coinciding points), so no ultraviolet pole multiplies $d_{\perp}=\delta_{ii}$ and $d_{\perp}=2$ may be set.]\end{notebox}
\needspace{3\baselineskip}\noindent\rule{\linewidth}{0.3pt}\par\noindent\textbf{Lines 84} \quad {\small (in \texttt{align}, lines 83--86)}\par
\begin{itemize}[leftmargin=2.2em,itemsep=1pt,topsep=2pt]
\item[\textbf{l.84}] \textsf{\small [typo/LaTeX]} $\displaystyle \, \frac{dp^{+}}{2\pi}\,\int_{w,w',y',z,y,x}\, e^{-ik\cdot (y'-w)}\,e^{-ik\cdot \frac{p^{+}}{k^{+}}(y'-z)}$\ $\longrightarrow$\ $\displaystyle \, \frac{dp^{+}}{2\pi}\,\int_{w,x',y',z,x}\, e^{-ik\cdot (w-y')}\,e^{-ik\cdot \frac{p^{+}}{k^{+}}(z-y')}$
\item[\textbf{l.84}] \textsf{\small [typo/LaTeX]} \emph{removed:} $\displaystyle y,$\quad{\color{gray}\small(original context: }$\displaystyle w,w',y',z,y,x${\color{gray}\small)}
\item[\textbf{l.84}] \textsf{\small [phase convention]} $\displaystyle (y'-w)$\ $\longrightarrow$\ $\displaystyle (w-y')$
\item[\textbf{l.84}] \textsf{\small [phase convention]} $\displaystyle (y'-z)$\ $\longrightarrow$\ $\displaystyle (z-y')$
\end{itemize}
\needspace{3\baselineskip}\noindent\rule{\linewidth}{0.3pt}\par\noindent\textbf{Lines 89--94} \quad {\small (in \texttt{align}, lines 88--92)}\par
\begin{itemize}[leftmargin=2.2em,itemsep=1pt,topsep=2pt]
\item[\textbf{l.89}] \textsf{\small [typo/LaTeX]} $\displaystyle \, \frac{dp^{+}}{2\pi}\,\int_{w,w',y',z,y,x}\, e^{-ik\cdot (y'-w)}\,e^{-ik\cdot \frac{p^{+}}{k^{+}}(y'-z)}$\ $\longrightarrow$\ $\displaystyle \, \frac{dp^{+}}{2\pi}\,\int_{w,x',y',z,x}\, e^{-ik\cdot (w-y')}\,e^{-ik\cdot \frac{p^{+}}{k^{+}}(z-y')}$
\item[\textbf{l.89}] \textsf{\small [typo/LaTeX]} \emph{removed:} $\displaystyle y,$\quad{\color{gray}\small(original context: }$\displaystyle w,w',y',z,y,x${\color{gray}\small)}
\item[\textbf{l.89}] \textsf{\small [phase convention]} $\displaystyle (y'-w)$\ $\longrightarrow$\ $\displaystyle (w-y')$
\item[\textbf{l.89}] \textsf{\small [phase convention]} $\displaystyle (y'-z)$\ $\longrightarrow$\ $\displaystyle (z-y')$
\item[\textbf{l.90}] \textsf{\small [typo/LaTeX]} \emph{added:} $\displaystyle x',$\quad{\color{gray}\small(in context: }$\displaystyle \Big)\Big]\,\int_{w,x',y',z,x}\, e^{-ik\cdot (w-y')}\,${\color{gray}\small)}
\item[\textbf{l.90}] \textsf{\small [typo/LaTeX]} \emph{removed:} $\displaystyle y,$\quad{\color{gray}\small(original context: }$\displaystyle \Big)\Big]\,\int_{w,y',z,y,x}\, e^{-ik\cdot (y'-w)}\,${\color{gray}\small)}
\item[\textbf{l.90}] \textsf{\small [phase convention]} $\displaystyle )\Big]\,\int_{w,y',z,y,x}\, e^{-ik\cdot (y'-w)}\,$\ $\longrightarrow$\ $\displaystyle )\Big]\,\int_{w,x',y',z,x}\, e^{-ik\cdot (w-y')}\,$
\end{itemize}
\begin{notebox}\textsf{\small Note, l.92:} [Checked as in the first part; the logarithmic term is the $p^{+}\to 0$ limit of the $1/p^{+}$ term, including $1/[p^{+}(x-z)^{2}+k^{+}(x-w)^{2}]\to 1/[k^{+}(x-w)^{2}]$. In both parts the kernel of the logarithmic term decreases only as $1/z^{2}$ at large $z$: for Wilson lines that do not make the colour factor vanish there, the $z$ integrals of the logarithmic term and of the plus-prescription term diverge logarithmically one by one and only their sum is defined.]\end{notebox}
\begin{notebox}\textsf{\small Note, l.94:} [Checked: this is the cross term $\frac14[\bar{\mf{A}}^{\dagger}\bar{\mf{A}}^{\dagger}U-\mf{A}^{\dagger}\bar{\mf{A}}^{\dagger}]\times[U^{\dagger}\bar{\mf{A}}\bar{\mf{A}}-U\bar{\mf{A}}U\mf{A}]$ of the two orderings; each $\mf{A}^{(1)}$ carries exactly one $\rho$, so only four-$\rho$ terms arise before reordering.]\end{notebox}
\needspace{3\baselineskip}\noindent\rule{\linewidth}{0.3pt}\par\noindent\textbf{Lines 97--98} \quad \textsf{Group III (with c.c.) $\rangle$ Row V { $\checkmark$]}} \quad {\small (in \texttt{align}, lines 96--98)}\par
\begin{itemize}[leftmargin=2.2em,itemsep=1pt,topsep=2pt]
\item[\textbf{l.97}] \textsf{\small [phase convention]} $\displaystyle \,\frac{2}{(2\pi)^{3}} \int\, \frac{dp^{+}}{2\pi}\int_{z,w,w^{\prime}} e^{-ik(w^{\prime}-w)}\,$\ $\longrightarrow$\ $\displaystyle \,\frac{2}{(2\pi)^{3}} \int\, \frac{dp^{+}}{2\pi}\int_{z,w,w^{\prime}} e^{-ik(w-w^{\prime})}\,$
\end{itemize}
\begin{notebox}\textsf{\small Note, l.98:} [Phase: with the conjugated bracket at $w'$ and the other one at $w$, as written, the definition gives $e^{-ik(w-w')}$ (as in the corrected Eq.~(CS2) of the main file); factor $-2=4\times(-\frac12)$ and both brackets agree with the corrected (CS2). Range, not written here: $\Lambda<p^{+}<\vee-k^{+}$, because the parent gluon of $\mf{C}_1$ is itself a soft mode.]\end{notebox}
\needspace{3\baselineskip}\noindent\rule{\linewidth}{0.3pt}\par\noindent\textbf{Lines 100} \quad {\small (in \texttt{align}, lines 99--101)}\par
\begin{itemize}[leftmargin=2.2em,itemsep=1pt,topsep=2pt]
\item[\textbf{l.100}] \textsf{\small [phase convention]} $\displaystyle \,\int_{\Lambda}^{\vee-k^{+}}\, \frac{dp^{+}}{2\pi}\, \int_{w,w',y',z,y,x,x'}\, e^{-ik\cdot (w'-w)}\,$\ $\longrightarrow$\ $\displaystyle \,\int_{\Lambda}^{\vee-k^{+}}\, \frac{dp^{+}}{2\pi}\, \int_{w,w',y',z,y,x,x'}\, e^{-ik\cdot (w-w')}\,$
\end{itemize}
\needspace{3\baselineskip}\noindent\rule{\linewidth}{0.3pt}\par\noindent\textbf{Lines 104--105} \quad {\small (in \texttt{align}, lines 103--105)}\par
\begin{itemize}[leftmargin=2.2em,itemsep=1pt,topsep=2pt]
\item[\textbf{l.104}] \textsf{\small [phase convention]} $\displaystyle \,\int_{\Lambda}^{\vee-k^{+}}\, \frac{dp^{+}}{2\pi}\, \int_{w,y',z,y,x,x'}\, e^{-ik\cdot (y'-w)}\,$\ $\longrightarrow$\ $\displaystyle \,\int_{\Lambda}^{\vee-k^{+}}\, \frac{dp^{+}}{2\pi}\, \int_{w,y',z,y,x,x'}\, e^{-ik\cdot (w-y')}\,$
\item[\textbf{l.104}] \textsf{\small [phase convention]} $\displaystyle \, \frac{dp^{+}}{2\pi}\, \int_{w,y',z,y,x,x'}\, e^{-ik\cdot (y'-w)}\, e^{-ik\cdot \frac{p^{+}}{k^{+}}\,(y'-z)}\,$\ $\longrightarrow$\ $\displaystyle \, \frac{dp^{+}}{2\pi}\, \int_{w,y',z,y,x,x'}\, e^{-ik\cdot (w-y')}\, e^{-ik\cdot \frac{p^{+}}{k^{+}}\,(z-y')}\,$
\end{itemize}
\begin{notebox}\textsf{\small Note, l.105:} [After the $w'$ integration $w'=y'+\frac{p^{+}}{k^{+}}(y'-z)$, so the corrected phase becomes $e^{-ik\cdot(w-y')}e^{-ik\cdot\frac{p^{+}}{k^{+}}(z-y')}$. The reduction $U^{d'b}(z)U^{cb}(z)=\delta^{d'c}$ used in the next equation holds already here and in the previous equation.]\end{notebox}
\needspace{3\baselineskip}\noindent\rule{\linewidth}{0.3pt}\par\noindent\textbf{Lines 108} \quad {\small (in \texttt{align}, lines 107--109)}\par
\begin{itemize}[leftmargin=2.2em,itemsep=1pt,topsep=2pt]
\item[\textbf{l.108}] \textsf{\small [phase convention]} $\displaystyle \,\int_{0}^{\vee-k^{+}}\, \frac{dp^{+}}{2\pi}\, \int_{w,y',z,y,x,x'}\, e^{-ik\cdot (y'-w)}\,$\ $\longrightarrow$\ $\displaystyle \,\int_{0}^{\vee-k^{+}}\, \frac{dp^{+}}{2\pi}\, \int_{w,y',z,y,x,x'}\, e^{-ik\cdot (w-y')}\,$
\item[\textbf{l.108}] \textsf{\small [phase convention]} $\displaystyle \, \frac{dp^{+}}{2\pi}\, \int_{w,y',z,y,x,x'}\, e^{-ik\cdot (y'-w)}\, e^{-ik\cdot \frac{p^{+}}{k^{+}}\,(y'-z)}\,$\ $\longrightarrow$\ $\displaystyle \, \frac{dp^{+}}{2\pi}\, \int_{w,y',z,y,x,x'}\, e^{-ik\cdot (w-y')}\, e^{-ik\cdot \frac{p^{+}}{k^{+}}\,(z-y')}\,$
\item[\textbf{l.108}] \textsf{\small [phase convention]} $\displaystyle \log\Big(1-\frac{k^{+}}{\vee}\Big)\Big] \int_{w,y',z,y,x,x'}\, e^{-ik\cdot (y'-w)}\,$\ $\longrightarrow$\ $\displaystyle \log\Big(1-\frac{k^{+}}{\vee}\Big)\Big] \int_{w,y',z,y,x,x'}\, e^{-ik\cdot (w-y')}\,$
\end{itemize}
\needspace{3\baselineskip}\noindent\rule{\linewidth}{0.3pt}\par\noindent\textbf{Lines 112--113} \quad \textsf{Group III (with c.c.) $\rangle$ Row VI { $\checkmark$]}} \quad {\small (in \texttt{align}, lines 111--113)}\par
\begin{itemize}[leftmargin=2.2em,itemsep=1pt,topsep=2pt]
\item[\textbf{l.112}] \textsf{\small [phase convention]} $\displaystyle \, \frac{2}{(2\pi)^{3}} \int\, \frac{dp^{+}}{2\pi}\int_{z,w,w^{\prime}} e^{-ik(w^{\prime}-w)}\,$\ $\longrightarrow$\ $\displaystyle \, \frac{2}{(2\pi)^{3}} \int\, \frac{dp^{+}}{2\pi}\int_{z,w,w^{\prime}} e^{-ik(w-w^{\prime})}\,$
\end{itemize}
\begin{notebox}\textsf{\small Note, l.113:} [Phase and range as in Row V: $e^{-ik(w-w')}$, $\Lambda<p^{+}<\vee-k^{+}$.]\end{notebox}
\needspace{3\baselineskip}\noindent\rule{\linewidth}{0.3pt}\par\noindent\textbf{Lines 115} \quad {\small (in \texttt{align}, lines 114--116)}\par
\begin{itemize}[leftmargin=2.2em,itemsep=1pt,topsep=2pt]
\item[\textbf{l.115}] \textsf{\small [phase convention]} $\displaystyle \,\int_{\Lambda}^{\vee-k^{+}}\, \frac{dp^{+}}{2\pi}\, \int_{w,w',y',x',z,y,x}\, e^{-ik\cdot (w'-w)}\,$\ $\longrightarrow$\ $\displaystyle \,\int_{\Lambda}^{\vee-k^{+}}\, \frac{dp^{+}}{2\pi}\, \int_{w,w',y',x',z,y,x}\, e^{-ik\cdot (w-w')}\,$
\end{itemize}
\needspace{3\baselineskip}\noindent\rule{\linewidth}{0.3pt}\par\noindent\textbf{Lines 119} \quad {\small (in \texttt{align}, lines 118--120)}\par
\begin{itemize}[leftmargin=2.2em,itemsep=1pt,topsep=2pt]
\item[\textbf{l.119}] \textsf{\small [phase convention]} $\displaystyle \,\int_{\Lambda}^{\vee-k^{+}}\, \frac{dp^{+}}{2\pi}\, \int_{w,y',z,y,x,x'}\, e^{-ik\cdot (y'-w)}\,$\ $\longrightarrow$\ $\displaystyle \,\int_{\Lambda}^{\vee-k^{+}}\, \frac{dp^{+}}{2\pi}\, \int_{w,y',z,y,x,x'}\, e^{-ik\cdot (w-y')}\,$
\item[\textbf{l.119}] \textsf{\small [phase convention]} $\displaystyle \, \frac{dp^{+}}{2\pi}\, \int_{w,y',z,y,x,x'}\, e^{-ik\cdot (y'-w)}\, e^{-ik\cdot \frac{p^{+}}{k^{+}}\,(y'-z)}\,$\ $\longrightarrow$\ $\displaystyle \, \frac{dp^{+}}{2\pi}\, \int_{w,y',z,y,x,x'}\, e^{-ik\cdot (w-y')}\, e^{-ik\cdot \frac{p^{+}}{k^{+}}\,(z-y')}\,$
\end{itemize}
\needspace{3\baselineskip}\noindent\rule{\linewidth}{0.3pt}\par\noindent\textbf{Lines 123--124} \quad {\small (in \texttt{align}, lines 122--124)}\par
\begin{itemize}[leftmargin=2.2em,itemsep=1pt,topsep=2pt]
\item[\textbf{l.123}] \textsf{\small [phase convention]} $\displaystyle \,\int_{0}^{\vee-k^{+}}\, \frac{dp^{+}}{2\pi}\, \int_{w,y',z,y,x,x'}\, e^{-ik\cdot (y'-w)}\,$\ $\longrightarrow$\ $\displaystyle \,\int_{0}^{\vee-k^{+}}\, \frac{dp^{+}}{2\pi}\, \int_{w,y',z,y,x,x'}\, e^{-ik\cdot (w-y')}\,$
\item[\textbf{l.123}] \textsf{\small [phase convention]} $\displaystyle \, \frac{dp^{+}}{2\pi}\, \int_{w,y',z,y,x,x'}\, e^{-ik\cdot (y'-w)}\, e^{-ik\cdot \frac{p^{+}}{k^{+}}\,(y'-z)}\,$\ $\longrightarrow$\ $\displaystyle \, \frac{dp^{+}}{2\pi}\, \int_{w,y',z,y,x,x'}\, e^{-ik\cdot (w-y')}\, e^{-ik\cdot \frac{p^{+}}{k^{+}}\,(z-y')}\,$
\item[\textbf{l.123}] \textsf{\small [phase convention]} $\displaystyle \log\Big(1-\frac{k^{+}}{\vee}\Big)\Big]\, \int_{w,y',z,y,x,x'}\, e^{-ik\cdot (y'-w)}\,$\ $\longrightarrow$\ $\displaystyle \log\Big(1-\frac{k^{+}}{\vee}\Big)\Big]\, \int_{w,y',z,y,x,x'}\, e^{-ik\cdot (w-y')}\,$
\end{itemize}
\begin{notebox}\textsf{\small Note, l.124:} [Checked with explicit matrices: the order of the $\rho$'s in Rows V and VI is as written. Here $U^{d'b}(z)U^{cb}(z)=\delta^{d'c}$ applies only to the first term of the last bracket; the colour structure is $f^{c'd'a}\big[U^{c'e'}(y')-U^{c'e'}(x')\big]U^{ad}(y)\big[U^{d'e}(x)-U^{d'e}(z)\big]\rho^{e'}(x')\rho^{d}(y)\rho^{e}(x)$, and that of Row V is $f^{c'd'a}\big[U^{c'e'}(y')-U^{c'e'}(x')\big]U^{d'e}(x)\big[U^{ad}(y)-U^{ad}(w)\big]\rho^{e'}(x')\rho^{e}(x)\rho^{d}(y)$, with identical kernels. Hence Row V $+$ Row VI has $\rho^{e'}(x')$ times $\big[U^{d'e}(x)U^{ad}(y)-\frac12U^{ad}(w)U^{d'e}(x)-\frac12U^{ad}(y)U^{d'e}(z)\big]\{\rho^{e}(x),\rho^{d}(y)\}-\frac12\big[U^{ad}(w)U^{d'e}(x)-U^{ad}(y)U^{d'e}(z)\big][\rho^{e}(x),\rho^{d}(y)]$.]\end{notebox}
\needspace{15\baselineskip}\bigskip
\subsection{Four-rho}\label{sec:file-four_rho}
\noindent\textbf{Original:} \texttt{/root/.claude/uploads/0ab97fcf-94a2-5aec-8bad-586400481685/43078eec-4rho.tex}\\\textbf{Checked copy:} \texttt{/home/user/correlations/nlo\_review/four\_rho\_checked/4rho\_checked.tex}\\
Lines: 227 (original), 227 (checked copy). Coloured groups of the check: 53 (53 top-level); author\textquotesingle s own coloured groups left unchanged: 0.\\
Records by type: phase convention 17; index/Wilson-line orientation 8; note only 3.
\par\noindent Unmarked layout-only edits (line breaks \texttt{\textbackslash notag\textbackslash\textbackslash\ \&}, \texttt{\textbackslash par\textbackslash noindent} around inserted notes, spacing), checked-copy lines: 69, 124.
\subsubsection*{Changes, block by block}
\needspace{3\baselineskip}\noindent\rule{\linewidth}{0.3pt}\par\noindent\textbf{Lines 65--69} \quad \textsf{Terms proportional to $\log{\frac{\vee}{k^{+}}}$ and $\log{\frac{k^{+}}{\Lambda}}$ $\rangle$ Four $\rho$ contributions} \quad {\small (in \texttt{align}, lines 64--69, label \texttt{eq:4rho})}\par
\begin{itemize}[leftmargin=2.2em,itemsep=1pt,topsep=2pt]
\item[\textbf{l.65}] \textsf{\small [phase convention]} $\displaystyle (w^{\prime}-w)$\ $\longrightarrow$\ $\displaystyle (w-w^{\prime})$
\item[\textbf{l.67}] \textsf{\small [phase convention]} $\displaystyle (w^{\prime}-w)$\ $\longrightarrow$\ $\displaystyle (w-w^{\prime})$
\item[\textbf{l.67}] \textsf{\small [index/Wilson-line orientation]} $\displaystyle (x^{\prime})\}U^{c^{\prime}a}(w^{\prime})$\ $\longrightarrow$\ $\displaystyle (x^{\prime})\}U^{ac^{\prime}}(w^{\prime})$
\item[\textbf{l.67}] \textsf{\small [index/Wilson-line orientation]} $\displaystyle )U^{bc^{\prime}}(z)$\ $\longrightarrow$\ $\displaystyle )U^{c^{\prime}b}(z)$
\item[\textbf{l.67}] \textsf{\small [index/Wilson-line orientation]} $\displaystyle \rho^{b}(y)\}-U^{bc}(z)$\ $\longrightarrow$\ $\displaystyle \rho^{b}(y)\}-U^{cb}(z)$
\item[\textbf{l.67}] \textsf{\small [index/Wilson-line orientation]} $\displaystyle (x^{\prime})U^{bd^{\prime}}(z)U^{d^{\prime}a^{\prime}}(x$\ $\longrightarrow$\ $\displaystyle (x^{\prime})U^{d^{\prime}b}(z)U^{d^{\prime}a^{\prime}}(x$
\item[\textbf{l.67}] \textsf{\small [index/Wilson-line orientation]} $\displaystyle x)\rho^{a^{\prime}}(x)U^{c^{\prime}a}(w^{\prime})$\ $\longrightarrow$\ $\displaystyle x)\rho^{a^{\prime}}(x)U^{ac^{\prime}}(w^{\prime})$
\item[\textbf{l.67}] \textsf{\small [index/Wilson-line orientation]} $\displaystyle [U^{bc}(z)U^{cd}(y$\ $\longrightarrow$\ $\displaystyle [U^{cb}(z)U^{cd}(y$
\item[\textbf{l.67}] \textsf{\small [index/Wilson-line orientation]} $\displaystyle (x^{\prime})U^{bc^{\prime}}(z)$\ $\longrightarrow$\ $\displaystyle (x^{\prime})U^{c^{\prime}b}(z)$
\item[\textbf{l.67}] \textsf{\small [index/Wilson-line orientation]} $\displaystyle [U^{bc}(z)U^{ad}(u$\ $\longrightarrow$\ $\displaystyle [U^{cb}(z)U^{ad}(u$
\item[\textbf{l.68}] \textsf{\small [phase convention]} $\displaystyle (w^{\prime}-w)$\ $\longrightarrow$\ $\displaystyle (w-w^{\prime})$
\end{itemize}
\begin{notebox}\textsf{\small Note, l.69:} [Compared term by term with the four-$\rho$ part of Eqs.~(CS1) and (CS2) of the main file, keeping the order of the $\rho$'s (matrix check with SU(2) and SU(3) generators for $\rho$): prefactor, kernels, signs, factors $\frac12$, $2$ and operator ordering agree. Changes: (i) with the conjugated LO bracket at $(x^{\prime},w^{\prime})$ and the $g^{3}$ bracket at $(u,w)$ the phase is $e^{-ik(w-w^{\prime})}$; this is irrelevant for the symmetrised four-$\rho$ terms but reverses the sign of the three-$\rho$ terms generated by reordering, so it is changed in every equation of this file (the three-$\rho$ equations then remain the reordering remainder of this one); (ii) in the second block eight index pairs of $U(z)$ and $U(w^{\prime})$ were transposed (they are written correctly in the third block, and the two- and four-Wilson-line equations below correspond to the corrected form); (iii) operator ordering of the Group I block (second and third lines): the $\mf{N}$-type terms of the amplitude are in total $-\bar{\mf{A}}\mf{N}_{2}+U\bar{\mf{N}}_{2}\mf{A}+U\mf{A}\mf{N}_{2}-\bar{\mf{N}}_{2}\bar{\mf{A}}$ ($\mf{A}\equiv\mf{A}^{(1)}$; the three-$\rho$ term of $\mf{A}^{(3)}$ is $\mf{A}^{(1)}\mf{N}_{2}$, and the commutator $-[\bar{\mf{N}}_{2},\bar{\mf{A}}]$ of Row I cancels against the difference between $-\bar{\mf{A}}\bar{\mf{N}}_{2}$ and $-\bar{\mf{N}}_{2}\bar{\mf{A}}$). The Row II piece is written here as $U\mf{N}_{2}\mf{A}-\bar{\mf{N}}_{2}\bar{\mf{A}}$ (third line). This equation with its ordering is the correct four-$\rho$ content provided the term $-2k^{i}\log(\vee/\Lambda)$ of Group I Row II is dropped for the unbarred colour structure $f^{hbc}U^{ah}(w)\{\rho^{b},\rho^{c}\}$, as is done in the Group I file (since $U\mf{A}\mf{N}_{2}=U\mf{N}_{2}\mf{A}-U[\mf{N}_{2},\mf{A}]$ and $-[\mf{N}_{2},\mf{A}]$ is minus that term); equivalently, the unbarred term of the third line is reordered to $U\mf{A}\mf{N}_{2}$, i.e. $U^{ab}(w)\rho^{b}(u)\rho^{c}(x)\rho^{c}(y)$, and that term is kept. No commutator term is to be added; (iv) a common $\log(\vee/\Lambda)$ requires $\Lambda<p^{+}<\vee$ for the $\{\rho,\rho\}$ part of $\mf{B}_{2}$ (two independent emissions, no condition on $k^{+}+p^{+}$); only then do the four-$\rho$ terms cancel for commuting $\rho$.]\end{notebox}
\needspace{3\baselineskip}\noindent\rule{\linewidth}{0.3pt}\par\noindent\textbf{Lines 72} \quad {\small (in \texttt{align}, lines 71--73)}\par
\begin{itemize}[leftmargin=2.2em,itemsep=1pt,topsep=2pt]
\item[\textbf{l.72}] \textsf{\small [phase convention]} $\displaystyle (w^{\prime}-w)$\ $\longrightarrow$\ $\displaystyle (w-w^{\prime})$
\end{itemize}
\needspace{3\baselineskip}\noindent\rule{\linewidth}{0.3pt}\par\noindent\textbf{Lines 76--77} \quad {\small (in \texttt{align}, lines 75--79)}\par
\begin{itemize}[leftmargin=2.2em,itemsep=1pt,topsep=2pt]
\item[\textbf{l.76}] \textsf{\small [phase convention]} $\displaystyle (w^{\prime}-w)$\ $\longrightarrow$\ $\displaystyle (w-w^{\prime})$
\item[\textbf{l.77}] \textsf{\small [phase convention]} $\displaystyle (w^{\prime}-w)$\ $\longrightarrow$\ $\displaystyle (w-w^{\prime})$
\item[\textbf{l.77}] \textsf{\small [phase convention]} $\displaystyle (w^{\prime}-w)$\ $\longrightarrow$\ $\displaystyle (w-w^{\prime})$
\end{itemize}
\needspace{3\baselineskip}\noindent\rule{\linewidth}{0.3pt}\par\noindent\textbf{Lines 82} \quad {\small (in \texttt{align}, lines 81--83)}\par
\begin{itemize}[leftmargin=2.2em,itemsep=1pt,topsep=2pt]
\item[\textbf{l.82}] \textsf{\small [phase convention]} $\displaystyle (w^{\prime}-w)$\ $\longrightarrow$\ $\displaystyle (w-w^{\prime})$
\end{itemize}
\needspace{3\baselineskip}\noindent\rule{\linewidth}{0.3pt}\par\noindent\textbf{Lines 89} \quad \textsf{Terms proportional to $\log{\frac{\vee}{k^{+}}}$ and $\log{\frac{k^{+}}{\Lambda}}$ $\rangle$ Three $\rho$ contributions} \quad {\small (in \texttt{align}, lines 88--91)}\par
\begin{itemize}[leftmargin=2.2em,itemsep=1pt,topsep=2pt]
\item[\textbf{l.89}] \textsf{\small [phase convention]} $\displaystyle \,\,\int_{x,y,z,x^{\prime},w,w^{\prime}}e^{-ik(w^{\prime}-w)}\,$\ $\longrightarrow$\ $\displaystyle \,\,\int_{x,y,z,x^{\prime},w,w^{\prime}}e^{-ik(w-w^{\prime})}\,$
\end{itemize}
\needspace{3\baselineskip}\noindent\rule{\linewidth}{0.3pt}\par\noindent\textbf{Lines 94} \quad {\small (in \texttt{align}, lines 93--96)}\par
\begin{itemize}[leftmargin=2.2em,itemsep=1pt,topsep=2pt]
\item[\textbf{l.94}] \textsf{\small [phase convention]} $\displaystyle \,\,\int_{x,y,z,x^{\prime},w,w^{\prime}}e^{-ik(w^{\prime}-w)}\,$\ $\longrightarrow$\ $\displaystyle \,\,\int_{x,y,z,x^{\prime},w,w^{\prime}}e^{-ik(w-w^{\prime})}\,$
\end{itemize}
\needspace{3\baselineskip}\noindent\rule{\linewidth}{0.3pt}\par\noindent\textbf{Lines 99} \quad {\small (in \texttt{align}, lines 98--103)}\par
\begin{itemize}[leftmargin=2.2em,itemsep=1pt,topsep=2pt]
\item[\textbf{l.99}] \textsf{\small [phase convention]} $\displaystyle \,\,\int_{x,y,z,x^{\prime},w,w^{\prime}}e^{-ik(w^{\prime}-w)}\,$\ $\longrightarrow$\ $\displaystyle \,\,\int_{x,y,z,x^{\prime},w,w^{\prime}}e^{-ik(w-w^{\prime})}\,$
\end{itemize}
\needspace{3\baselineskip}\noindent\rule{\linewidth}{0.3pt}\par\noindent\textbf{Lines 106} \quad {\small (in \texttt{align}, lines 105--110)}\par
\begin{itemize}[leftmargin=2.2em,itemsep=1pt,topsep=2pt]
\item[\textbf{l.106}] \textsf{\small [phase convention]} $\displaystyle \,\,\int_{x,y,z,x^{\prime},w,w^{\prime}}e^{-ik(w^{\prime}-w)}\,$\ $\longrightarrow$\ $\displaystyle \,\,\int_{x,y,z,x^{\prime},w,w^{\prime}}e^{-ik(w-w^{\prime})}\,$
\end{itemize}
\needspace{3\baselineskip}\noindent\rule{\linewidth}{0.3pt}\par\noindent\textbf{Lines 113} \quad {\small (in \texttt{align}, lines 112--117)}\par
\begin{itemize}[leftmargin=2.2em,itemsep=1pt,topsep=2pt]
\item[\textbf{l.113}] \textsf{\small [phase convention]} $\displaystyle \,\,\int_{x,y,z,x^{\prime},w,w^{\prime}}e^{-ik(w^{\prime}-w)}\,$\ $\longrightarrow$\ $\displaystyle \,\,\int_{x,y,z,x^{\prime},w,w^{\prime}}e^{-ik(w-w^{\prime})}\,$
\end{itemize}
\needspace{3\baselineskip}\noindent\rule{\linewidth}{0.3pt}\par\noindent\textbf{Lines 120--124} \quad {\small (in \texttt{align}, lines 119--124)}\par
\begin{itemize}[leftmargin=2.2em,itemsep=1pt,topsep=2pt]
\item[\textbf{l.120}] \textsf{\small [phase convention]} $\displaystyle \,\,\int_{x,y,z,x^{\prime},w,w^{\prime}}e^{-ik(w^{\prime}-w)}\,$\ $\longrightarrow$\ $\displaystyle \,\,\int_{x,y,z,x^{\prime},w,w^{\prime}}e^{-ik(w-w^{\prime})}\,$
\end{itemize}
\begin{notebox}\textsf{\small Note, l.124:} [The two- and three-Wilson-line expressions of this subsection are exactly the three-$\rho$ remainder obtained by symmetrising Eq.~(\texttt{eq:4rho}) (checked numerically). No commutator term from Row I is to be added: with the three-$\rho$ term of $\mf{A}^{(3)}$ ordered as $\mf{A}^{(1)}\mf{N}_{2}$ it cancels, and Eq.~(\texttt{eq:4rho}) with its ordering is the complete four-$\rho$ content once the term $-2k^{i}\log(\vee/\Lambda)$ of Group I Row II is dropped for the unbarred colour structure (or, equivalently, the unbarred Row II term is reordered to $U\mf{A}\mf{N}_{2}$ and that term kept). The symmetrised three-$\rho$ terms then cancel against the three-$\rho$ terms of the other rows for both logarithms (evolution files).]\end{notebox}
\needspace{3\baselineskip}\noindent\rule{\linewidth}{0.3pt}\par\noindent\textbf{Lines 130} \quad \textsf{Terms proportional to $\log{\frac{\vee}{k^{+}}}$ and $\log{\frac{k^{+}}{\Lambda}}$ $\rangle$ Two $\rho$ contribution}\par
\begin{notebox}\textsf{\small Note, l.130:} [Not confirmed. As written these terms do not vanish at $U=1$: the terms with $\mathcal{K}(x^{\prime}-z)\cdot\mathcal{K}(y-z)$ leave $N_{c}\,\rho^{a}(x^{\prime})\rho^{a}(y)\,[3\,\mathcal{K}(x^{\prime}-w^{\prime})\cdot\mathcal{K}(y-w)-3\,\mathcal{K}(y-w^{\prime})\cdot\mathcal{K}(x^{\prime}-w)-8\,\mathcal{K}(x^{\prime}-w^{\prime})\cdot\mathcal{K}(x^{\prime}-w)]$; the second term of $N^{(1)}$ carries the index $b^{\prime}$ four times; and the sum does not reproduce the two-$\rho$ remainder of Eq.~(\texttt{eq:4rho}). The first four terms of the first part contain the scaleless integral $\int_{z}1/(x^{\prime}-z)^{2}$.]\end{notebox}
\needspace{3\baselineskip}\noindent\rule{\linewidth}{0.3pt}\par\noindent\textbf{Lines 133} \quad {\small (in \texttt{align}, lines 132--139)}\par
\begin{itemize}[leftmargin=2.2em,itemsep=1pt,topsep=2pt]
\item[\textbf{l.133}] \textsf{\small [phase convention]} $\displaystyle \,{\color{blue}[\log{[\frac{k^{+}}{\Lambda}]}+\log{[\frac{\vee}{k^{+}}]}]}\,\int_{x,y,z,x^{\prime},w,w^{\prime}}e^{-ik(w^{\prime}-w)}\,$\ $\longrightarrow$\ $\displaystyle \,{\color{blue}[\log{[\frac{k^{+}}{\Lambda}]}+\log{[\frac{\vee}{k^{+}}]}]}\,\int_{x,y,z,x^{\prime},w,w^{\prime}}e^{-ik(w-w^{\prime})}\,$
\end{itemize}
\needspace{3\baselineskip}\noindent\rule{\linewidth}{0.3pt}\par\noindent\textbf{Lines 142} \quad {\small (in \texttt{align}, lines 141--163)}\par
\begin{itemize}[leftmargin=2.2em,itemsep=1pt,topsep=2pt]
\item[\textbf{l.142}] \textsf{\small [phase convention]} $\displaystyle \,\,\int_{x,y,z,x^{\prime},w,w^{\prime}}e^{-ik(w^{\prime}-w)}\,$\ $\longrightarrow$\ $\displaystyle \,\,\int_{x,y,z,x^{\prime},w,w^{\prime}}e^{-ik(w-w^{\prime})}\,$
\end{itemize}
\needspace{3\baselineskip}\noindent\rule{\linewidth}{0.3pt}\par\noindent\textbf{Lines 165} \quad {\small (in \texttt{align}, lines 164--169)}\par
\begin{itemize}[leftmargin=2.2em,itemsep=1pt,topsep=2pt]
\item[\textbf{l.165}] \textsf{\small [phase convention]} $\displaystyle \,\,\int_{x,y,z,x^{\prime},w,w^{\prime}}e^{-ik(w^{\prime}-w)}\,$\ $\longrightarrow$\ $\displaystyle \,\,\int_{x,y,z,x^{\prime},w,w^{\prime}}e^{-ik(w-w^{\prime})}\,$
\end{itemize}
\needspace{15\baselineskip}\bigskip
\subsection{Evolution two-rho}\label{sec:file-ev_two_rho}
\noindent\textbf{Original:} \texttt{/root/.claude/uploads/0ab97fcf-94a2-5aec-8bad-586400481685/294d9dfb-Evolution\_two\_rho.tex}\\\textbf{Checked copy:} \texttt{/home/user/correlations/nlo\_review/evolution\_checked/Evolution\_two\_rho\_checked.tex}\\
Lines: 108 (original), 108 (checked copy). Coloured groups of the check: 37 (37 top-level); author\textquotesingle s own coloured groups left unchanged: 0.\\
Records by type: phase convention 15; note only 7.
\par\noindent Unmarked layout-only edits (line breaks \texttt{\textbackslash notag\textbackslash\textbackslash\ \&}, \texttt{\textbackslash par\textbackslash noindent} around inserted notes, spacing), checked-copy lines: 35, 45, 65, 73, 84, 101.
\subsubsection*{Changes, block by block}
\needspace{3\baselineskip}\noindent\rule{\linewidth}{0.3pt}\par\noindent\textbf{Lines 35}\par
\begin{notebox}\textsf{\small Note, l.35:} \textbf{Checked copy.} Marking: corrected or added material is in red; material removed from the author's text is kept in orange (cancelled); notes of the check are in red in square brackets. Phase: $e^{-ik\cdot(w-w')}$ everywhere (conjugated bracket, with the primed coordinates and standing left, at $w'$); this does not change the symmetrised two-$\rho$ parts, while the three-$\rho$ and one-$\rho$ parts change sign together. Result: every entry was compared with the coefficient of $\log\frac{k^{+}}{\Lambda}$ in the checked Group I, II (Row IV file) and III files (numerically, with $SU(2)$ and $SU(3)$ matrices for the $\rho$'s), and all eight equations are correct. The only change is the phase.\end{notebox}
\needspace{3\baselineskip}\noindent\rule{\linewidth}{0.3pt}\par\noindent\textbf{Lines 38} \quad \textsf{Group I (with c.c.) $\rangle$ Row I}\par
\begin{notebox}\textsf{\small Note, l.38:} [Checked, nothing else: the commutator term $-[\bar N_{2},\bar A^{(1)}]$ of Row I of the Group I file cancels against the difference between the correctly ordered three-$\rho$ term $-\bar A^{(1)}\bar N_{2}$ of $-\bar A^{(3)}$ (Row II) and $-\bar N_{2}\bar A^{(1)}$ (see the resolution of the operator ordering), so no three-$\rho$ or two-$\rho$ term comes from it.]\end{notebox}
\needspace{3\baselineskip}\noindent\rule{\linewidth}{0.3pt}\par\noindent\textbf{Lines 41} \quad \textsf{Group I (with c.c.) $\rangle$ Row II} \quad {\small (in \texttt{align}, lines 40--42)}\par
\begin{itemize}[leftmargin=2.2em,itemsep=1pt,topsep=2pt]
\item[\textbf{l.41}] \textsf{\small [phase convention]} $\displaystyle \int_{w, w^\prime} e^{-ik(w^\prime-w)}\,$\ $\longrightarrow$\ $\displaystyle \int_{w, w^\prime} e^{-ik(w-w^\prime)}\,$
\end{itemize}
\needspace{3\baselineskip}\noindent\rule{\linewidth}{0.3pt}\par\noindent\textbf{Lines 44--45} \quad {\small (in \texttt{align}, lines 43--45)}\par
\begin{itemize}[leftmargin=2.2em,itemsep=1pt,topsep=2pt]
\item[\textbf{l.44}] \textsf{\small [phase convention]} $\displaystyle (w^\prime-w)$\ $\longrightarrow$\ $\displaystyle (w-w^\prime)$
\end{itemize}
\begin{notebox}\textsf{\small Note, l.45:} [Agrees with $N_{1}$ of the Group I file: $\frac{g^{4}N_{c}}{32\pi^{4}}\times(-2\log\frac{k^{+}}{\Lambda})$.]\end{notebox}
\needspace{3\baselineskip}\noindent\rule{\linewidth}{0.3pt}\par\noindent\textbf{Lines 52} \quad \textsf{Group I (with c.c.) $\rangle$ Row V} \quad {\small (in \texttt{align}, lines 51--53)}\par
\begin{itemize}[leftmargin=2.2em,itemsep=1pt,topsep=2pt]
\item[\textbf{l.52}] \textsf{\small [phase convention]} $\displaystyle \int_{w, w^\prime} e^{-ik(w^\prime-w)}\, \int\, \frac{dp^{+}}{2\pi}$\ $\longrightarrow$\ $\displaystyle \int_{w, w^\prime} e^{-ik(w-w^\prime)}\, \int\, \frac{dp^{+}}{2\pi}$
\end{itemize}
\needspace{3\baselineskip}\noindent\rule{\linewidth}{0.3pt}\par\noindent\textbf{Lines 56} \quad {\small (in \texttt{align}, lines 55--57)}\par
\begin{itemize}[leftmargin=2.2em,itemsep=1pt,topsep=2pt]
\item[\textbf{l.56}] \textsf{\small [phase convention]} $\displaystyle \,{\color{blue}\log{[\frac{k^{+}}{\Lambda}]}} \int_{x,y,z,x',w'}\, e^{-ik\cdot(w'-z)}\,\frac{(x'-w')^{i}\,(z-x)^{m}}{(x'-w')^{2}\,(z-x)^{2}\, (y-z)^{2}}\,$\ $\longrightarrow$\ $\displaystyle \,{\color{blue}\log{[\frac{k^{+}}{\Lambda}]}} \int_{x,y,z,x',w'}\, e^{-ik\cdot(z-w')}\,\frac{(x'-w')^{i}\,(z-x)^{m}}{(x'-w')^{2}\,(z-x)^{2}\, (y-z)^{2}}\,$
\item[\textbf{l.56}] \textsf{\small [phase convention]} $\displaystyle \, \int_{x,y,z,x',w'}\,e^{-ik\cdot(w'-x)}\,\frac{(x'-w')^{i}\,(z-x)^{m}}{(x'-w')^{2}\,(z-x)^{2}\,(y-x)^{2}}$\ $\longrightarrow$\ $\displaystyle \, \int_{x,y,z,x',w'}\,e^{-ik\cdot(x-w')}\,\frac{(x'-w')^{i}\,(z-x)^{m}}{(x'-w')^{2}\,(z-x)^{2}\,(y-x)^{2}}$
\end{itemize}
\needspace{3\baselineskip}\noindent\rule{\linewidth}{0.3pt}\par\noindent\textbf{Lines 60} \quad {\small (in \texttt{align}, lines 58--62)}\par
\begin{itemize}[leftmargin=2.2em,itemsep=1pt,topsep=2pt]
\item[\textbf{l.60}] \textsf{\small [phase convention]} $\displaystyle \,\int_{y,z,x',w'}e^{-ik\cdot(w'-z)}\, [\frac{2}{\epsilon}$\ $\longrightarrow$\ $\displaystyle \,\int_{y,z,x',w'}e^{-ik\cdot(z-w')}\, [\frac{2}{\epsilon}$
\end{itemize}
\needspace{3\baselineskip}\noindent\rule{\linewidth}{0.3pt}\par\noindent\textbf{Lines 64--65} \quad {\small (in \texttt{align}, lines 63--65)}\par
\begin{itemize}[leftmargin=2.2em,itemsep=1pt,topsep=2pt]
\item[\textbf{l.64}] \textsf{\small [phase convention]} $\displaystyle \,{\color{blue}\log{(\frac{k^{+}}{\Lambda})}}\,\int_{x,y,z,x',w'}e^{-ik\cdot(w'-x)}\,$\ $\longrightarrow$\ $\displaystyle \,{\color{blue}\log{(\frac{k^{+}}{\Lambda})}}\,\int_{x,y,z,x',w'}e^{-ik\cdot(x-w')}\,$
\item[\textbf{l.64}] \textsf{\small [phase convention]} $\displaystyle \,{\color{blue}\log{(\frac{k^{+}}{\Lambda})}}\,\int_{x,y,z,x',w'}e^{-ik\cdot(w'-z)}\,$\ $\longrightarrow$\ $\displaystyle \,{\color{blue}\log{(\frac{k^{+}}{\Lambda})}}\,\int_{x,y,z,x',w'}e^{-ik\cdot(z-w')}\,$
\item[\textbf{l.64}] \textsf{\small [phase convention]} $\displaystyle (w'-z)$\ $\longrightarrow$\ $\displaystyle (z-w')$
\end{itemize}
\begin{notebox}\textsf{\small Note, l.65:} [$2a$, $2b$ agree with the Group I file (Row V); there the UV term and the last term are written with $x\leftrightarrow z$, i.e. with $e^{-ik\cdot(x-w')}$.]\end{notebox}
\needspace{3\baselineskip}\noindent\rule{\linewidth}{0.3pt}\par\noindent\textbf{Lines 69} \quad \textsf{Group II $\rangle$ Row I} \quad {\small (in \texttt{align}, lines 68--70)}\par
\begin{itemize}[leftmargin=2.2em,itemsep=1pt,topsep=2pt]
\item[\textbf{l.69}] \textsf{\small [phase convention]} $\displaystyle \int_{z,w,w^{\prime}} e^{-ik(w^{\prime}-w)}\,$\ $\longrightarrow$\ $\displaystyle \int_{z,w,w^{\prime}} e^{-ik(w-w^{\prime})}\,$
\end{itemize}
\needspace{3\baselineskip}\noindent\rule{\linewidth}{0.3pt}\par\noindent\textbf{Lines 72--73} \quad {\small (in \texttt{align}, lines 71--73)}\par
\begin{itemize}[leftmargin=2.2em,itemsep=1pt,topsep=2pt]
\item[\textbf{l.72}] \textsf{\small [phase convention]} $\displaystyle \,\int_{z,w,w^{\prime},x,x'}\, e^{-ik(w^{\prime}-w)}\, \frac{1}{[(x'-w')^{2}\,(x-z)^{2}\,(x-w)^{2} - (x-w)^{4}\, (x'-z)^{2}]}$\ $\longrightarrow$\ $\displaystyle \,\int_{z,w,w^{\prime},x,x'}\, e^{-ik(w-w^{\prime})}\, \frac{1}{[(x'-w')^{2}\,(x-z)^{2}\,(x-w)^{2} - (x-w)^{4}\, (x'-z)^{2}]}$
\end{itemize}
\begin{notebox}\textsf{\small Note, l.73:} [Agrees with $N_{2b}$ of the Group II file; the kernel equals $\mathcal{K}(x-w)\cdot\mathcal{K}(x'-w')\big[K(w,w',z)-\frac12K(x',w,z)-\frac12K(x,w',z)+\frac14K(x,x',z)\big]$, $K(a,b,z)=\frac{(a-z)\cdot(b-z)}{(a-z)^{2}(b-z)^{2}}$.]\end{notebox}
\needspace{3\baselineskip}\noindent\rule{\linewidth}{0.3pt}\par\noindent\textbf{Lines 80} \quad \textsf{Group II $\rangle$ Row IV} \quad {\small (in \texttt{align}, lines 79--81)}\par
\begin{itemize}[leftmargin=2.2em,itemsep=1pt,topsep=2pt]
\item[\textbf{l.80}] \textsf{\small [phase convention]} $\displaystyle \int_{z,w,w^{\prime}} e^{-ik(w^{\prime}-w)}\, \int_{y,y^{\prime}}$\ $\longrightarrow$\ $\displaystyle \int_{z,w,w^{\prime}} e^{-ik(w-w^{\prime})}\, \int_{y,y^{\prime}}$
\end{itemize}
\needspace{3\baselineskip}\noindent\rule{\linewidth}{0.3pt}\par\noindent\textbf{Lines 83--84} \quad {\small (in \texttt{align}, lines 82--84)}\par
\begin{itemize}[leftmargin=2.2em,itemsep=1pt,topsep=2pt]
\item[\textbf{l.83}] \textsf{\small [phase convention]} $\displaystyle \,\int_{x,y,x',y'}\,e^{-ik\cdot (y'-y)}\,\frac{(x-y)\cdot (x'-y')}{(x-y)^{2}\,(x'-y')^{2}}\,$\ $\longrightarrow$\ $\displaystyle \,\int_{x,y,x',y'}\,e^{-ik\cdot (y-y')}\,\frac{(x-y)\cdot (x'-y')}{(x-y)^{2}\,(x'-y')^{2}}\,$
\end{itemize}
\begin{notebox}\textsf{\small Note, l.84:} [With $d_{\perp}=2-\epsilon$ the bracket is $-\big[\frac{2}{\epsilon}+2\gamma_{E}+\log\frac{(y'-y)^{2}\mu^{2}_{\overline{MS}}}{4}\big]+{\cal O}(\epsilon)$ (the $-1$ cancels), as in the Row IV file; equivalently $+\frac{g^{4}N_{c}}{8\pi^{5}}\int_{z}K(y,y',z)$ in place of $\frac{g^{4}N_{c}}{8\pi^{4}}[\dots]$, up to scaleless integrals.]\end{notebox}
\needspace{3\baselineskip}\noindent\rule{\linewidth}{0.3pt}\par\noindent\textbf{Lines 92} \quad \textsf{Group III (with c.c.) $\rangle$ Row III} \quad {\small (in \texttt{align}, lines 91--93)}\par
\begin{itemize}[leftmargin=2.2em,itemsep=1pt,topsep=2pt]
\item[\textbf{l.92}] \textsf{\small [phase convention]} $\displaystyle \int_{z,w,w^{\prime}} e^{-ik(w^{\prime}-w)}\,\int_{y^{\prime}}$\ $\longrightarrow$\ $\displaystyle \int_{z,w,w^{\prime}} e^{-ik(w-w^{\prime})}\,\int_{y^{\prime}}$
\end{itemize}
\needspace{3\baselineskip}\noindent\rule{\linewidth}{0.3pt}\par\noindent\textbf{Lines 99--101} \quad {\small (in \texttt{align}, lines 98--101)}\par
\begin{itemize}[leftmargin=2.2em,itemsep=1pt,topsep=2pt]
\item[\textbf{l.99}] \textsf{\small [phase convention]} $\displaystyle \,\int_{w,x',y',z,x}\, e^{-ik\cdot (y'-w)}\,$\ $\longrightarrow$\ $\displaystyle \,\int_{w,x',y',z,x}\, e^{-ik\cdot (w-y')}\,$
\end{itemize}
\begin{notebox}\textsf{\small Note, l.101:} [Agrees with the Group III file (Row III, second part).]\end{notebox}
\needspace{15\baselineskip}\bigskip
\subsection{Evolution three-rho}\label{sec:file-ev_three_rho}
\noindent\textbf{Original:} \texttt{/root/.claude/uploads/0ab97fcf-94a2-5aec-8bad-586400481685/6a6a9fd3-Evolution\_three\_rho.tex}\\\textbf{Checked copy:} \texttt{/home/user/correlations/nlo\_review/evolution\_checked/Evolution\_three\_rho\_checked.tex}\\
Lines: 316 (original), 316 (checked copy). Coloured groups of the check: 204 (204 top-level); author\textquotesingle s own coloured groups left unchanged: 260.\\
Records by type: sign/factor correction 2; phase convention 61; index/Wilson-line orientation 26; term removed/replaced 9; typo/LaTeX 1; note only 10.
\par\noindent Unmarked layout-only edits (line breaks \texttt{\textbackslash notag\textbackslash\textbackslash\ \&}, \texttt{\textbackslash par\textbackslash noindent} around inserted notes, spacing), checked-copy lines: 35, 47, 68, 103, 107, 150.
\subsubsection*{Changes, block by block}
\needspace{3\baselineskip}\noindent\rule{\linewidth}{0.3pt}\par\noindent\textbf{Lines 35}\par
\begin{notebox}\textsf{\small Note, l.35:} \textbf{Checked copy.} Marking: corrected or added material is in red; material removed from the author's text is kept in orange (cancelled); notes of the check are in red in square brackets. The coloured $U$'s and the coloured labels $\alpha$, $\alpha'$, $\beta$, $\beta'$ of the list ``All three $\rho$ term with c.c.'' and of the cancellation section are the author's own. Phase: $e^{-ik\cdot(w-w')}$ everywhere (conjugated bracket, with the primed coordinates and standing left, at $w'$); this does not change the symmetrised two-$\rho$ parts, while the three-$\rho$ and one-$\rho$ parts change sign together. Result (checked with explicit $SU(2)$ matrices for $\rho$ on lattice sites, so that contact terms are included): the collected list and the cancellation hold, ``$=0$'' meaning that the totally symmetrised three-$\rho$ part vanishes identically, and with the operator ordering $A^{(1)}N_{2}$ of the three-$\rho$ part of $A^{(3)}$ no term has to be added. Corrections are confined to the row-by-row equations (two signs and one $U(z)\to U(w)$ in Group I Row III, transposed colour indices of some $U$'s) and to misprints (kernels of the $\beta$, $\beta'$ headers, a bracket, a sign); one point is open (Group I Row III).\end{notebox}
\needspace{3\baselineskip}\noindent\rule{\linewidth}{0.3pt}\par\noindent\textbf{Lines 38} \quad \textsf{Group I (with c.c.) $\rangle$ Row I}\par
\begin{notebox}\textsf{\small Note, l.38:} [Checked, nothing else: the commutator term $-[\bar N_{2},\bar A^{(1)}]$ of Row I of the Group I file cancels against the difference between the correctly ordered three-$\rho$ term $-\bar A^{(1)}\bar N_{2}$ of $-\bar A^{(3)}$ (Row II) and $-\bar N_{2}\bar A^{(1)}$ (see the resolution of the operator ordering), so no three-$\rho$ or two-$\rho$ term comes from it.]\end{notebox}
\needspace{3\baselineskip}\noindent\rule{\linewidth}{0.3pt}\par\noindent\textbf{Lines 41} \quad \textsf{Group I (with c.c.) $\rangle$ Row II} \quad {\small (in \texttt{align}, lines 40--42)}\par
\begin{itemize}[leftmargin=2.2em,itemsep=1pt,topsep=2pt]
\item[\textbf{l.41}] \textsf{\small [phase convention]} $\displaystyle \int_{w, w^\prime} e^{-ik(w^\prime-w)}\,$\ $\longrightarrow$\ $\displaystyle \int_{w, w^\prime} e^{-ik(w-w^\prime)}\,$
\end{itemize}
\needspace{3\baselineskip}\noindent\rule{\linewidth}{0.3pt}\par\noindent\textbf{Lines 45--47} \quad {\small (in \texttt{align}, lines 43--47)}\par
\begin{itemize}[leftmargin=2.2em,itemsep=1pt,topsep=2pt]
\item[\textbf{l.45}] \textsf{\small [phase convention]} $\displaystyle \int_{w,w'} e^{-ik\cdot(w'-w)}\int_{x,y,x',z}\frac{(x'-w')^{i}}{(x'-w')^{2}}$\ $\longrightarrow$\ $\displaystyle \int_{w,w'} e^{-ik\cdot(w-w')}\int_{x,y,x',z}\frac{(x'-w')^{i}}{(x'-w')^{2}}$
\item[\textbf{l.46}] \textsf{\small [phase convention]} $\displaystyle \int_{w,w'} e^{-ik\cdot(w'-w)}\int_{x,y,x',z}\frac{(x'-w')^{i}}{(x'-w')^{2}}$\ $\longrightarrow$\ $\displaystyle \int_{w,w'} e^{-ik\cdot(w-w')}\int_{x,y,x',z}\frac{(x'-w')^{i}}{(x'-w')^{2}}$
\end{itemize}
\begin{notebox}\textsf{\small Note, l.47:} [Checked against $N_{2a}$ of the Group I file, where $\log\frac{(w-y)^{2}}{(x-y)^{2}}=\frac{1}{\pi}\int_{z}[\mathcal{K}(x-z)\cdot\mathcal{K}(y-z)-\mathcal{K}(w-z)\cdot\mathcal{K}(y-z)]$: the second term agrees; the first term has only the $\mathcal{K}(w-z)$ part for the colour structure $f^{hbc}U^{ah}(w)\{\rho^{b}(x),\rho^{c}(y)\}$, and this is correct. Its $\mathcal{K}(x-z)\cdot\mathcal{K}(y-z)$ part is absent because of the ordering $A^{(1)}N_{2}$ of the three-$\rho$ part of $A^{(3)}$ (the term $-2k^{i}\log\frac{\vee}{\Lambda}$ of Row II is absent for this structure); the same structure appears instead with $-\log\frac{\vee}{k^{+}}$ in the BFKL file.]\end{notebox}
\needspace{3\baselineskip}\noindent\rule{\linewidth}{0.3pt}\par\noindent\textbf{Lines 50} \quad \textsf{Group I (with c.c.) $\rangle$ Row III} \quad {\small (in \texttt{align}, lines 49--51)}\par
\begin{itemize}[leftmargin=2.2em,itemsep=1pt,topsep=2pt]
\item[\textbf{l.50}] \textsf{\small [phase convention]} $\displaystyle \int_{w, w^\prime} e^{-ik(w^\prime-w)}\, \int\,\frac{dp^{+}}{2\pi}$\ $\longrightarrow$\ $\displaystyle \int_{w, w^\prime} e^{-ik(w-w^\prime)}\, \int\,\frac{dp^{+}}{2\pi}$
\end{itemize}
\needspace{3\baselineskip}\noindent\rule{\linewidth}{0.3pt}\par\noindent\textbf{Lines 54} \quad {\small (in \texttt{align}, lines 53--55)}\par
\begin{itemize}[leftmargin=2.2em,itemsep=1pt,topsep=2pt]
\item[\textbf{l.54}] \textsf{\small [phase convention]} $\displaystyle \int_{w, w^\prime} e^{-ik(w^\prime-w)}\,$\ $\longrightarrow$\ $\displaystyle \int_{w, w^\prime} e^{-ik(w-w^\prime)}\,$
\end{itemize}
\needspace{3\baselineskip}\noindent\rule{\linewidth}{0.3pt}\par\noindent\textbf{Lines 57} \quad {\small (in \texttt{align}, lines 56--61)}\par
\begin{itemize}[leftmargin=2.2em,itemsep=1pt,topsep=2pt]
\item[\textbf{l.57}] \textsf{\small [phase convention]} $\displaystyle \int_{w, w^\prime} e^{-ik(w^\prime-w)}\,$\ $\longrightarrow$\ $\displaystyle \int_{w, w^\prime} e^{-ik(w-w^\prime)}\,$
\item[\textbf{l.57}] \textsf{\small [sign/factor correction]} \emph{added:} $\displaystyle -$\quad{\color{gray}\small(in context: }$\displaystyle -\frac{ig^{4}\, f^{abd}}{8\pi^{5}}\, \frac{1}{k^{+}}\,{\color{blue}\log{[\frac{k^{+}}{\Lambda}]}}\, \int_{x,y,z,x'}${\color{gray}\small)}
\item[\textbf{l.57}] \textsf{\small [sign/factor correction]} $\displaystyle - \frac{ig^{4}\, f^{abd}}{8\pi^{5}}\, \frac{1}{k^{+}}\,{\color{blue}\log{[\frac{k^{+}}{\Lambda}]}}\, \int_{x,y,z,x'}$\ $\longrightarrow$\ $\displaystyle + \frac{ig^{4}\, f^{abd}}{8\pi^{5}}\, \frac{1}{k^{+}}\,{\color{blue}\log{[\frac{k^{+}}{\Lambda}]}}\, \int_{x,y,z,x'}$
\end{itemize}
\needspace{3\baselineskip}\noindent\rule{\linewidth}{0.3pt}\par\noindent\textbf{Lines 64} \quad {\small (in \texttt{align}, lines 63--65)}\par
\begin{itemize}[leftmargin=2.2em,itemsep=1pt,topsep=2pt]
\item[\textbf{l.64}] \textsf{\small [phase convention]} $\displaystyle \,\int_{x,y,z,w,w',x'}\,e^{-ik(w^\prime-w)}\, \mathcal{K}(x'-w$\ $\longrightarrow$\ $\displaystyle \,\int_{x,y,z,w,w',x'}\,e^{-ik(w-w^\prime)}\, \mathcal{K}(x'-w$
\end{itemize}
\needspace{3\baselineskip}\noindent\rule{\linewidth}{0.3pt}\par\noindent\textbf{Lines 67--68} \quad {\small (in \texttt{align}, lines 66--68)}\par
\begin{itemize}[leftmargin=2.2em,itemsep=1pt,topsep=2pt]
\item[\textbf{l.67}] \textsf{\small [phase convention]} $\displaystyle \,\int_{x,y,z,w,w',x'}\,e^{-ik(w^\prime-w)}\, \mathcal{K}(x'-w$\ $\longrightarrow$\ $\displaystyle \,\int_{x,y,z,w,w',x'}\,e^{-ik(w-w^\prime)}\, \mathcal{K}(x'-w$
\item[\textbf{l.67}] \textsf{\small [term removed/replaced]} $\displaystyle \,[U^{de}(x)\,U^{bc}(z)\, \rho^{e}(x)\,\rho^{c}(y)$\ $\longrightarrow$\ $\displaystyle \,[U^{de}(x)\,U^{bc}(w)\, \rho^{e}(x)\,\rho^{c}(y)$
\end{itemize}
\begin{notebox}\textsf{\small Note, l.68:} [Signs of the first two terms of the long equation corrected ($N_{2}$ of the Group I file has an overall minus sign), so that they cancel the last two terms, as used in the last equation. In the last equation $U^{bc}(z)\to U^{bc}(w)$, as in the third term of the long equation and in the collected list (checked: only then do the two forms agree). Open (not decided): at fixed coordinates the checked Group I row gives the kernel $-\mathcal{K}(x'-w')\cdot\mathcal{K}(y-z)/(z-w)^{2}$ with $U^{bc}(z)$, whereas the evolution files use $-\mathcal{K}(x'-w')\cdot\mathcal{K}(y-w)\,\mathcal{K}(x-z)\cdot\mathcal{K}(w-z)$ with $U^{bc}(w)$; the step between them ($z$ integration with $\Lambda$ kept; the region $z\to w$) was not verified.]\end{notebox}
\needspace{3\baselineskip}\noindent\rule{\linewidth}{0.3pt}\par\noindent\textbf{Lines 71} \quad \textsf{Group I (with c.c.) $\rangle$ Row IV} \quad {\small (in \texttt{align}, lines 70--72)}\par
\begin{itemize}[leftmargin=2.2em,itemsep=1pt,topsep=2pt]
\item[\textbf{l.71}] \textsf{\small [phase convention]} $\displaystyle \int_{w, w^\prime} e^{-ik(w^\prime-w)}\, \int\,\frac{dp^{+}}{2\pi}$\ $\longrightarrow$\ $\displaystyle \int_{w, w^\prime} e^{-ik(w-w^\prime)}\, \int\,\frac{dp^{+}}{2\pi}$
\end{itemize}
\needspace{3\baselineskip}\noindent\rule{\linewidth}{0.3pt}\par\noindent\textbf{Lines 75} \quad {\small (in \texttt{align}, lines 73--76)}\par
\begin{itemize}[leftmargin=2.2em,itemsep=1pt,topsep=2pt]
\item[\textbf{l.75}] \textsf{\small [phase convention]} $\displaystyle (w^\prime-w)$\ $\longrightarrow$\ $\displaystyle (w-w^\prime)$
\end{itemize}
\needspace{3\baselineskip}\noindent\rule{\linewidth}{0.3pt}\par\noindent\textbf{Lines 80} \quad {\small (in \texttt{align}, lines 78--81)}\par
\begin{itemize}[leftmargin=2.2em,itemsep=1pt,topsep=2pt]
\item[\textbf{l.80}] \textsf{\small [phase convention]} $\displaystyle (w^\prime-w)$\ $\longrightarrow$\ $\displaystyle (w-w^\prime)$
\end{itemize}
\needspace{3\baselineskip}\noindent\rule{\linewidth}{0.3pt}\par\noindent\textbf{Lines 84} \quad \textsf{Group I (with c.c.) $\rangle$ Row V} \quad {\small (in \texttt{align}, lines 83--85)}\par
\begin{itemize}[leftmargin=2.2em,itemsep=1pt,topsep=2pt]
\item[\textbf{l.84}] \textsf{\small [phase convention]} $\displaystyle \int_{w, w^\prime} e^{-ik(w^\prime-w)}\, \int\, \frac{dp^{+}}{2\pi}$\ $\longrightarrow$\ $\displaystyle \int_{w, w^\prime} e^{-ik(w-w^\prime)}\, \int\, \frac{dp^{+}}{2\pi}$
\end{itemize}
\needspace{3\baselineskip}\noindent\rule{\linewidth}{0.3pt}\par\noindent\textbf{Lines 87} \quad {\small (in \texttt{align}, lines 86--88)}\par
\begin{itemize}[leftmargin=2.2em,itemsep=1pt,topsep=2pt]
\item[\textbf{l.87}] \textsf{\small [phase convention]} $\displaystyle (w^\prime-w)$\ $\longrightarrow$\ $\displaystyle (w-w^\prime)$
\item[\textbf{l.87}] \textsf{\small [phase convention]} $\displaystyle (w^\prime-w)$\ $\longrightarrow$\ $\displaystyle (w-w^\prime)$
\end{itemize}
\needspace{3\baselineskip}\noindent\rule{\linewidth}{0.3pt}\par\noindent\textbf{Lines 91} \quad {\small (in \texttt{align}, lines 90--92)}\par
\begin{itemize}[leftmargin=2.2em,itemsep=1pt,topsep=2pt]
\item[\textbf{l.91}] \textsf{\small [phase convention]} $\displaystyle \, {\color{blue}\log{[\frac{k^{+}}{\Lambda}]}}\,\int_{x,y,z,x^{\prime},w,w^{\prime}}e^{-ik(w^{\prime}-w)}\,$\ $\longrightarrow$\ $\displaystyle \, {\color{blue}\log{[\frac{k^{+}}{\Lambda}]}}\,\int_{x,y,z,x^{\prime},w,w^{\prime}}e^{-ik(w-w^{\prime})}\,$
\end{itemize}
\needspace{3\baselineskip}\noindent\rule{\linewidth}{0.3pt}\par\noindent\textbf{Lines 94} \quad {\small (in \texttt{align}, lines 93--95)}\par
\begin{itemize}[leftmargin=2.2em,itemsep=1pt,topsep=2pt]
\item[\textbf{l.94}] \textsf{\small [phase convention]} $\displaystyle \, {\color{blue}\log{[\frac{k^{+}}{\Lambda}]}}\,\int_{x,y,z,x^{\prime},w,w^{\prime}}e^{-ik(w^{\prime}-w)}\,$\ $\longrightarrow$\ $\displaystyle \, {\color{blue}\log{[\frac{k^{+}}{\Lambda}]}}\,\int_{x,y,z,x^{\prime},w,w^{\prime}}e^{-ik(w-w^{\prime})}\,$
\end{itemize}
\needspace{3\baselineskip}\noindent\rule{\linewidth}{0.3pt}\par\noindent\textbf{Lines 99} \quad \textsf{Group II $\rangle$ Row I} \quad {\small (in \texttt{align}, lines 98--100)}\par
\begin{itemize}[leftmargin=2.2em,itemsep=1pt,topsep=2pt]
\item[\textbf{l.99}] \textsf{\small [phase convention]} $\displaystyle \int_{z,w,w^{\prime}} e^{-ik(w^{\prime}-w)}\,$\ $\longrightarrow$\ $\displaystyle \int_{z,w,w^{\prime}} e^{-ik(w-w^{\prime})}\,$
\end{itemize}
\needspace{3\baselineskip}\noindent\rule{\linewidth}{0.3pt}\par\noindent\textbf{Lines 102--103} \quad {\small (in \texttt{align}, lines 101--103)}\par
\begin{itemize}[leftmargin=2.2em,itemsep=1pt,topsep=2pt]
\item[\textbf{l.102}] \textsf{\small [phase convention]} $\displaystyle (w^{\prime}-w)$\ $\longrightarrow$\ $\displaystyle (w-w^{\prime})$
\item[\textbf{l.102}] \textsf{\small [index/Wilson-line orientation]} $\displaystyle w')\,\{\rho^{c'}(x'),\rho^{b}(y)\}-U^{bc'}(z)\,U^{ad'}(x')\,U^{c'e'}(y)\, \{\rho^{d'}$\ $\longrightarrow$\ $\displaystyle w')\,\{\rho^{c'}(x'),\rho^{b}(y)\}-U^{c'b}(z)\,U^{ad'}(x')\,U^{c'e'}(y)\, \{\rho^{d'}$
\item[\textbf{l.102}] \textsf{\small [index/Wilson-line orientation]} $\displaystyle f^{acd}\, U^{bc}(z)\,U^{de}(x)\,\rho^{e}(x)$\ $\longrightarrow$\ $\displaystyle f^{acd}\, U^{cb}(z)\,U^{de}(x)\,\rho^{e}(x)$
\item[\textbf{l.102}] \textsf{\small [phase convention]} $\displaystyle (w^{\prime}-w)$\ $\longrightarrow$\ $\displaystyle (w-w^{\prime})$
\item[\textbf{l.102}] \textsf{\small [index/Wilson-line orientation]} $\displaystyle f^{ac'd'}\, U^{bc'}(z)\,U^{d'e'}(x)\,\rho^{e'}(x)$\ $\longrightarrow$\ $\displaystyle f^{ac'd'}\, U^{c'b}(z)\,U^{d'e'}(x)\,\rho^{e'}(x)$
\item[\textbf{l.102}] \textsf{\small [index/Wilson-line orientation]} $\displaystyle (w)\,\{\rho^{c}(x'),\rho^{b}(y)\}-U^{bc}(z)\,U^{ad}(x')\,U^{ce}(y)\, \{\rho^{d}$\ $\longrightarrow$\ $\displaystyle (w)\,\{\rho^{c}(x'),\rho^{b}(y)\}-U^{cb}(z)\,U^{ad}(x')\,U^{ce}(y)\, \{\rho^{d}$
\end{itemize}
\begin{notebox}\textsf{\small Note, l.103:} [In the equations of Group II Row I and of Group III Rows I, II, III, V, VI of this part the indices of the marked $U(z)$, $U(w')$, $U(y')$ were transposed with respect to the Group II and Group III files; corrected. The collected list below already has the correct order.]\end{notebox}
\needspace{3\baselineskip}\noindent\rule{\linewidth}{0.3pt}\par\noindent\textbf{Lines 106--107} \quad {\small (in \texttt{align}, lines 105--107)}\par
\begin{itemize}[leftmargin=2.2em,itemsep=1pt,topsep=2pt]
\item[\textbf{l.106}] \textsf{\small [phase convention]} $\displaystyle \,{\color{blue}\log{[\frac{k^{+}}{\Lambda}]}}\, \int_{x,y,z,x^{\prime},w,w^{\prime}}e^{-ik(w^{\prime}-w)}\,$\ $\longrightarrow$\ $\displaystyle \,{\color{blue}\log{[\frac{k^{+}}{\Lambda}]}}\, \int_{x,y,z,x^{\prime},w,w^{\prime}}e^{-ik(w-w^{\prime})}\,$
\end{itemize}
\begin{notebox}\textsf{\small Note, l.107:} [This intermediate form is superseded by the collected equation ``Group II, Row I'' below, which is the correct one (checked against the Group II file): here each product $\rho\rho$ coming from an anticommutator stands for $\frac{1}{2}\{\rho,\rho\}$, and several $U$ carry transposed indices (those printed in red in the collected equation).]\end{notebox}
\needspace{3\baselineskip}\noindent\rule{\linewidth}{0.3pt}\par\noindent\textbf{Lines 117} \quad \textsf{Group III (with c.c.) $\rangle$ Row I} \quad {\small (in \texttt{align}, lines 116--119)}\par
\begin{itemize}[leftmargin=2.2em,itemsep=1pt,topsep=2pt]
\item[\textbf{l.117}] \textsf{\small [phase convention]} $\displaystyle \int_{z,w,w^{\prime}} e^{-ik(w^{\prime}-w)}\,$\ $\longrightarrow$\ $\displaystyle \int_{z,w,w^{\prime}} e^{-ik(w-w^{\prime})}\,$
\end{itemize}
\needspace{3\baselineskip}\noindent\rule{\linewidth}{0.3pt}\par\noindent\textbf{Lines 121} \quad {\small (in \texttt{align}, lines 120--122)}\par
\begin{itemize}[leftmargin=2.2em,itemsep=1pt,topsep=2pt]
\item[\textbf{l.121}] \textsf{\small [phase convention]} $\displaystyle (w^{\prime}-w)$\ $\longrightarrow$\ $\displaystyle (w-w^{\prime})$
\item[\textbf{l.121}] \textsf{\small [index/Wilson-line orientation]} $\displaystyle (z)\,\rho^{b'}(x')\rho^{d'}(y)-U^{bd'}(z)\,U^{d'e'}(y)U^{ac'}(w')\,\rho^{c'}($\ $\longrightarrow$\ $\displaystyle (z)\,\rho^{b'}(x')\rho^{d'}(y)-U^{d'b}(z)\,U^{d'e'}(y)U^{ac'}(w')\,\rho^{c'}($
\item[\textbf{l.121}] \textsf{\small [index/Wilson-line orientation]} $\displaystyle f^{acd}\, U^{bc}(z)\,U^{de}(x)\,\rho^{e}(x)$\ $\longrightarrow$\ $\displaystyle f^{acd}\, U^{cb}(z)\,U^{de}(x)\,\rho^{e}(x)$
\end{itemize}
\needspace{3\baselineskip}\noindent\rule{\linewidth}{0.3pt}\par\noindent\textbf{Lines 125} \quad {\small (in \texttt{align}, lines 124--126)}\par
\begin{itemize}[leftmargin=2.2em,itemsep=1pt,topsep=2pt]
\item[\textbf{l.125}] \textsf{\small [phase convention]} $\displaystyle (w^{\prime}-w)$\ $\longrightarrow$\ $\displaystyle (w-w^{\prime})$
\item[\textbf{l.125}] \textsf{\small [index/Wilson-line orientation]} $\displaystyle (z)\,\rho^{b'}(x')\rho^{d'}(y)-U^{bd'}(z)\,U^{d'e'}(y)U^{ac'}(w')\,\rho^{c'}($\ $\longrightarrow$\ $\displaystyle (z)\,\rho^{b'}(x')\rho^{d'}(y)-U^{d'b}(z)\,U^{d'e'}(y)U^{ac'}(w')\,\rho^{c'}($
\item[\textbf{l.125}] \textsf{\small [index/Wilson-line orientation]} $\displaystyle f^{acd}\, U^{bc}(z)\,U^{de}(x)\,\rho^{e}(x)$\ $\longrightarrow$\ $\displaystyle f^{acd}\, U^{cb}(z)\,U^{de}(x)\,\rho^{e}(x)$
\end{itemize}
\needspace{3\baselineskip}\noindent\rule{\linewidth}{0.3pt}\par\noindent\textbf{Lines 129} \quad \textsf{Group III (with c.c.) $\rangle$ Row II} \quad {\small (in \texttt{align}, lines 128--130)}\par
\begin{itemize}[leftmargin=2.2em,itemsep=1pt,topsep=2pt]
\item[\textbf{l.129}] \textsf{\small [phase convention]} $\displaystyle \int_{z,w,w^{\prime}} e^{-ik(w^{\prime}-w)}\,$\ $\longrightarrow$\ $\displaystyle \int_{z,w,w^{\prime}} e^{-ik(w-w^{\prime})}\,$
\end{itemize}
\needspace{3\baselineskip}\noindent\rule{\linewidth}{0.3pt}\par\noindent\textbf{Lines 132} \quad {\small (in \texttt{align}, lines 131--133)}\par
\begin{itemize}[leftmargin=2.2em,itemsep=1pt,topsep=2pt]
\item[\textbf{l.132}] \textsf{\small [phase convention]} $\displaystyle (w^{\prime}-w)$\ $\longrightarrow$\ $\displaystyle (w-w^{\prime})$
\item[\textbf{l.132}] \textsf{\small [index/Wilson-line orientation]} $\displaystyle [U^{c'd'}(y)\,U^{ab'}(x')\,U^{bc'}(z)\,\rho^{d'}(y)\rho^{b'}(x')$\ $\longrightarrow$\ $\displaystyle [U^{c'd'}(y)\,U^{ab'}(x')\,U^{c'b}(z)\,\rho^{d'}(y)\rho^{b'}(x')$
\item[\textbf{l.132}] \textsf{\small [index/Wilson-line orientation]} $\displaystyle f^{acd}\, U^{bc}(z)\,U^{de}(x)\,\rho^{e}(x)$\ $\longrightarrow$\ $\displaystyle f^{acd}\, U^{cb}(z)\,U^{de}(x)\,\rho^{e}(x)$
\end{itemize}
\needspace{3\baselineskip}\noindent\rule{\linewidth}{0.3pt}\par\noindent\textbf{Lines 136} \quad {\small (in \texttt{align}, lines 135--137)}\par
\begin{itemize}[leftmargin=2.2em,itemsep=1pt,topsep=2pt]
\item[\textbf{l.136}] \textsf{\small [phase convention]} $\displaystyle (w^{\prime}-w)$\ $\longrightarrow$\ $\displaystyle (w-w^{\prime})$
\item[\textbf{l.136}] \textsf{\small [index/Wilson-line orientation]} $\displaystyle [U^{c'd'}(y)\,U^{ab'}(x')\,U^{bc'}(z)\,\rho^{d'}(y)\rho^{b'}(x')$\ $\longrightarrow$\ $\displaystyle [U^{c'd'}(y)\,U^{ab'}(x')\,U^{c'b}(z)\,\rho^{d'}(y)\rho^{b'}(x')$
\item[\textbf{l.136}] \textsf{\small [index/Wilson-line orientation]} $\displaystyle f^{acd}\, U^{bc}(z)\,U^{de}(x)\,\rho^{e}(x)$\ $\longrightarrow$\ $\displaystyle f^{acd}\, U^{cb}(z)\,U^{de}(x)\,\rho^{e}(x)$
\end{itemize}
\needspace{3\baselineskip}\noindent\rule{\linewidth}{0.3pt}\par\noindent\textbf{Lines 140} \quad \textsf{Group III (with c.c.) $\rangle$ Row III} \quad {\small (in \texttt{align}, lines 139--141)}\par
\begin{itemize}[leftmargin=2.2em,itemsep=1pt,topsep=2pt]
\item[\textbf{l.140}] \textsf{\small [phase convention]} $\displaystyle \int_{z,w,w^{\prime}} e^{-ik(w^{\prime}-w)}\,\int_{y^{\prime}}$\ $\longrightarrow$\ $\displaystyle \int_{z,w,w^{\prime}} e^{-ik(w-w^{\prime})}\,\int_{y^{\prime}}$
\end{itemize}
\needspace{3\baselineskip}\noindent\rule{\linewidth}{0.3pt}\par\noindent\textbf{Lines 145} \quad {\small (in \texttt{align}, lines 144--146)}\par
\begin{itemize}[leftmargin=2.2em,itemsep=1pt,topsep=2pt]
\item[\textbf{l.145}] \textsf{\small [phase convention]} $\displaystyle \, \int_{w,x',y',z,y,x}\, e^{-ik\cdot (y'-w)}\,\frac{(y'-z)\cdot (x-z)\,(x'-y')\cdot (y-w)}{(y'-z)^{2}\,(x'-y')^{2}\,(y-w)^{2}\,(x-z)^{2}}$\ $\longrightarrow$\ $\displaystyle \, \int_{w,x',y',z,y,x}\, e^{-ik\cdot (w-y')}\,\frac{(y'-z)\cdot (x-z)\,(x'-y')\cdot (y-w)}{(y'-z)^{2}\,(x'-y')^{2}\,(y-w)^{2}\,(x-z)^{2}}$
\item[\textbf{l.145}] \textsf{\small [index/Wilson-line orientation]} $\displaystyle [U^{e'c'}(y')\,U^{d'b}(z)\,\rho^{e'}(x')$\ $\longrightarrow$\ $\displaystyle [U^{c'e'}(y')\,U^{d'b}(z)\,\rho^{e'}(x')$
\item[\textbf{l.145}] \textsf{\small [index/Wilson-line orientation]} $\displaystyle U^{e'c'}(y')\,U^{d'b}(z)\,\rho^{e'}(x')-U^{bd'}(z)\,U^{c'e'}(x')\,\rho^{e'}(x')]$\ $\longrightarrow$\ $\displaystyle U^{c'e'}(y')\,U^{d'b}(z)\,\rho^{e'}(x')-U^{d'b}(z)\,U^{c'e'}(x')\,\rho^{e'}(x')]$
\item[\textbf{l.145}] \textsf{\small [index/Wilson-line orientation]} $\displaystyle (w)\, \{\rho^{c}(y),\rho^{b}(x)\}-U^{bc}(z)\,U^{ad}(y)\,U^{ce}(x)\,\{\rho^{d}$\ $\longrightarrow$\ $\displaystyle (w)\, \{\rho^{c}(y),\rho^{b}(x)\}-U^{cb}(z)\,U^{ad}(y)\,U^{ce}(x)\,\{\rho^{d}$
\end{itemize}
\needspace{3\baselineskip}\noindent\rule{\linewidth}{0.3pt}\par\noindent\textbf{Lines 149--150} \quad {\small (in \texttt{align}, lines 148--150)}\par
\begin{itemize}[leftmargin=2.2em,itemsep=1pt,topsep=2pt]
\item[\textbf{l.149}] \textsf{\small [phase convention]} $\displaystyle \, \int_{w,x',w',z,y,x}\, e^{-ik\cdot (w'-w)}\,\mathcal{K}(w'-z$\ $\longrightarrow$\ $\displaystyle \, \int_{w,x',w',z,y,x}\, e^{-ik\cdot (w-w')}\,\mathcal{K}(w'-z$
\item[\textbf{l.149}] \textsf{\small [index/Wilson-line orientation]} $\displaystyle [U^{e'c'}(w')\,U^{d'b}(z)\,\rho^{e'}(x')$\ $\longrightarrow$\ $\displaystyle [U^{c'e'}(w')\,U^{d'b}(z)\,\rho^{e'}(x')$
\item[\textbf{l.149}] \textsf{\small [index/Wilson-line orientation]} $\displaystyle U^{e'c'}(w')\,U^{d'b}(z)\,\rho^{e'}(x')-U^{bd'}(z)\,U^{c'e'}(x')\,\rho^{e'}(x')]$\ $\longrightarrow$\ $\displaystyle U^{c'e'}(w')\,U^{d'b}(z)\,\rho^{e'}(x')-U^{d'b}(z)\,U^{c'e'}(x')\,\rho^{e'}(x')]$
\item[\textbf{l.149}] \textsf{\small [index/Wilson-line orientation]} $\displaystyle (w)\, \{\rho^{c}(x),\rho^{b}(y)\}-U^{bc}(z)\,U^{ad}(x)\,U^{ce}(y)\,\{\rho^{d}$\ $\longrightarrow$\ $\displaystyle (w)\, \{\rho^{c}(x),\rho^{b}(y)\}-U^{cb}(z)\,U^{ad}(x)\,U^{ce}(y)\,\{\rho^{d}$
\end{itemize}
\begin{notebox}\textsf{\small Note, l.150:} [The explicit logarithm of this row (and of Rows V, VI) is $\log\frac{\vee}{\Lambda}$ in the Group III file; this is its $\log\frac{k^{+}}{\Lambda}$ half. The complete $\log\frac{\vee}{k^{+}}$ of these rows is extracted with the author's split of the $p^{+}$ integral (BFKL file, Other\_terms). The sign change of Row III in the Group III file concerns only its first line; the expression here is unchanged.]\end{notebox}
\needspace{3\baselineskip}\noindent\rule{\linewidth}{0.3pt}\par\noindent\textbf{Lines 157} \quad \textsf{Group III (with c.c.) $\rangle$ Row V} \quad {\small (in \texttt{align}, lines 156--158)}\par
\begin{itemize}[leftmargin=2.2em,itemsep=1pt,topsep=2pt]
\item[\textbf{l.157}] \textsf{\small [phase convention]} $\displaystyle \int_{z,w,w^{\prime}} e^{-ik(w^{\prime}-w)}\,\int_{y^{\prime}}$\ $\longrightarrow$\ $\displaystyle \int_{z,w,w^{\prime}} e^{-ik(w-w^{\prime})}\,\int_{y^{\prime}}$
\end{itemize}
\needspace{3\baselineskip}\noindent\rule{\linewidth}{0.3pt}\par\noindent\textbf{Lines 160} \quad {\small (in \texttt{align}, lines 159--161)}\par
\begin{itemize}[leftmargin=2.2em,itemsep=1pt,topsep=2pt]
\item[\textbf{l.160}] \textsf{\small [phase convention]} $\displaystyle \, \int_{w,w',z,y,x,x'}\, e^{-ik\cdot (w'-w)}\,\frac{(w'-z)\cdot (y-z)\,(x'-w')\cdot (x-w)}{(w'-z)^{2}\,(x'-w')^{2}\,(x-w)^{2}\,(y-z)^{2}}$\ $\longrightarrow$\ $\displaystyle \, \int_{w,w',z,y,x,x'}\, e^{-ik\cdot (w-w')}\,\frac{(w'-z)\cdot (y-z)\,(x'-w')\cdot (x-w)}{(w'-z)^{2}\,(x'-w')^{2}\,(x-w)^{2}\,(y-z)^{2}}$
\item[\textbf{l.160}] \textsf{\small [index/Wilson-line orientation]} $\displaystyle [U^{e'c'}(w')\,\rho^{e'}(x')$\ $\longrightarrow$\ $\displaystyle [U^{c'e'}(w')\,\rho^{e'}(x')$
\end{itemize}
\needspace{3\baselineskip}\noindent\rule{\linewidth}{0.3pt}\par\noindent\textbf{Lines 164} \quad {\small (in \texttt{align}, lines 163--165)}\par
\begin{itemize}[leftmargin=2.2em,itemsep=1pt,topsep=2pt]
\item[\textbf{l.164}] \textsf{\small [phase convention]} $\displaystyle \, \int_{w,w',z,y,x,x'}\, e^{-ik\cdot (w'-w)}\,\mathcal{K}(x'-w$\ $\longrightarrow$\ $\displaystyle \, \int_{w,w',z,y,x,x'}\, e^{-ik\cdot (w-w')}\,\mathcal{K}(x'-w$
\item[\textbf{l.164}] \textsf{\small [index/Wilson-line orientation]} $\displaystyle [U^{e'c'}(w')\,\rho^{e'}(x')$\ $\longrightarrow$\ $\displaystyle [U^{c'e'}(w')\,\rho^{e'}(x')$
\end{itemize}
\needspace{3\baselineskip}\noindent\rule{\linewidth}{0.3pt}\par\noindent\textbf{Lines 168} \quad \textsf{Group III (with c.c.) $\rangle$ Row VI} \quad {\small (in \texttt{align}, lines 167--169)}\par
\begin{itemize}[leftmargin=2.2em,itemsep=1pt,topsep=2pt]
\item[\textbf{l.168}] \textsf{\small [phase convention]} $\displaystyle \int_{z,w,w^{\prime}} e^{-ik(w^{\prime}-w)}\, \int_{y^{\prime}}$\ $\longrightarrow$\ $\displaystyle \int_{z,w,w^{\prime}} e^{-ik(w-w^{\prime})}\, \int_{y^{\prime}}$
\end{itemize}
\needspace{3\baselineskip}\noindent\rule{\linewidth}{0.3pt}\par\noindent\textbf{Lines 171} \quad {\small (in \texttt{align}, lines 170--172)}\par
\begin{itemize}[leftmargin=2.2em,itemsep=1pt,topsep=2pt]
\item[\textbf{l.171}] \textsf{\small [phase convention]} $\displaystyle \, \int_{w,w',z,y,x,x'}\, e^{-ik\cdot (w'-w)}\,\frac{(w'-z)\cdot (y-z)\,(x'-w')\cdot (x-w)}{(w'-z)^{2}\,(x'-w')^{2}\,(x-w)^{2}\,(y-z)^{2}}$\ $\longrightarrow$\ $\displaystyle \, \int_{w,w',z,y,x,x'}\, e^{-ik\cdot (w-w')}\,\frac{(w'-z)\cdot (y-z)\,(x'-w')\cdot (x-w)}{(w'-z)^{2}\,(x'-w')^{2}\,(x-w)^{2}\,(y-z)^{2}}$
\item[\textbf{l.171}] \textsf{\small [index/Wilson-line orientation]} $\displaystyle [U^{e'c'}(w')\,U^{d'b}(z)\,\rho^{e'}(x')$\ $\longrightarrow$\ $\displaystyle [U^{c'e'}(w')\,U^{d'b}(z)\,\rho^{e'}(x')$
\item[\textbf{l.171}] \textsf{\small [index/Wilson-line orientation]} $\displaystyle U^{e'c'}(w')\,U^{d'b}(z)\,\rho^{e'}(x')-U^{bd'}(z)\,U^{c'e'}(x')\,\rho^{e'}(x')]$\ $\longrightarrow$\ $\displaystyle U^{c'e'}(w')\,U^{d'b}(z)\,\rho^{e'}(x')-U^{d'b}(z)\,U^{c'e'}(x')\,\rho^{e'}(x')]$
\item[\textbf{l.171}] \textsf{\small [index/Wilson-line orientation]} $\displaystyle [U^{bc}(z)\,U^{ad}(x)\,U^{ce}(y)\,\rho^{d}($\ $\longrightarrow$\ $\displaystyle [U^{cb}(z)\,U^{ad}(x)\,U^{ce}(y)\,\rho^{d}($
\end{itemize}
\needspace{3\baselineskip}\noindent\rule{\linewidth}{0.3pt}\par\noindent\textbf{Lines 175--178} \quad {\small (in \texttt{align}, lines 174--176)}\par
\begin{itemize}[leftmargin=2.2em,itemsep=1pt,topsep=2pt]
\item[\textbf{l.175}] \textsf{\small [phase convention]} $\displaystyle \, \int_{w,w',z,y,x,x'}\, e^{-ik\cdot (w'-w)}\,\mathcal{K}(x'-w$\ $\longrightarrow$\ $\displaystyle \, \int_{w,w',z,y,x,x'}\, e^{-ik\cdot (w-w')}\,\mathcal{K}(x'-w$
\item[\textbf{l.175}] \textsf{\small [index/Wilson-line orientation]} $\displaystyle [U^{e'c'}(w')\,U^{d'b}(z)\,\rho^{e'}(x')$\ $\longrightarrow$\ $\displaystyle [U^{c'e'}(w')\,U^{d'b}(z)\,\rho^{e'}(x')$
\item[\textbf{l.175}] \textsf{\small [index/Wilson-line orientation]} $\displaystyle U^{e'c'}(w')\,U^{d'b}(z)\,\rho^{e'}(x')-U^{bd'}(z)\,U^{c'e'}(x')\,\rho^{e'}(x')]$\ $\longrightarrow$\ $\displaystyle U^{c'e'}(w')\,U^{d'b}(z)\,\rho^{e'}(x')-U^{d'b}(z)\,U^{c'e'}(x')\,\rho^{e'}(x')]$
\item[\textbf{l.175}] \textsf{\small [index/Wilson-line orientation]} $\displaystyle [U^{bc}(z)\,U^{ad}(x)\,U^{ce}(y)\,\rho^{d}($\ $\longrightarrow$\ $\displaystyle [U^{cb}(z)\,U^{ad}(x)\,U^{ce}(y)\,\rho^{d}($
\end{itemize}
\begin{notebox}\textsf{\small Note, l.178:} [These are the three-$\rho$ remainder of the reordered four-$\rho$ terms (Eq.~(4rho) of the four-$\rho$ file); their logarithm is $\log\frac{\vee}{\Lambda}$, so the same expressions occur with $\log\frac{\vee}{k^{+}}$. No c.c. is to be added to $N_{a}$, $N_{b}$.]\end{notebox}
\needspace{3\baselineskip}\noindent\rule{\linewidth}{0.3pt}\par\noindent\textbf{Lines 180} \quad \textsf{All three $\rho$ term with c.c. $\rangle$ From four $\rho$'s} \quad {\small (in \texttt{align}, lines 179--184)}\par
\begin{itemize}[leftmargin=2.2em,itemsep=1pt,topsep=2pt]
\item[\textbf{l.180}] \textsf{\small [phase convention]} $\displaystyle \,{\color{blue}\log{[\frac{k^{+}}{\Lambda}]}}\,\int_{x,y,z,x^{\prime},w,w^{\prime}}e^{-ik(w^{\prime}-w)}\,$\ $\longrightarrow$\ $\displaystyle \,{\color{blue}\log{[\frac{k^{+}}{\Lambda}]}}\,\int_{x,y,z,x^{\prime},w,w^{\prime}}e^{-ik(w-w^{\prime})}\,$
\end{itemize}
\needspace{3\baselineskip}\noindent\rule{\linewidth}{0.3pt}\par\noindent\textbf{Lines 186} \quad {\small (in \texttt{align}, lines 185--190)}\par
\begin{itemize}[leftmargin=2.2em,itemsep=1pt,topsep=2pt]
\item[\textbf{l.186}] \textsf{\small [phase convention]} $\displaystyle \,{\color{blue}\log{[\frac{k^{+}}{\Lambda}]}}\,\int_{x,y,z,x^{\prime},w,w^{\prime}}e^{-ik(w^{\prime}-w)}\,$\ $\longrightarrow$\ $\displaystyle \,{\color{blue}\log{[\frac{k^{+}}{\Lambda}]}}\,\int_{x,y,z,x^{\prime},w,w^{\prime}}e^{-ik(w-w^{\prime})}\,$
\end{itemize}
\needspace{3\baselineskip}\noindent\rule{\linewidth}{0.3pt}\par\noindent\textbf{Lines 193} \quad \textsf{All three $\rho$ term with c.c. $\rangle$ Genuine terms}\par
\begin{notebox}\textsf{\small Note, l.193:} [Checked: every equation of this list equals the corresponding row of the checked cross-section files plus its c.c. (Group II: as it is), as ordered operators; for Group I Row II with the ordering $A^{(1)}N_{2}$ (note in Row II above), and for Group I Row III in the form with $U^{bc}(w)$ ($\beta$, $\beta'$ terms; open point noted there). Under the change of the phase the totally symmetrised three-$\rho$ part and the one-$\rho$ remainder of every equation change sign, the two-$\rho$ remainder is unchanged.]\end{notebox}
\needspace{3\baselineskip}\noindent\rule{\linewidth}{0.3pt}\par\noindent\textbf{Lines 198} \quad {\small (in \texttt{align}, lines 197--200)}\par
\begin{itemize}[leftmargin=2.2em,itemsep=1pt,topsep=2pt]
\item[\textbf{l.198}] \textsf{\small [phase convention]} $\displaystyle \, {\color{blue}\log{[\frac{k^{+}}{\Lambda}]}}\,\int_{x,y,z,x^{\prime},w,w^{\prime}}e^{-ik(w^{\prime}-w)}\,$\ $\longrightarrow$\ $\displaystyle \, {\color{blue}\log{[\frac{k^{+}}{\Lambda}]}}\,\int_{x,y,z,x^{\prime},w,w^{\prime}}e^{-ik(w-w^{\prime})}\,$
\end{itemize}
\needspace{3\baselineskip}\noindent\rule{\linewidth}{0.3pt}\par\noindent\textbf{Lines 205} \quad {\small (in \texttt{align}, lines 204--206)}\par
\begin{itemize}[leftmargin=2.2em,itemsep=1pt,topsep=2pt]
\item[\textbf{l.205}] \textsf{\small [phase convention]} $\displaystyle \, {\color{blue}\log{[\frac{k^{+}}{\Lambda}]}}\, \int_{x,y,z,x^{\prime},w,w^{\prime}}e^{-ik(w^{\prime}-w)}\,$\ $\longrightarrow$\ $\displaystyle \, {\color{blue}\log{[\frac{k^{+}}{\Lambda}]}}\, \int_{x,y,z,x^{\prime},w,w^{\prime}}e^{-ik(w-w^{\prime})}\,$
\item[\textbf{l.205}] \textsf{\small [typo/LaTeX]} \emph{added:} $\displaystyle \Bigg]$\quad{\color{gray}\small(in context: }$\displaystyle w)\,U^{be}(y)\,U^{ch}(x')\,{\color{green}\beta'}\Bigg]${\color{gray}\small)}
\end{itemize}
\needspace{3\baselineskip}\noindent\rule{\linewidth}{0.3pt}\par\noindent\textbf{Lines 209} \quad {\small (in \texttt{align}, lines 208--211)}\par
\begin{itemize}[leftmargin=2.2em,itemsep=1pt,topsep=2pt]
\item[\textbf{l.209}] \textsf{\small [phase convention]} $\displaystyle \, {\color{blue}\log{[\frac{k^{+}}{\Lambda}]}}\, \int_{x,y,z,x^{\prime},w,w^{\prime}}e^{-ik(w^{\prime}-w)}\,$\ $\longrightarrow$\ $\displaystyle \, {\color{blue}\log{[\frac{k^{+}}{\Lambda}]}}\, \int_{x,y,z,x^{\prime},w,w^{\prime}}e^{-ik(w-w^{\prime})}\,$
\end{itemize}
\needspace{3\baselineskip}\noindent\rule{\linewidth}{0.3pt}\par\noindent\textbf{Lines 214} \quad {\small (in \texttt{align}, lines 213--219)}\par
\begin{itemize}[leftmargin=2.2em,itemsep=1pt,topsep=2pt]
\item[\textbf{l.214}] \textsf{\small [phase convention]} $\displaystyle \, {\color{blue}\log{[\frac{k^{+}}{\Lambda}]}}\, \int_{x,y,z,x^{\prime},w,w^{\prime}}e^{-ik(w^{\prime}-w)}\,$\ $\longrightarrow$\ $\displaystyle \, {\color{blue}\log{[\frac{k^{+}}{\Lambda}]}}\, \int_{x,y,z,x^{\prime},w,w^{\prime}}e^{-ik(w-w^{\prime})}\,$
\end{itemize}
\needspace{3\baselineskip}\noindent\rule{\linewidth}{0.3pt}\par\noindent\textbf{Lines 222} \quad {\small (in \texttt{align}, lines 221--223)}\par
\begin{itemize}[leftmargin=2.2em,itemsep=1pt,topsep=2pt]
\item[\textbf{l.222}] \textsf{\small [phase convention]} $\displaystyle \, {\color{blue}\log{[\frac{k^{+}}{\Lambda}]}}\,\int_{x,y,z,x^{\prime},w,w^{\prime}}e^{-ik(w^{\prime}-w)}\,$\ $\longrightarrow$\ $\displaystyle \, {\color{blue}\log{[\frac{k^{+}}{\Lambda}]}}\,\int_{x,y,z,x^{\prime},w,w^{\prime}}e^{-ik(w-w^{\prime})}\,$
\end{itemize}
\needspace{3\baselineskip}\noindent\rule{\linewidth}{0.3pt}\par\noindent\textbf{Lines 225} \quad {\small (in \texttt{align}, lines 224--228)}\par
\begin{itemize}[leftmargin=2.2em,itemsep=1pt,topsep=2pt]
\item[\textbf{l.225}] \textsf{\small [phase convention]} $\displaystyle \, {\color{blue}\log{[\frac{k^{+}}{\Lambda}]}}\, \int_{x,y,z,x^{\prime},w,w^{\prime}}e^{-ik(w^{\prime}-w)}\,$\ $\longrightarrow$\ $\displaystyle \, {\color{blue}\log{[\frac{k^{+}}{\Lambda}]}}\, \int_{x,y,z,x^{\prime},w,w^{\prime}}e^{-ik(w-w^{\prime})}\,$
\end{itemize}
\needspace{3\baselineskip}\noindent\rule{\linewidth}{0.3pt}\par\noindent\textbf{Lines 231} \quad {\small (in \texttt{align}, lines 230--232)}\par
\begin{itemize}[leftmargin=2.2em,itemsep=1pt,topsep=2pt]
\item[\textbf{l.231}] \textsf{\small [phase convention]} $\displaystyle \, {\color{blue}\log{[\frac{k^{+}}{\Lambda}]}}\, \int_{x,y,z,x^{\prime},w,w^{\prime}}e^{-ik(w^{\prime}-w)}\,$\ $\longrightarrow$\ $\displaystyle \, {\color{blue}\log{[\frac{k^{+}}{\Lambda}]}}\, \int_{x,y,z,x^{\prime},w,w^{\prime}}e^{-ik(w-w^{\prime})}\,$
\end{itemize}
\needspace{3\baselineskip}\noindent\rule{\linewidth}{0.3pt}\par\noindent\textbf{Lines 235} \quad {\small (in \texttt{align}, lines 234--236)}\par
\begin{itemize}[leftmargin=2.2em,itemsep=1pt,topsep=2pt]
\item[\textbf{l.235}] \textsf{\small [phase convention]} $\displaystyle \, {\color{blue}\log{[\frac{k^{+}}{\Lambda}]}}\, \int_{x,y,z,x^{\prime},w,w^{\prime}}e^{-ik(w^{\prime}-w)}\,$\ $\longrightarrow$\ $\displaystyle \, {\color{blue}\log{[\frac{k^{+}}{\Lambda}]}}\, \int_{x,y,z,x^{\prime},w,w^{\prime}}e^{-ik(w-w^{\prime})}\,$
\end{itemize}
\needspace{3\baselineskip}\noindent\rule{\linewidth}{0.3pt}\par\noindent\textbf{Lines 239} \quad {\small (in \texttt{align}, lines 238--240)}\par
\begin{itemize}[leftmargin=2.2em,itemsep=1pt,topsep=2pt]
\item[\textbf{l.239}] \textsf{\small [phase convention]} $\displaystyle \, {\color{blue}\log{[\frac{k^{+}}{\Lambda}]}}\,\int_{x,y,z,x^{\prime},w,w^{\prime}}e^{-ik(w^{\prime}-w)}\,$\ $\longrightarrow$\ $\displaystyle \, {\color{blue}\log{[\frac{k^{+}}{\Lambda}]}}\,\int_{x,y,z,x^{\prime},w,w^{\prime}}e^{-ik(w-w^{\prime})}\,$
\end{itemize}
\needspace{3\baselineskip}\noindent\rule{\linewidth}{0.3pt}\par\noindent\textbf{Lines 243--244} \quad {\small (in \texttt{align}, lines 242--244)}\par
\begin{itemize}[leftmargin=2.2em,itemsep=1pt,topsep=2pt]
\item[\textbf{l.243}] \textsf{\small [phase convention]} $\displaystyle \, {\color{blue}\log{[\frac{k^{+}}{\Lambda}]}}\, \int_{x,y,z,x^{\prime},w,w^{\prime}}e^{-ik(w^{\prime}-w)}\,$\ $\longrightarrow$\ $\displaystyle \, {\color{blue}\log{[\frac{k^{+}}{\Lambda}]}}\, \int_{x,y,z,x^{\prime},w,w^{\prime}}e^{-ik(w-w^{\prime})}\,$
\item[\textbf{l.244}] \textsf{\small [term removed/replaced]} $\displaystyle \-\mathcal{K}(x-w)\cdot\mathcal{K}(x'-w')\,\mathcal{K}$\ $\longrightarrow$\ $\displaystyle -\mathcal{K}(x-w)\cdot\mathcal{K}(x'-w')\,\mathcal{K}$
\item[\textbf{l.244}] \textsf{\small [term removed/replaced]} $\displaystyle \-\mathcal{K}(x-w)\cdot\mathcal{K}(x'-w')\,\mathcal{K}$\ $\longrightarrow$\ $\displaystyle -\mathcal{K}(x-w)\cdot\mathcal{K}(x'-w')\,\mathcal{K}$
\end{itemize}
\needspace{3\baselineskip}\noindent\rule{\linewidth}{0.3pt}\par\noindent\textbf{Lines 247} \quad {\small (in \texttt{align}, lines 246--248)}\par
\begin{itemize}[leftmargin=2.2em,itemsep=1pt,topsep=2pt]
\item[\textbf{l.247}] \textsf{\small [phase convention]} $\displaystyle \, {\color{blue}\log{[\frac{k^{+}}{\Lambda}]}}\,\int_{x,y,z,x^{\prime},w,w^{\prime}}e^{-ik(w^{\prime}-w)}\,$\ $\longrightarrow$\ $\displaystyle \, {\color{blue}\log{[\frac{k^{+}}{\Lambda}]}}\,\int_{x,y,z,x^{\prime},w,w^{\prime}}e^{-ik(w-w^{\prime})}\,$
\end{itemize}
\needspace{3\baselineskip}\noindent\rule{\linewidth}{0.3pt}\par\noindent\textbf{Lines 251--253} \quad {\small (in \texttt{align}, lines 250--252)}\par
\begin{itemize}[leftmargin=2.2em,itemsep=1pt,topsep=2pt]
\item[\textbf{l.251}] \textsf{\small [phase convention]} $\displaystyle \, {\color{blue}\log{[\frac{k^{+}}{\Lambda}]}}\,\int_{x,y,z,x^{\prime},w,w^{\prime}}e^{-ik(w^{\prime}-w)}\,$\ $\longrightarrow$\ $\displaystyle \, {\color{blue}\log{[\frac{k^{+}}{\Lambda}]}}\,\int_{x,y,z,x^{\prime},w,w^{\prime}}e^{-ik(w-w^{\prime})}\,$
\end{itemize}
\begin{notebox}\textsf{\small Note, l.253:} [What is shown (checked for each of the twelve equations): the \emph{totally symmetrised} three-$\rho$ part of each bracket vanishes identically, for any function of $(w,w')$ in place of the phase; the twelve equations are the regrouping of the collected list by kernel and by number of Wilson lines. As ordered operators the brackets do not vanish: each leaves a two-$\rho$ and a one-$\rho$ remainder from the commutators, so ``$=0$'' means ``$=0$ up to terms with fewer $\rho$'s''; the two-$\rho$ remainders are what the JIMWLK file collects. With the ordering $A^{(1)}N_{2}$ nothing has to be added to the list, and the phase changes the sign of all three-$\rho$ terms together. In the $\beta$ and $\beta'$ equations the kernel in front is misprinted; corrected to $\mathcal{K}(w-z)\cdot\mathcal{K}(y-z)$ and $\mathcal{K}(w'-z)\cdot\mathcal{K}(y-z)$.]\end{notebox}
\needspace{3\baselineskip}\noindent\rule{\linewidth}{0.3pt}\par\noindent\textbf{Lines 257} \quad \textsf{Cancellation of all three $\rho$'s $\rangle$ $\alpha$ terms} \quad {\small (in \texttt{align}, lines 256--259)}\par
\begin{itemize}[leftmargin=2.2em,itemsep=1pt,topsep=2pt]
\item[\textbf{l.257}] \textsf{\small [phase convention]} $\displaystyle \,{\color{blue}\log{[\frac{k^{+}}{\Lambda}]}}\,\int_{x,y,z,x',w,w'}e^{-ik(w'-w)}\,\mathcal{K}(x'-w$\ $\longrightarrow$\ $\displaystyle \,{\color{blue}\log{[\frac{k^{+}}{\Lambda}]}}\,\int_{x,y,z,x',w,w'}e^{-ik(w-w')}\,\mathcal{K}(x'-w$
\end{itemize}
\needspace{3\baselineskip}\noindent\rule{\linewidth}{0.3pt}\par\noindent\textbf{Lines 262} \quad {\small (in \texttt{align}, lines 261--265)}\par
\begin{itemize}[leftmargin=2.2em,itemsep=1pt,topsep=2pt]
\item[\textbf{l.262}] \textsf{\small [phase convention]} $\displaystyle \,{\color{blue}\log{[\frac{k^{+}}{\Lambda}]}}\,\int_{x,y,z,x',w,w'}e^{-ik(w'-w)}\,\mathcal{K}(x'-w$\ $\longrightarrow$\ $\displaystyle \,{\color{blue}\log{[\frac{k^{+}}{\Lambda}]}}\,\int_{x,y,z,x',w,w'}e^{-ik(w-w')}\,\mathcal{K}(x'-w$
\end{itemize}
\needspace{3\baselineskip}\noindent\rule{\linewidth}{0.3pt}\par\noindent\textbf{Lines 268} \quad {\small (in \texttt{align}, lines 267--271)}\par
\begin{itemize}[leftmargin=2.2em,itemsep=1pt,topsep=2pt]
\item[\textbf{l.268}] \textsf{\small [phase convention]} $\displaystyle \,{\color{blue}\log{[\frac{k^{+}}{\Lambda}]}}\,\int_{x,y,z,x',w,w'}e^{-ik(w'-w)}\,\mathcal{K}(x'-w$\ $\longrightarrow$\ $\displaystyle \,{\color{blue}\log{[\frac{k^{+}}{\Lambda}]}}\,\int_{x,y,z,x',w,w'}e^{-ik(w-w')}\,\mathcal{K}(x'-w$
\end{itemize}
\needspace{3\baselineskip}\noindent\rule{\linewidth}{0.3pt}\par\noindent\textbf{Lines 275} \quad \textsf{Cancellation of all three $\rho$'s $\rangle$ $\alpha'$ terms} \quad {\small (in \texttt{align}, lines 274--277)}\par
\begin{itemize}[leftmargin=2.2em,itemsep=1pt,topsep=2pt]
\item[\textbf{l.275}] \textsf{\small [phase convention]} $\displaystyle \,{\color{blue}\log{[\frac{k^{+}}{\Lambda}]}}\,\int_{x,y,z,x',w,w'}e^{-ik(w'-w)}\,\mathcal{K}(x'-w$\ $\longrightarrow$\ $\displaystyle \,{\color{blue}\log{[\frac{k^{+}}{\Lambda}]}}\,\int_{x,y,z,x',w,w'}e^{-ik(w-w')}\,\mathcal{K}(x'-w$
\end{itemize}
\needspace{3\baselineskip}\noindent\rule{\linewidth}{0.3pt}\par\noindent\textbf{Lines 280} \quad {\small (in \texttt{align}, lines 279--283)}\par
\begin{itemize}[leftmargin=2.2em,itemsep=1pt,topsep=2pt]
\item[\textbf{l.280}] \textsf{\small [phase convention]} $\displaystyle \,{\color{blue}\log{[\frac{k^{+}}{\Lambda}]}}\,\int_{x,y,z,x',w,w'}e^{-ik(w'-w)}\,\mathcal{K}(x'-w$\ $\longrightarrow$\ $\displaystyle \,{\color{blue}\log{[\frac{k^{+}}{\Lambda}]}}\,\int_{x,y,z,x',w,w'}e^{-ik(w-w')}\,\mathcal{K}(x'-w$
\end{itemize}
\needspace{3\baselineskip}\noindent\rule{\linewidth}{0.3pt}\par\noindent\textbf{Lines 286} \quad {\small (in \texttt{align}, lines 285--287)}\par
\begin{itemize}[leftmargin=2.2em,itemsep=1pt,topsep=2pt]
\item[\textbf{l.286}] \textsf{\small [phase convention]} $\displaystyle \,{\color{blue}\log{[\frac{k^{+}}{\Lambda}]}}\,\int_{x,y,z,x',w,w'}e^{-ik(w'-w)}\,\mathcal{K}(x'-w$\ $\longrightarrow$\ $\displaystyle \,{\color{blue}\log{[\frac{k^{+}}{\Lambda}]}}\,\int_{x,y,z,x',w,w'}e^{-ik(w-w')}\,\mathcal{K}(x'-w$
\end{itemize}
\needspace{3\baselineskip}\noindent\rule{\linewidth}{0.3pt}\par\noindent\textbf{Lines 291} \quad \textsf{Cancellation of all three $\rho$'s $\rangle$ $\beta$ terms} \quad {\small (in \texttt{align}, lines 290--292)}\par
\begin{itemize}[leftmargin=2.2em,itemsep=1pt,topsep=2pt]
\item[\textbf{l.291}] \textsf{\small [phase convention]} $\displaystyle \,{\color{blue}\log{[\frac{k^{+}}{\Lambda}]}}\,\int_{x,y,z,x',w,w'}e^{-ik(w'-w)}\,\mathcal{K}(x'-w$\ $\longrightarrow$\ $\displaystyle \,{\color{blue}\log{[\frac{k^{+}}{\Lambda}]}}\,\int_{x,y,z,x',w,w'}e^{-ik(w-w')}\,\mathcal{K}(x'-w$
\item[\textbf{l.291}] \textsf{\small [term removed/replaced]} $\displaystyle \,\mathcal{K}(x'-w')\cdot\mathcal{K}(x-w)\,\mathcal{K}(x'-z)\cdot\mathcal{K}(y-z)$\ $\longrightarrow$\ $\displaystyle \,\mathcal{K}(x'-w')\cdot\mathcal{K}(x-w)\,\mathcal{K}(w-z)\cdot\mathcal{K}(y-z)$
\end{itemize}
\needspace{3\baselineskip}\noindent\rule{\linewidth}{0.3pt}\par\noindent\textbf{Lines 295} \quad {\small (in \texttt{align}, lines 294--297)}\par
\begin{itemize}[leftmargin=2.2em,itemsep=1pt,topsep=2pt]
\item[\textbf{l.295}] \textsf{\small [phase convention]} $\displaystyle \,{\color{blue}\log{[\frac{k^{+}}{\Lambda}]}}\,\int_{x,y,z,x',w,w'}e^{-ik(w'-w)}\,\mathcal{K}(x'-w$\ $\longrightarrow$\ $\displaystyle \,{\color{blue}\log{[\frac{k^{+}}{\Lambda}]}}\,\int_{x,y,z,x',w,w'}e^{-ik(w-w')}\,\mathcal{K}(x'-w$
\item[\textbf{l.295}] \textsf{\small [term removed/replaced]} $\displaystyle \,\mathcal{K}(x'-w')\cdot\mathcal{K}(x-w)\,\mathcal{K}(x'-z)\cdot\mathcal{K}(y-z)$\ $\longrightarrow$\ $\displaystyle \,\mathcal{K}(x'-w')\cdot\mathcal{K}(x-w)\,\mathcal{K}(w-z)\cdot\mathcal{K}(y-z)$
\end{itemize}
\needspace{3\baselineskip}\noindent\rule{\linewidth}{0.3pt}\par\noindent\textbf{Lines 300} \quad {\small (in \texttt{align}, lines 299--301)}\par
\begin{itemize}[leftmargin=2.2em,itemsep=1pt,topsep=2pt]
\item[\textbf{l.300}] \textsf{\small [phase convention]} $\displaystyle \,{\color{blue}\log{[\frac{k^{+}}{\Lambda}]}}\,\int_{x,y,z,x',w,w'}e^{-ik(w'-w)}\,\mathcal{K}(x'-w$\ $\longrightarrow$\ $\displaystyle \,{\color{blue}\log{[\frac{k^{+}}{\Lambda}]}}\,\int_{x,y,z,x',w,w'}e^{-ik(w-w')}\,\mathcal{K}(x'-w$
\item[\textbf{l.300}] \textsf{\small [term removed/replaced]} $\displaystyle \,\mathcal{K}(x'-w')\cdot\mathcal{K}(x-w)\,\mathcal{K}(x'-z)\cdot\mathcal{K}(y-z)$\ $\longrightarrow$\ $\displaystyle \,\mathcal{K}(x'-w')\cdot\mathcal{K}(x-w)\,\mathcal{K}(w-z)\cdot\mathcal{K}(y-z)$
\end{itemize}
\needspace{3\baselineskip}\noindent\rule{\linewidth}{0.3pt}\par\noindent\textbf{Lines 305} \quad \textsf{Cancellation of all three $\rho$'s $\rangle$ $\beta'$ terms} \quad {\small (in \texttt{align}, lines 304--306)}\par
\begin{itemize}[leftmargin=2.2em,itemsep=1pt,topsep=2pt]
\item[\textbf{l.305}] \textsf{\small [phase convention]} $\displaystyle \,{\color{blue}\log{[\frac{k^{+}}{\Lambda}]}}\,\int_{x,y,z,x',w,w'}e^{-ik(w'-w)}\,\mathcal{K}(x'-w$\ $\longrightarrow$\ $\displaystyle \,{\color{blue}\log{[\frac{k^{+}}{\Lambda}]}}\,\int_{x,y,z,x',w,w'}e^{-ik(w-w')}\,\mathcal{K}(x'-w$
\item[\textbf{l.305}] \textsf{\small [term removed/replaced]} $\displaystyle \,\mathcal{K}(x'-w')\cdot\mathcal{K}(x-w)\,\mathcal{K}(x'-z)\cdot\mathcal{K}(y-z)$\ $\longrightarrow$\ $\displaystyle \,\mathcal{K}(x'-w')\cdot\mathcal{K}(x-w)\,\mathcal{K}(w'-z)\cdot\mathcal{K}(y-z)$
\end{itemize}
\needspace{3\baselineskip}\noindent\rule{\linewidth}{0.3pt}\par\noindent\textbf{Lines 309} \quad {\small (in \texttt{align}, lines 308--310)}\par
\begin{itemize}[leftmargin=2.2em,itemsep=1pt,topsep=2pt]
\item[\textbf{l.309}] \textsf{\small [phase convention]} $\displaystyle \,{\color{blue}\log{[\frac{k^{+}}{\Lambda}]}}\,\int_{x,y,z,x',w,w'}e^{-ik(w'-w)}\,\mathcal{K}(x'-w$\ $\longrightarrow$\ $\displaystyle \,{\color{blue}\log{[\frac{k^{+}}{\Lambda}]}}\,\int_{x,y,z,x',w,w'}e^{-ik(w-w')}\,\mathcal{K}(x'-w$
\item[\textbf{l.309}] \textsf{\small [term removed/replaced]} $\displaystyle \,\mathcal{K}(x'-w')\cdot\mathcal{K}(x-w)\,\mathcal{K}(x'-z)\cdot\mathcal{K}(y-z)$\ $\longrightarrow$\ $\displaystyle \,\mathcal{K}(x'-w')\cdot\mathcal{K}(x-w)\,\mathcal{K}(w'-z)\cdot\mathcal{K}(y-z)$
\end{itemize}
\needspace{3\baselineskip}\noindent\rule{\linewidth}{0.3pt}\par\noindent\textbf{Lines 313} \quad {\small (in \texttt{align}, lines 312--314)}\par
\begin{itemize}[leftmargin=2.2em,itemsep=1pt,topsep=2pt]
\item[\textbf{l.313}] \textsf{\small [phase convention]} $\displaystyle \,{\color{blue}\log{[\frac{k^{+}}{\Lambda}]}}\,\int_{x,y,z,x',w,w'}e^{-ik(w'-w)}\,\mathcal{K}(x'-w$\ $\longrightarrow$\ $\displaystyle \,{\color{blue}\log{[\frac{k^{+}}{\Lambda}]}}\,\int_{x,y,z,x',w,w'}e^{-ik(w-w')}\,\mathcal{K}(x'-w$
\item[\textbf{l.313}] \textsf{\small [term removed/replaced]} $\displaystyle \,\mathcal{K}(x'-w')\cdot\mathcal{K}(x-w)\,\mathcal{K}(x'-z)\cdot\mathcal{K}(y-z)$\ $\longrightarrow$\ $\displaystyle \,\mathcal{K}(x'-w')\cdot\mathcal{K}(x-w)\,\mathcal{K}(w'-z)\cdot\mathcal{K}(y-z)$
\end{itemize}
\needspace{15\baselineskip}\bigskip
\subsection{Evolution (JIMWLK) terms}\label{sec:file-ev_jimwlk}
\noindent\textbf{Original:} \texttt{/root/.claude/uploads/0ab97fcf-94a2-5aec-8bad-586400481685/4af0d15a-Evolution\_JIMWLK\_\_terms.tex}\\\textbf{Checked copy:} \texttt{/home/user/correlations/nlo\_review/evolution\_checked/Evolution\_JIMWLK\_terms\_checked.tex}\\
Lines: 764 (original), 764 (checked copy). Coloured groups of the check: 87 (87 top-level); author\textquotesingle s own coloured groups left unchanged: 0.\\
Records by type: phase convention 28; operator ordering 8; term removed/replaced 1; typo/LaTeX 1; note only 12.
\par\noindent Unmarked layout-only edits (line breaks \texttt{\textbackslash notag\textbackslash\textbackslash\ \&}, \texttt{\textbackslash par\textbackslash noindent} around inserted notes, spacing), checked-copy lines: 35, 36, 194, 228, 284, 313, 364, 707, 735, 762.
\subsubsection*{Changes, block by block}
\needspace{3\baselineskip}\noindent\rule{\linewidth}{0.3pt}\par\noindent\textbf{Lines 35--36}\par
\begin{notebox}\textsf{\small Note, l.35:} \textbf{Checked copy.} Marking: corrected or added material is in red; material removed from the author's text is kept in orange (cancelled); notes of the check are in red in square brackets. Phase: $e^{-ik\cdot(w-w')}$ everywhere (conjugated bracket, with the primed coordinates and standing left, at $w'$); this does not change the symmetrised two-$\rho$ parts, while the three-$\rho$ and one-$\rho$ parts change sign together. Result: the file is internally exact (the lists of Secs.~1 and 2 add up to the final formula of Sec.~3; the symmetrised three- and four-$\rho$ parts cancel identically), and with the operator ordering $A^{(1)}N_{2}$ of the three-$\rho$ part of $A^{(3)}$ the cross section reproduces it without any additional term; the final formula is the JIMWLK evolution of the target acting on the LO cross section (derived from first principles). Corrections: in the final list ten terms are written $\rho(x')\rho(x)$ instead of $\rho(x)\rho(x')$, and two misprints; two points are stated as open (Group I Row III of Sec.~2; Group II Row IV of Sec.~1).\end{notebox}
\begin{notebox}\textsf{\small Note, l.36:} [\textbf{Check of Secs.~1 and 2.} Notation: $P_{L}\equiv\frac{1}{(2\pi)^{3}}\frac{g^{4}}{16\pi^{5}}\frac{1}{k^{+}}\log\frac{k^{+}}{\Lambda}$, $K(a,b,z)\equiv\mathcal{K}(a-z)\cdot\mathcal{K}(b-z)$. (i) Every list was compared, pointwise in $x,x',z,w,w'$ and with $SU(2)$, $SU(3)$ matrices for the $\rho$'s, with the coefficient of $\log\frac{k^{+}}{\Lambda}$ of the checked Group files and of the four-$\rho$ file: the 8 genuine lists and the 16 reordering lists agree (Group I Rows I and II of Sec.~2: with the operator ordering $A^{(1)}N_{2}$, notes there); one point is left open at Group I Row III of Sec.~2. (ii) Phase: with the conjugated bracket at the primed points the phase is $e^{-ik\cdot(w-w')}$; it is changed in every header. The lists are not affected: each list is invariant under $w\leftrightarrow w'$ in the integrand (a Hermitian two-$\rho$ term with real coefficient is even in $k$); the symmetrised three-$\rho$ terms are odd. (iii) The UV brackets of the Group files are written here as $\big[\frac{2}{\epsilon}+2\gamma_{E}+\log\frac{r^{2}\mu^{2}_{\overline{MS}}}{4}\big]\to\frac{1}{\pi}\int_{z}\big[\frac{1}{(w-z)^{2}}-K(x,w,z)\big]$, $r=x-w$, which equals $\frac{1}{2\pi}\int_{z}\frac{r^{2}}{(x-z)^{2}(w-z)^{2}}$ up to scaleless integrals; pole and logarithm agree, the constant was not checked.]\end{notebox}
\needspace{3\baselineskip}\noindent\rule{\linewidth}{0.3pt}\par\noindent\textbf{Lines 39} \quad \textsf{Genuine contributions $\rangle$ Group I}\par
\begin{notebox}\textsf{\small Note, l.39:} [Checked, nothing else: the commutator term $-[\bar N_{2},\bar A^{(1)}]$ of Row I of the Group I file cancels against the difference between the correctly ordered three-$\rho$ term $-\bar A^{(1)}\bar N_{2}$ of $-\bar A^{(3)}$ (Row II) and $-\bar N_{2}\bar A^{(1)}$ (see the resolution of the operator ordering), so no three-$\rho$ or two-$\rho$ term comes from it.]\end{notebox}
\needspace{3\baselineskip}\noindent\rule{\linewidth}{0.3pt}\par\noindent\textbf{Lines 42} \quad {\small (in \texttt{align}, lines 41--59)}\par
\begin{itemize}[leftmargin=2.2em,itemsep=1pt,topsep=2pt]
\item[\textbf{l.42}] \textsf{\small [phase convention]} $\displaystyle \,\, \int_{x,x',z,w,w'}e^{-ik\cdot(w'-w)}$\ $\longrightarrow$\ $\displaystyle \,\, \int_{x,x',z,w,w'}e^{-ik\cdot(w-w')}$
\end{itemize}
\needspace{3\baselineskip}\noindent\rule{\linewidth}{0.3pt}\par\noindent\textbf{Lines 67} \quad {\small (in \texttt{align}, lines 66--92)}\par
\begin{itemize}[leftmargin=2.2em,itemsep=1pt,topsep=2pt]
\item[\textbf{l.67}] \textsf{\small [phase convention]} $\displaystyle \,\,\int_{x,x',z,w,w'}e^{-ik\cdot(w'-w)}$\ $\longrightarrow$\ $\displaystyle \,\,\int_{x,x',z,w,w'}e^{-ik\cdot(w-w')}$
\end{itemize}
\needspace{3\baselineskip}\noindent\rule{\linewidth}{0.3pt}\par\noindent\textbf{Lines 94} \quad {\small (in \texttt{align}, lines 93--110)}\par
\begin{itemize}[leftmargin=2.2em,itemsep=1pt,topsep=2pt]
\item[\textbf{l.94}] \textsf{\small [phase convention]} $\displaystyle \,\,\int_{x,x',z,w,w'}e^{-ik\cdot(w'-w)}$\ $\longrightarrow$\ $\displaystyle \,\,\int_{x,x',z,w,w'}e^{-ik\cdot(w-w')}$
\end{itemize}
\needspace{3\baselineskip}\noindent\rule{\linewidth}{0.3pt}\par\noindent\textbf{Lines 112} \quad {\small (in \texttt{align}, lines 111--136)}\par
\begin{itemize}[leftmargin=2.2em,itemsep=1pt,topsep=2pt]
\item[\textbf{l.112}] \textsf{\small [phase convention]} $\displaystyle \,\,\int_{x,x',z,w,w'}e^{-ik\cdot(w'-w)}$\ $\longrightarrow$\ $\displaystyle \,\,\int_{x,x',z,w,w'}e^{-ik\cdot(w-w')}$
\end{itemize}
\needspace{3\baselineskip}\noindent\rule{\linewidth}{0.3pt}\par\noindent\textbf{Lines 138} \quad {\small (in \texttt{align}, lines 137--154)}\par
\begin{itemize}[leftmargin=2.2em,itemsep=1pt,topsep=2pt]
\item[\textbf{l.138}] \textsf{\small [phase convention]} $\displaystyle \,\,\int_{x,x',z,w,w'}e^{-ik\cdot(w'-w)}$\ $\longrightarrow$\ $\displaystyle \,\,\int_{x,x',z,w,w'}e^{-ik\cdot(w-w')}$
\end{itemize}
\needspace{3\baselineskip}\noindent\rule{\linewidth}{0.3pt}\par\noindent\textbf{Lines 158} \quad \textsf{Genuine contributions $\rangle$ Group II} \quad {\small (in \texttt{align}, lines 157--175)}\par
\begin{itemize}[leftmargin=2.2em,itemsep=1pt,topsep=2pt]
\item[\textbf{l.158}] \textsf{\small [phase convention]} $\displaystyle \,\,\int_{x,x',z,w,w'}e^{-ik\cdot(w'-w)}$\ $\longrightarrow$\ $\displaystyle \,\,\int_{x,x',z,w,w'}e^{-ik\cdot(w-w')}$
\end{itemize}
\needspace{3\baselineskip}\noindent\rule{\linewidth}{0.3pt}\par\noindent\textbf{Lines 182--194} \quad {\small (in \texttt{align}, lines 181--194)}\par
\begin{itemize}[leftmargin=2.2em,itemsep=1pt,topsep=2pt]
\item[\textbf{l.182}] \textsf{\small [phase convention]} $\displaystyle \,\,\int_{x,x',z,w,w'}\,e^{-ik\cdot(w'-w)}\,$\ $\longrightarrow$\ $\displaystyle \,\,\int_{x,x',z,w,w'}\,e^{-ik\cdot(w-w')}\,$
\end{itemize}
\begin{notebox}\textsf{\small Note, l.194:} [This list is $P_{L}N_{c}\big[2K(w,w',z)-\frac{1}{(w-z)^{2}}-\frac{1}{(w'-z)^{2}}\big]\mathcal{K}(x-w)\cdot\mathcal{K}(x'-w')[U(w')-U(x')][U(w)-U(x)]\rho\rho$, i.e. the $\frac{2}{\epsilon}$ form of the Row IV file. It agrees with the $\int_{z}K(w,w',z)$ form of that file up to the two terms $\int_{z}\frac{1}{(w-z)^{2}}$, $\int_{z}\frac{1}{(w'-z)^{2}}$ with $z$-independent coefficient (scaleless); whether they are to be kept is left open.]\end{notebox}
\needspace{3\baselineskip}\noindent\rule{\linewidth}{0.3pt}\par\noindent\textbf{Lines 206} \quad \textsf{Genuine contributions $\rangle$ Group III (with c.c.)} \quad {\small (in \texttt{align}, lines 205--221)}\par
\begin{itemize}[leftmargin=2.2em,itemsep=1pt,topsep=2pt]
\item[\textbf{l.206}] \textsf{\small [phase convention]} $\displaystyle \,\,\int_{x,x',z,w,w'}e^{-ik\cdot(w'-w)}$\ $\longrightarrow$\ $\displaystyle \,\,\int_{x,x',z,w,w'}e^{-ik\cdot(w-w')}$
\end{itemize}
\needspace{3\baselineskip}\noindent\rule{\linewidth}{0.3pt}\par\noindent\textbf{Lines 228} \quad \textsf{Symmetrizing three and four $\rho$'s}\par
\begin{notebox}\textsf{\small Note, l.228:} [Convention found in the lists: every product of $\rho$'s is written as its totally symmetrised (Weyl) product plus remainders with fewer $\rho$'s, $\rho_{1}\rho_{2}\rho_{3}=S(\rho_{1}\rho_{2}\rho_{3})+\frac12\sum_{i<j}S([\rho_{i},\rho_{j}]\rho_{k})+(\text{one }\rho)$ and correspondingly for four $\rho$'s; the lists are the two-$\rho$ remainders, with the two $\rho$'s treated as commuting and the one-$\rho$ terms dropped; Groups I and III include the c.c. Result of an independent reordering: with the lists as they stand the symmetrised four-$\rho$ part vanishes identically, the symmetrised three-$\rho$ parts with $\log\frac{k^{+}}{\Lambda}$ (genuine rows plus the three-$\rho$ remainder of the four-$\rho$ terms) cancel pointwise, and the sum of all lists of Secs.~1 and 2 equals the final formula of Sec.~3 exactly. With the operator ordering $A^{(1)}N_{2}$ of the three-$\rho$ part of $A^{(3)}$ nothing has to be added to these lists (notes at Group I Rows I and II).]\end{notebox}
\needspace{3\baselineskip}\noindent\rule{\linewidth}{0.3pt}\par\noindent\textbf{Lines 231} \quad \textsf{Symmetrizing three and four $\rho$'s $\rangle$ From $4 \rho's$} \quad {\small (in \texttt{align}, lines 230--251)}\par
\begin{itemize}[leftmargin=2.2em,itemsep=1pt,topsep=2pt]
\item[\textbf{l.231}] \textsf{\small [phase convention]} $\displaystyle \,{\color{blue}\log{(\frac{k^{+}}{\Lambda})}}\,\int_{x,x',z,w,w'}e^{-ik(w'-w)}\,$\ $\longrightarrow$\ $\displaystyle \,{\color{blue}\log{(\frac{k^{+}}{\Lambda})}}\,\int_{x,x',z,w,w'}e^{-ik(w-w')}\,$
\end{itemize}
\needspace{3\baselineskip}\noindent\rule{\linewidth}{0.3pt}\par\noindent\textbf{Lines 253--284} \quad {\small (in \texttt{align}, lines 252--284)}\par
\begin{itemize}[leftmargin=2.2em,itemsep=1pt,topsep=2pt]
\item[\textbf{l.253}] \textsf{\small [phase convention]} $\displaystyle \,{\color{blue}\log{(\frac{k^{+}}{\Lambda})}}\,\int_{x,x',z,w,w'}e^{-ik(w'-w)}\,$\ $\longrightarrow$\ $\displaystyle \,{\color{blue}\log{(\frac{k^{+}}{\Lambda})}}\,\int_{x,x',z,w,w'}e^{-ik(w-w')}\,$
\end{itemize}
\begin{notebox}\textsf{\small Note, l.284:} [$N_{1}+N_{2}$ agrees with the two-$\rho$ remainder of the corrected Eq.~(4rho) of the four-$\rho$ file, obtained by an independent reordering.]\end{notebox}
\needspace{3\baselineskip}\noindent\rule{\linewidth}{0.3pt}\par\noindent\textbf{Lines 288} \quad \textsf{Symmetrizing three and four $\rho$'s $\rangle$ Group I (with c.c.)}\par
\begin{notebox}\textsf{\small Note, l.288:} [Checked, nothing else: the commutator term $-[\bar N_{2},\bar A^{(1)}]$ of Row I of the Group I file cancels against the difference between the correctly ordered three-$\rho$ term $-\bar A^{(1)}\bar N_{2}$ of $-\bar A^{(3)}$ (Row II) and $-\bar N_{2}\bar A^{(1)}$ (see the resolution of the operator ordering), so no three-$\rho$ or two-$\rho$ term comes from it.]\end{notebox}
\needspace{3\baselineskip}\noindent\rule{\linewidth}{0.3pt}\par\noindent\textbf{Lines 291--313} \quad {\small (in \texttt{align}, lines 290--313)}\par
\begin{itemize}[leftmargin=2.2em,itemsep=1pt,topsep=2pt]
\item[\textbf{l.291}] \textsf{\small [phase convention]} $\displaystyle \,{\color{blue}\log{(\frac{k^{+}}{\Lambda})}}\,\int_{x,x',z,w,w'}e^{-ik(w'-w)}\,$\ $\longrightarrow$\ $\displaystyle \,{\color{blue}\log{(\frac{k^{+}}{\Lambda})}}\,\int_{x,x',z,w,w'}e^{-ik(w-w')}\,$
\end{itemize}
\begin{notebox}\textsf{\small Note, l.313:} [Checked: this list is the two-$\rho$ remainder of $N_{2a}$ of the Group I file with the ordering $A^{(1)}N_{2}$ of the three-$\rho$ part of $A^{(3)}$, i.e. with the term $-2k^{i}\log\frac{\vee}{\Lambda}$ of the two-$\rho$ bracket of Row II kept only for the barred colour structure $f^{abc}U^{bd}(x)U^{ce}(y)\{\rho^{d}(x),\rho^{e}(y)\}$. The unbarred structure $f^{hbc}U^{ah}(w)\{\rho^{b}(x),\rho^{c}(y)\}$ therefore has no $K(x,y,z)$ term at $\log\frac{k^{+}}{\Lambda}$; the absence of that term here is correct.]\end{notebox}
\needspace{3\baselineskip}\noindent\rule{\linewidth}{0.3pt}\par\noindent\textbf{Lines 317} \quad {\small (in \texttt{align}, lines 316--338)}\par
\begin{itemize}[leftmargin=2.2em,itemsep=1pt,topsep=2pt]
\item[\textbf{l.317}] \textsf{\small [phase convention]} $\displaystyle \,{\color{blue}\log{(\frac{k^{+}}{\Lambda})}}\,\int_{x,x',z,w,w'}e^{-ik(w'-w)}\,$\ $\longrightarrow$\ $\displaystyle \,{\color{blue}\log{(\frac{k^{+}}{\Lambda})}}\,\int_{x,x',z,w,w'}e^{-ik(w-w')}\,$
\end{itemize}
\needspace{3\baselineskip}\noindent\rule{\linewidth}{0.3pt}\par\noindent\textbf{Lines 340--364} \quad {\small (in \texttt{align}, lines 339--364)}\par
\begin{itemize}[leftmargin=2.2em,itemsep=1pt,topsep=2pt]
\item[\textbf{l.340}] \textsf{\small [phase convention]} $\displaystyle \,{\color{blue}\log{(\frac{k^{+}}{\Lambda})}}\,\int_{x,x',z,w,w'}e^{-ik(w'-w)}\,$\ $\longrightarrow$\ $\displaystyle \,{\color{blue}\log{(\frac{k^{+}}{\Lambda})}}\,\int_{x,x',z,w,w'}e^{-ik(w-w')}\,$
\end{itemize}
\begin{notebox}\textsf{\small Note, l.364:} [$N_{1}$ agrees with the Group I file. $N_{2}$ is the remainder of the $\beta$ term in the form used in the three-$\rho$ file (with $U^{bc}(w)$). Open (not decided): at fixed coordinates the checked Group I row gives the kernel $-\mathcal{K}(x'-w')\cdot\mathcal{K}(y-z)/(z-w)^{2}$ with $U^{bc}(z)$, whereas the evolution files use $-\mathcal{K}(x'-w')\cdot\mathcal{K}(y-w)\,\mathcal{K}(x-z)\cdot\mathcal{K}(w-z)$ with $U^{bc}(w)$; the step between them ($z$ integration with $\Lambda$ kept; the region $z\to w$) was not verified.]\end{notebox}
\needspace{3\baselineskip}\noindent\rule{\linewidth}{0.3pt}\par\noindent\textbf{Lines 368} \quad {\small (in \texttt{align}, lines 367--402)}\par
\begin{itemize}[leftmargin=2.2em,itemsep=1pt,topsep=2pt]
\item[\textbf{l.368}] \textsf{\small [phase convention]} $\displaystyle \,{\color{blue}\log{(\frac{k^{+}}{\Lambda})}}\,\int_{x,x',z,w,w'}e^{-ik(w'-w)}\,$\ $\longrightarrow$\ $\displaystyle \,{\color{blue}\log{(\frac{k^{+}}{\Lambda})}}\,\int_{x,x',z,w,w'}e^{-ik(w-w')}\,$
\end{itemize}
\needspace{3\baselineskip}\noindent\rule{\linewidth}{0.3pt}\par\noindent\textbf{Lines 404} \quad {\small (in \texttt{align}, lines 403--416)}\par
\begin{itemize}[leftmargin=2.2em,itemsep=1pt,topsep=2pt]
\item[\textbf{l.404}] \textsf{\small [phase convention]} $\displaystyle \,{\color{blue}\log{(\frac{k^{+}}{\Lambda})}}\,\int_{x,x',z,w,w'}e^{-ik(w'-w)}\,$\ $\longrightarrow$\ $\displaystyle \,{\color{blue}\log{(\frac{k^{+}}{\Lambda})}}\,\int_{x,x',z,w,w'}e^{-ik(w-w')}\,$
\end{itemize}
\needspace{3\baselineskip}\noindent\rule{\linewidth}{0.3pt}\par\noindent\textbf{Lines 419} \quad {\small (in \texttt{align}, lines 418--435)}\par
\begin{itemize}[leftmargin=2.2em,itemsep=1pt,topsep=2pt]
\item[\textbf{l.419}] \textsf{\small [phase convention]} $\displaystyle \,{\color{blue}\log{(\frac{k^{+}}{\Lambda})}}\,\int_{x,x',z,w,w'}e^{-ik(w'-w)}\,$\ $\longrightarrow$\ $\displaystyle \,{\color{blue}\log{(\frac{k^{+}}{\Lambda})}}\,\int_{x,x',z,w,w'}e^{-ik(w-w')}\,$
\end{itemize}
\needspace{3\baselineskip}\noindent\rule{\linewidth}{0.3pt}\par\noindent\textbf{Lines 439} \quad \textsf{Symmetrizing three and four $\rho$'s $\rangle$ Group II} \quad {\small (in \texttt{align}, lines 438--470)}\par
\begin{itemize}[leftmargin=2.2em,itemsep=1pt,topsep=2pt]
\item[\textbf{l.439}] \textsf{\small [phase convention]} $\displaystyle \,{\color{blue}\log{(\frac{k^{+}}{\Lambda})}}\,\int_{x,x',z,w,w'}e^{-ik(w'-w)}\,$\ $\longrightarrow$\ $\displaystyle \,{\color{blue}\log{(\frac{k^{+}}{\Lambda})}}\,\int_{x,x',z,w,w'}e^{-ik(w-w')}\,$
\end{itemize}
\needspace{3\baselineskip}\noindent\rule{\linewidth}{0.3pt}\par\noindent\textbf{Lines 481} \quad \textsf{Symmetrizing three and four $\rho$'s $\rangle$ Group III (with c.c.)} \quad {\small (in \texttt{align}, lines 480--511)}\par
\begin{itemize}[leftmargin=2.2em,itemsep=1pt,topsep=2pt]
\item[\textbf{l.481}] \textsf{\small [phase convention]} $\displaystyle \,{\color{blue}\log{(\frac{k^{+}}{\Lambda})}}\,\int_{x,x',z,w,w'}e^{-ik(w'-w)}\,$\ $\longrightarrow$\ $\displaystyle \,{\color{blue}\log{(\frac{k^{+}}{\Lambda})}}\,\int_{x,x',z,w,w'}e^{-ik(w-w')}\,$
\end{itemize}
\needspace{3\baselineskip}\noindent\rule{\linewidth}{0.3pt}\par\noindent\textbf{Lines 513} \quad {\small (in \texttt{align}, lines 512--528)}\par
\begin{itemize}[leftmargin=2.2em,itemsep=1pt,topsep=2pt]
\item[\textbf{l.513}] \textsf{\small [phase convention]} $\displaystyle \,{\color{blue}\log{(\frac{k^{+}}{\Lambda})}}\,\int_{x,x',z,w,w'}e^{-ik(w'-w)}\,$\ $\longrightarrow$\ $\displaystyle \,{\color{blue}\log{(\frac{k^{+}}{\Lambda})}}\,\int_{x,x',z,w,w'}e^{-ik(w-w')}\,$
\end{itemize}
\needspace{3\baselineskip}\noindent\rule{\linewidth}{0.3pt}\par\noindent\textbf{Lines 532} \quad {\small (in \texttt{align}, lines 531--566)}\par
\begin{itemize}[leftmargin=2.2em,itemsep=1pt,topsep=2pt]
\item[\textbf{l.532}] \textsf{\small [phase convention]} $\displaystyle \,{\color{blue}\log{(\frac{k^{+}}{\Lambda})}}\,\int_{x,x',z,w,w'}e^{-ik(w'-w)}\,$\ $\longrightarrow$\ $\displaystyle \,{\color{blue}\log{(\frac{k^{+}}{\Lambda})}}\,\int_{x,x',z,w,w'}e^{-ik(w-w')}\,$
\end{itemize}
\needspace{3\baselineskip}\noindent\rule{\linewidth}{0.3pt}\par\noindent\textbf{Lines 568} \quad {\small (in \texttt{align}, lines 567--579)}\par
\begin{itemize}[leftmargin=2.2em,itemsep=1pt,topsep=2pt]
\item[\textbf{l.568}] \textsf{\small [phase convention]} $\displaystyle \,{\color{blue}\log{(\frac{k^{+}}{\Lambda})}}\,\int_{x,x',z,w,w'}e^{-ik(w'-w)}\,$\ $\longrightarrow$\ $\displaystyle \,{\color{blue}\log{(\frac{k^{+}}{\Lambda})}}\,\int_{x,x',z,w,w'}e^{-ik(w-w')}\,$
\end{itemize}
\needspace{3\baselineskip}\noindent\rule{\linewidth}{0.3pt}\par\noindent\textbf{Lines 582} \quad {\small (in \texttt{align}, lines 581--598)}\par
\begin{itemize}[leftmargin=2.2em,itemsep=1pt,topsep=2pt]
\item[\textbf{l.582}] \textsf{\small [phase convention]} $\displaystyle \,{\color{blue}\log{(\frac{k^{+}}{\Lambda})}}\,\int_{x,x',z,w,w'}e^{-ik(w'-w)}\,$\ $\longrightarrow$\ $\displaystyle \,{\color{blue}\log{(\frac{k^{+}}{\Lambda})}}\,\int_{x,x',z,w,w'}e^{-ik(w-w')}\,$
\end{itemize}
\needspace{3\baselineskip}\noindent\rule{\linewidth}{0.3pt}\par\noindent\textbf{Lines 603} \quad {\small (in \texttt{align}, lines 602--627)}\par
\begin{itemize}[leftmargin=2.2em,itemsep=1pt,topsep=2pt]
\item[\textbf{l.603}] \textsf{\small [phase convention]} $\displaystyle \,{\color{blue}\log{(\frac{k^{+}}{\Lambda})}}\,\int_{x,x',z,w,w'}e^{-ik(w'-w)}\,$\ $\longrightarrow$\ $\displaystyle \,{\color{blue}\log{(\frac{k^{+}}{\Lambda})}}\,\int_{x,x',z,w,w'}e^{-ik(w-w')}\,$
\end{itemize}
\needspace{3\baselineskip}\noindent\rule{\linewidth}{0.3pt}\par\noindent\textbf{Lines 630} \quad {\small (in \texttt{align}, lines 629--654)}\par
\begin{itemize}[leftmargin=2.2em,itemsep=1pt,topsep=2pt]
\item[\textbf{l.630}] \textsf{\small [phase convention]} $\displaystyle \,{\color{blue}\log{(\frac{k^{+}}{\Lambda})}}\,\int_{x,x',z,w,w'}e^{-ik(w'-w)}\,$\ $\longrightarrow$\ $\displaystyle \,{\color{blue}\log{(\frac{k^{+}}{\Lambda})}}\,\int_{x,x',z,w,w'}e^{-ik(w-w')}\,$
\end{itemize}
\needspace{3\baselineskip}\noindent\rule{\linewidth}{0.3pt}\par\noindent\textbf{Lines 657--677} \quad \textsf{JIMWLK contribution} \quad {\small (in \texttt{align}, lines 656--686)}\par
\begin{itemize}[leftmargin=2.2em,itemsep=1pt,topsep=2pt]
\item[\textbf{l.657}] \textsf{\small [phase convention]} $\displaystyle \,\int_{x,x',z,w,w'}\,e^{-ik(w'-w)}\,$\ $\longrightarrow$\ $\displaystyle \,\int_{x,x',z,w,w'}\,e^{-ik(w-w')}\,$
\item[\textbf{l.668}] \textsf{\small [operator ordering]} $\displaystyle (w)\,U^{ma}(z)\,U^{md}(x')\,\rho^{l}(x)\,\rho^{e}(x')$\ $\longrightarrow$\ $\displaystyle (w)\,U^{ma}(z)\,U^{md}(x')\,\rho^{e}(x')\,\rho^{l}(x)$
\item[\textbf{l.669}] \textsf{\small [operator ordering]} $\displaystyle (x)\,U^{ma}(z)\,U^{md}(x')\,\rho^{l}(x)\,\rho^{e}(x')$\ $\longrightarrow$\ $\displaystyle (x)\,U^{ma}(z)\,U^{md}(x')\,\rho^{e}(x')\,\rho^{l}(x)$
\item[\textbf{l.670}] \textsf{\small [operator ordering]} $\displaystyle (w)\,U^{ma}(z)\,U^{mb}(x')\,\rho^{l}(x)\,\rho^{e}(x')$\ $\longrightarrow$\ $\displaystyle (w)\,U^{ma}(z)\,U^{mb}(x')\,\rho^{e}(x')\,\rho^{l}(x)$
\item[\textbf{l.671}] \textsf{\small [operator ordering]} $\displaystyle (x)\,U^{ma}(z)\,U^{mb}(x')\,\rho^{l}(x)\,\rho^{e}(x')$\ $\longrightarrow$\ $\displaystyle (x)\,U^{ma}(z)\,U^{mb}(x')\,\rho^{e}(x')\,\rho^{l}(x)$
\item[\textbf{l.676}] \textsf{\small [operator ordering]} $\displaystyle (x')\,U^{hl}(x)\,U^{da}(z)\,\rho^{l}(x)\,\rho^{c}(x')$\ $\longrightarrow$\ $\displaystyle (x')\,U^{hl}(x)\,U^{da}(z)\,\rho^{c}(x')\,\rho^{l}(x)$
\item[\textbf{l.677}] \textsf{\small [operator ordering]} $\displaystyle (z)\,U^{cl}(x')\,U^{be}(x)\,\rho^{h}(x)\,\rho^{l}(x')$\ $\longrightarrow$\ $\displaystyle (z)\,U^{cl}(x')\,U^{be}(x)\,\rho^{l}(x')\,\rho^{h}(x)$
\end{itemize}
\needspace{3\baselineskip}\noindent\rule{\linewidth}{0.3pt}\par\noindent\textbf{Lines 688--707} \quad {\small (in \texttt{align}, lines 687--707)}\par
\begin{itemize}[leftmargin=2.2em,itemsep=1pt,topsep=2pt]
\item[\textbf{l.688}] \textsf{\small [phase convention]} $\displaystyle \,{\color{blue}\log{(\frac{k^{+}}{\Lambda})}}\,\int_{x,x',z,w,w'}e^{-ik(w'-w)}\,$\ $\longrightarrow$\ $\displaystyle \,{\color{blue}\log{(\frac{k^{+}}{\Lambda})}}\,\int_{x,x',z,w,w'}e^{-ik(w-w')}\,$
\item[\textbf{l.705}] \textsf{\small [operator ordering]} $\displaystyle f^{cde}\,U^{hd}(x')\,U^{ha}(w)\,\rho^{b}(x)\,\rho^{e}(x')$\ $\longrightarrow$\ $\displaystyle f^{cde}\,U^{hd}(x')\,U^{ha}(w)\,\rho^{e}(x')\,\rho^{b}(x)$
\item[\textbf{l.706}] \textsf{\small [operator ordering]} $\displaystyle (x')\,U^{bd}(z)\,U^{cl}(w)\,\rho^{l}(x)\,\rho^{h}(x')$\ $\longrightarrow$\ $\displaystyle (x')\,U^{bd}(z)\,U^{cl}(w)\,\rho^{h}(x')\,\rho^{l}(x)$
\end{itemize}
\begin{notebox}\textsf{\small Note, l.707:} [Phase as in the corrected cross section, $e^{-ik(w-w')}$ with the primed (conjugated) bracket at $w'$. In the terms marked the order of the two $\rho$'s is set to $\rho(x')\rho(x)$, so that $N_{1}+N_{2}$ equals the grouped form below also in the contact terms at $x=x'$]\end{notebox}
\needspace{3\baselineskip}\noindent\rule{\linewidth}{0.3pt}\par\noindent\textbf{Lines 709--735} \quad {\small (in \texttt{align}, lines 708--735)}\par
\begin{itemize}[leftmargin=2.2em,itemsep=1pt,topsep=2pt]
\item[\textbf{l.709}] \textsf{\small [phase convention]} $\displaystyle \,{\color{blue}\log{(\frac{k^{+}}{\Lambda})}}\,\int_{x,x',z,w,w'}e^{-ik(w'-w)}\,$\ $\longrightarrow$\ $\displaystyle \,{\color{blue}\log{(\frac{k^{+}}{\Lambda})}}\,\int_{x,x',z,w,w'}e^{-ik(w-w')}\,$
\end{itemize}
\begin{notebox}\textsf{\small Note, l.735:} [Checked: this equals $-H_{\rm JIMWLK}$, $H=\frac{\alpha_s}{2\pi^{2}}\int_{u,v,z}K(u,v,z)\,[J^{a}_{L}(u)J^{a}_{L}(v)+J^{a}_{R}(u)J^{a}_{R}(v)-2J^{a}_{L}(u)U^{ab}(z)J^{b}_{R}(v)]$, acting on $[U(x')-U(w')]^{ab'}[U(x)-U(w)]^{ab}$ of the LO cross section, with $\rho^{b'}(x')$ to the left of $\rho^{b}(x)$ and the phase $e^{-ik(w-w')}$, including the contact terms. The same expression follows from the emission of one gluon with $\Lambda<p^{+}<k^{+}$ by the valence charges and by the gluon $k$, which also fixes the normalisation. It vanishes at $U=1$ and is finite at $z\to x,x',w,w'$]\end{notebox}
\needspace{3\baselineskip}\noindent\rule{\linewidth}{0.3pt}\par\noindent\textbf{Lines 738--763} \quad {\small (in \texttt{align}, lines 737--762)}\par
\begin{itemize}[leftmargin=2.2em,itemsep=1pt,topsep=2pt]
\item[\textbf{l.738}] \textsf{\small [typo/LaTeX]} \emph{added:} $\displaystyle y$\quad{\color{gray}\small(in context: }$\displaystyle \,{\color{blue}\log{(\frac{k^{+}}{\Lambda})}}\,\int_{x,y,z,w,w'}e^{-ik(w-w')}\,\langle\rho^{f}${\color{gray}\small)}
\item[\textbf{l.738}] \textsf{\small [phase convention]} $\displaystyle \,{\color{blue}\log{(\frac{k^{+}}{\Lambda})}}\,\int_{x,,z,w,w'}e^{-ik(w'-w)}\,\langle\rho^{f}(x)$\ $\longrightarrow$\ $\displaystyle \,{\color{blue}\log{(\frac{k^{+}}{\Lambda})}}\,\int_{x,y,z,w,w'}e^{-ik(w-w')}\,\langle\rho^{f}(x)$
\item[\textbf{l.763}] \textsf{\small [term removed/replaced]} $\displaystyle \mathcal{K}$\ $\longrightarrow$\ $\displaystyle K$
\end{itemize}
\begin{notebox}\textsf{\small Note, l.762:} [Follows from the previous form with $\langle\rho^{a}(x)\rho^{b}(y)\rangle=\delta^{ab}\langle\rho^{f}(x)\rho^{f}(y)\rangle/(N_c^{2}-1)$ (symmetric under $x\leftrightarrow y$); $U$ and the traces are in the adjoint representation, $(T^{e})_{ab}=-if^{eab}$; independent of the sign of the phase]\end{notebox}
\needspace{15\baselineskip}\bigskip
\subsection{BFKL}\label{sec:file-bfkl}
\noindent\textbf{Original:} \texttt{/root/.claude/uploads/0ab97fcf-94a2-5aec-8bad-586400481685/9da5dc76-BFKL.tex}\\\textbf{Checked copy:} \texttt{/home/user/correlations/nlo\_review/evolution\_checked/BFKL\_checked.tex}\\
Lines: 646 (original), 646 (checked copy). Coloured groups of the check: 163 (163 top-level); author\textquotesingle s own coloured groups left unchanged: 8.\\
Records by type: phase convention 59; operator ordering 3; typo/LaTeX 6; note only 33.
\par\noindent Unmarked layout-only edits (line breaks \texttt{\textbackslash notag\textbackslash\textbackslash\ \&}, \texttt{\textbackslash par\textbackslash noindent} around inserted notes, spacing), checked-copy lines: 31, 39, 45, 68, 101, 111, 122, 143, 185, 193, 237, 257, 292, 315, 337, 346, 353, 361, 403, 410, 456, 464, 484, 493, 516, 525, 544, 551, 568, 605, 633, 644.
\subsubsection*{Changes, block by block}
\needspace{3\baselineskip}\noindent\rule{\linewidth}{0.3pt}\par\noindent\textbf{Lines 31}\par
\begin{notebox}\textsf{\small Note, l.31:} \textbf{Checked copy.} Marking: corrected or added material is in red; material removed from the author's text is kept in orange (cancelled); notes of the check are in red in square brackets. The red labels $\alpha$, $\alpha'$, $\beta$, $\beta'$ of the first two equations are the author's own. Phase: $e^{-ik\cdot(w-w')}$ everywhere (conjugated bracket, with the primed coordinates and standing left, at $w'$); this does not change the symmetrised two-$\rho$ parts, while the three-$\rho$ and one-$\rho$ parts change sign together. Result: the file is internally exact (the lists add up to the final formula; the symmetrised three- and four-$\rho$ parts cancel identically), and with the operator ordering $A^{(1)}N_{2}$ of the three-$\rho$ part of $A^{(3)}$ the cross section reproduces it without any additional term; the final formula is the BFKL evolution of the projectile acting on the LO cross section (derived from first principles). Corrections: four missing brackets, and the order of the two $\rho$'s in three terms of the final list; one remark on the grouped form (note after it).\end{notebox}
\needspace{3\baselineskip}\noindent\rule{\linewidth}{0.3pt}\par\noindent\textbf{Lines 35--39} \quad \textsf{From $4$ $\rho$'s $\rangle$ From four $\rho$'s} \quad {\small (in \texttt{align}, lines 34--39)}\par
\begin{itemize}[leftmargin=2.2em,itemsep=1pt,topsep=2pt]
\item[\textbf{l.35}] \textsf{\small [phase convention]} $\displaystyle \,{\color{blue}\log{[\frac{\vee}{k^{+}}]}}\,\int_{x,y,z,x^{\prime},w,w^{\prime}}e^{-ik(w^{\prime}-w)}\,$\ $\longrightarrow$\ $\displaystyle \,{\color{blue}\log{[\frac{\vee}{k^{+}}]}}\,\int_{x,y,z,x^{\prime},w,w^{\prime}}e^{-ik(w-w^{\prime})}\,$
\end{itemize}
\begin{notebox}\textsf{\small Note, l.39:} [Phase: with the conjugated bracket at $w'$ and the other bracket at $w$ the definition gives $e^{-ik\cdot(w-w')}$; changed in every equation of this file (in Group III, Rows III, V, VI: $e^{-ik\cdot(w-y')}$). Checked with non-commuting $\rho$ (explicit SU(2), SU(3) matrices and the algebra): for every row the symmetrised two-$\rho$ part is the same for both phases, the symmetrised three-$\rho$ and one-$\rho$ parts change sign.]\end{notebox}
\needspace{3\baselineskip}\noindent\rule{\linewidth}{0.3pt}\par\noindent\textbf{Lines 41--45} \quad {\small (in \texttt{align}, lines 40--45)}\par
\begin{itemize}[leftmargin=2.2em,itemsep=1pt,topsep=2pt]
\item[\textbf{l.41}] \textsf{\small [phase convention]} $\displaystyle \,{\color{blue}\log{[\frac{\vee}{k^{+}}]}}\,\int_{x,y,z,x^{\prime},w,w^{\prime}}e^{-ik(w^{\prime}-w)}\,$\ $\longrightarrow$\ $\displaystyle \,{\color{blue}\log{[\frac{\vee}{k^{+}}]}}\,\int_{x,y,z,x^{\prime},w,w^{\prime}}e^{-ik(w-w^{\prime})}\,$
\item[\textbf{l.41}] \textsf{\small [typo/LaTeX]} \emph{added:} $\displaystyle )$\quad{\color{gray}\small(in context: }$\displaystyle w') U^{ce}(x${\color{gray}\small)}
\item[\textbf{l.41}] \textsf{\small [typo/LaTeX]} \emph{added:} $\displaystyle )$\quad{\color{gray}\small(in context: }$\displaystyle x) U^{a h}(y)\,{\color{red}\alpha}${\color{gray}\small)}
\item[\textbf{l.42}] \textsf{\small [typo/LaTeX]} \emph{added:} $\displaystyle )$\quad{\color{gray}\small(in context: }$\displaystyle x') U^{be}(w${\color{gray}\small)}
\item[\textbf{l.42}] \textsf{\small [typo/LaTeX]} \emph{added:} $\displaystyle )$\quad{\color{gray}\small(in context: }$\displaystyle w) U^{a h}(y)\, {\color{orange} \alpha'}${\color{gray}\small)}
\item[\textbf{l.43}] \textsf{\small [typo/LaTeX]} \emph{added:} $\displaystyle )$\quad{\color{gray}\small(in context: }$\displaystyle x') U^{b e}(x) U^{a h}(y)\, {\color{orange} \alpha'}${\color{gray}\small)}
\item[\textbf{l.44}] \textsf{\small [typo/LaTeX]} \emph{added:} $\displaystyle )$\quad{\color{gray}\small(in context: }$\displaystyle x) U^{b e}(x') U^{a h}(y)\, {\color{red}\alpha}${\color{gray}\small)}
\end{itemize}
\begin{notebox}\textsf{\small Note, l.45:} [Checked: $N_{a}+N_{b}$ is the three-$\rho$ remainder of the corrected four-$\rho$ equation of the four-$\rho$ file; with the corrected phase the same expression holds without an extra sign. Four missing brackets restored.]\end{notebox}
\needspace{3\baselineskip}\noindent\rule{\linewidth}{0.3pt}\par\noindent\textbf{Lines 48--68} \quad \textsf{From $4$ $\rho$'s $\rangle$ From $4 \rho's$} \quad {\small (in \texttt{align}, lines 47--68)}\par
\begin{itemize}[leftmargin=2.2em,itemsep=1pt,topsep=2pt]
\item[\textbf{l.48}] \textsf{\small [phase convention]} $\displaystyle \,{\color{blue}\log{(\frac{\vee}{k^{+}})}}\,\int_{x,x',z,w,w'}e^{-ik(w'-w)}\,$\ $\longrightarrow$\ $\displaystyle \,{\color{blue}\log{(\frac{\vee}{k^{+}})}}\,\int_{x,x',z,w,w'}e^{-ik(w-w')}\,$
\end{itemize}
\begin{notebox}\textsf{\small Note, l.68:} [Checked: $N_{1}$ is the symmetrised two-$\rho$ part of $N_{a}+N_{b}$ above.]\end{notebox}
\needspace{3\baselineskip}\noindent\rule{\linewidth}{0.3pt}\par\noindent\textbf{Lines 70--101} \quad {\small (in \texttt{align}, lines 69--101)}\par
\begin{itemize}[leftmargin=2.2em,itemsep=1pt,topsep=2pt]
\item[\textbf{l.70}] \textsf{\small [phase convention]} $\displaystyle \,{\color{blue}\log{(\frac{\vee}{k^{+}})}}\,\int_{x,x',z,w,w'}e^{-ik(w'-w)}\,$\ $\longrightarrow$\ $\displaystyle \,{\color{blue}\log{(\frac{\vee}{k^{+}})}}\,\int_{x,x',z,w,w'}e^{-ik(w-w')}\,$
\end{itemize}
\begin{notebox}\textsf{\small Note, l.101:} [Checked by an independent reordering of the corrected four-$\rho$ equation (the two-$\rho$ section of the four-$\rho$ file was not used): its symmetrised four-$\rho$ part vanishes identically and $N_{1}+N_{2}$ is its complete symmetrised two-$\rho$ part, for both phases (only the sum is independent of the phase). Here and in all lists below $\rho\rho$ has to be read as the symmetrised product: as ordered products the lists differ from the operator identity by one-$\rho$ terms, which are not tracked.]\end{notebox}
\needspace{3\baselineskip}\noindent\rule{\linewidth}{0.3pt}\par\noindent\textbf{Lines 104} \quad \textsf{Group I (with c.c.) $\rangle$ Row I}\par
\begin{notebox}\textsf{\small Note, l.104:} [Checked: remains NA; the commutator term $-[\bar N_{2},\bar A^{(1)}]$ of Row I of the Group I file cancels against the difference between the correctly ordered three-$\rho$ term $-\bar A^{(1)}\bar N_{2}$ of $-\bar A^{(3)}$ (Row II) and $-\bar N_{2}\bar A^{(1)}$ (see the resolution of the operator ordering), so no three-$\rho$ or two-$\rho$ term comes from it.]\end{notebox}
\needspace{3\baselineskip}\noindent\rule{\linewidth}{0.3pt}\par\noindent\textbf{Lines 108--111} \quad \textsf{Group I (with c.c.) $\rangle$ Row II} \quad {\small (in \texttt{align}, lines 106--111)}\par
\begin{itemize}[leftmargin=2.2em,itemsep=1pt,topsep=2pt]
\item[\textbf{l.108}] \textsf{\small [phase convention]} $\displaystyle \int_{w,w'} e^{-ik\cdot(w'-w)}\int_{x,y,x',z}\frac{(x'-w')^{i}}{(x'-w')^{2}}$\ $\longrightarrow$\ $\displaystyle \int_{w,w'} e^{-ik\cdot(w-w')}\int_{x,y,x',z}\frac{(x'-w')^{i}}{(x'-w')^{2}}$
\item[\textbf{l.110}] \textsf{\small [phase convention]} $\displaystyle \int_{w,w'} e^{-ik\cdot(w'-w)}\int_{x,y,x'}\frac{(x'-w')^{i}}{(x'-w')^{2}}$\ $\longrightarrow$\ $\displaystyle \int_{w,w'} e^{-ik\cdot(w-w')}\int_{x,y,x'}\frac{(x'-w')^{i}}{(x'-w')^{2}}$
\end{itemize}
\begin{notebox}\textsf{\small Note, l.111:} [Kept; it is correct. Origin: the three-$\rho$ part of $A^{(3)}$ has the ordering $A^{(1)}N_{2}$; with the four-$\rho$ equation as it stands this means that the term $-2k^{i}\log\frac{\vee}{\Lambda}$ of the two-$\rho$ bracket of Row II of the Group I file is present for the barred colour structure only and absent for the unbarred one, $f^{hbc}U^{ah}(w)\{\rho^{b}(x),\rho^{c}(y)\}$. The unbarred structure thereby acquires this $\log\frac{\vee}{k^{+}}$ term and loses the corresponding $\log\frac{k^{+}}{\Lambda}$ term (JIMWLK and three-$\rho$ files).]\end{notebox}
\needspace{3\baselineskip}\noindent\rule{\linewidth}{0.3pt}\par\noindent\textbf{Lines 114--122} \quad {\small (in \texttt{align}, lines 113--122)}\par
\begin{itemize}[leftmargin=2.2em,itemsep=1pt,topsep=2pt]
\item[\textbf{l.114}] \textsf{\small [phase convention]} $\displaystyle \log{(\frac{\vee}{k^{+}})}\, \int_{x,x',z,w,w'}e^{-ik\cdot(w'-w)}$\ $\longrightarrow$\ $\displaystyle \log{(\frac{\vee}{k^{+}})}\, \int_{x,x',z,w,w'}e^{-ik\cdot(w-w')}$
\end{itemize}
\begin{notebox}\textsf{\small Note, l.122:} [Checked: this list is the symmetrised two-$\rho$ part of the expression above plus its c.c.]\end{notebox}
\needspace{3\baselineskip}\noindent\rule{\linewidth}{0.3pt}\par\noindent\textbf{Lines 126} \quad \textsf{Group I (with c.c.) $\rangle$ Row III} \quad {\small (in \texttt{align}, lines 125--131)}\par
\begin{itemize}[leftmargin=2.2em,itemsep=1pt,topsep=2pt]
\item[\textbf{l.126}] \textsf{\small [phase convention]} $\displaystyle \int_{w, w^\prime} e^{-ik(w^\prime-w)}\,$\ $\longrightarrow$\ $\displaystyle \int_{w, w^\prime} e^{-ik(w-w^\prime)}\,$
\end{itemize}
\needspace{3\baselineskip}\noindent\rule{\linewidth}{0.3pt}\par\noindent\textbf{Lines 133--143} \quad {\small (in \texttt{align}, lines 132--143)}\par
\begin{itemize}[leftmargin=2.2em,itemsep=1pt,topsep=2pt]
\item[\textbf{l.133}] \textsf{\small [phase convention]} $\displaystyle \int_{w, w^\prime} e^{-ik(w^\prime-w)}\,$\ $\longrightarrow$\ $\displaystyle \int_{w, w^\prime} e^{-ik(w-w^\prime)}\,$
\end{itemize}
\begin{notebox}\textsf{\small Note, l.143:} [Checked: the expressions above are the $\log\frac{\vee}{k^{+}}$ terms of the checked Group I file for this row (up to the phase); the terms added there in red cancel among themselves.]\end{notebox}
\needspace{3\baselineskip}\noindent\rule{\linewidth}{0.3pt}\par\noindent\textbf{Lines 146} \quad {\small (in \texttt{align}, lines 145--175)}\par
\begin{itemize}[leftmargin=2.2em,itemsep=1pt,topsep=2pt]
\item[\textbf{l.146}] \textsf{\small [phase convention]} $\displaystyle \log{(\frac{\vee}{k^{+}})}\, \int_{x,x',z,w,w'}e^{-ik\cdot(w'-w)}$\ $\longrightarrow$\ $\displaystyle \log{(\frac{\vee}{k^{+}})}\, \int_{x,x',z,w,w'}e^{-ik\cdot(w-w')}$
\end{itemize}
\needspace{3\baselineskip}\noindent\rule{\linewidth}{0.3pt}\par\noindent\textbf{Lines 177--185} \quad {\small (in \texttt{align}, lines 176--185)}\par
\begin{itemize}[leftmargin=2.2em,itemsep=1pt,topsep=2pt]
\item[\textbf{l.177}] \textsf{\small [phase convention]} $\displaystyle \log{(\frac{\vee}{k^{+}})}\, \int_{x,x',z,w,w'}e^{-ik\cdot(w'-w)}$\ $\longrightarrow$\ $\displaystyle \log{(\frac{\vee}{k^{+}})}\, \int_{x,x',z,w,w'}e^{-ik\cdot(w-w')}$
\end{itemize}
\begin{notebox}\textsf{\small Note, l.185:} [Checked: $N_{1}+N_{2}$ is the symmetrised two-$\rho$ part of the expressions above plus their c.c.]\end{notebox}
\needspace{3\baselineskip}\noindent\rule{\linewidth}{0.3pt}\par\noindent\textbf{Lines 189--193} \quad \textsf{Group I (with c.c.) $\rangle$ Row IV} \quad {\small (in \texttt{align}, lines 187--193)}\par
\begin{itemize}[leftmargin=2.2em,itemsep=1pt,topsep=2pt]
\item[\textbf{l.189}] \textsf{\small [phase convention]} $\displaystyle \int_{x,y,z,w,x',w'}\,e^{-ik\cdot (w^\prime-w)}\,$\ $\longrightarrow$\ $\displaystyle \int_{x,y,z,w,x',w'}\,e^{-ik\cdot (w-w^\prime)}\,$
\end{itemize}
\begin{notebox}\textsf{\small Note, l.193:} [Checked: the expressions above are the $\log\frac{\vee}{k^{+}}$ terms of the checked Group I file for this row (up to the phase).]\end{notebox}
\needspace{3\baselineskip}\noindent\rule{\linewidth}{0.3pt}\par\noindent\textbf{Lines 196} \quad {\small (in \texttt{align}, lines 195--228)}\par
\begin{itemize}[leftmargin=2.2em,itemsep=1pt,topsep=2pt]
\item[\textbf{l.196}] \textsf{\small [phase convention]} $\displaystyle \log{(\frac{\vee}{k^{+}})}\, \int_{x,x',z,w,w'}e^{-ik\cdot(w'-w)}$\ $\longrightarrow$\ $\displaystyle \log{(\frac{\vee}{k^{+}})}\, \int_{x,x',z,w,w'}e^{-ik\cdot(w-w')}$
\end{itemize}
\needspace{3\baselineskip}\noindent\rule{\linewidth}{0.3pt}\par\noindent\textbf{Lines 230--237} \quad {\small (in \texttt{align}, lines 229--237)}\par
\begin{itemize}[leftmargin=2.2em,itemsep=1pt,topsep=2pt]
\item[\textbf{l.230}] \textsf{\small [phase convention]} $\displaystyle \log{(\frac{\vee}{k^{+}})}\, \int_{x,x',z,w,w'}e^{-ik\cdot(w'-w)}$\ $\longrightarrow$\ $\displaystyle \log{(\frac{\vee}{k^{+}})}\, \int_{x,x',z,w,w'}e^{-ik\cdot(w-w')}$
\end{itemize}
\begin{notebox}\textsf{\small Note, l.237:} [Checked: $N_{1}+N_{2}$ is the symmetrised two-$\rho$ part of the expression above plus its c.c.]\end{notebox}
\needspace{3\baselineskip}\noindent\rule{\linewidth}{0.3pt}\par\noindent\textbf{Lines 247} \quad \textsf{Group II $\rangle$ Row I} \quad {\small (in \texttt{align}, lines 246--251)}\par
\begin{itemize}[leftmargin=2.2em,itemsep=1pt,topsep=2pt]
\item[\textbf{l.247}] \textsf{\small [phase convention]} $\displaystyle \,\int_{x,y,z,w,x',w^{\prime}} e^{-ik(w^{\prime}-w)}$\ $\longrightarrow$\ $\displaystyle \,\int_{x,y,z,w,x',w^{\prime}} e^{-ik(w-w^{\prime})}$
\end{itemize}
\needspace{3\baselineskip}\noindent\rule{\linewidth}{0.3pt}\par\noindent\textbf{Lines 253--257} \quad {\small (in \texttt{align}, lines 252--257)}\par
\begin{itemize}[leftmargin=2.2em,itemsep=1pt,topsep=2pt]
\item[\textbf{l.253}] \textsf{\small [phase convention]} $\displaystyle \,\int_{x,y,z,w,x',w^{\prime}} e^{-ik(w^{\prime}-w)}$\ $\longrightarrow$\ $\displaystyle \,\int_{x,y,z,w,x',w^{\prime}} e^{-ik(w-w^{\prime})}$
\end{itemize}
\begin{notebox}\textsf{\small Note, l.257:} [Checked: the two expressions above together are the $\log\frac{\vee}{k^{+}}$ terms of the checked Group II file, Row I, first part (up to the phase).]\end{notebox}
\needspace{3\baselineskip}\noindent\rule{\linewidth}{0.3pt}\par\noindent\textbf{Lines 260--292} \quad {\small (in \texttt{align}, lines 259--292)}\par
\begin{itemize}[leftmargin=2.2em,itemsep=1pt,topsep=2pt]
\item[\textbf{l.260}] \textsf{\small [phase convention]} $\displaystyle \log{(\frac{\vee}{k^{+}})}\, \int_{x,x',z,w,w'}e^{-ik\cdot(w'-w)}$\ $\longrightarrow$\ $\displaystyle \log{(\frac{\vee}{k^{+}})}\, \int_{x,x',z,w,w'}e^{-ik\cdot(w-w')}$
\end{itemize}
\begin{notebox}\textsf{\small Note, l.292:} [Checked: symmetrised two-$\rho$ part of the two expressions above (no c.c. in Group II).]\end{notebox}
\needspace{3\baselineskip}\noindent\rule{\linewidth}{0.3pt}\par\noindent\textbf{Lines 295--297} \quad {\small (in \texttt{align}, lines 294--298)}\par
\begin{itemize}[leftmargin=2.2em,itemsep=1pt,topsep=2pt]
\item[\textbf{l.295}] \textsf{\small [phase convention]} $\displaystyle \,\int_{z,w,w^{\prime},x,x'}\, e^{-ik(w^{\prime}-w)}\, \frac{1}{[(x'-w')^{2}\,(x-z)^{2}\,(x-w)^{2} - (x-w)^{4}\, (x'-z)^{2}]}$\ $\longrightarrow$\ $\displaystyle \,\int_{z,w,w^{\prime},x,x'}\, e^{-ik(w-w^{\prime})}\, \frac{1}{[(x'-w')^{2}\,(x-z)^{2}\,(x-w)^{2} - (x-w)^{4}\, (x'-z)^{2}]}$
\item[\textbf{l.296}] \textsf{\small [phase convention]} $\displaystyle \,\int_{z,w,w^{\prime},x,x'}\, e^{-ik(w^{\prime}-w)}\, \frac{(x-z)^{2}}{[(x'-w')^{2}\,(x-z)^{2}\,(x-w)^{2} - (x-w)^{4}\, (x'-z)^{2}]}\,$\ $\longrightarrow$\ $\displaystyle \,\int_{z,w,w^{\prime},x,x'}\, e^{-ik(w-w^{\prime})}\, \frac{(x-z)^{2}}{[(x'-w')^{2}\,(x-z)^{2}\,(x-w)^{2} - (x-w)^{4}\, (x'-z)^{2}]}\,$
\item[\textbf{l.297}] \textsf{\small [phase convention]} $\displaystyle \,\int_{z,w,w^{\prime},x,x'}\, e^{-ik(w^{\prime}-w)}\, \frac{(x'-z)^{2}}{[(x'-w')^{4}\,(x-z)^{2}\, - (x-w)^{2}\, (x'-z)^{2}\,(x'-w')^{2}]}\,$\ $\longrightarrow$\ $\displaystyle \,\int_{z,w,w^{\prime},x,x'}\, e^{-ik(w-w^{\prime})}\, \frac{(x'-z)^{2}}{[(x'-w')^{4}\,(x-z)^{2}\, - (x-w)^{2}\, (x'-z)^{2}\,(x'-w')^{2}]}\,$
\end{itemize}
\needspace{3\baselineskip}\noindent\rule{\linewidth}{0.3pt}\par\noindent\textbf{Lines 300--301} \quad {\small (in \texttt{align}, lines 299--302)}\par
\begin{itemize}[leftmargin=2.2em,itemsep=1pt,topsep=2pt]
\item[\textbf{l.300}] \textsf{\small [phase convention]} $\displaystyle (w^{\prime}-w)$\ $\longrightarrow$\ $\displaystyle (w-w^{\prime})$
\item[\textbf{l.301}] \textsf{\small [phase convention]} $\displaystyle (w^{\prime}-w)$\ $\longrightarrow$\ $\displaystyle (w-w^{\prime})$
\end{itemize}
\needspace{3\baselineskip}\noindent\rule{\linewidth}{0.3pt}\par\noindent\textbf{Lines 304--306} \quad {\small (in \texttt{align}, lines 303--307)}\par
\begin{itemize}[leftmargin=2.2em,itemsep=1pt,topsep=2pt]
\item[\textbf{l.304}] \textsf{\small [phase convention]} $\displaystyle (w^{\prime}-w)$\ $\longrightarrow$\ $\displaystyle (w-w^{\prime})$
\item[\textbf{l.305}] \textsf{\small [phase convention]} $\displaystyle (w^{\prime}-w)$\ $\longrightarrow$\ $\displaystyle (w-w^{\prime})$
\item[\textbf{l.306}] \textsf{\small [phase convention]} $\displaystyle (w^{\prime}-w)$\ $\longrightarrow$\ $\displaystyle (w-w^{\prime})$
\end{itemize}
\needspace{3\baselineskip}\noindent\rule{\linewidth}{0.3pt}\par\noindent\textbf{Lines 309--310} \quad {\small (in \texttt{align}, lines 308--311)}\par
\begin{itemize}[leftmargin=2.2em,itemsep=1pt,topsep=2pt]
\item[\textbf{l.309}] \textsf{\small [phase convention]} $\displaystyle (w^{\prime}-w)$\ $\longrightarrow$\ $\displaystyle (w-w^{\prime})$
\item[\textbf{l.309}] \textsf{\small [phase convention]} $\displaystyle (w^{\prime}-w)$\ $\longrightarrow$\ $\displaystyle (w-w^{\prime})$
\item[\textbf{l.310}] \textsf{\small [phase convention]} $\displaystyle (w^{\prime}-w)$\ $\longrightarrow$\ $\displaystyle (w-w^{\prime})$
\end{itemize}
\needspace{3\baselineskip}\noindent\rule{\linewidth}{0.3pt}\par\noindent\textbf{Lines 313--315} \quad {\small (in \texttt{align}, lines 312--315)}\par
\begin{itemize}[leftmargin=2.2em,itemsep=1pt,topsep=2pt]
\item[\textbf{l.313}] \textsf{\small [phase convention]} $\displaystyle \,\int_{z,w,w^{\prime},x,x'} e^{-ik(w^{\prime}-w)}\,\frac{[(x'-w')^{2}\, (x-w)^{2} - (x-z)^{2}\, (x'-z)^{2}]}{(z-w)^{2}\, (z-w')^{2}\,[(x-z)^{2} - (x-w)^{2}]\, [(x'-z)^{2} - (x'-w')^{2}]}$\ $\longrightarrow$\ $\displaystyle \,\int_{z,w,w^{\prime},x,x'} e^{-ik(w-w^{\prime})}\,\frac{[(x'-w')^{2}\, (x-w)^{2} - (x-z)^{2}\, (x'-z)^{2}]}{(z-w)^{2}\, (z-w')^{2}\,[(x-z)^{2} - (x-w)^{2}]\, [(x'-z)^{2} - (x'-w')^{2}]}$
\item[\textbf{l.314}] \textsf{\small [phase convention]} $\displaystyle \log{(\frac{\vee}{k^{+}})}\,\int_{z,w,w^{\prime},x,x'} e^{-ik(w^{\prime}-w)}\,\frac{(x-z)^{2}\, (x-w)^{2}}{(z-w)^{2}\, (z-w')^{2}\,[(x-z)^{2} - (x-w)^{2}]}\,$\ $\longrightarrow$\ $\displaystyle \log{(\frac{\vee}{k^{+}})}\,\int_{z,w,w^{\prime},x,x'} e^{-ik(w-w^{\prime})}\,\frac{(x-z)^{2}\, (x-w)^{2}}{(z-w)^{2}\, (z-w')^{2}\,[(x-z)^{2} - (x-w)^{2}]}\,$
\item[\textbf{l.314}] \textsf{\small [phase convention]} $\displaystyle \log{(\frac{\vee}{k^{+}})}\,\int_{z,w,w^{\prime},x,x'} e^{-ik(w^{\prime}-w)}\,\frac{(x'-w')^{2}\, (x'-z)^{2}}{(z-w)^{2}\, (z-w')^{2}}$\ $\longrightarrow$\ $\displaystyle \log{(\frac{\vee}{k^{+}})}\,\int_{z,w,w^{\prime},x,x'} e^{-ik(w-w^{\prime})}\,\frac{(x'-w')^{2}\, (x'-z)^{2}}{(z-w)^{2}\, (z-w')^{2}}$
\end{itemize}
\begin{notebox}\textsf{\small Note, l.315:} [Checked: the sum of $N_{2b},\dots,N_{2f}$ agrees with the $\log\frac{\vee}{k^{+}}$ terms of the checked Group II file, Row I, second part (up to the phase), where the whole logarithm sits in $N_{2c}$.]\end{notebox}
\needspace{3\baselineskip}\noindent\rule{\linewidth}{0.3pt}\par\noindent\textbf{Lines 318--337} \quad {\small (in \texttt{align}, lines 317--337)}\par
\begin{itemize}[leftmargin=2.2em,itemsep=1pt,topsep=2pt]
\item[\textbf{l.318}] \textsf{\small [phase convention]} $\displaystyle \log{(\frac{\vee}{k^{+}})}\, \int_{x,x',z,w,w'}e^{-ik\cdot(w'-w)}$\ $\longrightarrow$\ $\displaystyle \log{(\frac{\vee}{k^{+}})}\, \int_{x,x',z,w,w'}e^{-ik\cdot(w-w')}$
\end{itemize}
\begin{notebox}\textsf{\small Note, l.337:} [Checked: equal to the sum of $N_{2b},\dots,N_{2f}$ above at all points where these are regular (they have removable $0/0$ points at $x=x'$, $w=w'$).]\end{notebox}
\needspace{3\baselineskip}\noindent\rule{\linewidth}{0.3pt}\par\noindent\textbf{Lines 345--346} \quad \textsf{Group II $\rangle$ Row IV} \quad {\small (in \texttt{align}, lines 344--346)}\par
\begin{itemize}[leftmargin=2.2em,itemsep=1pt,topsep=2pt]
\item[\textbf{l.345}] \textsf{\small [phase convention]} $\displaystyle (w'-w)$\ $\longrightarrow$\ $\displaystyle (w-w')$
\end{itemize}
\begin{notebox}\textsf{\small Note, l.346:} [Checked: this is the $\log\frac{\vee}{k^{+}}$ term of the checked Row IV file (up to the phase); the term $-(1-\frac{k^{+}}{\vee})$ standing with it there belongs to the left-over terms.]\end{notebox}
\needspace{3\baselineskip}\noindent\rule{\linewidth}{0.3pt}\par\noindent\textbf{Lines 349--353} \quad {\small (in \texttt{align}, lines 348--353)}\par
\begin{itemize}[leftmargin=2.2em,itemsep=1pt,topsep=2pt]
\item[\textbf{l.349}] \textsf{\small [phase convention]} $\displaystyle \log{(\frac{\vee}{k^{+}})}\, \int_{x,x',z,w,w'}e^{-ik\cdot(w'-w)}$\ $\longrightarrow$\ $\displaystyle \log{(\frac{\vee}{k^{+}})}\, \int_{x,x',z,w,w'}e^{-ik\cdot(w-w')}$
\end{itemize}
\begin{notebox}\textsf{\small Note, l.353:} [Checked: equal to the expression above.]\end{notebox}
\needspace{3\baselineskip}\noindent\rule{\linewidth}{0.3pt}\par\noindent\textbf{Lines 357--361} \quad \textsf{Group III (with c.c.) $\rangle$ Row I} \quad {\small (in \texttt{align}, lines 356--361)}\par
\begin{itemize}[leftmargin=2.2em,itemsep=1pt,topsep=2pt]
\item[\textbf{l.357}] \textsf{\small [phase convention]} $\displaystyle \, \int_{x,y,x',z,w,w^{\prime}} e^{-ik(w^{\prime}-w)}\,$\ $\longrightarrow$\ $\displaystyle \, \int_{x,y,x',z,w,w^{\prime}} e^{-ik(w-w^{\prime})}\,$
\end{itemize}
\begin{notebox}\textsf{\small Note, l.361:} [Checked: equal to the $\log\frac{\vee}{k^{+}}$ terms of the checked Group III file, Row I (up to the phase).]\end{notebox}
\needspace{3\baselineskip}\noindent\rule{\linewidth}{0.3pt}\par\noindent\textbf{Lines 364} \quad {\small (in \texttt{align}, lines 363--392)}\par
\begin{itemize}[leftmargin=2.2em,itemsep=1pt,topsep=2pt]
\item[\textbf{l.364}] \textsf{\small [phase convention]} $\displaystyle \log{(\frac{\vee}{k^{+}})}\, \int_{x,x',z,w,w'}e^{-ik\cdot(w'-w)}$\ $\longrightarrow$\ $\displaystyle \log{(\frac{\vee}{k^{+}})}\, \int_{x,x',z,w,w'}e^{-ik\cdot(w-w')}$
\end{itemize}
\needspace{3\baselineskip}\noindent\rule{\linewidth}{0.3pt}\par\noindent\textbf{Lines 394--403} \quad {\small (in \texttt{align}, lines 393--403)}\par
\begin{itemize}[leftmargin=2.2em,itemsep=1pt,topsep=2pt]
\item[\textbf{l.394}] \textsf{\small [phase convention]} $\displaystyle \log{(\frac{\vee}{k^{+}})}\, \int_{x,x',z,w,w'}e^{-ik\cdot(w'-w)}$\ $\longrightarrow$\ $\displaystyle \log{(\frac{\vee}{k^{+}})}\, \int_{x,x',z,w,w'}e^{-ik\cdot(w-w')}$
\end{itemize}
\begin{notebox}\textsf{\small Note, l.403:} [Checked: the sum of the two lists is the symmetrised two-$\rho$ part of the expression above plus its c.c.]\end{notebox}
\needspace{3\baselineskip}\noindent\rule{\linewidth}{0.3pt}\par\noindent\textbf{Lines 406--410} \quad \textsf{Group III (with c.c.) $\rangle$ Row II} \quad {\small (in \texttt{align}, lines 405--410)}\par
\begin{itemize}[leftmargin=2.2em,itemsep=1pt,topsep=2pt]
\item[\textbf{l.406}] \textsf{\small [phase convention]} $\displaystyle \, \int_{x,y,x',z,w,w^{\prime}} e^{-ik(w^{\prime}-w)}\,$\ $\longrightarrow$\ $\displaystyle \, \int_{x,y,x',z,w,w^{\prime}} e^{-ik(w-w^{\prime})}\,$
\end{itemize}
\begin{notebox}\textsf{\small Note, l.410:} [Checked: equal to the $\log\frac{\vee}{k^{+}}$ terms of the checked Group III file, Row II (up to the phase).]\end{notebox}
\needspace{3\baselineskip}\noindent\rule{\linewidth}{0.3pt}\par\noindent\textbf{Lines 413} \quad {\small (in \texttt{align}, lines 412--444)}\par
\begin{itemize}[leftmargin=2.2em,itemsep=1pt,topsep=2pt]
\item[\textbf{l.413}] \textsf{\small [phase convention]} $\displaystyle \log{(\frac{\vee}{k^{+}})}\, \int_{x,x',z,w,w'}e^{-ik\cdot(w'-w)}$\ $\longrightarrow$\ $\displaystyle \log{(\frac{\vee}{k^{+}})}\, \int_{x,x',z,w,w'}e^{-ik\cdot(w-w')}$
\end{itemize}
\needspace{3\baselineskip}\noindent\rule{\linewidth}{0.3pt}\par\noindent\textbf{Lines 446--456} \quad {\small (in \texttt{align}, lines 445--456)}\par
\begin{itemize}[leftmargin=2.2em,itemsep=1pt,topsep=2pt]
\item[\textbf{l.446}] \textsf{\small [phase convention]} $\displaystyle \log{(\frac{\vee}{k^{+}})}\, \int_{x,x',z,w,w'}e^{-ik\cdot(w'-w)}$\ $\longrightarrow$\ $\displaystyle \log{(\frac{\vee}{k^{+}})}\, \int_{x,x',z,w,w'}e^{-ik\cdot(w-w')}$
\end{itemize}
\begin{notebox}\textsf{\small Note, l.456:} [Checked: the sum of the two lists is the symmetrised two-$\rho$ part of the expression above plus its c.c.]\end{notebox}
\needspace{3\baselineskip}\noindent\rule{\linewidth}{0.3pt}\par\noindent\textbf{Lines 460--464} \quad \textsf{Group III (with c.c.) $\rangle$ Row III} \quad {\small (in \texttt{align}, lines 459--464)}\par
\begin{itemize}[leftmargin=2.2em,itemsep=1pt,topsep=2pt]
\item[\textbf{l.460}] \textsf{\small [phase convention]} $\displaystyle \, \int_{w,x',y',z,y,x}\, e^{-ik\cdot (y'-w)}\,\frac{(y'-z)\cdot (x-z)\,(x'-y')\cdot (y-w)}{(y'-z)^{2}\,(x'-y')^{2}\,(y-w)^{2}\,(x-z)^{2}}$\ $\longrightarrow$\ $\displaystyle \, \int_{w,x',y',z,y,x}\, e^{-ik\cdot (w-y')}\,\frac{(y'-z)\cdot (x-z)\,(x'-y')\cdot (y-w)}{(y'-z)^{2}\,(x'-y')^{2}\,(y-w)^{2}\,(x-z)^{2}}$
\item[\textbf{l.461}] \textsf{\small [phase convention]} $\displaystyle \, \int_{w,x',y',z,y,x}\, e^{-ik\cdot (y'-w)}\,\frac{(y'-z)\cdot (x-z)\,(x'-y')\cdot (y-w)}{(y'-z)^{2}\,(x'-y')^{2}\,(y-w)^{2}\,(x-z)^{2}}$\ $\longrightarrow$\ $\displaystyle \, \int_{w,x',y',z,y,x}\, e^{-ik\cdot (w-y')}\,\frac{(y'-z)\cdot (x-z)\,(x'-y')\cdot (y-w)}{(y'-z)^{2}\,(x'-y')^{2}\,(y-w)^{2}\,(x-z)^{2}}$
\item[\textbf{l.462}] \textsf{\small [phase convention]} $\displaystyle \,\log{(\frac{\vee}{k^{+}})}\int_{w,x',w',z,y,x} e^{-ik\cdot (w'-w)}\,\frac{(w'-z)\cdot (y-w)\,(x'-z)\cdot (x-z)}{(w'-z)^{2}\,(y-w)^{2}\,(x'-z)^{2}\,(x-z)^{2}}$\ $\longrightarrow$\ $\displaystyle \,\log{(\frac{\vee}{k^{+}})}\int_{w,x',w',z,y,x} e^{-ik\cdot (w-w')}\,\frac{(w'-z)\cdot (y-w)\,(x'-z)\cdot (x-z)}{(w'-z)^{2}\,(y-w)^{2}\,(x'-z)^{2}\,(x-z)^{2}}$
\end{itemize}
\begin{notebox}\textsf{\small Note, l.464:} [All three terms are kept as the author wrote them (phase corrected): the second and third terms are the $\log\frac{\vee}{k^{+}}$ content of the $p^{+}$ integral in the author's split (kernel $\big[\frac{k^{+}}{p^{+}(p^{+}+k^{+})}\big]_{+}+\frac{1}{p^{+}+k^{+}}$ and subtracted $y'\to z$ term, see Other\_terms), which is the correct way to extract the complete $\log\frac{\vee}{k^{+}}$ of this row. The first two terms cancel each other, so the net coefficient is the third term.]\end{notebox}
\needspace{3\baselineskip}\noindent\rule{\linewidth}{0.3pt}\par\noindent\textbf{Lines 467--484} \quad {\small (in \texttt{align}, lines 466--484)}\par
\begin{itemize}[leftmargin=2.2em,itemsep=1pt,topsep=2pt]
\item[\textbf{l.467}] \textsf{\small [phase convention]} $\displaystyle \log{(\frac{\vee}{k^{+}})}\, \int_{x,x',z,w,w'}e^{-ik\cdot(w'-w)}$\ $\longrightarrow$\ $\displaystyle \log{(\frac{\vee}{k^{+}})}\, \int_{x,x',z,w,w'}e^{-ik\cdot(w-w')}$
\end{itemize}
\begin{notebox}\textsf{\small Note, l.484:} [Checked: this list is the symmetrised two-$\rho$ part of the third term above plus its c.c.]\end{notebox}
\needspace{3\baselineskip}\noindent\rule{\linewidth}{0.3pt}\par\noindent\textbf{Lines 487--493} \quad {\small (in \texttt{align}, lines 486--493)}\par
\begin{itemize}[leftmargin=2.2em,itemsep=1pt,topsep=2pt]
\item[\textbf{l.487}] \textsf{\small [phase convention]} $\displaystyle \,\int_{w,x',y',z,x}\, e^{-ik\cdot (y'-w)}\,$\ $\longrightarrow$\ $\displaystyle \,\int_{w,x',y',z,x}\, e^{-ik\cdot (w-y')}\,$
\item[\textbf{l.489}] \textsf{\small [phase convention]} $\displaystyle \,\int_{w,x',y',z,x}\, e^{-ik\cdot (y'-w)}\,$\ $\longrightarrow$\ $\displaystyle \,\int_{w,x',y',z,x}\, e^{-ik\cdot (w-y')}\,$
\item[\textbf{l.491}] \textsf{\small [phase convention]} $\displaystyle \,\log{(\frac{\vee}{k^{+}})}\int_{w,x',w',z,x} e^{-ik\cdot (w'-w)}\,\frac{(x'-z)\cdot (x-z)}{(x'-z)^{2}\,(x-z)^{2}}\,$\ $\longrightarrow$\ $\displaystyle \,\log{(\frac{\vee}{k^{+}})}\int_{w,x',w',z,x} e^{-ik\cdot (w-w')}\,\frac{(x'-z)\cdot (x-z)}{(x'-z)^{2}\,(x-z)^{2}}\,$
\end{itemize}
\begin{notebox}\textsf{\small Note, l.493:} [All three terms are kept as the author wrote them (phase corrected): the second and third terms are the $\log\frac{\vee}{k^{+}}$ content of the $p^{+}$ integral in the author's split (kernel $\big[\frac{k^{+}}{p^{+}(p^{+}+k^{+})}\big]_{+}+\frac{1}{p^{+}+k^{+}}$ and subtracted $y'\to z$ term, see Other\_terms), which is the correct way to extract the complete $\log\frac{\vee}{k^{+}}$ of this row. The first two terms cancel each other, so the net coefficient is the third term.]\end{notebox}
\needspace{3\baselineskip}\noindent\rule{\linewidth}{0.3pt}\par\noindent\textbf{Lines 496--516} \quad {\small (in \texttt{align}, lines 495--516)}\par
\begin{itemize}[leftmargin=2.2em,itemsep=1pt,topsep=2pt]
\item[\textbf{l.496}] \textsf{\small [phase convention]} $\displaystyle \log{(\frac{\vee}{k^{+}})}\, \int_{x,x',z,w,w'}e^{-ik\cdot(w'-w)}$\ $\longrightarrow$\ $\displaystyle \log{(\frac{\vee}{k^{+}})}\, \int_{x,x',z,w,w'}e^{-ik\cdot(w-w')}$
\end{itemize}
\begin{notebox}\textsf{\small Note, l.516:} [Checked: this list is the symmetrised form of the third term above plus its c.c.]\end{notebox}
\needspace{3\baselineskip}\noindent\rule{\linewidth}{0.3pt}\par\noindent\textbf{Lines 521--525} \quad \textsf{Group III (with c.c.) $\rangle$ Row V} \quad {\small (in \texttt{align}, lines 520--525)}\par
\begin{itemize}[leftmargin=2.2em,itemsep=1pt,topsep=2pt]
\item[\textbf{l.521}] \textsf{\small [phase convention]} $\displaystyle \log{(\frac{\vee}{k^{+}})} \int_{w,y',z,y,x,x'}\, e^{-ik\cdot (y'-w)}\,\frac{(y'-z)\cdot (x-z)\,(x'-y')\cdot (y-w)}{(y'-z)^{2}\,(x'-y')^{2}\,(y-w)^{2}\,(x-z)^{2}}$\ $\longrightarrow$\ $\displaystyle \log{(\frac{\vee}{k^{+}})} \int_{w,y',z,y,x,x'}\, e^{-ik\cdot (w-y')}\,\frac{(y'-z)\cdot (x-z)\,(x'-y')\cdot (y-w)}{(y'-z)^{2}\,(x'-y')^{2}\,(y-w)^{2}\,(x-z)^{2}}$
\item[\textbf{l.522}] \textsf{\small [phase convention]} $\displaystyle \log{(\frac{\vee}{k^{+}})} \int_{w,y',z,y,x,x'}\, e^{-ik\cdot (y'-w)}\,\frac{(y'-z)\cdot (x-z)\,(x'-y')\cdot (y-w)}{(y'-z)^{2}\,(x'-y')^{2}\,(y-w)^{2}\,(x-z)^{2}}$\ $\longrightarrow$\ $\displaystyle \log{(\frac{\vee}{k^{+}})} \int_{w,y',z,y,x,x'}\, e^{-ik\cdot (w-y')}\,\frac{(y'-z)\cdot (x-z)\,(x'-y')\cdot (y-w)}{(y'-z)^{2}\,(x'-y')^{2}\,(y-w)^{2}\,(x-z)^{2}}$
\item[\textbf{l.523}] \textsf{\small [phase convention]} $\displaystyle \,\log{(\frac{\vee}{k^{+}})}\int_{w,w',z,y,x,x'} e^{-ik\cdot (w'-w)}\,\frac{(w'-z)\cdot (y-w)\,(x'-z)\cdot (x-z)}{(w'-z)^{2}\,(y-w)^{2}\,(x'-z)^{2}\,(x-z)^{2}}$\ $\longrightarrow$\ $\displaystyle \,\log{(\frac{\vee}{k^{+}})}\int_{w,w',z,y,x,x'} e^{-ik\cdot (w-w')}\,\frac{(w'-z)\cdot (y-w)\,(x'-z)\cdot (x-z)}{(w'-z)^{2}\,(y-w)^{2}\,(x'-z)^{2}\,(x-z)^{2}}$
\end{itemize}
\begin{notebox}\textsf{\small Note, l.525:} [All three terms are kept as the author wrote them (phase corrected): the second and third terms are the $\log\frac{\vee}{k^{+}}$ content of the $p^{+}$ integral in the author's split (kernel $\big[\frac{k^{+}}{p^{+}(p^{+}+k^{+})}\big]_{+}+\frac{1}{p^{+}+k^{+}}$ and subtracted $y'\to z$ term, see Other\_terms), which is the correct way to extract the complete $\log\frac{\vee}{k^{+}}$ of this row. The first two terms cancel each other, so the net coefficient is the third term.]\end{notebox}
\needspace{3\baselineskip}\noindent\rule{\linewidth}{0.3pt}\par\noindent\textbf{Lines 528--544} \quad {\small (in \texttt{align}, lines 527--544)}\par
\begin{itemize}[leftmargin=2.2em,itemsep=1pt,topsep=2pt]
\item[\textbf{l.528}] \textsf{\small [phase convention]} $\displaystyle \log{(\frac{\vee}{k^{+}})}\, \int_{x,x',z,w,w'}e^{-ik\cdot(w'-w)}$\ $\longrightarrow$\ $\displaystyle \log{(\frac{\vee}{k^{+}})}\, \int_{x,x',z,w,w'}e^{-ik\cdot(w-w')}$
\end{itemize}
\begin{notebox}\textsf{\small Note, l.544:} [Checked: this list is the symmetrised two-$\rho$ part of the third term above plus its c.c.]\end{notebox}
\needspace{3\baselineskip}\noindent\rule{\linewidth}{0.3pt}\par\noindent\textbf{Lines 547--551} \quad \textsf{Group III (with c.c.) $\rangle$ Row VI} \quad {\small (in \texttt{align}, lines 546--551)}\par
\begin{itemize}[leftmargin=2.2em,itemsep=1pt,topsep=2pt]
\item[\textbf{l.547}] \textsf{\small [phase convention]} $\displaystyle \, \int_{w,y',z,y,x,x'}\, e^{-ik\cdot (y'-w)}\,\frac{(y'-z)\cdot (x-z)\,(x'-y')\cdot (y-w)}{(y'-z)^{2}\,(x'-y')^{2}\,(y-w)^{2}\,(x-z)^{2}}$\ $\longrightarrow$\ $\displaystyle \, \int_{w,y',z,y,x,x'}\, e^{-ik\cdot (w-y')}\,\frac{(y'-z)\cdot (x-z)\,(x'-y')\cdot (y-w)}{(y'-z)^{2}\,(x'-y')^{2}\,(y-w)^{2}\,(x-z)^{2}}$
\item[\textbf{l.548}] \textsf{\small [phase convention]} $\displaystyle \, \int_{w,y',z,y,x,x'}\, e^{-ik\cdot (y'-w)}\,\frac{(y'-z)\cdot (x-z)\,(x'-y')\cdot (y-w)}{(y'-z)^{2}\,(x'-y')^{2}\,(y-w)^{2}\,(x-z)^{2}}$\ $\longrightarrow$\ $\displaystyle \, \int_{w,y',z,y,x,x'}\, e^{-ik\cdot (w-y')}\,\frac{(y'-z)\cdot (x-z)\,(x'-y')\cdot (y-w)}{(y'-z)^{2}\,(x'-y')^{2}\,(y-w)^{2}\,(x-z)^{2}}$
\item[\textbf{l.549}] \textsf{\small [phase convention]} $\displaystyle \,\log{(\frac{\vee}{k^{+}})}\int_{w,w',z,y,x,x'} e^{-ik\cdot (w'-w)}\,\frac{(w'-z)\cdot (y-w)\,(x'-z)\cdot (x-z)}{(w'-z)^{2}\,(y-w)^{2}\,(x'-z)^{2}\,(x-z)^{2}}$\ $\longrightarrow$\ $\displaystyle \,\log{(\frac{\vee}{k^{+}})}\int_{w,w',z,y,x,x'} e^{-ik\cdot (w-w')}\,\frac{(w'-z)\cdot (y-w)\,(x'-z)\cdot (x-z)}{(w'-z)^{2}\,(y-w)^{2}\,(x'-z)^{2}\,(x-z)^{2}}$
\end{itemize}
\begin{notebox}\textsf{\small Note, l.551:} [All three terms are kept as the author wrote them (phase corrected): the second and third terms are the $\log\frac{\vee}{k^{+}}$ content of the $p^{+}$ integral in the author's split (kernel $\big[\frac{k^{+}}{p^{+}(p^{+}+k^{+})}\big]_{+}+\frac{1}{p^{+}+k^{+}}$ and subtracted $y'\to z$ term, see Other\_terms), which is the correct way to extract the complete $\log\frac{\vee}{k^{+}}$ of this row. The first two terms cancel each other, so the net coefficient is the third term.]\end{notebox}
\needspace{3\baselineskip}\noindent\rule{\linewidth}{0.3pt}\par\noindent\textbf{Lines 554--568} \quad {\small (in \texttt{align}, lines 553--568)}\par
\begin{itemize}[leftmargin=2.2em,itemsep=1pt,topsep=2pt]
\item[\textbf{l.554}] \textsf{\small [phase convention]} $\displaystyle \log{(\frac{\vee}{k^{+}})}\, \int_{x,x',z,w,w'}e^{-ik\cdot(w'-w)}$\ $\longrightarrow$\ $\displaystyle \log{(\frac{\vee}{k^{+}})}\, \int_{x,x',z,w,w'}e^{-ik\cdot(w-w')}$
\end{itemize}
\begin{notebox}\textsf{\small Note, l.568:} [Checked: this list is the symmetrised two-$\rho$ part of the third term above plus its c.c.]\end{notebox}
\needspace{3\baselineskip}\noindent\rule{\linewidth}{0.3pt}\par\noindent\textbf{Lines 571--605} \quad \textsf{BFKL contribution} \quad {\small (in \texttt{align}, lines 570--605)}\par
\begin{itemize}[leftmargin=2.2em,itemsep=1pt,topsep=2pt]
\item[\textbf{l.571}] \textsf{\small [phase convention]} $\displaystyle \Big]\,\int_{x,x',z,w,w'}e^{-ik\cdot(w'-w)}$\ $\longrightarrow$\ $\displaystyle \Big]\,\int_{x,x',z,w,w'}e^{-ik\cdot(w-w')}$
\item[\textbf{l.579}] \textsf{\small [operator ordering]} $\displaystyle f^{cde}\,U^{hd}(x')\,U^{ha}(w)\,\rho^{b}(x)\,\rho^{e}(x')$\ $\longrightarrow$\ $\displaystyle f^{cde}\,U^{hd}(x')\,U^{ha}(w)\,\rho^{e}(x')\,\rho^{b}(x)$
\item[\textbf{l.581}] \textsf{\small [operator ordering]} $\displaystyle f^{cde}\,U^{hd}(w')\,U^{ha}(w)\,\rho^{b}(x)\,\rho^{e}(x')$\ $\longrightarrow$\ $\displaystyle f^{cde}\,U^{hd}(w')\,U^{ha}(w)\,\rho^{e}(x')\,\rho^{b}(x)$
\item[\textbf{l.594}] \textsf{\small [operator ordering]} $\displaystyle f^{bde}\,U^{hd}(w')\,U^{ha}(w)\,\rho^{e}(x')\,\rho^{c}(x)$\ $\longrightarrow$\ $\displaystyle f^{bde}\,U^{hd}(w')\,U^{ha}(w)\,\rho^{c}(x)\,\rho^{e}(x')$
\end{itemize}
\begin{notebox}\textsf{\small Note, l.605:} [Phase: with the conjugated bracket (primed coordinates, standing left) at $w'$ the phase is $e^{-ik\cdot(w-w')}$, as in the corrected cross section; only the part antisymmetric under primed $\leftrightarrow$ unprimed (the contact terms at $x=x'$) is sensitive to it. In three terms the two $\rho$'s stood in the opposite order to the grouped form below; exchanged, so that this list equals the grouped form also in the contact terms]\end{notebox}
\needspace{3\baselineskip}\noindent\rule{\linewidth}{0.3pt}\par\noindent\textbf{Lines 607--633} \quad {\small (in \texttt{align}, lines 606--633)}\par
\begin{itemize}[leftmargin=2.2em,itemsep=1pt,topsep=2pt]
\item[\textbf{l.607}] \textsf{\small [phase convention]} $\displaystyle \Big]\,\int_{x,x',z,w,w'}e^{-ik\cdot(w'-w)}$\ $\longrightarrow$\ $\displaystyle \Big]\,\int_{x,x',z,w,w'}e^{-ik\cdot(w-w')}$
\end{itemize}
\begin{notebox}\textsf{\small Note, l.633:} [Checked against one step of projectile evolution of the LO cross section: one gluon with $k^{+}<p^{+}<\vee$ at $z$ added to the projectile state with the LO coefficients $\mathfrak{A}^{(1)}$, $\mathfrak{N}$, $\rho^{a}(x)\to\rho^{a}(x)+T^{a}\delta^{2}(x-z)$, the gluon scattering with $U(z)$. Prefactor, kernels, signs and colour structures agree term by term, and the result vanishes at $U=1$ and is finite at $z\to x,x'$. Remark (flagged, not changed): with non-commuting $\rho$ the evolution gives $\tfrac12\{\rho(x'),\rho(x)\}$ in the four blocks with $\mathcal{K}(x'-w')\cdot\mathcal{K}(w-z)$ and $\mathcal{K}(x-w)\cdot\mathcal{K}(w'-z)$ (the orders written in the other blocks are the ones it gives); the fixed orders written in these four blocks differ from that by one-$\rho$ contact terms $\propto f^{abc}\rho^{c}(x)\,\delta^{2}(x-x')$, which drop out of the colour-singlet form]\end{notebox}
\needspace{3\baselineskip}\noindent\rule{\linewidth}{0.3pt}\par\noindent\textbf{Lines 636--644} \quad {\small (in \texttt{align}, lines 635--644)}\par
\begin{itemize}[leftmargin=2.2em,itemsep=1pt,topsep=2pt]
\item[\textbf{l.636}] \textsf{\small [phase convention]} $\displaystyle \Big]\,\int_{x,y,z,w,w'}e^{-ik(w'-w)}\,\langle\rho^{f}(x)$\ $\longrightarrow$\ $\displaystyle \Big]\,\int_{x,y,z,w,w'}e^{-ik(w-w')}\,\langle\rho^{f}(x)$
\end{itemize}
\begin{notebox}\textsf{\small Note, l.644:} [Follows from the previous form with $\langle\rho^{a}(x)\rho^{b}(y)\rangle=\delta^{ab}\langle\rho^{f}(x)\rho^{f}(y)\rangle/(N_c^{2}-1)$, which is symmetric under $x\leftrightarrow y$ by the algebra; $U$ and the traces are in the adjoint representation; the result does not depend on the sign of the phase]\end{notebox}
\needspace{15\baselineskip}\bigskip
\subsection{Other terms (v1)}\label{sec:file-other_terms}
\noindent\textbf{Original:} \texttt{/root/.claude/uploads/0ab97fcf-94a2-5aec-8bad-586400481685/99883d04-Other\_terms.tex}\\\textbf{Checked copy:} \texttt{/home/user/correlations/nlo\_review/evolution\_checked/Other\_terms\_checked.tex}\\
Lines: 237 (original), 237 (checked copy). Coloured groups of the check: 147 (147 top-level); author\textquotesingle s own coloured groups left unchanged: 0.\\
Records by type: phase convention 61; limits/range 1; missing term added 14; term removed/replaced 2; typo/LaTeX 1; note only 10.
\par\noindent Unmarked layout-only edits (line breaks \texttt{\textbackslash notag\textbackslash\textbackslash\ \&}, \texttt{\textbackslash par\textbackslash noindent} around inserted notes, spacing), checked-copy lines: 35, 45, 46, 111, 112, 130, 133, 143, 192, 212, 223, 231, 236.
\par\noindent Unmarked edits that are not layout only (listed with their block below): 2.
\subsubsection*{Changes, block by block}
\needspace{3\baselineskip}\noindent\rule{\linewidth}{0.3pt}\par\noindent\textbf{Lines 35}\par
\begin{notebox}\textsf{\small Note, l.35:} \textbf{Checked copy.} Marking: corrected or added material is in red; material removed from the author's text is kept in orange (cancelled); notes of the check are in red in square brackets. Phase: $e^{-ik\cdot(w-w')}$ everywhere (conjugated bracket, with the primed coordinates and standing left, at $w'$); this does not change the symmetrised two-$\rho$ parts, while the three-$\rho$ and one-$\rho$ parts change sign together. Result: the file was compared, row by row, with what remains of the checked cross-section files (Group I, II, II Row IV, III) when exactly the terms with the coefficients $\log\frac{\vee}{k^{+}}$ and $\log\frac{k^{+}}{\Lambda}$ are taken out; it agrees once the terms marked in red are added (Group I Rows II, III, V; Group II Row IV) and the lower limit of Group II Row IV is corrected. Group III Rows III, V, VI are kept in the author's original split of the $p^{+}$ integral, which is the correct one (notes there). Further conventions: As the title says, $\Lambda\to0$ is taken here; the terms of the checked rows that are dropped by this are $\log(1+\frac{\Lambda}{k^{+}})$, $\log(1-\frac{\Lambda}{k^{+}})$, $\log\frac{k^{+}+\Lambda}{k^{+}-\Lambda}$ times the kernels of the corresponding $\log\frac{\vee}{k^{+}}$ or $\log\frac{k^{+}}{\Lambda}$ terms, and the $\Lambda$ inside the transverse logarithms; the $1/\Lambda$ terms of Group I Row III cancel exactly. $\Lambda$ is kept only where the termwise limit does not exist (last term of $N_{3b}$ of Group I Row III). All logarithms with a difference of transverse distances in the argument are $\log|\dots|$ (principal value). By the author's convention the $\frac{2}{\epsilon}$ poles are not written at this stage. The pole terms dropped in this way are: Group I Row II, $-\frac{2}{\epsilon}\big[\frac{11}{3}+\log(1-\frac{k^{+}}{\vee})\big]$ times the LO-type structure of $N_{1}$; Group I Row V, $-\frac{11}{6}\cdot\frac{2}{\epsilon}$ in $N_{2b\,\rm UV}$ (it cancels the $\frac{11}{3}$ part of the former); Group II Row IV, the $\frac{2}{\epsilon}$ of the DGLAP, RC and $\log(1-\frac{k^{+}}{\vee})$ terms.\end{notebox}
\needspace{3\baselineskip}\noindent\rule{\linewidth}{0.3pt}\par\noindent\textbf{Lines 38} \quad \textsf{Group I (with c.c.) $\rangle$ Row I}\par
\begin{notebox}\textsf{\small Note, l.38:} [Checked: remains NA. Row I is a four-$\rho$ term multiplied by $\log\frac{\vee}{\Lambda}=\log\frac{\vee}{k^{+}}+\log\frac{k^{+}}{\Lambda}$ only; its commutator term cancels (see the resolution of the operator ordering), so nothing is left over.]\end{notebox}
\needspace{3\baselineskip}\noindent\rule{\linewidth}{0.3pt}\par\noindent\textbf{Lines 45--46} \quad \textsf{Group I (with c.c.) $\rangle$ Row II} \quad {\small (in \texttt{align}, lines 44--46)}\par
\begin{itemize}[leftmargin=2.2em,itemsep=1pt,topsep=2pt]
\item[\textbf{l.45}] \textsf{\small [phase convention]} $\displaystyle w^\prime-w$\ $\longrightarrow$\ $\displaystyle w-w^\prime$
\item[\textbf{l.45}] \textsf{\small [missing term added]} \emph{added:} $\displaystyle [-2\gamma-\log{(\frac{(x-w)^{2}\,\mu^{2}_{\overline{MS}}}{4})}][\frac{11}{3}+\log{(1-\frac{k^{+}}{\vee})}]$\quad{\color{gray}\small(in context: }$\displaystyle [-2\gamma-\log{(\frac{(x-w)^{2}\,\mu^{2}_{\overline{MS}}}{4})}][\frac{11}{3}+\log{(1-\frac{k^{+}}{\vee})}]-\frac{67}{9}${\color{gray}\small)}
\end{itemize}
\begin{notebox}\textsf{\small Note, l.46:} [Added: the first bracket times $\frac{11}{3}+\log(1-\frac{k^{+}}{\vee})$ is what remains of the corresponding term of the checked Row II when $-2\log\frac{k^{+}}{\Lambda}$ is taken out; it was commented out in the original. The $-\frac{2}{\epsilon}$ of the bracket is not written (see the note at the top). Its $\frac{11}{3}$ part cancels against the $-\frac{11}{6}$ term of Row V ($2b$, UV) below. The last term is the product $\log(1-\frac{k^{+}}{\vee})\,\log\frac{\vee}{k^{+}}$, i.e. it is linear in the BFKL logarithm; by the author's decision it stays here and not in the BFKL file.]\end{notebox}
\needspace{3\baselineskip}\noindent\rule{\linewidth}{0.3pt}\par\noindent\textbf{Lines 48} \quad {\small (in \texttt{align}, lines 47--52)}\par
\begin{itemize}[leftmargin=2.2em,itemsep=1pt,topsep=2pt]
\item[\textbf{l.48}] \textsf{\small [typo/LaTeX]} \emph{added:} $\displaystyle ,x'$\quad{\color{gray}\small(in context: }$\displaystyle \,\frac{1}{k^{+}}\int_{w,w',x'} e^{-ik\cdot(w-w')}${\color{gray}\small)}
\item[\textbf{l.48}] \textsf{\small [phase convention]} $\displaystyle w'-w$\ $\longrightarrow$\ $\displaystyle w-w'$
\end{itemize}
\needspace{3\baselineskip}\noindent\rule{\linewidth}{0.3pt}\par\noindent\textbf{Lines 63} \quad \textsf{Group I (with c.c.) $\rangle$ Row III} \quad {\small (in \texttt{align}, lines 62--68)}\par
\begin{itemize}[leftmargin=2.2em,itemsep=1pt,topsep=2pt]
\item[\textbf{l.63}] \textsf{\small [phase convention]} $\displaystyle w^\prime-w$\ $\longrightarrow$\ $\displaystyle w-w^\prime$
\end{itemize}
\needspace{3\baselineskip}\noindent\rule{\linewidth}{0.3pt}\par\noindent\textbf{Lines 70} \quad {\small (in \texttt{align}, lines 69--77)}\par
\begin{itemize}[leftmargin=2.2em,itemsep=1pt,topsep=2pt]
\item[\textbf{l.70}] \textsf{\small [phase convention]} $\displaystyle w^\prime-w$\ $\longrightarrow$\ $\displaystyle w-w^\prime$
\end{itemize}
\needspace{3\baselineskip}\noindent\rule{\linewidth}{0.3pt}\par\noindent\textbf{Lines 79} \quad {\small (in \texttt{align}, lines 78--101)}\par
\begin{itemize}[leftmargin=2.2em,itemsep=1pt,topsep=2pt]
\item[\textbf{l.79}] \textsf{\small [phase convention]} $\displaystyle w'-w$\ $\longrightarrow$\ $\displaystyle w-w'$
\end{itemize}
\needspace{3\baselineskip}\noindent\rule{\linewidth}{0.3pt}\par\noindent\textbf{Lines 103--112} \quad {\small (in \texttt{align}, lines 102--112)}\par
\begin{itemize}[leftmargin=2.2em,itemsep=1pt,topsep=2pt]
\item[\textbf{l.103}] \textsf{\small [phase convention]} $\displaystyle w'-w$\ $\longrightarrow$\ $\displaystyle w-w'$
\item[\textbf{l.111}] \textsf{\small [term removed/replaced]} \emph{removed:} $\displaystyle -\frac{ig^{4}f^{abd}}{16\pi^{5}}\,\frac{1}{k^{+}}\,\int_{x,y,z,x'}\log{[\frac{(z-w)^{2}}{(x-w)^{2}}]}\,\frac{(x'-w')\cdot(y-w)\, (x-z)\cdot (w-z)}{(x'-w')^{2}(y-w)^{2}(x-z)^{2}(w-z)^{2}}$
\item[\textbf{l.111}] \textsf{\small [term removed/replaced]} \emph{removed:} $\displaystyle [U^{ab'}(x')\,\rho^{b'}(x')-U^{ab'}(w')\, \rho^{b'}(x')]\times [U^{de}(x)\,U^{bc}(w)\, \rho^{e}(x)\,\rho^{c}(y)-U^{de}(x)\,U^{bc}(y)\,\rho^{e}(x)\,\rho^{c}(y)]$
\item[\textbf{l.111}] \textsf{\small [missing term added]} \emph{added:} $\displaystyle -\frac{ig^{4}f^{abd}}{16\pi^{5}}\,\frac{1}{k^{+}}\,\int_{x,y,z,x'}\log{[\frac{[k^{+}\,(x-w)-(k^{+}-\Lambda)\,(x-z)]^{2}}{[k^{+}\,(x-w)-\Lambda\,(x-z)]^{2}}]}\,\frac{(x'-w')\cdot(y-z)}{(x'-w')^{2}(y-z)^{2}\,(z-w)^{2}}$
\item[\textbf{l.111}] \textsf{\small [missing term added]} \emph{added:} $\displaystyle [U^{ab'}(x')\,\rho^{b'}(x')-U^{ab'}(w')\, \rho^{b'}(x')]\times [U^{de}(x)\,U^{bc}(z)\, \rho^{e}(x)\,\rho^{c}(y)-U^{de}(x)\,U^{bc}(y)\,\rho^{e}(x)\,\rho^{c}(y)]$
\end{itemize}
\begin{notebox}\textsf{\small Note, l.112:} [Last term of $N_{3b}$: the orange term is not what the checked Row III gives; the red term is (last term of $N_{3b}$ there). $\Lambda$ is kept in it: for $\Lambda\to0$ its logarithm becomes $\log\frac{(z-w)^{2}}{(x-w)^{2}}$ only for $|z-w|\gg\frac{\Lambda}{k^{+}}|x-w|$, and with the kernel $1/(z-w)^{2}$ neither this term nor the $\log\frac{k^{+}}{\Lambda}$ term of $N_{3a}$ with the same kernel is integrable at $z=w$ in that limit; at $z=w$ the logarithm equals $-2\log\frac{k^{+}-\Lambda}{\Lambda}$ and the two terms cancel there.]\end{notebox}
\needspace{3\baselineskip}\noindent\rule{\linewidth}{0.3pt}\par\noindent\textbf{Lines 119} \quad \textsf{Group I (with c.c.) $\rangle$ Row IV} \quad {\small (in \texttt{align}, lines 117--122)}\par
\begin{itemize}[leftmargin=2.2em,itemsep=1pt,topsep=2pt]
\item[\textbf{l.119}] \textsf{\small [phase convention]} $\displaystyle w^\prime-w$\ $\longrightarrow$\ $\displaystyle w-w^\prime$
\end{itemize}
\needspace{3\baselineskip}\noindent\rule{\linewidth}{0.3pt}\par\noindent\textbf{Lines 126} \quad \textsf{Group I (with c.c.) $\rangle$ Row V} \quad {\small (in \texttt{align}, lines 125--127)}\par
\begin{itemize}[leftmargin=2.2em,itemsep=1pt,topsep=2pt]
\item[\textbf{l.126}] \textsf{\small [phase convention]} $\displaystyle w^\prime-z$\ $\longrightarrow$\ $\displaystyle z-w^\prime$
\item[\textbf{l.126}] \textsf{\small [phase convention]} $\displaystyle z-x$\ $\longrightarrow$\ $\displaystyle x-z$
\item[\textbf{l.126}] \textsf{\small [phase convention]} $\displaystyle w^\prime-x$\ $\longrightarrow$\ $\displaystyle x-w^\prime$
\item[\textbf{l.126}] \textsf{\small [phase convention]} $\displaystyle x-z$\ $\longrightarrow$\ $\displaystyle z-x$
\end{itemize}
\needspace{3\baselineskip}\noindent\rule{\linewidth}{0.3pt}\par\noindent\textbf{Lines 130} \quad {\small (in \texttt{align}, lines 129--131)}\par
\begin{itemize}[leftmargin=2.2em,itemsep=1pt,topsep=2pt]
\item[\textbf{l.130}] \textsf{\small [phase convention]} $\displaystyle w'-z$\ $\longrightarrow$\ $\displaystyle z-w'$
\item[\textbf{l.130}] \textsf{\small [phase convention]} $\displaystyle z-x$\ $\longrightarrow$\ $\displaystyle x-z$
\item[\textbf{l.130}] \textsf{\small [phase convention]} $\displaystyle w'-x$\ $\longrightarrow$\ $\displaystyle x-w'$
\item[\textbf{l.130}] \textsf{\small [phase convention]} $\displaystyle x-z$\ $\longrightarrow$\ $\displaystyle z-x$
\end{itemize}
\needspace{3\baselineskip}\noindent\rule{\linewidth}{0.3pt}\par\noindent\textbf{Lines 133--141} \quad {\small (in \texttt{align}, lines 132--143)}\par
\begin{itemize}[leftmargin=2.2em,itemsep=1pt,topsep=2pt]
\item[\textbf{l.133}] \textsf{\small [phase convention]} $\displaystyle w'-z$\ $\longrightarrow$\ $\displaystyle z-w'$
\item[\textbf{l.133}] \textsf{\small [phase convention]} $\displaystyle z-x$\ $\longrightarrow$\ $\displaystyle x-z$
\item[\textbf{l.141}] \textsf{\small [phase convention]} $\displaystyle w'-x$\ $\longrightarrow$\ $\displaystyle x-w'$
\end{itemize}
\needspace{3\baselineskip}\noindent\rule{\linewidth}{0.3pt}\par\noindent\textbf{Lines 143} \quad {\small (in \texttt{align}, lines 143--143)}\par
\begin{itemize}[leftmargin=2.2em,itemsep=1pt,topsep=2pt]
\item[\textbf{l.143}] \textsf{\small [missing term added]} \emph{added:} $\displaystyle \frac{d^{3}N_{2b\,\rm UV}}{d^{3}k}$
\item[\textbf{l.143}] \textsf{\small [phase convention]} \emph{added:} $\displaystyle =-\frac{1}{(2\pi)^{3}}\,\frac{g^{4}N_{c}}{16\pi^{4}}\,\frac{1}{k^{+}}\,\frac{11}{6}\int_{x,y,x',w'}e^{-ik\cdot(x-w')}\,[2\gamma_{E}+\log\frac{(y-x)^{2}\,\bar\mu^{2}_{\overline{MS}}}{4}]\frac{(x'-w')\cdot(y-x)}{(x'-w')^{2}\,(y-x)^{2}}$
\item[\textbf{l.143}] \textsf{\small [missing term added]} \emph{added:} $\displaystyle \times\Big[U^{ab'}(x')-U^{ab'}(w')\Big]\Big[U^{af}(y)-U^{af}(x)\Big]\rho^{b'}(x')\,\rho^{f}(y)$
\end{itemize}
\begin{notebox}\textsf{\small Note, l.143:} [Added: what remains of $N_{2b\,\rm UV}$ of the checked Row V, $\big[2\log\frac{k^{+}}{\Lambda}-\frac{11}{6}\big]\times$$\big[\frac{2}{\epsilon}+2\gamma_{E}+\log\dots\big]$, when the $\log\frac{k^{+}}{\Lambda}$ term is taken out; the $\frac{2}{\epsilon}$ is not written. With $x\to w$, $y\to x$ it is minus the $\frac{11}{3}$ part of the pole term of Row II.]\end{notebox}
\needspace{3\baselineskip}\noindent\rule{\linewidth}{0.3pt}\par\noindent\textbf{Lines 148} \quad \textsf{Group II $\rangle$ Row I} \quad {\small (in \texttt{align}, lines 147--151)}\par
\begin{itemize}[leftmargin=2.2em,itemsep=1pt,topsep=2pt]
\item[\textbf{l.148}] \textsf{\small [phase convention]} $\displaystyle w^{\prime}-w$\ $\longrightarrow$\ $\displaystyle w-w^{\prime}$
\end{itemize}
\needspace{3\baselineskip}\noindent\rule{\linewidth}{0.3pt}\par\noindent\textbf{Lines 153} \quad {\small (in \texttt{align}, lines 152--156)}\par
\begin{itemize}[leftmargin=2.2em,itemsep=1pt,topsep=2pt]
\item[\textbf{l.153}] \textsf{\small [phase convention]} $\displaystyle w^{\prime}-w$\ $\longrightarrow$\ $\displaystyle w-w^{\prime}$
\end{itemize}
\needspace{3\baselineskip}\noindent\rule{\linewidth}{0.3pt}\par\noindent\textbf{Lines 159} \quad {\small (in \texttt{align}, lines 158--161)}\par
\begin{itemize}[leftmargin=2.2em,itemsep=1pt,topsep=2pt]
\item[\textbf{l.159}] \textsf{\small [phase convention]} $\displaystyle w^{\prime}-w$\ $\longrightarrow$\ $\displaystyle w-w^{\prime}$
\end{itemize}
\needspace{3\baselineskip}\noindent\rule{\linewidth}{0.3pt}\par\noindent\textbf{Lines 163--165} \quad {\small (in \texttt{align}, lines 162--166)}\par
\begin{itemize}[leftmargin=2.2em,itemsep=1pt,topsep=2pt]
\item[\textbf{l.163}] \textsf{\small [phase convention]} $\displaystyle w^{\prime}-w$\ $\longrightarrow$\ $\displaystyle w-w^{\prime}$
\item[\textbf{l.164}] \textsf{\small [phase convention]} $\displaystyle w^{\prime}-w$\ $\longrightarrow$\ $\displaystyle w-w^{\prime}$
\item[\textbf{l.165}] \textsf{\small [phase convention]} $\displaystyle w^{\prime}-w$\ $\longrightarrow$\ $\displaystyle w-w^{\prime}$
\end{itemize}
\needspace{3\baselineskip}\noindent\rule{\linewidth}{0.3pt}\par\noindent\textbf{Lines 168--169} \quad {\small (in \texttt{align}, lines 167--170)}\par
\begin{itemize}[leftmargin=2.2em,itemsep=1pt,topsep=2pt]
\item[\textbf{l.168}] \textsf{\small [phase convention]} $\displaystyle w^{\prime}-w$\ $\longrightarrow$\ $\displaystyle w-w^{\prime}$
\item[\textbf{l.169}] \textsf{\small [phase convention]} $\displaystyle w^{\prime}-w$\ $\longrightarrow$\ $\displaystyle w-w^{\prime}$
\end{itemize}
\needspace{3\baselineskip}\noindent\rule{\linewidth}{0.3pt}\par\noindent\textbf{Lines 172--173} \quad {\small (in \texttt{align}, lines 171--174)}\par
\begin{itemize}[leftmargin=2.2em,itemsep=1pt,topsep=2pt]
\item[\textbf{l.172}] \textsf{\small [phase convention]} $\displaystyle w^{\prime}-w$\ $\longrightarrow$\ $\displaystyle w-w^{\prime}$
\item[\textbf{l.173}] \textsf{\small [phase convention]} $\displaystyle w^{\prime}-w$\ $\longrightarrow$\ $\displaystyle w-w^{\prime}$
\end{itemize}
\needspace{3\baselineskip}\noindent\rule{\linewidth}{0.3pt}\par\noindent\textbf{Lines 176--177} \quad {\small (in \texttt{align}, lines 175--178)}\par
\begin{itemize}[leftmargin=2.2em,itemsep=1pt,topsep=2pt]
\item[\textbf{l.176}] \textsf{\small [phase convention]} $\displaystyle w^{\prime}-w$\ $\longrightarrow$\ $\displaystyle w-w^{\prime}$
\item[\textbf{l.177}] \textsf{\small [phase convention]} $\displaystyle w^{\prime}-w$\ $\longrightarrow$\ $\displaystyle w-w^{\prime}$
\end{itemize}
\needspace{3\baselineskip}\noindent\rule{\linewidth}{0.3pt}\par\noindent\textbf{Lines 180--182} \quad {\small (in \texttt{align}, lines 179--183)}\par
\begin{itemize}[leftmargin=2.2em,itemsep=1pt,topsep=2pt]
\item[\textbf{l.180}] \textsf{\small [phase convention]} $\displaystyle w^{\prime}-w$\ $\longrightarrow$\ $\displaystyle w-w^{\prime}$
\item[\textbf{l.181}] \textsf{\small [phase convention]} $\displaystyle w^{\prime}-w$\ $\longrightarrow$\ $\displaystyle w-w^{\prime}$
\item[\textbf{l.182}] \textsf{\small [phase convention]} $\displaystyle w^{\prime}-w$\ $\longrightarrow$\ $\displaystyle w-w^{\prime}$
\end{itemize}
\needspace{3\baselineskip}\noindent\rule{\linewidth}{0.3pt}\par\noindent\textbf{Lines 190--192} \quad \textsf{Group II $\rangle$ Row IV} \quad {\small (in \texttt{align}, lines 189--192)}\par
\begin{itemize}[leftmargin=2.2em,itemsep=1pt,topsep=2pt]
\item[\textbf{l.190}] \textsf{\small [limits/range]} $\displaystyle 0$\ $\longrightarrow$\ $\displaystyle \frac{k^{+}}{\vee}$
\item[\textbf{l.190}] \textsf{\small [phase convention]} $\displaystyle y'-y$\ $\longrightarrow$\ $\displaystyle y-y'$
\end{itemize}
\begin{notebox}\textsf{\small Note, l.192:} [The lower limit is $\frac{k^{+}}{\vee}$. In the checked Row IV file the coefficients of the evolution terms are $\log\frac{\vee}{k^{+}}-\big(1-\frac{k^{+}}{\vee}\big)$ and $\log\frac{k^{+}}{\Lambda}+\log\big(1-\frac{k^{+}}{\vee}\big)$; when only the two logarithms are taken out, the following terms are left in addition to the finite term above (in the checked file each bracket $[2\gamma_{E}+\log\dots]$ below is $[\frac{2}{\epsilon}+2\gamma_{E}+\log\dots]$; the $\frac{2}{\epsilon}$ is not written here; $\mathcal P_{gg}(\xi)=2N_{c}\big[\xi(1-\xi)+\frac{1-\xi}{\xi}+\frac{\xi}{(1-\xi)_{+}}\big]+\frac{\beta_{0}}{2}\delta(1-\xi)$, $\int_{k^{+}/\vee}^{1}d\xi\,\frac{\xi}{(1-\xi)_{+}}f(\xi)\equiv\int_{k^{+}/\vee}^{1}d\xi\,\frac{\xi f(\xi)-f(1)}{1-\xi}$):]\end{notebox}
\needspace{3\baselineskip}\noindent\rule{\linewidth}{0.3pt}\par\noindent\textbf{Lines 192} \quad {\small (in \texttt{align}, lines 192--192)}\par
\begin{itemize}[leftmargin=2.2em,itemsep=1pt,topsep=2pt]
\item[\textbf{l.192}] \textsf{\small [missing term added]} \emph{added:} $\displaystyle \frac{d^{3}N_{-(1-k^{+}/\vee)}}{d^{3}k}$
\item[\textbf{l.192}] \textsf{\small [missing term added]} \emph{added:} $\displaystyle = -\frac{1}{(2\pi)^{3}}\,\frac{g^{4}\,N_{c}}{2\pi^{3}}\,\frac{1}{k^{2}}\Big[\frac{1}{k^{+}}-\frac{1}{\vee}\Big]\int_{x,x',z}\frac{(x-z)\cdot(x'-z)}{(x-z)^{2}(x'-z)^{2}}[U^{ce'}(z) - U^{ce'}(x')][U^{ce}(z) - U^{ce}(x)]\rho^{e'}(x')\rho^{e}(x)$
\item[\textbf{l.192}] \textsf{\small [missing term added]} \emph{added:} $\displaystyle \frac{d^{3}N_{\log(1-k^{+}/\vee)}}{d^{3}k}$
\item[\textbf{l.192}] \textsf{\small [phase convention]} \emph{added:} $\displaystyle = -\frac{1}{(2\pi)^{3}}\,\frac{g^{4}\,N_{c}}{8\pi^{4}}\,\frac{1}{k^{+}}\,\log\Big(1-\frac{k^{+}}{\vee}\Big)\int_{x,y,x',y'}e^{-ik\cdot(y-y')}\Big[2\gamma_{E}+\log\frac{(y'-y)^{2}\mu^{2}_{\overline{MS}}}{4}\Big]\frac{(x-y)\cdot(x'-y')}{(x-y)^{2}(x'-y')^{2}}$
\item[\textbf{l.192}] \textsf{\small [missing term added]} \emph{added:} $\displaystyle \times[U^{ce'}(y') - U^{ce'}(x')][U^{ce}(y) - U^{ce}(x)]\rho^{e'}(x')\rho^{e}(x)$
\item[\textbf{l.192}] \textsf{\small [missing term added]} \emph{added:} $\displaystyle \frac{d^{3}N_{\rm DGLAP}}{d^{3}k}$
\item[\textbf{l.192}] \textsf{\small [phase convention]} \emph{added:} $\displaystyle = -\frac{1}{(2\pi)^{3}}\frac{g^{2}}{2\pi^{2}}\frac{g^{2}}{8\pi^{2}}\frac{1}{k^{+}}\int_{\frac{k^{+}}{\vee}}^{1}d\xi\,\Bigg\{\frac{\mathcal P_{gg}(\xi)}{\xi^{2}}\int_{x,y,x',y'}e^{-ik\cdot(y-y')\frac{1}{\xi}}\Big[2\gamma_{E}+\log\frac{(y'-y)^{2}\mu^{2}_{\overline{MS}}}{4\,\xi^{2}}\Big]$
\item[\textbf{l.192}] \textsf{\small [missing term added]} \emph{added:} $\displaystyle \times\frac{(x-y)\cdot(x'-y')}{(x-y)^{2}(x'-y')^{2}}[U^{ce'}(y') - U^{ce'}(x')][U^{ce}(y) - U^{ce}(x)]\rho^{e'}(x')\rho^{e}(x)$
\item[\textbf{l.192}] \textsf{\small [missing term added]} \emph{added:} $\displaystyle +2N_{c}\,\frac{1-\xi}{\xi}\,\frac{4\pi}{k^{2}}\int_{x,x',z}\frac{(x-z)\cdot(x'-z)}{(x-z)^{2}(x'-z)^{2}}[U^{ce'}(z) - U^{ce'}(x')][U^{ce}(z) - U^{ce}(x)]\rho^{e'}(x')\rho^{e}(x)\Bigg\}$
\item[\textbf{l.192}] \textsf{\small [missing term added]} \emph{added:} $\displaystyle \frac{d^{3}N_{\rm RC}}{d^{3}k}$
\item[\textbf{l.192}] \textsf{\small [phase convention]} \emph{added:} $\displaystyle = +\frac{1}{(2\pi)^{3}}\frac{g^{2}}{2\pi^{2}}\frac{g^{2}\beta_{0}}{16\pi^{2}}\frac{1}{k^{+}}\int_{x,y,x',y'}e^{-ik\cdot(y-y')}\Big[2\gamma_{E}+\log\frac{(y'-y)^{2}\mu^{2}_{\overline{MS}}}{4}\Big]\frac{(x-y)\cdot(x'-y')}{(x-y)^{2}(x'-y')^{2}}$
\item[\textbf{l.192}] \textsf{\small [missing term added]} \emph{added:} $\displaystyle \times[U^{ce'}(y') - U^{ce'}(x')][U^{ce}(y) - U^{ce}(x)]\rho^{e'}(x')\rho^{e}(x)$
\end{itemize}
\needspace{3\baselineskip}\noindent\rule{\linewidth}{0.3pt}\par\noindent\textbf{Lines 196} \quad \textsf{Group III (with c.c.) $\rangle$ Row I} \quad {\small (in \texttt{align}, lines 195--199)}\par
\begin{itemize}[leftmargin=2.2em,itemsep=1pt,topsep=2pt]
\item[\textbf{l.196}] \textsf{\small [phase convention]} $\displaystyle w^{\prime}-w$\ $\longrightarrow$\ $\displaystyle w-w^{\prime}$
\end{itemize}
\needspace{3\baselineskip}\noindent\rule{\linewidth}{0.3pt}\par\noindent\textbf{Lines 202} \quad \textsf{Group III (with c.c.) $\rangle$ Row II} \quad {\small (in \texttt{align}, lines 201--205)}\par
\begin{itemize}[leftmargin=2.2em,itemsep=1pt,topsep=2pt]
\item[\textbf{l.202}] \textsf{\small [phase convention]} $\displaystyle w^{\prime}-w$\ $\longrightarrow$\ $\displaystyle w-w^{\prime}$
\end{itemize}
\needspace{3\baselineskip}\noindent\rule{\linewidth}{0.3pt}\par\noindent\textbf{Lines 209--212} \quad \textsf{Group III (with c.c.) $\rangle$ Row III} \quad {\small (in \texttt{align}, lines 208--212)}\par
\begin{itemize}[leftmargin=2.2em,itemsep=1pt,topsep=2pt]
\item[\textbf{l.209}] \textsf{\small [phase convention]} $\displaystyle y'-w$\ $\longrightarrow$\ $\displaystyle w-y'$
\item[\textbf{l.209}] \textsf{\small [phase convention]} $\displaystyle y'-z$\ $\longrightarrow$\ $\displaystyle z-y'$
\item[\textbf{l.210}] \textsf{\small [phase convention]} $\displaystyle y'-w$\ $\longrightarrow$\ $\displaystyle w-y'$
\item[\textbf{l.210}] \textsf{\small [phase convention]} $\displaystyle y'-z$\ $\longrightarrow$\ $\displaystyle z-y'$
\item[\textbf{l.211}] \textsf{\small [phase convention]} $\displaystyle y'-w$\ $\longrightarrow$\ $\displaystyle w-y'$
\end{itemize}
\begin{notebox}\textsf{\small Note, l.212:} [Kept in the author's original form (phase corrected only). The split of the $p^{+}$ integral with the kernel $\big[\frac{k^{+}}{p^{+}(p^{+}+k^{+})}\big]_{+}+\frac{1}{p^{+}+k^{+}}$ and the subtracted $y'\to z$ term is the correct way to take out the complete $\log\frac{\vee}{k^{+}}$ of this row (second and third terms of this row in the BFKL file): with the plain $\big[\frac{1}{p^{+}}\big]_{+}$ split of the checked Group III file the left-over $p^{+}$ integral up to $\vee-k^{+}$ would still depend logarithmically on $\vee$.]\end{notebox}
\needspace{3\baselineskip}\noindent\rule{\linewidth}{0.3pt}\par\noindent\textbf{Lines 215--223} \quad {\small (in \texttt{align}, lines 214--223)}\par
\begin{itemize}[leftmargin=2.2em,itemsep=1pt,topsep=2pt]
\item[\textbf{l.215}] \textsf{\small [phase convention]} $\displaystyle y'-w$\ $\longrightarrow$\ $\displaystyle w-y'$
\item[\textbf{l.215}] \textsf{\small [phase convention]} $\displaystyle y'-z$\ $\longrightarrow$\ $\displaystyle z-y'$
\item[\textbf{l.216}] \textsf{\small [phase convention]} $\displaystyle y'-w$\ $\longrightarrow$\ $\displaystyle w-y'$
\item[\textbf{l.216}] \textsf{\small [phase convention]} $\displaystyle y'-z$\ $\longrightarrow$\ $\displaystyle z-y'$
\item[\textbf{l.219}] \textsf{\small [phase convention]} $\displaystyle w'-w$\ $\longrightarrow$\ $\displaystyle w-w'$
\item[\textbf{l.221}] \textsf{\small [phase convention]} $\displaystyle y'-w$\ $\longrightarrow$\ $\displaystyle w-y'$
\end{itemize}
\begin{notebox}\textsf{\small Note, l.223:} [Kept in the author's original form (phase corrected only). The split of the $p^{+}$ integral with the kernel $\big[\frac{k^{+}}{p^{+}(p^{+}+k^{+})}\big]_{+}+\frac{1}{p^{+}+k^{+}}$ and the subtracted $y'\to z$ term is the correct way to take out the complete $\log\frac{\vee}{k^{+}}$ of this row (second and third terms of this row in the BFKL file): with the plain $\big[\frac{1}{p^{+}}\big]_{+}$ split of the checked Group III file the left-over $p^{+}$ integral up to $\vee-k^{+}$ would still depend logarithmically on $\vee$.]\end{notebox}
\needspace{3\baselineskip}\noindent\rule{\linewidth}{0.3pt}\par\noindent\textbf{Lines 228--231} \quad \textsf{Group III (with c.c.) $\rangle$ Row V} \quad {\small (in \texttt{align}, lines 227--231)}\par
\begin{itemize}[leftmargin=2.2em,itemsep=1pt,topsep=2pt]
\item[\textbf{l.228}] \textsf{\small [phase convention]} $\displaystyle y'-w$\ $\longrightarrow$\ $\displaystyle w-y'$
\item[\textbf{l.228}] \textsf{\small [phase convention]} $\displaystyle y'-z$\ $\longrightarrow$\ $\displaystyle z-y'$
\item[\textbf{l.229}] \textsf{\small [phase convention]} $\displaystyle y'-w$\ $\longrightarrow$\ $\displaystyle w-y'$
\item[\textbf{l.229}] \textsf{\small [phase convention]} $\displaystyle y'-z$\ $\longrightarrow$\ $\displaystyle z-y'$
\item[\textbf{l.230}] \textsf{\small [phase convention]} $\displaystyle y'-w$\ $\longrightarrow$\ $\displaystyle w-y'$
\end{itemize}
\begin{notebox}\textsf{\small Note, l.231:} [Kept in the author's original form (phase corrected only). The split of the $p^{+}$ integral with the kernel $\big[\frac{k^{+}}{p^{+}(p^{+}+k^{+})}\big]_{+}+\frac{1}{p^{+}+k^{+}}$ and the subtracted $y'\to z$ term is the correct way to take out the complete $\log\frac{\vee}{k^{+}}$ of this row (second and third terms of this row in the BFKL file): with the plain $\big[\frac{1}{p^{+}}\big]_{+}$ split of the checked Group III file the left-over $p^{+}$ integral up to $\vee-k^{+}$ would still depend logarithmically on $\vee$.]\end{notebox}
\needspace{3\baselineskip}\noindent\rule{\linewidth}{0.3pt}\par\noindent\textbf{Lines 234--236} \quad \textsf{Group III (with c.c.) $\rangle$ Row VI} \quad {\small (in \texttt{align}, lines 233--236)}\par
\begin{itemize}[leftmargin=2.2em,itemsep=1pt,topsep=2pt]
\item[\textbf{l.234}] \textsf{\small [phase convention]} $\displaystyle y'-w$\ $\longrightarrow$\ $\displaystyle w-y'$
\item[\textbf{l.234}] \textsf{\small [phase convention]} $\displaystyle y'-z$\ $\longrightarrow$\ $\displaystyle z-y'$
\item[\textbf{l.234}] \textsf{\small [phase convention]} $\displaystyle y'-w$\ $\longrightarrow$\ $\displaystyle w-y'$
\item[\textbf{l.234}] \textsf{\small [phase convention]} $\displaystyle y'-z$\ $\longrightarrow$\ $\displaystyle z-y'$
\item[\textbf{l.235}] \textsf{\small [phase convention]} $\displaystyle y'-w$\ $\longrightarrow$\ $\displaystyle w-y'$
\end{itemize}
\begin{notebox}\textsf{\small Note, l.236:} [Kept in the author's original form (phase corrected only). The split of the $p^{+}$ integral with the kernel $\big[\frac{k^{+}}{p^{+}(p^{+}+k^{+})}\big]_{+}+\frac{1}{p^{+}+k^{+}}$ and the subtracted $y'\to z$ term is the correct way to take out the complete $\log\frac{\vee}{k^{+}}$ of this row (second and third terms of this row in the BFKL file): with the plain $\big[\frac{1}{p^{+}}\big]_{+}$ split of the checked Group III file the left-over $p^{+}$ integral up to $\vee-k^{+}$ would still depend logarithmically on $\vee$.]\end{notebox}
\subsubsection*{Other unmarked edits}
\noindent{\small\textsf{l.42 (inserted):} \texttt{\%    \textbackslash{}\allowbreak{}frac\{d\^{}\{3\}\allowbreak{}N\_\{1\textbackslash{}\allowbreak{},\allowbreak{}\textbackslash{}\allowbreak{}rm NLO\}\allowbreak{}\^{}\{\textbackslash{}\allowbreak{}rm\}\allowbreak{}\}\allowbreak{}\{d\^{}\{3\}\allowbreak{}k\}\allowbreak{} \&  =\allowbreak{}  \textbackslash{}\allowbreak{}frac\{1\}\allowbreak{}\{(2\textbackslash{}\allowbreak{}pi)\allowbreak{}\^{}\{3\}\allowbreak{}\}\allowbreak{}\textbackslash{}\allowbreak{},\allowbreak{} \textbackslash{}\allowbreak{}frac\{g\^{}\{4\}\allowbreak{}\textbackslash{}\allowbreak{},\allowbreak{}N\_\{c\}\allowbreak{}\}\allowbreak{}\{32\textbackslash{}\allowbreak{}pi\^{}\{4\}\allowbreak{}\}\allowbreak{}\textbackslash{}\allowbreak{},\allowbreak{}\textbackslash{}\allowbreak{}frac\{1\}\allowbreak{}\{k\^{}\{+\allowbreak{}\}\allowbreak{}\}\allowbreak{}\textbackslash{}\allowbreak{},\allowbreak{} \textbackslash{}\allowbreak{}int\_\{w,\allowbreak{} w',\allowbreak{}x,\allowbreak{}x',\allowbreak{}\}\allowbreak{} e\^{}\{-ik(\{\textbackslash{}\allowbreak{}color\{orange\}\allowbreak{}\textbackslash{}\allowbreak{}cancel\{w\^{}\textbackslash{}\allowbreak{}prime-w\}\allowbreak{}\}\allowbreak{}\{\textbackslash{}\allowbreak{}color\{red\}\allowbreak{}w-w\^{}\textbackslash{}\allowbreak{}prime\}\allowbreak{})\allowbreak{}\}\allowbreak{}\textbackslash{}\allowbreak{},\allowbreak{} \textbackslash{}\allowbreak{}frac\{(x'-w')\allowbreak{}\textbackslash{}\allowbreak{}cdot (x-w)\allowbreak{}\}\allowbreak{}\{(x'-w')\allowbreak{}\^{}\{2\}\allowbreak{}\textbackslash{}\allowbreak{},\allowbreak{} (x-w)\allowbreak{}\^{}\{2\}\allowbreak{}\}\allowbreak{}\textbackslash{}\allowbreak{},\allowbreak{} \textbackslash{}\allowbreak{}Bigg\textbackslash{}\allowbreak{}\{\{\textbackslash{}\allowbreak{}color\{orange\}\allowbreak{}\textbackslash{}\allowbreak{}left[-\textbackslash{}\allowbreak{}frac\{2\}\allowbreak{}\{\textbackslash{}\allowbreak{}epsilon\}\allowbreak{}-2\textbackslash{}\allowbreak{}gamma-\textbackslash{}\allowbreak{}log\{\textbackslash{}\allowbreak{}left(\textbackslash{}\allowbreak{}frac\{(x-w)\allowbreak{}\^{}\{2\}\allowbreak{}\textbackslash{}\allowbreak{},\allowbreak{}\textbackslash{}\allowbreak{}mu\^{}\{2\}\allowbreak{}\_\{\textbackslash{}\allowbreak{}overline\{MS\}\allowbreak{}\}\allowbreak{}\}\allowbreak{}\{4\}\allowbreak{}\textbackslash{}\allowbreak{}right)\allowbreak{}\}\allowbreak{}\textbackslash{}\allowbreak{}right]\}\allowbreak{}\textbackslash{}\allowbreak{}times}}\par
\noindent{\small\textsf{l.42 (deleted):} \texttt{\%    \textbackslash{}\allowbreak{}frac\{d\^{}\{3\}\allowbreak{}N\_\{1\textbackslash{}\allowbreak{},\allowbreak{}\textbackslash{}\allowbreak{}rm NLO\}\allowbreak{}\^{}\{\textbackslash{}\allowbreak{}rm\}\allowbreak{}\}\allowbreak{}\{d\^{}\{3\}\allowbreak{}k\}\allowbreak{} \&  =\allowbreak{}  \textbackslash{}\allowbreak{}frac\{1\}\allowbreak{}\{(2\textbackslash{}\allowbreak{}pi)\allowbreak{}\^{}\{3\}\allowbreak{}\}\allowbreak{}\textbackslash{}\allowbreak{},\allowbreak{} \textbackslash{}\allowbreak{}frac\{g\^{}\{4\}\allowbreak{}\textbackslash{}\allowbreak{},\allowbreak{}N\_\{c\}\allowbreak{}\}\allowbreak{}\{32\textbackslash{}\allowbreak{}pi\^{}\{4\}\allowbreak{}\}\allowbreak{}\textbackslash{}\allowbreak{},\allowbreak{}\textbackslash{}\allowbreak{}frac\{1\}\allowbreak{}\{k\^{}\{+\allowbreak{}\}\allowbreak{}\}\allowbreak{}\textbackslash{}\allowbreak{},\allowbreak{} \textbackslash{}\allowbreak{}int\_\{w,\allowbreak{} w',\allowbreak{}x,\allowbreak{}x',\allowbreak{}\}\allowbreak{} e\^{}\{-ik(w\^{}\textbackslash{}\allowbreak{}prime-w)\allowbreak{}\}\allowbreak{}\textbackslash{}\allowbreak{},\allowbreak{} \textbackslash{}\allowbreak{}frac\{(x'-w')\allowbreak{}\textbackslash{}\allowbreak{}cdot (x-w)\allowbreak{}\}\allowbreak{}\{(x'-w')\allowbreak{}\^{}\{2\}\allowbreak{}\textbackslash{}\allowbreak{},\allowbreak{} (x-w)\allowbreak{}\^{}\{2\}\allowbreak{}\}\allowbreak{}\textbackslash{}\allowbreak{},\allowbreak{} \textbackslash{}\allowbreak{}Bigg\textbackslash{}\allowbreak{}\{\{\textbackslash{}\allowbreak{}color\{orange\}\allowbreak{}\textbackslash{}\allowbreak{}left[-\textbackslash{}\allowbreak{}frac\{2\}\allowbreak{}\{\textbackslash{}\allowbreak{}epsilon\}\allowbreak{}-2\textbackslash{}\allowbreak{}gamma-\textbackslash{}\allowbreak{}log\{\textbackslash{}\allowbreak{}left(\textbackslash{}\allowbreak{}frac\{(x-w)\allowbreak{}\^{}\{2\}\allowbreak{}\textbackslash{}\allowbreak{},\allowbreak{}\textbackslash{}\allowbreak{}mu\^{}\{2\}\allowbreak{}\_\{\textbackslash{}\allowbreak{}overline\{MS\}\allowbreak{}\}\allowbreak{}\}\allowbreak{}\{4\}\allowbreak{}\textbackslash{}\allowbreak{}right)\allowbreak{}\}\allowbreak{}\textbackslash{}\allowbreak{}right]\}\allowbreak{}\textbackslash{}\allowbreak{}times \textbackslash{}\allowbreak{}notag \textbackslash{}\allowbreak{}\textbackslash{}\allowbreak{} \& \{\textbackslash{}\allowbreak{}color\{orange\}\allowbreak{}\textbackslash{}\allowbreak{}left[\textbackslash{}\allowbreak{}frac\{11\}\allowbreak{}\{3\}\allowbreak{}+\allowbreak{}\textbackslash{}\allowbreak{}}}\par
\needspace{15\baselineskip}\bigskip
\subsection{Other terms (v2)}\label{sec:file-other_terms_v2}
\noindent\textbf{Original:} \texttt{/root/.claude/uploads/0ab97fcf-94a2-5aec-8bad-586400481685/ff9ece59-Other\_terms.tex}\\\textbf{Checked copy:} \texttt{/home/user/correlations/nlo\_review/other\_terms\_v2/Other\_terms\_v2\_checked.tex}\\
Lines: 239 (original), 239 (checked copy). Coloured groups of the check: 47 (47 top-level); author\textquotesingle s own coloured groups left unchanged: 0.\\
Records by type: phase convention 19; typo/LaTeX 5; note only 23.
\par\noindent Unmarked layout-only edits (line breaks \texttt{\textbackslash notag\textbackslash\textbackslash\ \&}, \texttt{\textbackslash par\textbackslash noindent} around inserted notes, spacing), checked-copy lines: 35, 43, 49, 109, 112, 119, 124, 128, 140, 142, 154, 155, 163, 164, 182, 194, 201, 207, 214, 220, 225, 233, 238.
\par\noindent Unmarked edits that are not layout only (listed with their block below): 19.
\subsubsection*{Changes, block by block}
\needspace{3\baselineskip}\noindent\rule{\linewidth}{0.3pt}\par\noindent\textbf{Lines 35}\par
\begin{notebox}\textsf{\small Note, l.35:} \noindent[\textbf{Checked copy.} Marking: red $=$ corrected or added, part of the expression; red text in square brackets $=$ notes; orange, struck out $=$ removed. Checked, row by row: [this file] $=$ [checked row with the cut-offs $\Lambda$, $\vee$] $-\log\frac{\vee}{k^{+}}\times$[terms of that row in the BFKL file] $-\log\frac{k^{+}}{\Lambda}\times$[terms of that row in the JIMWLK, two-$\rho$, three-$\rho$ files], up to terms that vanish for $\Lambda\to0$. The comparison is with the \emph{checked} (corrected) Group I, Group II, Group II Row IV and Group III files, textually and numerically (numbers in the notes; ``residual'' is the largest component of [row] $-\log\frac{\vee}{k^{+}}$[BFKL] $-\log\frac{k^{+}}{\Lambda}$[JIMWLK] $-$ [this file] at random points with explicit SU(3) matrices, for $\Lambda/k^{+}=10^{-3},10^{-5},10^{-7}$). Result: all rows agree; open point at Group I Row III ($N_{3b}$). The only corrections made in this file: the closing brace of the phases was misplaced (the red parenthesis now closes inside the exponent; in Group I Row V the factors $\frac{p^{+}}{k^{+}}$, $\frac{q^{+}}{k^{+}}$ had fallen out of the exponent) and commas in two lists of integration variables; four displayed lines that reached or passed the edge of the page are broken (Group II Row I, $1b$ and $2b$; Group III Row III, second part), without change of content. Conventions as decided: phase $e^{-ik\cdot(w-w')}$ with the primed (conjugated) bracket at $w'$; $\frac{2}{\epsilon}$ poles are not written; $\int dp^{+}\,[h(p^{+})]_{+}\,f(p^{+})\equiv\int dp^{+}\,h(p^{+})\,[f(p^{+})-f(0)]$, where $f$ is everything else that depends on $p^{+}$.]\end{notebox}
\needspace{3\baselineskip}\noindent\rule{\linewidth}{0.3pt}\par\noindent\textbf{Lines 38} \quad \textsf{Group I (with c.c.) $\rangle$ Row I}\par
\begin{notebox}\textsf{\small Note, l.38:} [Not re-checked in this pass. Earlier check: the row carries $\log\frac{\vee}{\Lambda}$ only and nothing is left over.]\end{notebox}
\needspace{3\baselineskip}\noindent\rule{\linewidth}{0.3pt}\par\noindent\textbf{Lines 42--43} \quad \textsf{Group I (with c.c.) $\rangle$ Row II} \quad {\small (in \texttt{align}, lines 41--43)}\par
\begin{itemize}[leftmargin=2.2em,itemsep=1pt,topsep=2pt]
\item[\textbf{l.42}] \textsf{\small [typo/LaTeX]} \emph{removed:} $\displaystyle ,$
\item[\textbf{l.42}] \textsf{\small [phase convention]} \emph{added:} $\displaystyle )$\quad{\color{gray}\small(in context: }$\displaystyle w^\prime)${\color{gray}\small)}
\end{itemize}
\begin{notebox}\textsf{\small Note, l.43:} [Checked: this is $N_{1}$ of the checked Row II without $-2\log\frac{k^{+}}{\Lambda}$ in the second bracket (Row II term of the two-$\rho$ list) and without the $-\frac{2}{\epsilon}$ of the first bracket; the two forms of the second bracket of the checked file agree numerically. The row has no term linear in $\log\frac{\vee}{k^{+}}$ except the product $\log(1-\frac{k^{+}}{\vee})\log\frac{\vee}{k^{+}}$, kept here by decision. The $\frac{11}{3}$ term cancels against $N_{2b\,\rm UV}$ of Row V (same structure after $w\to x$, $x\to y$, coefficient $-\frac{11}{6}\cdot\frac{g^{4}N_{c}}{16\pi^{4}}$). Removed: a trailing comma in the list of integration variables]\end{notebox}
\noindent{\small\textsf{Unmarked edit, l.42 (deleted):} \texttt{)\allowbreak{}}}\par
\needspace{3\baselineskip}\noindent\rule{\linewidth}{0.3pt}\par\noindent\textbf{Lines 45--49} \quad {\small (in \texttt{align}, lines 44--49)}\par
\begin{itemize}[leftmargin=2.2em,itemsep=1pt,topsep=2pt]
\item[\textbf{l.45}] \textsf{\small [phase convention]} \emph{added:} $\displaystyle )$\quad{\color{gray}\small(in context: }$\displaystyle w')${\color{gray}\small)}
\end{itemize}
\begin{notebox}\textsf{\small Note, l.49:} [Checked: literal first term of $N_{2b}$ of the checked Row II ($q$ here is the integration variable called $k$ there). Numerically, at random $p$, $q$ and three sets $(k^{+},\Lambda,\vee)$: the two-$\rho$ bracket of the checked $N_{2}$ equals $C_{J}\log\frac{k^{+}}{\Lambda}+R^{i}$ for the colour structure with $U^{bd}(x)U^{ce}(y)$, and $(C_{J}+2q^{i})\log\frac{k^{+}}{\Lambda}+2q^{i}\log\frac{\vee}{k^{+}}+R^{i}$ for the one with $U^{ah}(w)$, with $R^{i}$ as below (residuals below $10^{-14}$; $C_{J}$ is the coefficient of $\log\frac{k^{+}}{\Lambda}$). The term $2q^{i}\log\frac{\vee}{k^{+}}$ is the Row II term of the BFKL file. It comes from the operator-ordering correction of the Group I file; in the uncorrected Group I file it is absent, and the BFKL Row II term and the form of the $\log\frac{k^{+}}{\Lambda}$ term in the three-$\rho$ list then have no counterpart there]\end{notebox}
\noindent{\small\textsf{Unmarked edit, l.45 (deleted):} \texttt{)\allowbreak{}}}\par
\needspace{3\baselineskip}\noindent\rule{\linewidth}{0.3pt}\par\noindent\textbf{Lines 60} \quad \textsf{Group I (with c.c.) $\rangle$ Row III} \quad {\small (in \texttt{align}, lines 59--65)}\par
\begin{itemize}[leftmargin=2.2em,itemsep=1pt,topsep=2pt]
\item[\textbf{l.60}] \textsf{\small [phase convention]} \emph{added:} $\displaystyle )$\quad{\color{gray}\small(in context: }$\displaystyle w^\prime)${\color{gray}\small)}
\end{itemize}
\noindent{\small\textsf{Unmarked edit, l.60 (deleted):} \texttt{)\allowbreak{}}}\par
\needspace{3\baselineskip}\noindent\rule{\linewidth}{0.3pt}\par\noindent\textbf{Lines 67} \quad {\small (in \texttt{align}, lines 66--74)}\par
\begin{itemize}[leftmargin=2.2em,itemsep=1pt,topsep=2pt]
\item[\textbf{l.67}] \textsf{\small [phase convention]} \emph{added:} $\displaystyle )$\quad{\color{gray}\small(in context: }$\displaystyle w^\prime)${\color{gray}\small)}
\end{itemize}
\noindent{\small\textsf{Unmarked edit, l.67 (deleted):} \texttt{)\allowbreak{}}}\par
\needspace{3\baselineskip}\noindent\rule{\linewidth}{0.3pt}\par\noindent\textbf{Lines 76} \quad {\small (in \texttt{align}, lines 75--98)}\par
\begin{itemize}[leftmargin=2.2em,itemsep=1pt,topsep=2pt]
\item[\textbf{l.76}] \textsf{\small [phase convention]} \emph{added:} $\displaystyle )$\quad{\color{gray}\small(in context: }$\displaystyle w')${\color{gray}\small)}
\end{itemize}
\noindent{\small\textsf{Unmarked edit, l.76 (deleted):} \texttt{)\allowbreak{}}}\par
\needspace{3\baselineskip}\noindent\rule{\linewidth}{0.3pt}\par\noindent\textbf{Lines 100--109} \quad {\small (in \texttt{align}, lines 99--109)}\par
\begin{itemize}[leftmargin=2.2em,itemsep=1pt,topsep=2pt]
\item[\textbf{l.100}] \textsf{\small [phase convention]} \emph{added:} $\displaystyle )$\quad{\color{gray}\small(in context: }$\displaystyle w')${\color{gray}\small)}
\end{itemize}
\begin{notebox}\textsf{\small Note, l.109:} [Checked, Row III ($N_{1}$, $N_{2}$, $N_{3a}$, $N_{3b}$): every term is a term of the checked Row III with the explicit $\log\frac{\vee}{k^{+}}$ and $\log\frac{k^{+}}{\Lambda}$ removed and $\Lambda\to0$ in the remaining logarithms; the last term of $N_{3b}$ is the term of the checked row itself, with $\Lambda$ kept. Dropped for $\Lambda\to0$: the terms $\log(1\pm\frac{\Lambda}{k^{+}})$ and $\log\frac{k^{+}+\Lambda}{k^{+}-\Lambda}$ times kernels, the $\Lambda$ inside the transverse logarithms and in $\mathcal{A}\to\mathcal{A}_{0}$. The $\frac{1}{\Lambda}$ terms cancel in the sum $N_{1}+N_{2}+N_{3a}$, and two $\log(1-\frac{k^{+}}{\vee})$ terms of the checked $N_{2}$ with equal kernels cancel. Numerically, for the sum of the four parts (terms of size $10^{-7}$ to $10^{-6}$): residual $4\cdot10^{-10}$, $4\cdot10^{-12}$, $4\cdot10^{-14}$, \emph{if} the $\log\frac{k^{+}}{\Lambda}$ term subtracted for the kernel $1/(z-w)^{2}$ is the one of the checked row, $-\frac{ig^{4}f^{abd}}{8\pi^{5}}\frac{1}{k^{+}}\log\frac{k^{+}}{\Lambda}\frac{(x'-w')\cdot(y-z)}{(x'-w')^{2}(y-z)^{2}(z-w)^{2}}$ with $U^{bc}(z)$. \textbf{Open point:} the three-$\rho$ list (Evolution\_three\_rho, Row III, third term and last equation) has this term in another form, $-\frac{ig^{4}f^{abd}}{8\pi^{5}}\frac{1}{k^{+}}\log\frac{k^{+}}{\Lambda}\frac{(x'-w')\cdot(y-w)\,(x-z)\cdot(w-z)}{(x'-w')^{2}(y-w)^{2}(x-z)^{2}(w-z)^{2}}$ with $U^{bc}(w)$, which is not equal to it at fixed coordinates: with the list as written the residual is $4\cdot10^{-7}$ to $9\cdot10^{-7}$ and does not decrease with $\Lambda$. The last term of $N_{3b}$ here is the partner of the first form. Not decided here]\end{notebox}
\noindent{\small\textsf{Unmarked edit, l.100 (deleted):} \texttt{)\allowbreak{}}}\par
\needspace{3\baselineskip}\noindent\rule{\linewidth}{0.3pt}\par\noindent\textbf{Lines 112}\par
\begin{notebox}\textsf{\small Note, l.112:} [Checked numerically: the two forms of $\mathcal{A}_{0}$ agree (12 digits), and $\mathcal{A}$ of the checked file tends to $\mathcal{A}_{0}$ linearly in $\Lambda/k^{+}$]\end{notebox}
\needspace{3\baselineskip}\noindent\rule{\linewidth}{0.3pt}\par\noindent\textbf{Lines 116--119} \quad \textsf{Group I (with c.c.) $\rangle$ Row IV} \quad {\small (in \texttt{align}, lines 114--119)}\par
\begin{itemize}[leftmargin=2.2em,itemsep=1pt,topsep=2pt]
\item[\textbf{l.116}] \textsf{\small [phase convention]} \emph{added:} $\displaystyle )$\quad{\color{gray}\small(in context: }$\displaystyle w^\prime)${\color{gray}\small)}
\end{itemize}
\begin{notebox}\textsf{\small Note, l.119:} [Checked: sub-expression of the checked Row IV. Removed with the evolution logarithms: $\log\frac{\vee}{k^{+}}$ times the first and third kernels and times $-\frac{ig^{4}}{16\pi^{5}k^{+}}\frac{(x'-w')\cdot(y-z)}{(x'-w')^{2}(y-z)^{2}(z-w)^{2}}$, and $\log\frac{k^{+}}{\Lambda}$ times the third kernel (BFKL file Row IV, three-$\rho$ list Row IV). Dropped for $\Lambda\to0$: $\frac{\Lambda}{k^{+}}(x-z)^{2}$ in the denominator of the first logarithm and the term with $\log(1+\frac{\Lambda}{k^{+}})$. Numerically: residual $2\cdot10^{-11}$, $2\cdot10^{-13}$, $2\cdot10^{-15}$ (row $2\cdot10^{-7}$ to $5\cdot10^{-7}$), also for $\vee/k^{+}=40$]\end{notebox}
\noindent{\small\textsf{Unmarked edit, l.116 (deleted):} \texttt{)\allowbreak{}}}\par
\needspace{3\baselineskip}\noindent\rule{\linewidth}{0.3pt}\par\noindent\textbf{Lines 123--124} \quad \textsf{Group I (with c.c.) $\rangle$ Row V} \quad {\small (in \texttt{align}, lines 122--124)}\par
\begin{itemize}[leftmargin=2.2em,itemsep=1pt,topsep=2pt]
\item[\textbf{l.123}] \textsf{\small [typo/LaTeX]} \emph{added:} $\displaystyle ,$\quad{\color{gray}\small(in context: }$\displaystyle \, \frac{dp^{+}}{2\pi}\,\int_{x,z,y,x',y',w'}\,e^{-ik\cdot (z-w^\prime)}\,e^{-ik\cdot (x-z)\frac{p^{+}}{k^{+}}}${\color{gray}\small)}
\item[\textbf{l.123}] \textsf{\small [phase convention]} \emph{added:} $\displaystyle )$\quad{\color{gray}\small(in context: }$\displaystyle w^\prime)${\color{gray}\small)}
\item[\textbf{l.123}] \textsf{\small [typo/LaTeX]} {\small\ttfamily \}\allowbreak{})\allowbreak{}\textbackslash{}\allowbreak{}frac\{p\^{}\{+\allowbreak{}\}\allowbreak{}\}\allowbreak{}\{k\^{}\{+\allowbreak{}\}\allowbreak{}}\ $\longrightarrow$\ $\displaystyle x-z)\frac{p^{+}}{k^{+}}$
\item[\textbf{l.123}] \textsf{\small [typo/LaTeX]} \emph{added:} $\displaystyle ,$\quad{\color{gray}\small(in context: }$\displaystyle \, \frac{dq^{+}}{2\pi}\,\int_{x,z,y,x',y',w'}\,e^{-ik\cdot (x-w^\prime)}\,e^{-ik\cdot (z-x)\frac{q^{+}}{k^{+}}}${\color{gray}\small)}
\item[\textbf{l.123}] \textsf{\small [phase convention]} \emph{added:} $\displaystyle )$\quad{\color{gray}\small(in context: }$\displaystyle w^\prime)${\color{gray}\small)}
\item[\textbf{l.123}] \textsf{\small [typo/LaTeX]} {\small\ttfamily \}\allowbreak{})\allowbreak{}\textbackslash{}\allowbreak{}frac\{q\^{}\{+\allowbreak{}\}\allowbreak{}\}\allowbreak{}\{k\^{}\{+\allowbreak{}\}\allowbreak{}}\ $\longrightarrow$\ $\displaystyle z-x)\frac{q^{+}}{k^{+}}$
\end{itemize}
\begin{notebox}\textsf{\small Note, l.124:} [Checked: sub-expression of the checked Row V, first part. Removed: its two terms $-\frac{ig^{4}}{32\pi^{5}k^{+}}\big[\log\frac{k^{+}}{\Lambda}+\log(1-\frac{\Lambda}{k^{+}})\big]\dots$ (three-$\rho$ list, Row V; $\log(1-\frac{\Lambda}{k^{+}})$ dropped). The row has no $\log\frac{\vee}{k^{+}}$. Corrected: the braces of the four phases ($\frac{p^{+}}{k^{+}}$, $\frac{q^{+}}{k^{+}}$ belong to the exponents, as in the checked file) and the missing commas]\end{notebox}
\noindent{\small\textsf{Unmarked edit, l.123 (deleted):} \texttt{)\allowbreak{}}}\par
\noindent{\small\textsf{Unmarked edit, l.123 (deleted):} \texttt{)\allowbreak{}}}\par
\needspace{3\baselineskip}\noindent\rule{\linewidth}{0.3pt}\par\noindent\textbf{Lines 127--128} \quad {\small (in \texttt{align}, lines 126--128)}\par
\begin{itemize}[leftmargin=2.2em,itemsep=1pt,topsep=2pt]
\item[\textbf{l.127}] \textsf{\small [phase convention]} \emph{added:} $\displaystyle )$\quad{\color{gray}\small(in context: }$\displaystyle w')${\color{gray}\small)}
\item[\textbf{l.127}] \textsf{\small [phase convention]} \emph{added:} $\displaystyle )$\quad{\color{gray}\small(in context: }$\displaystyle w')${\color{gray}\small)}
\end{itemize}
\begin{notebox}\textsf{\small Note, l.128:} [Checked: sub-expression of the checked $N_{2a}$. Removed: its two terms with $-\frac{g^{4}}{16\pi^{5}k^{+}}\big[\log\frac{k^{+}}{\Lambda}+\log(1-\frac{\Lambda}{k^{+}})\big]$ (two-$\rho$ list, Row V, $2a$). As in the checked file $\int^{0}_{k^{+}}dq^{+}\big[\frac{k^{+}}{q^{+}}\big]_{+}f$ stands for $-\int_{0}^{k^{+}}dq^{+}\frac{k^{+}}{q^{+}}(f(q^{+})-f(0))$, and $d_{\perp}=2$]\end{notebox}
\noindent{\small\textsf{Unmarked edit, l.127 (deleted):} \texttt{)\allowbreak{}}}\par
\noindent{\small\textsf{Unmarked edit, l.127 (deleted):} \texttt{)\allowbreak{}}}\par
\needspace{3\baselineskip}\noindent\rule{\linewidth}{0.3pt}\par\noindent\textbf{Lines 130--142} \quad {\small (in \texttt{align}, lines 129--140)}\par
\begin{itemize}[leftmargin=2.2em,itemsep=1pt,topsep=2pt]
\item[\textbf{l.130}] \textsf{\small [phase convention]} \emph{added:} $\displaystyle )$\quad{\color{gray}\small(in context: }$\displaystyle w')${\color{gray}\small)}
\item[\textbf{l.138}] \textsf{\small [phase convention]} \emph{added:} $\displaystyle )$\quad{\color{gray}\small(in context: }$\displaystyle w')${\color{gray}\small)}
\end{itemize}
\begin{notebox}\textsf{\small Note, l.140:} [Checked: sub-expression of the checked $N_{2b\,\rm rest}$; removed: its three $\log\frac{k^{+}}{\Lambda}$ terms (two-$\rho$ list, Row V, $2b$ finite)]\end{notebox}
\begin{notebox}\textsf{\small Note, l.142:} [Checked: $N_{2b\,\rm UV}$ of the checked Row V, $\big[2\log\frac{k^{+}}{\Lambda}-\frac{11}{6}\big]\big[\frac{2}{\epsilon}+2\gamma_{E}+\log\dots\big]$, without the $\log\frac{k^{+}}{\Lambda}$ term (two-$\rho$ list, Row V, $2b$ UV) and without $\frac{2}{\epsilon}$. Notation: $\bar\mu^{2}_{\overline{MS}}$ here, $\mu^{2}_{\overline{MS}}$ in the other rows]\end{notebox}
\noindent{\small\textsf{Unmarked edit, l.130 (deleted):} \texttt{)\allowbreak{}}}\par
\noindent{\small\textsf{Unmarked edit, l.138 (deleted):} \texttt{)\allowbreak{}}}\par
\needspace{3\baselineskip}\noindent\rule{\linewidth}{0.3pt}\par\noindent\textbf{Lines 155} \quad \textsf{Group II $\rangle$ Row I}\par
\begin{notebox}\textsf{\small Note, l.155:} [Checked, first part ($1a+1b$): equal to the first part of the checked Group II Row I minus $\log\frac{\vee}{k^{+}}\times$(BFKL file, Row I, first part) minus $\log\frac{k^{+}}{\Lambda}\times$(three-$\rho$ list, Group II Row I); dropped: $\frac{\Lambda}{k^{+}}(x-z)^{2}$ in the logarithms and $\log(1+\frac{\Lambda}{k^{+}})$. Numerically: residual $1\cdot10^{-11}$, $1\cdot10^{-13}$, $1\cdot10^{-15}$ (row $1\cdot10^{-7}$ to $3\cdot10^{-7}$)]\end{notebox}
\needspace{3\baselineskip}\noindent\rule{\linewidth}{0.3pt}\par\noindent\textbf{Lines 182}\par
\begin{notebox}\textsf{\small Note, l.182:} [Checked, second part ($2a$ to $2f$, compared as a sum): equal to the sum of $2a$ to $2f$ of the checked Group II Row I minus $\log\frac{\vee}{k^{+}}\times$(BFKL file, Row I, second part, $2b$ to $2f$) minus $\log\frac{k^{+}}{\Lambda}\times$(two-$\rho$ list, Group II Row I). Numerically: residual $3\cdot10^{-11}$, $3\cdot10^{-13}$, $3\cdot10^{-15}$ (row $5\cdot10^{-8}$ to $1\cdot10^{-7}$), also for $\vee/k^{+}=40$]\end{notebox}
\needspace{3\baselineskip}\noindent\rule{\linewidth}{0.3pt}\par\noindent\textbf{Lines 184} \quad \textsf{Group II $\rangle$ Row II}\par
\begin{notebox}\textsf{\small Note, l.184:} [Not re-checked in this pass (four-$\rho$ only in the checked file).]\end{notebox}
\needspace{3\baselineskip}\noindent\rule{\linewidth}{0.3pt}\par\noindent\textbf{Lines 186} \quad \textsf{Group II $\rangle$ Row III}\par
\begin{notebox}\textsf{\small Note, l.186:} [Not re-checked in this pass (four-$\rho$ only in the checked file).]\end{notebox}
\needspace{3\baselineskip}\noindent\rule{\linewidth}{0.3pt}\par\noindent\textbf{Lines 194} \quad \textsf{Group II $\rangle$ Row IV}\par
\begin{notebox}\textsf{\small Note, l.194:} [Checked, Row IV (five equations): the $\xi$ integral called $N_{\rm Finite}$, $N_{\rm DGLAP}$ and $N_{\rm RC}$ are literal copies of the final list of the checked Row IV file, with the lower limit $\frac{k^{+}}{\vee}$ and without $\frac{2}{\epsilon}$. The second equation is $-(1-\frac{k^{+}}{\vee})\times$ the BFKL structure of that file and the third is $\log(1-\frac{k^{+}}{\vee})\times$ its JIMWLK structure (in the forms with $\frac{1}{k^{2}}$ and with $[2\gamma_{E}+\log\dots]$); both are left over by decision, the terms with $\log\frac{\vee}{k^{+}}$ and $\log\frac{k^{+}}{\Lambda}$ being those of the BFKL file and of the two-$\rho$ list. Not defined in this file (taken from the checked file): $\mathcal P_{gg}(\xi)=2N_{c}\big[\xi(1-\xi)+\frac{1-\xi}{\xi}+\frac{\xi}{(1-\xi)_{+}}\big]+\frac{\beta_{0}}{2}\delta(1-\xi)$, with $\int_{k^{+}/\vee}^{1}d\xi\,\frac{\xi}{(1-\xi)_{+}}f(\xi)\equiv\int_{k^{+}/\vee}^{1}d\xi\,\frac{\xi f(\xi)-f(1)}{1-\xi}$. Three different terms carry the name $N_{\rm Finite}$]\end{notebox}
\needspace{3\baselineskip}\noindent\rule{\linewidth}{0.3pt}\par\noindent\textbf{Lines 201} \quad \textsf{Group III (with c.c.) $\rangle$ Row I}\par
\begin{notebox}\textsf{\small Note, l.201:} [Checked: sub-expression of the checked Group III Row I (BFKL file Row I and three-$\rho$ list Row I removed; $\frac{\Lambda}{k^{+}}(x-z)^{2}$ and $\log(1+\frac{\Lambda}{k^{+}})$ dropped). Numerically: residual $2\cdot10^{-11}$, $2\cdot10^{-13}$, $2\cdot10^{-15}$ (row $2\cdot10^{-7}$ to $5\cdot10^{-7}$)]\end{notebox}
\needspace{3\baselineskip}\noindent\rule{\linewidth}{0.3pt}\par\noindent\textbf{Lines 207} \quad \textsf{Group III (with c.c.) $\rangle$ Row II}\par
\begin{notebox}\textsf{\small Note, l.207:} [Checked: as Row I, with the order of the two $\rho$'s of the first bracket as in the checked Row II. Numerically: residual $2\cdot10^{-11}$, $2\cdot10^{-13}$, $2\cdot10^{-15}$ (row $2\cdot10^{-7}$ to $6\cdot10^{-7}$)]\end{notebox}
\needspace{3\baselineskip}\noindent\rule{\linewidth}{0.3pt}\par\noindent\textbf{Lines 211--214} \quad \textsf{Group III (with c.c.) $\rangle$ Row III} \quad {\small (in \texttt{align}, lines 210--214)}\par
\begin{itemize}[leftmargin=2.2em,itemsep=1pt,topsep=2pt]
\item[\textbf{l.211}] \textsf{\small [phase convention]} \emph{added:} $\displaystyle )$\quad{\color{gray}\small(in context: }$\displaystyle y')${\color{gray}\small)}
\item[\textbf{l.211}] \textsf{\small [phase convention]} \emph{added:} $\displaystyle )$\quad{\color{gray}\small(in context: }$\displaystyle y')${\color{gray}\small)}
\item[\textbf{l.212}] \textsf{\small [phase convention]} \emph{added:} $\displaystyle )$\quad{\color{gray}\small(in context: }$\displaystyle y')${\color{gray}\small)}
\item[\textbf{l.212}] \textsf{\small [phase convention]} \emph{added:} $\displaystyle )$\quad{\color{gray}\small(in context: }$\displaystyle y')${\color{gray}\small)}
\item[\textbf{l.213}] \textsf{\small [phase convention]} \emph{added:} $\displaystyle )$\quad{\color{gray}\small(in context: }$\displaystyle y')${\color{gray}\small)}
\end{itemize}
\begin{notebox}\textsf{\small Note, l.214:} [Checked: \textbf{this form matches.} [Cut-off integral of the checked Group III file over $\Lambda<p^{+}<\vee-k^{+}$] $-\log\frac{k^{+}}{\Lambda}\times$[three-$\rho$ list, Row III] $-\log\frac{\vee}{k^{+}}\times$[\emph{all three} terms of this row in the BFKL file] $=$ this equation $+\mathcal{O}(\Lambda/k^{+})$. With the common prefactor set to 1, the first colour bracket replaced by a smooth test function of $y'$ (its value at $z$ where the bracket is taken at $y'\to z$), the second bracket by 1, and the $p^{+}$ and $y'$ integrals done by quadrature, for $\vee/k^{+}=6$: left-hand side $0.0044667+0.0226868\,i$, $0.0051661+0.0230997\,i$, $0.0051731+0.0231038\,i$ for $\Lambda/k^{+}=10^{-2},10^{-4},10^{-6}$; this equation: $0.0051732+0.0231038\,i$ (the same agreement for $\vee/k^{+}=2,15$ and for a second set of points). If only the explicit $\log\frac{\vee}{k^{+}}$ of the Group III file (first BFKL term) is subtracted, the two sides differ by $\log\frac{\vee}{k^{+}}\times$(second $+$ third BFKL term), $0.16$ in the same units: this is the mismatch seen in a direct comparison with the Group III file. Reason: there $-\frac{1}{p^{+}}\to-\big[\frac{1}{p^{+}}\big]_{+}$ with the coefficient $\log\frac{k^{+}}{\Lambda}+\log\frac{\vee}{k^{+}}+\log(1-\frac{k^{+}}{\vee})$; here $-\frac{1}{p^{+}}=-\frac{k^{+}}{p^{+}(p^{+}+k^{+})}-\frac{1}{p^{+}+k^{+}}$ and $\int_{\Lambda}^{\vee-k^{+}}\frac{k^{+}dp^{+}}{p^{+}(p^{+}+k^{+})}=\log\frac{k^{+}}{\Lambda}+\log(1-\frac{k^{+}}{\vee})$, so the coefficient has no $\log\frac{\vee}{k^{+}}$ (second BFKL term, which cancels the first); and the $-\frac{1}{k^{+}}\delta_{im}\delta_{jk'}$ term is subtracted at $y'\to z$, the subtracted piece being exactly $\log\frac{\vee}{k^{+}}\times$(third BFKL term), because $\int_{y'}e^{-ik\cdot(w-y')}e^{-ik\cdot\frac{p^{+}}{k^{+}}(z-y')}\frac{(y'-z)^{m}}{(y'-z)^{2}}=\frac{k^{+}}{p^{+}+k^{+}}\int_{w'}e^{-ik\cdot(w-w')}\frac{(w'-z)^{m}}{(w'-z)^{2}}$ (checked numerically to $10^{-11}$). The identity for the kernel holds at fixed $y'$, the one for the subtracted term after the $y'$ integration; that integral is only conditionally convergent (Gaussian regulator in the check). The plus prescription acts on the phase $e^{-ik\cdot\frac{p^{+}}{k^{+}}(z-y')}$]\end{notebox}
\noindent{\small\textsf{Unmarked edit, l.211 (deleted):} \texttt{)\allowbreak{}}}\par
\noindent{\small\textsf{Unmarked edit, l.211 (deleted):} \texttt{)\allowbreak{}}}\par
\noindent{\small\textsf{Unmarked edit, l.212 (deleted):} \texttt{)\allowbreak{}}}\par
\noindent{\small\textsf{Unmarked edit, l.212 (deleted):} \texttt{)\allowbreak{}}}\par
\noindent{\small\textsf{Unmarked edit, l.213 (deleted):} \texttt{)\allowbreak{}}}\par
\needspace{3\baselineskip}\noindent\rule{\linewidth}{0.3pt}\par\noindent\textbf{Lines 217--227} \quad {\small (in \texttt{align}, lines 216--225)}\par
\begin{itemize}[leftmargin=2.2em,itemsep=1pt,topsep=2pt]
\item[\textbf{l.217}] \textsf{\small [phase convention]} \emph{added:} $\displaystyle )$\quad{\color{gray}\small(in context: }$\displaystyle y')${\color{gray}\small)}
\end{itemize}
\begin{notebox}\textsf{\small Note, l.225:} [Checked: \textbf{this form matches}, in the same sense and with the same check as the first part (three-$\rho$ list replaced by the two-$\rho$ list, Group III Row III). For $\vee/k^{+}=6$: left-hand side $-0.0048036-0.0070435\,i$, $-0.0048166-0.0071848\,i$, $-0.0048167-0.0071862\,i$ for $\Lambda/k^{+}=10^{-2},10^{-4},10^{-6}$; this equation: $-0.0048167-0.0071862\,i$; without the second and third BFKL terms the difference is $0.020$. Here the plus prescription acts on the phase and on $1/[p^{+}(x-z)^{2}+k^{+}(x-w)^{2}]$. The term $-\frac{p^{+}}{k^{+2}}[\dots]$ is subtracted at $y'\to z$; with $\frac{p^{+}}{k^{+}(p^{+}+k^{+})[p^{+}a+k^{+}b]}=\frac{1}{k^{+}a\,(p^{+}+k^{+})}-\frac{b}{a\,(p^{+}+k^{+})[p^{+}a+k^{+}b]}$, $a=(x-z)^{2}$, $b=(x-w)^{2}$, the subtracted piece is $\log\frac{\vee}{k^{+}}\times$(third BFKL term) minus the block with $e^{-ik\cdot(w-w')}$ (third term here), which is therefore needed and has the correct sign and kernel. The term $-\frac{1}{k^{+}}[\dots]$ needs no subtraction]\end{notebox}
\begin{notebox}\textsf{\small Note, l.227:} [Not re-checked in this pass (four-$\rho$ only in the checked file).]\end{notebox}
\noindent{\small\textsf{Unmarked edit, l.217 (deleted):} \texttt{)\allowbreak{}}}\par
\needspace{3\baselineskip}\noindent\rule{\linewidth}{0.3pt}\par\noindent\textbf{Lines 233} \quad \textsf{Group III (with c.c.) $\rangle$ Row V}\par
\begin{notebox}\textsf{\small Note, l.233:} [Checked: \textbf{this form matches}; kernel, subtraction and logarithms as in Row III, first part, with twice the prefactor; colour structure as in the checked Row V ($f^{c'ca}$ without $U(z)$, after $U^{d'b}(z)U^{cb}(z)=\delta^{d'c}$; the subtracted bracket is the first bracket at $y'\to z$). Same check: for $\vee/k^{+}=6$ the left-hand side is $0.0089333+0.0453736\,i$, $0.0103322+0.0461993\,i$, $0.0103462+0.0462076\,i$ for $\Lambda/k^{+}=10^{-2},10^{-4},10^{-6}$; this equation: $0.0103463+0.0462076\,i$; without the second and third BFKL terms of Row V the difference is $0.32$]\end{notebox}
\needspace{3\baselineskip}\noindent\rule{\linewidth}{0.3pt}\par\noindent\textbf{Lines 238} \quad \textsf{Group III (with c.c.) $\rangle$ Row VI}\par
\begin{notebox}\textsf{\small Note, l.238:} [Checked: \textbf{this form matches}; it has the same three terms as Rows III and V (the first two are written on one source line), the same kernels and numbers as Row V, and the colour structure of the checked Row VI ($f^{c'd'a}$ with $U^{d'b}(z)$ in the first bracket and $U^{cb}(z)$ in the second; the subtracted bracket is the first bracket at $y'\to z$). The three terms of Row VI in the BFKL file are needed in the same way]\end{notebox}
\needspace{15\baselineskip}\bigskip
\subsection{Divergent terms}\label{sec:file-divergent}
\noindent\textbf{Original:} \texttt{/root/.claude/uploads/0ab97fcf-94a2-5aec-8bad-586400481685/371e60d2-Divergent\_terms.tex}\\\textbf{Checked copy:} \texttt{/home/user/correlations/nlo\_review/limits\_checked/Divergent\_terms\_checked.tex}\\
Lines: 113 (original), 113 (checked copy). Coloured groups of the check: 36 (36 top-level); author\textquotesingle s own coloured groups left unchanged: 8.\\
Records by type: phase convention 11; missing term added 2; typo/LaTeX 1; note only 11.
\par\noindent Unmarked layout-only edits (line breaks \texttt{\textbackslash notag\textbackslash\textbackslash\ \&}, \texttt{\textbackslash par\textbackslash noindent} around inserted notes, spacing), checked-copy lines: 35, 42, 53, 60, 64, 68, 86, 90, 112.
\subsubsection*{Changes, block by block}
\needspace{3\baselineskip}\noindent\rule{\linewidth}{0.3pt}\par\noindent\textbf{Lines 35}\par
\begin{itemize}[leftmargin=2.2em,itemsep=1pt,topsep=2pt]
\item[\textbf{l.35}] \textsf{\small [typo/LaTeX]} \emph{added:} \textbf{Checked copy.} Corrections and additions of the check are in red, notes of the check in red in square brackets;
\end{itemize}
\begin{notebox}\textsf{\small Note, l.35:} Orange is the author's own highlighting and is left as it is. Phase: $e^{-ik\cdot(w-w')}$ (conjugated bracket, with the primed coordinates and standing left, at $w'$), and correspondingly $e^{-ik\cdot(x-w')}$, $e^{-ik\cdot(y-y')}$, $e^{-ik\cdot(y-y')/\xi}$; c.c.\ is implied for Group I, not for Group II. $L\equiv\log\frac{k^{+}}{\Lambda}$, $\ell\equiv\log\big(1-\frac{k^{+}}{\vee}\big)$ in the notes.\end{notebox}
\needspace{3\baselineskip}\noindent\rule{\linewidth}{0.3pt}\par\noindent\textbf{Lines 41--42} \quad \textsf{Group I (with c.c.) $\rangle$ Row II} \quad {\small (in \texttt{align}, lines 40--42)}\par
\begin{itemize}[leftmargin=2.2em,itemsep=1pt,topsep=2pt]
\item[\textbf{l.41}] \textsf{\small [phase convention]} $\displaystyle \, \int_{w, w',x,x',} e^{-ik(w^\prime-w)}\, \frac{(x'-w')\cdot (x-w)}{(x'-w')^{2}\, (x-w)^{2}}\, {\color{orange}[-\frac{2}{\epsilon}-2\gamma-\log{(\frac{(x-w)^{2}\,\mu^{2}_{\overline{MS}}}{4})}]}$\ $\longrightarrow$\ $\displaystyle \, \int_{w, w',x,x',} e^{-ik(w-w^\prime)}\, \frac{(x'-w')\cdot (x-w)}{(x'-w')^{2}\, (x-w)^{2}}\, {\color{orange}[-\frac{2}{\epsilon}-2\gamma-\log{(\frac{(x-w)^{2}\,\mu^{2}_{\overline{MS}}}{4})}]}$
\end{itemize}
\begin{notebox}\textsf{\small Note, l.42:} [Checked against $N_{1}$ of the checked Group I file: prefactor, sign of the bracket and $\frac{11}{3}+\ell$ agree. There the bracket multiplies $\frac{11}{3}-2L+\ell$; the $-2L$ part is in the JIMWLK terms. The finite terms of that row, $-\frac{67}{9}+\frac{\pi^{2}}{3}+\ell^{2}+\ell\,\log\frac{\vee}{k^{+}}$, carry no pole and stay in the left-over terms. $\gamma\equiv\gamma_{E}$.]\end{notebox}
\needspace{3\baselineskip}\noindent\rule{\linewidth}{0.3pt}\par\noindent\textbf{Lines 52--53} \quad \textsf{Group I (with c.c.) $\rangle$ Row V} \quad {\small (in \texttt{align}, lines 51--53)}\par
\begin{itemize}[leftmargin=2.2em,itemsep=1pt,topsep=2pt]
\item[\textbf{l.52}] \textsf{\small [phase convention]} $\displaystyle ]\int_{x,y,x',w'}e^{-ik\cdot(w'-x)}\,[\frac{2}{\epsilon}$\ $\longrightarrow$\ $\displaystyle ]\int_{x,y,x',w'}e^{-ik\cdot(x-w')}\,[\frac{2}{\epsilon}$
\end{itemize}
\begin{notebox}\textsf{\small Note, l.53:} [Checked against $N_{2b\,\rm UV}$ of the checked Group I file: the full coefficient there is $2L-\frac{11}{6}$ ($\frac{1}{k^{+}}\int_{\Lambda}^{k^{+}-\Lambda}dp^{+}\big[\frac{p^{+}(k^{+}-p^{+})}{k^{+2}}+\frac{p^{+}}{k^{+}-p^{+}}+\frac{k^{+}-p^{+}}{p^{+}}\big]$, recomputed); the $2L$ part is in the JIMWLK terms. $\bar\mu_{\overline{MS}}$ has to be the same scale as $\mu_{\overline{MS}}$ of Row II for the cancellation below.]\end{notebox}
\needspace{3\baselineskip}\noindent\rule{\linewidth}{0.3pt}\par\noindent\textbf{Lines 56} \quad \textsf{Adding Group I} \quad {\small (in \texttt{align}, lines 55--57)}\par
\begin{itemize}[leftmargin=2.2em,itemsep=1pt,topsep=2pt]
\item[\textbf{l.56}] \textsf{\small [phase convention]} $\displaystyle \, \int_{w, w',x,x',} e^{-ik(w^\prime-w)}\, \frac{(x'-w')\cdot (x-w)}{(x'-w')^{2}\, (x-w)^{2}}\, {\color{orange}[-\frac{2}{\epsilon}-2\gamma-\log{(\frac{(x-w)^{2}\,\mu^{2}_{\overline{MS}}}{4})}]}$\ $\longrightarrow$\ $\displaystyle \, \int_{w, w',x,x',} e^{-ik(w-w^\prime)}\, \frac{(x'-w')\cdot (x-w)}{(x'-w')^{2}\, (x-w)^{2}}\, {\color{orange}[-\frac{2}{\epsilon}-2\gamma-\log{(\frac{(x-w)^{2}\,\mu^{2}_{\overline{MS}}}{4})}]}$
\item[\textbf{l.56}] \textsf{\small [phase convention]} $\displaystyle ]\int_{x,y,x',w'}e^{-ik\cdot(w'-x)}\,[\frac{2}{\epsilon}$\ $\longrightarrow$\ $\displaystyle ]\int_{x,y,x',w'}e^{-ik\cdot(x-w')}\,[\frac{2}{\epsilon}$
\end{itemize}
\needspace{3\baselineskip}\noindent\rule{\linewidth}{0.3pt}\par\noindent\textbf{Lines 59--60} \quad {\small (in \texttt{align}, lines 58--60)}\par
\begin{itemize}[leftmargin=2.2em,itemsep=1pt,topsep=2pt]
\item[\textbf{l.59}] \textsf{\small [phase convention]} $\displaystyle \, \int_{w, w',x,x',} e^{-ik(w^\prime-w)}\, \frac{(x'-w')\cdot (x-w)}{(x'-w')^{2}\, (x-w)^{2}}\, {\color{orange}[-\frac{2}{\epsilon}-2\gamma-\log{(\frac{(x-w)^{2}\,\mu^{2}_{\overline{MS}}}{4})}]}$\ $\longrightarrow$\ $\displaystyle \, \int_{w, w',x,x',} e^{-ik(w-w^\prime)}\, \frac{(x'-w')\cdot (x-w)}{(x'-w')^{2}\, (x-w)^{2}}\, {\color{orange}[-\frac{2}{\epsilon}-2\gamma-\log{(\frac{(x-w)^{2}\,\mu^{2}_{\overline{MS}}}{4})}]}$
\item[\textbf{l.59}] \textsf{\small [phase convention]} $\displaystyle ]\int_{x,y,x',w'}e^{-ik\cdot(w'-x)}\,[\frac{2}{\epsilon}$\ $\longrightarrow$\ $\displaystyle ]\int_{x,y,x',w'}e^{-ik\cdot(x-w')}\,[\frac{2}{\epsilon}$
\end{itemize}
\begin{notebox}\textsf{\small Note, l.60:} [Checked. For the next step rename $x\to w$, $y\to x$ in the second term: the phase becomes $e^{-ik\cdot(w-w')}$, the kernel $\frac{(x'-w')\cdot(x-w)}{(x'-w')^{2}(x-w)^{2}}$, the bracket $\big[\frac{2}{\epsilon}+2\gamma_{E}+\log\frac{(x-w)^{2}\mu^{2}}{4}\big]$, and $[U(y)-U(x)]\rho(y)\to-[U(w)-U(x)]\rho(x)$; the order $\rho(x')\rho(x)$ is the same in both terms. The $\frac{11}{3}$ terms then cancel completely (pole, $2\gamma_{E}$ and logarithm), provided $\bar\mu_{\overline{MS}}=\mu_{\overline{MS}}$.]\end{notebox}
\needspace{3\baselineskip}\noindent\rule{\linewidth}{0.3pt}\par\noindent\textbf{Lines 63--64} \quad {\small (in \texttt{align}, lines 62--64)}\par
\begin{itemize}[leftmargin=2.2em,itemsep=1pt,topsep=2pt]
\item[\textbf{l.63}] \textsf{\small [phase convention]} $\displaystyle \, \int_{w, w',x,x',} e^{-ik(w^\prime-w)}\, \frac{(x'-w')\cdot (x-w)}{(x'-w')^{2}\, (x-w)^{2}}\, {\color{orange}[-\frac{2}{\epsilon}-2\gamma-\log{(\frac{(x-w)^{2}\,\mu^{2}_{\overline{MS}}}{4})}]}$\ $\longrightarrow$\ $\displaystyle \, \int_{w, w',x,x',} e^{-ik(w-w^\prime)}\, \frac{(x'-w')\cdot (x-w)}{(x'-w')^{2}\, (x-w)^{2}}\, {\color{orange}[-\frac{2}{\epsilon}-2\gamma-\log{(\frac{(x-w)^{2}\,\mu^{2}_{\overline{MS}}}{4})}]}$
\end{itemize}
\begin{notebox}\textsf{\small Note, l.64:} [Checked: this is what remains.]\end{notebox}
\needspace{3\baselineskip}\noindent\rule{\linewidth}{0.3pt}\par\noindent\textbf{Lines 68}\par
\begin{notebox}\textsf{\small Note, l.68:} [Correct at fixed $\epsilon$: $\Lambda$ does not occur and $\ell\to0$ for $\vee\to\infty$. Order of limits: the term is $\big(\frac{2}{\epsilon}+{\rm finite}\big)\times\ell$, a pole times a vanishing coefficient; it is zero only if $\vee\to\infty$ is taken before $\epsilon\to0$ (or after the pole has been removed); at finite $\vee$ the pole is there. Pole bookkeeping, coefficient of $\frac{1}{(2\pi)^{3}}\frac{g^{4}N_{c}}{8\pi^{4}}\frac{1}{k^{+}}\,\frac{2}{\epsilon}\int_{w,w',x,x'}e^{-ik\cdot(w-w')}\frac{(x'-w')\cdot(x-w)}{(x'-w')^{2}(x-w)^{2}}[U(w')-U(x')][U(w)-U(x)]\rho(x')\rho(x)$ (checked rows, c.c.\ included for Group I): Group I Row II: $\frac{11}{6}-L+\frac{\ell}{2}$; Group I Row V: $-\frac{11}{6}+2L$; Group II Row IV, term with $L+\ell$: $-L-\ell$; Group II Row IV, RC: $+\frac{\beta_{0}}{4N_{c}}=\frac{11}{12}$; Group II Row IV, $\delta(1-\xi)$ term of DGLAP: $-\frac{11}{12}$. Sum: $-\frac{\ell}{2}$. So the $\frac{11}{3}$ poles and the $L$ poles cancel, the two $\beta_{0}$ terms cancel, and the pole proportional to $\ell$ does \emph{not} cancel between Group I and Group II Row IV (it vanishes only for $\vee\to\infty$). The finite logarithms of the three $L$, $\ell$ terms have different arguments, $(x-w)^{2}$ (ultraviolet, Group I) and $(w-w')^{2}$ (collinear, Group II Row IV), and do not cancel. In addition there is the pole of the remaining part of the DGLAP term, with the integrand at $k/\xi$ (below).]\end{notebox}
\needspace{3\baselineskip}\noindent\rule{\linewidth}{0.3pt}\par\noindent\textbf{Lines 81} \quad \textsf{Group II $\rangle$ Row IV} \quad {\small (in \texttt{align}, lines 80--83)}\par
\begin{itemize}[leftmargin=2.2em,itemsep=1pt,topsep=2pt]
\item[\textbf{l.81}] \textsf{\small [phase convention]} $\displaystyle (\xi)\,\int_{x,y,x',y'}e^{-ik\cdot(y'-y)\frac{1}{\xi}}\Big[2\gamma_{E}$\ $\longrightarrow$\ $\displaystyle (\xi)\,\int_{x,y,x',y'}e^{-ik\cdot(y-y')\frac{1}{\xi}}\Big[2\gamma_{E}$
\end{itemize}
\needspace{3\baselineskip}\noindent\rule{\linewidth}{0.3pt}\par\noindent\textbf{Lines 86}\par
\begin{notebox}\textsf{\small Note, l.86:} [Checked against the checked Row IV file. (i) This is its DGLAP term for $\vee\to\infty$ (there the lower limit is $\frac{k^{+}}{\vee}$), if the plus prescriptions act on the whole function that multiplies them, $f(\xi)\equiv\frac{1}{\xi^{2}}\int_{x,y,x',y'}e^{-ik\cdot(y-y')\frac{1}{\xi}}[\dots]\dots$: $\int_{0}^{1}d\xi\,\big[\frac{1-\xi}{\xi}\big]_{+}f(\xi)\equiv\int_{0}^{1}d\xi\,\frac{1-\xi}{\xi}\,[f(\xi)-f(0)]$ and $\int_{0}^{1}d\xi\,\frac{\xi}{(1-\xi)_{+}}f(\xi)\equiv\int_{0}^{1}d\xi\,\frac{\xi f(\xi)-f(1)}{1-\xi}$. Here $f(0)=\lim_{\xi\to0}f(\xi)=-\frac{4\pi}{k^{2}}\int_{x,x',z}\frac{(x-z)\cdot(x'-z)}{(x-z)^{2}(x'-z)^{2}}[U^{ce'}(z)-U^{ce'}(x')][U^{ce}(z)-U^{ce}(x)]\rho^{e'}(x')\rho^{e}(x)$ (from $y'-y=\xi s$ and $\int d^{2}s\,e^{-ik\cdot s}\log s^{2}=-\frac{4\pi}{k^{2}}$ for $k\neq0$; verified numerically), so the subtraction is exactly the term $+2N_{c}\frac{1-\xi}{\xi}\frac{4\pi}{k^{2}}\int_{x,x',z}\dots$ of the checked file. Note that $f(0)\neq0$ although the integrand without the $\frac{1}{\xi^{2}}$ vanishes at $\xi=0$; without the subtraction the integral diverges as $\log\frac{\vee}{k^{+}}$. (ii) The term $-\big(1-\frac{k^{+}}{\vee}\big)\times$(BFKL structure) of the checked file is not contained in this form; it is in the left-over terms and must be kept there. (iii) The $\xi\to0$ end converges if $f(\xi)-f(0)$ vanishes like a power of $\xi$ (found for smooth test functions; for the Wilson-line structure this is an assumption). (iv) In the checked file the bracket is $\big[\frac{2}{\epsilon}+2\gamma_{E}+\log\dots\big]$, here and in $N_{\rm RC}$; the pole of the DGLAP term, $\frac{2}{\epsilon}\int d\xi\,\frac{\mathcal P_{gg}(\xi)}{\xi^{2}}\times$(LO integrand with $k\to k/\xi$), is not cancelled by any other term. (v) The $\delta(1-\xi)$ term is exactly $-N_{\rm RC}$: $N_{\rm DGLAP}+N_{\rm RC}$ does not depend on $\beta_{0}$ (added and subtracted), so $N_{\rm RC}$ may be used only together with the $\delta(1-\xi)$ term here.]\end{notebox}
\needspace{3\baselineskip}\noindent\rule{\linewidth}{0.3pt}\par\noindent\textbf{Lines 88--90} \quad {\small (in \texttt{align}, lines 87--90)}\par
\begin{itemize}[leftmargin=2.2em,itemsep=1pt,topsep=2pt]
\item[\textbf{l.88}] \textsf{\small [phase convention]} $\displaystyle \int_{x,y,x',y'}e^{-ik\cdot(y'-y)}\Big[2\gamma_{E}$\ $\longrightarrow$\ $\displaystyle \int_{x,y,x',y'}e^{-ik\cdot(y-y')}\Big[2\gamma_{E}$
\end{itemize}
\begin{notebox}\textsf{\small Note, l.90:} [$N_{\rm RC}$ agrees with the checked file (sign, factor, bracket without the $\frac{2}{\epsilon}$). Missing in this list: the third term of the checked Row IV with a pole bracket, left over when only $L$ is taken out of the coefficient $L+\ell$:]\end{notebox}
\needspace{3\baselineskip}\noindent\rule{\linewidth}{0.3pt}\par\noindent\textbf{Lines 90} \quad {\small (in \texttt{align}, lines 90--90)}\par
\begin{itemize}[leftmargin=2.2em,itemsep=1pt,topsep=2pt]
\item[\textbf{l.90}] \textsf{\small [missing term added]} \emph{added:} $\displaystyle \frac{d^{3}N_{\log(1-k^{+}/\vee)}}{d^{3}k}=$
\item[\textbf{l.90}] \textsf{\small [phase convention]} \emph{added:} $\displaystyle -\frac{1}{(2\pi)^{3}}\,\frac{g^{4}\,N_{c}}{8\pi^{4}}\,\frac{1}{k^{+}}\,\log\Big(1-\frac{k^{+}}{\vee}\Big)\int_{x,y,x',y'}e^{-ik\cdot(y-y')}\Big[\frac{2}{\epsilon}+2\gamma_{E}+\log\frac{(y'-y)^{2}\mu^{2}_{\overline{MS}}}{4}\Big]\frac{(x-y)\cdot(x'-y')}{(x-y)^{2}(x'-y')^{2}}$
\item[\textbf{l.90}] \textsf{\small [missing term added]} \emph{added:} $\displaystyle \times[U^{ce'}(y') - U^{ce'}(x')][U^{ce}(y) - U^{ce}(x)]\rho^{e'}(x')\rho^{e}(x)\;\longrightarrow\;0\quad(\vee\to\infty\ \text{at fixed }\epsilon)$
\end{itemize}
\begin{notebox}\textsf{\small Note, l.90:} [It has the same order-of-limits caveat as the Group I term above and does not cancel it (coefficients $-\ell$ and $+\frac{\ell}{2}$ in the units of the note there).]\end{notebox}
\needspace{3\baselineskip}\noindent\rule{\linewidth}{0.3pt}\par\noindent\textbf{Lines 110--112} \quad \textsf{Running coupling} \quad {\small (in \texttt{align}, lines 109--112)}\par
\begin{itemize}[leftmargin=2.2em,itemsep=1pt,topsep=2pt]
\item[\textbf{l.110}] \textsf{\small [phase convention]} $\displaystyle \int_{x,y,x',y'}e^{-ik\cdot(y'-y)}[1$\ $\longrightarrow$\ $\displaystyle \int_{x,y,x',y'}e^{-ik\cdot(y-y')}[1$
\end{itemize}
\begin{notebox}\textsf{\small Note, l.112:} [Checked: $N_{\rm RC}$ has the prefactor $\frac{1}{(2\pi)^{3}}\frac{g^{2}}{2\pi^{2}}\frac{1}{k^{+}}$ and the integrand of the LO cross section of the main file ($|\frac{ig}{\sqrt{2k^{+}}\pi}|^{2}$; there the points are $w,w'$, here $y,y'$), times $\frac{g^{2}\beta_{0}}{16\pi^{2}}[\dots]$; $2\gamma_{E}+\log\frac{r^{2}\mu^{2}}{4}=\log\frac{r^{2}\mu^{2}e^{2\gamma_{E}}}{4}$. Sign: to this order the bracket is $\alpha_{s}(Q^{2})/\alpha_{s}(\mu^{2})$ with $Q^{2}=\frac{4e^{-2\gamma_{E}}}{(y'-y)^{2}}$ and $\alpha_{s}(Q^{2})=\alpha_{s}(\mu^{2})\big/\big[1+\frac{\alpha_{s}\beta_{0}}{4\pi}\log\frac{Q^{2}}{\mu^{2}}\big]$, i.e.\ asymptotically free for $\beta_{0}>0$. $g^{2}\to g^{2}(\mu)$ presupposes that the $\frac{2}{\epsilon}$ of the checked $N_{\rm RC}$ is removed by the coupling renormalisation; in this row that pole is a collinear pole created by adding and subtracting the $\delta(1-\xi)$ term, so the statement holds only together with the DGLAP term with $\delta(1-\xi)$ and the same $\mu$.]\end{notebox}
\begin{notebox}\textsf{\small Note, l.112:} [Decisions recorded: no $1/\epsilon$ poles are written (rule $\int d^{2}p/p^{2}=0$), so the order of the limits $\vee\to\infty$ and $\epsilon\to0$ does not arise; the running-coupling term is the $\delta(1-\xi)$ part of DGLAP, with one scale $\mu$; both plus prescriptions of $\widetilde{\mathcal{P}}_{gg}$ act on the whole function of $\xi$; $\bar\mu_{\overline{MS}}=\mu_{\overline{MS}}$, so the $11/3$ terms of Rows II and V of Group I cancel. In $N_{2b\,\rm UV}$ the factor $d_{\perp}$ of the $p^{+}(k^{+}-p^{+})/k^{+2}$ term multiplies the angular average $\delta^{im}/d_{\perp}$, i.e. $d_{\perp}\cdot\frac{1}{d_{\perp}}=1$ exactly, so no extra finite constant ($-1/6$) arises.]\end{notebox}
\needspace{15\baselineskip}\bigskip
\subsection{scales}\label{sec:file-scales}
\noindent\textbf{Original:} \texttt{/root/.claude/uploads/0ab97fcf-94a2-5aec-8bad-586400481685/8999d151-scales.tex}\\\textbf{Checked copy:} \texttt{/home/user/correlations/nlo\_review/limits\_checked/scales\_checked.tex}\\
Lines: 82 (original), 82 (checked copy). Coloured groups of the check: 19 (19 top-level); author\textquotesingle s own coloured groups left unchanged: 3.\\
Records by type: phase convention 4; typo/LaTeX 2; note only 8.
\par\noindent Unmarked layout-only edits (line breaks \texttt{\textbackslash notag\textbackslash\textbackslash\ \&}, \texttt{\textbackslash par\textbackslash noindent} around inserted notes, spacing), checked-copy lines: 44, 47, 51, 80, 81.
\subsubsection*{Changes, block by block}
\needspace{3\baselineskip}\noindent\rule{\linewidth}{0.3pt}\par\noindent\textbf{Lines 35}\par
\begin{itemize}[leftmargin=2.2em,itemsep=1pt,topsep=2pt]
\item[\textbf{l.35}] \textsf{\small [typo/LaTeX]} \emph{added:} \ = removed. Green, orange and blue are the author\textquotesingle s own highlighting and are unchanged.]
\end{itemize}
\begin{notebox}\textsf{\small Note, l.35:} [Marks of the check: red = corrected or added (part of the expression), red text in square brackets = notes,\end{notebox}
\needspace{3\baselineskip}\noindent\rule{\linewidth}{0.3pt}\par\noindent\textbf{Lines 43--44} \quad \textsf{Group I (with c.c.) $\rangle$ Row III} \quad {\small (in \texttt{align}, lines 42--44)}\par
\begin{itemize}[leftmargin=2.2em,itemsep=1pt,topsep=2pt]
\item[\textbf{l.43}] \textsf{\small [phase convention]} $\displaystyle \int_{w, w^\prime} e^{-ik(w^\prime-w)}\,$\ $\longrightarrow$\ $\displaystyle \int_{w, w^\prime} e^{-ik\cdot(w-w^\prime)}\,$
\end{itemize}
\begin{notebox}\textsf{\small Note, l.44:} [Checked: both green brackets are the double-pole parts of the first two terms of the $p^{+}$ integrand of Row III: with $D=p^{+}(x-z)^{2}-k^{+}(x-w)^{2}$, $\frac{p^{+}+k^{+}}{(k^{+}-p^{+})^{2}\,D}$ $=\frac{2}{(x-z)^{2}-(x-w)^{2}}\,\frac{1}{(p^{+}-k^{+})^{2}}$ $+$ simple poles (exact partial fractions), $\int_{k^{+}+\Lambda}^{\vee}\frac{dp^{+}}{(p^{+}-k^{+})^{2}}=\frac{1}{\Lambda}-\frac{1}{\vee-k^{+}}$, $\int_{\Lambda}^{k^{+}-\Lambda}\frac{dp^{+}}{(p^{+}-k^{+})^{2}}=\frac{1}{\Lambda}-\frac{1}{k^{+}-\Lambda}$. Agrees with $N_{1}$ of the checked Group I file. Phase: convention $e^{-ik\cdot(w-w^\prime)}$ (here and below).]\end{notebox}
\needspace{3\baselineskip}\noindent\rule{\linewidth}{0.3pt}\par\noindent\textbf{Lines 46--47} \quad {\small (in \texttt{align}, lines 45--47)}\par
\begin{itemize}[leftmargin=2.2em,itemsep=1pt,topsep=2pt]
\item[\textbf{l.46}] \textsf{\small [phase convention]} $\displaystyle \int_{w, w^\prime} e^{-ik(w^\prime-w)}\,$\ $\longrightarrow$\ $\displaystyle \int_{w, w^\prime} e^{-ik\cdot(w-w^\prime)}\,$
\end{itemize}
\begin{notebox}\textsf{\small Note, l.47:} [Checked: the $p^{+}>k^{+}$ non-instantaneous term has the double pole $-\frac{1}{(x-z)^{2}-(x-w)^{2}}\frac{1}{(p^{+}-k^{+})^{2}}$ (only the $\delta_{ij}\delta_{km}$ structure contributes; $(z-w)\cdot[(x-w)-(x-z)]=(z-w)^{2}$), range $k^{+}+\Lambda<p^{+}<\vee$; with the overall minus sign this is the expression above.]\end{notebox}
\needspace{3\baselineskip}\noindent\rule{\linewidth}{0.3pt}\par\noindent\textbf{Lines 49--51} \quad {\small (in \texttt{align}, lines 48--51)}\par
\begin{itemize}[leftmargin=2.2em,itemsep=1pt,topsep=2pt]
\item[\textbf{l.49}] \textsf{\small [phase convention]} $\displaystyle \int_{w,w'}e^{-ik(w'-w)}$\ $\longrightarrow$\ $\displaystyle \int_{w,w'}e^{-ik\cdot(w-w')}$
\end{itemize}
\begin{notebox}\textsf{\small Note, l.51:} [Checked: the $p^{+}<k^{+}$ non-instantaneous term has the same double pole $-\frac{1}{(x-z)^{2}-(x-w)^{2}}\frac{1}{(p^{+}-k^{+})^{2}}$, range $\Lambda<p^{+}<k^{+}-\Lambda$. $N_{3b}$ and the instantaneous (third) term have only simple poles: no power terms.]\end{notebox}
\needspace{3\baselineskip}\noindent\rule{\linewidth}{0.3pt}\par\noindent\textbf{Lines 55} \quad \textsf{Group I (with c.c.) $\rangle$ Row V}\par
\begin{notebox}\textsf{\small Note, l.55:} [No inverse power of a cut-off. For completeness: $N_{2b\,\rm UV}$ has the coefficient $\frac{1}{k^{+}}\big[2\log\frac{k^{+}-\Lambda}{\Lambda}-\frac{11}{6}+\frac{4\Lambda}{k^{+}}-\frac{\Lambda^{2}}{k^{+2}}+\frac{2\Lambda^{3}}{3k^{+3}}\big]$: positive powers of $\Lambda$, which vanish for $\Lambda\to0$.]\end{notebox}
\needspace{3\baselineskip}\noindent\rule{\linewidth}{0.3pt}\par\noindent\textbf{Lines 58} \quad \textsf{Group II $\rangle$ Row I}\par
\begin{notebox}\textsf{\small Note, l.58:} [No $1/\Lambda$. But $N_{2f}$ contains a pure power without logarithm, $\frac{1}{(2\pi)^{3}}\frac{g^{4}d_{\perp}}{32\pi^{5}}\big[\frac{1}{\vee}-\frac{1}{k^{+}+\Lambda}\big]\int\frac{e^{-ik\cdot(w-w^\prime)}}{(z-w)^{2}(z-w^\prime)^{2}}[\dots]$ (from the double pole at $p^{+}=-k^{+}$, $\int_{\Lambda}^{\vee-k^{+}}\frac{dp^{+}}{(p^{+}+k^{+})^{2}}$); it tends to $-\frac{1}{k^{+}}$, a finite term, and the $\frac{1}{\vee}$ part vanishes. To be listed here or not: for the author.]\end{notebox}
\needspace{3\baselineskip}\noindent\rule{\linewidth}{0.3pt}\par\noindent\textbf{Lines 64} \quad \textsf{Group II $\rangle$ Row IV}\par
\begin{notebox}\textsf{\small Note, l.64:} [No $1/\Lambda$. The BFKL coefficient is $\frac{1}{k^{+}}\log\frac{\vee}{k^{+}}-\frac{1}{k^{+}}+\frac{1}{\vee}$ ($=\frac{1}{k^{+}}\int_{k^{+}/\vee}^{1}d\xi\,\frac{1-\xi}{\xi}$): the $\frac{1}{\vee}$ vanishes for $\vee\to\infty$ and $-\frac{1}{k^{+}}$ is a finite term. To be listed here or not: for the author.]\end{notebox}
\needspace{3\baselineskip}\noindent\rule{\linewidth}{0.3pt}\par\noindent\textbf{Lines 80--81} \quad \textsf{Algebra} \quad {\small (in \texttt{align}, lines 79--81)}\par
\begin{itemize}[leftmargin=2.2em,itemsep=1pt,topsep=2pt]
\item[\textbf{l.80}] \textsf{\small [typo/LaTeX]} \emph{added:} $\displaystyle \int_{w,w'}$
\item[\textbf{l.80}] \textsf{\small [phase convention]} $\displaystyle (w^\prime-w)$\ $\longrightarrow$\ $\displaystyle \cdot(w-w^\prime)$
\end{itemize}
\begin{notebox}\textsf{\small Note, l.81:} [Checked: the left-hand side is the sum of the power terms of $N_{1}+N_{2}+N_{3a}$. With the common denominator $(x-z)^{2}-(x-w)^{2}$ the factors are $N_{1}$: $\frac{1}{\Lambda}+\frac{1}{k^{+}-\vee}+\frac{1}{\Lambda}-\frac{1}{k^{+}-\Lambda}$; $N_{2}$: $\frac{1}{\vee-k^{+}}-\frac{1}{\Lambda}$; $N_{3a}$: $-\frac{1}{\Lambda}+\frac{1}{k^{+}-\Lambda}$. The sum of the eight is identically zero at finite $\Lambda$ and $\vee$ (not only in the limit). The integral over $w,w^\prime$ was missing in front of the phase.]\end{notebox}
\needspace{15\baselineskip}\bigskip
\subsection{Limiting cases (flags)}\label{sec:file-limiting_checked}
\noindent\textbf{Original:} \texttt{/root/.claude/uploads/0ab97fcf-94a2-5aec-8bad-586400481685/f164e8ee-Limiting\_cases.tex}\\\textbf{Checked copy:} \texttt{/home/user/correlations/nlo\_review/limits\_checked/Limiting\_cases\_checked.tex}\\
Lines: 215 (original), 215 (checked copy). Coloured groups of the check: 52 (52 top-level); author\textquotesingle s own coloured groups left unchanged: 0.\\
Records by type: limits/range 1; typo/LaTeX 5; note only 44.
\par\noindent Unmarked layout-only edits (line breaks \texttt{\textbackslash notag\textbackslash\textbackslash\ \&}, \texttt{\textbackslash par\textbackslash noindent} around inserted notes, spacing), checked-copy lines: 35, 44, 50, 56, 64, 72, 96, 104, 106, 107, 110, 116, 121, 125, 137, 144, 148, 174, 182, 183, 189, 194, 198, 199, 202, 204, 209, 210, 213.
\subsubsection*{Changes, block by block}
\needspace{3\baselineskip}\noindent\rule{\linewidth}{0.3pt}\par\noindent\textbf{Lines 30}\par
\begin{itemize}[leftmargin=2.2em,itemsep=1pt,topsep=2pt]
\item[\textbf{l.30}] \textsf{\small [typo/LaTeX]} \emph{added:} $\displaystyle ,$
\end{itemize}
\needspace{3\baselineskip}\noindent\rule{\linewidth}{0.3pt}\par\noindent\textbf{Lines 35}\par
\begin{itemize}[leftmargin=2.2em,itemsep=1pt,topsep=2pt]
\item[\textbf{l.35}] \textsf{\small [typo/LaTeX]} \emph{added:} \noindent[\textbf{Checked copy.} Red $=$ corrected or added (part of the expression); red text in square brackets $=$ notes of the check; orange, struck out $=$ removed.]
\end{itemize}
\begin{notebox}\textsf{\small Flag, l.35:} Magenta \fbox{FLAG $n$} $=$ the expression is not finite as written, or a limit is not taken; the flag says what fails and where the finite form is (sections and equation numbers of \texttt{Finite\_terms\_final}).\end{notebox}
\begin{notebox}\textsf{\small Note, l.35:} [Compared, equation by equation, with \texttt{Other\_terms} (v2, checked copy: the left-over terms before $\vee\to\infty$), textually and numerically (random points, explicit SU(3) matrices; $p^{+}$, $\xi$ and transverse integrals by quadrature where stated). \textbf{Result.} Correct limits, finite as written under the conditions stated in the notes: Group I Row II, Row IV; Group II Row I; Group III Rows I, II. Correct at fixed coordinates but not integrable as written: Group I Row III (FLAG 1--3) and Row V (FLAG 4). Group II Row IV: the lower limit can be set to $0$; the second term depends on how the DGLAP term is written (FLAG 5). Group III Rows III, V, VI: $\vee\to\infty$ is not taken and cannot be taken in the form written; the form does not match the BFKL list (FLAG 6--9). \textbf{General remark on the dropped terms.} All terms $\log\big(1-\frac{k^{+}}{\vee}\big)\times$(kernel) of \texttt{Other\_terms} are dropped here. At fixed coordinates this is right, but each of them is $\log\big(1-\frac{k^{+}}{\vee}\big)$ times the JIMWLK structure of its row, whose $z$ integral diverges logarithmically at large $z$ if the colour factor does not vanish there (projectile not colour neutral): $0\times\infty$. They vanish only if the large-distance regulator is kept while $\vee\to\infty$; they are best written with the JIMWLK term, as $\big[\log\frac{k^{+}}{\Lambda}+\log\big(1-\frac{k^{+}}{\vee}\big)\big]\times$(structure). Those of Group I Row IV and Group III Row I cancel each other identically. Conventions: phase $e^{-ik\cdot(w-w')}$; c.c.\ implied for Groups I and III; logarithms are $\log|\dots|$; $\int dp^{+}\big[\frac{1}{p^{+}}\big]_{+}f\equiv\int dp^{+}\frac{f(p^{+})-f(0)}{p^{+}}$, $f$ = everything that depends on $p^{+}$, phase included; no $1/\epsilon$ written.]\end{notebox}
\needspace{3\baselineskip}\noindent\rule{\linewidth}{0.3pt}\par\noindent\textbf{Lines 39} \quad \textsf{Group I (with c.c.) $\rangle$ Row I}\par
\begin{notebox}\textsf{\small Note, l.39:} [Not re-checked; as in \texttt{Other\_terms}]\end{notebox}
\needspace{3\baselineskip}\noindent\rule{\linewidth}{0.3pt}\par\noindent\textbf{Lines 44} \quad \textsf{Group I (with c.c.) $\rangle$ Row II}\par
\begin{notebox}\textsf{\small Note, l.44:} [Checked: correct. This is $N_{1}$ of \texttt{Other\_terms} without (a) $-\big[2\gamma_{E}+\log\frac{(x-w)^{2}\mu^{2}_{\overline{MS}}}{4}\big]\big[\frac{11}{3}+\log\big(1-\frac{k^{+}}{\vee}\big)\big]$, which is in \texttt{Divergent\_terms} (its $\frac{11}{3}$ part cancels $N_{2b\,\rm UV}$ of Row V), and (b) $\log^{2}(1-\frac{k^{+}}{\vee})+\log\big(1-\frac{k^{+}}{\vee}\big)\log\frac{\vee}{k^{+}}\to0$ (numbers times the LO-type integrand). Numerically, ([\texttt{Other\_terms}] $-$ [this])/(integrand) $+\frac{11}{3}\big[2\gamma_{E}+\log\frac{(x-w)^{2}\mu^{2}}{4}\big]$ $=-1.5\times10^{-2}$, $-6.1\times10^{-4}$, $-1.1\times10^{-5}$ for $\frac{\vee}{k^{+}}=10^{2},10^{4},10^{6}$]\end{notebox}
\needspace{3\baselineskip}\noindent\rule{\linewidth}{0.3pt}\par\noindent\textbf{Lines 50}\par
\begin{notebox}\textsf{\small Note, l.50:} [Checked: identical to $N_{2b}$ of \texttt{Other\_terms}]\end{notebox}
\needspace{3\baselineskip}\noindent\rule{\linewidth}{0.3pt}\par\noindent\textbf{Lines 56}\par
\begin{notebox}\textsf{\small Note, l.56:} [Checked: correct limit of $R^{i}$ of \texttt{Other\_terms}: $|R^{i}(\vee)-R^{i}|=0.20$, $2.1\times10^{-3}$, $2.1\times10^{-5}$, $2.1\times10^{-7}$ for $\frac{k^{+}}{\vee}=10^{-2},10^{-4},10^{-6},10^{-8}$ at a random point ($|R|=5.0$). The three dropped pieces have to be dropped together: with $a\equiv\frac{k^{+}}{\vee}$, $R^{i}(\vee)-R^{i}=2q^{i}\log\big(1+a\frac{p^{2}}{q^{2}}\big)+\frac{(2p^{2}-2q\cdot p+q^{2})q^{i}+q^{2}p^{i}}{(q-p)^{2}}\log\big[1+\frac{a\,(q-p)^{2}}{(1-a)\,q^{2}}\big]+q^{i}\log(1-a)$ (identity checked to $10^{-15}$), which is finite at $q\to p$ ($1.8\times10^{-3}$ for $a=10^{-3}$, $|q-p|=10^{-1}\dots10^{-5}$), whereas dropping the third term of $R^{i}$ alone, or the $\vee$ of the second logarithm alone, leaves $a\,\frac{2p^{2}p^{i}}{(q-p)^{2}}$ ($4.3\times10^{-2}\dots4.3\times10^{6}$ at the same points). Unchanged with respect to \texttt{Other\_terms}: $R^{i}/(p^{2}q^{2})$ is not integrable at $p\to0$ and at $q\to p$ by itself; it is defined by the rule $\int d^{2}p/p^{2}=0$, i.e.\ as the kernel antisymmetrised in $p\to q-p$, Eq.~(1.4) of \texttt{Finite\_terms\_final} (earlier check, not repeated)]\end{notebox}
\needspace{3\baselineskip}\noindent\rule{\linewidth}{0.3pt}\par\noindent\textbf{Lines 64} \quad \textsf{Group I (with c.c.) $\rangle$ Row III}\par
\begin{notebox}\textsf{\small Note, l.64:} [Checked: correct at fixed coordinates. Dropped from $N_{1}$ of \texttt{Other\_terms}: its first term, $\log\big(1-\frac{k^{+}}{\vee}\big)$ times the kernel $\frac{(x-z)^{2}+(x-w)^{2}}{[(x-z)^{2}-(x-w)^{2}]^{2}}$, and $\frac{k^{+}}{\vee}(x-w)^{2}$ in two logarithms; relative difference [\texttt{Other\_terms}] $-$ [this]: $1.0\times10^{-2}$, $1.0\times10^{-4}$, $1.0\times10^{-6}$ for $\frac{k^{+}}{\vee}=10^{-2},10^{-4},10^{-6}$. See FLAG 1]\end{notebox}
\needspace{3\baselineskip}\noindent\rule{\linewidth}{0.3pt}\par\noindent\textbf{Lines 72}\par
\begin{notebox}\textsf{\small Note, l.72:} [Checked: correct at fixed coordinates. Dropped from $N_{2}$ of \texttt{Other\_terms}: its first two terms ($\log\big(1-\frac{k^{+}}{\vee}\big)$ times the two kernels with $[(x-w)^{2}-(x-z)^{2}]^{-2}$) and $\frac{k^{+}}{\vee}(x-w)^{2}$ in six logarithms; relative difference $9.8\times10^{-4}$, $9.7\times10^{-6}$, $9.7\times10^{-8}$ for $\frac{k^{+}}{\vee}=10^{-2},10^{-4},10^{-6}$. Not integrable by itself: FLAG 1]\end{notebox}
\needspace{3\baselineskip}\noindent\rule{\linewidth}{0.3pt}\par\noindent\textbf{Lines 96}\par
\begin{notebox}\textsf{\small Note, l.96:} [Checked: identical to $N_{3a}$ of \texttt{Other\_terms} (no $\vee$ in it). Not integrable by itself: FLAG 1]\end{notebox}
\needspace{3\baselineskip}\noindent\rule{\linewidth}{0.3pt}\par\noindent\textbf{Lines 104--107} \quad {\small (in \texttt{align}, lines 97--107)}\par
\begin{notebox}\textsf{\small Flag, l.104:} \fbox{FLAG 3}\end{notebox}
\begin{notebox}\textsf{\small Flag, l.106:} \fbox{FLAG 2}\end{notebox}
\begin{notebox}\textsf{\small Note, l.107:} [Checked: the first three terms are identical to those of $N_{3b}$ of \texttt{Other\_terms} (the sign of the third term is the same there and in the checked row). The last term is changed: there the logarithm is $\log\frac{[k^{+}(x-w)-(k^{+}-\Lambda)(x-z)]^{2}}{[k^{+}(x-w)-\Lambda(x-z)]^{2}}$, here its limit at fixed $z\neq w$]\end{notebox}
\begin{notebox}\textsf{\small Flag, l.107:} \fbox{FLAG 2} (last term.) $\log\frac{(z-w)^{2}}{(x-w)^{2}}$ is the limit $\Lambda\to0$ at fixed $z\neq w$ only. With the kernel $\frac{1}{(z-w)^{2}}$ and $U^{bc}(z)$, which nothing softens at $z=w$, the $z$ integral diverges at $z\to w$: $I(\varepsilon)\equiv\int_{\varepsilon<|z-w|}\!d^{2}z\,\frac{1}{(z-w)^{2}}\log\frac{(z-w)^{2}}{(x-w)^{2}}\,\Phi(z)=-2\pi\,\Phi(w)\log^{2}\frac{\varepsilon}{|x-w|}+{\rm const}$ ($\Phi$: rest of the integrand; numerically $I$ goes from $13.7$ to $976$ for $\varepsilon=10^{-1}\dots10^{-8}$ and $I+2\pi\Phi(w)\log^{2}\frac{\varepsilon}{|x-w|}=-1.386305$ for all of them). With $\Lambda$ kept and together with its partner in the row, $-\frac{ig^{4}f^{abd}}{8\pi^{5}k^{+}}\big[\log\frac{k^{+}}{\Lambda}+\log(1-\frac{\Lambda}{k^{+}})\big]\frac{(x'-w')\cdot(y-z)}{(x'-w')^{2}(y-z)^{2}(z-w)^{2}}$, the integral converges at $z\to w$ (stable to $10^{-4}$ for $\varepsilon=10^{-4}\dots10^{-10}$) but grows like $2\pi\Phi(w)\log^{2}\frac{k^{+}}{\Lambda}$ at fixed $w$ ($-6.4$, $-43.9$, $-111.4$, $-209.3$ for $\frac{\Lambda}{k^{+}}=10^{-1}\dots10^{-4}$). The finite form is the third line of Eq.~(1.5) of \texttt{Finite\_terms\_final}: $-\frac{ig^{4}f^{abd}}{8\pi^{5}k^{+}}\,\mathcal{E}\big(k\cdot(z-w)\big)\frac{(x-z)\cdot(w-z)}{(x-z)^{2}(w-z)^{2}}\frac{(x'-w')\cdot(y-w)}{(x'-w')^{2}(y-w)^{2}}$ with $U^{bc}(z)\to U^{bc}(w)$, $\mathcal{E}(s)=\int_{0}^{1}\frac{dt}{t}(e^{-ist}-1)$; it goes with the JIMWLK term in the form of the three-$\rho$ list (derivation: Eqs.~(1.9)--(1.11) there; not repeated in this check).\end{notebox}
\begin{notebox}\textsf{\small Flag, l.107:} \fbox{FLAG 3} (third term.) Correct limit, but at large $z$ the kernel tends to $-\frac{1}{z^{2}}\log\frac{z^{2}}{(x-w)^{2}}\,\frac{(x'-w')\cdot(x-w)}{(x'-w')^{2}(x-w)^{2}}$ ($z^{2}\times$angular mean$\,/\log=-1.000000$ at $|z|=10^{2},10^{3},10^{4}$) and the part $-U^{de}(x)U^{bc}(y)\rho^{e}(x)\rho^{c}(y)$ of the colour factor does not depend on $z$: $\log^{2}$ divergence at large $z$. It is finite only under one integral with the other terms that have a $\frac{1}{z^{2}}\log$ tail (FLAG 2, 4, 6--9): second line of Eq.~(1.5), Sec.~1.6 and OPEN I-2 of \texttt{Finite\_terms\_final}.\end{notebox}
\needspace{3\baselineskip}\noindent\rule{\linewidth}{0.3pt}\par\noindent\textbf{Lines 110}\par
\begin{notebox}\textsf{\small Note, l.110:} [Checked: the two forms of $\mathcal{A}_{0}$ agree (12 digits, four random points)]\end{notebox}
\begin{notebox}\textsf{\small Flag, l.110:} \fbox{FLAG 1} (Row III as a whole.) $N_{1}$, $N_{2}$, $N_{3a}$, $N_{3b}$ are not integrable one by one. On the line $(x-z)^{2}=(x-w)^{2}$ ($\delta$: distance of $x$ from it): $N_{2}$ and $N_{3a}$ grow like $\log\delta/\delta^{2}$ ($|N_{2}|=0.89$, $1.2\times10^{2}$, $1.6\times10^{4}$ at $\delta=10^{-3},10^{-4},10^{-5}$), $N_{1}$ and $N_{2}+N_{3a}$ like $1/\delta$ ($3.0\times10^{-4}$, $3.0\times10^{-3}$, $3.0\times10^{-2}$ each), the sum of the four is regular ($6.0\times10^{-7}$ at all $\delta$). At $z\to w$: $|N_{3a}|=8.0\times10^{-2}$, $7.8$, $7.6\times10^{2}$ at $|z-w|=10^{-2},10^{-3},10^{-4}$, while the sum without the last term of $N_{3b}$ is $3.9\times10^{-4}$, $4.7\times10^{-3}$, $5.6\times10^{-2}$ ($\propto\log/|z-w|$). The dropped $\log\big(1-\frac{k^{+}}{\vee}\big)$ terms of $N_{1}$, $N_{2}$ have the double pole as well: the limit $\vee\to\infty$ exists for the sum only. The sum of the four equations is identical (random points, $10^{-15}$ relative) to the combined kernel $\mathcal{K}^{ij}$ of Eqs.~(1.5), (1.6) of \texttt{Finite\_terms\_final} plus the two terms of FLAG 2, 3; that is the form to use. Its $z$ integral converges at large $z$ only conditionally (OPEN I-1 there).\end{notebox}
\needspace{3\baselineskip}\noindent\rule{\linewidth}{0.3pt}\par\noindent\textbf{Lines 116} \quad \textsf{Group I (with c.c.) $\rangle$ Row IV}\par
\begin{notebox}\textsf{\small Note, l.116:} [Checked: correct limit of Row IV of \texttt{Other\_terms} (relative difference $2.5\times10^{-2}$, $2.4\times10^{-4}$, $2.4\times10^{-6}$ for $\frac{k^{+}}{\vee}=10^{-2},10^{-4},10^{-6}$). Dropped: $\frac{k^{+}}{\vee}$ in the logarithm and the term $+\frac{ig^{4}}{8\pi^{5}k^{+}}\log\big(1-\frac{k^{+}}{\vee}\big)\frac{(x'-w')^{i}(y-z)^{j}}{(x'-w')^{2}(y-z)^{2}}\big[\frac{(x-w)^{i}(z-w)^{j}}{(z-w)^{2}(x-w)^{2}}+\frac{(x-z)^{j}(x-w)^{i}}{2(x-z)^{2}(x-w)^{2}}\big]$ (general remark; its kernel tends to $-\frac{1}{2z^{2}}\frac{(x'-w')\cdot(x-w)}{(x'-w')^{2}(x-w)^{2}}$ at large $z$). Observation: this row and Group III Row I are equal and opposite, identically in the colour algebra ($U^{T}U=1$), here, in \texttt{Other\_terms} with the $\log\big(1-\frac{k^{+}}{\vee}\big)$ terms, and in the checked rows with the cut-offs: [Row IV] $+$ [Group III Row I] $=3\times10^{-24}$, $7\times10^{-24}$, $7\times10^{-23}$ against $9\times10^{-9}$, $2\times10^{-8}$, $1\times10^{-7}$ (same order of the $\rho$'s); the two rows drop out of the sum]\end{notebox}
\needspace{3\baselineskip}\noindent\rule{\linewidth}{0.3pt}\par\noindent\textbf{Lines 120--121} \quad \textsf{Group I (with c.c.) $\rangle$ Row V} \quad {\small (in \texttt{align}, lines 119--121)}\par
\begin{itemize}[leftmargin=2.2em,itemsep=1pt,topsep=2pt]
\item[\textbf{l.120}] \textsf{\small [typo/LaTeX]} \emph{added:} $\displaystyle ,$\quad{\color{gray}\small(in context: }$\displaystyle \, \frac{dp^{+}}{2\pi}\,\int_{x,z,y,x',y',w'}\,e^{-ik\cdot (z-w^\prime)}\,e^{-ik\cdot (x-z)\frac{p^{+}}{k^{+}}}${\color{gray}\small)}
\item[\textbf{l.120}] \textsf{\small [typo/LaTeX]} \emph{added:} $\displaystyle ,$\quad{\color{gray}\small(in context: }$\displaystyle \, \frac{dq^{+}}{2\pi}\,\int_{x,z,y,x',y',w'}\,e^{-ik\cdot (x-w^\prime)}\,e^{-ik\cdot (z-x)\frac{q^{+}}{k^{+}}}${\color{gray}\small)}
\end{itemize}
\begin{notebox}\textsf{\small Note, l.121:} [Checked: identical to the first part of Row V of \texttt{Other\_terms} (no $\vee$ in this row); two commas added in the lists of integration variables]\end{notebox}
\begin{notebox}\textsf{\small Flag, l.121:} \fbox{FLAG 4} (all plus-prescription terms of Row V: this equation, $N_{2a}$, $N_{2b\,\rm rest}$.) Correct at fixed coordinates, not integrable as written. In the subtracted term the position of the soft gluon is neither in the phase nor in a denominator; e.g.\ the $\big[\frac{1}{q^{+}}\big]_{+}$ term here is $\frac{(z-x)\cdot(y-z)}{(z-x)^{2}(y-z)^{2}}\,\mathcal{E}\big(k\cdot(z-x)\big)$, $\mathcal{E}(s)=\int_{0}^{1}\frac{dt}{t}(e^{-ist}-1)$, with the tail $\frac{1}{z^{2}}\big[\gamma_{E}+\log\frac{|k|\,|z|}{2}\big]$ ($z^{2}\times$angular mean $=4.5176$, $6.8210$, $9.1235$ against $4.5184$, $6.8210$, $9.1235$ at $|z-x|=10^{2},10^{3},10^{4}$), and part of the colour factor does not depend on $z$: $\log^{2}$ divergence at large $z$. The tails cancel only in the sum over Group I Rows III, V and Group III Rows III, V, VI under one integral (Sec.~1.6 and OPEN I-2 of \texttt{Finite\_terms\_final}). At fixed coordinates this equation is Eq.~(1.13) there (verified), and $N_{2a}$, $N_{2b\,\rm rest}$ are Eqs.~(1.14)--(1.16) (earlier check, not repeated).\end{notebox}
\needspace{3\baselineskip}\noindent\rule{\linewidth}{0.3pt}\par\noindent\textbf{Lines 125}\par
\begin{notebox}\textsf{\small Note, l.125:} [Checked: identical to $N_{2a}$ of \texttt{Other\_terms}. FLAG 4 applies to the terms with $\big[\frac{k^{+}}{p^{+}}\big]_{+}$ and $\big[\frac{k^{+}}{q^{+}}\big]_{+}$]\end{notebox}
\needspace{3\baselineskip}\noindent\rule{\linewidth}{0.3pt}\par\noindent\textbf{Lines 127--137} \quad {\small (in \texttt{align}, lines 126--137)}\par
\begin{itemize}[leftmargin=2.2em,itemsep=1pt,topsep=2pt]
\item[\textbf{l.127}] \textsf{\small [typo/LaTeX]} {\small\ttfamily \}\allowbreak{})\allowbreak{}}\ $\longrightarrow$\ {\small\ttfamily )\allowbreak{}\}\allowbreak{}}
\end{itemize}
\begin{notebox}\textsf{\small Note, l.137:} [Checked: identical to $N_{2b\,\rm rest}$ of \texttt{Other\_terms}; the closing brace of the first phase is put right. $N_{2b\,\rm UV}$ of \texttt{Other\_terms} (bracket with $-\frac{11}{6}$) is in \texttt{Divergent\_terms} and cancels the $\frac{11}{3}$ term of Row II. FLAG 4 applies to the terms with $\big[\frac{k^{+}}{p^{+}}\big]_{+}$ and $\big[\frac{k^{+}}{k^{+}-p^{+}}\big]_{+}$]\end{notebox}
\needspace{3\baselineskip}\noindent\rule{\linewidth}{0.3pt}\par\noindent\textbf{Lines 144} \quad \textsf{Group II $\rangle$ Row I}\par
\begin{notebox}\textsf{\small Note, l.144:} [Checked: correct limit of $N_{1a}$ of \texttt{Other\_terms} (relative difference $2.5\times10^{-2}$, $2.4\times10^{-4}$, $2.4\times10^{-6}$ for $\frac{k^{+}}{\vee}=10^{-2},10^{-4},10^{-6}$). Dropped: $\frac{k^{+}}{\vee}$ in the logarithm and the term $-\frac{ig^{4}}{16\pi^{5}k^{+}}\log\big(1-\frac{k^{+}}{\vee}\big)\times[\dots]$ (general remark)]\end{notebox}
\needspace{3\baselineskip}\noindent\rule{\linewidth}{0.3pt}\par\noindent\textbf{Lines 148}\par
\begin{notebox}\textsf{\small Note, l.148:} [Checked: correct limit of $N_{1b}$ of \texttt{Other\_terms} (relative difference $1.2\times10^{-2}$, $1.2\times10^{-4}$, $1.2\times10^{-6}$); dropped: as in $N_{1a}$ with $w\leftrightarrow w'$ and the opposite sign]\end{notebox}
\needspace{3\baselineskip}\noindent\rule{\linewidth}{0.3pt}\par\noindent\textbf{Lines 174}\par
\begin{notebox}\textsf{\small Note, l.174:} [Checked, $N_{2a}$ to $N_{2f}$, each one separately: correct limits. Relative difference [\texttt{Other\_terms}] $-$ [this] at $\frac{k^{+}}{\vee}=10^{-2}$: $5.0\times10^{-3}$ ($2a$), $6.1\times10^{-6}$ ($2b$), $3.2\times10^{-3}$ ($2c$), $4.1\times10^{-5}$ ($2d$), $1.3\times10^{-2}$ ($2e$), $1.6\times10^{-4}$ ($2f$), decreasing like $\frac{k^{+}}{\vee}$ ($2a$, $2c$, $2e$) or like its square ($2b$, $2d$, $2f$). Changes: $\frac{k^{+}}{\vee}\to0$ in all logarithms; the first term of $N_{2b}$ of \texttt{Other\_terms} ($\log\big(1-\frac{k^{+}}{\vee}\big)\times\dots$) is dropped; in $N_{2f}$, $+\big[\frac{1}{\vee}-\frac{1}{k^{+}}\big]\to-\frac{1}{k^{+}}$. As before the limit (earlier check, not repeated): the terms of each equation, and the six equations, have to stay under one $z$ integral (the simple poles on $(x'-w')^{2}(x-z)^{2}=(x-w)^{2}(x'-z)^{2}$ and the $\frac{1}{(z-w)^{2}(z-w')^{2}}$ of $N_{2f}$ cancel only within each equation, the $\frac{\log z^{2}}{z^{2}}$ tails only in the sum); then $d_{\perp}=2$]\end{notebox}
\needspace{3\baselineskip}\noindent\rule{\linewidth}{0.3pt}\par\noindent\textbf{Lines 176} \quad \textsf{Group II $\rangle$ Row II}\par
\begin{notebox}\textsf{\small Note, l.176:} [As in \texttt{Other\_terms}]\end{notebox}
\needspace{3\baselineskip}\noindent\rule{\linewidth}{0.3pt}\par\noindent\textbf{Lines 178} \quad \textsf{Group II $\rangle$ Row III}\par
\begin{notebox}\textsf{\small Note, l.178:} [As in \texttt{Other\_terms}]\end{notebox}
\needspace{3\baselineskip}\noindent\rule{\linewidth}{0.3pt}\par\noindent\textbf{Lines 181--183} \quad \textsf{Group II $\rangle$ Row IV} \quad {\small (in \texttt{align}, lines 180--183)}\par
\begin{itemize}[leftmargin=2.2em,itemsep=1pt,topsep=2pt]
\item[\textbf{l.181}] \textsf{\small [limits/range]} $\displaystyle \frac{k^{+}}{\vee}$\ $\longrightarrow$\ $\displaystyle \,0$
\end{itemize}
\begin{notebox}\textsf{\small Flag, l.182:} \fbox{FLAG 5}\end{notebox}
\begin{notebox}\textsf{\small Note, l.183:} [Checked. First term: $N_{\rm Finite}$ of \texttt{Other\_terms}, unchanged. The lower limit can be set to $0$ (done, red): at fixed coordinates $I(a)=\int_{a}^{1}d\xi\,\frac{1-\xi}{\xi}e^{-is/\xi}$, $s=k\cdot(y-y')$, converges for $a=\frac{k^{+}}{\vee}\to0$ ($|I(a)-I(0)|=2.7\times10^{-2}$, $2.7\times10^{-4}$, $2.7\times10^{-6}$ at $s=0.37$, i.e.\ $a/|s|$), with an integrable $-\log|s|$ at $s\to0$; after the integral over $y-y'$ the integrand is ${\cal O}(\xi)$ (the square bracket is traceless). Second term: it is the limit of $-\frac{1}{(2\pi)^{3}}\frac{g^{4}N_{c}}{2\pi^{3}}\frac{1}{k^{2}}\big[\frac{1}{k^{+}}-\frac{1}{\vee}\big]\int_{x,x',z}\dots$ of \texttt{Other\_terms}; sign and prefactor are right. $N_{\rm DGLAP}$, $N_{\rm RC}$ and the $\log\big(1-\frac{k^{+}}{\vee}\big)$ term of \texttt{Other\_terms} are in \texttt{Divergent\_terms}]\end{notebox}
\begin{notebox}\textsf{\small Flag, l.183:} \fbox{FLAG 5} (second term.) Whether this term has to be there depends on how $\big[\frac{1-\xi}{\xi}\big]_{+}$ in $\widetilde{\mathcal{P}}_{gg}$ of $N_{\rm DGLAP}$ (\texttt{Divergent\_terms}, $\int_{0}^{1}d\xi$) is read. With $F(\xi)$ the function of $\xi$ that multiplies $\widetilde{\mathcal{P}}_{gg}$, $F(0)=-\frac{4\pi}{k^{2}}\int_{x,x',z}\frac{(x-z)\cdot(x'-z)}{(x-z)^{2}(x'-z)^{2}}[\dots][\dots]\rho\rho$: (A) $\frac{1-\xi}{\xi}\,[F(\xi)-F(0)]$, the form of \texttt{Other\_terms} and the reading stated in the note of \texttt{Divergent\_terms}: this term is needed, as written; with this reading the present file and \texttt{Divergent\_terms} are consistent. (B) $\frac{(1-\xi)F(\xi)-F(0)}{\xi}$, the form of Sec.~2.4.1 of \texttt{Finite\_terms\_final}: this term is already contained and must be left out, otherwise it is counted twice. [(A) $+$ this term $=$ (B) exactly, since $\int_{0}^{1}d\xi\,\big[\frac{1}{\xi}-\frac{1-\xi}{\xi}\big]=1$; model check $10^{-16}$.] The $z$ integral of this term has no phase and a kernel $\sim\frac{1}{z^{2}}$: it converges at large $z$ only if the colour factor vanishes there, as for the BFKL term itself (OPEN II/III-1 of \texttt{Finite\_terms\_final}).\end{notebox}
\needspace{3\baselineskip}\noindent\rule{\linewidth}{0.3pt}\par\noindent\textbf{Lines 189} \quad \textsf{Group III (with c.c.) $\rangle$ Row I}\par
\begin{notebox}\textsf{\small Note, l.189:} [Checked: correct limit of Row I of \texttt{Other\_terms} (relative difference $2.5\times10^{-2}$, $2.4\times10^{-4}$, $2.4\times10^{-6}$ for $\frac{k^{+}}{\vee}=10^{-2},10^{-4},10^{-6}$); dropped: $\frac{k^{+}}{\vee}$ in the logarithm and the $\log\big(1-\frac{k^{+}}{\vee}\big)$ term. Equal and opposite to Group I Row IV (note there)]\end{notebox}
\needspace{3\baselineskip}\noindent\rule{\linewidth}{0.3pt}\par\noindent\textbf{Lines 194} \quad \textsf{Group III (with c.c.) $\rangle$ Row II}\par
\begin{notebox}\textsf{\small Note, l.194:} [Checked: correct limit of Row II of \texttt{Other\_terms} (relative difference $2.5\times10^{-2}$, $2.4\times10^{-4}$, $2.4\times10^{-6}$); dropped: $\frac{k^{+}}{\vee}$ in the logarithm and the $\log\big(1-\frac{k^{+}}{\vee}\big)$ term (general remark)]\end{notebox}
\needspace{3\baselineskip}\noindent\rule{\linewidth}{0.3pt}\par\noindent\textbf{Lines 198--199} \quad \textsf{Group III (with c.c.) $\rangle$ Row III} \quad {\small (in \texttt{align}, lines 197--199)}\par
\begin{notebox}\textsf{\small Flag, l.198:} \fbox{FLAG 6}\end{notebox}
\begin{notebox}\textsf{\small Note, l.199:} [Checked: this is \emph{not} the first part of Row III of \texttt{Other\_terms} in the limit $\vee\to\infty$. It is the $p^{+}$ integral of the checked Group III row with $\frac{1}{p^{+}}\to\big[\frac{1}{p^{+}}\big]_{+}$, lower limit $0$, upper limit $\vee-k^{+}$: the other split. Prefactor, kernels, phases and colour structure are those of the checked row]\end{notebox}
\begin{notebox}\textsf{\small Flag, l.199:} \fbox{FLAG 6} $\vee$ is still there, and $\vee\to\infty$ cannot be taken in this form. Exactly, [this form] $=$ [form of \texttt{Other\_terms} without its $\log\big(1-\frac{k^{+}}{\vee}\big)$ term] $+\log\frac{\vee}{k^{+}}\times$[second $+$ third term of this row in the BFKL list]. Numerically (test function in place of the first colour bracket, $p^{+}$ and $y'$ integrals by quadrature, prefactor $1$), for $\frac{\vee}{k^{+}}=3,10,30,100$: this form $=0.455-0.328i$, $1.035-0.193i$, $1.535+0.029i$, $2.084+0.296i$; form of \texttt{Other\_terms} $=-0.047-0.581i$, $-0.017-0.724i$, $-0.018-0.756i$, $-0.019-0.766i$ (converges); difference $-\log\frac{\vee}{k^{+}}\times[\dots]$: $10^{-12}$. Origin: (a) $-\big[\frac{1}{p^{+}}\big]_{+}\delta_{jm}\delta_{ik'}$: for $p^{+}\to\infty$, $-\frac{f(p^{+})-f(0)}{p^{+}}\to+\frac{f(0)}{p^{+}}$, which gives $+\log\frac{\vee}{k^{+}}f(0)$ (second BFKL term); (b) $-\frac{1}{k^{+}}\delta_{im}\delta_{jk'}$: constant kernel times the phase; at fixed coordinates $\int_{0}^{\vee-k^{+}}dp^{+}e^{-i\frac{p^{+}}{k^{+}}k\cdot(z-y')}$ oscillates without a limit, and after the $y'$ integral the integrand is $\propto\frac{k^{+}}{p^{+}+k^{+}}$, which gives $\log\frac{\vee}{k^{+}}\times$(third BFKL term). So this form belongs to a BFKL list with only the first term of the row; with the three terms of the BFKL list the left-over is the form of \texttt{Other\_terms}, and at $\vee=\infty$ it is $N_{1}^{\widehat{\mathcal{E}}}+N_{1}^{\Psi}+N_{1}^{\rm sub}$ of Sec.~3.3.1 of \texttt{Finite\_terms\_final} (subtraction at $y'\to z$ under the same $p^{+}$ integral, $y'$ integral first). The soft subtraction of $\big[\frac{1}{p^{+}}\big]_{+}$ has a $\frac{1}{z^{2}}\log$ tail as in FLAG 4 (OPEN II/III-2 there).\end{notebox}
\needspace{3\baselineskip}\noindent\rule{\linewidth}{0.3pt}\par\noindent\textbf{Lines 202--206} \quad {\small (in \texttt{align}, lines 201--204)}\par
\begin{notebox}\textsf{\small Flag, l.202:} \fbox{FLAG 7}\end{notebox}
\begin{notebox}\textsf{\small Note, l.204:} [Checked: as the first part: the checked Group III row with $\frac{1}{p^{+}}\to\big[\frac{1}{p^{+}}\big]_{+}$, not the second part of Row III of \texttt{Other\_terms}]\end{notebox}
\begin{notebox}\textsf{\small Flag, l.204:} \fbox{FLAG 7} As FLAG 6. [This form] $=$ [form of \texttt{Other\_terms} without its $\log\big(1-\frac{k^{+}}{\vee}\big)$ term] $+\log\frac{\vee}{k^{+}}\times$[second $+$ third term of the second part in the BFKL list] (residual $10^{-12}$): this form $=-0.057-0.030i$, $-0.071-0.012i$, $-0.068+0.030i$, $-0.055+0.085i$ for $\frac{\vee}{k^{+}}=3,10,30,100$; form of \texttt{Other\_terms} $=-0.075-0.084i$, $-0.109-0.126i$, $-0.124-0.138i$, $-0.131-0.143i$. The logarithm comes from $-\big[\frac{1}{p^{+}}\big]_{+}[\dots]$ and from $-\frac{p^{+}}{k^{+\,2}}[\dots]$, whose kernel $\frac{p^{+}}{p^{+}(x-z)^{2}+k^{+}(x-w)^{2}}$ tends to the constant $\frac{1}{(x-z)^{2}}$ at $p^{+}\to\infty$. Finite form: $N_{2}^{\rm reg}+N_{2}^{\Psi}+N_{2}^{\rm sub}+N_{2}^{\log}$ of Sec.~3.3.2 of \texttt{Finite\_terms\_final}.\end{notebox}
\begin{notebox}\textsf{\small Note, l.206:} [As in \texttt{Other\_terms}]\end{notebox}
\needspace{3\baselineskip}\noindent\rule{\linewidth}{0.3pt}\par\noindent\textbf{Lines 209--210} \quad \textsf{Group III (with c.c.) $\rangle$ Row V} \quad {\small (in \texttt{align}, lines 208--210)}\par
\begin{notebox}\textsf{\small Flag, l.209:} \fbox{FLAG 8}\end{notebox}
\begin{notebox}\textsf{\small Note, l.210:} [Checked: the checked Group III Row V with $\frac{1}{p^{+}}\to\big[\frac{1}{p^{+}}\big]_{+}$; not Row V of \texttt{Other\_terms}]\end{notebox}
\begin{notebox}\textsf{\small Flag, l.210:} \fbox{FLAG 8} As FLAG 6 (same kernels and $p^{+}$ dependence as Row III, first part; prefactor $\times2$, other colour structure): $\vee\to\infty$ cannot be taken; [this form] $-$ [form of \texttt{Other\_terms}] $=\log\frac{\vee}{k^{+}}\times$[second $+$ third term of Row V in the BFKL list]. Finite form: $N^{\widehat{\mathcal{E}}}+N^{\Psi}+N^{\rm sub}$ of Sec.~3.5 of \texttt{Finite\_terms\_final}.\end{notebox}
\needspace{3\baselineskip}\noindent\rule{\linewidth}{0.3pt}\par\noindent\textbf{Lines 213--214} \quad \textsf{Group III (with c.c.) $\rangle$ Row VI} \quad {\small (in \texttt{align}, lines 212--214)}\par
\begin{notebox}\textsf{\small Flag, l.213:} \fbox{FLAG 9}\end{notebox}
\begin{notebox}\textsf{\small Note, l.214:} [Checked: the checked Group III Row VI with $\frac{1}{p^{+}}\to\big[\frac{1}{p^{+}}\big]_{+}$; not Row VI of \texttt{Other\_terms}]\end{notebox}
\begin{notebox}\textsf{\small Flag, l.214:} \fbox{FLAG 9} As FLAG 6 and FLAG 8: $\vee\to\infty$ cannot be taken; [this form] $-$ [form of \texttt{Other\_terms}] $=\log\frac{\vee}{k^{+}}\times$[second $+$ third term of Row VI in the BFKL list]. Finite form: $N^{\widehat{\mathcal{E}}}+N^{\Psi}+N^{\rm sub}$ of Sec.~3.6 of \texttt{Finite\_terms\_final}.\end{notebox}
\needspace{15\baselineskip}\bigskip
\subsection{Limiting cases (decisions implemented)}\label{sec:file-limiting_final}
\noindent\textbf{Original:} \texttt{/root/.claude/uploads/0ab97fcf-94a2-5aec-8bad-586400481685/f164e8ee-Limiting\_cases.tex}\\\textbf{Checked copy:} \texttt{/home/user/correlations/nlo\_review/limits\_checked/Limiting\_cases\_final.tex}\\
Lines: 215 (original), 214 (checked copy). Coloured groups of the check: 96 (96 top-level); author\textquotesingle s own coloured groups left unchanged: 0.\\
Records by type: phase convention 11; limits/range 1; missing term added 20; term removed/replaced 17; typo/LaTeX 16; note only 27.
\par\noindent Unmarked layout-only edits (line breaks \texttt{\textbackslash notag\textbackslash\textbackslash\ \&}, \texttt{\textbackslash par\textbackslash noindent} around inserted notes, spacing), checked-copy lines: 36, 70, 71, 78, 85, 91, 95, 107, 122, 153, 158, 160, 183, 184, 186, 191, 194, 197, 203, 204, 206, 210.
\par\noindent Unmarked edits that are not layout only (listed with their block below): 19.
\subsubsection*{Equations replaced as a whole}
In this file several equations of the author were removed and replaced (the orange notes in the checked copy name the removed lines of the author\textquotesingle s file). They are shown here in full: first the author\textquotesingle s equations (lines of the original), then the equations of the checked copy (red text accepted). The individual coloured groups inside these ranges are also listed in the blocks below and in the CSV file.
\paragraph{Group I, Row III: the four separate equations $N_{1}$, $N_{2}$, $N_{3a}$, $N_{3b}$ replaced by their sum (combined kernel $\mathcal{K}^{ij}$) with the last term of $N_{3b}$ in its recoil form.}
\noindent\textbf{Original, lines 58--107:}\par\noindent{\small We write it in three parts:
\begin{align*}
   \frac{d^{3}N_{1\,\rm NLO}^{\rm}}{d^{3}k}  & =  \frac{1}{(2\pi)^{3}}  \int_{w, w^\prime} e^{-ik(w-w^\prime)}\, \Bigg\{\notag \\
& \frac{ig^{4}\,f^{abd}}{16\pi^{5}}\, \frac{1}{k^{+}}\,\int_{x,y,z,x'}\, \frac{(x'-w')\cdot (y-z)}{(x'-w')^{2}\,(y-z)^{2}}\, \frac{\left[(x-w)^{2}+(x-z)^{2}\right]}{\left[(x-z)^{2}-(x-w)^{2}\right]^{2}}\, \times\, \log{\left[\frac{(x-z)^{2}}{(x-z)^{2}-(x-w)^{2}}\right]}\notag \\ & \left[U^{ab'}(x')\,\rho^{b'}(x')-U^{ab'}(w')\, \rho^{b'}(x')\right]\times\, \left[U^{de}(x)\,U^{bc}(z)\, \rho^{e}(x)\,\rho^{c}(y)-U^{de}(x)\,U^{bc}(y)\,\rho^{e}(x)\,\rho^{c}(y)\right] \notag \\
&  +\frac{ig^{4}\,f^{abd}}{16\pi^{5}}\, \frac{1}{k^{+}}\,\int_{x,y,z,x'}\, \frac{(x'-w')\cdot (y-z)}{(x'-w')^{2}\,(y-z)^{2}}\, \frac{\left[(x-w)^{2}+(x-z)^{2}\right]}{\left[(x-z)^{2}-(x-w)^{2}\right]^{2}}\, \times\, \log{\left[\frac{(x-w)^{2}-(x-z)^{2}}{(x-w)^{2}}\right]}\notag \\ & \left[U^{ab'}(x')\,\rho^{b'}(x')-U^{ab'}(w')\, \rho^{b'}(x')\right]\times\, \left[U^{de}(x)\,U^{bc}(z)\, \rho^{e}(x)\,\rho^{c}(y)-U^{de}(x)\,U^{bc}(y)\,\rho^{e}(x)\,\rho^{c}(y)\right] \notag \\
&-\, \frac{ig^{4}\,f^{abd}}{8\pi^{5}}\,\frac{1}{k^{+}}\,\int_{x,y,z,x'}\frac{(x'-w')\cdot (x-w)\, (x-z)\cdot (y-z)}{(x'-w')^{2}\, (x-w)^{2}\, (x-z)^{2}\, (y-z)^{2}}\,\log{\left[\frac{(x-z)^{2}}{(x-w)^{2}}\right]} \notag \\ & \left[U^{ab'}(x')\,\rho^{b'}(x')-U^{ab'}(w')\, \rho^{b'}(x')\right]\times\, \left[U^{de}(x)\,U^{bc}(z)\, \rho^{e}(x)\,\rho^{c}(y)-U^{de}(x)\,U^{bc}(y)\,\rho^{e}(x)\,\rho^{c}(y)\right]\Bigg\}
\end{align*}
\begin{align*}
   \frac{d^{3}N_{2\,\rm NLO}^{\rm}}{d^{3}k}  & =  -\,\frac{1}{(2\pi)^{3}}  \int_{w, w^\prime} e^{-ik(w-w^\prime)}\, \Bigg\{\notag \\ &   + \frac{ig^{4}\, f^{abd}}{8\pi^{5}}\, \frac{1}{k^{+}}\, \int_{x,y,z,x'}\,\log{\left[\frac{(x-z)^{2}}{(x-w)^{2}-(x-z)^{2}}\right]}\,\frac{(x'-w')\cdot (z-w)\, (y-z)\cdot (x-w)}{(x'-w')^{2}\, (y-z)^{2}\, (z-w)^{2}\,\left[(x-w)^{2}-(x-z)^{2}\right]}\, \notag \\ & \left[U^{ab'}(x')\,\rho^{b'}(x')-U^{ab'}(w')\, \rho^{b'}(x')\right]\times\,\left[U^{de}(x)\,U^{bc}(z)\, \rho^{e}(x)\,\rho^{c}(y)-U^{de}(x)\,U^{bc}(y)\,\rho^{e}(x)\,\rho^{c}(y)\right] \notag \\
& -\,\frac{ig^{4}\, f^{abd}}{8\pi^{5}}\, \frac{1}{k^{+}}\, \int_{x,y,z,x'}\,\log{\left[\frac{(x-z)^{2}}{(x-w)^{2}-(x-z)^{2}}\right]}\,\frac{(x'-w')\cdot (x-w)\, (y-z)\cdot(z-w)\,(x-z)^{2}}{(x'-w')^{2}\, (y-z)^{2}\, (z-w)^{2}\,\left[(x-w)^{2}-(x-z)^{2}\right]\,(x-w)^{2}}\, \notag \\ & \left[U^{ab'}(x')\,\rho^{b'}(x')-U^{ab'}(w')\, \rho^{b'}(x')\right]\times\,\left[U^{de}(x)\,U^{bc}(z)\, \rho^{e}(x)\,\rho^{c}(y)-U^{de}(x)\,U^{bc}(y)\,\rho^{e}(x)\,\rho^{c}(y)\right] \notag \\
& +\frac{ig^{4}\, f^{abd}}{8\pi^{5}}\, \frac{1}{k^{+}}\, \int_{x,y,z,x'}\,\log{\left[\frac{(x-z)^{2}}{(x-w)^{2}-(x-z)^{2}}\right]}\,\frac{(x'-w')\cdot  (y-z)\, (z-w)\cdot (x-w)\, (x-z)^{2}}{(x'-w')^{2}\, (y-z)^{2}\, (z-w)^{2}\,\left[(x-w)^{2}-(x-z)^{2}\right]^{2}}\, \notag \\ & \left[U^{ab'}(x')\,\rho^{b'}(x')-U^{ab'}(w')\, \rho^{b'}(x')\right]\times\,\left[U^{de}(x)\,U^{bc}(z)\, \rho^{e}(x)\,\rho^{c}(y)-U^{de}(x)\,U^{bc}(y)\,\rho^{e}(x)\,\rho^{c}(y)\right] \notag \\
& - \frac{ig^{4}\, f^{abd}}{8\pi^{5}}\, \frac{1}{k^{+}}\, \int_{x,y,z,x'}\,\log{\left[\frac{(x-z)^{2}}{(x-w)^{2}-(x-z)^{2}}\right]}\,\frac{(x'-w')\cdot (z-w)\, (y-z)\cdot (x-z)\, (x-w)^{2}}{(x'-w')^{2}\, (y-z)^{2}\, (z-w)^{2}\,\left[(x-w)^{2} - (x-z)^{2}\right]\, (x-z)^{2}}\, \notag \\ & \left[U^{ab'}(x')\,\rho^{b'}(x')-U^{ab'}(w')\, \rho^{b'}(x')\right]\times\,\left[U^{de}(x)\,U^{bc}(z)\, \rho^{e}(x)\,\rho^{c}(y)-U^{de}(x)\,U^{bc}(y)\,\rho^{e}(x)\,\rho^{c}(y)\right] \notag \\
& +\frac{ig^{4}\, f^{abd}}{8\pi^{5}}\, \frac{1}{k^{+}}\, \int_{x,y,z,x'}\,\log{\left[\frac{(x-z)^{2}}{(x-w)^{2}-(x-z)^{2}}\right]}\,\frac{(x'-w')\cdot (x-z)\, (y-z)\cdot (z-w)}{(x'-w')^{2}\, (y-z)^{2}\, (z-w)^{2}\,\left[(x-w)^{2} -(x-z)^{2}\right]}\, \notag \\ & \left[U^{ab'}(x')\,\rho^{b'}(x')-U^{ab'}(w')\, \rho^{b'}(x')\right]\times\,\left[U^{de}(x)\,U^{bc}(z)\, \rho^{e}(x)\,\rho^{c}(y)-U^{de}(x)\,U^{bc}(y)\,\rho^{e}(x)\,\rho^{c}(y)\right] \notag \\
& -\frac{ig^{4}\, f^{abd}}{8\pi^{5}}\, \frac{1}{k^{+}}\, \int_{x,y,z,x'}\,\log{\left[\frac{(x-z)^{2}}{(x-w)^{2}-(x-z)^{2}}\right]}\,\frac{(x'-w')\cdot (y-z)\, (z-w)\cdot (x-z)\,(x-w)^{2}}{(x'-w')^{2}\, (y-z)^{2}\, (z-w)^{2}\,\left[(x-w)^{2}-(x-z)^{2}\right]^{2}}\, \notag \\ & \left[U^{ab'}(x')\,\rho^{b'}(x')-U^{ab'}(w')\, \rho^{b'}(x')\right]\times\,\left[U^{de}(x)\,U^{bc}(z)\, \rho^{e}(x)\,\rho^{c}(y)-U^{de}(x)\,U^{bc}(y)\,\rho^{e}(x)\,\rho^{c}(y)\right]\Bigg\}
\end{align*}
\begin{align*}
\frac{d^{3}N_{3a\,\rm NLO}}{d^{3}k} & = \frac{1}{(2\pi)^{3}}\int_{w,w'}e^{-ik(w-w')}\Bigg\{\notag \\
& -\,\frac{ig^{4}f^{abd}}{8\pi^{5}}\,\frac{1}{k^{+}}\,\int_{x,y,z,x'}\log{\left[\frac{(x-w)^{2}-(x-z)^{2}}{(x-w)^{2}}\right]}\,
\frac{(x'-w')\cdot(y-z)\,(z-w)\cdot(x-w)\,(x-z)^{2}}{(x'-w')^{2}(y-z)^{2}(z-w)^{2}\,\left[(x-w)^{2}-(x-z)^{2}\right]^{2}}\notag \\ & \left[U^{ab'}(x')\,\rho^{b'}(x')-U^{ab'}(w')\, \rho^{b'}(x')\right]\times \left[U^{de}(x)\,U^{bc}(z)\, \rho^{e}(x)\,\rho^{c}(y)-U^{de}(x)\,U^{bc}(y)\,\rho^{e}(x)\,\rho^{c}(y)\right] \notag \\
& +\frac{ig^{4}f^{abd}}{8\pi^{5}}\,\frac{1}{k^{+}}\,\int_{x,y,z,x'}\log{\left[\frac{(x-w)^{2}-(x-z)^{2}}{(x-w)^{2}}\right]}\,
\frac{(x'-w')\cdot(y-z)\,(z-w)\cdot(x-z)\,(x-w)^{2}}{(x'-w')^{2}(y-z)^{2}(z-w)^{2}\,\left[(x-w)^{2}-(x-z)^{2}\right]^{2}}\notag \\ 
&\left[U^{ab'}(x')\,\rho^{b'}(x')-U^{ab'}(w')\, \rho^{b'}(x')\right]\times \left[U^{de}(x)\,U^{bc}(z)\, \rho^{e}(x)\,\rho^{c}(y)-U^{de}(x)\,U^{bc}(y)\,\rho^{e}(x)\,\rho^{c}(y)\right] \notag \\
& -\frac{ig^{4}f^{abd}}{8\pi^{5}}\,\frac{1}{k^{+}}\,\int_{x,y,z,x'}\log{\left[\frac{(x-w)^{2}-(x-z)^{2}}{(x-w)^{2}}\right]}\,
\frac{(x'-w')\cdot(z-w)\,(y-z)\cdot(x-w)}{(x'-w')^{2}(y-z)^{2}(z-w)^{2}\,\left[(x-w)^{2}-(x-z)^{2}\right]}\notag \\ & \left[U^{ab'}(x')\,\rho^{b'}(x')-U^{ab'}(w')\, \rho^{b'}(x')\right]\times \left[U^{de}(x)\,U^{bc}(z)\, \rho^{e}(x)\,\rho^{c}(y)-U^{de}(x)\,U^{bc}(y)\,\rho^{e}(x)\,\rho^{c}(y)\right] \notag \\
&  +\frac{ig^{4}f^{abd}}{8\pi^{5}}\,\frac{1}{k^{+}}\,\int_{x,y,z,x'}\log{\left[\frac{(x-w)^{2}-(x-z)^{2}}{(x-w)^{2}}\right]}\,
\frac{(x'-w')\cdot(z-w)\,(y-z)\cdot(x-z)\,(x-w)^{2}}{(x'-w')^{2}(y-z)^{2}(z-w)^{2}\,\left[(x-w)^{2}-(x-z)^{2}\right](x-z)^{2}}\notag \\ & \left[U^{ab'}(x')\,\rho^{b'}(x')-U^{ab'}(w')\, \rho^{b'}(x')\right]\times \left[U^{de}(x)\,U^{bc}(z)\, \rho^{e}(x)\,\rho^{c}(y)-U^{de}(x)\,U^{bc}(y)\,\rho^{e}(x)\,\rho^{c}(y)\right] \notag \\
& +\frac{ig^{4}f^{abd}}{8\pi^{5}}\,\frac{1}{k^{+}}\,\int_{x,y,z,x'}\log{\left[\frac{(x-w)^{2}-(x-z)^{2}}{(x-w)^{2}}\right]}\,
\frac{(y-z)\cdot(z-w)\,(x'-w')\cdot(x-w)\,(x-z)^{2}}{(x'-w')^{2}(y-z)^{2}(z-w)^{2}\,\left[(x-w)^{2}-(x-z)^{2}\right](x-w)^{2}}\notag \\ 
& \left[U^{ab'}(x')\,\rho^{b'}(x')-U^{ab'}(w')\, \rho^{b'}(x')\right]\times \left[U^{de}(x)\,U^{bc}(z)\, \rho^{e}(x)\,\rho^{c}(y)-U^{de}(x)\,U^{bc}(y)\,\rho^{e}(x)\,\rho^{c}(y)\right] \notag \\
& -\frac{ig^{4}f^{abd}}{8\pi^{5}}\,\frac{1}{k^{+}}\,\int_{x,y,z,x'}\log{\left[\frac{(x-w)^{2}-(x-z)^{2}}{(x-w)^{2}}\right]}\,
\frac{(y-z)\cdot(z-w)\,(x'-w')\cdot(x-z)}{(x'-w')^{2}(y-z)^{2}(z-w)^{2}\,\left[(x-w)^{2}-(x-z)^{2}\right]}\notag \\ & \left[U^{ab'}(x')\,\rho^{b'}(x')-U^{ab'}(w')\, \rho^{b'}(x')\right]\times \left[U^{de}(x)\,U^{bc}(z)\, \rho^{e}(x)\,\rho^{c}(y)-U^{de}(x)\,U^{bc}(y)\,\rho^{e}(x)\,\rho^{c}(y)\right] \notag \\
& +\frac{ig^{4}f^{abd}}{8\pi^{5}}\,\frac{1}{k^{+}}\,\int_{x,y,z,x'}\mathcal{A}_{0}\,
\frac{(x'-w')\cdot(z-w)\,(y-z)\cdot(x-w)}{(x'-w')^{2}(y-z)^{2}(z-w)^{2}}\notag \\
& \left[U^{ab'}(x')\,\rho^{b'}(x')-U^{ab'}(w')\, \rho^{b'}(x')\right]\times \left[U^{de}(x)\,U^{bc}(z)\, \rho^{e}(x)\,\rho^{c}(y)-U^{de}(x)\,U^{bc}(y)\,\rho^{e}(x)\,\rho^{c}(y)\right] \notag \\
& +\frac{ig^{4}f^{abd}}{8\pi^{5}}\,\frac{1}{k^{+}}\,\int_{x,y,z,x'}\mathcal{A}_{0}\,
\frac{(y-z)\cdot(z-w)\,(x'-w')\cdot(x-z)}{(x'-w')^{2}(y-z)^{2}(z-w)^{2}}\notag \\ & \left[U^{ab'}(x')\,\rho^{b'}(x')-U^{ab'}(w')\, \rho^{b'}(x')\right]\times \left[U^{de}(x)\,U^{bc}(z)\, \rho^{e}(x)\,\rho^{c}(y)-U^{de}(x)\,U^{bc}(y)\,\rho^{e}(x)\,\rho^{c}(y)\right] \notag \\
& -\frac{ig^{4}f^{abd}}{8\pi^{5}}\frac{1}{k^{+}}\,\int_{x,y,z,x'}\mathcal{A}_{0}\,
\frac{(x-w)\cdot(x-z)\,(x'-w')\cdot(z-w)\,(y-z)\cdot(x-z)}{(x'-w')^{2}(y-z)^{2}(z-w)^{2}(x-z)^{2}}\notag \\ & \left[U^{ab'}(x')\,\rho^{b'}(x')-U^{ab'}(w')\, \rho^{b'}(x')\right]\times \left[U^{de}(x)\,U^{bc}(z)\, \rho^{e}(x)\,\rho^{c}(y)-U^{de}(x)\,U^{bc}(y)\,\rho^{e}(x)\,\rho^{c}(y)\right]\Bigg\}
\end{align*}
\begin{align*}
\frac{d^{3}N_{3b\,\rm NLO}}{d^{3}k} & = \frac{1}{(2\pi)^{3}}\int_{w,w'}e^{-ik(w-w')}\,\Bigg\{\notag \\
&-\frac{ig^{4}f^{abd}}{8\pi^{5}}\,\frac{1}{k^{+}}\,\int_{x,y,z,x'}\mathcal{A}_{0}\,
\frac{(x-w)\cdot(x-z)\,(y-z)\cdot(z-w)\,(x'-w')\cdot(x-w)}{(x'-w')^{2}(y-z)^{2}(z-w)^{2}(x-w)^{2}}\notag \\ & \left[U^{ab'}(x')\,\rho^{b'}(x')-U^{ab'}(w')\, \rho^{b'}(x')\right]\times \left[U^{de}(x)\,U^{bc}(z)\, \rho^{e}(x)\,\rho^{c}(y)-U^{de}(x)\,U^{bc}(y)\,\rho^{e}(x)\,\rho^{c}(y)\right] \notag \\
& -\frac{ig^{4}f^{abd}}{16\pi^{5}}\,\frac{1}{k^{+}}\,\int_{x,y,z,x'}\log{\left[\frac{(z-w)^{2}}{(x-w)^{2}}\right]}\,
\frac{(x'-w')\cdot(z-w)\,(y-z)\cdot(x-z)}{(x'-w')^{2}(y-z)^{2}(z-w)^{2}(x-z)^{2}}\notag \\
&\left[U^{ab'}(x')\,\rho^{b'}(x')-U^{ab'}(w')\, \rho^{b'}(x')\right]\times \left[U^{de}(x)\,U^{bc}(z)\, \rho^{e}(x)\,\rho^{c}(y)-U^{de}(x)\,U^{bc}(y)\,\rho^{e}(x)\,\rho^{c}(y)\right] \notag \\
&+ \frac{ig^{4}f^{abd}}{16\pi^{5}}\,\frac{1}{k^{+}}\,\int_{x,y,z,x'}\log{\left[\frac{(z-w)^{2}}{(x-w)^{2}}\right]}\,
\frac{(y-z)\cdot(z-w)\,(x'-w')\cdot(x-w)}{(x'-w')^{2}(y-z)^{2}(z-w)^{2}(x-w)^{2}}\notag \\ 
& \left[U^{ab'}(x')\,\rho^{b'}(x')-U^{ab'}(w')\, \rho^{b'}(x')\right]\times \left[U^{de}(x)\,U^{bc}(z)\, \rho^{e}(x)\,\rho^{c}(y)-U^{de}(x)\,U^{bc}(y)\,\rho^{e}(x)\,\rho^{c}(y)\right]\notag \\ &  -\frac{ig^{4}f^{abd}}{16\pi^{5}}\,\frac{1}{k^{+}}\,\int_{x,y,z,x'}\log{\left[\frac{(z-w)^{2}}{(x-w)^{2}}\right]}\,\frac{(x'-w')\cdot(y-z)}{(x'-w')^{2}(y-z)^{2}\,(z-w)^{2}}\notag \\ & \left[U^{ab'}(x')\,\rho^{b'}(x')-U^{ab'}(w')\, \rho^{b'}(x')\right]\times \left[U^{de}(x)\,U^{bc}(z)\, \rho^{e}(x)\,\rho^{c}(y)-U^{de}(x)\,U^{bc}(y)\,\rho^{e}(x)\,\rho^{c}(y)\right]\Bigg\}
\end{align*}}\par
\noindent\textbf{Checked copy, lines 59--78:}\par\noindent{\small \begin{align*}
\frac{d^{3}N_{\rm NLO}}{d^{3}k} & = \frac{1}{(2\pi)^{3}}\int_{w,w'}e^{-ik\cdot(w-w')}\,\Bigg\{\notag \\
& \frac{ig^{4}f^{abd}}{8\pi^{5}}\,\frac{1}{k^{+}}\,\int_{x,y,z,x'}\frac{(x'-w')^{i}\,(y-z)^{j}}{(x'-w')^{2}\,(y-z)^{2}}\;\mathcal{K}^{ij}\notag \\ & \left[U^{ab'}(x')\,\rho^{b'}(x')-U^{ab'}(w')\, \rho^{b'}(x')\right]\times \left[U^{de}(x)\,U^{bc}(z)\, \rho^{e}(x)\,\rho^{c}(y)-U^{de}(x)\,U^{bc}(y)\,\rho^{e}(x)\,\rho^{c}(y)\right]\notag \\
&+ \frac{ig^{4}f^{abd}}{16\pi^{5}}\,\frac{1}{k^{+}}\,\int_{x,y,z,x'}\log{\left[\frac{(z-w)^{2}}{(x-w)^{2}}\right]}\,
\frac{(y-z)\cdot(z-w)\,(x'-w')\cdot(x-w)}{(x'-w')^{2}(y-z)^{2}(z-w)^{2}(x-w)^{2}}\notag \\ 
& \left[U^{ab'}(x')\,\rho^{b'}(x')-U^{ab'}(w')\, \rho^{b'}(x')\right]\times \left[U^{de}(x)\,U^{bc}(z)\, \rho^{e}(x)\,\rho^{c}(y)-U^{de}(x)\,U^{bc}(y)\,\rho^{e}(x)\,\rho^{c}(y)\right]\notag \\
& -\frac{ig^{4}f^{abd}}{8\pi^{5}}\,\frac{1}{k^{+}}\,\int_{x,y,z,x'}\mathcal{E}\big(k\cdot(z-w)\big)\,\frac{(x-z)\cdot(w-z)\,(x'-w')\cdot(y-w)}{(x-z)^{2}\,(w-z)^{2}\,(x'-w')^{2}\,(y-w)^{2}}\notag \\ & \left[U^{ab'}(x')\,\rho^{b'}(x')-U^{ab'}(w')\, \rho^{b'}(x')\right]\times \left[U^{de}(x)\,U^{bc}(w)\, \rho^{e}(x)\,\rho^{c}(y)-U^{de}(x)\,U^{bc}(y)\,\rho^{e}(x)\,\rho^{c}(y)\right]\Bigg\}
\end{align*}
with
\begin{align*}
\mathcal{K}^{ij} & = \log\left[\frac{(x-z)^{2}}{(x-w)^{2}}\right]\Bigg\{\frac{\delta^{ij}}{2\,(z-w)^{2}}-\frac{(x-w)^{i}\,(x-z)^{j}}{(x-w)^{2}\,(x-z)^{2}}+\frac{1}{(z-w)^{2}\left[(x-w)^{2}-(x-z)^{2}\right]}\Big[(x-w)^{i}(z-w)^{j}\,\frac{(x-z)^{2}}{(x-w)^{2}}\notag \\
& +(z-w)^{i}(x-z)^{j}\,\frac{(x-w)^{2}}{(x-z)^{2}}-(z-w)^{i}(x-w)^{j}-(x-z)^{i}(z-w)^{j}\Big]\Bigg\}\notag \\
& +\frac{\mathcal{A}_{0}}{(z-w)^{2}}\Bigg\{(z-w)^{i}(x-w)^{j}+(x-z)^{i}(z-w)^{j}-(x-w)\cdot(x-z)\left[\frac{(z-w)^{i}(x-z)^{j}}{(x-z)^{2}}+\frac{(x-w)^{i}(z-w)^{j}}{(x-w)^{2}}\right]\Bigg\}\notag \\
& -\frac{1}{2\,(z-w)^{2}}\log\left[\frac{(z-w)^{2}}{(x-w)^{2}}\right]\frac{(z-w)^{i}(x-z)^{j}}{(x-z)^{2}}\,,\notag \\
\mathcal{E}(s) & \equiv\int_{0}^{1}\frac{dt}{t}\,\big(e^{-ist}-1\big)=\mathrm{Ci}(|s|)-\gamma_{E}-\log|s|-i\,\mathrm{Si}(s)
\end{align*}
where: \begin{align*}
\mathcal{A}_{0} & = \frac{1}{\sqrt{(x-w)^{2}\,(x-z)^{2}-\left[(x-w)\cdot (x-z)\right]^{2}}}\, \Bigg[\mathrm{arc\,tan}\,\left(\frac{(x-z)^{2}-(x-w)\cdot (x-z)}{\sqrt{(x-w)^{2}\,(x-z)^{2}-\left[(x-w)\cdot (x-z)\right]^{2}}}\right)\notag \\ & + \mathrm{arc\,tan}\,\left(\frac{(x-w)\cdot (x-z)}{\sqrt{(x-w)^{2}\,(x-z)^{2}-\left[(x-w)\cdot (x-z)\right]^{2}}}\right) \Bigg]\notag \\ & = \frac{1}{\sqrt{(x-w)^{2}\,(x-z)^{2}-\left[(x-w)\cdot (x-z)\right]^{2}}}\, \Bigg[\mathrm{arc\,cos}\,\left(\frac{(x-w)^{2}-(x-w)\cdot (x-z)}{\sqrt{(x-w)^{2}\,(w-z)^{2}}}\right)\Bigg]
\end{align*} [\textbf{First term}: the 24 terms of $N_{1}$, $N_{2}$, $N_{3a}$ and the first two terms of $N_{3b}$, added (Eqs.~(1.5), (1.6) of \texttt{Finite\_terms\_final}). \textbf{Second term} (black): the third term of $N_{3b}$, unchanged. Verified at random points with explicit SU(3) matrices: [$N_{1}+N_{2}+N_{3a}+$ first three terms of $N_{3b}$ of the author's file] $-$ [first $+$ second term here] $=$ $2\times10^{-22}$ to $7\times10^{-18}$ against $7\times10^{-9}$ to $5\times10^{-6}$ at eight points (relative $5\times10^{-16}$ to $1.5\times10^{-12}$). The separate equations at distance $\delta=10^{-3},10^{-4},10^{-5}$ from the line $(x-z)^{2}=(x-w)^{2}$: $|N_{2}|\simeq|N_{3a}|=0.89,\ 1.2\times10^{2},\ 1.6\times10^{4}$, $|N_{1}|\simeq|N_{2}+N_{3a}|=3.0\times10^{-4},\ 3.0\times10^{-3},\ 3.0\times10^{-2}$, sum of the four $=6.0\times10^{-7}$ at all three. \textbf{Third term}: it replaces the last term of $N_{3b}$ (orange sentence above). That term is the limit $\Lambda\to0$ at fixed $z\neq w$ only; with the kernel $\frac{1}{(z-w)^{2}}$ and $U^{bc}(z)$, which nothing softens at $z=w$, its $z$ integral diverges like $\log^{2}$ at $z\to w$ ($|z-w|^{2}\,|{\rm term}|/\log\frac{1}{|z-w|}=5.1\times10^{-7},\ 4.6\times10^{-7},\ 4.5\times10^{-7}$ at $|z-w|=10^{-2},10^{-3},10^{-4}$, i.e.\ constant). Taken together with its $\log\frac{k^{+}}{\Lambda}$ partner in the row at finite $\Lambda$ and after the change of variable to the position of the soft gluon, it is the term of the JIMWLK list with $U^{bc}(w)$ (three-$\rho$ list, Row III) plus the third term here (third line of Eq.~(1.5) of \texttt{Finite\_terms\_final}, derivation Eqs.~(1.9)--(1.11) there; the identity was checked there with a Gaussian regulator, re-run now: both sides $-10.896402+0.822556\,i$, difference $3\times10^{-13}$). The JIMWLK term of this row is therefore the one of the evolution lists, with $U^{bc}(w)$. $\mathcal{E}$: closed form against quadrature, $10^{-16}$. \textbf{Integrability}: the first term is integrable only as the combined kernel (not term by term), and at large $z$ only conditionally (angular integral first). The second and the third term have a $\frac{1}{z^{2}}\log$ tail at large $z$: they are finite only in the sum over Group I Rows III, V and Group III Rows III, V, VI under one integral.]\par\noindent}\par
\paragraph{Group III, Row III, first part: the equation in the plain $[1/p^{+}]_{+}$ split replaced by the form with the large logarithm removed completely.}
\noindent\textbf{Original, lines 197--199:}\par\noindent{\small \begin{align*}
     \frac{d^{3}N_{1\,\rm NLO}}{d^{3}k} & = \frac{1}{(2\pi)^{3}}\, \frac{ig^{4}\,f^{c'd'a}}{8\pi^{4}}\, \frac{1}{k^{+}}\,\int_{0}^{\vee-k^{+}}\,  \frac{dp^{+}}{2\pi}\, \int_{w,x',y',z,y,x}\, e^{-ik\cdot (w-y')}\, e^{-ik\cdot \frac{p^{+}}{k^{+}}\,(z-y')}\,\frac{(y'-z)^{m}\,(x'-y')^{k'}\,(y-w)^{i}\,(x-z)^{j}}{(y'-z)^{2}\,(x'-y')^{2}\,(y-w)^{2}\,(x-z)^{2}}\notag \\ & \left[\frac{\delta_{k'm}\,\delta_{ij}}{(p^{+}+k^{+})}-\left[\frac{1}{p^{+}}\right]_{+}\,\delta_{jm}\,\delta_{ik'}-\frac{1}{k^{+}}\, \delta_{im}\, \delta_{jk'}\right]\, \left[U^{c'e'}(y')\,U^{d'b}(z)\,\rho^{e'}(x')-U^{d'b}(z)\,U^{c'e'}(x')\,\rho^{e'}(x')\right]\times\notag \\ &  \left[U^{ac}(w)\, \{\rho^{c}(y),\rho^{b}(x)\}-U^{cb}(z)\,U^{ad}(y)\,U^{ce}(x)\,\{\rho^{d}(y),\rho^{e}(x)\}\right]
\end{align*}}\par
\noindent\textbf{Checked copy, lines 181--186:}\par\noindent{\small \begin{align*}
     \frac{d^{3}N_{1\,\rm NLO}}{d^{3}k} & = \frac{1}{(2\pi)^{3}}\, \frac{ig^{4}\,f^{c'd'a}}{8\pi^{4}}\, \frac{1}{k^{+}}\,\int_{0}^{\infty}\,  \frac{dp^{+}}{2\pi}\, \int_{w,x',y',z,y,x}\, e^{-ik\cdot (w-y')}\, e^{-ik\cdot \frac{p^{+}}{k^{+}}\,(z-y')}\,\frac{(y'-z)^{m}\,(x'-y')^{k'}\,(y-w)^{i}\,(x-z)^{j}}{(y'-z)^{2}\,(x'-y')^{2}\,(y-w)^{2}\,(x-z)^{2}}\notag \\ & \left[\frac{\delta_{k'm}\,\delta_{ij}}{(p^{+}+k^{+})}-\left(\left[\frac{k^{+}}{p^{+}\,(p^{+}+k^{+})}\right]_{+}+\frac{1}{(p^{+}+k^{+})}\right)\,\delta_{jm}\,\delta_{ik'}\right]\, \left[U^{c'e'}(y')\,U^{d'b}(z)\,\rho^{e'}(x')-U^{d'b}(z)\,U^{c'e'}(x')\,\rho^{e'}(x')\right]\times\notag \\ &  \left[U^{ac}(w)\, \{\rho^{c}(y),\rho^{b}(x)\}-U^{cb}(z)\,U^{ad}(y)\,U^{ce}(x)\,\{\rho^{d}(y),\rho^{e}(x)\}\right] \notag \\
&  +\,\frac{1}{(2\pi)^{3}}\, \frac{ig^{4}\,f^{c'd'a}}{8\pi^{4}}\, \frac{1}{k^{+}}\,\int_{0}^{\infty}\,  \frac{dp^{+}}{2\pi}\, \int_{w,x',y',z,y,x}\, e^{-ik\cdot (w-y')}\, e^{-ik\cdot \frac{p^{+}}{k^{+}}\,(z-y')}\,\frac{(y'-z)^{m}\,(y-w)^{i}\,(x-z)^{j}}{(y'-z)^{2}\,(y-w)^{2}\,(x-z)^{2}}\notag \\ &   \left[-\frac{1}{k^{+}}\, \delta_{im}\, \delta_{jk'}\right]\, \Bigg[\frac{(x'-y')^{k'}}{(x'-y')^{2}}\,\left[U^{c'e'}(y')\,U^{d'b}(z)\,\rho^{e'}(x')-U^{d'b}(z)\,U^{c'e'}(x')\,\rho^{e'}(x')\right]\notag \\ & -\frac{(x'-z)^{k'}}{(x'-z)^{2}}\,\left[U^{c'e'}(z)\,U^{d'b}(z)\,\rho^{e'}(x')-U^{d'b}(z)\,U^{c'e'}(x')\,\rho^{e'}(x')\right]\Bigg]\times\notag \\ &   \left[U^{ac}(w)\, \{\rho^{c}(y),\rho^{b}(x)\}-U^{cb}(z)\,U^{ad}(y)\,U^{ce}(x)\,\{\rho^{d}(y),\rho^{e}(x)\}\right]
\end{align*}
[This is the first part of Row III of \texttt{Other\_terms} (v2) at $\vee\to\infty$: upper limit of the $p^{+}$ integrals $\infty$, its last term ($\log\big(1-\frac{k^{+}}{\vee}\big)\times\dots$) dropped; nothing else is changed. All three terms of this row in the BFKL list are removed from it (verified for the v2 form against the row with the cut-offs). \textbf{Condition for the limit}: in the term with the constant kernel (the one whose square bracket contains the same colour bracket at $y'$ and at $y'\to z$) the $y'$ integral has to be done before the $p^{+}$ integral, with the two terms of that bracket together. At fixed coordinates $\int_{0}^{P}dp^{+}e^{-i\frac{p^{+}}{k^{+}}s}=\frac{k^{+}}{is}\big(1-e^{-i\frac{P}{k^{+}}s}\big)$, $s=k\cdot(z-y')$, has no limit for $P\to\infty$; after the $y'$ integral the integrand is ${\cal O}(k^{+2}/p^{+2})$. The $y'$ integral of the subtracted piece converges only conditionally at large $y'$ (to be integrated symmetrically; Gaussian regulator in the check). \textbf{Numerical check} (smooth test function of $y'$ in place of the first colour bracket, second bracket $\to1$, prefactor $1$, $|k|=1$, $p^{+}$ and $y'$ integrals by quadrature): this form with the upper limit $\vee-k^{+}$, for $\frac{\vee}{k^{+}}=10,\,30,\,60,\,100$: $0.0201-1.0301\,i,\ 0.0389-1.1011\,i,\ 0.0419-1.1183\,i,\ 0.0430-1.1250\,i$; it converges like $\frac{k^{+}}{\vee}$: extrapolated with $a+b\,\frac{k^{+}}{\vee}$ from the last two values, limit $0.0445-1.1352\,i$, and $\frac{\vee}{k^{+}}\times|{\rm value}-{\rm limit}|=1.08,\ 1.04,\ 1.03$ at $\frac{\vee}{k^{+}}=10,30,60$. The removed form at the same points: $1.048-0.434\,i,\ 1.557-0.220\,i,\ 1.869-0.058\,i,\ 2.098+0.068\,i$ (grows like $\log\frac{\vee}{k^{+}}$). Second set of points, $\frac{\vee}{k^{+}}=10,\,30,\,60$: $-1.5656-0.3693\,i,\ -1.7886-0.4655\,i,\ -1.8452-0.4835\,i$ (successive differences $0.243,\ 0.059$, ratio $4.1$; $4$ for $\frac{k^{+}}{\vee}$). The subtracted term alone, with an analytic test function: $|J(\vee)-J(\infty)|=0.64$, $7.0\times10^{-2}$, $7.0\times10^{-3}$, $7.0\times10^{-4}$ for $\frac{\vee}{k^{+}}-1=10,\,10^{2},\,10^{3},\,10^{4}$, and $J(\infty)$ agrees with the form in which the $y'$ integral is done first to $7\times10^{-8}$. The plus prescription acts on the phase $e^{-ik\cdot\frac{p^{+}}{k^{+}}(z-y')}$. The soft subtraction terms of the plus prescriptions in this equation are finite only in the sum over Group I Rows III, V and Group III Rows III, V, VI under one integral.]\par\noindent}\par
\paragraph{Group III, Row III, second part: as for the first part.}
\noindent\textbf{Original, lines 201--204:}\par\noindent{\small \begin{align*}
 \frac{d^{3}N_{2\,\rm NLO}}{d^{3}k} & = \frac{1}{(2\pi)^{3}}\,\frac{g^{4}}{4\pi^{4}}\,\int_{0}^{\vee-k^{+}}\, \frac{dp^{+}}{2\pi}\,\int_{w,x',y',z,x}\, e^{-ik\cdot (w-y')}\,e^{-ik\cdot \frac{p^{+}}{k^{+}}\left(z-y'\right)}\, \frac{(y'-z)^{m}\,(x'-y')^{k'}\,}{(y'-z)^{2}\,(x'-y')^{2}\,}\, \frac{1}{\left[p^{+}\,(x-z)^{2}+k^{+}\,(x-w)^{2}\right]}\notag \\ & \Bigg[\frac{\delta_{k'm}\,d_{\perp}\,p^{+}}{2\,(z-w)^{2}\,(p^{+}+k^{+})^{2}}\, \left[(x-z)^{2}-(x-w)^{2}\right]+\frac{p^{+}\,\delta_{k'm}}{k^{+}\,(p^{+}+k^{+})}\,\bigg[\frac{(x-z)\cdot (z-w)}{(z-w)^{2}} -\frac{(x-z)\cdot (x-w)}{2\,(x-w)^{2}}\notag \\ & -\frac{\left[(x-z)^{2}-(x-w)^{2}\right]}{2\,(z-w)^{2}}\bigg] +\frac{\delta_{k'm}}{(p^{+}+k^{+})}\,\bigg[\frac{(x-w)\cdot (z-w)}{(z-w)^{2}}+\frac{(x-z)\cdot (x-w)}{2\,(x-z)^{2}} -\frac{\left[(x-z)^{2}-(x-w)^{2}\right]}{2\,(z-w)^{2}}\bigg]\notag \\ & -\left[\frac{1}{p^{+}}\right]_{+}\,\bigg[\frac{(x-w)^{k'}\,(z-w)^{m}}{(z-w)^{2}}+\frac{(x-z)^{m}\,(x-w)^{k'}}{2\,(x-z)^{2}}\bigg]-\frac{p^{+}}{k^{+\,2}}\,\bigg[\frac{(x-z)^{k'}\,(z-w)^{m}}{(z-w)^{2}} -\,\frac{(x-z)^{k'}\,(x-w)^{m}}{2\,(x-w)^{2}}\bigg]\notag \\ & -\frac{1}{k^{+}}\,\bigg[\frac{(x-z)^{m}\,(z-w)^{k'}}{(z-w)^{2}}-\frac{(x-z)^{m}\,(x-w)^{k'}}{2\,(x-w)^{2}}+\frac{(x-w)^{m}\,(z-w)^{k'}}{(z-w)^{2}} +\frac{(x-z)^{k'}\,(x-w)^{m}}{2\,(x-z)^{2}}\bigg]\Bigg]\times\,\notag \\ & \left[f^{c'd'a}\,U^{c'e'}(y')\,U^{d'b}(z)\,\rho^{e'}(x')-f^{c'd'a}\, 
    U^{d'b}(z)\,U^{c'e'}(x')\,\rho^{e'}(x')\right]\times \left[f^{cbd}\,U^{ac}(w)\,\rho^{d}(x) -f^{acd}\,U^{cb}(z)\,U^{de}(x)\,\rho^{e}(x)\right]
\end{align*}}\par
\noindent\textbf{Checked copy, lines 188--197:}\par\noindent{\small \begin{align*}
 \frac{d^{3}N_{2\,\rm NLO}}{d^{3}k} & = \frac{1}{(2\pi)^{3}}\,\frac{g^{4}}{4\pi^{4}}\,\int_{0}^{\infty}\, \frac{dp^{+}}{2\pi}\,\int_{w,x',y',z,x}\, e^{-ik\cdot (w-y')}\,e^{-ik\cdot \frac{p^{+}}{k^{+}}\left(z-y'\right)}\, \frac{(y'-z)^{m}\,(x'-y')^{k'}\,}{(y'-z)^{2}\,(x'-y')^{2}\,}\, \frac{1}{\left[p^{+}\,(x-z)^{2}+k^{+}\,(x-w)^{2}\right]}\notag \\ & \Bigg[\frac{\delta_{k'm}\,d_{\perp}\,p^{+}}{2\,(z-w)^{2}\,(p^{+}+k^{+})^{2}}\, \left[(x-z)^{2}-(x-w)^{2}\right]+\frac{p^{+}\,\delta_{k'm}}{k^{+}\,(p^{+}+k^{+})}\,\bigg[\frac{(x-z)\cdot (z-w)}{(z-w)^{2}} -\frac{(x-z)\cdot (x-w)}{2\,(x-w)^{2}}\notag \\ & -\frac{\left[(x-z)^{2}-(x-w)^{2}\right]}{2\,(z-w)^{2}}\bigg] +\frac{\delta_{k'm}}{(p^{+}+k^{+})}\,\bigg[\frac{(x-w)\cdot (z-w)}{(z-w)^{2}}+\frac{(x-z)\cdot (x-w)}{2\,(x-z)^{2}} -\frac{\left[(x-z)^{2}-(x-w)^{2}\right]}{2\,(z-w)^{2}}\bigg]\notag \\ & -\left(\left[\frac{k^{+}}{p^{+}\,(p^{+}+k^{+})}\right]_{+}+\frac{1}{(p^{+}+k^{+})}\right)\,\bigg[\frac{(x-w)^{k'}\,(z-w)^{m}}{(z-w)^{2}}+\frac{(x-z)^{m}\,(x-w)^{k'}}{2\,(x-z)^{2}}\bigg]\notag \\ & -\frac{1}{k^{+}}\,\bigg[\frac{(x-z)^{m}\,(z-w)^{k'}}{(z-w)^{2}}-\frac{(x-z)^{m}\,(x-w)^{k'}}{2\,(x-w)^{2}}+\frac{(x-w)^{m}\,(z-w)^{k'}}{(z-w)^{2}} +\frac{(x-z)^{k'}\,(x-w)^{m}}{2\,(x-z)^{2}}\bigg]\Bigg]\times\,\notag \\ & \left[f^{c'd'a}\,U^{c'e'}(y')\,U^{d'b}(z)\,\rho^{e'}(x')-f^{c'd'a}\, 
    U^{d'b}(z)\,U^{c'e'}(x')\,\rho^{e'}(x')\right]\times \left[f^{cbd}\,U^{ac}(w)\,\rho^{d}(x) -f^{acd}\,U^{cb}(z)\,U^{de}(x)\,\rho^{e}(x)\right] \notag \\ &   +\,\frac{1}{(2\pi)^{3}}\,\frac{g^{4}}{4\pi^{4}}\,\int_{0}^{\infty}\, \frac{dp^{+}}{2\pi}\,\int_{w,x',y',z,x}\, e^{-ik\cdot (w-y')}\,e^{-ik\cdot \frac{p^{+}}{k^{+}}\left(z-y'\right)}\, \frac{(y'-z)^{m}}{(y'-z)^{2}}\, \frac{1}{\left[p^{+}\,(x-z)^{2}+k^{+}\,(x-w)^{2}\right]}\notag \\ &   \Bigg[ -\frac{p^{+}}{k^{+\,2}}\,\bigg[\frac{(x-z)^{k'}\,(z-w)^{m}}{(z-w)^{2}} -\,\frac{(x-z)^{k'}\,(x-w)^{m}}{2\,(x-w)^{2}}\bigg] \Bigg]\times\,\notag \\ &  \Bigg[\frac{(x'-y')^{k'}}{(x'-y')^{2}}\,\left[f^{c'd'a}\,U^{c'e'}(y')\,U^{d'b}(z)\,\rho^{e'}(x')-f^{c'd'a}\, 
    U^{d'b}(z)\,U^{c'e'}(x')\,\rho^{e'}(x')\right]\notag \\ & -\frac{(x'-z)^{k'}}{(x'-z)^{2}}\,\left[f^{c'd'a}\,U^{c'e'}(z)\,U^{d'b}(z)\,\rho^{e'}(x')-f^{c'd'a}\, 
    U^{d'b}(z)\,U^{c'e'}(x')\,\rho^{e'}(x')\right]\Bigg]\times \notag \\ & \left[f^{cbd}\,U^{ac}(w)\,\rho^{d}(x) -f^{acd}\,U^{cb}(z)\,U^{de}(x)\,\rho^{e}(x)\right] \notag \\
&  +\,\frac{1}{(2\pi)^{3}}\,\frac{g^{4}}{4\pi^{4}}\,\int_{0}^{\infty}\, \frac{dp^{+}}{2\pi}\,\int_{w,x',w',z,x}\, e^{-ik\cdot (w-w')}\, \frac{(w'-z)^{m}\,(x'-z)^{k'}}{(w'-z)^{2}\,(x'-z)^{2}}\, \notag \\ &  \Bigg[ \frac{(x-w)^{2}}{(x-z)^{2}\,(p^{+}+k^{+})\,\left[p^{+}\,(x-z)^{2}+k^{+}\,(x-w)^{2}\right]}\,\bigg[\frac{(x-z)^{k'}\,(z-w)^{m}}{(z-w)^{2}} -\,\frac{(x-z)^{k'}\,(x-w)^{m}}{2\,(x-w)^{2}}\bigg] \Bigg]\times\,\notag \\ &   \left[f^{c'd'a}\,U^{c'e'}(z)\,U^{d'b}(z)\,\rho^{e'}(x')-f^{c'd'a}\, 
    U^{d'b}(z)\,U^{c'e'}(x')\,\rho^{e'}(x')\right]\times \left[f^{cbd}\,U^{ac}(w)\,\rho^{d}(x) -f^{acd}\,U^{cb}(z)\,U^{de}(x)\,\rho^{e}(x)\right]
\end{align*}
[This is the second part of Row III of \texttt{Other\_terms} (v2) at $\vee\to\infty$: upper limit $\infty$, last term ($\log\big(1-\frac{k^{+}}{\vee}\big)\times\dots$) dropped; nothing else is changed. The term $-\frac{p^{+}}{k^{+\,2}}[\dots]$, whose kernel tends to the constant $-\frac{1}{k^{+\,2}(x-z)^{2}}$ at large $p^{+}$, is the one that is subtracted at $y'\to z$ (second term); the third term is what remains of the piece added back, and its $p^{+}$ integral can be done: $\int_{0}^{\infty}\frac{dp^{+}}{2\pi}\frac{(x-w)^{2}}{(x-z)^{2}(p^{+}+k^{+})[p^{+}(x-z)^{2}+k^{+}(x-w)^{2}]}=\frac{1}{2\pi k^{+}}\frac{(x-w)^{2}}{(x-z)^{2}[(x-w)^{2}-(x-z)^{2}]}\log\frac{(x-w)^{2}}{(x-z)^{2}}$ (checked by quadrature). \textbf{Condition for the limit}: in the term with the constant kernel (the one whose square bracket contains the same colour bracket at $y'$ and at $y'\to z$) the $y'$ integral has to be done before the $p^{+}$ integral, with the two terms of that bracket together. At fixed coordinates $\int_{0}^{P}dp^{+}e^{-i\frac{p^{+}}{k^{+}}s}=\frac{k^{+}}{is}\big(1-e^{-i\frac{P}{k^{+}}s}\big)$, $s=k\cdot(z-y')$, has no limit for $P\to\infty$; after the $y'$ integral the integrand is ${\cal O}(k^{+2}/p^{+2})$. The $y'$ integral of the subtracted piece converges only conditionally at large $y'$ (to be integrated symmetrically; Gaussian regulator in the check). \textbf{Numerical check} (as in the first part): this form with the upper limit $\vee-k^{+}$, for $\frac{\vee}{k^{+}}=10,\,30,\,60,\,100$: $-0.0960-0.2091\,i,\ -0.1086-0.2357\,i,\ -0.1127-0.2427\,i,\ -0.1144-0.2455\,i$; it converges like $\frac{k^{+}}{\vee}$: extrapolated limit $-0.1170-0.2497\,i$, $\frac{\vee}{k^{+}}\times|{\rm value}-{\rm limit}|=0.46,\ 0.49,\ 0.49$ at $\frac{\vee}{k^{+}}=10,30,60$. The removed form: $-0.0732-0.0467\,i,\ -0.0749+0.0042\,i,\ -0.0720+0.0462\,i,\ -0.0687+0.0794\,i$. Second set of points, $\frac{\vee}{k^{+}}=10,\,30,\,60$: $0.4870+0.0491\,i,\ 0.5728+0.0643\,i,\ 0.5966+0.0625\,i$ (successive differences $0.087,\ 0.024$, ratio $3.7$). The plus prescription acts on the phase and on $1/[p^{+}(x-z)^{2}+k^{+}(x-w)^{2}]$. The soft subtraction terms of the plus prescriptions in this equation are finite only in the sum over Group I Rows III, V and Group III Rows III, V, VI under one integral.]\par\noindent}\par
\paragraph{Group III, Row V: as for Row III.}
\noindent\textbf{Original, lines 208--210:}\par\noindent{\small \begin{align*}
    \frac{d^{3}N_{\rm NLO}}{d^{3}k} & = \frac{1}{(2\pi)^{3}}\, \frac{ig^{4}\,f^{c'ca}}{4\pi^{4}}\, \frac{1}{k^{+}}\,\int_{0}^{\vee-k^{+}}\, \frac{dp^{+}}{2\pi}\, \int_{w,y',z,y,x,x'}\, e^{-ik\cdot (w-y')}\, e^{-ik\cdot \frac{p^{+}}{k^{+}}\,(z-y')}\, \frac{(y'-z)^{m}\,(x'-y')^{k'}\,(y-w)^{i}\,(x-z)^{j}}{(y'-z)^{2}\,(x'-y')^{2}\,(y-w)^{2}\,(x-z)^{2}}\notag \\ & \left[\frac{\delta_{k'm}\,\delta_{ij}}{(p^{+}+k^{+})}-\left[\frac{1}{p^{+}}\right]_{+}\,\delta_{jm}\,\delta_{ik'}-\frac{1}{k^{+}}\, \delta_{im}\, \delta_{jk'}\right]\, \left[U^{c'e'}(y')\,\rho^{e'}(x')-U^{c'e'}(x')\,\rho^{e'}(x')\right]\times \notag \\ &  \left[U^{ce}(x)\,U^{ad}(y)\,\rho^{e}(x)\rho^{d}(y)-U^{ad}(w)\,U^{ce}(x)\,\rho^{e}(x)\rho^{d}(y)\right]
\end{align*}}\par
\noindent\textbf{Checked copy, lines 201--206:}\par\noindent{\small \begin{align*}
    \frac{d^{3}N_{\rm NLO}}{d^{3}k} & = \frac{1}{(2\pi)^{3}}\, \frac{ig^{4}\,f^{c'ca}}{4\pi^{4}}\, \frac{1}{k^{+}}\,\int_{0}^{\infty}\, \frac{dp^{+}}{2\pi}\, \int_{w,y',z,y,x,x'}\, e^{-ik\cdot (w-y')}\, e^{-ik\cdot \frac{p^{+}}{k^{+}}\,(z-y')}\,\frac{(y'-z)^{m}\,(x'-y')^{k'}\,(y-w)^{i}\,(x-z)^{j}}{(y'-z)^{2}\,(x'-y')^{2}\,(y-w)^{2}\,(x-z)^{2}}\notag \\ & \left[\frac{\delta_{k'm}\,\delta_{ij}}{(p^{+}+k^{+})}-\left(\left[\frac{k^{+}}{p^{+}\,(p^{+}+k^{+})}\right]_{+}+\frac{1}{(p^{+}+k^{+})}\right)\,\delta_{jm}\,\delta_{ik'}\right]\, \left[U^{c'e'}(y')\,\rho^{e'}(x')-U^{c'e'}(x')\,\rho^{e'}(x')\right]\times \notag \\ &  \left[U^{ce}(x)\,U^{ad}(y)\,\rho^{e}(x)\rho^{d}(y)-U^{ad}(w)\,U^{ce}(x)\,\rho^{e}(x)\rho^{d}(y)\right] \notag \\
&   +\,\frac{1}{(2\pi)^{3}}\, \frac{ig^{4}\,f^{c'ca}}{4\pi^{4}}\, \frac{1}{k^{+}}\,\int_{0}^{\infty}\, \frac{dp^{+}}{2\pi}\, \int_{w,y',z,y,x,x'}\, e^{-ik\cdot (w-y')}\, e^{-ik\cdot \frac{p^{+}}{k^{+}}\,(z-y')}\,\frac{(y'-z)^{m}\,(y-w)^{i}\,(x-z)^{j}}{(y'-z)^{2}\,(y-w)^{2}\,(x-z)^{2}}\notag \\ &   \left[-\frac{1}{k^{+}}\, \delta_{im}\, \delta_{jk'}\right]\, \Bigg[\frac{(x'-y')^{k'}}{(x'-y')^{2}}\,\left[U^{c'e'}(y')\,\rho^{e'}(x')-U^{c'e'}(x')\,\rho^{e'}(x')\right]\notag \\ & -\frac{(x'-z)^{k'}}{(x'-z)^{2}}\,\left[U^{c'e'}(z)\,\rho^{e'}(x')-U^{c'e'}(x')\,\rho^{e'}(x')\right]\Bigg]\times \notag \\ &    \left[U^{ce}(x)\,U^{ad}(y)\,\rho^{e}(x)\rho^{d}(y)-U^{ad}(w)\,U^{ce}(x)\,\rho^{e}(x)\rho^{d}(y)\right]
\end{align*}
[This is Row V of \texttt{Other\_terms} (v2) at $\vee\to\infty$: upper limit $\infty$, last term ($\log\big(1-\frac{k^{+}}{\vee}\big)\times\dots$) dropped; nothing else is changed. Kernels and $p^{+}$ dependence are those of Row III, first part, with twice the prefactor: the condition for the limit ($y'$ integral before the $p^{+}$ integral in the second term; conditional convergence) and the numerical check are the same (numbers of Row III, first part, $\times2$). The soft subtraction terms of the plus prescriptions in this equation are finite only in the sum over Group I Rows III, V and Group III Rows III, V, VI under one integral.]\par\noindent}\par
\paragraph{Group III, Row VI: as for Row III.}
\noindent\textbf{Original, lines 212--214:}\par\noindent{\small \begin{align*}
    \frac{d^{3}N_{\rm NLO}}{d^{3}k} & = \frac{1}{(2\pi)^{3}}\, \frac{ig^{4}\,f^{c'd'a}}{4\pi^{4}}\, \frac{1}{k^{+}}\,\int_{0}^{\vee-k^{+}}\, \frac{dp^{+}}{2\pi}\, \int_{w,y',z,y,x,x'}\, e^{-ik\cdot (w-y')}\, e^{-ik\cdot \frac{p^{+}}{k^{+}}\,(z-y')}\, \frac{(y'-z)^{m}\,(x'-y')^{k'}\,(y-w)^{i}\,(x-z)^{j}}{(y'-z)^{2}\,(x'-y')^{2}\,(y-w)^{2}\,(x-z)^{2}}\notag \\ & \left[\frac{\delta_{k'm}\,\delta_{ij}}{(p^{+}+k^{+})}-\left[\frac{1}{p^{+}}\right]_{+}\,\delta_{jm}\,\delta_{ik'}-\frac{1}{k^{+}}\, \delta_{im}\, \delta_{jk'}\right]\, \left[U^{c'e'}(y')\,U^{d'b}(z)\,\rho^{e'}(x')-U^{d'b}(z)\,U^{c'e'}(x')\,\rho^{e'}(x')\right]\times\notag \\ &  \left[U^{cb}(z)\,U^{ad}(y)\,U^{ce}(x)\,\rho^{d}(y)\,\rho^{e}(x)-U^{ad}(y)\,\rho^{d}(y)\,\rho^{b}(x)\right]
\end{align*}}\par
\noindent\textbf{Checked copy, lines 208--212:}\par\noindent{\small \begin{align*}
    \frac{d^{3}N_{\rm NLO}}{d^{3}k} & = \frac{1}{(2\pi)^{3}}\, \frac{ig^{4}\,f^{c'd'a}}{4\pi^{4}}\, \frac{1}{k^{+}}\,\int_{0}^{\infty}\, \frac{dp^{+}}{2\pi}\, \int_{w,y',z,y,x,x'}\, e^{-ik\cdot (w-y')}\, e^{-ik\cdot \frac{p^{+}}{k^{+}}\,(z-y')}\,\frac{(y'-z)^{m}\,(x'-y')^{k'}\,(y-w)^{i}\,(x-z)^{j}}{(y'-z)^{2}\,(x'-y')^{2}\,(y-w)^{2}\,(x-z)^{2}}\notag \\ & \left[\frac{\delta_{k'm}\,\delta_{ij}}{(p^{+}+k^{+})}-\left(\left[\frac{k^{+}}{p^{+}\,(p^{+}+k^{+})}\right]_{+}+\frac{1}{(p^{+}+k^{+})}\right)\,\delta_{jm}\,\delta_{ik'}\right]\, \left[U^{c'e'}(y')\,U^{d'b}(z)\,\rho^{e'}(x')-U^{d'b}(z)\,U^{c'e'}(x')\,\rho^{e'}(x')\right]\times\notag \\ &  \left[U^{cb}(z)\,U^{ad}(y)\,U^{ce}(x)\,\rho^{d}(y)\,\rho^{e}(x)-U^{ad}(y)\,\rho^{d}(y)\,\rho^{b}(x)\right] \notag \\ &  +\,\frac{1}{(2\pi)^{3}}\, \frac{ig^{4}\,f^{c'd'a}}{4\pi^{4}}\, \frac{1}{k^{+}}\,\int_{0}^{\infty}\, \frac{dp^{+}}{2\pi}\, \int_{w,y',z,y,x,x'}\, e^{-ik\cdot (w-y')}\, e^{-ik\cdot \frac{p^{+}}{k^{+}}\,(z-y')}\,\frac{(y'-z)^{m}\,(y-w)^{i}\,(x-z)^{j}}{(y'-z)^{2}\,(y-w)^{2}\,(x-z)^{2}}\notag \\ &  \left[-\frac{1}{k^{+}}\, \delta_{im}\, \delta_{jk'}\right]\, \Bigg[\frac{(x'-y')^{k'}}{(x'-y')^{2}}\,\left[U^{c'e'}(y')\,U^{d'b}(z)\,\rho^{e'}(x')-U^{d'b}(z)\,U^{c'e'}(x')\,\rho^{e'}(x')\right]\notag \\ & -\frac{(x'-z)^{k'}}{(x'-z)^{2}}\,\left[U^{c'e'}(z)\,U^{d'b}(z)\,\rho^{e'}(x')-U^{d'b}(z)\,U^{c'e'}(x')\,\rho^{e'}(x')\right]\Bigg]\times\notag \\ &    \left[U^{cb}(z)\,U^{ad}(y)\,U^{ce}(x)\,\rho^{d}(y)\,\rho^{e}(x)-U^{ad}(y)\,\rho^{d}(y)\,\rho^{b}(x)\right]
\end{align*}
[This is Row VI of \texttt{Other\_terms} (v2) at $\vee\to\infty$: upper limit $\infty$, last term ($\log\big(1-\frac{k^{+}}{\vee}\big)\times\dots$) dropped; nothing else is changed. Kernels and $p^{+}$ dependence are those of Row III, first part, with twice the prefactor: the condition for the limit ($y'$ integral before the $p^{+}$ integral in the second term; conditional convergence) and the numerical check are the same (numbers of Row III, first part, $\times2$). The soft subtraction terms of the plus prescriptions in this equation are finite only in the sum over Group I Rows III, V and Group III Rows III, V, VI under one integral.]}\par
\subsubsection*{Changes, block by block}
\needspace{3\baselineskip}\noindent\rule{\linewidth}{0.3pt}\par\noindent\textbf{Lines 30}\par
\begin{itemize}[leftmargin=2.2em,itemsep=1pt,topsep=2pt]
\item[\textbf{l.30}] \textsf{\small [typo/LaTeX]} \emph{added:} $\displaystyle ,$
\end{itemize}
\needspace{3\baselineskip}\noindent\rule{\linewidth}{0.3pt}\par\noindent\textbf{Lines 36}\par
\begin{itemize}[leftmargin=2.2em,itemsep=1pt,topsep=2pt]
\item[\textbf{l.36}] \textsf{\small [note only]} \emph{added:} \noindent[\textbf{Final version of the author's file.} \textbf{Colours.} Red: changed or added with respect to the author's file (red text in square brackets: notes).  Green: the two rows that cancel each other.
\end{itemize}
\begin{notebox}\textsf{\small Note, l.36:} Black: the author's text, unchanged.\par\noindent\textbf{The author's decisions implemented here.} (1) Group III Rows III, V, VI: the large logarithm $\log\frac{\vee}{k^{+}}$ is removed completely (all three terms of each row in the BFKL list); what is left, in the split $\big[\frac{k^{+}}{p^{+}(p^{+}+k^{+})}\big]_{+}+\frac{1}{p^{+}+k^{+}}$ with the subtracted terms, is the finite term, with the upper limit of the $p^{+}$ integral $\infty$. (2) Group II Row IV: the term $-\frac{1}{k^{+}}\times$(BFKL structure) is the left-over of the BFKL term and is counted as finite; the subtraction of the DGLAP term stays inside the splitting function. (3) Group I Row III: the separate equations are added together (combined kernel $\mathcal{K}^{ij}$). (4) Group I Row IV and Group III Row I cancel each other identically; both are shown, in green. (5) All terms with $\log\big(1-\frac{k^{+}}{\vee}\big)$ are dropped: at $\vee\to\infty$ this logarithm is $0$. (6) The plus-prescription terms of Group I Row V and of Group III stay in the author's form; their soft subtraction terms are finite only in the sum over rows.\par\noindent\textbf{Not verified.} No complete numerical integration over the transverse plane with Wilson lines has been done; the cancellation of the large-distance tails in the sum over rows rests on the tail coefficients of the single terms, and nothing is assumed about $U$ at large distance.\par\noindent\textbf{Conventions.} Phase $e^{-ik\cdot(w-w')}$; c.c.\ implied for Groups I and III; the $\rho$'s do not commute and stand in the order written; logarithms are $\log|\dots|$ (principal value); $\int dp^{+}\big[\frac{1}{p^{+}}\big]_{+}f\equiv\int dp^{+}\,\frac{f(p^{+})-f(0)}{p^{+}}$, where $f$ is everything to the right that depends on $p^{+}$ (phase and denominators); no $1/\epsilon$ is written; one scale $\mu$. One displayed line of the author's file that ran past the edge of the page (Group II Row I, $N_{2a}$) is broken, without change of content.]\end{notebox}
\needspace{3\baselineskip}\noindent\rule{\linewidth}{0.3pt}\par\noindent\textbf{Lines 59} \quad \textsf{Group I (with c.c.) $\rangle$ Row III}\par
\begin{notebox}\textsf{\small Note on removed material, l.59:} [Removed here: the sentence ``We write it in three parts:'' (l.58 of the author's file) and the four separate equations for $\frac{d^{3}N_{1\,\rm NLO}}{d^{3}k}$ (l.59--64), $\frac{d^{3}N_{2\,\rm NLO}}{d^{3}k}$ (l.65--72), $\frac{d^{3}N_{3a\,\rm NLO}}{d^{3}k}$ (l.73--96) and $\frac{d^{3}N_{3b\,\rm NLO}}{d^{3}k}$ (l.97--107), because they are not integrable one by one (the poles on the line $(x-z)^{2}=(x-w)^{2}$ and the $\frac{1}{(z-w)^{2}}$ at $z\to w$ cancel only in their sum), and the last term of $N_{3b}$ (l.106), $-\frac{ig^{4}f^{abd}}{16\pi^{5}}\frac{1}{k^{+}}\int_{x,y,z,x'}\log\big[\frac{(z-w)^{2}}{(x-w)^{2}}\big]\frac{(x'-w')\cdot(y-z)}{(x'-w')^{2}(y-z)^{2}(z-w)^{2}}[\dots][\dots]$, whose $z$ integral diverges at $z\to w$. They are replaced by their sum, with the last term of $N_{3b}$ in its finite form.]\end{notebox}
\noindent{\small\textsf{Unmarked edit, l.59 (deleted):} \texttt{We write it in three parts:}}\par
\needspace{3\baselineskip}\noindent\rule{\linewidth}{0.3pt}\par\noindent\textbf{Lines 61--68} \quad {\small (in \texttt{align}, lines 60--67)}\par
\begin{itemize}[leftmargin=2.2em,itemsep=1pt,topsep=2pt]
\item[\textbf{l.61}] \textsf{\small [term removed/replaced]} $\displaystyle \frac{d^{3}N_{1\,\rm NLO}^{\rm}}{d^{3}k}$\ $\longrightarrow$\ $\displaystyle \frac{d^{3}N_{\rm NLO}}{d^{3}k}$
\item[\textbf{l.61}] \textsf{\small [phase convention]} \emph{added:} $\displaystyle \frac{1}{(2\pi)^{3}}\int_{w,w'}e^{-ik\cdot(w-w')}\,\Bigg\{$
\item[\textbf{l.62}] \textsf{\small [term removed/replaced]} \par\noindent\emph{original:} $\displaystyle +\frac{ig^{4}\,f^{abd}}{16\pi^{5}}\, \frac{1}{k^{+}}\,\int_{x,y,z,x'}\, \frac{(x'-w')\cdot (y-z)}{(x'-w')^{2}\,(y-z)^{2}}\, \frac{[(x-w)^{2}+(x-z)^{2}]}{[(x-z)^{2}-(x-w)^{2}]^{2}}\, \times\, \log{[\frac{(x-w)^{2}-(x-z)^{2}}{(x-w)^{2}}]}$\par\noindent\emph{checked:} $\displaystyle \frac{ig^{4}f^{abd}}{8\pi^{5}}\,\frac{1}{k^{+}}\,\int_{x,y,z,x'}\frac{(x'-w')^{i}\,(y-z)^{j}}{(x'-w')^{2}\,(y-z)^{2}}\;\mathcal{K}^{ij}$
\item[\textbf{l.62}] \textsf{\small [typo/LaTeX]} \par\noindent\emph{original:} $\displaystyle [U^{ab'}(x')\,\rho^{b'}(x')-U^{ab'}(w')\, \rho^{b'}(x')]\times\, [U^{de}(x)\,U^{bc}(z)\, \rho^{e}(x)\,\rho^{c}(y)-U^{de}(x)\,U^{bc}(y)\,\rho^{e}(x)\,\rho^{c}(y)]$\par\noindent\emph{checked:} $\displaystyle [U^{ab'}(x')\,\rho^{b'}(x')-U^{ab'}(w')\, \rho^{b'}(x')]\times [U^{de}(x)\,U^{bc}(z)\, \rho^{e}(x)\,\rho^{c}(y)-U^{de}(x)\,U^{bc}(y)\,\rho^{e}(x)\,\rho^{c}(y)]$
\item[\textbf{l.66}] \textsf{\small [term removed/replaced]} \par\noindent\emph{original:} $\displaystyle -\frac{ig^{4}f^{abd}}{16\pi^{5}}\,\frac{1}{k^{+}}\,\int_{x,y,z,x'}\log{[\frac{(z-w)^{2}}{(x-w)^{2}}]}\,\frac{(x'-w')\cdot(y-z)}{(x'-w')^{2}(y-z)^{2}\,(z-w)^{2}}$\par\noindent\emph{checked:} $\displaystyle -\frac{ig^{4}f^{abd}}{8\pi^{5}}\,\frac{1}{k^{+}}\,\int_{x,y,z,x'}\mathcal{E}\big(k\cdot(z-w)\big)\,\frac{(x-z)\cdot(w-z)\,(x'-w')\cdot(y-w)}{(x-z)^{2}\,(w-z)^{2}\,(x'-w')^{2}\,(y-w)^{2}}$
\item[\textbf{l.66}] \textsf{\small [term removed/replaced]} {\small\ttfamily \textbackslash{}\allowbreak{}left[U\^{}}\ $\longrightarrow$\ $\displaystyle [U^{ab'}(x')\,\rho^{b'}(x')-U^{ab'}(w')\, \rho^{b'}(x')]\times [U^{de}(x)\,U^{bc}(w)\, \rho^{e}(x)\,\rho^{c}(y)-U^{de}(x)\,U^{bc}(y)\,\rho^{e}(x)\,\rho^{c}(y)]\Bigg\}$
\item[\textbf{l.68}] \textsf{\small [typo/LaTeX]} {\small\ttfamily (x')\allowbreak{}\textbackslash{}\allowbreak{},\allowbreak{}\textbackslash{}\allowbreak{}rho\^{}\{b'\}\allowbreak{}(x')\allowbreak{}-U\^{}}\ $\longrightarrow$\ with
\end{itemize}
\noindent{\small\textsf{Unmarked edit, l.66 (inserted):} \texttt{ \textbackslash{}\allowbreak{}end}}\par
\noindent{\small\textsf{Unmarked edit, l.67 (inserted):} \texttt{lign}}\par
\noindent{\small\textsf{Unmarked edit, l.67 (inserted):} \texttt{  \textbackslash{}\allowbreak{}begin}}\par
\noindent{\small\textsf{Unmarked edit, l.69 (inserted):} \texttt{lign}}\par
\noindent{\small\textsf{Unmarked edit, l.61 (deleted):} \texttt{  \textbackslash{}\allowbreak{}frac\{1\}\allowbreak{}\{(2\textbackslash{}\allowbreak{}pi)\allowbreak{}\^{}\{3\}\allowbreak{}\}\allowbreak{}  \textbackslash{}\allowbreak{}int\_\{w,\allowbreak{} w\^{}\textbackslash{}\allowbreak{}prime\}\allowbreak{} e\^{}\{-ik(w-w\^{}\textbackslash{}\allowbreak{}prime)\allowbreak{}\}\allowbreak{}\textbackslash{}\allowbreak{},\allowbreak{} \textbackslash{}\allowbreak{}Bigg\textbackslash{}\allowbreak{}\{\textbackslash{}\allowbreak{}notag \textbackslash{}\allowbreak{}\textbackslash{}\allowbreak{} \& \textbackslash{}\allowbreak{}frac\{ig\^{}\{4\}\allowbreak{}\textbackslash{}\allowbreak{},\allowbreak{}f\^{}\{abd\}\allowbreak{}\}\allowbreak{}\{16\textbackslash{}\allowbreak{}pi\^{}\{5\}\allowbreak{}\}\allowbreak{}\textbackslash{}\allowbreak{},\allowbreak{} \textbackslash{}\allowbreak{}frac\{1\}\allowbreak{}\{k\^{}\{+\allowbreak{}\}\allowbreak{}\}\allowbreak{}\textbackslash{}\allowbreak{},\allowbreak{}\textbackslash{}\allowbreak{}int\_\{x,\allowbreak{}y,\allowbreak{}z,\allowbreak{}x'\}\allowbreak{}\textbackslash{}\allowbreak{},\allowbreak{} \textbackslash{}\allowbreak{}frac\{(x'-w')\allowbreak{}\textbackslash{}\allowbreak{}cdot (y-z)\allowbreak{}\}\allowbreak{}\{(x'-w')\allowbreak{}\^{}\{2\}\allowbreak{}\textbackslash{}\allowbreak{},\allowbreak{}(y-z)\allowbreak{}\^{}\{2\}\allowbreak{}\}\allowbreak{}\textbackslash{}\allowbreak{},\allowbreak{} \textbackslash{}\allowbreak{}frac\{\textbackslash{}\allowbreak{}left[(x-w)\allowbreak{}\^{}\{2\}\allowbreak{}+\allowbreak{}(x-z)\allowbreak{}\^{}\{2\}\allowbreak{}\textbackslash{}\allowbreak{}right]\}\allowbreak{}\{\textbackslash{}\allowbreak{}left[(x-z)\allowbreak{}\^{}\{2\}\allowbreak{}-(x-w)\allowbreak{}\^{}\{2\}\allowbreak{}\textbackslash{}\allowbreak{}right]\^{}\{2\}\allowbreak{}\}\allowbreak{}\textbackslash{}\allowbreak{},\allowbreak{} \textbackslash{}\allowbreak{}times\textbackslash{}\allowbreak{},\allowbreak{} \textbackslash{}\allowbreak{}log\{\textbackslash{}\allowbreak{}left[\textbackslash{}\allowbreak{}frac\{(x-z)\allowbreak{}\^{}\{2\}\allowbreak{}\}\allowbreak{}\{(x-z)\allowbreak{}\^{}\{2\}\allowbreak{}-(x-w)\allowbreak{}\^{}\{2\}\allowbreak{}\}\allowbreak{}\textbackslash{}\allowbreak{}right]\}\allowbreak{}\textbackslash{}\allowbreak{}notag \textbackslash{}\allowbreak{}\textbackslash{}\allowbreak{} \& \textbackslash{}\allowbreak{}left[U\^{}\{ab'\}\allowbreak{}(x')\allowbreak{}\textbackslash{}\allowbreak{},\allowbreak{}\textbackslash{}\allowbreak{}rho\^{}\{b'\}\allowbreak{}(x')\allowbreak{}-U\^{}\{}}\par
\noindent{\small\textsf{Unmarked edit, l.63 (deleted):} \texttt{\&-\textbackslash{}\allowbreak{},\allowbreak{} \textbackslash{}\allowbreak{}frac\{ig\^{}\{4\}\allowbreak{}\textbackslash{}\allowbreak{},\allowbreak{}f\^{}\{abd\}\allowbreak{}\}\allowbreak{}\{8\textbackslash{}\allowbreak{}pi\^{}\{5\}\allowbreak{}\}\allowbreak{}\textbackslash{}\allowbreak{},\allowbreak{}\textbackslash{}\allowbreak{}frac\{1\}\allowbreak{}\{k\^{}\{+\allowbreak{}\}\allowbreak{}\}\allowbreak{}\textbackslash{}\allowbreak{},\allowbreak{}\textbackslash{}\allowbreak{}int\_\{x,\allowbreak{}y,\allowbreak{}z,\allowbreak{}x'\}\allowbreak{}\textbackslash{}\allowbreak{}frac\{(x'-w')\allowbreak{}\textbackslash{}\allowbreak{}cdot (x-w)\allowbreak{}\textbackslash{}\allowbreak{},\allowbreak{} (x-z)\allowbreak{}\textbackslash{}\allowbreak{}cdot (y-z)\allowbreak{}\}\allowbreak{}\{(x'-w')\allowbreak{}\^{}\{2\}\allowbreak{}\textbackslash{}\allowbreak{},\allowbreak{} (x-w)\allowbreak{}\^{}\{2\}\allowbreak{}\textbackslash{}\allowbreak{},\allowbreak{} (x-z)\allowbreak{}\^{}\{2\}\allowbreak{}\textbackslash{}\allowbreak{},\allowbreak{} (y-z)\allowbreak{}\^{}\{2\}\allowbreak{}\}\allowbreak{}\textbackslash{}\allowbreak{},\allowbreak{}\textbackslash{}\allowbreak{}log\{\textbackslash{}\allowbreak{}left[\textbackslash{}\allowbreak{}frac\{(x-z)\allowbreak{}\^{}\{2\}\allowbreak{}\}\allowbreak{}\{(x-w)\allowbreak{}\^{}\{2\}\allowbreak{}\}\allowbreak{}\textbackslash{}\allowbreak{}right]\}\allowbreak{} \textbackslash{}\allowbreak{}notag \textbackslash{}\allowbreak{}\textbackslash{}\allowbreak{} \& \textbackslash{}\allowbreak{}left[U\^{}\{ab'\}\allowbreak{}(x')\allowbreak{}\textbackslash{}\allowbreak{},\allowbreak{}\textbackslash{}\allowbreak{}rho\^{}\{b'\}\allowbreak{}(x')\allowbreak{}-U\^{}\{ab'\}\allowbreak{}(w')\allowbreak{}\textbackslash{}\allowbreak{},\allowbreak{} \textbackslash{}\allowbreak{}rho\^{}\{b'\}\allowbreak{}(x')\allowbreak{}\textbackslash{}\allowbreak{}right]\textbackslash{}\allowbreak{}times\textbackslash{}\allowbreak{},\allowbreak{} \textbackslash{}\allowbreak{}left[U\^{}\{de\}\allowbreak{}(x)\allowbreak{}\textbackslash{}\allowbreak{},\allowbreak{}U\^{}\{bc\}\allowbreak{}(z)\allowbreak{}\textbackslash{}\allowbreak{},\allowbreak{} \textbackslash{}\allowbreak{}rho\^{}\{e\}\allowbreak{}(x)\allowbreak{}\textbackslash{}\allowbreak{},\allowbreak{}\textbackslash{}\allowbreak{}rho\^{}\{c\}\allowbreak{}(y)\allowbreak{}-U\^{}\{de\}\allowbreak{}(x)\allowbreak{}\textbackslash{}\allowbreak{},\allowbreak{}U\^{}\{bc\}\allowbreak{}(y)\allowbreak{}\textbackslash{}\allowbreak{},\allowbreak{}\textbackslash{}\allowbreak{}rho\^{}\{e\}\allowbreak{}(x)\allowbreak{}\textbackslash{}\allowbreak{},\allowbreak{}\textbackslash{}\allowbreak{}rho\^{}\{c\}\allowbreak{}(y}}\par
\noindent{\small\textsf{Unmarked edit, l.67 (deleted):} \texttt{b'}}\par
\noindent{\small\textsf{Unmarked edit, l.69 (deleted):} \texttt{b'}}\par
\needspace{3\baselineskip}\noindent\rule{\linewidth}{0.3pt}\par\noindent\textbf{Lines 70--74} \quad {\small (in \texttt{align}, lines 69--75)}\par
\begin{itemize}[leftmargin=2.2em,itemsep=1pt,topsep=2pt]
\item[\textbf{l.70}] \textsf{\small [term removed/replaced]} $\displaystyle (w')\,$\ $\longrightarrow$\ $\displaystyle \mathcal{K}^{ij}$
\item[\textbf{l.70}] \textsf{\small [term removed/replaced]} \par\noindent\emph{original:} $\displaystyle \rho^{e}(x)\,\rho^{c}(y)-U^{de}(x)\,U^{bc}(y)\,\rho^{e}(x)\,\rho^{c}(y)]\Bigg\}$\par\noindent\emph{checked:} $\displaystyle \log[\frac{(x-z)^{2}}{(x-w)^{2}}]\Bigg\{\frac{\delta^{ij}}{2\,(z-w)^{2}}-\frac{(x-w)^{i}\,(x-z)^{j}}{(x-w)^{2}\,(x-z)^{2}}+\frac{1}{(z-w)^{2}[(x-w)^{2}-(x-z)^{2}]}\Big[(x-w)^{i}(z-w)^{j}\,\frac{(x-z)^{2}}{(x-w)^{2}}$
\item[\textbf{l.71}] \textsf{\small [missing term added]} \emph{added:} $\displaystyle +(z-w)^{i}(x-z)^{j}\,\frac{(x-w)^{2}}{(x-z)^{2}}-(z-w)^{i}(x-w)^{j}-(x-z)^{i}(z-w)^{j}\Big]\Bigg\}$
\item[\textbf{l.72}] \textsf{\small [missing term added]} \emph{added:} $\displaystyle +\frac{\mathcal{A}_{0}}{(z-w)^{2}}\Bigg\{(z-w)^{i}(x-w)^{j}+(x-z)^{i}(z-w)^{j}-(x-w)\cdot(x-z)[\frac{(z-w)^{i}(x-z)^{j}}{(x-z)^{2}}+\frac{(x-w)^{i}(z-w)^{j}}{(x-w)^{2}}]\Bigg\}$
\item[\textbf{l.73}] \textsf{\small [missing term added]} \emph{added:} $\displaystyle -\frac{1}{2\,(z-w)^{2}}\log[\frac{(z-w)^{2}}{(x-w)^{2}}]\frac{(z-w)^{i}(x-z)^{j}}{(x-z)^{2}}\,,$
\item[\textbf{l.74}] \textsf{\small [missing term added]} \emph{added:} $\displaystyle \mathcal{E}(s)$
\item[\textbf{l.74}] \textsf{\small [phase convention]} \emph{added:} $\displaystyle \equiv\int_{0}^{1}\frac{dt}{t}\,\big(e^{-ist}-1\big)=\mathrm{Ci}(|s|)-\gamma_{E}-\log|s|-i\,\mathrm{Si}(s)$
\end{itemize}
\noindent{\small\textsf{Unmarked edit, l.69 (inserted):} \texttt{lign}}\par
\noindent{\small\textsf{Unmarked edit, l.70 (inserted):} \texttt{=\allowbreak{}}}\par
\noindent{\small\textsf{Unmarked edit, l.69 (deleted):} \texttt{b'}}\par
\noindent{\small\textsf{Unmarked edit, l.70 (deleted):} \texttt{\textbackslash{}\allowbreak{}rho\^{}\{b'\}\allowbreak{}(x')\allowbreak{}\textbackslash{}\allowbreak{}right]\textbackslash{}\allowbreak{}times}}\par
\noindent{\small\textsf{Unmarked edit, l.70 (deleted):} \texttt{\textbackslash{}\allowbreak{}left[U\^{}\{de\}\allowbreak{}(x)\allowbreak{}\textbackslash{}\allowbreak{},\allowbreak{}U\^{}\{bc\}\allowbreak{}(z)\allowbreak{}\textbackslash{}\allowbreak{},\allowbreak{}}}\par
\needspace{3\baselineskip}\noindent\rule{\linewidth}{0.3pt}\par\noindent\textbf{Lines 78}\par
\begin{notebox}\textsf{\small Note, l.78:} [\textbf{First term}: the 24 terms of $N_{1}$, $N_{2}$, $N_{3a}$ and the first two terms of $N_{3b}$, added (Eqs.~(1.5), (1.6) of \texttt{Finite\_terms\_final}). \textbf{Second term} (black): the third term of $N_{3b}$, unchanged. Verified at random points with explicit SU(3) matrices: [$N_{1}+N_{2}+N_{3a}+$ first three terms of $N_{3b}$ of the author's file] $-$ [first $+$ second term here] $=$ $2\times10^{-22}$ to $7\times10^{-18}$ against $7\times10^{-9}$ to $5\times10^{-6}$ at eight points (relative $5\times10^{-16}$ to $1.5\times10^{-12}$). The separate equations at distance $\delta=10^{-3},10^{-4},10^{-5}$ from the line $(x-z)^{2}=(x-w)^{2}$: $|N_{2}|\simeq|N_{3a}|=0.89,\ 1.2\times10^{2},\ 1.6\times10^{4}$, $|N_{1}|\simeq|N_{2}+N_{3a}|=3.0\times10^{-4},\ 3.0\times10^{-3},\ 3.0\times10^{-2}$, sum of the four $=6.0\times10^{-7}$ at all three. \textbf{Third term}: it replaces the last term of $N_{3b}$ (orange sentence above). That term is the limit $\Lambda\to0$ at fixed $z\neq w$ only; with the kernel $\frac{1}{(z-w)^{2}}$ and $U^{bc}(z)$, which nothing softens at $z=w$, its $z$ integral diverges like $\log^{2}$ at $z\to w$ ($|z-w|^{2}\,|{\rm term}|/\log\frac{1}{|z-w|}=5.1\times10^{-7},\ 4.6\times10^{-7},\ 4.5\times10^{-7}$ at $|z-w|=10^{-2},10^{-3},10^{-4}$, i.e.\ constant). Taken together with its $\log\frac{k^{+}}{\Lambda}$ partner in the row at finite $\Lambda$ and after the change of variable to the position of the soft gluon, it is the term of the JIMWLK list with $U^{bc}(w)$ (three-$\rho$ list, Row III) plus the third term here (third line of Eq.~(1.5) of \texttt{Finite\_terms\_final}, derivation Eqs.~(1.9)--(1.11) there; the identity was checked there with a Gaussian regulator, re-run now: both sides $-10.896402+0.822556\,i$, difference $3\times10^{-13}$). The JIMWLK term of this row is therefore the one of the evolution lists, with $U^{bc}(w)$. $\mathcal{E}$: closed form against quadrature, $10^{-16}$. \textbf{Integrability}: the first term is integrable only as the combined kernel (not term by term), and at large $z$ only conditionally (angular integral first). The second and the third term have a $\frac{1}{z^{2}}\log$ tail at large $z$: they are finite only in the sum over Group I Rows III, V and Group III Rows III, V, VI under one integral.]\end{notebox}
\needspace{3\baselineskip}\noindent\rule{\linewidth}{0.3pt}\par\noindent\textbf{Lines 81--85} \quad \textsf{Group I (with c.c.) $\rangle$ Row IV} \quad {\small (in \texttt{align}, lines 80--84)}\par
\begin{itemize}[leftmargin=2.2em,itemsep=1pt,topsep=2pt]
\item[\textbf{l.81}] \textsf{\small [note only]} $\displaystyle \frac{d^{3}N_{\rm NLO}^{\rm}}{d^{3}k}$\ $\longrightarrow$\ $\displaystyle \frac{d^{3}N_{\rm NLO}^{\rm}}{d^{3}k}$
\item[\textbf{l.81}] \textsf{\small [note only]} $\displaystyle \frac{1}{(2\pi)^{3}}\, \int_{x,y,z,w,x',w'}\,e^{-ik\cdot (w-w^\prime)}\,\Bigg\{$\ $\longrightarrow$\ $\displaystyle \frac{1}{(2\pi)^{3}}\, \int_{x,y,z,w,x',w'}\,e^{-ik\cdot (w-w^\prime)}\,\Bigg\{$
\item[\textbf{l.83}] \textsf{\small [note only]} \par\noindent\emph{original:} $\displaystyle \frac{ig^{4}}{8\pi^{5}}\,\frac{1}{k^{+}}\,\frac{(x'-w')^{i}\,(y-z)^{j}}{(x'-w')^{2}\,(y-z)^{2}}\times\,\log{[\frac{(x-z)^{2}}{(x-w)^{2}}]}\times\, \bigg[\frac{(x-z)^{j}\,(z-w)^{i}}{(z-w)^{2}\,(x-z)^{2}} -\frac{(x-w)^{i}\,(z-w)^{j}}{(z-w)^{2}\,(x-w)^{2}}-\frac{(x-z)^{j}\,(x-w)^{i}}{(x-z)^{2}\,(x-w)^{2}}$\par\noindent\emph{checked:} $\displaystyle \frac{ig^{4}}{8\pi^{5}}\,\frac{1}{k^{+}}\,\frac{(x'-w')^{i}\,(y-z)^{j}}{(x'-w')^{2}\,(y-z)^{2}}\times\,\log{[\frac{(x-z)^{2}}{(x-w)^{2}}]}\times\, \bigg[\frac{(x-z)^{j}\,(z-w)^{i}}{(z-w)^{2}\,(x-z)^{2}} -\frac{(x-w)^{i}\,(z-w)^{j}}{(z-w)^{2}\,(x-w)^{2}}-\frac{(x-z)^{j}\,(x-w)^{i}}{(x-z)^{2}\,(x-w)^{2}}$
\item[\textbf{l.83}] \textsf{\small [note only]} \par\noindent\emph{original:} $\displaystyle +\frac{\delta^{ij}}{2\,(z-w)^{2}}\bigg]\Bigg\} [U^{ab'}(x')\,\rho^{b'}(x')-U^{ab'}(w')\,\rho^{b'}(x')] [f^{cdc'}\,U^{ac}(w)\,U^{bd}(z)\,U^{be}(y)\,\rho^{e}(y)\rho^{c'}(x)-f^{abc}\,U^{be}(y)\,U^{cd}(x)\,\rho^{e}(y)\rho^{d}(x)]$\par\noindent\emph{checked:} $\displaystyle +\frac{\delta^{ij}}{2\,(z-w)^{2}}\bigg]\Bigg\} [U^{ab'}(x')\,\rho^{b'}(x')-U^{ab'}(w')\,\rho^{b'}(x')] [f^{cdc'}\,U^{ac}(w)\,U^{bd}(z)\,U^{be}(y)\,\rho^{e}(y)\rho^{c'}(x)-f^{abc}\,U^{be}(y)\,U^{cd}(x)\,\rho^{e}(y)\rho^{d}(x)]$
\end{itemize}
\begin{notebox}\textsf{\small Note (highlighted, cancelling rows), l.85:} [This row cancels identically against Group III Row I (shown there).]\end{notebox}
\needspace{3\baselineskip}\noindent\rule{\linewidth}{0.3pt}\par\noindent\textbf{Lines 90--91} \quad \textsf{Group I (with c.c.) $\rangle$ Row V} \quad {\small (in \texttt{align}, lines 89--91)}\par
\begin{itemize}[leftmargin=2.2em,itemsep=1pt,topsep=2pt]
\item[\textbf{l.90}] \textsf{\small [typo/LaTeX]} \emph{added:} $\displaystyle ,$\quad{\color{gray}\small(in context: }$\displaystyle \, \frac{dp^{+}}{2\pi}\,\int_{x,z,y,x',y',w'}\,e^{-ik\cdot (z-w^\prime)}\,e^{-ik\cdot (x-z)\frac{p^{+}}{k^{+}}}${\color{gray}\small)}
\item[\textbf{l.90}] \textsf{\small [typo/LaTeX]} \emph{added:} $\displaystyle ,$\quad{\color{gray}\small(in context: }$\displaystyle \, \frac{dq^{+}}{2\pi}\,\int_{x,z,y,x',y',w'}\,e^{-ik\cdot (x-w^\prime)}\,e^{-ik\cdot (z-x)\frac{q^{+}}{k^{+}}}${\color{gray}\small)}
\end{itemize}
\begin{notebox}\textsf{\small Note, l.91:} [The soft subtraction terms of the plus prescriptions in this equation are finite only in the sum over Group I Rows III, V and Group III Rows III, V, VI under one integral.]\end{notebox}
\needspace{3\baselineskip}\noindent\rule{\linewidth}{0.3pt}\par\noindent\textbf{Lines 95}\par
\begin{notebox}\textsf{\small Note, l.95:} [The soft subtraction terms of the plus prescriptions in this equation are finite only in the sum over Group I Rows III, V and Group III Rows III, V, VI under one integral.]\end{notebox}
\needspace{3\baselineskip}\noindent\rule{\linewidth}{0.3pt}\par\noindent\textbf{Lines 97--107} \quad {\small (in \texttt{align}, lines 96--107)}\par
\begin{itemize}[leftmargin=2.2em,itemsep=1pt,topsep=2pt]
\item[\textbf{l.97}] \textsf{\small [typo/LaTeX]} {\small\ttfamily \}\allowbreak{})\allowbreak{}}\ $\longrightarrow$\ {\small\ttfamily )\allowbreak{}\}\allowbreak{}}
\end{itemize}
\begin{notebox}\textsf{\small Note, l.107:} [The soft subtraction terms of the plus prescriptions in this equation are finite only in the sum over Group I Rows III, V and Group III Rows III, V, VI under one integral.]\end{notebox}
\needspace{3\baselineskip}\noindent\rule{\linewidth}{0.3pt}\par\noindent\textbf{Lines 151--153} \quad \textsf{Group II $\rangle$ Row IV} \quad {\small (in \texttt{align}, lines 150--153)}\par
\begin{itemize}[leftmargin=2.2em,itemsep=1pt,topsep=2pt]
\item[\textbf{l.151}] \textsf{\small [limits/range]} $\displaystyle \frac{k^{+}}{\vee}$\ $\longrightarrow$\ $\displaystyle \,0$
\end{itemize}
\begin{notebox}\textsf{\small Note, l.153:} [Lower limit of the $\xi$ integral: $\frac{k^{+}}{\vee}\to0$ (the integral converges there). The second term, $-\frac{1}{k^{+}}\times\frac{g^{4}N_{c}}{2\pi^{3}k^{2}}\times$(BFKL structure), is the left-over of the BFKL term ($\big[\frac{1}{k^{+}}-\frac{1}{\vee}\big]\to\frac{1}{k^{+}}$) and is counted as finite, as the author has it. It corresponds to the DGLAP term written with the subtraction inside the splitting function, $\int_{0}^{1}d\xi\,\big[\frac{1-\xi}{\xi}\big]_{+}F(\xi)\equiv\int_{0}^{1}d\xi\,\frac{1-\xi}{\xi}\,\big[F(\xi)-F(0)\big]$ ($F$: everything that multiplies the splitting function), the two together being termed DGLAP. With the other reading, $\int_{0}^{1}\frac{d\xi}{\xi}\big[(1-\xi)F(\xi)-F(0)\big]$, this second term would be counted twice: the two readings differ by exactly this term, since $\int_{0}^{1}d\xi\big[\frac{1}{\xi}-\frac{1-\xi}{\xi}\big]=1$ (checked with a model $F$: difference $0$ to machine precision; prefactor ratio $1$)]\end{notebox}
\needspace{3\baselineskip}\noindent\rule{\linewidth}{0.3pt}\par\noindent\textbf{Lines 157--160} \quad \textsf{Group III (with c.c.) $\rangle$ Row I} \quad {\small (in \texttt{align}, lines 156--159)}\par
\begin{itemize}[leftmargin=2.2em,itemsep=1pt,topsep=2pt]
\item[\textbf{l.157}] \textsf{\small [note only]} $\displaystyle \frac{d^{3}N_{\rm NLO}^{\rm}}{d^{3}k}$\ $\longrightarrow$\ $\displaystyle \frac{d^{3}N_{\rm NLO}^{\rm}}{d^{3}k}$
\item[\textbf{l.157}] \textsf{\small [note only]} $\displaystyle \frac{1}{(2\pi)^{3}}\, \int_{x,y,x',z,w,w^{\prime}} e^{-ik\cdot (w-w')}\, \Bigg\{$\ $\longrightarrow$\ $\displaystyle \frac{1}{(2\pi)^{3}}\, \int_{x,y,x',z,w,w^{\prime}} e^{-ik\cdot (w-w')}\, \Bigg\{$
\item[\textbf{l.158}] \textsf{\small [missing term added]} \emph{added:} $\displaystyle -\,\frac{ig^{4}}{8\pi^{5}}\,\frac{1}{k^{+}}\,\frac{(x'-w')^{i}(y-z)^{j}}{(x'-w')^{2}\,(y-z)^{2}}\times \log{[\frac{(x-z)^{2}}{(x-w)^{2}}]}\times \Bigg[\frac{\delta^{ij}}{2\,(z-w)^{2}}+\frac{(x-z)^{j}\,(z-w)^{i}}{(z-w)^{2}\,(x-z)^{2}}$
\item[\textbf{l.158}] \textsf{\small [missing term added]} \emph{added:} $\displaystyle -\frac{(x-z)^{j}\,(x-w)^{i}}{(x-w)^{2}\,(x-z)^{2}}-\frac{(x-w)^{i}\,(z-w)^{j}}{(z-w)^{2}\,(x-w)^{2}}\Bigg]\Bigg\}\times$
\item[\textbf{l.158}] \textsf{\small [term removed/replaced]} $\displaystyle U^{cb}(z)\,U^{de}(x)\,\rho^{e}(x)\Bigg]$\ $\longrightarrow$\ $\displaystyle \Bigg[U^{ab'}(x')\,U^{c'd'}(y)\,U^{c'b}(z)\,\rho^{b'}(x')\rho^{d'}(y)-U^{d'b}(z)\,U^{d'e'}(y)U^{ac'}(w')\,\rho^{c'}(x')\rho^{e'}(y)\Bigg]\times$
\item[\textbf{l.158}] \textsf{\small [term removed/replaced]} $\displaystyle U^{cb}(z)\,U^{de}(x)\,\rho^{e}(x)\Bigg]$\ $\longrightarrow$\ $\displaystyle \Bigg[f^{cbd}\,U^{ac}(w)\,\rho^{d}(x)-f^{acd}\, U^{cb}(z)\,U^{de}(x)\,\rho^{e}(x)\Bigg]$
\end{itemize}
\begin{notebox}\textsf{\small Note (highlighted, cancelling rows), l.160:} [This row cancels identically against Group I Row IV. Term by term: the first colour bracket here is]\end{notebox}
\noindent{\small\textsf{Unmarked edit, l.158 (deleted):} \texttt{ -\textbackslash{}\allowbreak{},\allowbreak{}\textbackslash{}\allowbreak{}frac\{ig\^{}\{4\}\allowbreak{}\}\allowbreak{}\{8\textbackslash{}\allowbreak{}pi\^{}\{5\}\allowbreak{}\}\allowbreak{}\textbackslash{}\allowbreak{},\allowbreak{}\textbackslash{}\allowbreak{}frac\{1\}\allowbreak{}\{k\^{}\{+\allowbreak{}\}\allowbreak{}\}\allowbreak{}\textbackslash{}\allowbreak{},\allowbreak{}\textbackslash{}\allowbreak{}frac\{(x'-w')\allowbreak{}\^{}\{i\}\allowbreak{}(y-z)\allowbreak{}\^{}\{j\}\allowbreak{}\}\allowbreak{}\{(x'-w')\allowbreak{}\^{}\{2\}\allowbreak{}\textbackslash{}\allowbreak{},\allowbreak{}(y-z)\allowbreak{}\^{}\{2\}\allowbreak{}\}\allowbreak{}\textbackslash{}\allowbreak{}times \textbackslash{}\allowbreak{}log\{\textbackslash{}\allowbreak{}left[\textbackslash{}\allowbreak{}frac\{(x-z)\allowbreak{}\^{}\{2\}\allowbreak{}\}\allowbreak{}\{(x-w)\allowbreak{}\^{}\{2\}\allowbreak{}\}\allowbreak{}\textbackslash{}\allowbreak{}right]\}\allowbreak{}\textbackslash{}\allowbreak{}times \textbackslash{}\allowbreak{}Bigg[\textbackslash{}\allowbreak{}frac\{\textbackslash{}\allowbreak{}delta\^{}\{ij\}\allowbreak{}\}\allowbreak{}\{2\textbackslash{}\allowbreak{},\allowbreak{}(z-w)\allowbreak{}\^{}\{2\}\allowbreak{}\}\allowbreak{}+\allowbreak{}\textbackslash{}\allowbreak{}frac\{(x-z)\allowbreak{}\^{}\{j\}\allowbreak{}\textbackslash{}\allowbreak{},\allowbreak{}(z-w)\allowbreak{}\^{}\{i\}\allowbreak{}\}\allowbreak{}\{(z-w)\allowbreak{}\^{}\{2\}\allowbreak{}\textbackslash{}\allowbreak{},\allowbreak{}(x-z)\allowbreak{}\^{}\{2\}\allowbreak{}\}\allowbreak{}\textbackslash{}\allowbreak{}notag \textbackslash{}\allowbreak{}\textbackslash{}\allowbreak{} \& -\textbackslash{}\allowbreak{}frac\{(x-z)\allowbreak{}\^{}\{j\}\allowbreak{}\textbackslash{}\allowbreak{},\allowbreak{}(x-w)\allowbreak{}\^{}\{i\}\allowbreak{}\}\allowbreak{}\{(x-w)\allowbreak{}\^{}\{2\}\allowbreak{}\textbackslash{}\allowbreak{},\allowbreak{}(x-z)\allowbreak{}\^{}\{2\}\allowbreak{}\}\allowbreak{}-\textbackslash{}\allowbreak{}frac\{(x-w)\allowbreak{}\^{}\{i\}\allowbreak{}\textbackslash{}\allowbreak{},\allowbreak{}(z-w)\allowbreak{}\^{}\{j\}\allowbreak{}\}\allowbreak{}\{(z-w)\allowbreak{}\^{}\{2\}\allowbreak{}\textbackslash{}\allowbreak{},\allowbreak{}(x-w)\allowbreak{}\^{}\{2\}\allowbreak{}\}\allowbreak{}\textbackslash{}\allowbreak{}Bigg]\textbackslash{}\allowbreak{}Bigg\textbackslash{}\allowbreak{}\}\allowbreak{}\textbackslash{}\allowbreak{}times}}\par
\noindent{\small\textsf{Unmarked edit, l.158 (deleted):} \texttt{\textbackslash{}\allowbreak{}Bigg[U\^{}\{ab'\}\allowbreak{}(x')\allowbreak{}\textbackslash{}\allowbreak{},\allowbreak{}U\^{}\{c'd'\}\allowbreak{}(y)\allowbreak{}\textbackslash{}\allowbreak{},\allowbreak{}U\^{}\{c'b\}\allowbreak{}(z)\allowbreak{}\textbackslash{}\allowbreak{},\allowbreak{}\textbackslash{}\allowbreak{}rho\^{}\{b'\}\allowbreak{}(x')\allowbreak{}\textbackslash{}\allowbreak{}rho\^{}\{d'\}\allowbreak{}(y)\allowbreak{}-U\^{}\{d'b\}\allowbreak{}(z)\allowbreak{}\textbackslash{}\allowbreak{},\allowbreak{}U\^{}\{d'e'\}\allowbreak{}(y)\allowbreak{}U\^{}\{ac'\}\allowbreak{}(w')\allowbreak{}\textbackslash{}\allowbreak{},\allowbreak{}\textbackslash{}\allowbreak{}rho\^{}\{c'\}\allowbreak{}(x')\allowbreak{}\textbackslash{}\allowbreak{}rho\^{}\{e'\}\allowbreak{}(y)\allowbreak{}\textbackslash{}\allowbreak{}Bigg]\textbackslash{}\allowbreak{}times}}\par
\noindent{\small\textsf{Unmarked edit, l.158 (deleted):} \texttt{\textbackslash{}\allowbreak{}Bigg[f\^{}\{cbd\}\allowbreak{}\textbackslash{}\allowbreak{},\allowbreak{}U\^{}\{ac\}\allowbreak{}(w)\allowbreak{}\textbackslash{}\allowbreak{},\allowbreak{}\textbackslash{}\allowbreak{}rho\^{}\{d\}\allowbreak{}(x)\allowbreak{}-f\^{}\{acd\}\allowbreak{}\textbackslash{}\allowbreak{},\allowbreak{}}}\par
\needspace{3\baselineskip}\noindent\rule{\linewidth}{0.3pt}\par\noindent\textbf{Lines 162--164} \quad {\small (in \texttt{align}, lines 161--163)}\par
\begin{itemize}[leftmargin=2.2em,itemsep=1pt,topsep=2pt]
\item[\textbf{l.162}] \textsf{\small [missing term added]} \emph{added:} $\displaystyle \big[U^{ab'}(x')-U^{ab'}(w')\big]\,\rho^{b'}(x')\;U^{c'b}(z)\,U^{c'd'}(y)\,\rho^{d'}(y)\,;$
\end{itemize}
\begin{notebox}\textsf{\small Note (highlighted, cancelling rows), l.164:} [multiplied by the second bracket, with $U^{cb}(z)\,U^{c'b}(z)=\delta^{cc'}$ in the second term, the colour factor of this row is]\end{notebox}
\needspace{3\baselineskip}\noindent\rule{\linewidth}{0.3pt}\par\noindent\textbf{Lines 166--168} \quad {\small (in \texttt{align}, lines 165--167)}\par
\begin{itemize}[leftmargin=2.2em,itemsep=1pt,topsep=2pt]
\item[\textbf{l.166}] \textsf{\small [missing term added]} \emph{added:} $\displaystyle \big[U^{ab'}(x')\,\rho^{b'}(x')-U^{ab'}(w')\,\rho^{b'}(x')\big]\,\Big[f^{cbd}\,U^{ac}(w)\,U^{c'b}(z)\,U^{c'd'}(y)\,\rho^{d'}(y)\,\rho^{d}(x)-f^{acd}\,U^{cd'}(y)\,U^{de}(x)\,\rho^{d'}(y)\,\rho^{e}(x)\Big]\,,$
\end{itemize}
\begin{notebox}\textsf{\small Note (highlighted, cancelling rows), l.168:} [which is the colour factor of Group I Row IV after renaming the summed indices there ($d,b,c',e\to b,c',d,d'$ in its first term, $b,c,e,d\to c,d,d',e$ in its second term), with the same order $\rho(x')\rho(y)\rho(x)$. The logarithm and the four terms of the kernel bracket are the same (written in another order), the integration variables and the phase are the same, and the prefactors are $+\frac{ig^{4}}{8\pi^{5}}\frac{1}{k^{+}}$ there and $-\frac{ig^{4}}{8\pi^{5}}\frac{1}{k^{+}}$ here. Hence]\end{notebox}
\needspace{3\baselineskip}\noindent\rule{\linewidth}{0.3pt}\par\noindent\textbf{Lines 170--172} \quad {\small (in \texttt{align}, lines 169--171)}\par
\begin{itemize}[leftmargin=2.2em,itemsep=1pt,topsep=2pt]
\item[\textbf{l.170}] \textsf{\small [missing term added]} \emph{added:} $\displaystyle \frac{d^{3}N_{\rm NLO}}{d^{3}k}\bigg|_{\text{Group I, Row IV}}+\frac{d^{3}N_{\rm NLO}}{d^{3}k}\bigg|_{\text{Group III, Row I}}=0\,.$
\end{itemize}
\begin{notebox}\textsf{\small Note (highlighted, cancelling rows), l.172:} [Numerically, at random points with explicit SU(3) matrices: [Group I Row IV] $+$ [Group III Row I] $=$ $3\times10^{-24}$, $1\times10^{-23}$, $3\times10^{-22}$, $2\times10^{-23}$ against $8.8\times10^{-9}$, $2.9\times10^{-8}$, $3.9\times10^{-7}$, $2.7\times10^{-9}$.]\end{notebox}
\needspace{3\baselineskip}\noindent\rule{\linewidth}{0.3pt}\par\noindent\textbf{Lines 181} \quad \textsf{Group III (with c.c.) $\rangle$ Row III}\par
\begin{notebox}\textsf{\small Note on removed material, l.181:} [Removed here: the author's equation for $\frac{d^{3}N_{1\,\rm NLO}}{d^{3}k}$ (l.197--199 of the author's file), written in the other split (plain $\big[\frac{1}{p^{+}}\big]_{+}$, unsubtracted term with the constant kernel, upper limit $\vee-k^{+}$): it grows like $\log\frac{\vee}{k^{+}}$ (it still contains $\log\frac{\vee}{k^{+}}\times$ the second and third term of this row in the BFKL list), so that $\vee\to\infty$ cannot be taken in it. It is replaced by the form in which the large logarithm is removed completely.]\end{notebox}
\needspace{3\baselineskip}\noindent\rule{\linewidth}{0.3pt}\par\noindent\textbf{Lines 183--186} \quad {\small (in \texttt{align}, lines 182--185)}\par
\begin{itemize}[leftmargin=2.2em,itemsep=1pt,topsep=2pt]
\item[\textbf{l.183}] \textsf{\small [typo/LaTeX]} $\displaystyle \frac{d^{3}N_{1\,\rm NLO}}{d^{3}k}$\ $\longrightarrow$\ $\displaystyle \frac{d^{3}N_{1\,\rm NLO}}{d^{3}k}$
\item[\textbf{l.183}] \textsf{\small [phase convention]} \par\noindent\emph{original:} $\displaystyle \frac{1}{(2\pi)^{3}}\, \frac{ig^{4}\,f^{c'd'a}}{8\pi^{4}}\, \frac{1}{k^{+}}\,\int_{0}^{\vee-k^{+}}\, \frac{dp^{+}}{2\pi}\, \int_{w,x',y',z,y,x}\, e^{-ik\cdot (w-y')}\, e^{-ik\cdot \frac{p^{+}}{k^{+}}\,(z-y')}\,\frac{(y'-z)^{m}\,(x'-y')^{k'}\,(y-w)^{i}\,(x-z)^{j}}{(y'-z)^{2}\,(x'-y')^{2}\,(y-w)^{2}\,(x-z)^{2}}$\par\noindent\emph{checked:} $\displaystyle \frac{1}{(2\pi)^{3}}\, \frac{ig^{4}\,f^{c'd'a}}{8\pi^{4}}\, \frac{1}{k^{+}}\,\int_{0}^{\infty}\, \frac{dp^{+}}{2\pi}\, \int_{w,x',y',z,y,x}\, e^{-ik\cdot (w-y')}\, e^{-ik\cdot \frac{p^{+}}{k^{+}}\,(z-y')}\,\frac{(y'-z)^{m}\,(x'-y')^{k'}\,(y-w)^{i}\,(x-z)^{j}}{(y'-z)^{2}\,(x'-y')^{2}\,(y-w)^{2}\,(x-z)^{2}}$
\item[\textbf{l.183}] \textsf{\small [term removed/replaced]} \par\noindent\emph{original:} $\displaystyle [\frac{\delta_{k'm}\,\delta_{ij}}{(p^{+}+k^{+})}-[\frac{1}{p^{+}}]_{+}\,\delta_{jm}\,\delta_{ik'}-\frac{1}{k^{+}}\, \delta_{im}\, \delta_{jk'}]\, [U^{c'e'}(y')\,U^{d'b}(z)\,\rho^{e'}(x')-U^{d'b}(z)\,U^{c'e'}(x')\,\rho^{e'}(x')]\times$\par\noindent\emph{checked:} $\displaystyle [\frac{\delta_{k'm}\,\delta_{ij}}{(p^{+}+k^{+})}-([\frac{k^{+}}{p^{+}\,(p^{+}+k^{+})}]_{+}+\frac{1}{(p^{+}+k^{+})})\,\delta_{jm}\,\delta_{ik'}]\, [U^{c'e'}(y')\,U^{d'b}(z)\,\rho^{e'}(x')-U^{d'b}(z)\,U^{c'e'}(x')\,\rho^{e'}(x')]\times$
\item[\textbf{l.183}] \textsf{\small [term removed/replaced]} $\displaystyle [U^{ac}(w)\,$\ $\longrightarrow$\ $\displaystyle [U^{ac}(w)\, \{\rho^{c}(y),\rho^{b}(x)\}-U^{cb}(z)\,U^{ad}(y)\,U^{ce}(x)\,\{\rho^{d}(y),\rho^{e}(x)\}]$
\item[\textbf{l.184}] \textsf{\small [phase convention]} \emph{added:} $\displaystyle +\,\frac{1}{(2\pi)^{3}}\, \frac{ig^{4}\,f^{c'd'a}}{8\pi^{4}}\, \frac{1}{k^{+}}\,\int_{0}^{\infty}\, \frac{dp^{+}}{2\pi}\, \int_{w,x',y',z,y,x}\, e^{-ik\cdot (w-y')}\, e^{-ik\cdot \frac{p^{+}}{k^{+}}\,(z-y')}\,\frac{(y'-z)^{m}\,(y-w)^{i}\,(x-z)^{j}}{(y'-z)^{2}\,(y-w)^{2}\,(x-z)^{2}}$
\item[\textbf{l.184}] \textsf{\small [missing term added]} \emph{added:} $\displaystyle [-\frac{1}{k^{+}}\, \delta_{im}\, \delta_{jk'}]\, \Bigg[\frac{(x'-y')^{k'}}{(x'-y')^{2}}\,[U^{c'e'}(y')\,U^{d'b}(z)\,\rho^{e'}(x')-U^{d'b}(z)\,U^{c'e'}(x')\,\rho^{e'}(x')]$
\item[\textbf{l.184}] \textsf{\small [missing term added]} \emph{added:} $\displaystyle -\frac{(x'-z)^{k'}}{(x'-z)^{2}}\,[U^{c'e'}(z)\,U^{d'b}(z)\,\rho^{e'}(x')-U^{d'b}(z)\,U^{c'e'}(x')\,\rho^{e'}(x')]\Bigg]\times$
\item[\textbf{l.184}] \textsf{\small [missing term added]} \emph{added:} $\displaystyle [U^{ac}(w)\, \{\rho^{c}(y),\rho^{b}(x)\}-U^{cb}(z)\,U^{ad}(y)\,U^{ce}(x)\,\{\rho^{d}(y),\rho^{e}(x)\}]$
\end{itemize}
\begin{notebox}\textsf{\small Note, l.186:} [This is the first part of Row III of \texttt{Other\_terms} (v2) at $\vee\to\infty$: upper limit of the $p^{+}$ integrals $\infty$, its last term ($\log\big(1-\frac{k^{+}}{\vee}\big)\times\dots$) dropped; nothing else is changed. All three terms of this row in the BFKL list are removed from it (verified for the v2 form against the row with the cut-offs). \textbf{Condition for the limit}: in the term with the constant kernel (the one whose square bracket contains the same colour bracket at $y'$ and at $y'\to z$) the $y'$ integral has to be done before the $p^{+}$ integral, with the two terms of that bracket together. At fixed coordinates $\int_{0}^{P}dp^{+}e^{-i\frac{p^{+}}{k^{+}}s}=\frac{k^{+}}{is}\big(1-e^{-i\frac{P}{k^{+}}s}\big)$, $s=k\cdot(z-y')$, has no limit for $P\to\infty$; after the $y'$ integral the integrand is ${\cal O}(k^{+2}/p^{+2})$. The $y'$ integral of the subtracted piece converges only conditionally at large $y'$ (to be integrated symmetrically; Gaussian regulator in the check). \textbf{Numerical check} (smooth test function of $y'$ in place of the first colour bracket, second bracket $\to1$, prefactor $1$, $|k|=1$, $p^{+}$ and $y'$ integrals by quadrature): this form with the upper limit $\vee-k^{+}$, for $\frac{\vee}{k^{+}}=10,\,30,\,60,\,100$: $0.0201-1.0301\,i,\ 0.0389-1.1011\,i,\ 0.0419-1.1183\,i,\ 0.0430-1.1250\,i$; it converges like $\frac{k^{+}}{\vee}$: extrapolated with $a+b\,\frac{k^{+}}{\vee}$ from the last two values, limit $0.0445-1.1352\,i$, and $\frac{\vee}{k^{+}}\times|{\rm value}-{\rm limit}|=1.08,\ 1.04,\ 1.03$ at $\frac{\vee}{k^{+}}=10,30,60$. The removed form at the same points: $1.048-0.434\,i,\ 1.557-0.220\,i,\ 1.869-0.058\,i,\ 2.098+0.068\,i$ (grows like $\log\frac{\vee}{k^{+}}$). Second set of points, $\frac{\vee}{k^{+}}=10,\,30,\,60$: $-1.5656-0.3693\,i,\ -1.7886-0.4655\,i,\ -1.8452-0.4835\,i$ (successive differences $0.243,\ 0.059$, ratio $4.1$; $4$ for $\frac{k^{+}}{\vee}$). The subtracted term alone, with an analytic test function: $|J(\vee)-J(\infty)|=0.64$, $7.0\times10^{-2}$, $7.0\times10^{-3}$, $7.0\times10^{-4}$ for $\frac{\vee}{k^{+}}-1=10,\,10^{2},\,10^{3},\,10^{4}$, and $J(\infty)$ agrees with the form in which the $y'$ integral is done first to $7\times10^{-8}$. The plus prescription acts on the phase $e^{-ik\cdot\frac{p^{+}}{k^{+}}(z-y')}$. The soft subtraction terms of the plus prescriptions in this equation are finite only in the sum over Group I Rows III, V and Group III Rows III, V, VI under one integral.]\end{notebox}
\noindent{\small\textsf{Unmarked edit, l.183 (deleted):} \texttt{\textbackslash{}\allowbreak{}\{\textbackslash{}\allowbreak{}rho\^{}\{c\}\allowbreak{}(y)\allowbreak{},\allowbreak{}\textbackslash{}\allowbreak{}rho\^{}\{b\}\allowbreak{}(x)\allowbreak{}\textbackslash{}\allowbreak{}\}\allowbreak{}-U\^{}\{cb\}\allowbreak{}(z)\allowbreak{}\textbackslash{}\allowbreak{},\allowbreak{}U\^{}\{ad\}\allowbreak{}(y)\allowbreak{}\textbackslash{}\allowbreak{},\allowbreak{}U\^{}\{ce\}\allowbreak{}(x)\allowbreak{}\textbackslash{}\allowbreak{},\allowbreak{}\textbackslash{}\allowbreak{}\{\textbackslash{}\allowbreak{}rho\^{}\{d\}\allowbreak{}(y)\allowbreak{},\allowbreak{}\textbackslash{}\allowbreak{}rho\^{}\{e\}\allowbreak{}(x)\allowbreak{}\textbackslash{}\allowbreak{}\}\allowbreak{}\textbackslash{}\allowbreak{}right]}}\par
\needspace{3\baselineskip}\noindent\rule{\linewidth}{0.3pt}\par\noindent\textbf{Lines 188}\par
\begin{notebox}\textsf{\small Note on removed material, l.188:} [Removed here: the author's equation for $\frac{d^{3}N_{2\,\rm NLO}}{d^{3}k}$ (l.201--204 of the author's file), written in the other split (plain $\big[\frac{1}{p^{+}}\big]_{+}$, unsubtracted term with the constant kernel, upper limit $\vee-k^{+}$): it grows like $\log\frac{\vee}{k^{+}}$ (it still contains $\log\frac{\vee}{k^{+}}\times$ the second and third term of this row in the BFKL list), so that $\vee\to\infty$ cannot be taken in it. It is replaced by the form in which the large logarithm is removed completely.]\end{notebox}
\needspace{3\baselineskip}\noindent\rule{\linewidth}{0.3pt}\par\noindent\textbf{Lines 190--197} \quad {\small (in \texttt{align}, lines 189--196)}\par
\begin{itemize}[leftmargin=2.2em,itemsep=1pt,topsep=2pt]
\item[\textbf{l.190}] \textsf{\small [typo/LaTeX]} $\displaystyle \frac{d^{3}N_{2\,\rm NLO}}{d^{3}k}$\ $\longrightarrow$\ $\displaystyle \frac{d^{3}N_{2\,\rm NLO}}{d^{3}k}$
\item[\textbf{l.190}] \textsf{\small [phase convention]} \par\noindent\emph{original:} $\displaystyle \frac{1}{(2\pi)^{3}}\,\frac{g^{4}}{4\pi^{4}}\,\int_{0}^{\vee-k^{+}}\, \frac{dp^{+}}{2\pi}\,\int_{w,x',y',z,x}\, e^{-ik\cdot (w-y')}\,e^{-ik\cdot \frac{p^{+}}{k^{+}}(z-y')}\, \frac{(y'-z)^{m}\,(x'-y')^{k'}\,}{(y'-z)^{2}\,(x'-y')^{2}\,}\, \frac{1}{[p^{+}\,(x-z)^{2}+k^{+}\,(x-w)^{2}]}$\par\noindent\emph{checked:} $\displaystyle \frac{1}{(2\pi)^{3}}\,\frac{g^{4}}{4\pi^{4}}\,\int_{0}^{\infty}\, \frac{dp^{+}}{2\pi}\,\int_{w,x',y',z,x}\, e^{-ik\cdot (w-y')}\,e^{-ik\cdot \frac{p^{+}}{k^{+}}(z-y')}\, \frac{(y'-z)^{m}\,(x'-y')^{k'}\,}{(y'-z)^{2}\,(x'-y')^{2}\,}\, \frac{1}{[p^{+}\,(x-z)^{2}+k^{+}\,(x-w)^{2}]}$
\item[\textbf{l.190}] \textsf{\small [typo/LaTeX]} \par\noindent\emph{original:} $\displaystyle \Bigg[\frac{\delta_{k'm}\,d_{\perp}\,p^{+}}{2\,(z-w)^{2}\,(p^{+}+k^{+})^{2}}\, [(x-z)^{2}-(x-w)^{2}]+\frac{p^{+}\,\delta_{k'm}}{k^{+}\,(p^{+}+k^{+})}\,\bigg[\frac{(x-z)\cdot (z-w)}{(z-w)^{2}} -\frac{(x-z)\cdot (x-w)}{2\,(x-w)^{2}}$\par\noindent\emph{checked:} $\displaystyle \Bigg[\frac{\delta_{k'm}\,d_{\perp}\,p^{+}}{2\,(z-w)^{2}\,(p^{+}+k^{+})^{2}}\, [(x-z)^{2}-(x-w)^{2}]+\frac{p^{+}\,\delta_{k'm}}{k^{+}\,(p^{+}+k^{+})}\,\bigg[\frac{(x-z)\cdot (z-w)}{(z-w)^{2}} -\frac{(x-z)\cdot (x-w)}{2\,(x-w)^{2}}$
\item[\textbf{l.190}] \textsf{\small [typo/LaTeX]} \par\noindent\emph{original:} $\displaystyle -\frac{[(x-z)^{2}-(x-w)^{2}]}{2\,(z-w)^{2}}\bigg] +\frac{\delta_{k'm}}{(p^{+}+k^{+})}\,\bigg[\frac{(x-w)\cdot (z-w)}{(z-w)^{2}}+\frac{(x-z)\cdot (x-w)}{2\,(x-z)^{2}} -\frac{[(x-z)^{2}-(x-w)^{2}]}{2\,(z-w)^{2}}\bigg]$\par\noindent\emph{checked:} $\displaystyle -\frac{[(x-z)^{2}-(x-w)^{2}]}{2\,(z-w)^{2}}\bigg] +\frac{\delta_{k'm}}{(p^{+}+k^{+})}\,\bigg[\frac{(x-w)\cdot (z-w)}{(z-w)^{2}}+\frac{(x-z)\cdot (x-w)}{2\,(x-z)^{2}} -\frac{[(x-z)^{2}-(x-w)^{2}]}{2\,(z-w)^{2}}\bigg]$
\item[\textbf{l.190}] \textsf{\small [term removed/replaced]} \par\noindent\emph{original:} $\displaystyle -[\frac{1}{p^{+}}]_{+}\,\bigg[\frac{(x-w)^{k'}\,(z-w)^{m}}{(z-w)^{2}}+\frac{(x-z)^{m}\,(x-w)^{k'}}{2\,(x-z)^{2}}\bigg]-\frac{p^{+}}{k^{+\,2}}\,\bigg[\frac{(x-z)^{k'}\,(z-w)^{m}}{(z-w)^{2}} -\,\frac{(x-z)^{k'}\,(x-w)^{m}}{2\,(x-w)^{2}}\bigg]$\par\noindent\emph{checked:} $\displaystyle -([\frac{k^{+}}{p^{+}\,(p^{+}+k^{+})}]_{+}+\frac{1}{(p^{+}+k^{+})})\,\bigg[\frac{(x-w)^{k'}\,(z-w)^{m}}{(z-w)^{2}}+\frac{(x-z)^{m}\,(x-w)^{k'}}{2\,(x-z)^{2}}\bigg]$
\item[\textbf{l.190}] \textsf{\small [typo/LaTeX]} \par\noindent\emph{original:} $\displaystyle -\frac{1}{k^{+}}\,\bigg[\frac{(x-z)^{m}\,(z-w)^{k'}}{(z-w)^{2}}-\frac{(x-z)^{m}\,(x-w)^{k'}}{2\,(x-w)^{2}}+\frac{(x-w)^{m}\,(z-w)^{k'}}{(z-w)^{2}} +\frac{(x-z)^{k'}\,(x-w)^{m}}{2\,(x-z)^{2}}\bigg]\Bigg]\times\,$\par\noindent\emph{checked:} $\displaystyle -\frac{1}{k^{+}}\,\bigg[\frac{(x-z)^{m}\,(z-w)^{k'}}{(z-w)^{2}}-\frac{(x-z)^{m}\,(x-w)^{k'}}{2\,(x-w)^{2}}+\frac{(x-w)^{m}\,(z-w)^{k'}}{(z-w)^{2}} +\frac{(x-z)^{k'}\,(x-w)^{m}}{2\,(x-z)^{2}}\bigg]\Bigg]\times\,$
\item[\textbf{l.190}] \textsf{\small [term removed/replaced]} \par\noindent\emph{original:} $\displaystyle [f^{c'd'a}\,U^{c'e'}(y')\,U^{d'b}(z)\,\rho^{e'}(x')-f^{c'd'a}\,$\par\noindent\emph{checked:} $\displaystyle [f^{c'd'a}\,U^{c'e'}(y')\,U^{d'b}(z)\,\rho^{e'}(x')-f^{c'd'a}\, U^{d'b}(z)\,U^{c'e'}(x')\,\rho^{e'}(x')]\times [f^{cbd}\,U^{ac}(w)\,\rho^{d}(x) -f^{acd}\,U^{cb}(z)\,U^{de}(x)\,\rho^{e}(x)]$
\item[\textbf{l.191}] \textsf{\small [phase convention]} \emph{added:} $\displaystyle +\,\frac{1}{(2\pi)^{3}}\,\frac{g^{4}}{4\pi^{4}}\,\int_{0}^{\infty}\, \frac{dp^{+}}{2\pi}\,\int_{w,x',y',z,x}\, e^{-ik\cdot (w-y')}\,e^{-ik\cdot \frac{p^{+}}{k^{+}}(z-y')}\, \frac{(y'-z)^{m}}{(y'-z)^{2}}\, \frac{1}{[p^{+}\,(x-z)^{2}+k^{+}\,(x-w)^{2}]}$
\item[\textbf{l.191}] \textsf{\small [missing term added]} \emph{added:} $\displaystyle \Bigg[ -\frac{p^{+}}{k^{+\,2}}\,\bigg[\frac{(x-z)^{k'}\,(z-w)^{m}}{(z-w)^{2}} -\,\frac{(x-z)^{k'}\,(x-w)^{m}}{2\,(x-w)^{2}}\bigg] \Bigg]\times\,$
\item[\textbf{l.191}] \textsf{\small [missing term added]} \emph{added:} $\displaystyle \Bigg[\frac{(x'-y')^{k'}}{(x'-y')^{2}}\,[f^{c'd'a}\,U^{c'e'}(y')\,U^{d'b}(z)\,\rho^{e'}(x')-f^{c'd'a}\, U^{d'b}(z)\,U^{c'e'}(x')\,\rho^{e'}(x')]$
\item[\textbf{l.192}] \textsf{\small [missing term added]} \emph{added:} $\displaystyle -\frac{(x'-z)^{k'}}{(x'-z)^{2}}\,[f^{c'd'a}\,U^{c'e'}(z)\,U^{d'b}(z)\,\rho^{e'}(x')-f^{c'd'a}\, U^{d'b}(z)\,U^{c'e'}(x')\,\rho^{e'}(x')]\Bigg]\times$
\item[\textbf{l.193}] \textsf{\small [missing term added]} \emph{added:} $\displaystyle [f^{cbd}\,U^{ac}(w)\,\rho^{d}(x) -f^{acd}\,U^{cb}(z)\,U^{de}(x)\,\rho^{e}(x)]$
\item[\textbf{l.194}] \textsf{\small [phase convention]} \emph{added:} $\displaystyle +\,\frac{1}{(2\pi)^{3}}\,\frac{g^{4}}{4\pi^{4}}\,\int_{0}^{\infty}\, \frac{dp^{+}}{2\pi}\,\int_{w,x',w',z,x}\, e^{-ik\cdot (w-w')}\, \frac{(w'-z)^{m}\,(x'-z)^{k'}}{(w'-z)^{2}\,(x'-z)^{2}}\,$
\item[\textbf{l.194}] \textsf{\small [term removed/replaced]} \par\noindent\emph{original:} $\displaystyle -f^{acd}\,U^{cb}(z)\,U^{de}(x)\,\rho^{e}(x)]$\par\noindent\emph{checked:} $\displaystyle \Bigg[ \frac{(x-w)^{2}}{(x-z)^{2}\,(p^{+}+k^{+})\,[p^{+}\,(x-z)^{2}+k^{+}\,(x-w)^{2}]}\,\bigg[\frac{(x-z)^{k'}\,(z-w)^{m}}{(z-w)^{2}} -\,\frac{(x-z)^{k'}\,(x-w)^{m}}{2\,(x-w)^{2}}\bigg] \Bigg]\times\,$
\item[\textbf{l.194}] \textsf{\small [term removed/replaced]} \par\noindent\emph{original:} $\displaystyle -f^{acd}\,U^{cb}(z)\,U^{de}(x)\,\rho^{e}(x)]$\par\noindent\emph{checked:} $\displaystyle [f^{c'd'a}\,U^{c'e'}(z)\,U^{d'b}(z)\,\rho^{e'}(x')-f^{c'd'a}\, U^{d'b}(z)\,U^{c'e'}(x')\,\rho^{e'}(x')]\times [f^{cbd}\,U^{ac}(w)\,\rho^{d}(x) -f^{acd}\,U^{cb}(z)\,U^{de}(x)\,\rho^{e}(x)]$
\end{itemize}
\begin{notebox}\textsf{\small Note, l.197:} [This is the second part of Row III of \texttt{Other\_terms} (v2) at $\vee\to\infty$: upper limit $\infty$, last term ($\log\big(1-\frac{k^{+}}{\vee}\big)\times\dots$) dropped; nothing else is changed. The term $-\frac{p^{+}}{k^{+\,2}}[\dots]$, whose kernel tends to the constant $-\frac{1}{k^{+\,2}(x-z)^{2}}$ at large $p^{+}$, is the one that is subtracted at $y'\to z$ (second term); the third term is what remains of the piece added back, and its $p^{+}$ integral can be done: $\int_{0}^{\infty}\frac{dp^{+}}{2\pi}\frac{(x-w)^{2}}{(x-z)^{2}(p^{+}+k^{+})[p^{+}(x-z)^{2}+k^{+}(x-w)^{2}]}=\frac{1}{2\pi k^{+}}\frac{(x-w)^{2}}{(x-z)^{2}[(x-w)^{2}-(x-z)^{2}]}\log\frac{(x-w)^{2}}{(x-z)^{2}}$ (checked by quadrature). \textbf{Condition for the limit}: in the term with the constant kernel (the one whose square bracket contains the same colour bracket at $y'$ and at $y'\to z$) the $y'$ integral has to be done before the $p^{+}$ integral, with the two terms of that bracket together. At fixed coordinates $\int_{0}^{P}dp^{+}e^{-i\frac{p^{+}}{k^{+}}s}=\frac{k^{+}}{is}\big(1-e^{-i\frac{P}{k^{+}}s}\big)$, $s=k\cdot(z-y')$, has no limit for $P\to\infty$; after the $y'$ integral the integrand is ${\cal O}(k^{+2}/p^{+2})$. The $y'$ integral of the subtracted piece converges only conditionally at large $y'$ (to be integrated symmetrically; Gaussian regulator in the check). \textbf{Numerical check} (as in the first part): this form with the upper limit $\vee-k^{+}$, for $\frac{\vee}{k^{+}}=10,\,30,\,60,\,100$: $-0.0960-0.2091\,i,\ -0.1086-0.2357\,i,\ -0.1127-0.2427\,i,\ -0.1144-0.2455\,i$; it converges like $\frac{k^{+}}{\vee}$: extrapolated limit $-0.1170-0.2497\,i$, $\frac{\vee}{k^{+}}\times|{\rm value}-{\rm limit}|=0.46,\ 0.49,\ 0.49$ at $\frac{\vee}{k^{+}}=10,30,60$. The removed form: $-0.0732-0.0467\,i,\ -0.0749+0.0042\,i,\ -0.0720+0.0462\,i,\ -0.0687+0.0794\,i$. Second set of points, $\frac{\vee}{k^{+}}=10,\,30,\,60$: $0.4870+0.0491\,i,\ 0.5728+0.0643\,i,\ 0.5966+0.0625\,i$ (successive differences $0.087,\ 0.024$, ratio $3.7$). The plus prescription acts on the phase and on $1/[p^{+}(x-z)^{2}+k^{+}(x-w)^{2}]$. The soft subtraction terms of the plus prescriptions in this equation are finite only in the sum over Group I Rows III, V and Group III Rows III, V, VI under one integral.]\end{notebox}
\noindent{\small\textsf{Unmarked edit, l.194 (deleted):} \texttt{U\^{}\{d'b\}\allowbreak{}(z)\allowbreak{}\textbackslash{}\allowbreak{},\allowbreak{}U\^{}\{c'e'\}\allowbreak{}(x')\allowbreak{}\textbackslash{}\allowbreak{},\allowbreak{}\textbackslash{}\allowbreak{}rho\^{}\{e'\}\allowbreak{}(x')\allowbreak{}\textbackslash{}\allowbreak{}right]\textbackslash{}\allowbreak{}times}}\par
\noindent{\small\textsf{Unmarked edit, l.194 (deleted):} \texttt{\textbackslash{}\allowbreak{}left[f\^{}\{cbd\}\allowbreak{}\textbackslash{}\allowbreak{},\allowbreak{}U\^{}\{ac\}\allowbreak{}(w)\allowbreak{}\textbackslash{}\allowbreak{},\allowbreak{}\textbackslash{}\allowbreak{}rho\^{}\{d\}\allowbreak{}(x)\allowbreak{}}}\par
\needspace{3\baselineskip}\noindent\rule{\linewidth}{0.3pt}\par\noindent\textbf{Lines 201} \quad \textsf{Group III (with c.c.) $\rangle$ Row V}\par
\begin{notebox}\textsf{\small Note on removed material, l.201:} [Removed here: the author's equation for $\frac{d^{3}N_{\rm NLO}}{d^{3}k}$ (l.208--210 of the author's file), written in the other split (plain $\big[\frac{1}{p^{+}}\big]_{+}$, unsubtracted term with the constant kernel, upper limit $\vee-k^{+}$): it grows like $\log\frac{\vee}{k^{+}}$ (it still contains $\log\frac{\vee}{k^{+}}\times$ the second and third term of this row in the BFKL list), so that $\vee\to\infty$ cannot be taken in it. It is replaced by the form in which the large logarithm is removed completely.]\end{notebox}
\needspace{3\baselineskip}\noindent\rule{\linewidth}{0.3pt}\par\noindent\textbf{Lines 203--206} \quad {\small (in \texttt{align}, lines 202--205)}\par
\begin{itemize}[leftmargin=2.2em,itemsep=1pt,topsep=2pt]
\item[\textbf{l.203}] \textsf{\small [typo/LaTeX]} $\displaystyle \frac{d^{3}N_{\rm NLO}}{d^{3}k}$\ $\longrightarrow$\ $\displaystyle \frac{d^{3}N_{\rm NLO}}{d^{3}k}$
\item[\textbf{l.203}] \textsf{\small [phase convention]} \emph{added:} $\displaystyle \frac{1}{(2\pi)^{3}}\, \frac{ig^{4}\,f^{c'ca}}{4\pi^{4}}\, \frac{1}{k^{+}}\,\int_{0}^{\infty}\, \frac{dp^{+}}{2\pi}\, \int_{w,y',z,y,x,x'}\, e^{-ik\cdot (w-y')}\, e^{-ik\cdot \frac{p^{+}}{k^{+}}\,(z-y')}\,\frac{(y'-z)^{m}\,(x'-y')^{k'}\,(y-w)^{i}\,(x-z)^{j}}{(y'-z)^{2}\,(x'-y')^{2}\,(y-w)^{2}\,(x-z)^{2}}$
\item[\textbf{l.203}] \textsf{\small [missing term added]} \emph{added:} $\displaystyle [\frac{\delta_{k'm}\,\delta_{ij}}{(p^{+}+k^{+})}-([\frac{k^{+}}{p^{+}\,(p^{+}+k^{+})}]_{+}+\frac{1}{(p^{+}+k^{+})})\,\delta_{jm}\,\delta_{ik'}]\, [U^{c'e'}(y')\,\rho^{e'}(x')-U^{c'e'}(x')\,\rho^{e'}(x')]\times$
\item[\textbf{l.203}] \textsf{\small [typo/LaTeX]} $\displaystyle [U^{ce}(x)\,U^{ad}(y)\,\rho^{e}(x)\rho^{d}(y)-U^{ad}(w)\,U^{ce}(x)\,\rho^{e}(x)\rho^{d}(y)]$\ $\longrightarrow$\ $\displaystyle [U^{ce}(x)\,U^{ad}(y)\,\rho^{e}(x)\rho^{d}(y)-U^{ad}(w)\,U^{ce}(x)\,\rho^{e}(x)\rho^{d}(y)]$
\item[\textbf{l.204}] \textsf{\small [phase convention]} \emph{added:} $\displaystyle +\,\frac{1}{(2\pi)^{3}}\, \frac{ig^{4}\,f^{c'ca}}{4\pi^{4}}\, \frac{1}{k^{+}}\,\int_{0}^{\infty}\, \frac{dp^{+}}{2\pi}\, \int_{w,y',z,y,x,x'}\, e^{-ik\cdot (w-y')}\, e^{-ik\cdot \frac{p^{+}}{k^{+}}\,(z-y')}\,\frac{(y'-z)^{m}\,(y-w)^{i}\,(x-z)^{j}}{(y'-z)^{2}\,(y-w)^{2}\,(x-z)^{2}}$
\item[\textbf{l.204}] \textsf{\small [missing term added]} \emph{added:} $\displaystyle [-\frac{1}{k^{+}}\, \delta_{im}\, \delta_{jk'}]\, \Bigg[\frac{(x'-y')^{k'}}{(x'-y')^{2}}\,[U^{c'e'}(y')\,\rho^{e'}(x')-U^{c'e'}(x')\,\rho^{e'}(x')]$
\item[\textbf{l.204}] \textsf{\small [missing term added]} \emph{added:} $\displaystyle -\frac{(x'-z)^{k'}}{(x'-z)^{2}}\,[U^{c'e'}(z)\,\rho^{e'}(x')-U^{c'e'}(x')\,\rho^{e'}(x')]\Bigg]\times$
\item[\textbf{l.204}] \textsf{\small [missing term added]} \emph{added:} $\displaystyle [U^{ce}(x)\,U^{ad}(y)\,\rho^{e}(x)\rho^{d}(y)-U^{ad}(w)\,U^{ce}(x)\,\rho^{e}(x)\rho^{d}(y)]$
\end{itemize}
\begin{notebox}\textsf{\small Note, l.206:} [This is Row V of \texttt{Other\_terms} (v2) at $\vee\to\infty$: upper limit $\infty$, last term ($\log\big(1-\frac{k^{+}}{\vee}\big)\times\dots$) dropped; nothing else is changed. Kernels and $p^{+}$ dependence are those of Row III, first part, with twice the prefactor: the condition for the limit ($y'$ integral before the $p^{+}$ integral in the second term; conditional convergence) and the numerical check are the same (numbers of Row III, first part, $\times2$). The soft subtraction terms of the plus prescriptions in this equation are finite only in the sum over Group I Rows III, V and Group III Rows III, V, VI under one integral.]\end{notebox}
\noindent{\small\textsf{Unmarked edit, l.203 (deleted):} \texttt{ \textbackslash{}\allowbreak{}frac\{1\}\allowbreak{}\{(2\textbackslash{}\allowbreak{}pi)\allowbreak{}\^{}\{3\}\allowbreak{}\}\allowbreak{}\textbackslash{}\allowbreak{},\allowbreak{} \textbackslash{}\allowbreak{}frac\{ig\^{}\{4\}\allowbreak{}\textbackslash{}\allowbreak{},\allowbreak{}f\^{}\{c'ca\}\allowbreak{}\}\allowbreak{}\{4\textbackslash{}\allowbreak{}pi\^{}\{4\}\allowbreak{}\}\allowbreak{}\textbackslash{}\allowbreak{},\allowbreak{} \textbackslash{}\allowbreak{}frac\{1\}\allowbreak{}\{k\^{}\{+\allowbreak{}\}\allowbreak{}\}\allowbreak{}\textbackslash{}\allowbreak{},\allowbreak{}\textbackslash{}\allowbreak{}int\_\{0\}\allowbreak{}\^{}\{\textbackslash{}\allowbreak{}vee-k\^{}\{+\allowbreak{}\}\allowbreak{}\}\allowbreak{}\textbackslash{}\allowbreak{},\allowbreak{} \textbackslash{}\allowbreak{}frac\{dp\^{}\{+\allowbreak{}\}\allowbreak{}\}\allowbreak{}\{2\textbackslash{}\allowbreak{}pi\}\allowbreak{}\textbackslash{}\allowbreak{},\allowbreak{} \textbackslash{}\allowbreak{}int\_\{w,\allowbreak{}y',\allowbreak{}z,\allowbreak{}y,\allowbreak{}x,\allowbreak{}x'\}\allowbreak{}\textbackslash{}\allowbreak{},\allowbreak{} e\^{}\{-ik\textbackslash{}\allowbreak{}cdot (w-y')\allowbreak{}\}\allowbreak{}\textbackslash{}\allowbreak{},\allowbreak{} e\^{}\{-ik\textbackslash{}\allowbreak{}cdot \textbackslash{}\allowbreak{}frac\{p\^{}\{+\allowbreak{}\}\allowbreak{}\}\allowbreak{}\{k\^{}\{+\allowbreak{}\}\allowbreak{}\}\allowbreak{}\textbackslash{}\allowbreak{},\allowbreak{}(z-y')\allowbreak{}\}\allowbreak{}\textbackslash{}\allowbreak{},\allowbreak{} \textbackslash{}\allowbreak{}frac\{(y'-z)\allowbreak{}\^{}\{m\}\allowbreak{}\textbackslash{}\allowbreak{},\allowbreak{}(x'-y')\allowbreak{}\^{}\{k'\}\allowbreak{}\textbackslash{}\allowbreak{},\allowbreak{}(y-w)\allowbreak{}\^{}\{i\}\allowbreak{}\textbackslash{}\allowbreak{},\allowbreak{}(x-z)\allowbreak{}\^{}\{j\}\allowbreak{}\}\allowbreak{}\{(y'-z)\allowbreak{}\^{}\{2\}\allowbreak{}\textbackslash{}\allowbreak{},\allowbreak{}(x'-y')\allowbreak{}\^{}\{2\}\allowbreak{}\textbackslash{}\allowbreak{},\allowbreak{}(y-w)\allowbreak{}\^{}\{2\}\allowbreak{}\textbackslash{}\allowbreak{},\allowbreak{}(x-z)\allowbreak{}\^{}\{2\}\allowbreak{}\}\allowbreak{}\textbackslash{}\allowbreak{}notag \textbackslash{}\allowbreak{}\textbackslash{}\allowbreak{} \& \textbackslash{}\allowbreak{}left[\textbackslash{}\allowbreak{}frac\{\textbackslash{}\allowbreak{}delta\_\{k'm\}\allowbreak{}\textbackslash{}\allowbreak{},\allowbreak{}\textbackslash{}\allowbreak{}delta\_\{ij\}\allowbreak{}\}\allowbreak{}\{(p\^{}\{+\allowbreak{}\}\allowbreak{}+\allowbreak{}k\^{}\{+\allowbreak{}\}\allowbreak{})\allowbreak{}\}\allowbreak{}-\textbackslash{}\allowbreak{}left[\textbackslash{}\allowbreak{}frac\{1\}\allowbreak{}\{p\^{}\{+\allowbreak{}\}\allowbreak{}\}\allowbreak{}\textbackslash{}\allowbreak{}r}}\par
\needspace{3\baselineskip}\noindent\rule{\linewidth}{0.3pt}\par\noindent\textbf{Lines 208} \quad \textsf{Group III (with c.c.) $\rangle$ Row VI}\par
\begin{notebox}\textsf{\small Note on removed material, l.208:} [Removed here: the author's equation for $\frac{d^{3}N_{\rm NLO}}{d^{3}k}$ (l.212--214 of the author's file), written in the other split (plain $\big[\frac{1}{p^{+}}\big]_{+}$, unsubtracted term with the constant kernel, upper limit $\vee-k^{+}$): it grows like $\log\frac{\vee}{k^{+}}$ (it still contains $\log\frac{\vee}{k^{+}}\times$ the second and third term of this row in the BFKL list), so that $\vee\to\infty$ cannot be taken in it. It is replaced by the form in which the large logarithm is removed completely.]\end{notebox}
\needspace{3\baselineskip}\noindent\rule{\linewidth}{0.3pt}\par\noindent\textbf{Lines 210--212} \quad {\small (in \texttt{align}, lines 209--211)}\par
\begin{itemize}[leftmargin=2.2em,itemsep=1pt,topsep=2pt]
\item[\textbf{l.210}] \textsf{\small [typo/LaTeX]} $\displaystyle \frac{d^{3}N_{\rm NLO}}{d^{3}k}$\ $\longrightarrow$\ $\displaystyle \frac{d^{3}N_{\rm NLO}}{d^{3}k}$
\item[\textbf{l.210}] \textsf{\small [phase convention]} \par\noindent\emph{original:} $\displaystyle \frac{1}{(2\pi)^{3}}\, \frac{ig^{4}\,f^{c'd'a}}{4\pi^{4}}\, \frac{1}{k^{+}}\,\int_{0}^{\vee-k^{+}}\, \frac{dp^{+}}{2\pi}\, \int_{w,y',z,y,x,x'}\, e^{-ik\cdot (w-y')}\, e^{-ik\cdot \frac{p^{+}}{k^{+}}\,(z-y')}\, \frac{(y'-z)^{m}\,(x'-y')^{k'}\,(y-w)^{i}\,(x-z)^{j}}{(y'-z)^{2}\,(x'-y')^{2}\,(y-w)^{2}\,(x-z)^{2}}$\par\noindent\emph{checked:} $\displaystyle \frac{1}{(2\pi)^{3}}\, \frac{ig^{4}\,f^{c'd'a}}{4\pi^{4}}\, \frac{1}{k^{+}}\,\int_{0}^{\infty}\, \frac{dp^{+}}{2\pi}\, \int_{w,y',z,y,x,x'}\, e^{-ik\cdot (w-y')}\, e^{-ik\cdot \frac{p^{+}}{k^{+}}\,(z-y')}\,\frac{(y'-z)^{m}\,(x'-y')^{k'}\,(y-w)^{i}\,(x-z)^{j}}{(y'-z)^{2}\,(x'-y')^{2}\,(y-w)^{2}\,(x-z)^{2}}$
\item[\textbf{l.210}] \textsf{\small [term removed/replaced]} \par\noindent\emph{original:} $\displaystyle [\frac{\delta_{k'm}\,\delta_{ij}}{(p^{+}+k^{+})}-[\frac{1}{p^{+}}]_{+}\,\delta_{jm}\,\delta_{ik'}-\frac{1}{k^{+}}\, \delta_{im}\, \delta_{jk'}]\, [U^{c'e'}(y')\,U^{d'b}(z)\,\rho^{e'}(x')-U^{d'b}(z)\,U^{c'e'}(x')\,\rho^{e'}(x')]\times$\par\noindent\emph{checked:} $\displaystyle [\frac{\delta_{k'm}\,\delta_{ij}}{(p^{+}+k^{+})}-([\frac{k^{+}}{p^{+}\,(p^{+}+k^{+})}]_{+}+\frac{1}{(p^{+}+k^{+})})\,\delta_{jm}\,\delta_{ik'}]\, [U^{c'e'}(y')\,U^{d'b}(z)\,\rho^{e'}(x')-U^{d'b}(z)\,U^{c'e'}(x')\,\rho^{e'}(x')]\times$
\item[\textbf{l.210}] \textsf{\small [typo/LaTeX]} $\displaystyle [U^{cb}(z)\,U^{ad}(y)\,U^{ce}(x)\,\rho^{d}(y)\,\rho^{e}(x)-U^{ad}(y)\,\rho^{d}(y)\,\rho^{b}(x)]$\ $\longrightarrow$\ $\displaystyle [U^{cb}(z)\,U^{ad}(y)\,U^{ce}(x)\,\rho^{d}(y)\,\rho^{e}(x)-U^{ad}(y)\,\rho^{d}(y)\,\rho^{b}(x)]$
\item[\textbf{l.210}] \textsf{\small [phase convention]} \par\noindent\emph{original:} $\displaystyle [U^{cb}(z)\,U^{ad}(y)\,U^{ce}(x)\,\rho^{d}(y)\,\rho^{e}(x)-U^{ad}(y)\,\rho^{d}(y)\,\rho^{b}(x)]$\par\noindent\emph{checked:} $\displaystyle +\,\frac{1}{(2\pi)^{3}}\, \frac{ig^{4}\,f^{c'd'a}}{4\pi^{4}}\, \frac{1}{k^{+}}\,\int_{0}^{\infty}\, \frac{dp^{+}}{2\pi}\, \int_{w,y',z,y,x,x'}\, e^{-ik\cdot (w-y')}\, e^{-ik\cdot \frac{p^{+}}{k^{+}}\,(z-y')}\,\frac{(y'-z)^{m}\,(y-w)^{i}\,(x-z)^{j}}{(y'-z)^{2}\,(y-w)^{2}\,(x-z)^{2}}$
\item[\textbf{l.210}] \textsf{\small [term removed/replaced]} \par\noindent\emph{original:} $\displaystyle [U^{cb}(z)\,U^{ad}(y)\,U^{ce}(x)\,\rho^{d}(y)\,\rho^{e}(x)-U^{ad}(y)\,\rho^{d}(y)\,\rho^{b}(x)]$\par\noindent\emph{checked:} $\displaystyle [-\frac{1}{k^{+}}\, \delta_{im}\, \delta_{jk'}]\, \Bigg[\frac{(x'-y')^{k'}}{(x'-y')^{2}}\,[U^{c'e'}(y')\,U^{d'b}(z)\,\rho^{e'}(x')-U^{d'b}(z)\,U^{c'e'}(x')\,\rho^{e'}(x')]$
\item[\textbf{l.210}] \textsf{\small [term removed/replaced]} $\displaystyle [U^{cb}(z)\,U^{ad}(y)\,U^{ce}(x)\,\rho^{d}(y)\,\rho^{e}(x)-U^{ad}(y)\,\rho^{d}(y)\,\rho^{b}(x)]$\ $\longrightarrow$\ $\displaystyle -\frac{(x'-z)^{k'}}{(x'-z)^{2}}\,[U^{c'e'}(z)\,U^{d'b}(z)\,\rho^{e'}(x')-U^{d'b}(z)\,U^{c'e'}(x')\,\rho^{e'}(x')]\Bigg]\times$
\item[\textbf{l.210}] \textsf{\small [typo/LaTeX]} $\displaystyle [U^{cb}(z)\,U^{ad}(y)\,U^{ce}(x)\,\rho^{d}(y)\,\rho^{e}(x)-U^{ad}(y)\,\rho^{d}(y)\,\rho^{b}(x)]$\ $\longrightarrow$\ $\displaystyle [U^{cb}(z)\,U^{ad}(y)\,U^{ce}(x)\,\rho^{d}(y)\,\rho^{e}(x)-U^{ad}(y)\,\rho^{d}(y)\,\rho^{b}(x)]$
\end{itemize}
\begin{notebox}\textsf{\small Note, l.212:} [This is Row VI of \texttt{Other\_terms} (v2) at $\vee\to\infty$: upper limit $\infty$, last term ($\log\big(1-\frac{k^{+}}{\vee}\big)\times\dots$) dropped; nothing else is changed. Kernels and $p^{+}$ dependence are those of Row III, first part, with twice the prefactor: the condition for the limit ($y'$ integral before the $p^{+}$ integral in the second term; conditional convergence) and the numerical check are the same (numbers of Row III, first part, $\times2$). The soft subtraction terms of the plus prescriptions in this equation are finite only in the sum over Group I Rows III, V and Group III Rows III, V, VI under one integral.]\end{notebox}
\needspace{15\baselineskip}\bigskip
\subsection{Limiting cases final 1}\label{sec:file-limiting_final_1}
\noindent\textbf{Original:} \texttt{/root/.claude/uploads/0ab97fcf-94a2-5aec-8bad-586400481685/7daa274b-Limiting\_cases\_final\_1.tex}\\\textbf{Checked copy:} \texttt{/home/user/correlations/nlo\_review/limits\_checked/Limiting\_cases\_final\_1\_checked.tex}\\
Lines: 222 (original), 223 (checked copy). Coloured groups of the check: 6 (6 top-level); author\textquotesingle s own coloured groups left unchanged: 0.\\
Records by type: term removed/replaced 2; note only 4.
\par\noindent Unmarked layout-only edits (line breaks \texttt{\textbackslash notag\textbackslash\textbackslash\ \&}, \texttt{\textbackslash par\textbackslash noindent} around inserted notes, spacing), checked-copy lines: 36, 108, 122.
\par\noindent Unmarked edits that are not layout only (listed with their block below): 4.
\subsubsection*{Changes, block by block}
\needspace{3\baselineskip}\noindent\rule{\linewidth}{0.3pt}\par\noindent\textbf{Lines 36}\par
\begin{notebox}\textsf{\small Note, l.36:} [\textbf{Checked copy of the author's file.} Red: changed or added. Orange: removed. The dark-red rows (Group I Row IV, Group III Row I) are the author's marking of the two rows that cancel each other. Group III Rows III, V, VI keep $\vee-k^{+}$ as the upper limit of the $p^{+}$ integral: in this form the integrals are independent of $\vee$ up to corrections of order $k^{+}/\vee$ (the large logarithm is removed completely), and all four equations use the same limit consistently.]\end{notebox}
\needspace{3\baselineskip}\noindent\rule{\linewidth}{0.3pt}\par\noindent\textbf{Lines 107--108} \quad \textsf{Group I (with c.c.) $\rangle$ Row III} \quad {\small (in \texttt{align}, lines 98--108)}\par
\begin{itemize}[leftmargin=2.2em,itemsep=1pt,topsep=2pt]
\item[\textbf{l.107}] \textsf{\small [term removed/replaced]} \par\noindent\emph{original:} $\displaystyle -\frac{ig^{4}f^{abd}}{16\pi^{5}}\,\frac{1}{k^{+}}\,\int_{x,y,z,x'}\log{[\frac{(z-w)^{2}}{(x-w)^{2}}]}\,\frac{(x'-w')\cdot(y-z)}{(x'-w')^{2}(y-z)^{2}\,(z-w)^{2}}$\par\noindent\emph{checked:} $\displaystyle -\frac{ig^{4}f^{abd}}{8\pi^{5}}\,\frac{1}{k^{+}}\,\int_{x,y,z,x'}\mathcal{E}\big(k\cdot(z-w)\big)\,\frac{(x-z)\cdot(w-z)\,(x'-w')\cdot(y-w)}{(x-z)^{2}\,(w-z)^{2}\,(x'-w')^{2}\,(y-w)^{2}}$
\item[\textbf{l.107}] \textsf{\small [term removed/replaced]} \par\noindent\emph{original:} $\displaystyle [U^{ab'}(x')\,\rho^{b'}(x')-U^{ab'}(w')\, \rho^{b'}(x')]\times [U^{de}(x)\,U^{bc}(z)\, \rho^{e}(x)\,\rho^{c}(y)-U^{de}(x)\,U^{bc}(y)\,\rho^{e}(x)\,\rho^{c}(y)]\Bigg\}$\par\noindent\emph{checked:} $\displaystyle [U^{ab'}(x')\,\rho^{b'}(x')-U^{ab'}(w')\, \rho^{b'}(x')]\times [U^{de}(x)\,U^{bc}(w)\, \rho^{e}(x)\,\rho^{c}(y)-U^{de}(x)\,U^{bc}(y)\,\rho^{e}(x)\,\rho^{c}(y)]\Bigg\}$
\end{itemize}
\begin{notebox}\textsf{\small Note, l.108:} [with $\mathcal{E}(s)\equiv\int_{0}^{1}\frac{dt}{t}\big(e^{-ist}-1\big)=\mathrm{Ci}(|s|)-\gamma_{E}-\log|s|-i\,\mathrm{Si}(s)$. \textbf{Last term of $N_{3b}$ replaced.}]\end{notebox}
\begin{notebox}\textsf{\small Note on removed material, l.108:} [Removed: $-\frac{ig^{4}f^{abd}}{16\pi^{5}}\frac{1}{k^{+}}\int_{x,y,z,x'}\log\big[\frac{(z-w)^{2}}{(x-w)^{2}}\big]\frac{(x'-w')\cdot(y-z)}{(x'-w')^{2}(y-z)^{2}(z-w)^{2}}\times$ (colour bracket with $U^{bc}(z)$): its $z$ integral diverges at $z\to w$ like $\log^{2}$; it is the $\Lambda\to0$ limit at fixed $z\neq w$ only.]\end{notebox}
\begin{notebox}\textsf{\small Note, l.108:} [The JIMWLK term of this row is then the one of the evolution lists, with $U^{bc}(w)$. The four equations $N_{1}$, $N_{2}$, $N_{3a}$, $N_{3b}$ are not integrable one by one (poles on the line $(x-z)^{2}=(x-w)^{2}$ and at $z\to w$ cancel only in their sum): they have to be read as one integrand under one integral; their sum is the combined kernel given in \texttt{Limiting\_cases\_final}.]\end{notebox}
\subsubsection*{Other unmarked edits}
\noindent{\small\textsf{l.122 (deleted):} \texttt{\{\textbackslash{}\allowbreak{}color\{red\}\allowbreak{}}}\par
\noindent{\small\textsf{l.122 (deleted):} \texttt{\{\textbackslash{}\allowbreak{}color\{red\}\allowbreak{}}}\par
\noindent{\small\textsf{l.129 (deleted):} \texttt{\{\textbackslash{}\allowbreak{}color\{red\}\allowbreak{}}}\par
\noindent{\small\textsf{l.129 (deleted):} \texttt{\}\allowbreak{}\}\allowbreak{}\{\textbackslash{}\allowbreak{}color\{orange\}\allowbreak{}\textbackslash{}\allowbreak{}cancel\{)\allowbreak{}\}\allowbreak{}}}\par
\needspace{15\baselineskip}\bigskip
\subsection{Final terms}\label{sec:file-final_terms}
\noindent\textbf{Original:} \texttt{/root/.claude/uploads/0ab97fcf-94a2-5aec-8bad-586400481685/b2466302-Final\_terms.tex}\\\textbf{Checked copy:} \texttt{/home/user/correlations/nlo\_review/final\_checked/Final\_terms\_checked.tex}\\
Lines: 416 (original), 416 (checked copy). Coloured groups of the check: 16 (16 top-level); author\textquotesingle s own coloured groups left unchanged: 2.\\
Records by type: term removed/replaced 2; typo/LaTeX 1; note only 12.
\par\noindent Unmarked layout-only edits (line breaks \texttt{\textbackslash notag\textbackslash\textbackslash\ \&}, \texttt{\textbackslash par\textbackslash noindent} around inserted notes, spacing), checked-copy lines: 35, 72, 100, 165, 229, 239, 313, 316, 384.
\subsubsection*{Changes, block by block}
\needspace{3\baselineskip}\noindent\rule{\linewidth}{0.3pt}\par\noindent\textbf{Lines 35}\par
\begin{notebox}\textsf{\small Note, l.35:} [\textbf{Checked copy: internal consistency with the verified files of this project.} Red: corrected or added (the two red commas in Group I Row V were already in the file). Orange: removed. Blue is the author's own marking and is left as it is. Result: the BFKL and JIMWLK sections, $N_{\rm RC}$, the LO+RC formula and all finite terms agree character by character with the verified versions, with two exceptions, corrected below: (1) the phase of $N_{\rm DGLAP}$ (the $\frac{1}{\xi}$ belongs in the exponent); (2) the last term of $N_{3b}$ of Group I Row III (divergent at $z\to w$; replaced by the recoil term with $\mathcal{E}$). The phase is $e^{-ik\cdot(\text{unprimed}-\text{primed})}$ in all sections. The LO term has no section of its own: it appears only in the LO+RC formula. Group I Row IV and Group III Row I are not listed because they cancel each other.]\end{notebox}
\needspace{3\baselineskip}\noindent\rule{\linewidth}{0.3pt}\par\noindent\textbf{Lines 72} \quad \textsf{BFKL contribution}\par
\begin{notebox}\textsf{\small Note, l.72:} [This list and the next equation are the same quantity (term-by-term list and grouped form; checked numerically for SU(2), SU(3)); they are not to be added]\end{notebox}
\needspace{3\baselineskip}\noindent\rule{\linewidth}{0.3pt}\par\noindent\textbf{Lines 100}\par
\begin{notebox}\textsf{\small Note, l.100:} [Equal to one step of projectile evolution of the LO cross section. Remark: with non-commuting $\rho$ the evolution gives $\tfrac12\{\rho(x'),\rho(x)\}$ in the four blocks with $\mathcal{K}(x'-w')\cdot\mathcal{K}(w-z)$ and $\mathcal{K}(x-w)\cdot\mathcal{K}(w'-z)$; the fixed orders written there differ from it by one-$\rho$ contact terms $\propto f^{abc}\rho^{c}(x)\,\delta^{2}(x-x')$, which drop out of the colour-singlet form]\end{notebox}
\needspace{3\baselineskip}\noindent\rule{\linewidth}{0.3pt}\par\noindent\textbf{Lines 165} \quad \textsf{JIMWLK contribution}\par
\begin{notebox}\textsf{\small Note, l.165:} [$N_{1}+N_{2}$ equals the next equation (grouped form, $=-H_{\rm JIMWLK}$ acting on the LO Wilson-line structure; checked numerically); they are not to be added. The phase $e^{-ik(w-w')}$ means $e^{-ik\cdot(w-w')}$ as in the BFKL section]\end{notebox}
\needspace{3\baselineskip}\noindent\rule{\linewidth}{0.3pt}\par\noindent\textbf{Lines 224} \quad \textsf{DGLAP contribution} \quad {\small (in \texttt{align}, lines 223--226)}\par
\begin{itemize}[leftmargin=2.2em,itemsep=1pt,topsep=2pt]
\item[\textbf{l.224}] \textsf{\small [typo/LaTeX]} {\small\ttfamily \}\allowbreak{}\textbackslash{}\allowbreak{}frac\{1\}\allowbreak{}\{\textbackslash{}\allowbreak{}xi\}\allowbreak{}}\ $\longrightarrow$\ {\small\ttfamily \textbackslash{}\allowbreak{}frac\{1\}\allowbreak{}\{\textbackslash{}\allowbreak{}xi\}\allowbreak{}\}\allowbreak{}}
\end{itemize}
\needspace{3\baselineskip}\noindent\rule{\linewidth}{0.3pt}\par\noindent\textbf{Lines 229}\par
\begin{notebox}\textsf{\small Note, l.229:} [Phase: the factor $\frac{1}{\xi}$ is in the exponent, $e^{-ik\cdot(y-y')\frac{1}{\xi}}$, as in the verified Row IV of Group II (the measure is $\frac{d\xi}{\xi^{2}}$ only). Both plus prescriptions act on the whole function $f(\xi)\equiv\frac{1}{\xi^{2}}\int_{x,y,x',y'}e^{-ik\cdot(y-y')\frac{1}{\xi}}[\dots]\dots$: $\int_{0}^{1}d\xi\,\big[\frac{1-\xi}{\xi}\big]_{+}f\equiv\int_{0}^{1}d\xi\,\frac{1-\xi}{\xi}[f(\xi)-f(0)]$, $\int_{0}^{1}d\xi\,\frac{\xi}{(1-\xi)_{+}}f\equiv\int_{0}^{1}d\xi\,\frac{\xi f(\xi)-f(1)}{1-\xi}$, with $f(0)=-\frac{4\pi}{k^{2}}\int_{x,x',z}\frac{(x-z)\cdot(x'-z)}{(x-z)^{2}(x'-z)^{2}}[U^{ce'}(z)-U^{ce'}(x')][U^{ce}(z)-U^{ce}(x)]\rho^{e'}(x')\rho^{e}(x)$. This is consistent with the rest of the file only in this combination: the subtraction is inside $\widetilde{\mathcal P}_{gg}$, the remaining piece of the $\frac{1}{\xi}$ pole is the term $-\frac{1}{k^{+}}\times$(BFKL structure) kept in the finite terms (Group II Row IV, second term), and the $\delta(1-\xi)$ term is exactly $-N_{\rm RC}$ (added and subtracted: $N_{\rm DGLAP}+N_{\rm RC}$ does not depend on $\beta_{0}$)]\end{notebox}
\needspace{3\baselineskip}\noindent\rule{\linewidth}{0.3pt}\par\noindent\textbf{Lines 239} \quad \textsf{Running coupling contribution}\par
\begin{notebox}\textsf{\small Note, l.239:} [This is the only place where the LO term appears in this file. $g^{2}\to g^{2}(\mu)$ holds only together with the $\delta(1-\xi)$ term of $N_{\rm DGLAP}$ and the same $\mu$]\end{notebox}
\needspace{3\baselineskip}\noindent\rule{\linewidth}{0.3pt}\par\noindent\textbf{Lines 241} \quad \textsf{Finite terms}\par
\begin{notebox}\textsf{\small Note, l.241:} [Group III Rows III, V, VI keep $\vee-k^{+}$ as the upper limit of the $p^{+}$ integral: in this form the integrals do not depend on $\vee$ up to corrections of order $k^{+}/\vee$.]\end{notebox}
\needspace{3\baselineskip}\noindent\rule{\linewidth}{0.3pt}\par\noindent\textbf{Lines 312--313} \quad \textsf{Group I (with c.c.) $\rangle$ Row III} \quad {\small (in \texttt{align}, lines 303--313)}\par
\begin{itemize}[leftmargin=2.2em,itemsep=1pt,topsep=2pt]
\item[\textbf{l.312}] \textsf{\small [term removed/replaced]} \par\noindent\emph{original:} $\displaystyle -\frac{ig^{4}f^{abd}}{16\pi^{5}}\,\frac{1}{k^{+}}\,\int_{x,y,z,x'}\log{[\frac{(z-w)^{2}}{(x-w)^{2}}]}\,\frac{(x'-w')\cdot(y-z)}{(x'-w')^{2}(y-z)^{2}\,(z-w)^{2}}$\par\noindent\emph{checked:} $\displaystyle -\frac{ig^{4}f^{abd}}{8\pi^{5}}\,\frac{1}{k^{+}}\,\int_{x,y,z,x'}\mathcal{E}\big(k\cdot(z-w)\big)\,\frac{(x-z)\cdot(w-z)\,(x'-w')\cdot(y-w)}{(x-z)^{2}\,(w-z)^{2}\,(x'-w')^{2}\,(y-w)^{2}}$
\item[\textbf{l.312}] \textsf{\small [term removed/replaced]} \par\noindent\emph{original:} $\displaystyle [U^{ab'}(x')\,\rho^{b'}(x')-U^{ab'}(w')\, \rho^{b'}(x')]\times [U^{de}(x)\,U^{bc}(z)\, \rho^{e}(x)\,\rho^{c}(y)-U^{de}(x)\,U^{bc}(y)\,\rho^{e}(x)\,\rho^{c}(y)]$\par\noindent\emph{checked:} $\displaystyle [U^{ab'}(x')\,\rho^{b'}(x')-U^{ab'}(w')\, \rho^{b'}(x')]\times [U^{de}(x)\,U^{bc}(w)\, \rho^{e}(x)\,\rho^{c}(y)-U^{de}(x)\,U^{bc}(y)\,\rho^{e}(x)\,\rho^{c}(y)]$
\end{itemize}
\begin{notebox}\textsf{\small Note, l.313:} [with $\mathcal{E}(s)\equiv\int_{0}^{1}\frac{dt}{t}\big(e^{-ist}-1\big)=\mathrm{Ci}(|s|)-\gamma_{E}-\log|s|-i\,\mathrm{Si}(s)$. \textbf{Last term of $N_{3b}$ replaced} (as in the verified finite terms).]\end{notebox}
\begin{notebox}\textsf{\small Note on removed material, l.313:} [Removed: $-\frac{ig^{4}f^{abd}}{16\pi^{5}}\frac{1}{k^{+}}\int_{x,y,z,x'}\log\big[\frac{(z-w)^{2}}{(x-w)^{2}}\big]\frac{(x'-w')\cdot(y-z)}{(x'-w')^{2}(y-z)^{2}(z-w)^{2}}\times$ (colour bracket with $U^{bc}(z)$)]\end{notebox}
\begin{notebox}\textsf{\small Note, l.313:} [Reason: its $z$ integral diverges at $z\to w$ like $\log^{2}$ (it is the $\Lambda\to0$ limit at fixed $z\neq w$ only). $N_{1}$, $N_{2}$, $N_{3a}$, $N_{3b}$ are to be read as one integrand under one integral (poles on $(x-z)^{2}=(x-w)^{2}$ and at $z\to w$ cancel only in the sum).]\end{notebox}
\needspace{3\baselineskip}\noindent\rule{\linewidth}{0.3pt}\par\noindent\textbf{Lines 316}\par
\begin{notebox}\textsf{\small Note, l.316:} [Row IV is not listed: it cancels against Group III Row I]\end{notebox}
\needspace{3\baselineskip}\noindent\rule{\linewidth}{0.3pt}\par\noindent\textbf{Lines 384} \quad \textsf{Group III (with c.c.)}\par
\begin{notebox}\textsf{\small Note, l.384:} [Row I is not listed: it cancels against Group I Row IV]\end{notebox}
\clearpage
\section{Calculations and verifications}\label{sec:calc}
\noindent This section collects the derivations behind the substantive changes. Each subsection gives the key equations and numbers and the document (and line) from which they are taken; the full derivations are in those documents. ``[A]'' = analytic, ``[N]'' = numerical, as marked there.

\subsection{Operator ordering of the three-$\rho$ term of $\mf{A}^{(3)}$}\label{sec:calc-order}
\noindent\emph{Source:} \texttt{Calculation\_main\_1\_checked.tex}, notes at l.421, l.555, l.1101, l.1253; \texttt{Group\_I\_checked.tex} l.35, l.38, l.43.

\noindent In third-order perturbation theory with operator-valued matrix elements, the single-emission matrix element $\langle g^{a}_{i}(k)|H_{g}|0\rangle\propto\rho^{a}(-\mathbf{k})$ stands to the \emph{left} of the norm factor $\propto\rho^{b}(\mathbf{p})\rho^{b}(-\mathbf{p})$ (l.421). The surviving $\log(\vee/\Lambda)$ term of the three-$\rho$ state comes from $|\Psi^{7}_{g\rho\rho\rho}\rangle$, so
\begin{align}
[\mf{A}^{(3)}_{2}]_{\rho\rho\rho}=\mf{A}^{(1)}_{2}\,\mf{N}_{2}\,,\qquad\text{coordinate space: }\rho^{a}(\mathbf{w})\,\rho^{b}(\mathbf{x})\,\rho^{b}(\mathbf{y})\,.\notag
\end{align}
\textbf{Evidence} (l.555): ``verified with explicit matrices against an exact construction of $\Omega$ and against exact diagonalisation (agreement to $3\times10^{-16}$; the ordering $\rho^{b}\rho^{b}\rho^{a}$ deviates at order one).'' This differs from the ordering written in the paper, Eqs.~(3.81)--(3.83) and (3.102)--(3.103). The $1/\Lambda$ and $1/\vee$ power terms still cancel between $|\Psi^{6c}_{g\rho\rho\rho}\rangle$ (ordering $\rho^{b}\rho^{b}\rho^{a}$) and $|\Psi^{7}_{g\rho\rho\rho}\rangle$ (ordering $\rho^{a}\rho^{b}\rho^{b}$), because $\int d^{2}\mathbf{p}\,\rho^{b}(\mathbf{p})\rho^{b}(-\mathbf{p})$ (no $1/\mathbf{p}^{2}$ weight) is local in coordinate space and commutes with $\rho^{a}$; the two-$\rho$ part, Eq.~(combinedgrhorho), including its $-2k^{i}\log(\vee/\Lambda)$ term, is correct as it stands (l.555). Changes made: $\rho^{a}(-\mathbf{k})$ moved to the left (and the matrix element $\langle g|H_{g}|0\rangle$ at l.420) in Eqs.\ at l.424, 428, 432, 557, 561, 1252, 1256 of the main file; the statement ``$\mf{A}^{(3)}_{2}$'' $\to$ ``$\mf{A}^{(1)}_{2}$'' at l.1240.

\subsection{The commutator $[\mf{N}_{2},\mf{A}^{(1)}]$ and its cancellation}\label{sec:calc-comm}
\noindent\emph{Source:} \texttt{Calculation\_main\_1\_checked.tex} l.1092--1101, 1146--1147, 1435--1451, 1531, 1537; \texttt{Group\_I\_checked.tex} l.38, 43, 53, 60, 67, 74; \texttt{4rho\_checked.tex} l.69, 124; evolution files (notes at Group I Rows I, II).

\noindent (i) The original statement $\mf{N}^{(2)}\mf{A}^{(1)}=\mf{A}^{(1)}\mf{N}^{(2)}$ (``$A^{(1)}$ can at most be linear in $\rho$'') was removed: $\mf{N}^{(2)}\propto\int\rho^{a}\rho^{a}/\mathbf{p}^{2}$ does not commute with a single $\rho^{b}$, since $[\rho^{a}(x)\rho^{a}(y),\rho^{b}(w)]\neq0$ for $x\neq y$, so $[\mf{N}^{(2)},\mf{A}^{(1)}]\propto g^{3}f^{abc}\{\rho^{a},\rho^{c}\}\neq0$ (l.1092). The relation for $\mf{A}^{(3)\dagger}_{1}$ keeps the commutator, $-[\mf{N}^{(2)},\mf{A}^{(1)b}_{2j}]$ (l.1100, l.1146).

\noindent (ii) With the ordering of Sec.~\ref{sec:calc-order}, $-[\mf{A}^{(3)}_{2}]_{\rho\rho\rho}-[\mf{N}^{(2)},\mf{A}^{(1)}_{2}]=-\mf{N}^{(2)}\mf{A}^{(1)}_{2}$ (l.1101, l.1147). In the cross section the term $[\bar{\mf{N}}_{2},\bar{\mf{A}}^{(1)\dagger}_{1}]$ is added (l.1435, l.1451, l.1531) and the $\mf{N}$-type terms of the order-$g^{3}$ amplitude add up to
\begin{align}
-\bar{\mf{A}}\mf{N}_{2}+U\bar{\mf{N}}_{2}\mf{A}-[\bar{\mf{N}}_{2},\bar{\mf{A}}]+U[\mf{A}^{(3)}]_{\rho\rho\rho}-[\bar{\mf{A}}^{(3)}]_{\rho\rho\rho}=-\bar{\mf{A}}\mf{N}_{2}+U\bar{\mf{N}}_{2}\mf{A}+U\mf{A}\mf{N}_{2}-\bar{\mf{N}}_{2}\bar{\mf{A}}\,,\qquad\mf{A}\equiv\mf{A}^{(1)}_{2}\,,\notag
\end{align}
which vanishes at $U=1$, as $\bar\Omega^{\dagger}\Omega=1$ requires; no three-$\rho$ term arises from the commutator (l.1448, l.1537).

\noindent (iii) Consequence for Group I Row II (Group\_I l.43): the remainder $-U[\mf{N}_{2},\mf{A}^{(1)}]$ is a $\{\rho,\rho\}$ term, and
\begin{align}
[\mf{N}_{2},\mf{A}^{(1)a}_{i}(k^{+},-w)]=-\frac{g^{3}}{8\sqrt{2}\pi^{4}\sqrt{k^{+}}}\log\frac{\vee}{\Lambda}\,f^{abc}\int_{x,y,z}\frac{(x-z)\cdot(y-z)}{(x-z)^{2}(y-z)^{2}}\frac{(y-w)^{i}}{(y-w)^{2}}\{\rho^{b}(x),\rho^{c}(y)\}\notag
\end{align}
is exactly the term $-2k^{i}\log\frac{\vee}{\Lambda}$ of the two-$\rho$ part of $\mf{A}^{(3)}$ (``same coefficient and sign, and with explicit SU(2), SU(3) matrices''). It removes that term for the unbarred colour structure $f^{hbc}U^{ah}(w)\{\rho^{b}(x),\rho^{c}(y)\}$; for the barred structure it stays. Implemented by the red terms added to $N_{2}$, $N_{2a}$, $N_{2b}$ (Group\_I l.52, 59, 66): their $\log\frac{k^{+}}{\Lambda}$ part is the term which the JIMWLK and three-$\rho$ lists do not contain, their $\log\frac{\vee}{k^{+}}$ part ($R^{i}\to R^{i}+2k^{i}\log\frac{\vee}{k^{+}}$) is the Row II term of the BFKL file. Equivalent alternative (Group\_I l.43; 4rho l.69): no red terms, and the unbarred four-$\rho$ term written as $U\mf{A}^{(1)}\mf{N}_{2}$. Numbers used: $\int d^{2}k\,e^{ik\cdot(w-y)}\frac{k^{i}}{k^{2}}=-2\pi i\frac{(y-w)^{i}}{(y-w)^{2}}$, $\int_{z}\frac{(x-z)\cdot(y-z)}{(x-z)^{2}(y-z)^{2}}=-\pi\log(x-y)^{2}$ + const (``the difference of the integral at two separations $r$, $r_{0}$ is $-\pi\log(r^{2}/r_{0}^{2})$'', checked numerically; Group\_I l.60). Other\_terms\_v2 l.49 confirms numerically (residuals below $10^{-14}$) that the two-$\rho$ bracket equals $C_{J}\log\frac{k^{+}}{\Lambda}+R^{i}$ for the barred and $(C_{J}+2q^{i})\log\frac{k^{+}}{\Lambda}+2q^{i}\log\frac{\vee}{k^{+}}+R^{i}$ for the unbarred structure.

\subsection{The sign of the non-instantaneous $p^{+}>k^{+}$ term of $\mf{B}_{3}$}\label{sec:calc-B3}
\noindent\emph{Source:} \texttt{Calculation\_main\_1\_checked.tex} l.1276--1303 (three notes); \texttt{Group\_I\_checked.tex} l.82; \texttt{scales\_checked.tex} l.44--81.

\noindent Change: $\frac{(k^{+}-p^{+})}{(k-p)^{2}}\to\frac{(p^{+}-k^{+})}{(k-p)^{2}}$ in momentum space (l.1276), and a minus sign in front of the corresponding coordinate-space terms (l.1289, l.1302). Reason (l.1277): the sign follows from Eqs.~(m2), (m3), from the cancellation of the $1/(1-\xi)^{2}$ singularity, and from the $a^{\dagger}a$ unitarity constraint; Eqs.~(onegrho12) and (Eq:onegrhoX1) carry the wrong overall sign (in the checked copy these two equations carry red sign changes at l.684 and l.706). Consistency: the $p^{+}>k^{+}$ term of Group I Row III already has the corrected sign, so that the $1/\Lambda$ terms cancel between $N_{1}$ and $N_{2}$, $N_{3a}$ (Group\_I l.82; Sec.~\ref{sec:calc-scales}).

\subsection{The support of $\mf{B}_{2}$}\label{sec:calc-B2}
\noindent\emph{Source:} \texttt{Calculation\_main\_1\_checked.tex} l.1260--1269; Group I l.128, l.136; Group II l.40; Group III l.41; 4rho l.69 (iv).

\noindent Change: $\Theta(k^{+}-\Lambda-p^{+})\to\Theta(\vee-k^{+}-p^{+})$ in the single-$\rho$ (triple-vertex and instantaneous) terms and $\to\Theta(\vee-k^{+})\Theta(\vee-p^{+})$ in the $\{\rho,\rho\}$ term (with $k^{+},p^{+}>\Lambda$). Reason: kinematics. In the triple vertex the parent gluon carries $k^{+}+p^{+}$ and is itself a soft mode; the $\{\rho,\rho\}$ term describes two independent emissions, $\Lambda<k^{+}<\vee$, $\Lambda<p^{+}<\vee$ with no condition on the sum, as for the same intermediate state in Eq.~(3rho). With $k^{+}+p^{+}<\vee$ the four-$\rho$ terms would not cancel for commuting $\rho$ (coherent state): only with a common $\log\frac{\vee}{\Lambda}$ do they cancel. $\Theta(k^{+}-\Lambda-p^{+})$ is not symmetric under the exchange of the two gluons and removes $p^{+}>k^{+}$. The first (commutator) term is left to the author (l.1261). Consequences in the cross section: $\Lambda<p^{+}<\vee-k^{+}$ for every product containing the single-$\rho$ part of $\mf{B}_{2}$ or $\mf{C}_{1}$, $\Lambda<p^{+}<\vee$ for its $\{\rho,\rho\}$ part (four-$\rho$ terms), $\Lambda<p^{+}<k^{+}-\Lambda$ for the $\bar{\mf{C}}^{\dagger}_{1}\mf{B}_{2}$ terms (l.1537, l.1542).

\subsection{The sign of Group III Row III}\label{sec:calc-G3R3}
\noindent\emph{Source:} \texttt{Group\_III\_checked.tex} l.62--63.

\noindent Change: a minus sign added to the prefactor of the first equation of Row III, i.e.\ $-\frac{4}{(2\pi)^{3}}$. Reason: in Eq.~(CS2) of the main file this cross term enters as $-\int_{y'}[\dots]\times[\dots]$ (the $\mf{C}_{1}\mf{A}$ part of the amplitude has the opposite sign to the $\mf{B}_{2}$ part); the expressions further down already correspond to this sign (``checked numerically from the coefficient functions''). The range $\Lambda<p^{+}<\vee-k^{+}$ holds for both parts, since the parent gluon of $\mf{C}_{1}$ carries $p^{+}+k^{+}<\vee$. The three-$\rho$ list is not affected (Evolution\_three\_rho l.150: the sign change concerns only the first line).

\subsection{Group I Row III: the excluded window and the cancellation of the power terms}\label{sec:calc-scales}
\noindent\emph{Source:} \texttt{Group\_I\_checked.tex} l.80--84; \texttt{scales\_checked.tex} l.44--81.

\noindent Change (Group\_I l.80): $\int_{\Lambda}^{\vee}dp^{+}\to\big(\int_{\Lambda}^{k^{+}-\Lambda}+\int_{k^{+}+\Lambda}^{\vee}\big)dp^{+}$ in the third (instantaneous) term: the window $|p^{+}-k^{+}|<\Lambda$ is excluded, as in the wave functions. This adds two logarithms in the last two terms of $N_{1}$ (l.84):
\begin{align}
-\log\frac{k^{+}+\Lambda}{k^{+}-\Lambda}\,,\qquad -\log\frac{(k^{+}+\Lambda)(x-z)^{2}-k^{+}(x-w)^{2}}{(k^{+}-\Lambda)(x-z)^{2}-k^{+}(x-w)^{2}}\,.\notag
\end{align}
The poles at $p^{+}(x-z)^{2}=k^{+}(x-w)^{2}$ are principal values, so every logarithm is $\log|\dots|$ (l.82). Power terms (scales\_checked): with $D=p^{+}(x-z)^{2}-k^{+}(x-w)^{2}$,
\begin{align}
\frac{p^{+}+k^{+}}{(k^{+}-p^{+})^{2}D}=\frac{2}{(x-z)^{2}-(x-w)^{2}}\frac{1}{(p^{+}-k^{+})^{2}}+\text{simple poles},\quad
\int_{k^{+}+\Lambda}^{\vee}\frac{dp^{+}}{(p^{+}-k^{+})^{2}}=\frac1\Lambda-\frac{1}{\vee-k^{+}},\quad\int_{\Lambda}^{k^{+}-\Lambda}\frac{dp^{+}}{(p^{+}-k^{+})^{2}}=\frac1\Lambda-\frac{1}{k^{+}-\Lambda}.\notag
\end{align}
The non-instantaneous terms have the double pole $-\frac{1}{(x-z)^{2}-(x-w)^{2}}\frac{1}{(p^{+}-k^{+})^{2}}$ in both ranges; $N_{3b}$ and the instantaneous term only simple poles. With the common denominator $(x-z)^{2}-(x-w)^{2}$ the factors are $N_{1}$: $\frac1\Lambda+\frac{1}{k^{+}-\vee}+\frac1\Lambda-\frac{1}{k^{+}-\Lambda}$; $N_{2}$: $\frac{1}{\vee-k^{+}}-\frac1\Lambda$; $N_{3a}$: $-\frac1\Lambda+\frac{1}{k^{+}-\Lambda}$; ``the sum of the eight is identically zero at finite $\Lambda$ and $\vee$ (not only in the limit)'' (scales l.81). Other power terms: $N_{2b\,\rm UV}$ of Group I Row V has the coefficient $\frac{1}{k^{+}}\big[2\log\frac{k^{+}-\Lambda}{\Lambda}-\frac{11}{6}+\frac{4\Lambda}{k^{+}}-\frac{\Lambda^{2}}{k^{+2}}+\frac{2\Lambda^{3}}{3k^{+3}}\big]$ (positive powers only, l.55); $N_{2f}$ of Group II Row I contains $\big[\frac1\vee-\frac{1}{k^{+}+\Lambda}\big]\to-\frac{1}{k^{+}}$ (finite, l.58); Group II Row IV has the BFKL coefficient $\frac{1}{k^{+}}\log\frac{\vee}{k^{+}}-\frac{1}{k^{+}}+\frac1\vee$ (l.64). Also in scales: the missing $\int_{w,w'}$ in front of the phase was added (l.80).

\subsection{Pole bookkeeping and the cancellation of the $\frac{11}{3}$ terms}\label{sec:calc-poles}
\noindent\emph{Source:} \texttt{Divergent\_terms\_checked.tex} l.42--112; \texttt{Other\_terms\_checked.tex} l.35, 46, 143.

\noindent The bracket $\big[\frac{2}{\epsilon}+2\gamma_{E}+\log\frac{(x-w)^{2}\mu^{2}}{4}\big]$ multiplies $\frac{11}{3}-2L+\ell$ in Group I Row II ($L\equiv\log\frac{k^{+}}{\Lambda}$, $\ell\equiv\log(1-\frac{k^{+}}{\vee})$) and $2L-\frac{11}{6}$ in $N_{2b\,\rm UV}$ of Group I Row V, where $\frac{1}{k^{+}}\int_{\Lambda}^{k^{+}-\Lambda}dp^{+}\big[\frac{p^{+}(k^{+}-p^{+})}{k^{+2}}+\frac{p^{+}}{k^{+}-p^{+}}+\frac{k^{+}-p^{+}}{p^{+}}\big]=2L-\frac{11}{6}$ (recomputed; l.53). After the renaming $x\to w$, $y\to x$ in the Row V term ($[U(y)-U(x)]\rho(y)\to-[U(w)-U(x)]\rho(x)$, same order $\rho(x')\rho(x)$), the $\frac{11}{3}$ terms cancel completely (pole, $2\gamma_{E}$ and logarithm), provided $\bar\mu_{\overline{MS}}=\mu_{\overline{MS}}$ (l.60). Pole coefficients in units of $\frac{1}{(2\pi)^{3}}\frac{g^{4}N_{c}}{8\pi^{4}}\frac{1}{k^{+}}\frac{2}{\epsilon}\int e^{-ik\cdot(w-w')}\frac{(x'-w')\cdot(x-w)}{(x'-w')^{2}(x-w)^{2}}[U(w')-U(x')][U(w)-U(x)]\rho(x')\rho(x)$ (l.68):
\begin{center}
\begin{tabular}{lc}
\toprule
Group I Row II & $\frac{11}{6}-L+\frac{\ell}{2}$\\
Group I Row V & $-\frac{11}{6}+2L$\\
Group II Row IV, term with $L+\ell$ & $-L-\ell$\\
Group II Row IV, RC & $+\frac{\beta_{0}}{4N_{c}}=\frac{11}{12}$\\
Group II Row IV, $\delta(1-\xi)$ term of DGLAP & $-\frac{11}{12}$\\
\midrule
sum & $-\frac{\ell}{2}$\\
\bottomrule
\end{tabular}
\end{center}
The $\frac{11}{3}$ poles and the $L$ poles cancel, the two $\beta_{0}$ terms cancel, and a pole $\propto\ell$ remains that vanishes only if $\vee\to\infty$ is taken before $\epsilon\to0$; in addition there is the pole of the remaining DGLAP term, with the integrand at $k/\xi$, not cancelled by any other term (l.68, l.86 (iv)). The missing $\log(1-\frac{k^{+}}{\vee})$ term with a pole bracket was added to the Divergent\_terms list (l.90). In $N_{2b\,\rm UV}$ the factor $d_{\perp}$ multiplies the angular average $\delta^{im}/d_{\perp}$, so no extra constant ($-\frac16$) arises (l.112). For the complete account (UV/IR character, coupling counterterm, fragmentation counterterm) see Sec.~\ref{sec:fix-poles}.

\subsection{DGLAP, running coupling and the $-1/k^{+}$ remnant}\label{sec:calc-dglap}
\noindent\emph{Source:} \texttt{Group\_II\_Row\_IV\_checked.tex} l.93--184; \texttt{Divergent\_terms\_checked.tex} l.86, 112; \texttt{Cancellation\_of\_flagged\_terms.tex} Sec.~2.3; \texttt{Finite\_terms\_final.tex} Sec.~2.4; \texttt{Limiting\_cases\_checked.tex} FLAG 5 (l.183); \texttt{Limiting\_cases\_final.tex} l.153; \texttt{Final\_terms\_checked.tex} l.224--239.

\noindent\textbf{Limits of the $\xi$ integral} (Group II Row IV l.93): $\int_{k^{+}/\vee}^{k^{+}/(k^{+}+\Lambda)}\frac{d\xi}{1-\xi}=L+\ell+{\cal O}(\frac{\Lambda}{k^{+}})$, $\int\frac{d\xi}{\xi}=\log\frac{\vee}{k^{+}}+{\cal O}(\frac{\Lambda}{k^{+}})$; the lower limit $\frac{k^{+}}{\vee}$ is kept in the intermediate files (``with $\int_{0}^{1}$ the equations of this section would differ from the starting expression by terms of order $k^{+}/\vee$''); checked numerically with a test function (l.175). The $\frac{2}{\epsilon}$ poles are kept in this file: their coefficient depends on $k$ and on the target, so they cannot be dropped as scaleless (l.143). $d_{\perp}=2-\epsilon$ in the polarisation sum; with $d_{\perp}=2$ the bracket would get an extra $+1$ (l.86).

\noindent\textbf{Plus prescriptions and the remnant.} With $f(\xi)\equiv\frac{1}{\xi^{2}}\int e^{-ik\cdot(y-y')/\xi}[\dots]$,
\begin{align}
f(0)=-\frac{4\pi}{k^{2}}\int_{x,x',z}\frac{(x-z)\cdot(x'-z)}{(x-z)^{2}(x'-z)^{2}}[U^{ce'}(z)-U^{ce'}(x')][U^{ce}(z)-U^{ce}(x)]\rho^{e'}(x')\rho^{e}(x)\notag
\end{align}
(from $y'-y=\xi s$ and $\int d^{2}s\,e^{-ik\cdot s}\log s^{2}=-\frac{4\pi}{k^{2}}$, verified numerically; Divergent l.86). $f(0)\neq0$ although the integrand without $\frac{1}{\xi^{2}}$ vanishes at $\xi=0$; without the subtraction the integral diverges as $\log\frac{\vee}{k^{+}}$. Since $\int_{a}^{1}d\xi\,\frac{1-\xi}{\xi}=\log\frac{\vee}{k^{+}}-(1-a)$, $a=\frac{k^{+}}{\vee}$, exactly (Cancellation Sec.~2.3(a)), the compensating term of $N_{\rm DGLAP}$ is minus (BFKL term + remnant) of the row; the two readings of $\big[\frac{1-\xi}{\xi}\big]_{+}$ differ by exactly the remnant term $-\frac{1}{k^{+}}\times\frac{g^{4}N_{c}}{2\pi^{3}k^{2}}\times$(BFKL structure), because $\int_{0}^{1}d\xi\,[\frac1\xi-\frac{1-\xi}{\xi}]=1$ (``model check $10^{-16}$'', Limiting\_cases\_checked l.183; ``difference $0$ to machine precision; prefactor ratio $1$'', Limiting\_cases\_final l.153). Decision: subtraction inside the splitting function, remnant counted as finite (Sec.~\ref{sec:decisions}, item 6). The lower limit can be set to $0$: $|I(a)-I(0)|=2.7\times10^{-2},\,2.7\times10^{-4},\,2.7\times10^{-6}$ for $I(a)=\int_{a}^{1}d\xi\frac{1-\xi}{\xi}e^{-is/\xi}$ at $s=0.37$ (Limiting\_cases\_checked l.183).

\noindent\textbf{Running coupling} (Divergent l.112; Final\_terms l.239; Finite\_terms\_final Sec.~2.4(ii)). $N_{\rm RC}$ has the prefactor $\frac{1}{(2\pi)^{3}}\frac{g^{2}}{2\pi^{2}}\frac{1}{k^{+}}$ and the LO integrand times $\frac{g^{2}\beta_{0}}{16\pi^{2}}\big[2\gamma_{E}+\log\frac{r^{2}\mu^{2}}{4}\big]$; it is exactly minus the $\delta(1-\xi)$ term of $N_{\rm DGLAP}$, so $N_{\rm DGLAP}+N_{\rm RC}$ does not depend on $\beta_{0}$ (Group II Row IV l.165: ``This real-emission term alone does not produce $\beta_{0}$''). Combined with LO:
\begin{align}
\frac{d^{3}N_{\rm LO}}{d^{3}k}+\frac{d^{3}N_{\rm RC}}{d^{3}k}=\frac{1}{(2\pi)^{3}}\frac{g^{2}(\mu)}{2\pi^{2}}\frac{1}{k^{+}}\int e^{-ik\cdot(y-y')}\Big[1+\frac{\alpha_{s}(\mu)\beta_{0}}{4\pi}\log\frac{r^{2}\mu^{2}e^{2\gamma_{E}}}{4}\Big]\dots\,,\notag
\end{align}
and to this order the bracket is $\alpha_{s}(Q^{2})/\alpha_{s}(\mu^{2})$ with $Q^{2}=4e^{-2\gamma_{E}}/r^{2}$, asymptotically free for $\beta_{0}>0$. $g^{2}\to g^{2}(\mu)$ holds only together with the $\delta(1-\xi)$ term and the same $\mu$; the RC pole is collinear (created by adding and subtracting the $\delta$ term).

\noindent\textbf{DGLAP phase} (Final\_terms l.224, l.229): the factor $\frac1\xi$ belongs in the exponent, $e^{-ik\cdot(y-y')\frac1\xi}$ (the measure is $\frac{d\xi}{\xi^{2}}$ only), as in the checked Group II Row IV. \textbf{Numerical check} of DGLAP, RC and the finite term of the row together (Finite\_terms\_final Sec.~2.4): [row with cut-offs, pole dropped] $-$ [evolution logs] $-$ (i) $-$ (iii) $=0.145,\,1.5\times10^{-2},\,1.5\times10^{-3},\,1.5\times10^{-4}$ for $\frac{\Lambda}{k^{+}}=\frac{k^{+}}{\vee}=10^{-2},\dots,10^{-5}$.

\subsection{The recoil term of Group I Row III}\label{sec:calc-recoil}
\noindent\emph{Source:} \texttt{Finite\_terms\_final.tex} Eqs.~(III), (IIIA1)--(IIIA3); \texttt{Cancellation\_of\_flagged\_terms.tex} Sec.~1.3; \texttt{Limiting\_cases\_checked.tex} FLAG 1--3; \texttt{Limiting\_cases\_final.tex} l.78; \texttt{Final\_terms\_checked.tex} l.312--313.

\noindent\textbf{The problem.} The last term of $N_{3b}$, $-\frac{ig^{4}f^{abd}}{16\pi^{5}k^{+}}\int\log\frac{(z-w)^{2}}{(x-w)^{2}}\frac{(x'-w')\cdot(y-z)}{(x'-w')^{2}(y-z)^{2}(z-w)^{2}}\times$[colour bracket with $U^{bc}(z)$], is the limit $\Lambda\to0$ at fixed $z\neq w$ only. Its $z$ integral diverges at $z\to w$: $I(\varepsilon)=\int_{\varepsilon<|z-w|}\frac{d^{2}z}{(z-w)^{2}}\log\frac{(z-w)^{2}}{(x-w)^{2}}\Phi(z)=-2\pi\Phi(w)\log^{2}\frac{\varepsilon}{|x-w|}+$ const; ``numerically $I$ goes from $13.7$ to $976$ for $\varepsilon=10^{-1}\dots10^{-8}$ and $I+2\pi\Phi(w)\log^{2}\frac{\varepsilon}{|x-w|}=-1.386305$ for all of them'' (FLAG 2). With $\Lambda$ kept and together with its partner $-\big[L+\log(1-\frac{\Lambda}{k^{+}})\big]\frac{(x'-w')\cdot(y-z)}{(x'-w')^{2}(y-z)^{2}(z-w)^{2}}$ of $N_{3a}$ it converges at $z\to w$ but grows like $2\pi\Phi(w)\log^{2}\frac{k^{+}}{\Lambda}$ ($-6.4,\,-43.9,\,-111.4,\,-209.3$ for $\frac{\Lambda}{k^{+}}=10^{-1}\dots10^{-4}$).

\noindent\textbf{The identity.} With $\lambda=\Lambda/k^{+}$ the two terms are $-\frac12\int_{z}\frac{\mathcal{G}_{\Lambda}}{(z-w)^{2}}\Phi(z)$, $\mathcal{G}_{\Lambda}=\log\frac{[(x-w)+\frac{1-\lambda}{\lambda}(z-w)]^{2}}{[(x-w)+\frac{\lambda}{1-\lambda}(z-w)]^{2}}$. The substitution $z=w-\xi(z_{s}-w)$ ($z_{s}$: position of the soft gluon) and, under $\int_{w}e^{-ik\cdot w}$, the change to the hard position $w_{h}=w-\xi(z_{s}-w)$, $t=\frac{\xi}{1+\xi}$, give, exactly in $\lambda$,
\begin{align}
\int_{w}e^{-ik\cdot w}\Big(-\frac12\Big)\int_{z}\frac{\mathcal{G}_{\Lambda}}{(z-w)^{2}}\Phi(z)=-\int_{w,z}e^{-ik\cdot w}K(x,w,z)\,\Phi(w)\int_{\lambda}^{1-\lambda}\frac{dt}{t}e^{-it\,k\cdot(z-w)}
=-\int_{w,z}e^{-ik\cdot w}K(x,w,z)\Phi(w)\Big[L+\mathcal{E}\big(k\cdot(z-w)\big)+{\cal O}(\lambda)\Big]\,,\notag
\end{align}
$\mathcal{E}(s)\equiv\int_{0}^{1}\frac{dt}{t}(e^{-ist}-1)=\mathrm{Ci}(|s|)-\gamma_{E}-\log|s|-i\,\mathrm{Si}(s)$. The term with $L$ is the JIMWLK-list term in the author's form (kernel $-K(x'-w')\cdot K(y-w)K(x,w,z)$, Wilson line $U^{bc}(w)$); the rest is the recoil term that replaces the last term of $N_{3b}$:
\begin{align}
-\frac{ig^{4}f^{abd}}{8\pi^{5}}\frac{1}{k^{+}}\int_{x,y,z,x'}\mathcal{E}\big(k\cdot(z-w)\big)\frac{(x-z)\cdot(w-z)\,(x'-w')\cdot(y-w)}{(x-z)^{2}(w-z)^{2}(x'-w')^{2}(y-w)^{2}}\big[\dots\big]\big[\dots U^{bc}(w)\dots\big]\,.\notag
\end{align}
\textbf{Evidence}: the first line checked with a Gaussian regulator: ``$-10.896402+0.822556\,i$ on both sides for $\lambda=0.2$ (difference $3\times10^{-13}$), $-34.188243+1.392118\,i$ for $\lambda=0.05$ (difference $3\times10^{-10}$)'' (Finite\_terms\_final, under (IIIA3)); re-run for Limiting\_cases\_final (l.78): difference $3\times10^{-13}$; $\mathcal{E}$ closed form against quadrature $10^{-16}$. This resolved the point left open in the evolution files (three-$\rho$ l.68, JIMWLK l.364, Other\_terms\_v2 l.109): the evolution lists' form with $U^{bc}(w)$ is the right $\log\frac{k^{+}}{\Lambda}$ coefficient; the pointwise form with $U^{bc}(z)$ is the same term written in the recoil variable (Cancellation Sec.~1.3).

\noindent\textbf{Combined kernel.} $N_{1}$, $N_{2}$, $N_{3a}$, $N_{3b}$ are not integrable one by one: at distance $\delta$ from the line $(x-z)^{2}=(x-w)^{2}$, $|N_{2}|\simeq|N_{3a}|=0.89,\,1.2\times10^{2},\,1.6\times10^{4}$ and $|N_{1}|\simeq|N_{2}+N_{3a}|=3.0\times10^{-4},\,3.0\times10^{-3},\,3.0\times10^{-2}$ for $\delta=10^{-3},10^{-4},10^{-5}$, while the sum of the four is $6.0\times10^{-7}$ at all three (FLAG 1; Limiting\_cases\_final l.78). Their sum (24 terms plus the first two terms of $N_{3b}$) is the combined kernel $\mathcal{K}^{ij}$ written in Limiting\_cases\_final (l.62--74); ``[$N_{1}+N_{2}+N_{3a}+$ first three terms of $N_{3b}$] $-$ [first $+$ second term] $=2\times10^{-22}$ to $7\times10^{-18}$ against $7\times10^{-9}$ to $5\times10^{-6}$ at eight points'' (l.78). The third term of $N_{3b}$ is [LO kernel] $-$ [recoil-shifted LO kernel] (Finite\_terms\_final Eq.~(IIIB)) with a $-\frac{1}{z^{2}}\log\frac{|z|}{|x-w|}$ tail; its cancellation is Step 5 of Sec.~\ref{sec:fix-soft}.

\subsection{Group III Rows III, V, VI: the three-term BFKL list}\label{sec:calc-G3}
\noindent\emph{Source:} \texttt{BFKL\_checked.tex} l.464--568; \texttt{Other\_terms\_checked.tex} l.212--236; \texttt{Other\_terms\_v2\_checked.tex} l.214--238; \texttt{Limiting\_cases\_checked.tex} FLAG 6--9; \texttt{Limiting\_cases\_final.tex} l.181--212.

\noindent The author's split $-\frac{1}{p^{+}}=-\frac{k^{+}}{p^{+}(p^{+}+k^{+})}-\frac{1}{p^{+}+k^{+}}$, with $\int_{\Lambda}^{\vee-k^{+}}\frac{k^{+}dp^{+}}{p^{+}(p^{+}+k^{+})}=L+\ell$, and the subtracted $y'\to z$ term with $\int_{y'}e^{-ik\cdot(w-y')}e^{-ik\cdot\frac{p^{+}}{k^{+}}(z-y')}\frac{(y'-z)^{m}}{(y'-z)^{2}}=\frac{k^{+}}{p^{+}+k^{+}}\int_{w'}e^{-ik\cdot(w-w')}\frac{(w'-z)^{m}}{(w'-z)^{2}}$ (checked to $10^{-11}$), takes out the complete $\log\frac{\vee}{k^{+}}$ of these rows: all three BFKL terms of the row (the first two cancel each other, the third is the net coefficient). With the plain $[\frac{1}{p^{+}}]_{+}$ split the left-over would still depend logarithmically on $\vee$. Numerical check (Other\_terms\_v2 l.214, Row III first part, $\frac{\vee}{k^{+}}=6$): left-hand side $0.0044667+0.0226868i$, $0.0051661+0.0230997i$, $0.0051731+0.0231038i$ for $\frac{\Lambda}{k^{+}}=10^{-2},10^{-4},10^{-6}$; the equation $0.0051732+0.0231038i$; without the second and third BFKL terms the difference is $0.16$. Second part (l.225): $-0.0048167-0.0071862i$ against the limit $-0.0048167-0.0071862i$; Row V (l.233): $0.0103463+0.0462076i$. The form of Limiting\_cases (other split) still contains $\log\frac{\vee}{k^{+}}\times$(second + third BFKL term): ``this form $=0.455-0.328i$, $1.035-0.193i$, $1.535+0.029i$, $2.084+0.296i$; form of Other\_terms $=-0.047-0.581i$, \dots, $-0.019-0.766i$ (converges)'' for $\frac{\vee}{k^{+}}=3,10,30,100$ (FLAG 6). The replacement in Limiting\_cases\_final converges like $\frac{k^{+}}{\vee}$: ``$0.0201-1.0301i$, $0.0389-1.1011i$, $0.0419-1.1183i$, $0.0430-1.1250i$ for $\frac{\vee}{k^{+}}=10,30,60,100$, extrapolated limit $0.0445-1.1352i$'' (l.186); condition: in the term with the constant kernel the $y'$ integral is done before the $p^{+}$ integral.

\subsection{Group I Row IV $+$ Group III Row I $=0$}\label{sec:calc-cancelIV}
\noindent\emph{Source:} \texttt{Limiting\_cases\_final.tex} l.160--172; \texttt{Limiting\_cases\_checked.tex} l.116; \texttt{Cancellation\_of\_flagged\_terms.tex} Sec.~2.2.

\noindent With $U^{cb}(z)U^{c'b}(z)=\delta^{cc'}$ the colour factor of Group III Row I becomes $[U^{ab'}(x')\rho^{b'}(x')-U^{ab'}(w')\rho^{b'}(x')][f^{cbd}U^{ac}(w)U^{c'b}(z)U^{c'd'}(y)\rho^{d'}(y)\rho^{d}(x)-f^{acd}U^{cd'}(y)U^{de}(x)\rho^{d'}(y)\rho^{e}(x)]$, which is that of Group I Row IV after renaming summed indices, with the same order $\rho(x')\rho(y)\rho(x)$; logarithm, kernel, variables and phase are the same, and the prefactors are $\pm\frac{ig^{4}}{8\pi^{5}}\frac{1}{k^{+}}$. Numerically (random points, SU(3)): sum $=3\times10^{-24}$, $1\times10^{-23}$, $3\times10^{-22}$, $2\times10^{-23}$ against $8.8\times10^{-9}$, $2.9\times10^{-8}$, $3.9\times10^{-7}$, $2.7\times10^{-9}$ (l.172). The identity holds also with the cut-offs and the $\log(1-\frac{k^{+}}{\vee})$ terms (Limiting\_cases\_checked l.116).

\subsection{Soft sector: $\Lambda\to0$ and the double logarithms}\label{sec:calc-soft}
\noindent\emph{Source:} \texttt{Cancellation\_of\_flagged\_terms.tex} Sec.~1; \texttt{Finite\_terms\_final.tex} Sec.~1.6.

\noindent Every flagged soft term is, row by row, $\int_{\Lambda}\frac{dp^{+}}{p^{+}}\int_{z}\frac{T_{\rm row}}{z^{2}}F_{\rm row}(\frac{p^{+}}{k^{+}}z)$ with a cut-off function produced by the recoil of the measured gluon, so each row alone has $\pi L^{2}T_{\rm row}+{\cal O}(L)$. The sum of the tails over the recoil rows vanishes, $\Delta\equiv\sum_{\rm recoil\ rows}T_{\rm row}=0$, for a general projectile in the symmetrised two-$\rho$ and three-$\rho$ sectors [N: SU(2), SU(3), $|\Delta|<10^{-7}$ against single-row tails of order 1--6; changing any one weight gives $|\Delta|={\cal O}(1)$]. Table (tails by degree in $U(z)$, SU(3)): Group III Row III first part $0/1.46/0.92$, second part $0/0.68/3.42$, Row V $3.21/0/0$, Row VI $0/4.25/4.34$, Group I Row V first part $0.92/1.46/0$, $N_{2a}$ $1.41/1.95/0$, $N_{2b\,\rm rest}$ $4.14/1.81/0$, Group I Row III $\frac12K(x,y,z)$ term $2.17/2.12/0$ (weight 2), $z\to w$ term $3.21/0/0$; sum $2\times10^{-8}$, $1\times10^{-8}$, $3\times10^{-8}$. Rows built only from $\mf{B}_{2}$ and $\mf{A}^{(1)}$ have tails that are not cut at all, by the identity $\frac{1}{p^{+}[p^{+}z^{2}+k^{+}r^{2}]}+\frac{z^{2}}{k^{+}r^{2}[p^{+}z^{2}+k^{+}r^{2}]}=\frac{1}{k^{+}p^{+}r^{2}}$. What remains at large $z$ is the infrared divergence of the JIMWLK kernel only, $L\int\frac{d^{2}z}{z^{2}}C_{A}$, minus the Casimir of the pair acting on the LO structure, zero for $\langle\rho^{a}\rho^{b}\rangle\propto\delta^{ab}$. The pairing of the remainders (and the correction of the statement that the $U(z)$-linear tails cancel only for direction-independent $U(z)$) is in Sec.~\ref{sec:fix-soft}.

\subsection{Energetic sector: $\vee\to\infty$ and large-distance tails}\label{sec:calc-energetic}
\noindent\emph{Source:} \texttt{Cancellation\_of\_flagged\_terms.tex} Sec.~2; \texttt{Finite\_terms\_final.tex} OPEN II/III-1.

\noindent Each $\log(1-\frac{k^{+}}{\vee})$ term is $\log(1-a)$ times the row's JIMWLK structure (the correction to the upper limit $\vee-k^{+}$ of the soft integral). At $\vee=\infty$ none of the rows has a $1/z^{2}$ tail at fixed $w,w'$; the $\log(1-a)$ tails do not cancel among themselves (e.g.\ SU(2): two-$\rho$ norm $1.6\times10^{-6}$ against $0.3$--$1.2\times10^{-6}$ for the rows); Group I Row IV and Group III Row I cancel each other (to $10^{-14}$ relative). Pure numbers times LO/BFKL/JIMWLK structures in the left-over terms: Group I Row II $N_{1}$: $-\frac{67}{9}+\frac{\pi^{2}}{3}$ (LO type); Group I Rows II + V: $\frac{11}{3}-2\cdot\frac{11}{6}=0$; Group II Row IV remnant $-1$ and compensator $+1$ (BFKL type); sum: BFKL type $0$, LO type $(-\frac{67}{9}+\frac{\pi^{2}}{3})\mathcal{S}_{\rm LO}$. Group II Row IV: with the slow-variation assumption the left-over terms have the large-distance tail $2N_{c}C\frac{4\pi}{k^{2}}[-\frac{11}{6}+\frac16\cos2\varphi_{kY}]\frac{G_{\infty}(Y)}{Y^{2}}$ ($\int_{0}^{1}d\xi[\xi(1-\xi)-1+\frac{\xi}{(1-\xi)_{+}}]=-\frac{11}{6}$); confirmed in a model in Finite\_terms\_final (the subtracted quantity $66.70$--$73.75$ converging like $1/R^{2}$ for $R/\rho=8\dots128$). Its cancellation at amplitude level is Step 14 of Sec.~\ref{sec:fix-poles}.

\subsection{Literature comparisons}\label{sec:calc-lit}
\noindent\emph{Source:} \texttt{paper/notes/README.md}, \texttt{lit\_lo\_bfkl\_jimwlk.tex}, \texttt{lit\_dglap\_rc.tex}, \texttt{lit\_finite\_dilute.tex}.
\begin{itemize}
\item \textbf{LO, BFKL, JIMWLK} (lit\_lo\_bfkl\_jimwlk, Summary): LO agrees exactly with Kovchegov--Tuchin hep-ph/0111362 (prefactor $\alpha_{s}C_{F}/4\pi^{4}$ for a $q\bar q$ projectile); dilute target reproduces $k_{T}$ factorisation and the Lipatov vertex (Horowitz--Kovchegov); Kovner--Lublinsky hep-ph/0609227: structure agrees, normalisation there internally inconsistent. BFKL term = one KLWMIJ step (hep-ph/0501198 Eq.~(5)) including normalisation at finite $N_{c}$ ($10^{-13}$ on explicit singlet states); dipole kernel $\alpha_{s}N_{c}/2\pi^{2}$. JIMWLK term $=-H_{\rm JIMWLK}$ (KLM 1310.0378) on the LO structure, all eight kernel blocks ($10^{-12}$; blocks $10^{-15}$); BK limit; split of the logarithms agrees with KL, GLV 0804.2630, Gelis--Laidet 1211.1191. No discrepancy found in LO, BFKL, JIMWLK.
\item \textbf{DGLAP and RC} (lit\_dglap\_rc, Summary table): $P_{gg}$ normalisation and $\frac{\beta_{0}}{2}\delta$ agree with Chirilli--Xiao--Yuan 1203.6139; the $\xi\to0$ subtraction with the remnant equals the KLSZ (2308.15545) ``double plus'' prescription; $\delta$ term $=-N_{\rm RC}$ (derived; $\beta_{0}$ added and subtracted); the argument of the collinear logarithm differs from CXY/AABKL $\overline{\rm MS}$ (measured-gluon separation $(y'-y)/\xi$ vs parent separation): finite scheme term $\int d\xi\,\widetilde{\mathcal P}\log\xi^{2}\otimes$LO; the author's bracket equals the AABKL cut-off scheme with $\mu_{\overline{MS}}^{2}=1.26\,\mu^{2}_{\rm CO}$ (README); LO+RC $=g^{2}(c_{0}^{2}/r^{2})$, the scale being the choice $\mu=c_{0}/r$; ``RC from DGLAP'' consistent only as ``$\beta_{0}$ of collinear origin''.
\item \textbf{Finite terms, dilute limit} (lit\_finite\_dilute, README): Group I Row II with c.c.\ gives $\frac{\alpha_{s}N_{c}}{2\pi}\big(\frac{67}{9}-\frac{\pi^{2}}{3}\big)$LO, equal to the bare-gluon NLO JIMWLK kernel of KLSZ and twice the CMW/BC/FL/BSV value, which is inclusive over the collinear pair (the other half should sit in the collinear finite terms: not verified). Dilute limit vs BSV jet vertex: structural mapping only; Row IV collinear, $\xi\to0$ (BFKL) and $\xi\to1$ (JIMWLK) limits agree; $\log\mu^{2}$ coefficient $\frac{\alpha_{s}\beta_{0}}{4\pi}$ agrees. Sanity checks ($U=1$, reality of Group II, scaling, commuting $\rho$) pass. Open there: the Group II Row IV large-distance tail (now Sec.~\ref{sec:fix-poles}, Step 14).
\end{itemize}
\clearpage
\section{Open points and the fixes of the problems P1--P6}\label{sec:open}

\subsection{Open points recorded in the checked files}
\noindent Status as written in the files; where a later document settles a point, this is said.
\begin{enumerate}
\item \textbf{Support of the commutator term of $\mf{B}_{2}$} (Calculation\_main l.1261, l.1269): ``it is left unchanged here for the author to decide''.
\item \textbf{Two-$\rho$ section of the four-$\rho$ file}: ``Not confirmed'' (4rho l.130: does not vanish at $U=1$, index $b'$ four times, scaleless $\int_{z}1/(x'-z)^{2}$). Not used; replaced by an independent reordering (decision 12).
\item \textbf{Form of the $\log\frac{k^{+}}{\Lambda}$ term of Group I Row III} ($U^{bc}(z)$, kernel $\frac{K(x'-w')\cdot K(y-z)}{(z-w)^{2}}$, versus $U^{bc}(w)$, kernel $K(x'-w')\cdot K(y-w)K(x,w,z)$): open in Evolution\_three\_rho l.68, JIMWLK l.364, Other\_terms\_v2 l.109 (residual $4$--$9\times10^{-7}$ with the list as written). \emph{Settled} by the recoil identity (Sec.~\ref{sec:calc-recoil}): the evolution lists are right, the difference is the recoil term.
\item \textbf{Two scaleless terms in the JIMWLK list of Group II Row IV} ($\int_{z}\frac{1}{(w-z)^{2}}$, $\int_{z}\frac{1}{(w'-z)^{2}}$, z-independent coefficient): ``whether they are to be kept is left open'' (JIMWLK l.194). The constant of the UV bracket in the JIMWLK file ``was not checked'' (l.36).
\item \textbf{BFKL remark} (BFKL l.633; Final\_terms l.100): with non-commuting $\rho$ the evolution gives $\frac12\{\rho(x'),\rho(x)\}$ in four blocks; the fixed orders written differ by one-$\rho$ contact terms $\propto f^{abc}\rho^{c}(x)\delta^{2}(x-x')$, which drop out of the colour-singlet form. Flagged, not changed.
\item \textbf{Power terms} (scales l.58, l.64): the $-\frac{1}{k^{+}}$ of $N_{2f}$ (Group II Row I) and of Group II Row IV: ``to be listed or not: for the author''. Decided for Row IV (decision 6); for $N_{2f}$ the limit $[\frac1\vee-\frac{1}{k^{+}}]\to-\frac{1}{k^{+}}$ is taken in the finite term (Limiting\_cases\_checked l.174).
\item \textbf{Order of limits for the $\ell$ poles} (Divergent l.68, l.90): a pole times $\ell\to0$ is zero only if $\vee\to\infty$ is taken first. Decision: no poles written (decision 3).
\item \textbf{Flags 1--9 of Limiting\_cases\_checked}: Group I Row III (FLAG 1--3), Row V (FLAG 4), Group II Row IV second term (FLAG 5), Group III Rows III, V, VI (FLAG 6--9). Settled by the decisions implemented in Limiting\_cases\_final (items 24--27 of Section~\ref{sec:summary}) except for the integrability of the soft subtraction terms, which are ``finite only in the sum over rows'' (Limiting\_cases\_final l.91, 95, 107, 186, 197, 206, 212): this is P1/P2.
\item \textbf{OPEN items of \texttt{Finite\_terms\_final}}: I-1 (combined kernel converges at large $z$ only conditionally), I-2 (recoil terms not integrable one by one; $U(z)$-linear parts cancel only between the $w$ and $w'$ sides), I-3 (dropped $\log(1-\frac{k^{+}}{\vee})$ term of Row IV), II/III-1 (Group II Row IV tail $\propto[-\frac{11}{6}+\frac16\cos2\varphi]G_{\infty}$), II/III-2 ($N^{\Psi}$ terms not integrable one by one), II/III-3 ($U(z)$-linear tails), II/III-4 ($\log(1-\frac{k^{+}}{\vee})$ terms, non-singlet JIMWLK term, poles not among the finite terms). Updated by the fix notes: I-1 is integrable after the $w$ integration (fix\_soft, condition (C3)); I-2(b) and II/III-3 are \emph{corrected} by the pairing of the recoil remainders (fix\_soft, Step 4: no condition on $U(z)$ is needed); II/III-1 is cancelled at amplitude level by Group II Row I (second part) and Group III Row III (fix\_poles\_dglap, Step 14; term-by-term check on the finite terms as written still open).
\item \textbf{Not established in \texttt{Cancellation\_of\_flagged\_terms}}: plateaux of the phase-cut rows analytic only; finite remainder of the combination not evaluated by a full integration; tails of rows with a $\frac{2}{\epsilon}$ bracket depend on $\int d^{2}p/p^{2}=0$.
\item \textbf{Literature} (lit\_finite\_dilute): the second half of the CMW constant in the collinear finite terms was not verified; the full dilute limit was not compared term by term with Fadin--Lipatov/BSV; scheme dependence through $d_{\perp}$ (Group I Row V, Group II Row I $N_{2f}$, Group III Row III).
\end{enumerate}

\subsection{The problems P1--P6 and where they are fixed}
\noindent From \texttt{paper/BRIEF\_fix.md}: P1 Group I Row III (non-integrable summed kernel; ``the one problematic term should cancel with something''); P2 Group III Rows III, V, VI ($\log|z|/z^{2}$ tails of the soft remainders; distributional form of the instantaneous terms); P3 virtual $\frac{11}{3}$ terms and poles; P4 DGLAP form ($\log\xi^{2}$, explicit $\xi\to0$ subtraction); P5 explicit $\log|\cdot|$, distributional form, one phase convention; P6 finite terms averaged over the projectile, colour-singlet form. The three notes reproduced below are \texttt{paper/notes/fix\_soft.tex} (P1, P2, P5 parts), \texttt{paper/notes/fix\_poles\_dglap.tex} (P3, P4, P5 phase, Group II Row IV tail) and \texttt{paper/notes/fix\_singlet.tex} (P6), copied into this document as they stood when it was compiled (\texttt{fix\_soft.tex} of 20 Nov 2026 06:40 UTC; \texttt{fix\_poles\_dglap.tex} of 20 Nov 2026 06:34 UTC; \texttt{fix\_singlet.tex} of 20 Nov 2026 06:29 UTC; \texttt{fix\_tail.tex} of 20 Nov 2026 13:45 UTC; \texttt{fix\_dabc.tex} of 20 Nov 2026 14:20 UTC). Their section levels are lowered by one; nothing else is changed.

\subsection{P1, P2 (and P5): soft-gluon remainders --- summary and resulting changes}\label{sec:fix-soft}
\noindent\textbf{Result} (fix\_soft, first paragraph and Steps 3--5). The soft remainders cancel in pairs, each Group I term (virtual, recoil on $w$) with a Group III partner (real, recoil on $w'$) of the same colour structure and kernel with $w\to w'$: (A) Group I Row III, last term of $N_{3b}$ (recoil form) $\leftrightarrow$ Group III Row V; (B) Group I Row V first part $\leftrightarrow$ Group III Row III first part; (C) Group I Row V second part ($N_{2a}+N_{2b\,\rm rest}$) $\leftrightarrow$ Group III Row III second part; (D) Group I Row III, third term of $N_{3b}$ $\leftrightarrow$ Group III Row VI. The two recoil functions $\mathcal{E}(k\cdot(z-w))$ and $\Psi(k\cdot(z-w'))$ have the same large-distance form $-\gamma_{E}-\log|k\cdot z|-\frac{i\pi}{2}{\rm sgn}(k\cdot z)$, so each pair is integrable for any Wilson line at the soft position:
\begin{align}
\kappa(w)\mathcal{E}\big(k\cdot(z-w)\big)-\kappa(w')\Psi\big(k\cdot(z-w')\big)=[\kappa(w)-\kappa(w')]\,\mathcal{E}\big(k\cdot(z-w)\big)+\kappa(w')\big[\mathcal{E}\big(k\cdot(z-w)\big)-\Psi\big(k\cdot(z-w')\big)\big]\,.\notag
\end{align}
This \emph{corrects} \texttt{Finite\_terms\_final} OPEN I-2(b) and II/III-3 (there the $w$-side terms of Group I were paired with the c.c.\ terms of Group III, which carry $\Psi^{*}$). \textbf{The problematic term} (answer to the author's wish, Step 5): the third term of $N_{3b}$, with tail $-\frac{T_{D}}{z^{2}}\log\frac{|z|}{|x-w|}$, is cancelled (1) in its $\log|z|$ part by the soft remainder of Group III Row VI (same coefficient, $w\to w'$), and (2) in its remaining $1/z^{2}$ part, $\frac{T_{D}}{z^{2}}[\gamma_{E}+\log(|k||x-w|)+c(\theta)]$, by the odd-angle (``saw'') tail of the other 24 terms of Row III after the $w$ integration: $\int d^{2}r\,e^{ik\cdot r}K^{i}(r)[\gamma_{E}+\log(|k||r|)+c(\theta)]=+\frac{2\pi i\hat k^{i}}{|k|}\chi(\theta)$ and $\int d^{2}r\,e^{ik\cdot r}\varphi\frac{e_{\perp}^{i}}{|r|}=-\frac{2\pi i\hat k^{i}}{|k|}\chi(\theta)$, $\chi=\log|2\cos\theta|+\frac{i\pi}{2}{\rm sgn}\cos\theta$ [N: $z^{2}\times|{\rm sum}|=0.72,\,0.26,\,0.087,\,0.027$ at $|z|=40,160,640,2560$, $\theta=0.4$]. The last term of $N_{3b}$ (recoil form) is cancelled by Group III Row V.

\noindent\textbf{Resulting changes to the author's tex files} (fix\_soft, Step 6 and Steps 8, 9; to be applied to \texttt{Final\_terms.tex}, equation numbers of the compiled author's file as used in the reviews):
\begin{itemize}
\item Group I Row III, Eqs.~(6.4)--(6.7): replaced by Eq.~\eqref{sf:GIIII}: combined kernel $\mathcal{K}^{ij}$ (only positive arguments in the logarithms) plus the bracket $\big[-\frac12\log\frac{(z-w)^{2}}{(x-w)^{2}}-\mathcal{E}(k\cdot(z-w))\big]K(y,w,z)K(x'-w')\cdot K(x-w)$, with the $w$ integral done before the $z$ integral; the two soft remainders go into the block $\mathcal{S}$.
\item Group I Row V, Eqs.~(6.9)--(6.11): replaced by Eqs.~\eqref{sf:V1}, \eqref{sf:V2} (regular terms with $\mathcal{J}(s)$, $\mathcal{F}^{im}_{\rm reg}$ and the $\frac{11}{6}$ term); soft remainders into $\mathcal{S}$.
\item Group III Rows III, V, VI, Eqs.~(8.2)--(8.5): replaced by Eqs.~\eqref{sf:G3R31}, \eqref{sf:G3R32}, \eqref{sf:G3R5} (Row VI from \eqref{sf:G3R5} with the colour factor changed), with $\widehat{\mathcal{E}}(s)$ and the instantaneous terms in distributional form $\int_{0}^{\infty}dp^{+}e^{-i\frac{p^{+}}{k^{+}}s}=k^{+}[\pi\delta(s)-i{\rm P}\frac1s]$ ($y'$ integral before $p^{+}$).
\item New block $\mathcal{S}$, Eq.~\eqref{sf:S}: all soft remainders of the five rows, pairs A--D, with $\mathfrak{F}[\kappa]=\kappa(w)\mathcal{E}-\kappa(w')\Psi$.
\item P5: in the separate equations $N_{1}$, $N_{2}$, $N_{3a}$ the logarithms whose argument changes sign are $\log|\cdot|$ (principal value of the energy denominators of $\mf{B}_{3}$, Step 8); in the combined kernel the prescription is no longer visible.
\end{itemize}
\noindent Open (fix\_soft, ``Not done / open''): no complete numerical integration over all transverse positions; regular terms taken from \texttt{Finite\_terms\_final} without re-check; non-singlet JIMWLK term and Group II Row IV tail outside this note; $\pm i\epsilon$ alternative of Step 8 not evaluated.

\subsection{P3, P4, P5: poles, DGLAP, running coupling --- summary and resulting changes}\label{sec:fix-poles}
\noindent\textbf{Result} (fix\_poles\_dglap, Summary). Every $1/\epsilon$ is accounted for (pole table, Step 6). The $\frac{11}{3}$ term of Group I Row II is exactly the virtual term $\frac{\alpha_{s}\beta_{0}}{4\pi}[\frac{4}{\epsilon}+\mathcal{B}(x-w)+\mathcal{B}(x'-w')]\times$LO that the reviewer asks for; it is cancelled UV against UV by the $-\frac{11}{6}\times2$ of $N_{2b\,\rm UV}$ of Row V after the relabelling $x\to w$, $y\to x$, so the reviewer's claim that it is missing is wrong; the reviewer is right that the RC pole of Row IV is collinear. The coupling is renormalised by the UV half of the scaleless normalisation of the dressed produced gluon ($Z_{3}$, $Z_{g}=Z_{3}^{-1/2}$ in light-cone gauge); its IR half is the virtual $\frac{\beta_{0}}{2}\delta(1-\xi)$ that completes $\mathcal{P}_{gg+}$ to the $\overline{\rm MS}$ $\mathcal{P}_{gg}$. No pole is left except $-\frac{\ell}{2}\frac{2}{\epsilon}$, which vanishes for $\vee\to\infty$ taken first. DGLAP: in full CDR the bare logarithm has no $\xi^{2}$; the author's $\log\xi^{2}$ reappears from the $d$-dimensional Jacobian $\xi^{-d}$ of the $\overline{\rm MS}$ counterterm for number-density fragmentation functions; author's and CXY form differ by $\Delta$ (Eq.~\eqref{eq:pd-Delta}). Group II Row IV large-distance tail: soft-parent region, cancelled by Group II Row I (second part) and Group III Row III through $\mf{B}_{2}-\mf{C}_{1}\mf{A}^{(1)}$, whose parent pole cancels [N: coefficients of $1/q^{2}$ of $|{\rm B}|^{2}$, $|{\rm C}|^{2}$, $2{\rm B}\cdot{\rm C}$: $58.2183$, $58.2165$, $2\times58.2174$]. Phase: all working files used the reversed phase (counts in Sec.~\ref{sec:pd-phase}); the checked copies are corrected.

\noindent\textbf{Resulting changes to the author's tex files}: DGLAP term written in the final form Eq.~\eqref{eq:pd-final} of the note (explicit $\xi\to0$ subtraction), or in the recommended local form of its Step 12, with the remainder $N_{\mathcal{R}}$ grouped with the finite terms of Group II Row I (second part) and Group III Row III under one integral over the pair position; the account of the poles (table) added to the text; one phase convention $e^{-ik\cdot(\text{unprimed}-\text{primed})}$ in all files (already so in \texttt{Final\_terms.tex} and \texttt{Limiting\_cases\_final.tex}). \textbf{Open}: the term-by-term check that the finite terms of Group II Row I (second part) and Group III Row III as written carry the compensating tail $+2N_{c}\mathcal{C}\frac{4\pi}{k^{2}}\int\frac{d^{2}Y}{Y^{2}}[\frac{11}{6}-\frac16\cos2\varphi]G_{\infty}$.

\subsection{P6: colour-singlet form --- summary and resulting changes}\label{sec:fix-singlet}
\noindent\textbf{Result} (fix\_singlet). Under (A1) global colour invariance, $\langle\rho^{a}(x)\rho^{b}(y)\rangle=\frac{\delta^{ab}}{N_{c}^{2}-1}n(x,y)$, and (A2) vanishing symmetrised three-point function (Gaussian projectile), the finite terms, DGLAP and RC are written with traces of adjoint Wilson lines. \textbf{Correction of the brief}: the cubic (three-$\rho$) rows do \emph{not} vanish after the average; only their symmetrised part does. The reordering commutators leave two-point contact terms $\propto\delta^{(2)}(p_{i}-p_{j})\,n$, of the same order as the quadratic terms, even in $k$, not removed by colour neutrality (Eqs.~\eqref{eq:s-ord}--\eqref{eq:s-a12}); the part odd in $k$ is exactly the $S^{abc}$ part. Numerical checks: colour algebra $9\times10^{-15}$ (21 structures); operator level with explicit density matrices on lattice sites, largest relative deviation $5.8\times10^{-15}$ over $7\times21$ cases. \textbf{Resulting changes}: the finite-term section of \texttt{Final\_terms.tex} gets the singlet form of every row (quadratic rows: singlet factor; cubic rows: contact terms written as quadratic terms at coincident points). \textbf{Open}: whether the contact terms cancel among rows; integrability at the contact configurations (checked for two kernels); tails in singlet form not recomputed; (A2) is an assumption.

\clearpage
\noindent\textbf{\large Full text of the fix notes (as compiled).}

% ---- inlined copy of paper/notes/fix_soft.tex (20 Nov 2026 06:40 UTC) ----
\subsection{Soft-gluon remainders: cancellation and integrable form}\label{sec:softfix}
\noindent\textbf{Result in one paragraph.} The soft remainders of Group~I Rows III, V and Group~III Rows III, V, VI cancel in pairs: each Group~I term (soft gluon absorbed by the measured gluon in the amplitude: virtual, recoil on $w$) has one Group~III partner (soft gluon emitted by the measured gluon in the conjugate amplitude: real, recoil on $w'$) with \emph{the same colour structure and the same kernel, except that $w$ is replaced by $w'$}. The pairs are\\[2pt]
\hspace*{1em}(A) Group~I Row III, last term of $N_{3b}$ (recoil form) $\;\leftrightarrow\;$ Group~III Row V;\\
\hspace*{1em}(B) Group~I Row V, first part $\;\leftrightarrow\;$ Group~III Row III, first part;\\
\hspace*{1em}(C) Group~I Row V, second part ($N_{2a}+N_{2b\,\rm rest}$) $\;\leftrightarrow\;$ Group~III Row III, second part;\\
\hspace*{1em}(D) Group~I Row III, third term of $N_{3b}$ (the problematic $\log\frac{(z-w)^{2}}{(x-w)^{2}}$ term) $\;\leftrightarrow\;$ Group~III Row VI.\\[2pt]
In every pair the two recoil functions, $\mathcal{E}(k\cdot(z-w))$ and $\Psi(k\cdot(z-w'))$, have the same large-distance form $-\gamma_{E}-\log|k\cdot z|-\frac{i\pi}{2}{\rm sgn}(k\cdot z)$ (the same sign of the odd part), so the combination is integrable for \emph{any} Wilson line at the soft position; no condition on $U(z)$ at large distance is needed. The third term of $N_{3b}$ is cut off by the shift of the LO kernel and not by the recoil phase; the difference between the two cut-offs leaves a $1/z^{2}$ tail which is cancelled, after the $w$ integration, by the $1/z^{2}$ tail (odd in the angle) of the other 24 terms of the same row. Statements are marked [A] (analytic) or [N] (numerical; scripts in \texttt{fix/soft/v2/}).

\subsubsection{Notation and conventions}
\noindent Phase $e^{-ik\cdot(w-w')}$; the conjugated LO bracket is at $w'$ (left), the amplitude at $w$; c.c.\ implied for Groups I and III; $\rho$'s are operators in the order written. Units and abbreviations:
\begin{align}
&c_{0}\equiv\frac{1}{(2\pi)^{3}}\,\frac{g^{4}}{16\pi^{5}}\,\frac{1}{k^{+}}\,,\qquad K^{i}(v)\equiv\frac{v^{i}}{v^{2}}\,,\qquad K(a,b,z)\equiv K(a-z)\cdot K(b-z)\quad\Big(K(w,w,z)=\frac{1}{(z-w)^{2}}\Big)\,,\notag\\
&\mathcal{L}^{a}\equiv\big[U^{ab'}(x')-U^{ab'}(w')\big]\rho^{b'}(x')\,,\qquad \Delta^{c'e'}\equiv U^{c'e'}(w')-U^{c'e'}(x')\,,\qquad P\equiv K(x'-w')\,.\notag
\end{align}
Recoil functions ($t=p^{+}/k^{+}$; $\mathrm{Ci}$, $\mathrm{Si}$, $E_{1}$: cosine, sine and exponential integrals):
\begin{align}
\mathcal{E}(s)&\equiv\int_{0}^{1}\frac{dt}{t}\big(e^{-ist}-1\big)=\mathrm{Ci}(|s|)-\gamma_{E}-\log|s|-i\,\mathrm{Si}(s)\,,\qquad
\Psi(s)\equiv\int_{0}^{\infty}\frac{dt}{t(1+t)}\big(e^{-ist}-1\big)=-\gamma_{E}-\log|s|-\frac{i\pi}{2}{\rm sgn}(s)-\widehat{\mathcal{E}}(s)\,,\notag\\
\widehat{\mathcal{E}}(s)&\equiv\int_{0}^{\infty}\frac{dt}{1+t}e^{-ist}=e^{is}E_{1}(is)\,,\qquad \mathcal{J}(s)\equiv\int_{0}^{1}dt\,e^{-ist}=\frac{1-e^{-is}}{is}\,,\qquad
\mathcal{E}_{\infty}(s)\equiv-\gamma_{E}-\log|s|-\frac{i\pi}{2}{\rm sgn}(s)\,,\label{sf:fun}\\
\mathcal{E}(s)&=\mathcal{E}_{\infty}(s)-\frac{e^{-is}}{is}+{\cal O}(s^{-2})\,,\qquad \Psi(s)=\mathcal{E}_{\infty}(s)-\frac{1}{is}+{\cal O}(s^{-2})\quad(|s|\to\infty)\,;\qquad \mathcal{E}(s)\simeq-is\,,\ \ \Psi(s)={\cal O}(s\log|s|)\quad(s\to0)\,.\notag
\end{align}
\textbf{``Integrable''} means here: integrable singularities at all coincidence points and, at large soft position $z$, decay faster than $1/z^{2}$ (in practice $\log|z|/|z|^{3}$); a term $\propto e^{-ik\cdot z}/z^{2}$ (the measured momentum flowing through $z$) converges as a Fourier integral, as $\int_{w}e^{-ik\cdot w}K(x-w)$ at LO does, and is not counted as a tail.

\subsubsection{Step 1: where the soft remainders come from}
\noindent In the recoil-carrying rows the measured gluon in the amplitude (Group~I) or in the conjugate amplitude (Group~III) is a two-gluon state, a hard gluon at $w$ (resp.\ $w'$) and a gluon with $p^{+}=tk^{+}$ at $z$. Its position is the $+$-momentum weighted centre, so the phase is $e^{-ik\cdot w}e^{-itk\cdot(z-w)}$ (Group~I) and $e^{+ik\cdot w'}e^{-itk\cdot(z-w')}$ (Group~III: parent at $w'$, measured gluon at $w'-t(z-w')$); in both cases the exponent is $-it\,s$ with $s=k\cdot(z-\text{hard position})$. The $p^{+}$ integral of a row term $\frac{dp^{+}}{p^{+}}f(p^{+})$ (Group~I) or $\frac{k^{+}dp^{+}}{p^{+}(p^{+}+k^{+})}f(p^{+})$ (Group~III) is split as [A]
\begin{align}
\int_{\Lambda}^{k^{+}}\frac{dp^{+}}{p^{+}}\,e^{-i\frac{p^{+}}{k^{+}}s}f(p^{+})=\log\frac{k^{+}}{\Lambda}\,f(0)+f(0)\,\mathcal{E}(s)+\int_{0}^{k^{+}}\frac{dp^{+}}{p^{+}}e^{-i\frac{p^{+}}{k^{+}}s}\big[f(p^{+})-f(0)\big]+{\cal O}\Big(\frac{\Lambda}{k^{+}}\Big)\,,\label{sf:split}
\end{align}
and, for Group~III, $\int_{\Lambda}\frac{k^{+}dp^{+}}{p^{+}(p^{+}+k^{+})}e^{-i\frac{p^{+}}{k^{+}}s}f(p^{+})$ in the same way with $\Psi(s)$ in place of $\mathcal{E}(s)$ (up to the $\log(1-\frac{k^{+}}{\vee})$ term, which goes with the JIMWLK structure). Here $f(0)$ is the row's coefficient in the JIMWLK list and the last term is the plus-prescription term with the phase kept in the subtraction (the ``regular term'', integrable at large $z$ by itself; checked in \texttt{Finite\_terms\_final}). The soft remainder of a row is $f(0)\Phi(s)$, $\Phi=\mathcal{E}$ (Group~I) or $\Psi$ (Group~III). Since $|\Phi|\simeq\log|s|$ for $|s|\gg1$, the recoil cuts the soft integral at $p^{+}\simeq k^{+}/|k\cdot(z-w)|$, i.e.\ at $\frac{p^{+}}{p_{\perp}}\simeq\frac{k^{+}}{|k|}$ for $p_{\perp}\sim|z-w|^{-1}$: the soft gluon is restricted to rapidities below that of the measured gluon. Each remainder behaves like $\frac{T}{z^{2}}\mathcal{E}_{\infty}(k\cdot z)$ at large $z$ ($T$: the $1/z^{2}$ tail of $f(0)$): not integrable alone.

The third term of $N_{3b}$ of Group~I Row III is of a different kind [A, \texttt{Finite\_terms\_final} Eq.~(IIIB)]: with $a_{t}=(x-w)-t(x-z)$,
\begin{align}
\ell(z,w,x)\equiv-\frac12\log\frac{(z-w)^{2}}{(x-w)^{2}}=-\int_{0}^{1}\frac{dt}{t}\Big[1-\frac{(x-w)\cdot a_{t}}{a_{t}^{2}}\Big]\,,\label{sf:ell}
\end{align}
the LO kernel minus the LO kernel with the source shifted by the recoil; the soft integral is cut at $t\simeq|x-w|/|z-w|$, no phase. At large $z$: $\ell\to-\log\frac{|z|}{|x-w|}$, isotropic, without the odd part of $\mathcal{E}_{\infty}$.

\subsubsection{Step 2: the soft remainder pieces (item (a))}
\noindent All pieces in units of $c_{0}$, under $\int_{w,w'}e^{-ik\cdot(w-w')}\int_{x,y,x',z}$, written half (recoil of Group~I on $w$, of Group~III on $w'$; in Group~III $y'$ is called $w'$). The Group~I Row V terms are relabelled from the row variables (soft $x$, hard $z$, phase $e^{-ik\cdot(z-w')}$; resp.\ soft $z$, hard $x$, phase $e^{-ik\cdot(x-w')}$) to soft $z$, hard $w$; in pair B the second term has also $y\leftrightarrow y'$ relabelled and both are then written with $y'\to x$; in pair D Row VI is written with $x\leftrightarrow y$. Colour structures of Group~I:
\begin{align}
C_{A}&\equiv f^{abd}\,\mathcal{L}^{a}\,U^{de}(x)\big[U^{bc}(w)-U^{bc}(y)\big]\rho^{e}(x)\rho^{c}(y)\,,\qquad
C_{D}\equiv f^{abd}\,\mathcal{L}^{a}\,U^{de}(x)\big[U^{bc}(z)-U^{bc}(y)\big]\rho^{e}(x)\rho^{c}(y)\,,\notag\\
C_{B}&\equiv \mathcal{L}^{a}\Big[f^{adc}\,U^{db}(z)\,U^{ce}(w)\,\{\rho^{e}(y),\rho^{b}(x)\}-f^{abc}\,U^{cd}(y)\,U^{be}(x)\,\{\rho^{d}(y),\rho^{e}(x)\}\Big]\,,\label{sf:col}\\
C_{C}&\equiv\big[U^{ab'}(x')-U^{ab'}(w')\big]\Big[N_{c}\,U^{ah}(x)-f^{adc}f^{hbe}\,U^{db}(z)\,U^{ce}(w)\Big]\rho^{b'}(x')\,\rho^{h}(x)\,.\notag
\end{align}
{\sloppy Group~III colour structures as in the author's rows, with $y'\to w'$: $\mathcal{L}^{c'd'b}(w')\equiv U^{c'e'}(w')U^{d'b}(z)\rho^{e'}(x')-U^{d'b}(z)U^{c'e'}(x')\rho^{e'}(x')$, $\mathcal{R}_{1}^{ab}\equiv U^{ac}(w)\{\rho^{c}(y),\rho^{b}(x)\}-U^{cb}(z)U^{ad}(y)U^{ce}(x)\{\rho^{d}(y),\rho^{e}(x)\}$, $\mathcal{R}_{2}^{ab}\equiv f^{cbd}U^{ac}(w)\rho^{d}(x)-f^{acd}U^{cb}(z)U^{de}(x)\rho^{e}(x)$, $\mathcal{L}_{\rm V}^{c'}\equiv\Delta^{c'e'}\rho^{e'}(x')$, $\mathcal{R}_{\rm V}^{ac}\equiv U^{ce}(x)U^{ad}(y)\rho^{e}(x)\rho^{d}(y)-U^{ad}(w)U^{ce}(x)\rho^{e}(x)\rho^{d}(y)$, $\mathcal{R}_{\rm VI}^{ab}\equiv U^{cb}(z)U^{ad}(y)U^{ce}(x)\rho^{d}(y)\rho^{e}(x)-U^{ad}(y)\rho^{d}(y)\rho^{b}(x)$.\par}

\noindent Each piece is (coefficient)$\times$(recoil function)$\times$(kernel)$\times$(colour); first as it comes from the row, then (``$=$'') in the common variables; $\mathcal{E}\equiv\mathcal{E}(k\cdot(z-w))$, $\Psi\equiv\Psi(k\cdot(z-w'))$.
\begin{align}
\text{A}^{\rm I}\ &\text{(I-III, last term of $N_{3b}$, recoil form):}\quad -2i\,\mathcal{E}\,K(x,w,z)\,P\!\cdot\!K(y-w)\,C_{A}\,,\notag\\
\text{D}^{\rm I}\ &\text{(I-III, third term of $N_{3b}$):}\quad +2i\,\ell(z,w,x)\,K(y,w,z)\,P\!\cdot\!K(x-w)\,C_{D}\,,\notag\\
\text{B}^{\rm I}\ &\text{(I-V first part, the $[\tfrac{k^{+}}{p^{+}}]_{+}$ and $[\tfrac{1}{q^{+}}]_{+}$ remainders; row variables, $s=k\cdot(x-z)$, phase $e^{-ik\cdot(z-w')}$):}\notag\\
&\quad-\frac{i}{2}\Big[\mathcal{E}(s)\,P\!\cdot\!K(y-z)\,K(z-x)\!\cdot\!K(y'-x)+e^{-is}\mathcal{E}(-s)\,P\!\cdot\!K(y'-x)\,K(z-x)\!\cdot\!K(y-z)\Big]\,[\text{V}]\notag\\
&\hspace{20em}=\;-i\,\mathcal{E}\,K(x,w,z)\,P\!\cdot\!K(y-w)\,C_{B}\,,\notag\\
\text{C}^{\rm I}\ &\text{(I-V, $N_{2a}+N_{2b\,\rm rest}$, the $[\tfrac{k^{+}}{p^{+}}]_{+}$ and $[\tfrac{k^{+}}{q^{+}}]_{+}$ remainders; row variables):}\notag\\
&\quad-\Big[\mathcal{E}(s)\,\frac{K^{i}K^{m}B_{2}^{im}}{(y-z)^{2}}+e^{-is}\mathcal{E}(-s)\,\frac{K^{i}K^{m}B_{1}^{im}}{(y-x)^{2}}\Big]\,[\dots]_{2a,2b}
\;=\;-\mathcal{E}\,P\!\cdot\!K(x-w)\big[K(x,w,z)-2K(w,w,z)\big]\,C_{C}\,,\notag\\
\text{A}^{\rm III}\ &\text{(III-V, $[\tfrac{k^{+}}{p^{+}(p^{+}+k^{+})}]_{+}$ remainder):}\quad-2i\,\Psi\,K(w',x,z)\,P\!\cdot\!K(y-w)\,f^{c'ca}\mathcal{L}_{\rm V}^{c'}\mathcal{R}_{\rm V}^{ac}\;=\;+2i\,\Psi\,K(x,w',z)\,P\!\cdot\!K(y-w)\,C_{A}\,,\label{sf:pieces}\\
\text{B}^{\rm III}\ &\text{(III-III first part, same):}\quad-i\,\Psi\,K(w',x,z)\,P\!\cdot\!K(y-w)\,f^{c'd'a}\mathcal{L}^{c'd'b}\mathcal{R}_{1}^{ab}\;=\;+i\,\Psi\,K(x,w',z)\,P\!\cdot\!K(y-w)\,C_{B}\,,\notag\\
\text{C}^{\rm III}\ &\text{(III-III second part, same):}\quad-2\,\Psi\,P\!\cdot\!K(x-w)\Big[\frac12K(x,w',z)-K(w,w',z)\Big]f^{c'd'a}\mathcal{L}^{c'd'b}\mathcal{R}_{2}^{ab}\notag\\
&\hspace{20em}=\;+\Psi\,P\!\cdot\!K(x-w)\big[K(x,w',z)-2K(w,w',z)\big]\,C_{C}\,,\notag\\
\text{D}^{\rm III}\ &\text{(III-VI, same):}\quad-2i\,\Psi\,K(w',x,z)\,P\!\cdot\!K(y-w)\,f^{c'd'a}\mathcal{L}^{c'd'b}\mathcal{R}_{\rm VI}^{ab}\;\overset{x\leftrightarrow y}{=}\;-2i\,\Psi\,K(y,w',z)\,P\!\cdot\!K(x-w)\,C_{D}\,.\notag
\end{align}
Here $[\text{V}]=\mathcal{L}^{a}[f^{adc}U^{db}(x)U^{ce}(z)\{\rho^{e}(y),\rho^{b}(y')\}-f^{abc}U^{cd}(y)U^{be}(y')\{\rho^{d}(y),\rho^{e}(y')\}]$; $K^{i}K^{m}\equiv K^{i}(x'-w')K^{m}(z-x)$; $B_{1}$, $B_{2}$ as in Eq.~\eqref{sf:V2} below; $[\dots]_{2a,2b}$: the colour factors of $N_{2a}$ and of $N_{2b\,\rm rest}$ (the latter with $N_{c}$), $[U(x')-U(w')]^{ab'}[-f^{adc}f^{fbe}U^{db}(x)U^{ce}(z)+N_{c}U^{af}(x)]\rho^{b'}(x')\rho^{f}(y)$ and $N_{c}[U(x')-U(w')]^{ab'}[U^{af}(y)-U^{af}(x)]\rho^{b'}(x')\rho^{f}(y)$. The last column uses $K(z-w)\cdot K(y-z)=-K(w,y,z)$, $K(z-w)\cdot K(w'-z)=-K(w,w',z)$, the identity $f^{adc}f^{hbe}U^{db}(w)U^{ce}(z)=f^{adc}f^{hbe}U^{db}(z)U^{ce}(w)$ (relabel $d\leftrightarrow c$, $b\leftrightarrow e$), the sum of the $N_{c}$ terms of $N_{2a}$ and $N_{2b\,\rm rest}$ ($N_{c}[U(z)+U(w)]+N_{c}[2U(x)-U(z)-U(w)]=2N_{c}U(x)$), and the colour identities of Step 3 for the Group~III entries. [N: every entry of the last column equals the row's own $f(0)$ at arbitrary $z$, two-$\rho$ (Weyl-reduced) and three-$\rho$ sectors, $SU(2)$ and $SU(3)$: differences $\le2\times10^{-15}$ for entries of size $0.2$--$8$ (\texttt{t10\_forms.py}).]

\subsubsection{Step 3: the four colour identities and the pair structure (item (b), analytic)}
\noindent With $U^{d'b}(z)U^{cb}(z)=\delta^{d'c}$ (orthogonality of the adjoint Wilson line), $f^{c'd'a}f^{dd'a}=N_{c}\delta^{c'd}$ and relabelling of colour indices only (no reordering of $\rho$'s), [A]
\begin{align}
&f^{c'ca}\,\mathcal{L}_{\rm V}^{c'}\,\mathcal{R}_{\rm V}^{ac}=-C_{A}\,,\qquad f^{c'd'a}\,\mathcal{L}^{c'd'b}(w')\,\mathcal{R}_{1}^{ab}=-C_{B}\,,\notag\\ &f^{c'd'a}\,\mathcal{L}^{c'd'b}(w')\,\mathcal{R}_{2}^{ab}=-C_{C}\,,\qquad f^{c'd'a}\,\mathcal{L}^{c'd'b}(w')\,\mathcal{R}_{\rm VI}^{ab}\big|_{x\leftrightarrow y}=+C_{D}\,.\label{sf:colid}
\end{align}
{\sloppy Example (pair D): the $U(z)$ part of $\mathcal{R}_{\rm VI}|_{x\leftrightarrow y}$ gives $-f^{c'd'a}\Delta^{c'e'}U^{d'b}(z)U^{ad}(x)\rho^{e'}(x')\rho^{d}(x)\rho^{b}(y)$, which becomes $-f^{abd}\Delta^{ab'}U^{bc}(z)U^{de}(x)\rho^{b'}(x')\rho^{e}(x)\rho^{c}(y)$ under $c'\to a$, $e'\to b'$, $d'\to b$, $a\to d$, $d\to e$, $b\to c$; with $\mathcal{L}^{a}=-\Delta^{ab'}\rho^{b'}(x')$ this is the $U(z)$ part of $C_{D}$. [N: Eq.~\eqref{sf:colid} to $\le2\times10^{-15}$ in the two-$\rho$ and three-$\rho$ sectors, $SU(2)$, $SU(3)$, random Wilson lines (\texttt{t9\_colour.py}).] Hence every pair has the form\par}
\begin{align}
\text{pair }X:\qquad \alpha_{X}\,V_{X}\,C_{X}\Big[\kappa_{X}(w)\,\mathcal{E}\big(k\cdot(z-w)\big)-\kappa_{X}(w')\,\Psi\big(k\cdot(z-w')\big)\Big]\,,\label{sf:pair}
\end{align}
\begin{center}
\renewcommand{\arraystretch}{1.3}
\begin{tabular}{|c|c|c|c|l|}\hline
pair & $\alpha_{X}$ & $V_{X}$ & $\kappa_{X}(u)$ & $T_{X}\equiv\lim_{|z|\to\infty} z^{2}\alpha_{X}V_{X}\kappa_{X}(w)C_{X}$\\\hline
A & $-2i$ & $P\cdot K(y-w)$ & $K(x,u,z)$ & $T_{A}=-2i\,P\!\cdot\!K(y-w)\,C_{A}$\\
B & $-i$ & $P\cdot K(y-w)$ & $K(x,u,z)$ & $T_{B}=-i\,P\!\cdot\!K(y-w)\,C_{B}$\\
C & $-1$ & $P\cdot K(x-w)$ & $K(x,u,z)-2K(w,u,z)$ & $T_{C}=+P\!\cdot\!K(x-w)\,C_{C}$\\
D & $+2i$ & $P\cdot K(x-w)$ & $K(y,u,z)$ & $T_{D}=+2i\,P\!\cdot\!K(x-w)\,C_{D}$\\\hline
\end{tabular}
\end{center}
The Group~III member is the Group~I member with the soft gluon attached to the measured gluon at $w'$ instead of $w$ (real instead of virtual). The tails of the two members are $\pm T_{X}/z^{2}$ exactly, as operators, for every value of $U(z)$: the parts linear in $U(z)$, the $z$-independent parts and the parts with $U(z)U^{T}(z)$ cancel together. In particular the third term of $N_{3b}$ (pair D) has the tail $\frac{T_{D}}{z^{2}}\ell$ with $T_{D}=2iP\cdot K(x-w)C_{D}$: its $U(z)$-independent part is $-2iP\!\cdot\!K(x-w)f^{abd}\mathcal{L}^{a}U^{de}(x)U^{bc}(y)\rho^{e}(x)\rho^{c}(y)$, its $U(z)$-linear part $+2iP\!\cdot\!K(x-w)f^{abd}\mathcal{L}^{a}U^{de}(x)U^{bc}(z)\rho^{e}(x)\rho^{c}(y)$, and both are cancelled by Group~III Row VI.

\noindent\textbf{Numbers} [N, \texttt{t4\_pairs.py}, \texttt{t4b.py}, \texttt{t7\_Uvar.py}]: maximum over 7 directions of $z^{2}\,|f^{\rm I}_{X}+f^{\rm III}_{X}|$ (coefficients of $\mathcal{E}$, $\Psi$; max norm of the Weyl two-$\rho$ matrix at two charge sites, the three-$\rho$ part with random classical $\rho$, and the singlet projection $\mathrm{tr}$), at $|z|=10^{2},10^{3},10^{4}$:
\begin{center}
\renewcommand{\arraystretch}{1.25}\small
\begin{tabular}{|c|c|c|c|c|}\hline
pair & $SU(2)$, two-$\rho$ & $SU(2)$, three-$\rho$ & $SU(3)$, two-$\rho$ & single members\\
 & \multicolumn{3}{c|}{(in units $10^{-2}$; $|z|=10^{2},\,10^{3},\,10^{4}$)} & (max; $SU(2)$, $SU(3)$)\\\hline
A & $5.7,\ 0.56,\ 0.056$ & $1.6,\ 0.16,\ 0.016$ & $4.1,\ 0.41,\ 0.041$ & $2.9$, $10.0$\\
B & $4.0,\ 0.40,\ 0.039$ & $4.2,\ 0.43,\ 0.043$ & $4.1,\ 0.41,\ 0.041$ & $2.2$, $15.5$\\
C & $4.3,\ 0.42,\ 0.041$ & (none) & $4.0,\ 0.39,\ 0.039$ & $2.1$, $1.6$\\
D & $5.5,\ 0.53,\ 0.053$ & $1.8,\ 0.18,\ 0.018$ & $7.0,\ 0.70,\ 0.070$ & $2.7$, $2.8$\\\hline
\end{tabular}
\end{center}
The decrease is $1/|z|$, i.e.\ $f^{\rm I}_{X}+f^{\rm III}_{X}={\cal O}(|z|^{-3})$ pointwise in the direction; the singlet trace and the c.c.\ halves behave the same way (e.g.\ pair A, $SU(2)$: trace $4.2\cdot10^{-2},\,4.3\cdot10^{-3},\,4.3\cdot10^{-4}$); the same holds with $U(z)=1$ (finite target) and with a different random $U(z)$ in every direction (pair D, $SU(3)$: $9.5\cdot10^{-2},\,1.0\cdot10^{-2},\,8.8\cdot10^{-4}$). The angular mean of $\kappa_{X}(w)-\kappa_{X}(w')$ times $C_{X}V_{X}$ falls like $|z|^{-4}$ ($|z|^{4}\times$ mean $=3.088,\,1.756,\,9.345,\,2.141$ for A--D, constant from $|z|=10^{2}$ to $10^{4}$; \texttt{t5\_angmean.py}). In the earlier bookkeeping of all nine pieces (both halves, \texttt{t1\_tails.log}) the total tail is $5\times10^{-8}$ ($SU(2)$) and $2\times10^{-8}$ ($SU(3)$) against single tails $0.4$--$3.3$.

\subsubsection{Step 4: the two recoil sides combine without any condition on $U$}
\noindent Write each pair \eqref{sf:pair} as [A]
\begin{align}
\kappa(w)\,\mathcal{E}\big(k\cdot(z-w)\big)-\kappa(w')\,\Psi\big(k\cdot(z-w')\big)=\underbrace{\big[\kappa(w)-\kappa(w')\big]\,\mathcal{E}\big(k\cdot(z-w)\big)}_{\text{real minus virtual}\,\times\,\text{kinematic log}}+\underbrace{\kappa(w')\Big[\mathcal{E}\big(k\cdot(z-w)\big)-\Psi\big(k\cdot(z-w')\big)\Big]}_{\text{mismatch of the recoil on $w$ and on $w'$}}\,.\label{sf:rmv}
\end{align}
First term: $\kappa(w)-\kappa(w')={\cal O}(|w-w'|/|z|^{3})$ and $|\mathcal{E}|\le\log|z|+{\rm const}$: integrable. Second term: by \eqref{sf:fun} $\mathcal{E}(k\cdot(z-w))-\Psi(k\cdot(z-w'))=\log\frac{|k\cdot(z-w')|}{|k\cdot(z-w)|}+{\cal O}(1/|k\cdot z|)$: the odd parts $-\frac{i\pi}{2}{\rm sgn}$ are the same on the two sides and cancel, the remainder is ${\cal O}(1/|z|)$ outside the strip $|k\cdot z|\lesssim|k||w-w'|$, and bounded by logarithms inside it (the strip contributes $\int dz_{\perp}/z_{\perp}^{2}$, convergent) [N: $|z|\times$median over directions of $|\mathcal{E}-\Psi|=2.49,\,2.53,\,2.51,\,2.46$ and $|z|\times$mean $=6.2,\,8.5,\,10.3,\,9.2$ at $|z|=10^{2},\dots,10^{5}$; \texttt{t7\_Uvar.py}]. The colour structure $C_{X}$, with all its dependence on $U(z)$, is a common factor of both terms. Therefore:
\begin{itemize}
\item[(i)] The $U(z)$-linear tails do cancel between the $w$ and the $w'$ recoil sides, but the two sides are Group~I written (function $\mathcal{E}$) and Group~III written (function $\Psi$), whose large-distance forms are equal, including the odd part. Nothing has to be assumed about $U(z)$ at large distance, beyond $|U^{ab}|\le1$. (\textbf{Correction} of \texttt{Finite\_terms\_final}, OPEN I-2(b) and OPEN II/III-3: there the $w$-side terms of Group~I were paired with the c.c.\ terms of Group~III, which also have their recoil on $w$ but carry $\Psi^{*}$; with that pairing a $-i\pi\,{\rm sgn}(k\cdot z)\,T^{(1)}/z^{2}$ seemed to be left. With the pairing \eqref{sf:pair} it is absent; the $z$-independent tails, which are self-conjugate in Group~III, cancel with either pairing, which is why the earlier tables were consistent.)
\item[(ii)] For a finite target, $U(z)\to1$ at large $|z|$: everything cancels in the same way; it is a special case of (i), nothing new is needed.
\item[(iii)] The c.c.\ half is the complex conjugate of each pair, with $\mathcal{E}^{*}$, $\Psi^{*}$ (recoil of Group~I on $w'$, of Group~III on $w$ after the relabelling $w\leftrightarrow w'$); it cancels pair by pair in the same way [N: table above, c.c.\ rows].
\end{itemize}

\subsubsection{Step 5: the problematic term of Group I Row III and what cancels it}
\noindent Split the third term of $N_{3b}$ as $\ell=\mathcal{E}+(\ell-\mathcal{E})$. The part $2i\,K(y,w,z)P\!\cdot\!K(x-w)C_{D}\,\mathcal{E}(k\cdot(z-w))$ is the Group~I member of pair D (cancelled by Group~III Row VI, Steps 3, 4). The rest is put together with the 24 other terms of the row (first line of \texttt{Finite\_terms\_final} Eq.~(III), combined kernel $\mathcal{K}^{ij}$, Eq.~\eqref{sf:Kij}):
\begin{align}
\widehat{\mathcal{K}}^{i}(x,y,z,w)\equiv K^{j}(y-z)\,\mathcal{K}^{ij}+\Big[\ell(z,w,x)-\mathcal{E}\big(k\cdot(z-w)\big)\Big]\,K(y,w,z)\,K^{i}(x-w)\,.\label{sf:Khat}
\end{align}
\textbf{Claim} [A+N]: at fixed $x,y,x',w'$ and $|z|\to\infty$,
\begin{align}
\int d^{2}w\;e^{-ik\cdot w}\,\widehat{\mathcal{K}}^{i}(x,y,z,w)={\cal O}\Big(\frac{\log|z|}{|z|^{3}}\Big)+\big(\text{terms}\propto e^{-ik\cdot z}\big)\,,\label{sf:claim}
\end{align}
i.e.\ no non-oscillating $1/z^{2}$ term in any direction, for any $U(z)$ (the colour factor $\mathcal{L}^{a}C_{\rm III}(z)$ does not depend on $w$ and is a spectator). Derivation:
\begin{enumerate}
\item The 24-term kernel has the tail $K^{j}(y-z)\mathcal{K}^{ij}\to\frac{\varphi}{z^{2}}\frac{e_{\perp}^{i}}{|x-w|}$ [\texttt{Finite\_terms\_final} Eq.~(saw), N], $\varphi\in(-\pi,\pi)$ the angle from $x-w$ to $z$, $e_{\perp}$ the unit vector $\frac{x-w}{|x-w|}$ rotated by $+\frac{\pi}{2}$ (from the $\mathcal{A}_{0}$ terms). Its angular mean is zero, but it is not zero pointwise.
\item By \eqref{sf:fun} and $K(y,w,z)\to1/z^{2}$: $\big[\ell-\mathcal{E}\big]K(y,w,z)K^{i}(x-w)\to\frac{K^{i}(x-w)}{z^{2}}\Big[\gamma_{E}+\log\big(|k||x-w|\big)+c(\theta)\Big]$ plus terms $\propto e^{-ik\cdot(z-w)}$, with $c(\theta)=\log|\cos\theta|+\frac{i\pi}{2}{\rm sgn}(\cos\theta)$, $\theta=\angle(k,z)$.
\item With $r=x-w=|r|(\cos\alpha,\sin\alpha)$ ($k$ along the first axis), $e^{-ik\cdot w}=e^{-ik\cdot x}e^{ik\cdot r}$, and the radial integrals $\int_{0}^{\infty}dr\,e^{i\kappa r}=i\,{\rm P}\frac1\kappa+\pi\delta(\kappa)$, $\int_{0}^{\infty}dr\,e^{i\kappa r}[\gamma_{E}+\log(|k|r)]=\frac{i}{\kappa}\big[-\log\frac{|\kappa|}{|k|}+\frac{i\pi}{2}{\rm sgn}\kappa\big]$ ($\kappa=|k|\cos\alpha$; with $e^{-\epsilon r}$, $\epsilon\to0$), the angular integrals give
\begin{align}
\int d^{2}r\,e^{ik\cdot r}\,K^{i}(r)\Big[\gamma_{E}+\log(|k||r|)+c(\theta)\Big]=+\frac{2\pi i\,\hat k^{i}}{|k|}\,\chi(\theta)\,,\qquad
\int d^{2}r\,e^{ik\cdot r}\,\varphi\,\frac{e_{\perp}^{i}}{|r|}=-\frac{2\pi i\,\hat k^{i}}{|k|}\,\chi(\theta)\,,\notag\\ \chi(\theta)\equiv\log|2\cos\theta|+\frac{i\pi}{2}{\rm sgn}(\cos\theta)\,.\label{sf:id}
\end{align}
For the first: $\int_{0}^{2\pi}d\alpha\,[-\log|\cos\alpha|]=2\pi\log2$, the $\sin\alpha$ component is odd. For the second: $\varphi=\theta-\alpha$ (mod $2\pi$) has a jump $2\pi$ at $\alpha=\theta+\pi$; ${\rm P}\!\int d\alpha\,\varphi\tan\alpha=2\pi\log|2\cos\theta|$ by parts, and the $\pi\delta(\kappa)$ terms at $\alpha=\pm\frac{\pi}{2}$ give $\frac{\pi^{2}}{|k|}{\rm sgn}(\cos\theta)$. The two cancel. $\chi$ has zero angular mean, which is why each of the two terms alone is only conditionally convergent (angular integral first), and not at all if $U(z)$ depends on the direction.
\end{enumerate}
[N, \texttt{t2\_identity.py}, \texttt{t2b.py}: Eq.~\eqref{sf:id} with $e^{-\epsilon r}$, six directions, $|k|=1.3$: each term up to $8.2$ in modulus, maximum of $|{\rm sum}|$ (both components) $7.0\times10^{-2},\,7.1\times10^{-3},\,7.1\times10^{-4}$ for $\epsilon=10^{-2},10^{-3},10^{-4}$; each term differs from $\pm\frac{2\pi i}{|k|}\chi$ by at most $9.7\times10^{-4}$ at $\epsilon=10^{-4}$ (error ${\cal O}(\epsilon)$). \texttt{t3c\_wint.py}: the full kernels, $\int d^{2}w\,e^{-ik\cdot w}e^{-(|w-x|/r_{c})^{2}}$ with $r_{c}=|z|/5$ (the cut removes only the region $w\sim z$, which gives terms $\propto e^{-ik\cdot z}$), $\mathcal{E}\to\mathcal{E}_{\infty}$ in \eqref{sf:Khat} (the difference is $\propto e^{-ik\cdot(z-w)}$), $|k|=1$, $P=(0.7,-0.4)$: $z^{2}\times$[24-term part] $=6.933-2.690i$, $z^{2}\times$[$(\ell-\mathcal{E}_{\infty})$ part] $=-6.906+2.690i$ at $|z|=2560$, $\theta=0.4$ (prediction $\mp(6.909-2.687i)$); $z^{2}\times|$sum$|=0.72,\,0.26,\,0.087,\,0.027$ at $|z|=40,160,640,2560$ ($\theta=0.4$), $0.31,\,0.020,\,0.0037,\,0.0015$ ($\theta=2.2$), $0.37,\,0.027,\,0.0011,\,0.0017$ ($\theta=-1.0$), $0.70,\,0.26,\,0.086,\,0.027$ ($\theta=-2.8$).]

\noindent\textbf{Answer to ``it should cancel with something''.} The third term of $N_{3b}$, $+2i\,\ell\,K(y,w,z)P\!\cdot\!K(x-w)C_{D}$ (units $c_{0}$), with tail $-\frac{T_{D}}{z^{2}}\log\frac{|z|}{|x-w|}$, is cancelled (1) in its $\log|z|$ part, for the $U(z)$-independent and the $U(z)$-linear colour parts alike, by the soft remainder of Group~III Row VI, $-2i\,\Psi\,K(y,w',z)P\!\cdot\!K(x-w)C_{D}$ (same coefficient, kernel with $w\to w'$); (2) in its remaining $1/z^{2}$ part, $\frac{T_{D}}{z^{2}}[\gamma_{E}+\log(|k||x-w|)+c(\theta)]$, by the $\mathcal{A}_{0}$ (``saw'') tail of the other 24 terms of Row III, after the $w$ integration, Eq.~\eqref{sf:id}. The last term of $N_{3b}$ (recoil form) is cancelled by Group~III Row V, with coefficient $-2i$ on both sides and kernels $K(x,w,z)$, $K(x,w',z)$. At $z\to w$ both terms are integrable ($\mathcal{E}\to0$; $\ell K(y,w,z)\sim\log|z-w|/|z-w|$).

\subsubsection{Step 6: the integrable form of the finite terms (complete formulas)}
\noindent Every term below is integrable over all transverse positions in the sense defined above; the only prescriptions needed are the order of integrations (C3) ($w$ before $z$ in Eq.~\eqref{sf:GIIII}) and (C7) ($y'$ before $p^{+}$ in the distributional terms of Group~III). c.c.\ implied throughout.

\noindent\textbf{Group I, Row III} (replaces (6.4)--(6.7); the two soft remainders are in the block $\mathcal{S}$, pairs A and D):
\begin{align}
\frac{d^{3}N^{\rm I,III}}{d^{3}k}&=\frac{1}{(2\pi)^{3}}\,\frac{ig^{4}f^{abd}}{8\pi^{5}}\,\frac{1}{k^{+}}\int_{w,w'}e^{-ik\cdot(w-w')}\int_{x,y,z,x'}\Bigg\{K^{i}(x'-w')\,K^{j}(y-z)\,\mathcal{K}^{ij}\notag\\
&\qquad+\Big[-\frac12\log\frac{(z-w)^{2}}{(x-w)^{2}}-\mathcal{E}\big(k\cdot(z-w)\big)\Big]\,K(y,w,z)\,K(x'-w')\cdot K(x-w)\Bigg\}\notag\\&\qquad\times\mathcal{L}^{a}\Big[U^{de}(x)U^{bc}(z)\rho^{e}(x)\rho^{c}(y)-U^{de}(x)U^{bc}(y)\rho^{e}(x)\rho^{c}(y)\Big]\,,\label{sf:GIIII}
\end{align}
with the $w$ integral done first, and the combined kernel (only positive arguments in the logarithms; no pole on $(x-z)^{2}=(x-w)^{2}$)
\begin{align}
\mathcal{K}^{ij}&=\log\frac{(x-z)^{2}}{(x-w)^{2}}\Bigg\{\frac{\delta^{ij}}{2\,(z-w)^{2}}-\frac{(x-w)^{i}\,(x-z)^{j}}{(x-w)^{2}\,(x-z)^{2}}
+\frac{1}{(z-w)^{2}\left[(x-w)^{2}-(x-z)^{2}\right]}\Big[(x-w)^{i}(z-w)^{j}\,\frac{(x-z)^{2}}{(x-w)^{2}}\notag\\
&\qquad\quad+(z-w)^{i}(x-z)^{j}\,\frac{(x-w)^{2}}{(x-z)^{2}}-(z-w)^{i}(x-w)^{j}-(x-z)^{i}(z-w)^{j}\Big]\Bigg\}\notag\\
&\quad+\frac{\mathcal{A}_{0}}{(z-w)^{2}}\Bigg\{(z-w)^{i}(x-w)^{j}+(x-z)^{i}(z-w)^{j}-(x-w)\cdot(x-z)\left[\frac{(z-w)^{i}(x-z)^{j}}{(x-z)^{2}}+\frac{(x-w)^{i}(z-w)^{j}}{(x-w)^{2}}\right]\Bigg\}\notag\\
&\quad-\frac{1}{2\,(z-w)^{2}}\log\frac{(z-w)^{2}}{(x-w)^{2}}\;\frac{(z-w)^{i}(x-z)^{j}}{(x-z)^{2}}\,,\qquad
\mathcal{A}_{0}=\frac{\arccos\Big(\frac{(x-w)\cdot(z-w)}{|x-w|\,|z-w|}\Big)}{\sqrt{(x-w)^{2}(x-z)^{2}-[(x-w)\cdot(x-z)]^{2}}}\,.\label{sf:Kij}
\end{align}
Integrability: $z\to w$: $\mathcal{K}^{ij}\sim\log|z-w|/|z-w|$ (angular mean of the $1/(z-w)^{2}$ coefficient zero), second line $\sim\log|z-w|/|z-w|$; $x\to w$: $\log|x-w|/|x-w|$; $z\to x$, $z\to y$: $1/|\cdot|$; large $z$: Eq.~\eqref{sf:claim}.

\noindent\textbf{Group I, Row V} (replaces (6.9)--(6.11); the soft remainders are in $\mathcal{S}$, pairs B and C). Row variables: $x$ = position of the gluon with $p^{+}$, $z$ = of the gluon with $k^{+}-p^{+}$, $s\equiv k\cdot(x-z)$, $\xi=p^{+}/k^{+}$:
\begin{align}
\frac{d^{3}N^{\rm I,V}_{1}}{d^{3}k}&=\frac{1}{(2\pi)^{3}}\frac{ig^{4}}{32\pi^{5}}\frac{1}{k^{+}}\int_{x,z,y,y',x',w'}e^{-ik\cdot(z-w')}\,\mathcal{J}(s)\,K(x'-w')\cdot K(z-x)\;K(y-z)\cdot K(y'-x)\notag\\
&\qquad\times\mathcal{L}^{a}\Big[f^{adc}U^{db}(x)U^{ce}(z)\{\rho^{e}(y),\rho^{b}(y')\}-f^{abc}U^{cd}(y)U^{be}(y')\{\rho^{d}(y),\rho^{e}(y')\}\Big]\,,\label{sf:V1}\\
\frac{d^{3}N^{\rm I,V}_{2}}{d^{3}k}&=\frac{1}{(2\pi)^{3}}\frac{g^{4}}{16\pi^{5}}\frac{1}{k^{+}}\int_{x,y,z,x',w'}e^{-ik\cdot(z-w')}\Bigg\{K^{i}(x'-w')K^{m}(z-x)\,\mathcal{F}^{im}_{\rm reg}\,\Big[-f^{adc}f^{fbe}U^{db}(x)U^{ce}(z)+N_{c}U^{af}(y)\Big]\notag\\
&\qquad+N_{c}\,\frac{11}{6}\,e^{-is}\,K(x'-w')\cdot K(y-x)\;\frac{K(y-z)\cdot(y-x)}{(x-z)^{2}}\,\Big[U^{af}(y)-U^{af}(x)\Big]\Bigg\}\Big[U^{ab'}(x')-U^{ab'}(w')\Big]\rho^{b'}(x')\rho^{f}(y)\,,\label{sf:V2}\\
\mathcal{F}^{im}_{\rm reg}&\equiv\int_{0}^{1}d\xi\,\frac{e^{-i\xi s}\,\mathcal{N}^{im}(\xi)}{\xi(y-x)^{2}+(1-\xi)(y-z)^{2}}\,,\notag\\
\mathcal{N}^{im}&\equiv\delta^{im}\xi(1-\xi)\frac{(y-x)^{2}-(y-z)^{2}}{(x-z)^{2}}+\delta^{im}\big[(1-\xi)T_{1}+\xi T_{2}\big]+\frac{(y-x)^{2}}{(y-z)^{2}}B_{2}^{im}+\frac{(y-z)^{2}}{(y-x)^{2}}B_{1}^{im}-Q^{im}\,,\notag
\end{align}
with $B_{1}^{im}=\frac{(y-x)^{i}(x-z)^{m}}{(x-z)^{2}}-\frac{(y-x)^{i}(y-z)^{m}}{2(y-z)^{2}}$, $B_{2}^{im}=\frac{(y-z)^{i}(x-z)^{m}}{(x-z)^{2}}+\frac{(y-x)^{m}(y-z)^{i}}{2(y-x)^{2}}$, $T_{1}=\frac{(y-z)\cdot(x-z)}{(x-z)^{2}}+\frac{(y-x)\cdot(y-z)}{2(y-x)^{2}}-\frac{(y-x)^{2}-(y-z)^{2}}{2(x-z)^{2}}$, $T_{2}=\frac{(y-x)\cdot(x-z)}{(x-z)^{2}}-\frac{(y-x)\cdot(y-z)}{2(y-z)^{2}}-\frac{(y-x)^{2}-(y-z)^{2}}{2(x-z)^{2}}$, $Q^{im}=\frac{(y-x)^{m}(x-z)^{i}}{(x-z)^{2}}-\frac{(y-x)^{m}(y-z)^{i}}{2(y-z)^{2}}+\frac{(y-z)^{m}(x-z)^{i}}{(x-z)^{2}}+\frac{(y-x)^{i}(y-z)^{m}}{2(y-x)^{2}}$ (the colour bracket of the first line of \eqref{sf:V2} is the sum of those of $N_{2a}$ and $N_{2b\,\rm rest}$; $N_{2b\,\rm UV}$ cancels against Row II). These are the regular terms of \texttt{Finite\_terms\_final} Eqs.~(V1), (F) without the recoil terms; integrable alone (no $1/R^{2}$ tail; the collinear point $z\to x$ of the brace of \eqref{sf:V2} with the angular integral first, as in the checked row).

\noindent\textbf{Group III, Rows III, V, VI} (replace (8.2)--(8.5); the soft remainders are in $\mathcal{S}$, pairs B, C, A, D). $s\equiv k\cdot(z-y')$, $\mathcal{L}_{2}^{ab}(y')\equiv f^{c'd'a}\mathcal{L}^{c'd'b}(y')$, $\mathcal{L}^{c'd'b}(z)$ is $\mathcal{L}^{c'd'b}(y')$ with $U^{c'e'}(y')\to U^{c'e'}(z)$:
\begin{align}
\frac{d^{3}N^{\rm III,III}_{1}}{d^{3}k}&=\frac{1}{(2\pi)^{3}}\frac{ig^{4}f^{c'd'a}}{16\pi^{5}}\frac{1}{k^{+}}\int_{w,x',y',z,y,x}e^{-ik\cdot(w-y')}\Big\{\widehat{\mathcal{E}}(s)\Big[K(y'-z)\cdot K(x'-y')\,K(y-w)\cdot K(x-z)\notag\\&\qquad-K(y',x,z)\,K(x'-y')\cdot K(y-w)\Big]\mathcal{L}^{c'd'b}(y')\notag\\
&\qquad-\Big[\pi\delta(s)-i\,{\rm P}\frac1s\Big]\,K(y'-z)\cdot K(y-w)\,K^{j}(x-z)\Big[K^{j}(x'-y')\,\mathcal{L}^{c'd'b}(y')-K^{j}(x'-z)\,\mathcal{L}^{c'd'b}(z)\Big]\Big\}\,\mathcal{R}_{1}^{ab}\,,\label{sf:G3R31}\\
\frac{d^{3}N^{\rm III,III}_{2}}{d^{3}k}&=\frac{1}{(2\pi)^{3}}\frac{g^{4}}{4\pi^{4}}\int_{0}^{\infty}\frac{dp^{+}}{2\pi}\int_{w,x',y',z,x}e^{-ik\cdot(w-y')}e^{-i\frac{p^{+}}{k^{+}}s}\,\frac{1}{D_{p}}\Bigg\{K(y'-z)\cdot K(x'-y')\bigg[\frac{p^{+}(B-A)}{C\,(p^{+}+k^{+})^{2}}\notag\\&\qquad+\frac{p^{+}}{k^{+}(p^{+}+k^{+})}\Big(\frac{b\cdot c}{C}-\frac{a\cdot b}{2A}-\frac{B-A}{2C}\Big)\notag\\
&\qquad+\frac{1}{p^{+}+k^{+}}\Big(\frac{a\cdot c}{C}+\frac{a\cdot b}{2B}-\frac{B-A}{2C}\Big)\bigg]+\frac{B-A}{(p^{+}+k^{+})A}\,\widehat{\mathcal{J}}-\frac{1}{k^{+}}\,\mathcal{V}\Bigg\}\,\mathcal{L}_{2}^{ab}(y')\,\mathcal{R}_{2}^{ab}\notag\\
&\quad-\frac{1}{(2\pi)^{3}}\frac{g^{4}}{4\pi^{4}}\frac{1}{k^{+\,2}}\int_{0}^{\infty}\frac{dp^{+}}{2\pi}\int_{w,x',y',z,x}e^{-ik\cdot(w-y')}e^{-i\frac{p^{+}}{k^{+}}s}\,\frac{p^{+}}{D_{p}}\,K^{m}(y'-z)\,\mathcal{W}^{k'm}\notag\\&\qquad\times\Big[K^{k'}(x'-y')\,\mathcal{L}_{2}^{ab}(y')-K^{k'}(x'-z)\,\mathcal{L}_{2}^{ab}(z)\Big]\mathcal{R}_{2}^{ab}\notag\\
&\quad+\frac{1}{(2\pi)^{3}}\frac{g^{4}}{8\pi^{5}}\frac{1}{k^{+}}\int_{w,x',w',z,x}e^{-ik\cdot(w-w')}\frac{(x-w)^{2}}{(x-z)^{2}\left[(x-w)^{2}-(x-z)^{2}\right]}\log\frac{(x-w)^{2}}{(x-z)^{2}}\notag\\&\qquad\times K^{m}(w'-z)K^{k'}(x'-z)\,\mathcal{W}^{k'm}\,\mathcal{L}_{2}^{ab}(z)\,\mathcal{R}_{2}^{ab}\,,\label{sf:G3R32}
\end{align}
with $a=x-w$, $b=x-z$, $c=z-w$, $A=a^{2}$, $B=b^{2}$, $C=c^{2}$, $D_{p}=p^{+}B+k^{+}A$, $\widehat{\mathcal{J}}=\frac{a\cdot K(x'-y')\;c\cdot K(y'-z)}{C}+\frac{b\cdot K(y'-z)\;a\cdot K(x'-y')}{2B}$, $\mathcal{V}=\frac{b\cdot K(y'-z)\,c\cdot K(x'-y')}{C}-\frac{b\cdot K(y'-z)\,a\cdot K(x'-y')}{2A}+\frac{a\cdot K(y'-z)\,c\cdot K(x'-y')}{C}+\frac{b\cdot K(x'-y')\,a\cdot K(y'-z)}{2B}$, $\mathcal{W}^{k'm}=\frac{b^{k'}c^{m}}{C}-\frac{b^{k'}a^{m}}{2A}$; in the second term the $y'$ integral is done before the $p^{+}$ integral (Step 9). (The first brace is $N_{2}^{\rm reg}$ of \texttt{Finite\_terms\_final}; its coefficient of $\widehat{\mathcal{J}}$ is the plus term with the phase kept, $-\big[\frac{k^{+}}{p^{+}(p^{+}+k^{+})}(1-\frac{D_{p}}{k^{+}A})+\frac{1}{p^{+}+k^{+}}\big]=\frac{B-A}{(p^{+}+k^{+})A}$: no $1/p^{+}$ is left; the brace stays under one $p^{+}$ and one $z$ integral.)
\begin{align}
\frac{d^{3}N^{\rm III,V}}{d^{3}k}&=\frac{1}{(2\pi)^{3}}\frac{ig^{4}f^{c'ca}}{8\pi^{5}}\frac{1}{k^{+}}\int_{w,y',z,y,x,x'}e^{-ik\cdot(w-y')}\Big\{\widehat{\mathcal{E}}(s)\Big[K(y'-z)\cdot K(x'-y')\,K(y-w)\cdot K(x-z)\notag\\&\qquad-K(y',x,z)\,K(x'-y')\cdot K(y-w)\Big]\mathcal{L}_{\rm V}^{c'}(y')\notag\\
&\qquad-\Big[\pi\delta(s)-i\,{\rm P}\frac1s\Big]K(y'-z)\cdot K(y-w)\,K^{j}(x-z)\Big[K^{j}(x'-y')\mathcal{L}_{\rm V}^{c'}(y')-K^{j}(x'-z)\mathcal{L}_{\rm V}^{c'}(z)\Big]\Big\}\,\mathcal{R}_{\rm V}^{ac}\,,\label{sf:G3R5}
\end{align}
with $\mathcal{L}_{\rm V}^{c'}(y')=U^{c'e'}(y')\rho^{e'}(x')-U^{c'e'}(x')\rho^{e'}(x')$; Row VI is Eq.~\eqref{sf:G3R5} with $f^{c'ca}\mathcal{L}_{\rm V}^{c'}\mathcal{R}^{ac}_{\rm V}\to f^{c'd'a}\mathcal{L}^{c'd'b}\mathcal{R}^{ab}_{\rm VI}$ (same prefactor $\frac{ig^{4}}{8\pi^{5}k^{+}}$). Integrability: $\widehat{\mathcal{E}}(s)\simeq\frac{1}{is}$ at large $z$ ($|z|^{-3}$), logarithm on $k\cdot(z-y')=0$; the distributional terms are finite (Step 9), no tail without the phase $e^{-ik\cdot(w-z)}$; $N_{2}^{\rm reg}$: $\log R/R^{3}$ [checked in \texttt{Finite\_terms\_final}].

\noindent\textbf{The soft-recoil block $\mathcal{S}$} (all soft remainders of the five rows; pairs A--D, Eq.~\eqref{sf:pair}):
\begin{align}
\frac{d^{3}N^{\rm soft}}{d^{3}k}&=c_{0}\int_{w,w'}e^{-ik\cdot(w-w')}\int_{x,y,x',z}\Bigg\{-i\,K(x'-w')\cdot K(y-w)\,\big[2C_{A}+C_{B}\big]\;\mathfrak{F}\big[K(x,\cdot,z)\big]\notag\\
&\qquad-K(x'-w')\cdot K(x-w)\;C_{C}\;\mathfrak{F}\big[K(x,\cdot,z)-2K(w,\cdot,z)\big]+2i\,K(x'-w')\cdot K(x-w)\;C_{D}\;\mathfrak{F}\big[K(y,\cdot,z)\big]\Bigg\}\,,\label{sf:S}\\
\mathfrak{F}[\kappa]&\equiv\kappa(w)\,\mathcal{E}\big(k\cdot(z-w)\big)-\kappa(w')\,\Psi\big(k\cdot(z-w')\big)=\big[\kappa(w)-\kappa(w')\big]\,\mathcal{E}\big(k\cdot(z-w)\big)+\kappa(w')\Big[\mathcal{E}\big(k\cdot(z-w)\big)-\Psi\big(k\cdot(z-w')\big)\Big]\,,\notag
\end{align}
where $\kappa(\cdot)$ is evaluated at $u=w$ (Group~I member) and $u=w'$ (Group~III member); $C_{C}$ has $U(w)$ also in the second term ($K(w,u,z)$ is $1/(z-w)^{2}$ at $u=w$ and $K(w,w',z)$ at $u=w'$). Both terms of $\mathfrak{F}$ are integrable (Step 4); at the coincidences: $z\to w$: $\mathcal{E}={\cal O}(|z-w|)$ against at most $(z-w)^{-2}$; $z\to w'$: $\Psi\to0$, $\kappa(w')\sim|z-w'|^{-1}$; $z\to x,y$: $|\cdot|^{-1}$. Eqs.~\eqref{sf:GIIII}--\eqref{sf:S} together are exactly the finite terms of these five rows [A; the regular terms are those of \texttt{Finite\_terms\_final}, checked there pointwise against the rows with $\Lambda,\vee$ kept; the soft remainders are their recoil terms, rewritten with Eq.~\eqref{sf:colid}].

\subsubsection{Step 7: equivalent kinematically constrained form (item (c))}
\noindent The combination can be made integrable with what is in the finite terms (Steps 3--6), so the JIMWLK extraction need not be changed. Equivalently [A], adding the $\log\frac{k^{+}}{\Lambda}$ terms of the same pairs, $L\,\alpha V C[\kappa(w)-\kappa(w')]$ (their sum is the JIMWLK-list contribution of these rows: soft gluon attached to the measured gluon at $w$ minus at $w'$, a real-minus-virtual combination falling like $|z|^{-3}$, angular mean $|z|^{-4}$, Step 3),
\begin{align}
L\,[\kappa(w)-\kappa(w')]+\mathfrak{F}[\kappa]=\kappa(w)\int_{\Lambda}^{k^{+}}\frac{dp^{+}}{p^{+}}\,e^{-i\frac{p^{+}}{k^{+}}k\cdot(z-w)}-\kappa(w')\int_{\Lambda}^{\infty}\frac{k^{+}dp^{+}}{p^{+}(p^{+}+k^{+})}\,e^{-i\frac{p^{+}}{k^{+}}k\cdot(z-w')}+{\cal O}\Big(\frac{\Lambda}{k^{+}}\Big)\,:\label{sf:kc}
\end{align}
each kernel of the JIMWLK term with its own rapidity interval $\simeq\log\frac{k^{+}}{\Lambda|k\cdot(z-w)|}-\gamma_{E}-\frac{i\pi}{2}{\rm sgn}$ (resp.\ $w'$) for $|k\cdot(z-w)|\gg1$ and $\log\frac{k^{+}}{\Lambda}$ for $|k\cdot(z-w)|\ll1$: the soft gluon has a rapidity below that of the measured gluon, as in the kinematically constrained formulations (Ioffe-time restriction of Altinoluk et al.\ 1411.2869, whose remainder $[A_{\xi}(y-z)-A_{\xi}(\bar y-z)]^{2}[s(y,z)s(z,\bar y)-s(y,\bar y)]$ has the same real-minus-virtual structure as $[\kappa(w)-\kappa(w')]\mathcal{E}$; the constraint of Ducloué--Lappi--Zhu 1604.00225 was not compared in detail). Read this way, the virtual (Group~I) members of these pairs get a $z$-dependent interval as well, while the JIMWLK-list terms that carry no recoil phase keep $\log\frac{k^{+}}{\Lambda}$. In the constrained reading the whole block $\mathcal{S}$ is absorbed into the JIMWLK term (which stays integrable at large $z$, since \eqref{sf:kc} is the sum of the integrable $L[\kappa(w)-\kappa(w')]$ and $\mathfrak{F}$), and of the soft remainders only the $(\ell-\mathcal{E})$ term of Row III, Eq.~\eqref{sf:GIIII}, stays among the finite terms. Both readings give the same cross section; the first one keeps the JIMWLK term exactly as in the JIMWLK section.

\subsubsection{Step 8: principal value and $\log|\cdot|$ in Group I Row III (item (d))}
\noindent In $\mf{B}_{3}$ the denominators $p^{+}\mathbf{k}^{2}-k^{+}\mathbf{p}^{2}$ (momentum space) and $p^{+}(x-u)^{2}-k^{+}(x-z)^{2}$ (coordinate space) are differences of the light-cone energies $\frac{\mathbf{k}^{2}}{2k^{+}}$ and $\frac{\mathbf{p}^{2}}{2p^{+}}$ of two one-gluon states connected by the static valence charge $\rho(\mathbf{k}-\mathbf{p})$, which transfers transverse momentum but no light-cone energy; they vanish on the energy shell. The dressing operator is constructed in stationary perturbation theory, where such denominators are principal values (the prescription of the setup and of all checked rows; \texttt{Calculation\_main\_1\_checked}, header; \texttt{Group\_I\_checked}, note at l.82). With it, $\int_{p_{1}}^{p_{2}}dp^{+}\,{\rm P}\frac{1}{p^{+}B-k^{+}A}=\frac1B\log\Big|\frac{p_{2}B-k^{+}A}{p_{1}B-k^{+}A}\Big|$ ($A=(x-w)^{2}$, $B=(x-z)^{2}$), and in (6.4)--(6.6) every logarithm whose argument changes sign is a $\log|\cdot|$:
\begin{align}
\log\Big|\frac{(x-z)^{2}}{(x-z)^{2}-(x-w)^{2}}\Big|\ \ (N_{1},\,N_{2})\,,\qquad \log\Big|\frac{(x-w)^{2}-(x-z)^{2}}{(x-w)^{2}}\Big|\ \ (N_{1},\,N_{3a})\,;\notag
\end{align}
the other logarithms ($\log\frac{(x-z)^{2}}{(x-w)^{2}}$, $\log\frac{(z-w)^{2}}{(x-w)^{2}}$) have positive arguments, and $\mathcal{A}_{0}$ is real. In the combined kernel \eqref{sf:Kij} the PV logarithms have combined into $\log\frac{(x-z)^{2}}{(x-w)^{2}}$ and the pole on $(x-z)^{2}=(x-w)^{2}$ has cancelled [checked in \texttt{Finite\_terms\_final}: $0.19948358$ and $0.19948353$ at $\pm10^{-7}$ from the line], so in the integrable form the prescription is no longer visible. A $\pm i\epsilon$ prescription would add $\mp i\pi\,\theta\big((x-z)^{2}-(x-w)^{2}\big)$-type terms (on-shell gluon--valence scattering) to the separate equations; this was not examined.

\subsubsection{Step 9: instantaneous terms of Group III, $p^{+}$ to infinity (item (e))}
\noindent The $-\frac{1}{k^{+}}\delta\delta$ terms of Rows III, V, VI carry no $p^{+}$ dependence except the phase. With $s=k\cdot(z-y')$ and the $y'$ integral done first [A],
\begin{align}
\int_{0}^{\infty}dp^{+}\,e^{-i\frac{p^{+}}{k^{+}}s}=k^{+}\Big[\pi\,\delta(s)-i\,{\rm P}\frac1s\Big]\qquad(\text{c.c.: }\pi\delta(s)+i\,{\rm P}\tfrac1s)\,,\label{sf:distr}
\end{align}
which gives the second lines of \eqref{sf:G3R31}, \eqref{sf:G3R5} (prefactor $\frac{ig^{4}f}{8\pi^{4}k^{+}}\frac{1}{2\pi}\cdot\frac{-1}{k^{+}}\cdot k^{+}=-\frac{ig^{4}f}{16\pi^{5}k^{+}}$ for Row III). Finiteness: with $r=y'-z$, the bracket $K^{j}(x'-y')\mathcal{L}(y')-K^{j}(x'-z)\mathcal{L}(z)$ vanishes at $r\to0$, because $\mathcal{L}^{c'd'b}(y')\to U^{c'e'}(z)U^{d'b}(z)\rho^{e'}(x')-U^{d'b}(z)U^{c'e'}(x')\rho^{e'}(x')=\mathcal{L}^{c'd'b}(z)$ and $K(x'-y')\to K(x'-z)$: the $U(y')U(z)$ and the $U(z)U(z)$ pieces cancel at $y'=z$, so $K^{m}(r)[\dots]$ is bounded there. Integrating $r_{\perp}$ (component perpendicular to $k$) first, symmetrically at large distance, gives a function $h(r_{\parallel})$ which is ${\cal O}(r_{\parallel})$ for the component $m=\parallel$ and finite at $r_{\parallel}=0$ for $m=\perp$; then $\pi\delta(s)$ and ${\rm P}\frac1s$ act on a continuous function of $r_{\parallel}$, and the phase $e^{ik\cdot r}$ of $e^{-ik\cdot(w-y')}$ makes the $r_{\parallel}$ integral converge at large distance. The two terms of the bracket must stay together: each alone grows like $\pm2\pi i\frac{k^{m}}{k^{2}}\log\vee$ (the BFKL term). [N, \texttt{t8\_distr.py} (copy of the earlier \texttt{g3int.py}), smooth test function in place of $K(x'-y')\mathcal{L}(y')$: $\int_{0}^{T}dt\int d^{2}r\,e^{i(1+t)k\cdot r}K^{m}(r)[b(r)-b(0)]$ differs from its limit by $0.638,\,6.95\times10^{-2},\,7.01\times10^{-3},\,7.02\times10^{-4}$ for $T=10,\dots,10^{4}$; the distributional form gives the limit to $7.3\times10^{-8}$; without the subtraction the integral grows by $11.248i$ per decade of $T$ against $2\pi i\frac{n\cdot k}{k^{2}}\log10=11.253i$. Row III second part, kernel $\frac{p^{+}}{D_{p}}$: difference to the limit $0.33,\,3.65\times10^{-2},\,3.69\times10^{-3},\,3.69\times10^{-4}$.]

\subsubsection{Conditions and assumptions}
\begin{enumerate}
\item[(C1)] $\Lambda\to0$, $\vee\to\infty$ are taken after the subtractions \eqref{sf:split}; terms ${\cal O}(\Lambda/k^{+})$, ${\cal O}(k^{+}/\vee)$ are dropped; the $\log(1-\frac{k^{+}}{\vee})$ terms go with the JIMWLK structures of their rows (as in the checked files).
\item[(C2)] Pair identities \eqref{sf:colid} use only colour relabelling and $U(z)U^{T}(z)=1$; they are operator identities (no reordering of $\rho$), valid for any $U$ at any $z$. In pair B the two Group~I recoil terms are combined after relabelling of the integrated valence positions ($y\leftrightarrow y'$); in pair D Row VI is written with $x\leftrightarrow y$. The statements about the $z$ integrand hold at fixed values of all other coordinates after these relabellings.
\item[(C3)] Group~I Row III, Eq.~\eqref{sf:GIIII}: the $w$ integral (the Fourier transform at the measured momentum) is done before the $z$ integral. Pointwise in $w$ the brace has a $1/z^{2}$ tail ($\frac{\varphi\,e_{\perp}}{|x-w|}+K(x-w)[\gamma_{E}+\log(|k||x-w|)+c(\theta)]$); only after the $w$ integration it is absent, in every direction, Eq.~\eqref{sf:id}.
\item[(C4)] Terms $\propto e^{-ik\cdot z}/z^{2}$ (from the end point $t\to1$ of $\mathcal{E}$, and from the regions where the measured gluon sits at $z$) are treated as convergent Fourier integrals; ``no tail'' means no non-oscillating $1/z^{2}$ term.
\item[(C5)] No assumption on $U(z)$ at large distance; for a finite target ($U\to1$) the same cancellation holds. Unitarity of the adjoint Wilson lines is used.
\item[(C6)] The energy denominators of $\mf{B}_{3}$ are principal values (Step 8).
\item[(C7)] Group~III distributional terms: $y'$ integral before the $p^{+}$ integral, $r_{\perp}$ first, symmetric at large distance (Step 9).
\item[(C8)] The large-$z$ form of the 24-term kernel of Row III (``saw'') is numerical (earlier check and \texttt{t3c\_wint.py}); its Fourier identity \eqref{sf:id} is analytic.
\item[(C9)] Verified on a two-site projectile (all assignments of the charge positions to two sites, Weyl reduction including the contact terms of the three-$\rho$ products) and with random classical $\rho$ for the three-$\rho$ part; general $\langle\rho\rho\rangle$ (the full two-$\rho$ coefficient matrix) and the singlet projection (its trace).
\end{enumerate}

\subsubsection{Not done / open}
\noindent (1) No complete numerical integration over all transverse positions with Wilson lines: the checks are exact identities, pointwise tails in all directions, and Fourier integrals at fixed other coordinates. (2) The regular terms (\eqref{sf:V1}, the brace of \eqref{sf:V2}, $\widehat{\mathcal{E}}$ and $N_{2}^{\rm reg}$ terms) are taken from \texttt{Finite\_terms\_final}, where they were checked; they were not re-checked here. (3) The non-singlet $\log\frac{k^{+}}{\Lambda}\int\frac{d^{2}z}{z^{2}}C_{A}$ of the JIMWLK formula and the Group~II Row IV tail are outside this section and unchanged. (4) Step 8: the alternative $\pm i\epsilon$ prescription was not evaluated.


% ---- inlined copy of paper/notes/fix_poles_dglap.tex (20 Nov 2026 06:34 UTC) ----
\subsection{Divergences, DGLAP and the running coupling}
\label{sec:poles-dglap}

\noindent\textbf{Scope and sources.} This section settles problems P3 (poles, the virtual $\frac{11}{3}$ terms, what renormalises the coupling), P4 (DGLAP term in $\overline{\rm MS}$ form), the phase convention part of P5, and the large-distance tail of Group~II Row~IV (OPEN~II/III-1 of \texttt{Finite\_terms\_final}). It uses the checked rows (\texttt{Group\_I\_checked}, \texttt{Group\_II\_Row\_IV\_checked}, \texttt{Group\_II\_checked}, \texttt{Group\_III\_checked}, \texttt{Divergent\_terms\_checked}, \texttt{Calculation\_main\_1\_checked}). Scripts (all in \texttt{fix/poles/}): \texttt{s1\_pole\_algebra.py} (exact pole algebra), \texttt{s2\_uv\_ir\_regions.py} (region of each divergence, cut-off integrals), \texttt{s3\_dglap\_ddim.py} (the collinear integral in $d=2-\epsilon$, numerically), \texttt{s4\_tail.py} (soft-parent tail, from the author's momentum-space amplitudes), \texttt{s5\_phases.py} (inventory of phase conventions), \texttt{s6\_misc.py} (Fourier identity, tail coefficients); outputs \texttt{s*.out}.

\medskip
\noindent\textbf{Notation.} $K(r)\equiv r/r^{2}$; $\mathcal{B}(r)\equiv2\gamma_{E}+\log\frac{r^{2}\mu^{2}}{4}=\log\frac{r^{2}\mu^{2}}{c_{0}^{2}}$, $c_{0}=2e^{-\gamma_{E}}$, $\mu\equiv\mu_{\overline{\rm MS}}$; $L\equiv\log\frac{k^{+}}{\Lambda}$, $\ell\equiv\log(1-\frac{k^{+}}{\vee})$; $\beta_{0}=\frac{11N_{c}}{3}$; $\mathfrak{a}\equiv\frac{\alpha_{s}N_{c}}{\pi}$, so that $\frac{\alpha_{s}\beta_{0}}{4\pi}=\frac{11}{12}\mathfrak{a}$. The LO structure and its Fourier transform are
\begin{align}
\mathcal{O}(y,y')&\equiv K(x-y)\cdot K(x'-y')\,\big[U(y')-U(x')\big]^{ce'}\big[U(y)-U(x)\big]^{ce}\rho^{e'}(x')\rho^{e}(x),\notag\\
\widehat{L}(q)&\equiv\int_{x,y,x',y'}e^{-iq\cdot(y-y')}\,\mathcal{O}(y,y'),\notag\\
\frac{d^{3}N_{\rm LO}}{d^{3}k}&=\mathcal{N}_{0}\,\widehat{L}(k),\qquad \mathcal{N}_{0}\equiv\frac{1}{(2\pi)^{3}}\frac{g^{2}}{2\pi^{2}}\frac{1}{k^{+}},\qquad
\mathcal{C}\equiv\mathcal{N}_{0}\,\frac{\alpha_{s}}{2\pi}=\frac{1}{(2\pi)^{3}}\frac{g^{2}}{2\pi^{2}}\frac{g^{2}}{8\pi^{2}}\frac{1}{k^{+}},
\label{eq:pd-not}
\end{align}
and, for the DGLAP term, $\Pi(\xi)\equiv\xi^{-2}\widehat{L}(k/\xi)$ (the LO structure of a parent gluon with momentum $k/\xi$), $\mathcal{X}\equiv\int_{x,x',z}\mathcal{O}(z,z)$. Phase: $e^{-ik\cdot(\text{unprimed}-\text{primed})}$ throughout (Sec.~\ref{sec:pd-phase}).

\medskip
\noindent\textbf{Regularisation.} Transverse dimension $d=2-\epsilon$ for the loop and collinear integrals. A scaleless integral is written as $\int d^{d}p/p^{2}=\pi\big(\frac{2}{\epsilon_{\rm UV}}-\frac{2}{\epsilon_{\rm IR}}\big)$ (\texttt{s2} [4]); it is zero in the author's single-$\epsilon$ bookkeeping, where $\epsilon_{\rm UV}=\epsilon_{\rm IR}$. Both bookkeepings are given below; they lead to the same finite result.

\subsubsection{Pole bookkeeping (P3)}

\noindent\textbf{Step 1: the $\frac{11}{3}$ part of the one-loop wave function, Group~I Row~II.} The one-$\rho$ part of $\mf{A}^{(3)}$ (Eq.~(gtwoloop1) of the main file, (I.89), LM Eq.~(3.97)) carries $\big(\frac{11}{3}-2L+\ell\big)\big[-\frac{2}{\epsilon}+\log\frac{k^{2}}{\mu^{2}}\big]$ times $\mf{A}^{(1)}$ (the $-2L$ instead of LM's $4\log\frac{\Lambda}{k^{+}}$ comes from the commutator parts of the two-$\rho$ states). In coordinate space $\log\frac{k^{2}}{\mu^{2}}\,\frac{k^{i}}{k^{2}}\to-\mathcal{B}(x-w)\,\frac{(x-w)^{i}}{(x-w)^{2}}$, since $\int_{0}^{\infty}dk\,J_{1}(kr)\log\frac{k^{2}}{\mu^{2}}=-\frac{1}{r}\mathcal{B}(r)$ (\texttt{s6} [1]: $-0.456741$ against $-0.456734$). Its $\frac{11}{3}$ part in the cross section (Group~I Row~II, $N_{1}$, with $[U(x')-U(w')][U(w)-U(x)]=-[\dots]$ of $\mathcal{O}$) is
\begin{align}
\frac{d^{3}N^{\rm II}_{11/3}}{d^{3}k}&=\frac{1}{(2\pi)^{3}}\frac{g^{4}N_{c}}{32\pi^{4}}\frac{11}{3}\frac{1}{k^{+}}\int_{x,w,x',w'}e^{-ik\cdot(w-w')}\Big[\frac{2}{\epsilon}+\mathcal{B}(x-w)\Big]\mathcal{O}(w,w')+{\rm c.c.}\notag\\
&=\frac{\alpha_{s}\beta_{0}}{4\pi}\,\mathcal{N}_{0}\!\int e^{-ik\cdot(w-w')}\Big[\frac{4}{\epsilon}+\mathcal{B}(x-w)+\mathcal{B}(x'-w')\Big]\mathcal{O}(w,w').
\label{eq:pd-II}
\end{align}
This is exactly the term the reviewer asks for, $\frac{\alpha_{s}\beta_{0}}{4\pi}[\frac{4}{\epsilon}+L(x-y)+L(x'-y')]\times$LO. It is present in the calculation; it is absent from \texttt{Final\_terms} because it cancels (Step~2). Region of the pole: large loop momentum $\tilde p$ of the $gg$ loop, $\int d^{2}\tilde p\,\frac{\tilde p^{2}}{\tilde p^{2}+\xi(1-\xi)k^{2}}$, finite at $\tilde p\to0$ (\texttt{s2} [3]): \textbf{UV}.

\medskip
\noindent\textbf{Step 2: Group~I Row~V, ``2b UV'' ($\bar{\mf{C}}_{1}^{\dagger}\mf{B}_{2}$ with the pair at one point).} From the checked row, with the coefficient $\frac{1}{k^{+}}\int_{\Lambda}^{k^{+}-\Lambda}dp^{+}[\dots]=2L-\frac{11}{6}$,
\begin{align}
\frac{d^{3}N_{2b\,\rm UV}}{d^{3}k}&=\frac{1}{(2\pi)^{3}}\frac{g^{4}N_{c}}{16\pi^{4}}\frac{1}{k^{+}}\Big[2L-\frac{11}{6}\Big]\int_{x,y,x',w'}e^{-ik\cdot(x-w')}\Big[\frac{2}{\epsilon}+\mathcal{B}(y-x)\Big]\notag\\&\quad\times K(x'-w')\cdot K(y-x)\,[U(x')-U(w')]^{ab'}[U(y)-U(x)]^{af}\rho^{b'}(x')\rho^{f}(y)+{\rm c.c.}
\notag
\end{align}
\emph{Relabelling} $x\to w$, $y\to x$: the phase becomes $e^{-ik\cdot(w-w')}$, the kernel $K(x'-w')\cdot K(x-w)$, the bracket $\frac{2}{\epsilon}+\mathcal{B}(x-w)$, and $[U(x')-U(w')][U(x)-U(w)]\rho(x')\rho(x)=+[U(w')-U(x')][U(w)-U(x)]\rho(x')\rho(x)$ (same order of the $\rho$'s). Since $\frac{g^{4}N_{c}}{16\pi^{4}}\frac{11}{6}=\frac{g^{2}}{2\pi^{2}}\cdot\frac{\alpha_{s}\beta_{0}}{4\pi}$,
\begin{align}
\frac{d^{3}N_{2b\,\rm UV}}{d^{3}k}\Big|_{-\frac{11}{6}}=-\frac{\alpha_{s}\beta_{0}}{4\pi}\,\mathcal{N}_{0}\!\int e^{-ik\cdot(w-w')}\Big[\frac{4}{\epsilon}+\mathcal{B}(x-w)+\mathcal{B}(x'-w')\Big]\mathcal{O}(w,w')=-\frac{d^{3}N^{\rm II}_{11/3}}{d^{3}k},
\label{eq:pd-V}
\end{align}
pole, $2\gamma_{E}$ and logarithm (provided the same $\mu$ is used). Region of the pole: $z\to x$, i.e.\ the two gluons of the pair at one point; the $z$ integral has no large-$z$ divergence (\texttt{s2} [1]: $\pi\log\frac{r^{2}}{\delta^{2}}$ to $10^{-9}$, contribution of $|z|>r$ below $10^{-14}$): \textbf{UV}. The value $+\pi[\frac{2}{\epsilon}+\dots]$ is the UV pole: in the checked row it is obtained as (scaleless $\int_{z}(z-x)^{-2}=\pi(\frac{2}{\epsilon_{\rm UV}}-\frac{2}{\epsilon_{\rm IR}})$) minus (an IR-divergent integral $=-\pi[\frac{2}{\epsilon_{\rm IR}}+\dots]$), i.e.\ $\pi[\frac{2}{\epsilon_{\rm UV}}+\dots]$, the same sign convention ($\log\Lambda_{\rm UV}^{2}\leftrightarrow\frac{2}{\epsilon_{\rm UV}}$) as Step~1.

\noindent\emph{So Row~II $+$ Row~V is a UV--UV cancellation, not a UV/IR mix-up.} Physical reason: a point-like $gg$ fluctuation of the produced gluon does not resolve the shock wave ($f^{adc}f^{fbe}U^{db}(x)U^{ce}(z)\to N_{c}U^{af}(x)$ at $z=x$). The observable counts \emph{dressed} gluons, $\bra{0}\Omega^{\dagger}S^{\dagger}\Omega\,a^{\dagger}a\,\Omega^{\dagger}S\Omega\ket{0}$, and where $S$ is locally constant $\Omega^{\dagger}S\Omega=S$: the UV part of the self-energy before the shock wave (Row~II) is undone by the same loop straddling it (Row~V). Consequently the emission coupling does \emph{not} run with the emitter--gluon distance $x-w$ in this observable (agrees in structure with Kovchegov--Weigert/Horowitz--Kovchegov, where the produced gluon runs at the collinear scale; \texttt{notes/lit\_dglap\_rc.tex}). In NLO JIMWLK the soft gluon is not measured and only one $\beta_{0}$ of the bare vertex is removed by the UV region of the two-gluon term (KLSZ, as quoted in \texttt{notes/lit\_finite\_dilute.tex}; not re-derived here), which is why the $b$-terms survive there.

\medskip
\noindent\textbf{Step 3: what is not written: the normalisation of the dressed measured gluon.} The single-gluon normalisation of the main file (Sec.~Normalization) contains $\braket{(\Psi)_{gg}|(\Psi)_{gg}}=\frac{g^{2}N_{c}}{8\pi^{3}}\big(2L-\frac{11}{6}\big)\int\frac{d^{2}\tilde p}{\tilde p^{2}}$, set to zero as scaleless. In the cross section it enters as $(\mf{N}^{2}-1)\times$LO:
\begin{align}
\frac{d^{3}N_{\rm norm}}{d^{3}k}=-\frac{g^{2}N_{c}}{8\pi^{3}}\Big(2L-\frac{11}{6}\Big)\int\frac{d^{d}\tilde p}{\tilde p^{2}}\;\frac{d^{3}N_{\rm LO}}{d^{3}k}
=\Big[\frac{\alpha_{s}\beta_{0}}{4\pi}-\mathfrak{a}L\Big]\Big(\frac{2}{\epsilon_{\rm UV}}-\frac{2}{\epsilon_{\rm IR}}\Big)\frac{d^{3}N_{\rm LO}}{d^{3}k}\;\;(=0\text{ for }\epsilon_{\rm UV}=\epsilon_{\rm IR}).
\label{eq:pd-norm}
\end{align}
Its UV half, $\frac{\alpha_{s}N_{c}}{4\pi}(\frac{11}{3}-4L)\frac{2}{\epsilon_{\rm UV}}$, is the light-cone-gauge $Z_{3}-1$ of the produced gluon; its IR half ($\tilde p\to0$, collinear pair) is the virtual partner of the real splitting of Row~IV.

\medskip
\noindent\textbf{Step 4: Group~II Row~IV (real $g\to gg$ after the shock wave).} The $r$ integral ($r=w-z$, daughter separation at fixed parent positions $y,y'$) diverges only at large $r$ (\texttt{s2} [2]: $\pi\log\frac{R^{2}}{u^{2}}$, independent of the short-distance cut $\delta$ to $10^{-5}$): \textbf{IR (collinear)}; in $d=2-\epsilon$ it exists for $\epsilon<0$. Its pole multiplies $\mathcal{P}_{gg+}(\xi)=2N_{c}[\xi(1-\xi)+\frac{\xi}{1-\xi}+\frac{1-\xi}{\xi}]$ (no $\delta(1-\xi)$: a real emission). After the extraction of the end points the poles are: JIMWLK term $-(L+\ell)\,\mathfrak{a}\frac{2}{\epsilon}\,\widehat{L}(k)\mathcal{N}_{0}$; DGLAP term $-\mathcal{C}\frac{2}{\epsilon}\int d\xi\,\mathcal{P}_{gg+}\Pi=-\mathcal{C}\frac{2}{\epsilon}[\mathcal{P}_{gg}\otimes\Pi]+\frac{\alpha_{s}\beta_{0}}{4\pi}\frac{2}{\epsilon}\,\mathcal{N}_{0}\widehat{L}(k)$, where $\mathcal{P}_{gg}=\mathcal{P}_{gg+}+\frac{\beta_{0}}{2}\delta(1-\xi)$ is the full $\overline{\rm MS}$ splitting function and $\mathcal{C}\frac{\beta_{0}}{2}\Pi(1)=\frac{\alpha_{s}\beta_{0}}{4\pi}\mathcal{N}_{0}\widehat{L}(k)$; the $\xi\to0$ (BFKL) subtraction has no pole, since the pole coefficient $\mathcal{G}_{0}(\xi)$ vanishes at $\xi=0$ for $k\neq0$. The ``RC'' term of the author is the second piece: its pole is collinear, as the reviewer says.

\medskip
\noindent\textbf{Step 5: counterterms.} Coupling: $g_{0}^{2}=\mu^{\epsilon}g^{2}(\mu)\big[1-\frac{\alpha_{s}\beta_{0}}{4\pi}\frac{2}{\epsilon}\big]$ in LO gives $-\frac{\alpha_{s}\beta_{0}}{4\pi}\frac{2}{\epsilon_{\rm UV}}\times$LO ($\rho$ and $U$ are not renormalised: the source is classical and $U$ is the given target field). Fragmentation ($\overline{\rm MS}$, full $\mathcal{P}_{gg}$): $+\mathcal{C}\frac{2}{\epsilon_{\rm IR}}[\mathcal{P}_{gg}\otimes\Pi]$ (with the $d$-dimensional Jacobian, Step~8).

\medskip
\noindent\textbf{Step 6: the table.} Coefficients in units of $\mathfrak{a}\,\frac{2}{\epsilon}\times\frac{d^{3}N_{\rm LO}}{d^{3}k}$ (resp.\ of $\mathcal{C}\frac{2}{\epsilon}[\,\cdot\otimes\Pi]$ for the DGLAP column); c.c.\ included for Group~I. Exact rational arithmetic in \texttt{s1}.
\begin{center}
\renewcommand{\arraystretch}{1.3}
\begin{tabular}{|p{4.3cm}|p{4.2cm}|p{1.1cm}|p{2.6cm}|p{4.9cm}|}\hline
pole & origin (region) & UV/IR & coefficient & cancelled by / absorbed into\\\hline
Group I Row II, $\mf{A}^{(3)}$ row & $gg$ loop, $\tilde p\to\infty$ & UV & $\frac{11}{6}-L+\frac{\ell}{2}$ & $\frac{11}{6}$: Row V (Step 2); $-L$: JIMWLK sector\\
Group I Row V, $2b$ UV & pair at one point, $z\to x$ & UV & $-\frac{11}{6}+2L$ & $-\frac{11}{6}$: Row II; $2L$: JIMWLK sector\\
dressed-gluon norm., UV half & $\tilde p\to\infty$ & UV & $\frac{11}{12}-L$ & coupling counterterm ($\frac{11}{12}$); $-L$: JIMWLK sector\\
coupling counterterm & $g_{0}\to g(\mu)$ & UV & $-\frac{11}{12}$ & UV half of the normalisation\\\hline
dressed-gluon norm., IR half & collinear pair, $\tilde p\to0$ & IR & $-\frac{11}{12}+L$ & $-\frac{11}{12}$: ``RC'' pole of Row IV; $L$: Row IV $\xi\to1$\\
Row IV, $\xi\to1$ end (JIMWLK term) & large $r$ & IR & $-L-\ell$ & $-L$: normalisation IR half; $-\ell$: left\\
Row IV, real $\mathcal{P}_{gg+}$ & large $r$ & IR & $-[\mathcal{P}_{gg}\otimes\Pi]$ $+\frac{11}{12}$ & $-[\mathcal{P}_{gg}\otimes\Pi]$: FF counterterm; $+\frac{11}{12}$: normalisation IR half\\
Row IV, $\xi\to0$ (BFKL) & --- & --- & $0$ & ($\mathcal{G}_{0}(0)=0$)\\
FF counterterm, full $\mathcal{P}_{gg}$ & $D_{g}\to D_{g}(\mu)$ & IR & $+[\mathcal{P}_{gg}\otimes\Pi]$ & real $\mathcal{P}_{gg+}$ $+$ virtual $\frac{\beta_{0}}{2}\delta$\\\hline
sum & & & UV: $\frac{\ell}{2}$; IR: $-\ell$ & $\to0$ for $\vee\to\infty$\\\hline
\end{tabular}
\end{center}
\noindent Single-$\epsilon$ bookkeeping (normalisation $=0$, author's rule): Row II $+$ Row V $+$ Row IV $+$ coupling CT $+$ FF CT $=(\frac{11}{6}-L+\frac{\ell}{2})+(-\frac{11}{6}+2L)+(-L-\ell+\frac{11}{12}-[\mathcal{P}\otimes\Pi])-\frac{11}{12}+[\mathcal{P}\otimes\Pi]=-\frac{\ell}{2}$, the same total (\texttt{s1} [f]). The $L$ poles cancel in the JIMWLK sector in both bookkeepings (UV: $-L+2L-L$; IR: $-L+L$). The remainder $-\frac{\ell}{2}$ is a pole times a coefficient that vanishes for $\vee\to\infty$: it is zero only if $\vee\to\infty$ is taken before $\epsilon\to0$ (or the poles are removed first); with the author's rule ``no poles written'' this is consistent, but it should be stated.

\medskip
\noindent\textbf{Step 7: physical account (answers to Q2 of the review).}
\begin{enumerate}
\item The virtual $\frac{11}{3}$ term of (I.89) is not missing: it is Eq.~\eqref{eq:pd-II} and is cancelled exactly by Eq.~\eqref{eq:pd-V} after the relabelling of Step~2. Hence no finite term $\frac{\alpha_{s}\beta_{0}}{4\pi}[\mathcal{B}(x-w)+\mathcal{B}(x'-w')-\mathcal{B}(w-w')]\times$LO, as proposed in the review, exists in this observable. The reviewer's claim is \textbf{wrong} on this point.
\item The reviewer is \textbf{right} that the RC pole, as it arises in Row~IV, is collinear (Step~4). What renormalises the coupling is the UV half of the dressed-gluon normalisation \eqref{eq:pd-norm}, i.e.\ $Z_{3}$ of the produced gluon. This is consistent with the light-cone-gauge Ward identity $Z_{g}=Z_{3}^{-1/2}$ ($g_{0}A_{0}=gA$, so the coupling of the classical source to the field needs no counterterm of its own): after the UV--UV cancellation of Rows~II and V the only UV pole left is that of $Z_{3}$, $+\frac{\alpha_{s}\beta_{0}}{4\pi}\frac{2}{\epsilon_{\rm UV}}$ (the $L$ part belongs to the JIMWLK sector), which is exactly what $g_{0}^{2}\to g^{2}(\mu)$ removes.
\item The IR half of the same normalisation is the virtual $\frac{\beta_{0}}{2}\delta(1-\xi)$ (and the $-2N_{c}L\,\delta(1-\xi)$) of the final-state splitting: real $\mathcal{P}_{gg+}$ $+$ virtual $=$ full $\overline{\rm MS}$ $\mathcal{P}_{gg}$, absorbed by the fragmentation counterterm. \textbf{No IR pole is left} once the FF counterterm with the full $\mathcal{P}_{gg}$ (including $\frac{\beta_{0}}{2}\delta(1-\xi)$) is subtracted, provided the normalisation is kept (or, in single-$\epsilon$ bookkeeping, provided the coupling counterterm is subtracted: it then cancels the RC pole formally, which is legitimate only because the normalisation integral is scaleless, UV $=$ IR).
\item ``Running coupling'': $N_{\rm RC}=-N_{\rm DGLAP}|_{\delta(1-\xi)}$ identically (add and subtract). Its content is the $\mu$ dependence: $\partial_{\log\mu^{2}}N_{\rm RC}=+\frac{\alpha_{s}\beta_{0}}{4\pi}$LO compensates the running of $g^{2}(\mu)$ in LO, and the $\delta$ part of $N_{\rm DGLAP}$ compensates the $\delta$ part of the evolution of $D_{g}(\mu)$. In the classification of KLSZ (``not all that is $\beta_{0}$ is $\beta$-function''): one $\beta_{0}$ (UV half of $Z_{3}$) is the $\beta$ function, the other (IR half, $\frac{\beta_{0}}{2}\delta(1-\xi)$ of $\mathcal{P}_{gg}$) is DGLAP; here both are tied to the produced gluon and are equal in magnitude because the integral is scaleless. Written locally in momentum space (Step~12), LO$+$RC $=$ LO$\times\frac{\alpha_{s}(k^{2})}{\alpha_{s}(\mu^{2})}$ to this order, the running coupling at the measured momentum, as in Fadin--Lipatov/BSV.
\end{enumerate}

\subsubsection{DGLAP in $\overline{\rm MS}$ form (P4)}

\noindent\textbf{Step 8: the Row~IV collinear integral in $d=2-\epsilon$, consistently.} Variables (checked row): $r=w-z$, $u=w'-w$, $Y=\frac{y+y'}{2}$, with $y=z+\xi(w-z)$, $y'=z+\xi(w'-z)$, so $y'-y=\xi u$; the Jacobian of $(z,w,w')\to(Y,r,u)$ is $1$ in any dimension. The $r$ integral gives $-\pi\Gamma(1-\frac{\epsilon}{2})\big(\frac{\mu^{2}u^{2}e^{\gamma_{E}}}{4}\big)^{\epsilon/2}\big[\frac{\delta^{mm'}}{\epsilon}+\frac{u^{m}u^{m'}}{u^{2}}\big]$. Continuing also $u$ (CDR) and rescaling $s\equiv y'-y=\xi u$ in $d$ dimensions,
\begin{align}
\int d^{d}u\,(u^{2})^{\epsilon/2}\,e^{iku}\,f(\xi u)=\xi^{-d}\xi^{-\epsilon}\int d^{d}s\,(s^{2})^{\epsilon/2}e^{i\frac{k}{\xi}s}f(s)=\xi^{-2}\int d^{d}s\,(s^{2})^{\epsilon/2}e^{i\frac{k}{\xi}s}f(s)
\label{eq:pd-resc}
\end{align}
(\texttt{s3} [a]: equal to $10^{-14}$ for $\epsilon=0.3,-0.2,0.05$; with a two-dimensional measure on $u$ the two sides differ by $\xi^{-\epsilon}$). So the \emph{bare} collinear term is $\xi^{-2}\times\big[\frac{2}{\epsilon}+\mathcal{B}(y'-y)\big]_{d}\times$(LO structure of the parent): in the bare result the $\log\xi^{2}$ cancels, as the reviewer says. The author's $\log\frac{(y'-y)^{2}\mu^{2}}{4\xi^{2}}$ arises from continuing only $r$.

\medskip
\noindent\textbf{Step 9: which counterterm.} What survives depends on the $\xi$ power that multiplies the counterterm. For a fragmentation function that is a number density in $z$ (operator definition, momentum sum rule $\sum_{h}\int dz\,zD_{h}=1$ in $d$ dimensions) collinear kinematics $P=z\,p$ in all $d$ transverse components gives
\begin{align}
\frac{dN_{h}}{dP^{+}d^{d}P}=\int\frac{dz}{z^{1+d}}\,D(z)\,\frac{dN_{g}}{dk^{+}d^{d}k}\Big|_{k=P/z},
\notag
\end{align}
and the $\overline{\rm MS}$ renormalisation $D_{0}=[\delta(1-\xi)+\frac{\alpha_{s}}{2\pi}\frac{2}{\epsilon}\mathcal{P}_{gg}(\xi)+\dots]\otimes D(\mu)$ (pure pole, Mellin convolution; $\frac{2}{\epsilon}$ in the normalisation of the $\mathcal{B}$ brackets) then gives the partonic counterterm $+\mathcal{C}\frac{2}{\epsilon}\int d\xi\,\mathcal{P}_{gg}(\xi)\,\xi^{-d}\,\widehat{L}_{d}(k/\xi)$, i.e.\ with $\xi^{-d}$, not $\xi^{-2}$. Expanding $\xi^{-d}=\xi^{-2}(1+\epsilon\log\xi)$:
\begin{align}
\text{bare}-\text{CT}(\xi^{-d})&\to\xi^{-2}\!\int e^{-i\frac{k}{\xi}\cdot(y-y')}\,\log\frac{(y'-y)^{2}\mu^{2}}{c_{0}^{2}\,\xi^{2}}\,\mathcal{O}\quad\text{(author's form: log of the \emph{measured} separation $u$)},\notag\\
\text{bare}-\text{CT}(\xi^{-2})&\to\xi^{-2}\!\int e^{-i\frac{k}{\xi}\cdot(y-y')}\,\log\frac{(y'-y)^{2}\mu^{2}}{c_{0}^{2}}\,\mathcal{O}\quad\text{(CXY form: log of the \emph{parent} separation)}.
\label{eq:pd-two}
\end{align}
Verified in $d=2\pm0.002$ with $d$-dimensional Fourier--Bessel transforms of a test function (\texttt{s3} [b]): $-4.38223062$ against the author form $-4.38221572$, $-4.04289809$ against the CXY form $-4.04288419$ (differences $O(\epsilon^{2})$); the difference of the two is $-\frac12\log\xi^{2}\,\Pi$ to 8 digits. So the $\log\xi^{2}$ does \emph{not} cancel in the subtracted result if the counterterm carries the $d$-dimensional Jacobian.

\medskip
\noindent\textbf{Step 10: scheme difference.} With $\mathcal{G}_{0}(\xi)\equiv\xi^{-2}\int e^{-i\frac{k}{\xi}\cdot(y-y')}\mathcal{O}=\Pi(\xi)$,
\begin{align}
\Delta\equiv\frac{d^{3}N_{\rm DGLAP}^{\rm author}}{d^{3}k}-\frac{d^{3}N_{\rm DGLAP}^{\rm CXY}}{d^{3}k}=+\,\mathcal{C}\int_{0}^{1}d\xi\,\mathcal{P}_{gg}(\xi)\,\log\xi^{2}\;\mathcal{G}_{0}(\xi),
\label{eq:pd-Delta}
\end{align}
the reviewer's $\Delta$ (same sign). Its $\delta(1-\xi)$ part vanishes, it is integrable at $\xi\to0$ ($\mathcal{G}_{0}(\xi)\to0$), and it does not depend on $\mu$. Equivalently $D^{\rm CXY}=\big[\delta(1-\xi)+\frac{\alpha_{s}}{2\pi}\mathcal{P}_{gg}(\xi)\log\xi^{2}\big]\otimes D$. Since $\int_{0}^{1}d\xi\,\xi\,\mathcal{P}_{gg}\log\xi^{2}=-14.3225$ for $N_{c}=3$, $N_{f}=0$ (\texttt{s3} [d], closed form $4N_{c}[\frac54-\zeta_{2}-\frac34-\frac{7}{144}]$), the two definitions cannot both satisfy the momentum sum rule; the number-density definition, renormalised with a pure pole, does.

\medskip
\noindent\textbf{Recommendation.} The author's form \eqref{eq:pd-two} (first line) is the one that corresponds to $\overline{\rm MS}$ fragmentation functions defined as number densities with the $d$-dimensional collinear kinematics; the CXY form corresponds to the counterterm with the two-dimensional Jacobian $\xi^{-2}$ (CXY write the factorisation formula with $dz/z^{2}$; AABKL App.~C use $\zeta^{-(D-2)}$ but compensate it by the factor $(1-\xi)^{-(5-D)}$ in their FF definition, which is the same choice). I recommend keeping the author's form, stating the definition of the counterterm (Step~9) explicitly, and giving $\Delta$ so that the CXY convention can be recovered. Not checked here: an independent comparison of the collinear finite terms with a standard NLO collinear calculation of $pp\to gX$ (that would settle the convention empirically).

\medskip
\noindent\textbf{Step 11: final DGLAP and RC terms} ($\overline{\rm MS}$, number-density FF, single scale $\mu$, phase $e^{-ik\cdot(y-y')/\xi}$, $\xi\to0$ subtraction explicit):
\begin{align}
\frac{d^{3}N_{\rm DGLAP}}{d^{3}k}&=-\mathcal{C}\int_{0}^{1}d\xi\,\Big\{\mathcal{P}_{gg}(\xi)\,\mathcal{F}(\xi)-\frac{2N_{c}}{\xi}\,\mathcal{F}(0)\Big\},\qquad
\mathcal{F}(\xi)\equiv\frac{1}{\xi^{2}}\int_{x,y,x',y'}e^{-ik\cdot(y-y')\frac{1}{\xi}}\,\log\frac{(y'-y)^{2}\mu^{2}}{c_{0}^{2}\,\xi^{2}}\;\mathcal{O}(y,y'),\notag\\
\mathcal{P}_{gg}(\xi)&=2N_{c}\Big[\frac{\xi}{(1-\xi)_{+}}+\frac{1-\xi}{\xi}+\xi(1-\xi)\Big]+\frac{\beta_{0}}{2}\delta(1-\xi),\qquad\int_{0}^{1}d\xi\,\frac{\xi\,f(\xi)}{(1-\xi)_{+}}\equiv\int_{0}^{1}d\xi\,\frac{\xi f(\xi)-f(1)}{1-\xi},\notag\\
\mathcal{F}(0)&=-\frac{4\pi}{k^{2}}\,\mathcal{X}=-\frac{4\pi}{k^{2}}\int_{x,x',z}\frac{(x-z)\cdot(x'-z)}{(x-z)^{2}(x'-z)^{2}}\big[U^{ce'}(z)-U^{ce'}(x')\big]\big[U^{ce}(z)-U^{ce}(x)\big]\rho^{e'}(x')\rho^{e}(x),\notag\\
\frac{d^{3}N_{\rm RC}}{d^{3}k}&=+\mathcal{C}\,\frac{\beta_{0}}{2}\,\mathcal{F}(1)=+\frac{\alpha_{s}\beta_{0}}{4\pi}\,\mathcal{N}_{0}\int_{x,y,x',y'}e^{-ik\cdot(y-y')}\,\log\frac{(y'-y)^{2}\mu^{2}}{c_{0}^{2}}\;\mathcal{O}(y,y').
\label{eq:pd-final}
\end{align}
The second term in the braces is the author's $[\frac{1-\xi}{\xi}]_{+}$ ($\mathcal{F}(0)$ is the same in the author's and in the CXY form: \texttt{s3} [c], both $\to-\frac{4\pi}{k^{2}}g(0)$, difference $4\times10^{-6}$ at $\xi=0.2$ and $<10^{-8}$ below); with $\int_{0}^{1}$ in place of $\int_{k^{+}/\vee}^{1}$ the remnant $-2N_{c}\mathcal{C}\frac{4\pi}{k^{2}}\mathcal{X}$ stays among the finite terms of Row~IV. CXY form: replace $\log\frac{(y'-y)^{2}\mu^{2}}{c_{0}^{2}\xi^{2}}\to\log\frac{(y'-y)^{2}\mu^{2}}{c_{0}^{2}}$ and add $\Delta$ of \eqref{eq:pd-Delta}. $\partial_{\log\mu^{2}}N_{\rm DGLAP}=-\mathcal{C}\int d\xi\,\mathcal{P}_{gg}\Pi$: standard DGLAP of $D_{g}(\mu)$.

\medskip
\noindent\textbf{Step 12: local (tail-free) form; recommended for the paper.} With $\log\frac{r^{2}\mu^{2}}{c_{0}^{2}}=-\frac{1}{\pi}\int d^{2}l\,\frac{e^{-il\cdot r}-\theta(\mu-l)}{l^{2}}$ one has exactly
\begin{align}
\mathcal{F}(\xi)=-\log\frac{k^{2}}{\mu^{2}}\,\Pi(\xi)+\frac{1}{\xi^{2}}\,\mathcal{R}\Big(\frac{k}{\xi}\Big),\qquad
\mathcal{R}(q)\equiv-\frac{1}{\pi}\int d^{2}l\,\frac{\widehat{L}(q+l)-\theta(|q|-|l|)\,\widehat{L}(q)}{l^{2}},\qquad\lim_{\xi\to0}\xi^{-2}\mathcal{R}(k/\xi)=\mathcal{F}(0),
\label{eq:pd-loc}
\end{align}
($\mathcal{R}$ does not depend on $\mu$; \texttt{s4} [b] checks the identity against the Hankel form, $-0.49034734$ both). Hence
\begin{align}
\frac{d^{3}N_{\rm LO}}{d^{3}k}+\frac{d^{3}N_{\rm RC}}{d^{3}k}&=\mathcal{N}_{0}\widehat{L}(k)\Big[1-\frac{\alpha_{s}\beta_{0}}{4\pi}\log\frac{k^{2}}{\mu^{2}}\Big]+\mathcal{C}\frac{\beta_{0}}{2}\mathcal{R}(k),\qquad
\frac{d^{3}N_{\rm DGLAP}}{d^{3}k}&=\mathcal{C}\log\frac{k^{2}}{\mu^{2}}\int_{0}^{1}d\xi\,\mathcal{P}_{gg}(\xi)\,\Pi(\xi)\notag\\&\quad-\mathcal{C}\!\int_{0}^{1}\!d\xi\Big\{\mathcal{P}_{gg}\frac{\mathcal{R}(k/\xi)}{\xi^{2}}-\frac{2N_{c}}{\xi}\mathcal{F}(0)\Big\}.
\notag
\end{align}
The local pieces carry all the $\mu$ dependence and no subtraction at $\xi=0$ is needed in them ($\Pi(\xi)\to0$). The non-local, $\mu$-independent remainder, whose $\delta(1-\xi)$ parts cancel between DGLAP and RC,
\begin{align}
\frac{d^{3}N_{\mathcal{R}}}{d^{3}k}=-\mathcal{C}\int_{0}^{1}d\xi\,\Big\{\mathcal{P}_{gg+}(\xi)\,\frac{\mathcal{R}(k/\xi)}{\xi^{2}}-\frac{2N_{c}}{\xi}\mathcal{F}(0)\Big\},\qquad\mathcal{P}_{gg+}=\mathcal{P}_{gg}-\frac{\beta_{0}}{2}\delta(1-\xi),
\label{eq:pd-R}
\end{align}
belongs to the finite terms of Row~IV. It contains the BFKL ($\xi\to0$) subtraction and the soft-parent region (Step~13).

\subsubsection{The large-distance tail of Group~II Row~IV}

\noindent\textbf{Step 13: the tail and its origin.} At large parent position $Y$ the LO structure tends to $G_{\infty}(Y)/Y^{2}$, $G_{\infty}=\int_{x,x'}[U(Y)-U(x')][U(Y)-U(x)]\rho\rho$, independent of $s=y'-y$; the logarithm then survives the Fourier transform, $\int d^{2}s\,e^{-iqs}\log s^{2}=-\frac{4\pi}{q^{2}}$, and $\mathcal{F}(\xi)\to-\frac{4\pi}{k^{2}}G_{\infty}/Y^{2}$ for every $\xi$. Tail coefficients in units $\mathcal{C}\frac{4\pi}{k^{2}}\int\frac{d^{2}Y}{Y^{2}}G_{\infty}$ (\texttt{s6} [2]): DGLAP $+\frac{N_{c}}{6}$, RC $-\frac{11N_{c}}{6}$, remnant $-2N_{c}$, traceless tensor term $+\frac{N_{c}}{3}\cos2\varphi$; total $2N_{c}[-\frac{11}{6}+\frac16\cos2\varphi]$ (OPEN~II/III-1). The $-\frac{11}{6}=\int_{0}^{1}\frac{d\xi}{\xi_{+}(1-\xi)_{+}}N(\xi)$ is the KLSZ double-plus integral of the real splitting function; the RC term and the $\delta$ part of DGLAP cancel each other also in the tail, so the tail is not ``tied to RC''; it is not related to the pole either ($\mathcal{G}_{0}$ has no tail for $k\neq0$). In momentum space it is the region $k/\xi+l\to0$ of $\mathcal{R}$: a parent much softer than the measured gluon, $|q|\ll|k|$, which splits into two hard gluons; since the LO spectrum behaves as $\widehat{L}(q)\propto1/q^{2}$ at small $q$ (Weizs\"acker--Williams field of the colour charge $\int[U_{\infty}-U(x)]\rho(x)$, non-zero also for a neutral projectile), $\int d^{2}q\,\widehat{L}(q)/k^{2}$ diverges logarithmically. \texttt{s4} [b] (model $\widehat{L}=\frac{1}{p^{2}+m^{2}}-\frac{1}{p^{2}+M^{2}}$): the coordinate form grows like $-\frac{2}{k^{2}}\log\frac1m$ ($-0.8889$ per unit $\log\frac1m$, expected $-0.888889$), the local form $-\log\frac{k^{2}}{\mu^{2}}\widehat{L}(k)$ is $m$-independent. So the DGLAP and RC terms \emph{as written in coordinate space} are each logarithmically IR divergent; the local terms of Step~12 are not.

\medskip
\noindent\textbf{Step 14: cancellation.} The two-gluon final state of (CS2) is $\mathcal{A}=[\text{B-part}]-[\text{C-part}]$ with B-part $=U(w)\mf{B}_{2}-U^{\dagger}(z)\bar{\mf{B}}_{2}$ (two gluons from $\Omega$ before the shock wave, and from $\Omega^{\dagger}$ after it) and C-part $=\bar{\mf{C}}_{1}[U(z)U(y)\mf{A}^{(1)}-U(z)\bar{\mf{A}}^{(1)}]$ (the LO gluon split by $\Omega^{\dagger}$ after the shock wave); $|\text{B}|^{2}$ is Group~II Row~I (second part, single-$\rho$ $\mf{B}_{2}$), $|\text{C}|^{2}$ is Group~II Row~IV, and $-2{\rm Re}\,\text{C}^{\dagger}\text{B}$ is Group~III Row~III ($+$c.c.). In LCPT, with $E_{1}$ the energy of the one-gluon (parent) state and $E_{2}$ that of the pair, $b=\frac{V_{\rho g}V_{ggg}}{E_{1}E_{2}}$ and $c\,A=\frac{V_{\rho g}V_{ggg}}{E_{1}(E_{2}-E_{1})}$, so
\begin{align}
b-c\,A=-\frac{V_{\rho g}V_{ggg}}{E_{2}(E_{2}-E_{1})}:
\notag
\end{align}
the parent pole $1/E_{1}\propto q^{+}/q^{2}$ cancels. The colour factors are $[U(w)U(z)-U(x)]$ (B, in the reduced form) and $[U(y)-U(x)]$ (C), so that $\mathcal{A}=[U(w)U(z)-U(y)]\,b+[U(y)-U(x)]\,(b-cA)$. With the author's explicit momentum-space amplitudes (Eq.~(psiggrhomom) for $\mf{B}_{2}$, Eqs.~(rhoindep), (lowav1mom) for $\mf{C}_{1}\mf{A}^{(1)}$, labels matched, $f^{acb}=-f^{abc}$) this is verified numerically (\texttt{s4} [a]): $|\mf{B}_{2}|,|\mf{C}_{1}\mf{A}^{(1)}|\propto1/|q|$ ($7.6\times10^{4}$ at $|q|=10^{-4}$) while $|\mf{B}_{2}-\mf{C}_{1}\mf{A}^{(1)}|\to3.636$, $1.562$, $7.248$ (three random points); the coefficients of $1/q^{2}$ of $|\text{B}|^{2}$, $|\text{C}|^{2}$, $2\text{B}\cdot\text{C}$ are equal ($58.2183$, $58.2165$, $2\times58.2174$ at $|q|=10^{-4}$, first point), i.e.\ weights $1,1,-2$. Consequences:
\begin{enumerate}
\item The large-distance (soft-parent) tail of Row~IV is cancelled by the large-distance region of Group~II Row~I (second part) and Group~III Row~III, \emph{where both gluons are far from the projectile together} ($w,w',z\to\infty$ at fixed $w-z$). This region was not analysed in \texttt{Finite\_terms\_final} (tails there were taken at fixed $w,w'$), which is why ``no compensating term'' was found. The cancellation is pointwise in $p^{+}$ (in $\xi$) and in the relative momentum, hence it holds separately in the $\log\frac{\vee}{k^{+}}$ (BFKL), $\log\frac{k^{+}}{\Lambda}$ (JIMWLK) and finite sectors, including the $\cos2\varphi$ part, provided all three rows are split with the same end-point subtractions.
\item It needs $U(w)U(z)f=fU(y)$ for the pair at large distance, i.e.\ $U$ slowly varying there (true for a finite target, $U\to1$). For a colour-neutral projectile no tail survives even without this: the remaining pole term $[U(w)U(z)-U(y)]\,b$ multiplies $\int_{x}K(x-Y)\rho(x)$, which falls like $1/Y^{2}$. For a coloured projectile the remainder is the non-singlet BFKL/JIMWLK-type term (OPEN~II/III-4(b)).
\item It does not require any change of the RC/DGLAP split for consistency of the sum, but the coordinate-space DGLAP and RC terms are separately IR divergent; written as in Step~12, DGLAP and RC are finite and the remainder $N_{\mathcal{R}}$ \eqref{eq:pd-R} must be grouped with the finite terms of Group~II Row~I (second part) and Group~III Row~III under one integral over the pair position.
\end{enumerate}
\noindent\textbf{What remains (item 4).} The cancellation is established for the rows as derived from (CS2) (amplitude identity, verified numerically from the author's formulas). It was \emph{not} checked term by term that the finite terms of Group~II Row~I (second part, $N_{2a}\dots N_{2f}$) and of Group~III Row~III, as written in \texttt{Final\_terms\_checked}/\texttt{Finite\_terms\_final}, carry the compensating tail $+2N_{c}\mathcal{C}\frac{4\pi}{k^{2}}\int\frac{d^{2}Y}{Y^{2}}[\frac{11}{6}-\frac16\cos2\varphi]G_{\infty}$ in the joint limit; this has to be checked because the scaleless rule $\int d^{2}z/z^{2}=0$ and the recoil rewritings used in those rows can move large-distance pieces between sectors.

\subsubsection{Phase convention (P5)}\label{sec:pd-phase}
\noindent\textbf{Inventory} (\texttt{s5}): \texttt{Final\_terms.tex} and \texttt{Limiting\_cases\_final.tex} use $e^{-ik\cdot(\text{unprimed}-\text{primed})}$ everywhere (47 resp.\ 39 phases; the recoil factors $e^{-ik\cdot(x-z)\frac{p^{+}}{k^{+}}}$, $e^{-ik\cdot(z-x)\frac{q^{+}}{k^{+}}}$ are consistent with their base phases). All working files in \texttt{Main files/} use the reversed form: \texttt{Calculation\_main} ($e^{-ik(w'-w)}$, $(w'-x)$, $(w'-z)$; 24), \texttt{Group\_I} (58), \texttt{Group I, Row V} (29), \texttt{Group\_II} (57, incl.\ $(y'-y)$), \texttt{Group II, Row IV} (31, $(w'-w)$, $(y'-y)$), \texttt{Group\_III} (25, $(w'-w)$, $(y'-w)$), \texttt{Divergent\_terms} (10), \texttt{BFKL} (59), \texttt{Evolution\_(JIMWLK)\_terms} (28), \texttt{Evolution\_two\_rho} (15), \texttt{Evolution\_three\_rho} (63), \texttt{4rho} (30), \texttt{scales} (4). In these files the conjugated bracket stands at the primed point, so the reversed phase gives each row at $-k$. The checked copies are corrected.

\noindent\textbf{Uniform form} (conjugated bracket at the primed point, to the left): $e^{-ik\cdot(w-w')}$; after integrating out the measured-gluon position, $e^{-ik\cdot(z-w')}e^{-ik\cdot(x-z)\frac{p^{+}}{k^{+}}}$ and $e^{-ik\cdot(x-w')}e^{-ik\cdot(z-x)\frac{q^{+}}{k^{+}}}$ (Group~I Row~V), $e^{-ik\cdot(w-y')}e^{-ik\cdot(z-y')\frac{p^{+}}{k^{+}}}$ (Group~III), $e^{-ik\cdot(y-y')}$ (LO, RC), $e^{-ik\cdot(y-y')/\xi}$ (DGLAP and Row~IV finite terms).

\noindent\textbf{Effect.} $k\to-k$ leaves the part even in $k$ unchanged (symmetrised two-$\rho$ terms, hence all Gaussian/colour-singlet averaged results) and flips the sign of the odd part: the commutator remnants $\propto f^{abc}\rho^{c}$ and all three-$\rho$ (cubic) terms, which are of the form $i f\times$(real kernel)$\times$phase $+$ c.c. Mixing conventions between rows would therefore give wrong relative signs of the cubic terms (e.g.\ in the three-$\rho$ tail cancellation between Groups I and III), and the overall sign of the odd azimuthal part for a non-Gaussian projectile ($\langle\rho\rho\rho\rangle\propto d^{abc}$). For a Gaussian projectile the cubic terms vanish and the convention does not matter.

\subsubsection{Summary}
\begin{itemize}
\item Poles: every $1/\epsilon$ is accounted for (table, Step~6). Row~II's $\frac{11}{3}$ (the reviewer's virtual term) and Row~V's $-\frac{11}{6}\times2$ cancel UV against UV. The coupling is renormalised by the UV half of the scaleless normalisation of the dressed produced gluon ($Z_{3}$, $Z_{g}=Z_{3}^{-1/2}$ in light-cone gauge); its IR half is the virtual $\frac{\beta_{0}}{2}\delta(1-\xi)$ that completes $\mathcal{P}_{gg+}$ to the $\overline{\rm MS}$ $\mathcal{P}_{gg}$. No pole is left except $-\frac{\ell}{2}\frac{2}{\epsilon}$, which vanishes for $\vee\to\infty$ taken first.
\item DGLAP: in full CDR the bare log has no $\xi^{2}$; the $\log\xi^{2}$ of the author's form reappears from the $d$-dimensional Jacobian $\xi^{-d}$ of the $\overline{\rm MS}$ counterterm for number-density FFs. Author's form and CXY form differ by $\Delta$, Eq.~\eqref{eq:pd-Delta}. Final form: Eq.~\eqref{eq:pd-final}, or (recommended) the local form of Step~12.
\item Row~IV tail: soft-parent region; cancelled by Group~II Row~I (second part) and Group~III Row~III through $\mf{B}_{2}\to\mf{C}_{1}\mf{A}^{(1)}$ at small parent momentum. Open: the term-by-term check on the finite terms of those two rows as written.
\end{itemize}


% ---- inlined copy of paper/notes/fix_singlet.tex (20 Nov 2026 06:29 UTC) ----
\subsection{Finite terms in colour-singlet form}\label{sec:singlet}
\noindent The finite terms, the DGLAP term and the running-coupling term are written here, like the BFKL and JIMWLK terms, averaged over the projectile, with traces of adjoint Wilson lines. Labels, kernels and conventions of every row are those of the finite-term section (Finite\_terms\_final.tex; for the soft-remainder rows Group I Rows III, V and Group III Rows III, V, VI the integrable forms given there are used as the current version): only the colour factor is replaced, so that the replacement can be repeated mechanically if a kernel changes. The common prefactor is
\begin{align}
\frac{1}{(2\pi)^{3}}\,\frac{g^{4}}{16\pi^{5}}\,\frac{1}{k^{+}}\,\frac{1}{N_{c}^{2}-1}\,,\notag
\end{align}
the numerical factors of the rows are written inside the braces. ``c.c.'' is implied for Groups I and III and is written explicitly.

\subsubsection{Projectile average: what is assumed}
\noindent The projectile enters every term as an ordered product of two or three colour-charge operators, $C^{a_{1}\dots a_{n}}\,\rho^{a_{1}}(p_{1})\cdots\rho^{a_{n}}(p_{n})$, with a $c$-number tensor $C$ made of adjoint Wilson lines, $f^{abc}$ and $N_{c}$, and $[\rho^{a}(x),\rho^{b}(y)]=if^{abc}\rho^{c}(x)\,\delta^{(2)}(x-y)$. We assume
\begin{itemize}
\item[(A1)] \emph{Global colour invariance} (colour-singlet projectile): the projectile density matrix commutes with the total charge $Q^{a}=\int_{x}\rho^{a}(x)$. Then $\langle\rho^{a}(x)\rangle=0$ (there is no invariant vector in the adjoint representation) and
\begin{align}
\langle\rho^{a}(x)\,\rho^{b}(y)\rangle=\frac{\delta^{ab}}{N_{c}^{2}-1}\,n(x,y)\,,\qquad n(x,y)\equiv\langle\rho^{f}(x)\rho^{f}(y)\rangle=n(y,x)=n(x,y)^{*}\,;\label{eq:s-A1}
\end{align}
the symmetry follows from $\langle[\rho^{a}(x),\rho^{b}(y)]\rangle=if^{abc}\langle\rho^{c}\rangle\delta^{(2)}(x-y)=0$, the reality from the hermiticity of $\rho^{f}$.
\item[(A2)] \emph{No symmetrised three-point function}: $S^{abc}(x,y,z)\equiv\langle\{\rho^{a}(x)\rho^{b}(y)\rho^{c}(z)\}_{\rm sym}\rangle=0$, where $\{\dots\}_{\rm sym}$ is the average over the six orderings.
\end{itemize}
(A2) holds for a Gaussian projectile (Gaussian Weyl symbol of the density matrix, $W[\rho]\propto e^{-\frac12\int\rho\,\mu^{-1}\rho}$ as in the MV model): the symmetrised products are the classical moments of $W$, and its odd moments vanish. In general, by (A1), $S^{abc}=d^{abc}g_{d}(x,y,z)+f^{abc}g_{f}(x,y,z)$ with real $g_{d}$ totally symmetric and real $g_{f}$ totally antisymmetric in the three points. $g_{d}$ is C-odd (the odderon). Both parts vanish for an ensemble invariant under an antiunitary map $\rho^{a}\to-\rho^{a}$ (it preserves the algebra; then $\langle\{\rho\rho\rho\}_{\rm sym}\rangle$ is real and equal to minus itself), e.g.\ adjoint charges, or $SU(2)$ charges, with a density matrix that is a function of $\rho_{p}\cdot\rho_{q}$. C-evenness alone does not remove $g_{f}$: a colour-invariant ensemble of three charges with a term $f^{abc}\rho^{a}_{1}\rho^{b}_{2}\rho^{c}_{3}$ in $\ln$ of the density matrix has $g_{f}\neq0$ ($SU(2)$, where $d^{abc}=0$; see the checks below). (A2) is therefore stated as an assumption about the projectile, satisfied by a Gaussian one; nothing else is assumed. No assumption about the four-point function is needed (four-$\rho$ rows contribute only to the evolution terms).

\noindent\textbf{Quadratic terms.} By \eqref{eq:s-A1} the order of the two operators does not matter (the commutator is a one-$\rho$ operator), and
\begin{align}
\langle C^{ab}\rho^{a}(p)\rho^{b}(q)\rangle=\frac{C^{aa}}{N_{c}^{2}-1}\,n(p,q)\,.\label{eq:s-quad}
\end{align}

\noindent\textbf{Cubic terms.} Reordering an ordered product into the symmetrised one produces single commutators, which are \emph{two}-$\rho$ operators, and double commutators, which are one-$\rho$ operators (zero average by (A1)). Exactly, for $A,B,C$ standing for $\rho^{a}(p_{1}),\rho^{b}(p_{2}),\rho^{c}(p_{3})$: $ABC=\{ABC\}_{\rm sym}+\frac12\big([A,B]C+A[B,C]+[A,C]B\big)+(\text{double commutators})$. With $\Phi\equiv C^{abc}f^{abc}$ (the Wilson lines evaluated at the coincident points where the $\delta$-function puts them) this gives
\begin{align}
\langle C^{abc}\rho^{a}(p_{1})\rho^{b}(p_{2})\rho^{c}(p_{3})\rangle&=C^{abc}S^{abc}+\frac{i\,\Phi}{2(N_{c}^{2}-1)}\Big[\delta^{(2)}(p_{1}-p_{2})\,n(p_{1},p_{3})+\delta^{(2)}(p_{2}-p_{3})\,n(p_{1},p_{2})-\delta^{(2)}(p_{1}-p_{3})\,n(p_{1},p_{2})\Big],\label{eq:s-ord}\\
\langle C^{abc}\rho^{a}(p_{1})\{\rho^{b}(p_{2}),\rho^{c}(p_{3})\}\rangle&=2C^{abc}S^{abc}+\frac{i\,\Phi}{N_{c}^{2}-1}\Big[\delta^{(2)}(p_{1}-p_{2})\,n(p_{1},p_{3})-\delta^{(2)}(p_{1}-p_{3})\,n(p_{1},p_{2})\Big],\label{eq:s-a23}\\
\langle C^{abc}\{\rho^{a}(p_{1}),\rho^{b}(p_{2})\}\rho^{c}(p_{3})\rangle&=2C^{abc}S^{abc}+\frac{i\,\Phi}{N_{c}^{2}-1}\Big[\delta^{(2)}(p_{2}-p_{3})\,n(p_{1},p_{2})-\delta^{(2)}(p_{1}-p_{3})\,n(p_{1},p_{2})\Big],\label{eq:s-a12}
\end{align}
($\Phi$ of the ordering written first). Under (A2) the first term vanishes, \emph{but the contact terms do not}: they are two-point terms $\propto\delta^{(2)}(p_{i}-p_{j})\,n$, of the same order as the quadratic terms, and a colour-neutral projectile does not remove them. The cubic rows therefore do \emph{not} vanish after the average; what vanishes is their symmetrised part. The contact terms are written below as quadratic terms at coincident points: $\mathfrak{K}\big|_{q\to p}$ denotes the kernel of the row with the point $q$ replaced by $p$ (the integral over $q$ is removed by the $\delta$-function); $n$ is written as $\langle\rho^{f}\rho^{f}\rangle$.

\noindent\textbf{Parity in $k$.} Write a term as $A(k)\langle O\rangle$, $A$ the $c$-number factor (prefactor and kernel), $O$ the colour operator; its c.c.\ is $A(k)^{*}\langle O^{\dagger}\rangle=A(k)^{*}\langle O\rangle^{*}$, so the singlet form of [term + c.c.] is [singlet form of the term] + its complex conjugate. All kernels depend on $k$ through $e^{-ik\cdot r}$, $\mathcal{E}(s)$, $\mathcal{J}(s)$, $\widehat{\mathcal{E}}(s)$, $\Psi(s)$, $\pi\delta(s)-i{\rm P}\frac1s$, all of the form $f(-s)=f(s)^{*}$; the quadratic rows have real prefactors, the cubic rows an explicit $i$ (in Group I Row II, $N_{2b}$, the $i$ is in the Fourier transform of the odd function $\mathcal{R}^{i}(\mathbf{q},\mathbf{p})$). Hence $A(-k)=A(k)^{*}$ for the quadratic and $A(-k)=-A(k)^{*}$ for the cubic rows (for Group II Row I, first part, which has no c.c., $N_{1b}$ is $N_{1a}$ with $w\leftrightarrow w'$ and the colour factors in the reverse order, which gives the same statement for $N_{1a}+N_{1b}$). With $\langle O\rangle=\mathsf{S}+i\mathsf{C}$ ($\mathsf{S}=C^{abc}S^{abc}$ and $\mathsf{C}$ real; $\mathsf{C}$ = contact part): term $+$ c.c.\ $=\mathsf{S}\,[A(k)-A(-k)]+i\mathsf{C}\,[A(k)+A(-k)]$ for the cubic rows. \emph{The symmetrised part of the cubic terms is odd under $k\to-k$, their contact part is even}, and the quadratic terms are even. The part of the finite terms odd in $k$ is exactly the $S^{abc}$ part; it needs a C-odd ($d^{abc}$) or a $g_{f}$ three-point function, and is zero under (A2).

\noindent\textbf{Rows not listed.} Group I Row I and Group II Rows II, III and Group III Row IV contain only four-$\rho$ terms (evolution only); Group I Row IV and Group III Row I cancel identically (operator identity, hence also in singlet form).

\subsubsection{Colour structures and their singlet factors}
\noindent Notation ($U$ adjoint, real, $U^{T}U=1$, $U^{\dagger}=U^{T}$; $(T^{e})_{bc}=-if^{ebc}$, ${\rm Tr}\,T^{a}T^{b}=N_{c}\delta^{ab}$; at one point $f^{abc}U^{bd}U^{ce}=U^{ah}f^{hde}$, equivalently $U\,T^{b}U^{\dagger}=T^{a}U^{ab}$):
\begin{align}
\mf{D}(a,b)&\equiv{\rm Tr}\{U^{\dagger}(a)U(b)\}=U^{ij}(a)U^{ij}(b)\,,\notag\\
\mf{F}(a,b,c)&\equiv{\rm Tr}\{T^{e}U^{\dagger}(a)U(b)T^{e}U^{\dagger}(b)U(c)\}=f^{ijk}f^{lmn}\,U^{il}(a)\,U^{jm}(b)\,U^{kn}(c)\,,\notag\\
\mf{Q}(a,b;c,d)&\equiv{\rm Tr}\{T^{e}U^{\dagger}(a)U(b)T^{e}U^{\dagger}(c)U(d)\}\,.\notag
\end{align}
$\mf{D}$ is symmetric, $\mf{F}$ totally symmetric; $\mf{Q}(a,b;c,d)=\mf{Q}(c,d;a,b)=\mf{Q}(d,c;b,a)$. At coincident points: $\mf{D}(a,a)=N_{c}^{2}-1$, $\mf{F}(a,a,c)=N_{c}\mf{D}(a,c)$, $\mf{F}(a,a,a)=N_{c}(N_{c}^{2}-1)$, $\mf{Q}(a,a;c,d)=N_{c}\mf{D}(c,d)$, $\mf{Q}(a,b;b,d)=\mf{F}(a,b,d)$, $\mf{Q}(a,b;c,a)=\mf{F}(a,b,c)$. [The second form of $\mf{F}$: $U(b)T^{e}U^{\dagger}(b)=T^{h}U^{he}(b)$ and $(T^{e})_{ij}=-if^{eij}$.] A trace with one $U$ at each of two $\rho$ points comes from $U^{ab}(p)U^{ab'}(q)\delta^{bb'}$; a product $f\,f\,UUU$ with the $\rho$ indices contracted by $\delta$ (quadratic) or by $f^{abc}$ (cubic, $\Phi$) is an $\mf{F}$; the structures with four Wilson lines of Group II Row I and Group III Row II give $\mf{Q}$. Example (Group I Row V, $N_{2a}$): $[U(x')-U(w')]^{ab'}\big[-f^{adc}f^{fbe}U^{db}(x)U^{ce}(z)+N_{c}U^{af}(x)\big]\delta^{b'f}=-\mf{F}(x',x,z)+\mf{F}(w',x,z)+N_{c}[\mf{D}(x',x)-\mf{D}(w',x)]$. The Wilson-line arguments are ordered by the fixed list $x,y,x',y',z,w,w'$.

\noindent All singlet factors below were obtained by machine from the colour structures as printed in the finite-term section and verified numerically (Sec.~\ref{sec:s-checks}).

\noindent\textbf{Quadratic structures} ($\langle C^{ab}\rho^{a}(p)\rho^{b}(q)\rangle=\mathcal{S}\,\langle\rho^{f}(p)\rho^{f}(q)\rangle/(N_{c}^{2}-1)$, either order):
\begin{align}
\mathcal{S}_{\rm I.II}&=-\mf{D}(x,x')+\mf{D}(x,w')+\mf{D}(x',w)-\mf{D}(w,w')\notag\\&\qquad[\mathcal{L}^{a}_{x'w'}\big[U^{ab}(w)-U^{ab}(x)\big]\rho^{b}(x);\ \ \rho\text{ at }x',\,x]\notag\\
\mathcal{S}_{\rm I.V}^{2a}&=N_{c}\mf{D}(x,x')-N_{c}\mf{D}(x,w')-\mf{F}(x,x',z)+\mf{F}(x,z,w')\notag\\&\qquad[\big[U^{ab'}(x')-U^{ab'}(w')\big]\big[-f^{adc}f^{fbe}U^{db}(x)U^{ce}(z)+N_{c}U^{af}(x)\big]\rho^{b'}(x')\rho^{f}(y);\ \ \rho\text{ at }x',\,y]\notag\\
\mathcal{S}_{\rm I.V}^{2b}&=-\mf{D}(x,x')+\mf{D}(x,w')+\mf{D}(y,x')-\mf{D}(y,w')\notag\\&\qquad[\big[U^{ab'}(x')-U^{ab'}(w')\big]\big[U^{af}(y)-U^{af}(x)\big]\rho^{b'}(x')\rho^{f}(y);\ \ \rho\text{ at }x',\,y]\notag\\
\mathcal{S}_{\rm II.I}^{(2)}&=N_{c}\mf{D}(x,x')+N_{c}\mf{D}(w,w')-\mf{F}(x,z,w')-\mf{F}(x',z,w)\notag\\&\qquad[\text{colour bracket of } N_{2}\text{ (Group II Row I, second part)};\ \ \rho\text{ at }x',\,x]\notag\\
\mathcal{S}_{\rm LO}&=\mf{D}(x,x')-\mf{D}(x,y')-\mf{D}(y,x')+\mf{D}(y,y')\notag\\&\qquad[\big[U^{ce'}(y')-U^{ce'}(x')\big]\big[U^{ce}(y)-U^{ce}(x)\big]\rho^{e'}(x')\rho^{e}(x);\ \ \rho\text{ at }x',\,x]\notag\\
\mathcal{S}_{\mathcal{X}}&=N_{c}^{2}-1+\mf{D}(x,x')-\mf{D}(x,z)-\mf{D}(x',z)\notag\\&\qquad[\big[U^{ce'}(z)-U^{ce'}(x')\big]\big[U^{ce}(z)-U^{ce}(x)\big]\rho^{e'}(x')\rho^{e}(x);\ \ \rho\text{ at }x',\,x]\notag\\
\mathcal{S}_{\rm III.III}^{(2)}(y')&=-N_{c}\mf{D}(x,x')+N_{c}\mf{D}(x,y')+\mf{F}(x',z,w)-\mf{F}(y',z,w)\notag\\&\qquad[\mathcal{L}_{2}^{ab}(y')\,\mathcal{R}_{2}^{ab};\ \ \rho\text{ at }x',\,x]\notag\\
\mathcal{S}_{\rm III.III}^{(2)}(z)&=-N_{c}\mf{D}(x,x')+N_{c}\mf{D}(x,z)-N_{c}\mf{D}(z,w)+\mf{F}(x',z,w)\notag\\&\qquad[\mathcal{L}_{2}^{ab}(z)\,\mathcal{R}_{2}^{ab};\ \ \rho\text{ at }x',\,x]\notag
\end{align}
with the trace forms $\mathcal{S}_{\rm I.II}={\rm Tr}\{(U^{\dagger}(x')-U^{\dagger}(w'))(U(w)-U(x))\}$, $\mathcal{S}_{\rm I.V}^{2b}={\rm Tr}\{(U^{\dagger}(x')-U^{\dagger}(w'))(U(y)-U(x))\}$, $\mathcal{S}_{\rm LO}={\rm Tr}\{(U^{\dagger}(y')-U^{\dagger}(x'))(U(y)-U(x))\}$, $\mathcal{S}_{\mathcal{X}}={\rm Tr}\{(U^{\dagger}(z)-U^{\dagger}(x'))(U(z)-U(x))\}$, $\mathcal{S}^{2a}_{\rm I.V}=-{\rm Tr}\{T^{e}(U^{\dagger}(x')-U^{\dagger}(w'))U(x)T^{e}U^{\dagger}(x)U(z)\}+N_{c}{\rm Tr}\{(U^{\dagger}(x')-U^{\dagger}(w'))U(x)\}$.\par
\noindent\textbf{Cubic structures.} $\Phi\equiv C^{abc}f^{abc}$ for the first ordering written; type: ord $=\rho(p_{1})\rho(p_{2})\rho(p_{3})$, A23 $=\rho(p_{1})\{\rho(p_{2}),\rho(p_{3})\}$, A12 $=\{\rho(p_{1}),\rho(p_{2})\}\rho(p_{3})$. Each contact is listed as [pair; point removed $\to$ point kept; density; coefficient $c$ in $\langle C\rho\rho\rho\rangle\ni c\,i\,\Phi|\,\delta^{(2)}(\text{pair})\,\langle\rho^{f}\rho^{f}\rangle/(N_{c}^{2}-1)$; $\Phi$ at the contact].
\begin{align}
\Phi_{\rm I.II}^{2b}&=N_{c}\mf{D}(x',w)-N_{c}\mf{D}(w,w')-\mf{F}(x,y,x')+\mf{F}(x,y,w')\notag\\&\qquad[\mathcal{L}^{a}_{x'w'}\big[f^{hbc}U^{ah}(w)\{\rho^{b}(x),\rho^{c}(y)\}-f^{abc}U^{bd}(x)U^{ce}(y)\{\rho^{d}(x),\rho^{e}(y)\}\big];\ \text{A23}\ (x',x,y)]\notag\\
&\qquad\cdot\ \{x',x\},\ x'\to x,\ \langle\rho^{f}(x)\rho^{f}(y)\rangle,\ c=1:\quad \Phi_{\rm I.II}^{2b}\big|=-N_{c}\mf{D}(x,y)+N_{c}\mf{D}(x,w)-N_{c}\mf{D}(w,w')+\mf{F}(x,y,w')\notag\\
&\qquad\cdot\ \{x',y\},\ x'\to y,\ \langle\rho^{f}(y)\rho^{f}(x)\rangle,\ c=-1:\quad \Phi_{\rm I.II}^{2b}\big|=-N_{c}\mf{D}(x,y)+N_{c}\mf{D}(y,w)-N_{c}\mf{D}(w,w')+\mf{F}(x,y,w')\notag
\end{align}
\begin{align}
\Phi_{\rm I.III}(z)&=\mf{F}(x,y,x')-\mf{F}(x,y,w')-\mf{F}(x,x',z)+\mf{F}(x,z,w')\notag\\&\qquad[f^{abd}\,\mathcal{C}_{\rm III}(z);\ \text{ord}\ (x',x,y)]\notag\\
&\qquad\cdot\ \{x',x\},\ x'\to x,\ \langle\rho^{f}(x)\rho^{f}(y)\rangle,\ c=\tfrac{1}{2}:\quad \Phi_{\rm I.III}(z)\big|=N_{c}\mf{D}(x,y)-N_{c}\mf{D}(x,z)-\mf{F}(x,y,w')+\mf{F}(x,z,w')\notag\\
&\qquad\cdot\ \{x,y\},\ y\to x,\ \langle\rho^{f}(x)\rho^{f}(x')\rangle,\ c=\tfrac{1}{2}:\quad \Phi_{\rm I.III}(z)\big|=N_{c}\mf{D}(x,x')-N_{c}\mf{D}(x,w')-\mf{F}(x,x',z)+\mf{F}(x,z,w')\notag\\
&\qquad\cdot\ \{x',y\},\ x'\to y,\ \langle\rho^{f}(y)\rho^{f}(x)\rangle,\ c=-\tfrac{1}{2}:\quad \Phi_{\rm I.III}(z)\big|=N_{c}\mf{D}(x,y)-\mf{F}(x,y,z)-\mf{F}(x,y,w')+\mf{F}(x,z,w')\notag
\end{align}
\begin{align}
\Phi_{\rm I.III}(w)&=\mf{F}(x,y,x')-\mf{F}(x,y,w')-\mf{F}(x,x',w)+\mf{F}(x,w,w')\notag\\&\qquad[f^{abd}\,\mathcal{C}_{\rm III}(w);\ \text{ord}\ (x',x,y)]\notag\\
&\qquad\cdot\ \{x',x\},\ x'\to x,\ \langle\rho^{f}(x)\rho^{f}(y)\rangle,\ c=\tfrac{1}{2}:\quad \Phi_{\rm I.III}(w)\big|=N_{c}\mf{D}(x,y)-N_{c}\mf{D}(x,w)-\mf{F}(x,y,w')+\mf{F}(x,w,w')\notag\\
&\qquad\cdot\ \{x,y\},\ y\to x,\ \langle\rho^{f}(x)\rho^{f}(x')\rangle,\ c=\tfrac{1}{2}:\quad \Phi_{\rm I.III}(w)\big|=N_{c}\mf{D}(x,x')-N_{c}\mf{D}(x,w')-\mf{F}(x,x',w)+\mf{F}(x,w,w')\notag\\
&\qquad\cdot\ \{x',y\},\ x'\to y,\ \langle\rho^{f}(y)\rho^{f}(x)\rangle,\ c=-\tfrac{1}{2}:\quad \Phi_{\rm I.III}(w)\big|=N_{c}\mf{D}(x,y)-\mf{F}(x,y,w)-\mf{F}(x,y,w')+\mf{F}(x,w,w')\notag
\end{align}
\begin{align}
\Phi_{\rm I.V}^{(1)}&=-\mf{F}(x,x',z)+\mf{F}(x,z,w')+\mf{F}(y,x',y')-\mf{F}(y,y',w')\notag\\&\qquad[\mathcal{L}^{a}_{x'w'}\big[f^{adc}U^{db}(x)U^{ce}(z)\{\rho^{e}(y),\rho^{b}(y')\}-f^{abc}U^{cd}(y)U^{be}(y')\{\rho^{d}(y),\rho^{e}(y')\}\big];\ \text{A23}\ (x',y,y')]\notag\\
&\qquad\cdot\ \{x',y\},\ x'\to y,\ \langle\rho^{f}(y)\rho^{f}(y')\rangle,\ c=1:\quad \Phi_{\rm I.V}^{(1)}\big|=N_{c}\mf{D}(y,y')-\mf{F}(x,y,z)+\mf{F}(x,z,w')-\mf{F}(y,y',w')\notag\\
&\qquad\cdot\ \{x',y'\},\ x'\to y',\ \langle\rho^{f}(y')\rho^{f}(y)\rangle,\ c=-1:\quad \Phi_{\rm I.V}^{(1)}\big|=N_{c}\mf{D}(y,y')-\mf{F}(x,y',z)+\mf{F}(x,z,w')-\mf{F}(y,y',w')\notag
\end{align}
\begin{align}
\Phi_{\rm II.I}^{1a}&=N_{c}\mf{D}(w,w')+\mf{F}(x,y,x')-\mf{F}(x,z,w')-\mf{Q}(y,z;w,x')\notag\\&\qquad[\text{colour factor of }N_{1a};\ \text{A12}\ (x',y,x)]\notag\\
&\qquad\cdot\ \{y,x\},\ y\to x,\ \langle\rho^{f}(x)\rho^{f}(x')\rangle,\ c=1:\quad \Phi_{\rm II.I}^{1a}\big|=N_{c}\mf{D}(x,x')+N_{c}\mf{D}(w,w')-\mf{F}(x,z,w')-\mf{Q}(x,z;w,x')\notag\\
&\qquad\cdot\ \{x',x\},\ x'\to x,\ \langle\rho^{f}(x)\rho^{f}(y)\rangle,\ c=-1:\quad \Phi_{\rm II.I}^{1a}\big|=N_{c}\mf{D}(x,y)+N_{c}\mf{D}(w,w')-\mf{F}(x,z,w')-\mf{Q}(x,w;z,y)\notag
\end{align}
\begin{align}
\Phi_{\rm II.I}^{1b}&=N_{c}\mf{D}(w,w')+\mf{F}(x,y,x')-\mf{F}(x,z,w)-\mf{Q}(y,z;w',x')\notag\\&\qquad[\text{colour factor of }N_{1b};\ \text{A23}\ (x,x',y)]\notag\\
&\qquad\cdot\ \{x,x'\},\ x'\to x,\ \langle\rho^{f}(x)\rho^{f}(y)\rangle,\ c=1:\quad \Phi_{\rm II.I}^{1b}\big|=N_{c}\mf{D}(x,y)+N_{c}\mf{D}(w,w')-\mf{F}(x,z,w)-\mf{Q}(x,w';z,y)\notag\\
&\qquad\cdot\ \{x,y\},\ y\to x,\ \langle\rho^{f}(x)\rho^{f}(x')\rangle,\ c=-1:\quad \Phi_{\rm II.I}^{1b}\big|=N_{c}\mf{D}(x,x')+N_{c}\mf{D}(w,w')-\mf{F}(x,z,w)-\mf{Q}(x,z;w',x')\notag
\end{align}
\begin{align}
\Phi_{\rm III.II}&=N_{c}\mf{D}(x',w)+\mf{F}(x,y,x')-\mf{F}(x,x',z)-\mf{Q}(y,z;w,x')\notag\\&\qquad[\text{colour factor of Group III Row II};\ \text{ord}\ (y,x',x)]\notag\\
&\qquad\cdot\ \{y,x'\},\ x'\to y,\ \langle\rho^{f}(y)\rho^{f}(x)\rangle,\ c=\tfrac{1}{2}:\quad \Phi_{\rm III.II}\big|=N_{c}\mf{D}(x,y)+N_{c}\mf{D}(y,w)-\mf{F}(x,y,z)-\mf{F}(y,z,w)\notag\\
&\qquad\cdot\ \{x',x\},\ x'\to x,\ \langle\rho^{f}(x)\rho^{f}(y)\rangle,\ c=\tfrac{1}{2}:\quad \Phi_{\rm III.II}\big|=N_{c}\mf{D}(x,y)-N_{c}\mf{D}(x,z)+N_{c}\mf{D}(x,w)-\mf{Q}(x,w;z,y)\notag\\
&\qquad\cdot\ \{y,x\},\ y\to x,\ \langle\rho^{f}(x)\rho^{f}(x')\rangle,\ c=-\tfrac{1}{2}:\quad \Phi_{\rm III.II}\big|=N_{c}\mf{D}(x,x')+N_{c}\mf{D}(x',w)-\mf{F}(x,x',z)-\mf{Q}(x,z;w,x')\notag
\end{align}
\begin{align}
\Phi_{\rm III.III}^{(1)}(y')&=-\mf{F}(x,y,x')+\mf{F}(x,y,y')+\mf{F}(x',z,w)-\mf{F}(y',z,w)\notag\\&\qquad[f^{c'd'a}\mathcal{L}^{c'd'b}(y')\mathcal{R}_{1}^{ab};\ \text{A23}\ (x',y,x)]\notag\\
&\qquad\cdot\ \{x',y\},\ x'\to y,\ \langle\rho^{f}(y)\rho^{f}(x)\rangle,\ c=1:\quad \Phi_{\rm III.III}^{(1)}(y')\big|=-N_{c}\mf{D}(x,y)+\mf{F}(x,y,y')+\mf{F}(y,z,w)-\mf{F}(y',z,w)\notag\\
&\qquad\cdot\ \{x',x\},\ x'\to x,\ \langle\rho^{f}(x)\rho^{f}(y)\rangle,\ c=-1:\quad \Phi_{\rm III.III}^{(1)}(y')\big|=-N_{c}\mf{D}(x,y)+\mf{F}(x,y,y')+\mf{F}(x,z,w)-\mf{F}(y',z,w)\notag
\end{align}
\begin{align}
\Phi_{\rm III.III}^{(1)}(z)&=-N_{c}\mf{D}(z,w)-\mf{F}(x,y,x')+\mf{F}(x,y,z)+\mf{F}(x',z,w)\notag\\&\qquad[f^{c'd'a}\mathcal{L}^{c'd'b}(z)\mathcal{R}_{1}^{ab};\ \text{A23}\ (x',y,x)]\notag\\
&\qquad\cdot\ \{x',y\},\ x'\to y,\ \langle\rho^{f}(y)\rho^{f}(x)\rangle,\ c=1:\quad \Phi_{\rm III.III}^{(1)}(z)\big|=-N_{c}\mf{D}(x,y)-N_{c}\mf{D}(z,w)+\mf{F}(x,y,z)+\mf{F}(y,z,w)\notag\\
&\qquad\cdot\ \{x',x\},\ x'\to x,\ \langle\rho^{f}(x)\rho^{f}(y)\rangle,\ c=-1:\quad \Phi_{\rm III.III}^{(1)}(z)\big|=-N_{c}\mf{D}(x,y)-N_{c}\mf{D}(z,w)+\mf{F}(x,y,z)+\mf{F}(x,z,w)\notag
\end{align}
\begin{align}
\Phi_{\rm III.V}(y')&=-\mf{F}(x,y,x')+\mf{F}(x,y,y')+\mf{F}(x,x',w)-\mf{F}(x,y',w)\notag\\&\qquad[f^{c'ca}\mathcal{L}_{\rm V}^{c'}(y')\mathcal{R}_{\rm V}^{ac};\ \text{ord}\ (x',x,y)]\notag\\
&\qquad\cdot\ \{x',x\},\ x'\to x,\ \langle\rho^{f}(x)\rho^{f}(y)\rangle,\ c=\tfrac{1}{2}:\quad \Phi_{\rm III.V}(y')\big|=-N_{c}\mf{D}(x,y)+N_{c}\mf{D}(x,w)+\mf{F}(x,y,y')-\mf{F}(x,y',w)\notag\\
&\qquad\cdot\ \{x,y\},\ y\to x,\ \langle\rho^{f}(x)\rho^{f}(x')\rangle,\ c=\tfrac{1}{2}:\quad \Phi_{\rm III.V}(y')\big|=-N_{c}\mf{D}(x,x')+N_{c}\mf{D}(x,y')+\mf{F}(x,x',w)-\mf{F}(x,y',w)\notag\\
&\qquad\cdot\ \{x',y\},\ x'\to y,\ \langle\rho^{f}(y)\rho^{f}(x)\rangle,\ c=-\tfrac{1}{2}:\quad \Phi_{\rm III.V}(y')\big|=-N_{c}\mf{D}(x,y)+\mf{F}(x,y,y')+\mf{F}(x,y,w)-\mf{F}(x,y',w)\notag
\end{align}
\begin{align}
\Phi_{\rm III.V}(z)&=-\mf{F}(x,y,x')+\mf{F}(x,y,z)+\mf{F}(x,x',w)-\mf{F}(x,z,w)\notag\\&\qquad[f^{c'ca}\mathcal{L}_{\rm V}^{c'}(z)\mathcal{R}_{\rm V}^{ac};\ \text{ord}\ (x',x,y)]\notag\\
&\qquad\cdot\ \{x',x\},\ x'\to x,\ \langle\rho^{f}(x)\rho^{f}(y)\rangle,\ c=\tfrac{1}{2}:\quad \Phi_{\rm III.V}(z)\big|=-N_{c}\mf{D}(x,y)+N_{c}\mf{D}(x,w)+\mf{F}(x,y,z)-\mf{F}(x,z,w)\notag\\
&\qquad\cdot\ \{x,y\},\ y\to x,\ \langle\rho^{f}(x)\rho^{f}(x')\rangle,\ c=\tfrac{1}{2}:\quad \Phi_{\rm III.V}(z)\big|=-N_{c}\mf{D}(x,x')+N_{c}\mf{D}(x,z)+\mf{F}(x,x',w)-\mf{F}(x,z,w)\notag\\
&\qquad\cdot\ \{x',y\},\ x'\to y,\ \langle\rho^{f}(y)\rho^{f}(x)\rangle,\ c=-\tfrac{1}{2}:\quad \Phi_{\rm III.V}(z)\big|=-N_{c}\mf{D}(x,y)+\mf{F}(x,y,z)+\mf{F}(x,y,w)-\mf{F}(x,z,w)\notag
\end{align}
\begin{align}
\Phi_{\rm III.VI}(y')&=\mf{F}(x,y,x')-\mf{F}(x,y,y')-\mf{F}(y,x',z)+\mf{F}(y,y',z)\notag\\&\qquad[f^{c'd'a}\mathcal{L}^{c'd'b}(y')\mathcal{R}_{\rm VI}^{ab};\ \text{ord}\ (x',y,x)]\notag\\
&\qquad\cdot\ \{x',y\},\ x'\to y,\ \langle\rho^{f}(y)\rho^{f}(x)\rangle,\ c=\tfrac{1}{2}:\quad \Phi_{\rm III.VI}(y')\big|=N_{c}\mf{D}(x,y)-N_{c}\mf{D}(y,z)-\mf{F}(x,y,y')+\mf{F}(y,y',z)\notag\\
&\qquad\cdot\ \{y,x\},\ y\to x,\ \langle\rho^{f}(x)\rho^{f}(x')\rangle,\ c=\tfrac{1}{2}:\quad \Phi_{\rm III.VI}(y')\big|=N_{c}\mf{D}(x,x')-N_{c}\mf{D}(x,y')-\mf{F}(x,x',z)+\mf{F}(x,y',z)\notag\\
&\qquad\cdot\ \{x',x\},\ x'\to x,\ \langle\rho^{f}(x)\rho^{f}(y)\rangle,\ c=-\tfrac{1}{2}:\quad \Phi_{\rm III.VI}(y')\big|=N_{c}\mf{D}(x,y)-\mf{F}(x,y,y')-\mf{F}(x,y,z)+\mf{F}(y,y',z)\notag
\end{align}
\begin{align}
\Phi_{\rm III.VI}(z)&=N_{c}\mf{D}(y,z)+\mf{F}(x,y,x')-\mf{F}(x,y,z)-\mf{F}(y,x',z)\notag\\&\qquad[f^{c'd'a}\mathcal{L}^{c'd'b}(z)\mathcal{R}_{\rm VI}^{ab};\ \text{ord}\ (x',y,x)]\notag\\
&\qquad\cdot\ \{x',y\},\ x'\to y,\ \langle\rho^{f}(y)\rho^{f}(x)\rangle,\ c=\tfrac{1}{2}:\quad \Phi_{\rm III.VI}(z)\big|=N_{c}\mf{D}(x,y)-\mf{F}(x,y,z)\notag\\
&\qquad\cdot\ \{y,x\},\ y\to x,\ \langle\rho^{f}(x)\rho^{f}(x')\rangle,\ c=\tfrac{1}{2}:\quad \Phi_{\rm III.VI}(z)\big|=N_{c}\mf{D}(x,x')-\mf{F}(x,x',z)\notag\\
&\qquad\cdot\ \{x',x\},\ x'\to x,\ \langle\rho^{f}(x)\rho^{f}(y)\rangle,\ c=-\tfrac{1}{2}:\quad \Phi_{\rm III.VI}(z)\big|=N_{c}\mf{D}(x,y)+N_{c}\mf{D}(y,z)-2\mf{F}(x,y,z)\notag
\end{align}

\subsubsection{Group I (with c.c.)}
\noindent Kernels (notation of Group I of the finite-term section: $K^{i}(v)=v^{i}/v^{2}$, $K(a,b,z)=K(a-z)\cdot K(b-z)$, $\mathcal{E}$, $\mathcal{J}$, $\mathcal{K}^{ij}$, $\mathcal{A}_{0}$, $\mathcal{R}^{i}$, $\mathcal{F}^{im}$ as defined there):
\begin{align}
\mathfrak{K}^{\rm II}_{2b}&\equiv\int d^{2}p\,d^{2}q\;e^{iq\cdot(w-y)}\,e^{ip\cdot(y-x)}\,K^{i}(x'-w')\,\frac{1}{\mathbf{q}^{2}}\,\mathcal{R}^{i}(\mathbf{q},\mathbf{p})\,,\notag\\
\mathfrak{K}^{\rm III}_{z}&\equiv K^{i}(x'-w')\,K^{j}(y-z)\,\mathcal{K}^{ij}-\frac12\log\left[\frac{(z-w)^{2}}{(x-w)^{2}}\right]K(y,w,z)\,K(x'-w')\cdot K(x-w)\,,\notag\\
\mathfrak{K}^{\rm III}_{w}&\equiv-\mathcal{E}\big(k\cdot(z-w)\big)\,K(x,w,z)\,K(x'-w')\cdot K(y-w)\,,\notag\\
\mathfrak{K}^{\rm V}_{1}&\equiv\mathcal{J}(s)\,K(x'-w')\cdot K(z-x)\,K(y-z)\cdot K(y'-x)-\mathcal{E}(s)\,K(x'-w')\cdot K(y-z)\,K(z-x)\cdot K(y'-x)\notag\\&\qquad-e^{-is}\mathcal{E}(-s)\,K(x'-w')\cdot K(y'-x)\,K(z-x)\cdot K(y-z)\,,\notag\\
\mathfrak{K}^{\rm V}_{2a}&\equiv K^{i}(x'-w')\,K^{m}(z-x)\,\mathcal{F}^{im}\,,\notag\\
\mathfrak{K}^{\rm V}_{2b}&\equiv K^{i}(x'-w')\,K^{m}(z-x)\,\mathcal{F}^{im}+\frac{11}{6}\,e^{-is}\,K(x'-w')\cdot K(y-x)\,\frac{K(y-z)\cdot(y-x)}{(x-z)^{2}}\,,\qquad s\equiv k\cdot(x-z)\,.\notag
\end{align}

\paragraph*{Row II}

\begin{align}
\frac{d^{3}N_{1}^{\rm II}}{d^{3}k}\bigg|_{\rm singlet}&=\frac{1}{(2\pi)^{3}}\,\frac{g^{4}}{16\pi^{5}}\,\frac{1}{k^{+}}\,\frac{1}{N_{c}^{2}-1}\,\int_{w,w',x,x'}e^{-ik\cdot(w-w')}\Bigg\{\notag\\
&\qquad \frac{1}{2}\,\pi N_{c}\Big\{-\frac{67}{9}+\frac{\pi^{2}}{3}\Big\}\,K(x'-w')\cdot K(x-w)\,\Big[-\mf{D}(x,x')+\mf{D}(x,w')+\mf{D}(x',w)-\mf{D}(w,w')\Big]\,\langle\rho^{f}(x')\rho^{f}(x)\rangle\Bigg\}\;+\;{\rm c.c.}\label{eq:s-I-II-1}
\end{align}

The c.c.\ equals the term itself (relabel $w\leftrightarrow w'$, $x\leftrightarrow x'$; $\mathcal{S}_{\rm I.II}$ is invariant), so with c.c.\ the factor is $\pi N_{c}\{-\frac{67}{9}+\frac{\pi^{2}}{3}\}$; since $\mathcal{S}_{\rm I.II}=-\mathcal{S}_{\rm LO}$ after the relabelling $y\to w$, $y'\to w'$, this is $\frac{\alpha_{s}N_{c}}{2\pi}\big(\frac{67}{9}-\frac{\pi^{2}}{3}\big)$ times the leading order in singlet form, Eq.~\eqref{eq:s-LO}. \emph{Quadratic.}

\begin{align}
\frac{d^{3}N_{2b}^{\rm II}}{d^{3}k}\bigg|_{\rm singlet}&=\frac{1}{(2\pi)^{3}}\,\frac{g^{4}}{16\pi^{5}}\,\frac{1}{k^{+}}\,\frac{1}{N_{c}^{2}-1}\,\int_{w,w'}e^{-ik\cdot(w-w')}\int_{x,y}\Bigg\{\notag\\
&\qquad -\frac{i}{4}\,\frac{1}{\pi}\,\mathfrak{K}^{\rm II}_{2b}\big|_{x'\to x}\,\Big[-N_{c}\mf{D}(x,y)+N_{c}\mf{D}(x,w)-N_{c}\mf{D}(w,w')+\mf{F}(x,y,w')\Big]\,\langle\rho^{f}(x)\rho^{f}(y)\rangle\notag\\
&\qquad +\frac{i}{4}\,\frac{1}{\pi}\,\mathfrak{K}^{\rm II}_{2b}\big|_{x'\to y}\,\Big[-N_{c}\mf{D}(x,y)+N_{c}\mf{D}(y,w)-N_{c}\mf{D}(w,w')+\mf{F}(x,y,w')\Big]\,\langle\rho^{f}(y)\rho^{f}(x)\rangle\Bigg\}\;+\;{\rm c.c.}\label{eq:s-I-II-2b}
\end{align}

\emph{Cubic} ($\mathcal{L}^{a}_{x'w'}$ times a two-$\rho$ anticommutator; type A23): only the two contact terms are left; they are quadratic, at $x'=x$ and $x'=y$. The kernel is antisymmetric under $x\leftrightarrow y$ ($\mathbf{p}\to\mathbf{q}-\mathbf{p}$, $\mathcal{R}^{i}\to-\mathcal{R}^{i}$) and $\Phi^{2b}_{\rm I.II}|_{x'\to x}$ goes into $\Phi^{2b}_{\rm I.II}|_{x'\to y}$: the two contact terms are equal.

\paragraph*{Row III}

\begin{align}
\frac{d^{3}N^{\rm III}}{d^{3}k}\bigg|_{\rm singlet}&=\frac{1}{(2\pi)^{3}}\,\frac{g^{4}}{16\pi^{5}}\,\frac{1}{k^{+}}\,\frac{1}{N_{c}^{2}-1}\,\int_{w,w'}e^{-ik\cdot(w-w')}\int_{x,y,z,x'}\Bigg\{\notag\\
&\qquad -\,\mathfrak{K}^{\rm III}_{z}\big|_{x'\to x}\,\Big[N_{c}\mf{D}(x,y)-N_{c}\mf{D}(x,z)-\mf{F}(x,y,w')+\mf{F}(x,z,w')\Big]\,\langle\rho^{f}(x)\rho^{f}(y)\rangle\notag\\
&\qquad -\,\mathfrak{K}^{\rm III}_{z}\big|_{y\to x}\,\Big[N_{c}\mf{D}(x,x')-N_{c}\mf{D}(x,w')-\mf{F}(x,x',z)+\mf{F}(x,z,w')\Big]\,\langle\rho^{f}(x)\rho^{f}(x')\rangle\notag\\
&\qquad +\,\mathfrak{K}^{\rm III}_{z}\big|_{x'\to y}\,\Big[N_{c}\mf{D}(x,y)-\mf{F}(x,y,z)-\mf{F}(x,y,w')+\mf{F}(x,z,w')\Big]\,\langle\rho^{f}(y)\rho^{f}(x)\rangle\notag\\
&\qquad -\,\mathfrak{K}^{\rm III}_{w}\big|_{x'\to x}\,\Big[N_{c}\mf{D}(x,y)-N_{c}\mf{D}(x,w)-\mf{F}(x,y,w')+\mf{F}(x,w,w')\Big]\,\langle\rho^{f}(x)\rho^{f}(y)\rangle\notag\\
&\qquad -\,\mathfrak{K}^{\rm III}_{w}\big|_{y\to x}\,\Big[N_{c}\mf{D}(x,x')-N_{c}\mf{D}(x,w')-\mf{F}(x,x',w)+\mf{F}(x,w,w')\Big]\,\langle\rho^{f}(x)\rho^{f}(x')\rangle\notag\\
&\qquad +\,\mathfrak{K}^{\rm III}_{w}\big|_{x'\to y}\,\Big[N_{c}\mf{D}(x,y)-\mf{F}(x,y,w)-\mf{F}(x,y,w')+\mf{F}(x,w,w')\Big]\,\langle\rho^{f}(y)\rho^{f}(x)\rangle\Bigg\}\;+\;{\rm c.c.}\label{eq:s-I-III}
\end{align}

\emph{Cubic} (type ord, $\rho(x')\rho(x)\rho(y)$; colour structures $f^{abd}\mathcal{C}_{\rm III}(z)$ with the kernel $\mathfrak{K}^{\rm III}_{z}$, $f^{abd}\mathcal{C}_{\rm III}(w)$ with $\mathfrak{K}^{\rm III}_{w}$). The symmetrised part is zero under (A2); the three contact terms per structure ($x'=x$, $y=x$, $x'=y$) remain. If the kernels of this soft-remainder row are changed, the colour factors $\Phi_{\rm I.III}(z)|$, $\Phi_{\rm I.III}(w)|$ and the coefficients $-1,-1,+1$ stay as they are, as long as the colour structures $f^{abd}\mathcal{C}_{\rm III}(z)$, $f^{abd}\mathcal{C}_{\rm III}(w)$ are kept.

\paragraph*{Row V}

\begin{align}
\frac{d^{3}N_{1}^{\rm V}}{d^{3}k}\bigg|_{\rm singlet}&=\frac{1}{(2\pi)^{3}}\,\frac{g^{4}}{16\pi^{5}}\,\frac{1}{k^{+}}\,\frac{1}{N_{c}^{2}-1}\,\int_{x,z,y,y',x',w'}e^{-ik\cdot(z-w')}\Bigg\{\notag\\
&\qquad -\frac{1}{2}\,\mathfrak{K}^{\rm V}_{1}\big|_{x'\to y}\,\Big[N_{c}\mf{D}(y,y')-\mf{F}(x,y,z)+\mf{F}(x,z,w')-\mf{F}(y,y',w')\Big]\,\langle\rho^{f}(y)\rho^{f}(y')\rangle\notag\\
&\qquad +\frac{1}{2}\,\mathfrak{K}^{\rm V}_{1}\big|_{x'\to y'}\,\Big[N_{c}\mf{D}(y,y')-\mf{F}(x,y',z)+\mf{F}(x,z,w')-\mf{F}(y,y',w')\Big]\,\langle\rho^{f}(y')\rho^{f}(y)\rangle\Bigg\}\;+\;{\rm c.c.}\label{eq:s-I-V-1}
\end{align}

\emph{Cubic} (type A23, $\rho(x')\{\rho(y),\rho(y')\}$): contact terms at $x'=y$ and $x'=y'$ only.

\begin{align}
\frac{d^{3}N_{2a+2b}^{\rm V}}{d^{3}k}\bigg|_{\rm singlet}&=\frac{1}{(2\pi)^{3}}\,\frac{g^{4}}{16\pi^{5}}\,\frac{1}{k^{+}}\,\frac{1}{N_{c}^{2}-1}\,\int_{x,y,z,x',w'}e^{-ik\cdot(z-w')}\Bigg\{\notag\\
&\qquad \,\mathfrak{K}^{\rm V}_{2a}\,\Big[N_{c}\mf{D}(x,x')-N_{c}\mf{D}(x,w')-\mf{F}(x,x',z)+\mf{F}(x,z,w')\Big]\,\langle\rho^{f}(x')\rho^{f}(y)\rangle\notag\\
&\qquad +\,N_{c}\,\mathfrak{K}^{\rm V}_{2b}\,\Big[-\mf{D}(x,x')+\mf{D}(x,w')+\mf{D}(y,x')-\mf{D}(y,w')\Big]\,\langle\rho^{f}(x')\rho^{f}(y)\rangle\Bigg\}\;+\;{\rm c.c.}\label{eq:s-I-V-2}
\end{align}

\emph{Quadratic.} $N_{2b\,\rm UV}$ of Row V is cancelled by Row II (finite-term section) and does not appear. Group I has no DGLAP and no running-coupling term.

\subsubsection{Group II (no c.c.)}
\noindent Kernels: $\mathfrak{K}_{\mathcal{T}}(w,w')\equiv\mathcal{K}^{i}(x'-w')\,\mathcal{K}^{j}(y-z)\,\log\left[\frac{(x-z)^{2}}{(x-w)^{2}}\right]\mathcal{T}^{ij}(x,z,w)$ and $\mathfrak{K}_{\mathcal{T}}(w',w)$ the same with $w\leftrightarrow w'$ (the kernel of $N_{1b}$); $\mathfrak{K}_{2}\equiv\mathfrak{K}_{2a}+\dots+\mathfrak{K}_{2f}$ of the finite-term section (to be integrated over $z$ together).
\paragraph*{Row I, first part}

\begin{align}
\frac{d^{3}N_{1a+1b}}{d^{3}k}\bigg|_{\rm singlet}&=\frac{1}{(2\pi)^{3}}\,\frac{g^{4}}{16\pi^{5}}\,\frac{1}{k^{+}}\,\frac{1}{N_{c}^{2}-1}\,\int_{x,y,z,w,x',w'}e^{-ik\cdot(w-w')}\Bigg\{\notag\\
&\qquad +\,\mathfrak{K}_{\mathcal{T}}(w,w')\big|_{y\to x}\,\Big[N_{c}\mf{D}(x,x')+N_{c}\mf{D}(w,w')-\mf{F}(x,z,w')-\mf{Q}(x,z;w,x')\Big]\,\langle\rho^{f}(x)\rho^{f}(x')\rangle\notag\\
&\qquad -\,\mathfrak{K}_{\mathcal{T}}(w,w')\big|_{x'\to x}\,\Big[N_{c}\mf{D}(x,y)+N_{c}\mf{D}(w,w')-\mf{F}(x,z,w')-\mf{Q}(x,w;z,y)\Big]\,\langle\rho^{f}(x)\rho^{f}(y)\rangle\notag\\
&\qquad -\,\mathfrak{K}_{\mathcal{T}}(w',w)\big|_{x'\to x}\,\Big[N_{c}\mf{D}(x,y)+N_{c}\mf{D}(w,w')-\mf{F}(x,z,w)-\mf{Q}(x,w';z,y)\Big]\,\langle\rho^{f}(x)\rho^{f}(y)\rangle\notag\\
&\qquad +\,\mathfrak{K}_{\mathcal{T}}(w',w)\big|_{y\to x}\,\Big[N_{c}\mf{D}(x,x')+N_{c}\mf{D}(w,w')-\mf{F}(x,z,w)-\mf{Q}(x,z;w',x')\Big]\,\langle\rho^{f}(x)\rho^{f}(x')\rangle\Bigg\}\label{eq:s-II-I-1}
\end{align}

\emph{Cubic} ($\{\rho(x'),\rho(y)\}\rho(x)$ in $N_{1a}$, type A12; $\rho(x)\{\rho(x'),\rho(y)\}$ in $N_{1b}$, type A23). $\Phi^{1b}_{\rm II.I}$ is $\Phi^{1a}_{\rm II.I}$ with $w\leftrightarrow w'$. After $w\leftrightarrow w'$ in the $N_{1b}$ terms the four contact terms combine into
\begin{align}
2\cos\big(k\cdot(w-w')\big)\big[\mathfrak{K}_{\mathcal{T}}(w,w')|_{y\to x}\,\Phi^{1a}_{\rm II.I}|_{y\to x}\langle\rho^{f}(x)\rho^{f}(x')\rangle-\mathfrak{K}_{\mathcal{T}}(w,w')|_{x'\to x}\,\Phi^{1a}_{\rm II.I}|_{x'\to x}\langle\rho^{f}(x)\rho^{f}(y)\rangle\big]\,,\notag
\end{align}
real and even in $k$, as it must be for Group II.
\paragraph*{Row I, second part}

\begin{align}
\frac{d^{3}N_{2}}{d^{3}k}\bigg|_{\rm singlet}&=\frac{1}{(2\pi)^{3}}\,\frac{g^{4}}{16\pi^{5}}\,\frac{1}{k^{+}}\,\frac{1}{N_{c}^{2}-1}\,\int_{z,w,w',x,x'}e^{-ik\cdot(w-w')}\Bigg\{\notag\\
&\qquad 2\,\mathfrak{K}_{2}\,\Big[N_{c}\mf{D}(x,x')+N_{c}\mf{D}(w,w')-\mf{F}(x,z,w')-\mf{F}(x',z,w)\Big]\,\langle\rho^{f}(x')\rho^{f}(x)\rangle\Bigg\}\label{eq:s-II-I-2}
\end{align}

\emph{Quadratic.} $\mathcal{S}^{(2)}_{\rm II.I}$ is symmetric under $(x,w)\leftrightarrow(x',w')$.
\paragraph*{Row IV: DGLAP, running coupling, finite term, remnant}
\noindent With the colour factor $\mathcal{S}_{\rm LO}={\rm Tr}\{(U^{\dagger}(y')-U^{\dagger}(x'))(U(y)-U(x))\}$ (the leading-order one) define
\begin{align}
\mathcal{F}_{\rm s}(\xi)&\equiv\frac{1}{N_{c}^{2}-1}\,\frac{1}{\xi^{2}}\int_{x,y,x',y'}e^{-ik\cdot(y-y')\frac{1}{\xi}}\Big[2\gamma_{E}+\log\frac{(y'-y)^{2}\mu^{2}_{\overline{MS}}}{4\,\xi^{2}}\Big]\,\mathcal{K}(x-y)\cdot\mathcal{K}(x'-y')\;\mathcal{S}_{\rm LO}\;\langle\rho^{f}(x')\rho^{f}(x)\rangle\,,\notag\\
\mathcal{X}_{\rm s}&\equiv\frac{1}{N_{c}^{2}-1}\int_{x,x',z}\mathcal{K}(x-z)\cdot\mathcal{K}(x'-z)\;{\rm Tr}\{(U^{\dagger}(z)-U^{\dagger}(x'))(U(z)-U(x))\}\;\langle\rho^{f}(x')\rho^{f}(x)\rangle\,,\qquad\mathcal{F}_{\rm s}(0)=-\frac{4\pi}{k^{2}}\,\mathcal{X}_{\rm s}\,.\notag
\end{align}
(The square bracket is that of the DGLAP section; if its form is changed, e.g.\ the $\xi^{2}$ in the logarithm, the colour factor stays $\mathcal{S}_{\rm LO}$.) Then, with $\mathcal{C}=\frac{1}{(2\pi)^{3}}\frac{g^{4}}{16\pi^{4}}\frac{1}{k^{+}}=\pi\times\frac{1}{(2\pi)^{3}}\frac{g^{4}}{16\pi^{5}}\frac{1}{k^{+}}$,
\begin{align}
\frac{d^{3}N_{\rm DGLAP}}{d^{3}k}\bigg|_{\rm singlet}&=-\frac{1}{(2\pi)^{3}}\,\frac{g^{4}}{16\pi^{5}}\,\frac{1}{k^{+}}\;\pi\int_{0}^{1}d\xi\;\widetilde{\mathcal P}_{gg}(\xi)\,\mathcal{F}_{\rm s}(\xi)\notag\\
&=-\frac{1}{(2\pi)^{3}}\,\frac{g^{4}}{16\pi^{5}}\,\frac{1}{k^{+}}\;\pi\Bigg\{2N_{c}\int_{0}^{1}d\xi\Big[\xi(1-\xi)\mathcal{F}_{\rm s}(\xi)+\frac{\xi\mathcal{F}_{\rm s}(\xi)-\mathcal{F}_{\rm s}(1)}{1-\xi}\notag\\&\qquad+\frac{1-\xi}{\xi}\big(\mathcal{F}_{\rm s}(\xi)-\mathcal{F}_{\rm s}(0)\big)\Big]+\frac{\beta_{0}}{2}\,\mathcal{F}_{\rm s}(1)\Bigg\}\,,\label{eq:s-DGLAP}\\
\frac{d^{3}N_{\rm RC}}{d^{3}k}\bigg|_{\rm singlet}&=+\frac{1}{(2\pi)^{3}}\,\frac{g^{4}}{16\pi^{5}}\,\frac{1}{k^{+}}\;\pi\,\frac{\beta_{0}}{2}\,\mathcal{F}_{\rm s}(1)\,,\label{eq:s-RC}\\
\frac{d^{3}N_{\rm Finite}}{d^{3}k}\bigg|_{\rm singlet}&=\frac{1}{(2\pi)^{3}}\,\frac{g^{4}}{16\pi^{5}}\,\frac{1}{k^{+}}\,\frac{1}{N_{c}^{2}-1}\,2\pi N_{c}\int_{0}^{1}d\xi\,\frac{1-\xi}{\xi}\int_{x,y,x',y'}e^{-ik\cdot(y-y')\frac{1}{\xi}}\notag\\&\qquad\times\bigg[\mathcal{K}(x-y)\cdot\mathcal{K}(x'-y')-2\,\frac{\mathcal{K}(x-y)\cdot(y'-y)\;\mathcal{K}(x'-y')\cdot(y'-y)}{(y'-y)^{2}}\bigg]\,\mathcal{S}_{\rm LO}\,\langle\rho^{f}(x')\rho^{f}(x)\rangle\,,\label{eq:s-fin}\\
\frac{d^{3}N_{\rm rem}}{d^{3}k}\bigg|_{\rm singlet}&=-\frac{1}{(2\pi)^{3}}\,\frac{g^{4}}{16\pi^{5}}\,\frac{1}{k^{+}}\;\frac{8\pi^{2}N_{c}}{k^{2}}\;\mathcal{X}_{\rm s}\notag\\&=-\frac{1}{(2\pi)^{3}}\,\frac{g^{4}N_{c}}{2\pi^{3}}\,\frac{1}{k^{2}}\,\frac{1}{k^{+}}\,\mathcal{X}_{\rm s}\,.\label{eq:s-rem}
\end{align}
The remnant \eqref{eq:s-rem} is the $-\frac{1}{k^{+}}\times$(BFKL structure) term that goes with the reading $\int_{0}^{1}d\xi\big[\frac{1-\xi}{\xi}\big]_{+}\mathcal{F}=\int_{0}^{1}d\xi\,\frac{1-\xi}{\xi}\big(\mathcal{F}(\xi)-\mathcal{F}(0)\big)$ used in \eqref{eq:s-DGLAP}; with the plus prescription acting on $(1-\xi)\mathcal{F}(\xi)$ it is contained in the DGLAP term (the two forms are identical). The $\beta_{0}$ terms of \eqref{eq:s-DGLAP} and \eqref{eq:s-RC} cancel. The leading order in the same form, and its sum with \eqref{eq:s-RC}:
\begin{align}
\frac{d^{3}N_{\rm LO}}{d^{3}k}\bigg|_{\rm singlet}&=\frac{1}{(2\pi)^{3}}\,\frac{g^{2}}{2\pi^{2}}\,\frac{1}{k^{+}}\,\frac{1}{N_{c}^{2}-1}\int_{x,y,x',y'}e^{-ik\cdot(y-y')}\notag\\&\qquad\times\mathcal{K}(x-y)\cdot\mathcal{K}(x'-y')\,{\rm Tr}\{(U^{\dagger}(y')-U^{\dagger}(x'))(U(y)-U(x))\}\,\langle\rho^{f}(x')\rho^{f}(x)\rangle\,,\label{eq:s-LO}\\
\frac{d^{3}N_{\rm LO+RC}}{d^{3}k}\bigg|_{\rm singlet}&=\frac{1}{(2\pi)^{3}}\,\frac{g^{2}(\mu)}{2\pi^{2}}\,\frac{1}{k^{+}}\,\frac{1}{N_{c}^{2}-1}\int_{x,y,x',y'}e^{-ik\cdot(y-y')}\notag\\&\qquad\times\Big[1+\frac{\alpha_{s}(\mu)\beta_{0}}{4\pi}\log\frac{(y'-y)^{2}\mu^{2}_{\overline{MS}}e^{2\gamma_{E}}}{4}\Big]\mathcal{K}(x-y)\cdot\mathcal{K}(x'-y')\,\mathcal{S}_{\rm LO}\,\langle\rho^{f}(x')\rho^{f}(x)\rangle\,.\notag
\end{align}
In the BFKL labelling ($x'\to y$, $y\to w$, $y'\to w'$): $\mathcal{S}_{\rm LO}={\rm Tr}\{(U^{\dagger}(x)-U^{\dagger}(w))(U(y)-U(w'))\}$ with $\langle\rho^{f}(x)\rho^{f}(y)\rangle$, as in the BFKL and JIMWLK sections. \emph{All of Row IV is quadratic.}

\subsubsection{Group III (with c.c.)}
\noindent Kernels (notation of Groups II, III of the finite-term section: $\mathcal{K}^{i}(r)=r^{i}/r^{2}$, $\mathcal{T}^{ij}$, $\widehat{\mathcal{E}}$, $\Psi$, $D_{p}$, $\widehat{\mathcal{J}}$, $\mathcal{V}$, $\mathcal{W}^{k'm}$; $s\equiv k\cdot(z-y')$):
\begin{align}
\mathfrak{K}_{\widehat{\mathcal{E}}}&\equiv\widehat{\mathcal{E}}(s)\Big[\mathcal{K}(y'-z)\cdot\mathcal{K}(x'-y')\;\mathcal{K}(y-w)\cdot\mathcal{K}(x-z)-\mathcal{K}(y'-z)\cdot\mathcal{K}(x-z)\;\mathcal{K}(x'-y')\cdot\mathcal{K}(y-w)\Big]\,,\notag\\
\mathfrak{K}_{\Psi}&\equiv\Psi(s)\,\mathcal{K}(y'-z)\cdot\mathcal{K}(x-z)\;\mathcal{K}(x'-y')\cdot\mathcal{K}(y-w)\,,\notag\\
\mathfrak{K}_{\delta}^{(y')}&\equiv\Big[\pi\delta(s)-i{\rm P}\frac{1}{s}\Big]\mathcal{K}(y'-z)\cdot\mathcal{K}(y-w)\;\mathcal{K}^{j}(x-z)\,\mathcal{K}^{j}(x'-y')\,,\notag\\
\mathfrak{K}_{\delta}^{(z)}&\equiv-\Big[\pi\delta(s)-i{\rm P}\frac{1}{s}\Big]\mathcal{K}(y'-z)\cdot\mathcal{K}(y-w)\;\mathcal{K}^{j}(x-z)\,\mathcal{K}^{j}(x'-z)\,,\notag\\
\mathfrak{K}_{\rm reg}&\equiv e^{-is\frac{p^{+}}{k^{+}}}\,\frac{1}{D_{p}}\,\{\dots\}_{N_{2}^{\rm reg}}\ \text{(the brace of $N_{2}^{\rm reg}$)},\qquad
\mathfrak{K}_{\Psi,2}\equiv\Psi(s)\,\frac{\widehat{\mathcal{J}}}{(x-w)^{2}}\,,\notag\\
\mathfrak{K}_{\rm sub}^{(y')}&\equiv e^{-is\frac{p^{+}}{k^{+}}}\,\frac{p^{+}}{D_{p}}\,\mathcal{K}^{m}(y'-z)\,\mathcal{W}^{k'm}\,\mathcal{K}^{k'}(x'-y')\,,\qquad
\mathfrak{K}_{\rm sub}^{(z)}\equiv-e^{-is\frac{p^{+}}{k^{+}}}\,\frac{p^{+}}{D_{p}}\,\mathcal{K}^{m}(y'-z)\,\mathcal{W}^{k'm}\,\mathcal{K}^{k'}(x'-z)\,,\notag\\
\mathfrak{K}_{\log}&\equiv\frac{(x-w)^{2}}{(x-z)^{2}\left[(x-w)^{2}-(x-z)^{2}\right]}\,\log\left[\frac{(x-w)^{2}}{(x-z)^{2}}\right]\mathcal{K}^{m}(w'-z)\,\mathcal{K}^{k'}(x'-z)\,\mathcal{W}^{k'm}\,.\notag
\end{align}
In the terms with $\mathfrak{K}_{\delta}^{(y')}$ and $\mathfrak{K}_{\delta}^{(z)}$ (and with $\mathfrak{K}_{\rm sub}^{(y')}$, $\mathfrak{K}_{\rm sub}^{(z)}$) the two kernels belong to one subtraction and stay under one $y'$ integral, done before the $p^{+}$ integral, as in the finite-term section; in singlet form the pairs of lines with the same contact stay together.
\paragraph*{Row II}

\begin{align}
\frac{d^{3}N^{\rm III}_{\rm II}}{d^{3}k}\bigg|_{\rm singlet}&=\frac{1}{(2\pi)^{3}}\,\frac{g^{4}}{16\pi^{5}}\,\frac{1}{k^{+}}\,\frac{1}{N_{c}^{2}-1}\,\int_{x,y,x',z,w,w'}e^{-ik\cdot(w-w')}\Bigg\{\notag\\
&\qquad +\,\mathfrak{K}_{\mathcal{T}}(w,w')\big|_{x'\to y}\,\Big[N_{c}\mf{D}(x,y)+N_{c}\mf{D}(y,w)-\mf{F}(x,y,z)-\mf{F}(y,z,w)\Big]\,\langle\rho^{f}(y)\rho^{f}(x)\rangle\notag\\
&\qquad +\,\mathfrak{K}_{\mathcal{T}}(w,w')\big|_{x'\to x}\,\Big[N_{c}\mf{D}(x,y)-N_{c}\mf{D}(x,z)+N_{c}\mf{D}(x,w)-\mf{Q}(x,w;z,y)\Big]\,\langle\rho^{f}(x)\rho^{f}(y)\rangle\notag\\
&\qquad -\,\mathfrak{K}_{\mathcal{T}}(w,w')\big|_{y\to x}\,\Big[N_{c}\mf{D}(x,x')+N_{c}\mf{D}(x',w)-\mf{F}(x,x',z)-\mf{Q}(x,z;w,x')\Big]\,\langle\rho^{f}(x)\rho^{f}(x')\rangle\Bigg\}\;+\;{\rm c.c.}\label{eq:s-III-II}
\end{align}

\emph{Cubic} (type ord, $\rho(y)\rho(x')\rho(x)$); three contact terms.
\paragraph*{Row III, first part}

\begin{align}
\frac{d^{3}N^{\rm III}_{{\rm III},1}}{d^{3}k}\bigg|_{\rm singlet}&=\frac{1}{(2\pi)^{3}}\,\frac{g^{4}}{16\pi^{5}}\,\frac{1}{k^{+}}\,\frac{1}{N_{c}^{2}-1}\,\int_{w,x',y',z,y,x}e^{-ik\cdot(w-y')}\Bigg\{\notag\\
&\qquad -\,\mathfrak{K}_{\widehat{\mathcal{E}}}\big|_{x'\to y}\,\Big[-N_{c}\mf{D}(x,y)+\mf{F}(x,y,y')+\mf{F}(y,z,w)-\mf{F}(y',z,w)\Big]\,\langle\rho^{f}(y)\rho^{f}(x)\rangle\notag\\
&\qquad +\,\mathfrak{K}_{\widehat{\mathcal{E}}}\big|_{x'\to x}\,\Big[-N_{c}\mf{D}(x,y)+\mf{F}(x,y,y')+\mf{F}(x,z,w)-\mf{F}(y',z,w)\Big]\,\langle\rho^{f}(x)\rho^{f}(y)\rangle\notag\\
&\qquad +\,\mathfrak{K}_{\Psi}\big|_{x'\to y}\,\Big[-N_{c}\mf{D}(x,y)+\mf{F}(x,y,y')+\mf{F}(y,z,w)-\mf{F}(y',z,w)\Big]\,\langle\rho^{f}(y)\rho^{f}(x)\rangle\notag\\
&\qquad -\,\mathfrak{K}_{\Psi}\big|_{x'\to x}\,\Big[-N_{c}\mf{D}(x,y)+\mf{F}(x,y,y')+\mf{F}(x,z,w)-\mf{F}(y',z,w)\Big]\,\langle\rho^{f}(x)\rho^{f}(y)\rangle\notag\\
&\qquad +\,\mathfrak{K}_{\delta}^{(y')}\big|_{x'\to y}\,\Big[-N_{c}\mf{D}(x,y)+\mf{F}(x,y,y')+\mf{F}(y,z,w)-\mf{F}(y',z,w)\Big]\,\langle\rho^{f}(y)\rho^{f}(x)\rangle\notag\\
&\qquad -\,\mathfrak{K}_{\delta}^{(y')}\big|_{x'\to x}\,\Big[-N_{c}\mf{D}(x,y)+\mf{F}(x,y,y')+\mf{F}(x,z,w)-\mf{F}(y',z,w)\Big]\,\langle\rho^{f}(x)\rho^{f}(y)\rangle\notag\\
&\qquad +\,\mathfrak{K}_{\delta}^{(z)}\big|_{x'\to y}\,\Big[-N_{c}\mf{D}(x,y)-N_{c}\mf{D}(z,w)+\mf{F}(x,y,z)+\mf{F}(y,z,w)\Big]\,\langle\rho^{f}(y)\rho^{f}(x)\rangle\notag\\
&\qquad -\,\mathfrak{K}_{\delta}^{(z)}\big|_{x'\to x}\,\Big[-N_{c}\mf{D}(x,y)-N_{c}\mf{D}(z,w)+\mf{F}(x,y,z)+\mf{F}(x,z,w)\Big]\,\langle\rho^{f}(x)\rho^{f}(y)\rangle\Bigg\}\;+\;{\rm c.c.}\label{eq:s-III-III-1}
\end{align}

\emph{Cubic} (type A23, $\rho(x')\{\rho(y),\rho(x)\}$): contact terms at $x'=y$ and $x'=x$. Colour structures $f^{c'd'a}\mathcal{L}^{c'd'b}(y')\mathcal{R}^{ab}_{1}$ (kernels $\mathfrak{K}_{\widehat{\mathcal{E}}}$, $\mathfrak{K}_{\Psi}$, $\mathfrak{K}_{\delta}^{(y')}$) and $f^{c'd'a}\mathcal{L}^{c'd'b}(z)\mathcal{R}^{ab}_{1}$ (kernel $\mathfrak{K}_{\delta}^{(z)}$).
\paragraph*{Row III, second part}

\begin{align}
\frac{d^{3}N^{\rm III}_{{\rm III},2}}{d^{3}k}\bigg|_{\rm singlet}&=\frac{1}{(2\pi)^{3}}\,\frac{g^{4}}{16\pi^{5}}\,\frac{1}{k^{+}}\,\frac{1}{N_{c}^{2}-1}\,\int_{w,x',y',z,x}e^{-ik\cdot(w-y')}\Bigg\{\notag\\
&\qquad 2\,k^{+}\!\int_{0}^{\infty}\!dp^{+}\,\mathfrak{K}_{\rm reg}\,\Big[-N_{c}\mf{D}(x,x')+N_{c}\mf{D}(x,y')+\mf{F}(x',z,w)-\mf{F}(y',z,w)\Big]\,\langle\rho^{f}(x')\rho^{f}(x)\rangle\notag\\
&\qquad -2\,\mathfrak{K}_{\Psi,2}\,\Big[-N_{c}\mf{D}(x,x')+N_{c}\mf{D}(x,y')+\mf{F}(x',z,w)-\mf{F}(y',z,w)\Big]\,\langle\rho^{f}(x')\rho^{f}(x)\rangle\notag\\
&\qquad -2\,\frac{1}{k^{+}}\!\int_{0}^{\infty}\!dp^{+}\,\mathfrak{K}_{\rm sub}^{(y')}\,\Big[-N_{c}\mf{D}(x,x')+N_{c}\mf{D}(x,y')+\mf{F}(x',z,w)-\mf{F}(y',z,w)\Big]\,\langle\rho^{f}(x')\rho^{f}(x)\rangle\notag\\
&\qquad -2\,\frac{1}{k^{+}}\!\int_{0}^{\infty}\!dp^{+}\,\mathfrak{K}_{\rm sub}^{(z)}\,\Big[-N_{c}\mf{D}(x,x')+N_{c}\mf{D}(x,z)-N_{c}\mf{D}(z,w)+\mf{F}(x',z,w)\Big]\,\langle\rho^{f}(x')\rho^{f}(x)\rangle\Bigg\}\;+\;{\rm c.c.}\label{eq:s-III-III-2}
\end{align}

\begin{align}
\frac{d^{3}N^{\rm III}_{{\rm III},2\log}}{d^{3}k}\bigg|_{\rm singlet}&=\frac{1}{(2\pi)^{3}}\,\frac{g^{4}}{16\pi^{5}}\,\frac{1}{k^{+}}\,\frac{1}{N_{c}^{2}-1}\,\int_{w,x',w',z,x}e^{-ik\cdot(w-w')}\Bigg\{\notag\\
&\qquad 2\,\mathfrak{K}_{\log}\,\Big[-N_{c}\mf{D}(x,x')+N_{c}\mf{D}(x,z)-N_{c}\mf{D}(z,w)+\mf{F}(x',z,w)\Big]\,\langle\rho^{f}(x')\rho^{f}(x)\rangle\Bigg\}\;+\;{\rm c.c.}\label{eq:s-III-III-2log}
\end{align}

\emph{Quadratic} ($\mathcal{L}^{ab}_{2}\mathcal{R}^{ab}_{2}$ has one $\rho$ in each factor). [The prefactor $\frac{g^{4}}{4\pi^{4}}\int_{0}^{\infty}\frac{dp^{+}}{2\pi}$ of $N_{2}^{\rm reg}$ has no $\frac{1}{k^{+}}$: it is $2k^{+}\int_{0}^{\infty}dp^{+}$ times $\frac{g^{4}}{16\pi^{5}}\frac{1}{k^{+}}$; that of $N_{2}^{\rm sub}$ is $-\frac{2}{k^{+}}\int_{0}^{\infty}dp^{+}$ times $\frac{g^{4}}{16\pi^{5}}\frac{1}{k^{+}}$.]
\paragraph*{Row V}

\begin{align}
\frac{d^{3}N^{\rm III}_{\rm V}}{d^{3}k}\bigg|_{\rm singlet}&=\frac{1}{(2\pi)^{3}}\,\frac{g^{4}}{16\pi^{5}}\,\frac{1}{k^{+}}\,\frac{1}{N_{c}^{2}-1}\,\int_{w,y',z,y,x,x'}e^{-ik\cdot(w-y')}\Bigg\{\notag\\
&\qquad -\,\mathfrak{K}_{\widehat{\mathcal{E}}}\big|_{x'\to x}\,\Big[-N_{c}\mf{D}(x,y)+N_{c}\mf{D}(x,w)+\mf{F}(x,y,y')-\mf{F}(x,y',w)\Big]\,\langle\rho^{f}(x)\rho^{f}(y)\rangle\notag\\
&\qquad -\,\mathfrak{K}_{\widehat{\mathcal{E}}}\big|_{y\to x}\,\Big[-N_{c}\mf{D}(x,x')+N_{c}\mf{D}(x,y')+\mf{F}(x,x',w)-\mf{F}(x,y',w)\Big]\,\langle\rho^{f}(x)\rho^{f}(x')\rangle\notag\\
&\qquad +\,\mathfrak{K}_{\widehat{\mathcal{E}}}\big|_{x'\to y}\,\Big[-N_{c}\mf{D}(x,y)+\mf{F}(x,y,y')+\mf{F}(x,y,w)-\mf{F}(x,y',w)\Big]\,\langle\rho^{f}(y)\rho^{f}(x)\rangle\notag\\
&\qquad +\,\mathfrak{K}_{\Psi}\big|_{x'\to x}\,\Big[-N_{c}\mf{D}(x,y)+N_{c}\mf{D}(x,w)+\mf{F}(x,y,y')-\mf{F}(x,y',w)\Big]\,\langle\rho^{f}(x)\rho^{f}(y)\rangle\notag\\
&\qquad +\,\mathfrak{K}_{\Psi}\big|_{y\to x}\,\Big[-N_{c}\mf{D}(x,x')+N_{c}\mf{D}(x,y')+\mf{F}(x,x',w)-\mf{F}(x,y',w)\Big]\,\langle\rho^{f}(x)\rho^{f}(x')\rangle\notag\\
&\qquad -\,\mathfrak{K}_{\Psi}\big|_{x'\to y}\,\Big[-N_{c}\mf{D}(x,y)+\mf{F}(x,y,y')+\mf{F}(x,y,w)-\mf{F}(x,y',w)\Big]\,\langle\rho^{f}(y)\rho^{f}(x)\rangle\notag\\
&\qquad +\,\mathfrak{K}_{\delta}^{(y')}\big|_{x'\to x}\,\Big[-N_{c}\mf{D}(x,y)+N_{c}\mf{D}(x,w)+\mf{F}(x,y,y')-\mf{F}(x,y',w)\Big]\,\langle\rho^{f}(x)\rho^{f}(y)\rangle\notag\\
&\qquad +\,\mathfrak{K}_{\delta}^{(y')}\big|_{y\to x}\,\Big[-N_{c}\mf{D}(x,x')+N_{c}\mf{D}(x,y')+\mf{F}(x,x',w)-\mf{F}(x,y',w)\Big]\,\langle\rho^{f}(x)\rho^{f}(x')\rangle\notag\\
&\qquad -\,\mathfrak{K}_{\delta}^{(y')}\big|_{x'\to y}\,\Big[-N_{c}\mf{D}(x,y)+\mf{F}(x,y,y')+\mf{F}(x,y,w)-\mf{F}(x,y',w)\Big]\,\langle\rho^{f}(y)\rho^{f}(x)\rangle\notag\\
&\qquad +\,\mathfrak{K}_{\delta}^{(z)}\big|_{x'\to x}\,\Big[-N_{c}\mf{D}(x,y)+N_{c}\mf{D}(x,w)+\mf{F}(x,y,z)-\mf{F}(x,z,w)\Big]\,\langle\rho^{f}(x)\rho^{f}(y)\rangle\notag\\
&\qquad +\,\mathfrak{K}_{\delta}^{(z)}\big|_{y\to x}\,\Big[-N_{c}\mf{D}(x,x')+N_{c}\mf{D}(x,z)+\mf{F}(x,x',w)-\mf{F}(x,z,w)\Big]\,\langle\rho^{f}(x)\rho^{f}(x')\rangle\notag\\
&\qquad -\,\mathfrak{K}_{\delta}^{(z)}\big|_{x'\to y}\,\Big[-N_{c}\mf{D}(x,y)+\mf{F}(x,y,z)+\mf{F}(x,y,w)-\mf{F}(x,z,w)\Big]\,\langle\rho^{f}(y)\rho^{f}(x)\rangle\Bigg\}\;+\;{\rm c.c.}\label{eq:s-III-V}
\end{align}

\emph{Cubic} (type ord, $\rho(x')\rho(x)\rho(y)$; colour structures $f^{c'ca}\mathcal{L}^{c'}_{\rm V}(y')\mathcal{R}^{ac}_{\rm V}$ and, with $\mathfrak{K}_{\delta}^{(z)}$, $f^{c'ca}\mathcal{L}^{c'}_{\rm V}(z)\mathcal{R}^{ac}_{\rm V}$).
\paragraph*{Row VI}

\begin{align}
\frac{d^{3}N^{\rm III}_{\rm VI}}{d^{3}k}\bigg|_{\rm singlet}&=\frac{1}{(2\pi)^{3}}\,\frac{g^{4}}{16\pi^{5}}\,\frac{1}{k^{+}}\,\frac{1}{N_{c}^{2}-1}\,\int_{w,y',z,y,x,x'}e^{-ik\cdot(w-y')}\Bigg\{\notag\\
&\qquad -\,\mathfrak{K}_{\widehat{\mathcal{E}}}\big|_{x'\to y}\,\Big[N_{c}\mf{D}(x,y)-N_{c}\mf{D}(y,z)-\mf{F}(x,y,y')+\mf{F}(y,y',z)\Big]\,\langle\rho^{f}(y)\rho^{f}(x)\rangle\notag\\
&\qquad -\,\mathfrak{K}_{\widehat{\mathcal{E}}}\big|_{y\to x}\,\Big[N_{c}\mf{D}(x,x')-N_{c}\mf{D}(x,y')-\mf{F}(x,x',z)+\mf{F}(x,y',z)\Big]\,\langle\rho^{f}(x)\rho^{f}(x')\rangle\notag\\
&\qquad +\,\mathfrak{K}_{\widehat{\mathcal{E}}}\big|_{x'\to x}\,\Big[N_{c}\mf{D}(x,y)-\mf{F}(x,y,y')-\mf{F}(x,y,z)+\mf{F}(y,y',z)\Big]\,\langle\rho^{f}(x)\rho^{f}(y)\rangle\notag\\
&\qquad +\,\mathfrak{K}_{\Psi}\big|_{x'\to y}\,\Big[N_{c}\mf{D}(x,y)-N_{c}\mf{D}(y,z)-\mf{F}(x,y,y')+\mf{F}(y,y',z)\Big]\,\langle\rho^{f}(y)\rho^{f}(x)\rangle\notag\\
&\qquad +\,\mathfrak{K}_{\Psi}\big|_{y\to x}\,\Big[N_{c}\mf{D}(x,x')-N_{c}\mf{D}(x,y')-\mf{F}(x,x',z)+\mf{F}(x,y',z)\Big]\,\langle\rho^{f}(x)\rho^{f}(x')\rangle\notag\\
&\qquad -\,\mathfrak{K}_{\Psi}\big|_{x'\to x}\,\Big[N_{c}\mf{D}(x,y)-\mf{F}(x,y,y')-\mf{F}(x,y,z)+\mf{F}(y,y',z)\Big]\,\langle\rho^{f}(x)\rho^{f}(y)\rangle\notag\\
&\qquad +\,\mathfrak{K}_{\delta}^{(y')}\big|_{x'\to y}\,\Big[N_{c}\mf{D}(x,y)-N_{c}\mf{D}(y,z)-\mf{F}(x,y,y')+\mf{F}(y,y',z)\Big]\,\langle\rho^{f}(y)\rho^{f}(x)\rangle\notag\\
&\qquad +\,\mathfrak{K}_{\delta}^{(y')}\big|_{y\to x}\,\Big[N_{c}\mf{D}(x,x')-N_{c}\mf{D}(x,y')-\mf{F}(x,x',z)+\mf{F}(x,y',z)\Big]\,\langle\rho^{f}(x)\rho^{f}(x')\rangle\notag\\
&\qquad -\,\mathfrak{K}_{\delta}^{(y')}\big|_{x'\to x}\,\Big[N_{c}\mf{D}(x,y)-\mf{F}(x,y,y')-\mf{F}(x,y,z)+\mf{F}(y,y',z)\Big]\,\langle\rho^{f}(x)\rho^{f}(y)\rangle\notag\\
&\qquad +\,\mathfrak{K}_{\delta}^{(z)}\big|_{x'\to y}\,\Big[N_{c}\mf{D}(x,y)-\mf{F}(x,y,z)\Big]\,\langle\rho^{f}(y)\rho^{f}(x)\rangle\notag\\
&\qquad +\,\mathfrak{K}_{\delta}^{(z)}\big|_{y\to x}\,\Big[N_{c}\mf{D}(x,x')-\mf{F}(x,x',z)\Big]\,\langle\rho^{f}(x)\rho^{f}(x')\rangle\notag\\
&\qquad -\,\mathfrak{K}_{\delta}^{(z)}\big|_{x'\to x}\,\Big[N_{c}\mf{D}(x,y)+N_{c}\mf{D}(y,z)-2\mf{F}(x,y,z)\Big]\,\langle\rho^{f}(x)\rho^{f}(y)\rangle\Bigg\}\;+\;{\rm c.c.}\label{eq:s-III-VI}
\end{align}

\emph{Cubic} (type ord, $\rho(x')\rho(y)\rho(x)$; colour structures $f^{c'd'a}\mathcal{L}^{c'd'b}(y')\mathcal{R}^{ab}_{\rm VI}$ and, with $\mathfrak{K}_{\delta}^{(z)}$, $f^{c'd'a}\mathcal{L}^{c'd'b}(z)\mathcal{R}^{ab}_{\rm VI}$).

\subsubsection{Which rows are cubic, which quadratic}
\begin{center}
\renewcommand{\arraystretch}{1.3}
\small\begin{tabular}{|l|l|l|l|}
\hline
row, term & $\rho$ content & after the average (A1), (A2) & singlet form \\ \hline
I.II, $N_{1}$ ($-\frac{67}{9}+\frac{\pi^{2}}{3}$) & quadratic & $\mathcal{S}_{\rm I.II}$ & \eqref{eq:s-I-II-1} \\
I.II, $N_{2b}$ & cubic, A23 & 2 contact terms ($x'{=}x$, $x'{=}y$) & \eqref{eq:s-I-II-2b} \\
I.III (all three lines) & cubic, ord & 3 contact terms per structure ($z$, $w$) & \eqref{eq:s-I-III} \\
I.V, $N_{1}$ & cubic, A23 & 2 contact terms ($x'{=}y$, $x'{=}y'$) & \eqref{eq:s-I-V-1} \\
I.V, $N_{2a}$, $N_{2b\,\rm rest}$ & quadratic & $\mathcal{S}^{2a}_{\rm I.V}$, $\mathcal{S}^{2b}_{\rm I.V}$ & \eqref{eq:s-I-V-2} \\
II.I, first part ($N_{1a}$, $N_{1b}$) & cubic, A12 / A23 & 2+2 contact terms & \eqref{eq:s-II-I-1} \\
II.I, second part ($\mathfrak{K}_{2a}\dots\mathfrak{K}_{2f}$) & quadratic & $\mathcal{S}^{(2)}_{\rm II.I}$ & \eqref{eq:s-II-I-2} \\
II.IV: DGLAP, RC, finite, remnant & quadratic & $\mathcal{S}_{\rm LO}$, $\mathcal{S}_{\mathcal{X}}$ & \eqref{eq:s-DGLAP}--\eqref{eq:s-rem} \\
III.II & cubic, ord & 3 contact terms & \eqref{eq:s-III-II} \\
III.III, first part & cubic, A23 & 2 contact terms per structure ($y'$, $z$) & \eqref{eq:s-III-III-1} \\
III.III, second part & quadratic & $\mathcal{S}^{(2)}_{\rm III.III}(y')$, $\mathcal{S}^{(2)}_{\rm III.III}(z)$ & \eqref{eq:s-III-III-2}, \eqref{eq:s-III-III-2log} \\
III.V & cubic, ord & 3 contact terms per structure ($y'$, $z$) & \eqref{eq:s-III-V} \\
III.VI & cubic, ord & 3 contact terms per structure ($y'$, $z$) & \eqref{eq:s-III-VI} \\ \hline
\end{tabular}
\end{center}
For every cubic row the symmetrised part ($S^{abc}$) is zero under (A2) and would be odd in $k$ otherwise; the contact terms are even in $k$, real after adding the c.c., and remain for a Gaussian, colour-neutral projectile. For the soft-remainder rows (I.III, I.V, III.III, III.V, III.VI) the split of the large-$z$ tails by degree in the Wilson line at the soft position can be read off the singlet factors: the contact points are $\rho$ points, never the soft position $z$, so after $U^{T}U=1$ has been used (as in the tail tables of the finite-term section) the terms of $\Phi|$ and $\mathcal{S}$ that contain $z$ are the part linear in $U(z)$ and the others are the $z$-independent part.

\subsubsection{Numerical checks}\label{sec:s-checks}
\noindent Scripts in \texttt{scratchpad/fix/singlet/} (\texttt{structs.py}, \texttt{sym.py}, \texttt{check\_final.py}, \texttt{check\_extra.py}, \texttt{check\_contact\_kernels.py}); explicit $SU(2)$ and $SU(3)$ generators, random adjoint Wilson lines $U^{ab}=2{\rm Tr}(t^{a}Vt^{b}V^{\dagger})$.
\begin{itemize}
\item \emph{Colour algebra.} The colour structures were parsed from the TeX source of the finite-term section and evaluated as tensors; the singlet factor of every quadratic structure ($\delta^{ab}$ contraction) and $\Phi$ with the contact coefficients of every cubic structure agree with the expressions printed above: largest relative deviation $9\times10^{-15}$ (21 structures, $SU(2)$ and $SU(3)$, 4 sets of Wilson lines each); the reduced forms at the contact points (identities $\mf{F}(a,a,c)=N_{c}\mf{D}(a,c)$ etc.) $9\times10^{-14}$; the identities themselves and the trace forms of $\mf{F}$, $\mathcal{S}_{\rm I.II}$, $\mathcal{S}^{2a,2b}_{\rm I.V}$, $\mathcal{S}_{\rm LO}$, $\mathcal{S}_{\mathcal{X}}$: $\le2\times10^{-15}$.
\item \emph{Operator level.} $\rho^{a}(x)$ = generators of colour charges on $P$ lattice sites (explicit matrices on the tensor-product space), projectile = globally colour-invariant density matrix $e^{-H}/{\rm Tr}\,e^{-H}$, $H=\sum_{p<q}\beta_{pq}\rho_{p}\cdot\rho_{q}$ with random $\beta$. For every structure, all assignments of its $\rho$ points to sites (with random complex weights in place of the kernel; coincident sites = contact), the exact ordered expectation value ${\rm Tr}(\varrho\,C\rho\rho[\rho])$ was compared with [quadratic: $\mathcal{S}\,n/(N_{c}^{2}-1)$; cubic: the $S^{abc}$ part plus the printed contact terms, with the reduced $\Phi|$]. Ensembles: $SU(2)$ adjoint ($P=3$), $SU(2)$ fundamental ($P=3$), $SU(2)$ fundamental with a chiral term $0.7f^{abc}\rho^{a}_{1}\rho^{b}_{2}\rho^{c}_{3}$ in $H$, $SU(3)$ adjoint ($P=2$), $SU(3)$ fundamental ($P=3$), the same averaged with its charge conjugate (C-even), and that with a chiral term $0.6f\rho\rho\rho$. Largest relative deviation over all $7\times21$ cases: $5.8\times10^{-15}$. (A1) in these ensembles: $|[Q^{a},\varrho]|\le3\times10^{-17}$, $|\langle\rho\rangle|\le3\times10^{-17}$, $|\langle\rho^{a}\rho^{b}\rangle-\delta^{ab}n/(N_{c}^{2}-1)|\le2\times10^{-16}$.
\item \emph{(A2) and parity.} $S^{abc}$: $d$-part $1.11$ for fundamental $SU(3)$ charges (C-odd), $0$ after averaging with the charge conjugate; $f$-part $7.5\times10^{-2}$ ($SU(2)$, chiral) and $5.1\times10^{-2}$ ($SU(3)$, C-even, chiral), $\le2\times10^{-17}$ in the other ensembles. The $S^{abc}$ part of every cubic structure is real (imaginary part $\le4\times10^{-14}$ relative in the three ensembles where it is non-zero) for real weights, which with $A(-k)=-A(k)^{*}$ gives the oddness in $k$.
\item \emph{Other statements:} $\Phi^{1b}_{\rm II.I}=\Phi^{1a}_{\rm II.I}|_{w\leftrightarrow w'}$ (also at the contacts), $\mathcal{S}_{\rm I.II}$ invariant under $w\leftrightarrow w'$, $x\leftrightarrow x'$, $\mathcal{S}_{\rm I.II}=-\mathcal{S}_{\rm LO}|_{y\to w,y'\to w'}$, $\mathcal{S}^{(2)}_{\rm II.I}$ symmetric under $(x,w)\leftrightarrow(x',w')$: exact (symbolic). $\mathcal{R}^{i}_{0}(-\mathbf{q},-\mathbf{p})=-\mathcal{R}^{i}_{0}(\mathbf{q},\mathbf{p})$. All conversions of the row prefactors to $\frac{1}{(2\pi)^{3}}\frac{g^{4}}{16\pi^{5}}\frac{1}{k^{+}}$: exact (sympy).
\item \emph{Kernels at two contacts} (spot check): $\mathfrak{K}^{\rm III}_{z}|_{y\to x}$ and $\mathfrak{K}_{\mathcal{T}}|_{y\to x}$ at the new coincidence $x\to z$ grow like $\log(1/r)/r$ ($r\max|\mathfrak{K}|=0.73,\,1.10,\,1.47,\,1.84$ and $0.88,\,1.25,\,1.62,\,1.99$ for $r=10^{-2}\dots10^{-5}$): integrable; the $1/r^{2}$ terms cancel between the terms of $\mathcal{K}^{ij}$ resp.\ $\mathcal{T}^{ij}$.
\end{itemize}

\subsubsection{What is not settled}
\begin{itemize}
\item[(1)] \emph{The cubic rows do not vanish.} Contrary to the statement that the cubic terms vanish for a Gaussian projectile, only their symmetrised part does; the reordering commutators leave two-point contact terms (Eqs.~\eqref{eq:s-ord}--\eqref{eq:s-a12}), of the same order as the quadratic finite terms, even in $k$, and not removed by colour neutrality. They are written above. Whether they cancel among the rows (or against the corresponding contact terms of the three-$\rho$ evolution terms) was not checked: that needs the kernels integrated, and the kernels differ from row to row.
\item[(2)] \emph{Integrability at the contact configurations} was checked only for the two kernels above. The other contact kernels put two $\rho$ points together in kernels that were checked for integrability only at separate points; no kernel has a factor singular at the contact itself (checked by inspection: no kernel depends on the difference of two $\rho$ points that are merged), but the new coincidences of the merged point with $z$, $w$, $w'$, $y'$ were not power-counted row by row.
\item[(3)] \emph{Large-distance tails in singlet form.} The tail sums of the finite-term section (OPEN I-2, II/III-2, II/III-3) were computed for the symmetrised two-$\rho$ parts of the rows; for the cubic rows these are the operator form of the contact terms written here, so the cancellations found there, with their restrictions, should carry over to the singlet form. They were not recomputed in singlet form.
\item[(4)] The kernels of the soft-remainder rows are those of Finite\_terms\_final.tex; if they are replaced, the colour factors and coefficients above stay valid as long as the colour structures listed in the table are kept (the coefficients are the row prefactor times $c\,i$).
\item[(5)] (A2) is an assumption on the projectile; a C-odd ($d^{abc}$) or chiral ($f^{abc}g_{f}$) three-point function would add the $S^{abc}$ parts, odd in $k$.
\end{itemize}


% ---- inlined copy of paper/notes/fix_tail.tex (20 Nov 2026 13:45 UTC) ----
\subsection{The joint large-distance tail of Group II Row IV: term-by-term check}\label{sec:tailfix}
\noindent\textbf{Result in one paragraph.} In the region where the measured gluon, the unmeasured gluon and their parent are all far from the projectile together ($w,w',z,y'\to\infty$ at fixed $w-z$, $w'-z$, $y'-z$; ``joint limit''), the finite terms \emph{as written} in \texttt{Final\_terms\_checked} / \texttt{Finite\_terms\_final} carry the tails, in units
\begin{align}
\mathsf{u}\equiv2N_{c}\,\mathcal{C}\,\frac{4\pi}{k^{2}}\int\frac{d^{2}Y}{Y^{2}}\,G_{\infty}(Y),\qquad G_{\infty}(Y)=\int_{x,x'}\big[U(Y)-U(x')\big]^{ce'}\big[U(Y)-U(x)\big]^{ce}\rho^{e'}(x')\rho^{e}(x),\qquad\varphi=\angle(k,Y),
\label{ft:unit}
\end{align}
\begin{center}
\renewcommand{\arraystretch}{1.3}
\begin{tabular}{|l|c|}\hline
Group II Row IV: DGLAP $+$ RC $+$ finite (remnant, tensor term) & $-\frac{11}{6}+\frac16\cos2\varphi$\\
Group II Row I, second part ($N_{2a}+\dots+N_{2f}$, i.e.\ $\sum\mathfrak{K}_{2\cdot}$) & $-\frac{11}{6}+\frac16\cos2\varphi$\\
Group III Row III, second part, with c.c.\ ($N_{2}^{\rm reg}+N_{2}^{\Psi}+N_{2}^{\rm sub}+N_{2}^{\log}$) & $+\frac{11}{3}-\frac13\cos2\varphi$\\
Group II Row I first part; Group III Row III first part & $0$\\\hline
sum & $0$\\\hline
\end{tabular}
\end{center}
and the $\log\frac{\vee}{k^{+}}$ (BFKL) and $\log\frac{k^{+}}{\Lambda}$ (JIMWLK) coefficients of the three rows carry $1+1-2=0$ times $\mathsf{u}$ each. \textbf{The tail cancels; no correction of the written finite terms is needed.} The scaleless rule and the recoil rewritings do move tail pieces between terms, but only inside the finite sector of Group III Row III (the term $N_{2}^{\log}$ and the extra term $N_{2}^{\Psi}$ each carry $+2\mathsf{u}$; the subtracted term $N_{2}^{\rm sub}$ carries none) and between the BFKL list and $N_{2}^{\log}$; the bookkeeping is complete in each sector. This holds for any target that is uniform at large distance ($U\to V$, e.g.\ a finite target, $V=1$), and for a colour-neutral projectile for \emph{any} target; for a coloured projectile and a target that is not uniform at large distance a remainder proportional to the total charge is left (the non-singlet term already listed as OPEN~II/III-4(b)). What changes are \emph{statements} of \texttt{Finite\_terms\_final} (red items in Sec.~\ref{ft:corr}). Scripts in \texttt{fix/tail/}: \texttt{t1\_density.py} (amplitude Fourier transforms, tail densities), \texttt{t2\_colour.py} (colour factors, $SU(2)$, $SU(3)$), \texttt{t3\_sectors.py} (sector integrals, Row~III split), \texttt{t4\_rowI\_written.py} (the written kernels of Row~I at large distance, 60 digits); outputs \texttt{t*.out}. [A]: analytic (sympy); [N]: numerical.

\subsubsection{Setup}
\noindent Phase $e^{-ik\cdot(w-w')}$, $k=(|k|,0)$, $t\equiv p^{+}/k^{+}$, $\xi=\frac{1}{1+t}$; $\mathcal{C}=\frac{1}{(2\pi)^{3}}\frac{g^{4}}{16\pi^{4}}\frac{1}{k^{+}}$ as in Sec.~\ref{sec:poles-dglap}. Joint limit: $x,x'$ fixed, $z=Zn$, $w=z+a$, $w'=z+a'$ (Group II) resp.\ $y'=z+b$ (Group III, conjugate amplitude after the $w'$ integration), $n=(\cos\varphi,\sin\varphi)$, $Z\to\infty$ at fixed $a,a',b$. In this limit every two-gluon amplitude is ${\cal O}(1/Z)$ with a coefficient of the form $c\cdot a/a^{2}$ (the $1/q$ Weizs\"acker--Williams pole of the soft parent, Sec.~\ref{sec:poles-dglap}, Step 14), so each row is $\frac{1}{Z^{2}}\times$(a function of $a,a'$ or $a,b$), whose Fourier transform is elementary: $\int d^{2}a\,e^{-ik\cdot a}\frac{a^{i}}{a^{2}}=-2\pi i\frac{k^{i}}{k^{2}}$, and for the conjugate $C$ amplitude of Group III $\int d^{2}b\,e^{i(1+t)k\cdot b}\frac{b^{m}}{b^{2}}=\frac{2\pi i}{1+t}\frac{k^{m}}{k^{2}}$. $\int d^{2}Z/Z^{2}$ is the logarithmic tail; ${\cal O}(1/Z^{3})$ and terms odd in $n$ at order $1/Z$ do not contribute. The expansion is uniform in $t$ (the scales in $t$ are set by $(x-w)^{2}/(x-z)^{2}\to1$), so the $p^{+}$ integral and the joint limit commute; this is checked on the written Row~I kernels in Step~4.

\noindent\textbf{Colour} [N, \texttt{t2}]. With the far Wilson lines equal, $U(w)=U(w')=U(z)=U(y)=U(y')=V$, the colour factors of the three rows, as written, are exactly $+N_{c}\Gamma$ (Group~II Row~I, all of $N_{2a}\dots N_{2f}$), $+N_{c}\Gamma$ (Group~II Row~IV before the reduction), $-N_{c}\Gamma$ (Group~III Row~III, second part), $\Gamma=[V-U(x')]^{ce'}[V-U(x)]^{ce}\rho^{e'}(x')\rho^{e}(x)$: deviations $\le9\times10^{-15}$ for $|\Gamma|\simeq0.5$--$2.7$, three random configurations each, $SU(2)$ and $SU(3)$.

\subsubsection{Step 1: the pointwise tail density}
\noindent [A, \texttt{t1}] With the amplitudes as written (Group~II Row~I: the bracket of the checked row before the $p^{+}$ integral, prefactor $\frac{g^{4}}{4\pi^{4}}\frac{dp^{+}}{2\pi}\frac{p^{+}k^{+}}{(p^{+}+k^{+})^{2}}\frac{1}{DD'}$; Row~IV: $\mf{C}_{1}\mf{A}^{(1)}$ with $y=z+\xi(w-z)$, prefactor $\frac{p^{+}k^{+}}{(p^{+}+k^{+})^{4}}$; Group~III Row~III, second part: the bracket of \texttt{Final\_terms\_checked} with $\frac{1}{p^{+}}=[\frac{k^{+}}{p^{+}(p^{+}+k^{+})}]_{+}+\frac{1}{p^{+}+k^{+}}$ before the subtraction and the $-\frac{p^{+}}{k^{+2}}$ term before its recoil subtraction), the Fourier coefficients of the $1/Z$ amplitudes of $\mf{B}_{2}$ and of $\mf{C}_{1}\mf{A}^{(1)}$ differ only by the normalisation factor $1+t$ of the two conventions, e.g.\ $(i,j)=(1,1)$: $-2\pi i\frac{(t^{2}+t+1)\cos\varphi}{t(t+1)}$ against $-2\pi i\frac{(t^{2}+t+1)\cos\varphi}{t}$. The three rows have the same tail density per $dt$ (in units $\mathcal{C}\,\Gamma/Z^{2}$, colour stripped, $|k|=k^{+}=1$), $8\pi\,\varrho(t)$,
\begin{align}
\varrho(t)=\frac{(t^{2}+t+1)^{2}+t^{2}\cos2\varphi}{t\,(1+t)^{4}}\,,\qquad \varrho(t)\,dt=\Big[\frac{1-\xi}{\xi}+\frac{\xi}{1-\xi}+\xi(1-\xi)\big(1+\cos2\varphi\big)\Big]d\xi\,,
\label{ft:rho}
\end{align}
so that with the colour factors, and the c.c.\ of Group~III (the density is real and even in $k$), the weights are $+1$, $+1$, $-2$ and the sum vanishes pointwise in $t$. In units $\mathsf{u}$ the density of each row is $(+1,+1,-2)\times\varrho(t)\,dt$. This is the coordinate-space form of the amplitude identity $b-cA=-\frac{V_{\rho g}V_{ggg}}{E_{2}(E_{2}-E_{1})}$.

\subsubsection{Step 2: Group II Row IV (reference)}
\noindent [A, \texttt{t3} [4]] $\int_{1/V}^{1/(1+\lambda)}\varrho\,d\xi=\log V+\log\frac1\lambda-\frac{11}{6}+\frac16\cos2\varphi$ ($V=\vee/k^{+}$, $\lambda=\Lambda/k^{+}$). Written: BFKL $\log V$ (pure log, Final\_terms form), JIMWLK $\log\frac1\lambda$ (coordinate form, tail from $\log(y'-y)^{2}$), DGLAP $+\frac1{12}$, RC $-\frac{11}{12}$, finite: remnant $-1$, tensor $+\frac16\cos2\varphi$. [In the form of \texttt{Finite\_terms\_final}, where $[\frac{1-\xi}{\xi}]_{+}$ acts on $(1-\xi)\mathcal{F}$, the remnant is inside DGLAP: $-\frac{11}{12}$, $-\frac{11}{12}$, $+\frac16\cos2\varphi$; same sum.]

\subsubsection{Step 3: Group II Row I, second part}
\noindent The written finite term is $\lim\big[\int_{\lambda}^{V-1}dt\,I(t)-c_{B}\log V-c_{J}\log\frac1\lambda\big]$, $c_{B}=\frac{1}{BB'}[\text{bracket of }\mathfrak{K}_{2c}]$, $c_{J}=\frac{1}{AA'}[\text{bracket of }\mathfrak{K}_{2b}]$ (Finite\_terms\_final, grey note), i.e.\ $\int_{0}^{1}dt\,(I-\frac{c_{J}}{t})+\int_{1}^{\infty}dt\,(I-\frac{c_{B}}{t})$. In the joint limit $c_{J},c_{B}\to\lim_{t\to0,\infty}t\varrho=1$ and [A, \texttt{t3} [2]]
\begin{align}
\text{tail}\Big(\sum\mathfrak{K}_{2\cdot}\Big)=\int_{0}^{\infty}dt\,\frac{t\cos2\varphi-(2t^{2}+3t+2)}{(1+t)^{4}}\;\mathsf{u}=\Big(-\frac{11}{6}+\frac16\cos2\varphi\Big)\mathsf{u}\,,\qquad \text{BFKL: }+\log V\,\mathsf{u},\quad\text{JIMWLK: }+\log\tfrac1\lambda\,\mathsf{u}.
\label{ft:rowI}
\end{align}
\textbf{Step 4: the written kernels} [N, \texttt{t4}, 60 digits]. $Z^{2}\sum_{2a}^{2f}\mathfrak{K}$, evaluated from the formulas as written (with $d_{\perp}=2$ and the $-1$ of $\mathfrak{K}_{2f}$; $x=(0.3,-0.2)$, $x'=(-0.1,0.4)$, $n$ and $-n$ averaged), against the leading product form $L(a,a')=\int dt\,[\frac{t}{(1+t)^{2}}T_{1}(a)\!\cdot\!T_{1}(a')-\text{subtractions}]$ whose Fourier transform gives \eqref{ft:rowI}: three random $(a,a',n)$, values $-1.99535$, $1.97289$, $-4.28034$; differences $1.3\times10^{-6},\,1.3\times10^{-8},\,1.3\times10^{-10}$ (resp.\ $4.8\times10^{-7\dots-11}$, $1.0\times10^{-6\dots-10}$) at $Z=10^{3},10^{4},10^{5}$. The single $\mathfrak{K}_{2f}$ is ${\cal O}(Z^{-2})$ as well (its ${\cal O}(Z^{0})$ parts, $\frac{d_{\perp}}{4CC'}\times$[$-\frac{n\cdot a'}{n\cdot(a'-a)}$, $-1$, $-\frac{n\cdot a}{n\cdot(a'-a)}$ with signs], cancel). The written form also equals the subtracted $p^{+}$ integral at these large-distance points to all digits shown. Since the $x,x'$ dependence enters only at ${\cal O}(1/Z^{3})$, the check with $x\ne x'$ confirms that the tail is $\propto G_{\infty}$.

\subsubsection{Step 5: Group III Row III, second part, as written}
\noindent The written second part (\texttt{Final\_terms\_checked}) has three terms: T1 (main bracket, plus prescription at $p^{+}\to0$), T2 (the $-\frac{p^{+}}{k^{+2}}$ term with the recoil subtraction $\mathcal{K}(x'-y')\mathcal{L}(y')-\mathcal{K}(x'-z)\mathcal{L}(z)$), T3 (the term with $\mathcal{K}(w'-z)\mathcal{K}(x'-z)\frac{(x-w)^{2}}{(x-z)^{2}(p^{+}+k^{+})D}$). T3 is what is left of the added-back subtraction of T2 after $-\log\frac{\vee}{k^{+}}\frac{1}{k^{+}(x-z)^{2}}\mathcal{W}\,\mathcal{K}\mathcal{K}\,\mathcal{L}(z)$ went to the BFKL list: with $\int_{y'}e^{i(1+t)k\cdot(y'-z)}\mathcal{K}(y'-z)=\frac{1}{1+t}\int_{w'}e^{ik\cdot(w'-z)}\mathcal{K}(w'-z)$ and $-\frac{p^{+}}{k^{+}(p^{+}+k^{+})D}=-\frac{1}{k^{+}B(p^{+}+k^{+})}+\frac{A}{B(p^{+}+k^{+})D}$. Per unit colour weight (multiply by $-2$) [A, \texttt{t3} [1], [3]]:
\begin{center}
\renewcommand{\arraystretch}{1.35}
\begin{tabular}{|l|l|c|c|}\hline
term & density in the joint limit & tail (finite) & $\times(-2)$\\\hline
T1 (main term), author's $[\,]_{+}$ & $\varrho_{A}=\frac{t^{2}\cos2\varphi+2t^{2}+2t+1}{t(1+t)^{4}}$, minus $\frac{1}{t(1+t)}$ & $-\frac56+\frac16\cos2\varphi$ & $\frac53-\frac13\cos2\varphi$\\
\quad $=N_{2}^{\rm reg}$ (phase kept in the subtraction) & minus $\frac{1}{t(1+t)^{2}}$ & $\frac16+\frac16\cos2\varphi$ & $-\frac13-\frac13\cos2\varphi$\\
\quad $+\,N_{2}^{\Psi}=f(0)\Psi$ & $\frac{1}{t(1+t)^{2}}-\frac{1}{t(1+t)}$ & $-1$ & $+2$\\
T2 $=N_{2}^{\rm sub}$ & $\mathcal{K}(x'-y')-\mathcal{K}(x'-z)={\cal O}(b/Z^{2})$ & $0$ & $0$\\
T3 $=N_{2}^{\log}$ & $-\frac{1}{(1+t)^{2}}$ & $-1$ & $+2$\\\hline
BFKL list (moved) & $\frac{1}{1+t}$ over $(0,V-1)$ & $\log V$ & $-2\log V$\\
JIMWLK list & $c_{J}=\lim t\varrho_{A}=1$ & $\log\frac1\lambda$ & $-2\log\frac1\lambda$\\\hline
\end{tabular}
\end{center}
Checks: $\varrho_{A}+\frac{t}{(1+t)^{2}}=\varrho$ (T1 plus the unsubtracted T2 is the whole row); T3 $+$ BFKL-moved $=\frac{t}{(1+t)^{2}}$, the density of the unsubtracted T2 (the recoil rewriting is exact). Total of the written finite part of Row~III with c.c.: $\big(\frac{11}{3}-\frac13\cos2\varphi\big)\mathsf{u}$. T1 has no $\log\vee$: $\varrho_{A}={\cal O}(t^{-3})$. The first parts of Group~II Row~I and Group~III Row~III have no joint tail: there the measured gluon of one side enters only through $\mathcal{K}(x'-w')$ (resp.\ $\mathcal{K}(y-w)$) and the phase, and $\int_{w}e^{-ik\cdot w}\mathcal{K}(y-w)$ ties it to the projectile [A]; their large-distance region is the soft one (\texttt{fix\_soft}).

\subsubsection{Step 6: sums by sector}
\noindent [A, \texttt{t3} [5]] Finite $+$ DGLAP $+$ RC: $\big(-\frac{11}{6}+\frac16\cos2\varphi\big)+\big(-\frac{11}{6}+\frac16\cos2\varphi\big)+\big(\frac{11}{3}-\frac13\cos2\varphi\big)=0$. BFKL: $(1+1-2)\log V=0$; JIMWLK: $(1+1-2)\log\frac1\lambda=0$. The sector-wise cancellation is not automatic (the three rows use different end-point subtractions: $\frac{1}{t}$ split at $t=1$ in Row~I, $\frac{1}{t(1+t)}$ and $\frac{1}{1+t}$ in Row~III, $\frac{1}{\xi},\frac{1}{1-\xi}$ in Row~IV); it holds because every row's split is an exact rewriting over the same range $\Lambda<p^{+}<\vee-k^{+}$ and the end-point coefficients are equal ($1,1,-2$) at both ends. The BFKL and JIMWLK \emph{formulas} are consistent with this: the only BFKL block with all of $w,w',z$ far is $\mathcal{K}(w-z)\cdot\mathcal{K}(w'-z)\,\mathcal{K}(x-z)\cdot\mathcal{K}(x'-z)$, whose colour factor vanishes for uniform far lines up to the commutator (contact) term [N, \texttt{t2}: $P_{1}+P_{2}\le1.2\times10^{-14}$] and contains no term bilinear in $U(x),U(x')$; the JIMWLK formula is $-H_{\rm JIMWLK}$ on the LO structure, and $(L_{w}-U(z)R_{w})U(w)=0$ for $U(z)=U(w)$, while $H$ does not act on the $U(x')U(x)$ part. So neither formula carries a joint tail; the $\pm\log$ tails of the single rows cancel inside them.

\noindent\textbf{Where the scaleless rule enters.} Nowhere in this region: the dropped poles $\frac{2}{\epsilon}$ of Row~IV multiply $\int d^{2}u\,e^{ik\cdot u}G_{\infty}/Y^{2}\propto\delta^{(2)}(k)$ (no tail at $k\ne0$); in Row~I no pole multiplies $d_{\perp}$ (the ${\cal O}(Z^{0})$ parts of $\mathfrak{K}_{2f}$ cancel at $d_{\perp}=2$, Step~4).

\subsubsection{Step 7: colour-neutral projectile and general target}
\noindent [N, \texttt{t2}] (a) Uniform far lines (finite target, $V=1$, or any constant $V$), coloured or not: the three colour factors are $N_{c}\Gamma,N_{c}\Gamma,-N_{c}\Gamma$ (Step~1): cancellation as above. (b) Colour-neutral projectile, arbitrary target ($U(w),U(w'),U(z),U(y),U(y')$ all different and depending on the relative coordinates): at order $1/Z^{2}$ the kernels do not depend on $x,x'$, so every part of a colour factor that is not bilinear in $U(x),U(x')$ multiplies $\int_{x}\rho(x)$ or $\int_{x'}\rho(x')$ and vanishes on a singlet state ($Q^{a}$ stands next to the state in the written order). The bilinear parts are far-line independent: $+N_{c}U^{T}(x')U(x)$, $+N_{c}U^{T}(x')U(x)$, $-N_{c}U^{T}(x')U(x)$ for Rows I, IV, III [deviations $\le3.3\times10^{-15}$, $SU(2)$, $SU(3)$]; T2 has no bilinear part in $\mathcal{L}(y')-\mathcal{L}(z)$. The tail cancels for any target. (c) Coloured projectile, target not uniform at large distance: the charge parts do not cancel (norm of the remainder of $I+IV+2\,III$: $3.6$ ($SU(2)$), $4.8$ ($SU(3)$), zero to $5\times10^{-15}$ with uniform lines); this is the term $[U(w)U(z)-U(y)]\,b\,Q$ of the amplitude (Sec.~\ref{sec:poles-dglap}, Step 14), the non-singlet large-distance term of OPEN~II/III-4(b), not a defect of the finite terms.

\subsubsection{Corrections to statements}\label{ft:corr}
\noindent No term of the written finite terms has to be changed. The following statements are corrected:
\begin{enumerate}
\item {\color{red}\texttt{Finite\_terms\_final}, OPEN~II/III-1, ``No term of Groups II, III cancels it'': the tail $2N_{c}\mathcal{C}\frac{4\pi}{k^{2}}\int\frac{d^{2}Y}{Y^{2}}[-\frac{11}{6}+\frac16\cos2\varphi]G_{\infty}$ of Group~II Row~IV (i)+(ii)+(iii) is cancelled by Group~II Row~I, second part ($\sum\mathfrak{K}_{2\cdot}$: $-\frac{11}{6}+\frac16\cos2\varphi$) and Group~III Row~III, second part with c.c.\ ($+\frac{11}{3}-\frac13\cos2\varphi$), in the region where $w,w',z$ (and $y'$) are far together. The terms are integrable at large distance only together: the $d^{2}Y$ (pair position) integral of these three rows has to be done under one integral, with $Y=\frac{y+y'}{2}$ of Row~IV identified with $z$ (or any point of the pair) of the other rows.}
\item {\color{red}\texttt{Finite\_terms\_final}, Sec.~``Large-distance tails'', list of terms ``without a $1/z^{2}$ tail'': true at fixed $w,w'$ only. In the joint region $\sum\mathfrak{K}_{2\cdot}$, $N_{2}^{\rm reg}$, $N_{2}^{\log}$ do have $1/Z^{2}$ tails (values in Step~5); $N_{2}^{\Psi}$ has, besides its soft tail of the table, a joint tail $+2\mathsf{u}$ (with c.c.).}
\item {\color{red}Sec.~\ref{sec:poles-dglap}, ``What remains (item 4)'': done; the compensating tail is present in the written finite terms, sector by sector, and requires no change of the RC/DGLAP split. With the local (tail-free) DGLAP and RC of Step~12, the remainder $N_{\mathcal{R}}$ carries the whole Row~IV tail and must be grouped with $\sum\mathfrak{K}_{2\cdot}$ and Group~III Row~III (second part) as stated there.}
\end{enumerate}

\subsubsection{Not done / open}
\noindent (1) The joint limit was taken with $x,x'$ fixed and the far lines uniform (or the neutral-projectile reduction); the transition region where only some of $w,w',z$ are far (e.g.\ $|a|\sim Z$) was not analysed beyond the convergence of the Fourier integrals of the leading terms (oscillation at $|a|\gg1/|k|$). (2) The commutativity of the $p^{+}$ integral and the joint limit is verified numerically for Row~I only (Step~4); for Row~III the terms are $p^{+}$ integrands and the limit was taken at fixed $t$ with $\varrho_{A}={\cal O}(t^{-3})$, $\frac{1}{(1+t)^{2}}$ integrable. (3) No full transverse integration with Wilson lines. (4) The statement about the BFKL and JIMWLK formulas is analytic (colour of the one far block; $-H_{\rm JIMWLK}$ form as checked in \texttt{Final\_terms\_checked}), not a term-by-term numerical evaluation of the written formulas in the joint limit. (5) The coloured-projectile remainder of Step~7(c) is not addressed here.


% ---- inlined copy of paper/notes/fix_dabc.tex (20 Nov 2026 14:20 UTC) ----
\subsection{Symmetrised three-point function $\propto d^{abc}$: the odd part of the finite terms}\label{sec:dabc}
\noindent\textbf{Result in one paragraph.} The author replaces (A2) by
\begin{align}
\langle\{\rho^{a}(x)\rho^{b}(y)\rho^{c}(z)\}_{\rm sym}\rangle=d^{abc}\,\frac{N_{c}}{(N_{c}^{2}-4)(N_{c}^{2}-1)}\,o(x,y,z)\,,\qquad o(x,y,z)\equiv\langle d^{efg}\{\rho^{e}(x)\rho^{f}(y)\rho^{g}(z)\}_{\rm sym}\rangle\,,\tag{A2$'$}
\end{align}
(normalisation confirmed by the author; $d^{abc}d^{abc}=(N_{c}^{2}-4)(N_{c}^{2}-1)/N_{c}$; $o$ is real and totally symmetric). Then every cubic finite term acquires, in addition to its contact terms (which are unchanged, \texttt{fix\_singlet}), a \emph{symmetrised part} $\propto o$. For every cubic colour structure of the finite terms we give it in colour-singlet form: it is the row's $c$-number factor times $\frac{N_{c}}{(N_{c}^{2}-4)(N_{c}^{2}-1)}\,\Delta\,o$, where $\Delta=\sum_{\rm orderings}C^{abc}d^{abc}$ is a combination of a single new invariant $\mf{G}(a,b,c)=f^{ijk}d^{lmn}U^{il}(a)U^{jm}(b)U^{kn}(c)$ (totally antisymmetric, C-odd) and, in Rows II.I and III.II, of its four-line analogue $\mf{Q}_{d}$. In the fundamental representation
\begin{align}
\mf{G}(a,b,c)={\rm Im}\big[{\rm tr}(V_{a}^{\dagger}V_{b})\,{\rm tr}(V_{b}^{\dagger}V_{c})\,{\rm tr}(V_{c}^{\dagger}V_{a})\big]+{\rm Im}\,{\rm tr}\big(V_{a}^{\dagger}V_{b}V_{c}^{\dagger}V_{a}V_{b}^{\dagger}V_{c}\big)\,,\notag
\end{align}
derived exactly (Fierz) and verified for $N_{c}=2,\dots,6$. \textbf{All thirteen cubic colour structures (Rows I.II $N_{2b}$, I.III, I.V first part, II.I first part, III.II, III.III first part, III.V, III.VI) have a non-zero symmetrised part}; none vanishes identically. With the c.c.\ (Groups I, III), resp.\ with $N_{1a}+N_{1b}$ (Group II), it is real and \emph{odd} under $k\to-k$: it is the odderon-type, $k$-odd part of the single inclusive spectrum. It is C-odd in the projectile ($o$) and in the target ($\mf{G}$, $\mf{Q}_{d}$ change sign under $V\to V^{*}$; for a weak target they start at three gluon exchanges), so it vanishes after the target average for a C-even target ensemble (MV), and it vanishes for $SU(2)$. \textbf{The three-$\rho$ parts of the logarithmic (BFKL, JIMWLK) terms still cancel} for this three-point function: the cancellation is an identity of the classical (totally symmetrised) colour tensors, verified now also for $SU(3)$ and $SU(4)$ and in the $d^{abc}$ projection; this is consistent with the LO term having no three-$\rho$ part (Sec.~\ref{sec:dabc-evol}). Numerical checks: colour algebra $2.4\times10^{-15}$; fundamental forms $2.3\times10^{-15}$; operator level (explicit $SU(3)$ charges, C-odd ensemble with $d$-type three-point function) $7.7\times10^{-16}$; parity exact; evolution cancellation $|{\rm Sym}_{3}|\le2\times10^{-15}$ against termwise $d$-projections up to $44$. Drop-in replacements for the paper and for \texttt{Final\_terms\_fixed} are in Sec.~\ref{sec:dabc-dropin}. Scripts in \texttt{scratchpad/fix/dabc/} (Sec.~\ref{sec:dabc-num}).

\subsubsection{The assumption and its normalisation}\label{sec:dabc-ass}
\noindent By (A1) (global colour invariance) the symmetrised three-point function $S^{abc}(x,y,z)\equiv\langle\{\rho^{a}(x)\rho^{b}(y)\rho^{c}(z)\}_{\rm sym}\rangle$ ($\{\dots\}_{\rm sym}$: average over the six orderings) is an invariant tensor, $S^{abc}=d^{abc}g_{d}+f^{abc}g_{f}$, with $g_{d}$ real and totally symmetric, $g_{f}$ real and totally antisymmetric in the three points (\texttt{fix\_singlet}, Sec.~``Projectile average''). Contracting with $d^{abc}$ and using $d^{abc}f^{abc}=0$, $d^{abc}d^{abc}=(N_{c}^{2}-4)(N_{c}^{2}-1)/N_{c}$,
\begin{align}
g_{d}(x,y,z)=\frac{N_{c}}{(N_{c}^{2}-4)(N_{c}^{2}-1)}\,o(x,y,z)\,,\qquad g_{f}(x,y,z)=\frac{1}{N_{c}(N_{c}^{2}-1)}\,\langle f^{efg}\{\rho^{e}(x)\rho^{f}(y)\rho^{g}(z)\}_{\rm sym}\rangle\,.\label{eq:dabc-gdgf}
\end{align}
(A2$'$) is therefore equivalent to (A1) together with $g_{f}=0$, and it reduces to the former (A2) when $o=0$ (Gaussian projectile, or any ensemble invariant under the antiunitary map $\rho^{a}\to-\rho^{a}$, \texttt{fix\_singlet}). $o$ is C-odd: under charge conjugation of the projectile, $\rho^{a}t^{a}\to-(\rho^{a}t^{a})^{T}$, the cubic Casimir density $d^{efg}\rho^{e}\rho^{f}\rho^{g}$ changes sign; $o$ is the projectile odderon. A model in which it is non-zero is the weight functional of Jeon and Venugopalan (hep-ph/0503219, Eq.~(17)), $\exp\{-\int_{x}[\rho^{a}\rho^{a}/(2\mu_{A}^{2})-d_{abc}\rho^{a}\rho^{b}\rho^{c}/\kappa_{A}]\}$, whose cubic term ``generates Odderon excitations'' (their abstract); to first order in $1/\kappa_{A}$ it gives (our evaluation by Wick's theorem, their normalisation of $\mu_{A}$, $\kappa_{A}$, classical $\rho$ = Weyl symbol) $S^{abc}=\frac{6\mu_{A}^{6}}{\kappa_{A}}d^{abc}\delta^{(2)}(x-y)\delta^{(2)}(x-z)$, i.e.\ a \emph{local} $o$ (see Sec.~\ref{sec:dabc-local}).\par
\noindent\emph{$SU(2)$:} $d^{abc}\equiv0$, so $o\equiv0$ and the normalisation factor is $0/0$; for $SU(2)$, (A2$'$) is to be read as $S^{abc}=0$ (the only invariant is $f^{abc}$, which (A2$'$) excludes). All results below vanish identically for $SU(2)$ ($\mf{G}=\mf{Q}_{d}=0$).

\subsubsection{General formula}
\noindent Every cubic finite term has the form $A(k)\,C^{abc}\rho^{a}(p_{1})\rho^{b}(p_{2})\rho^{c}(p_{3})$ (ordered, type ``ord''), or with an anticommutator, $\rho(p_{1})\{\rho(p_{2}),\rho(p_{3})\}$ (A23) or $\{\rho(p_{1}),\rho(p_{2})\}\rho(p_{3})$ (A12); $A(k)$ is the $c$-number factor (row prefactor and kernel), $C$ is real (adjoint $U$, $f^{abc}$). Expanding the anticommutators into orderings, the symmetrised part of the average is $\sum_{\rm orderings}C^{abc}_{\rm seq}S^{abc}(\text{points of seq})$; since $S$ is symmetric under simultaneous permutations of indices and points, and $g_{d}$ totally symmetric,
\begin{align}
\big\langle C^{abc}\rho^{a}(p_{1})\rho^{b}(p_{2})\rho^{c}(p_{3})\big\rangle\big|_{\rm sym}=\frac{N_{c}}{(N_{c}^{2}-4)(N_{c}^{2}-1)}\,\Delta\,o(p_{1},p_{2},p_{3})\,,\qquad\Delta\equiv\sum_{\rm orderings}C^{abc}_{\rm seq}\,d^{abc}\,,\label{eq:dabc-gen}
\end{align}
($\Delta=C^{abc}d^{abc}$ for ``ord'', $\Delta=2C^{abc}d^{abc}$ for A23, A12, with $C$ the coefficient of the first ordering). The full average is
\begin{align}
\langle C\rho\rho\rho\rangle=\frac{N_{c}}{(N_{c}^{2}-4)(N_{c}^{2}-1)}\,\Delta\,o(p_{1},p_{2},p_{3})+[\text{contact terms of \texttt{fix\_singlet}, Eqs.~(s-ord)--(s-a12), unchanged}]\,.\notag
\end{align}
\emph{If $g_{f}\neq0$} (not part of (A2$'$)), the same argument adds $t\,\Phi\,g_{f}(p_{1},p_{2},p_{3})$, $t=1$ (ord), $2$ (A12, A23), with $\Phi=C^{abc}f^{abc}$ the singlet factor of \texttt{fix\_singlet} at \emph{separate} points (verified at the operator level, Sec.~\ref{sec:dabc-num}).

\subsubsection{The invariants $\mf{G}$ and $\mf{Q}_{d}$}\label{sec:dabc-inv}
\noindent Every cubic colour structure has exactly one $f$; contracting its three $\rho$ indices with $d$ gives a network with one $f$, one $d$ and chains of Wilson lines. Chains joining two legs of $f$ to two legs of $d$ directly give $f^{abe}d^{abe'}=0$ (e.g.\ $f^{hbc}U^{ah}(w)\{\rho^{b},\rho^{c}\}$, $U^{ac'}(w')U^{ac}(w)f^{cbd}\rho^{c'}\rho^{b}\rho^{d}$); the others reduce, with $U^{T}U=1$, to (notation of \texttt{fix\_singlet}, $(T^{e})_{bc}=-if^{ebc}$, and $(D^{e})_{bc}\equiv d^{ebc}$)
\begin{align}
\mf{G}(a,b,c)&\equiv f^{ijk}\,d^{lmn}\,U^{il}(a)\,U^{jm}(b)\,U^{kn}(c)\,,\notag\\
\mf{Q}_{d}(a,b;c,e)&\equiv i\,{\rm Tr}\{T^{h}\,U^{\dagger}(a)U(b)\,D^{h}\,U^{\dagger}(c)U(e)\}=f^{hij}d^{hkl}\,[U^{\dagger}(a)U(b)]^{jk}\,[U^{\dagger}(c)U(e)]^{li}\,.\label{eq:dabc-GQ}
\end{align}
\textbf{Properties} (all verified numerically, Sec.~\ref{sec:dabc-num}): $\mf{G}$ is real and totally antisymmetric; $\mf{G}(a,a,c)=0$; $\mf{Q}_{d}$ is real, $\mf{Q}_{d}(a,b;c,e)=-\mf{Q}_{d}(e,c;b,a)$, $\mf{Q}_{d}(a,a;c,e)=\mf{Q}_{d}(a,b;c,c)=0$, $\mf{Q}_{d}(a,b;c,a)=-\mf{G}(a,b,c)$, and $\mf{Q}_{d}(a,b;b,c)=\mf{G}(a,b,c)|_{U\to U^{T}}$ (i.e.\ $V\to V^{\dagger}$; not reducible to $\mf{G}$ itself). Both are invariant under $U(p)\to RU(p)R'$ for all $p$ (global rotations $R,R'$), hence vanish for constant $U$, and both are \emph{C-odd}: under $V\to V^{*}$, $U\to\eta U\eta$ with $\eta={\rm diag}(\pm1)$, and $\eta_{a}\eta_{b}\eta_{c}f^{abc}=-f^{abc}$, $\eta_{a}\eta_{b}\eta_{c}d^{abc}=+d^{abc}$, so $\mf{G}\to-\mf{G}$, $\mf{Q}_{d}\to-\mf{Q}_{d}$ (whereas $\mf{D},\mf{F},\mf{Q}$ are C-even). For a weak target, $V_{p}=e^{i\epsilon\alpha_{p}^{a}t^{a}}$, $\mf{G}$ and every $\Delta$ below are $O(\epsilon^{3})$: the odd part needs three gluon exchanges with the target (a C-odd, $d^{abc}$ exchange). For $SU(2)$, $\mf{G}=\mf{Q}_{d}=0$.

\noindent\textbf{Fundamental-trace forms (exact).} With $U^{ab}(p)=2{\rm tr}(t^{a}V_{p}t^{b}V_{p}^{\dagger})$ and $A\equiv V_{a}^{\dagger}V_{b}$, $B\equiv V_{c}^{\dagger}V_{e}$,
\begin{align}
\mf{Q}_{d}(a,b;c,e)&=-{\rm Im}\big[{\rm tr}A\,{\rm tr}B\,{\rm tr}(A^{\dagger}B^{\dagger})\big]-{\rm Im}\,{\rm tr}\big(ABA^{\dagger}B^{\dagger}\big)\,,\label{eq:dabc-Qfund}\\
\mf{G}(a,b,c)=-\mf{Q}_{d}(a,b;c,a)&={\rm Im}\big[{\rm tr}(V_{a}^{\dagger}V_{b})\,{\rm tr}(V_{b}^{\dagger}V_{c})\,{\rm tr}(V_{c}^{\dagger}V_{a})\big]+{\rm Im}\,{\rm tr}\big(V_{a}^{\dagger}V_{b}V_{c}^{\dagger}V_{a}V_{b}^{\dagger}V_{c}\big)\,.\label{eq:dabc-Gfund}
\end{align}
\emph{Derivation.} $\sum_{h}f^{hij}d^{hkl}=-2i\,{\rm tr}([t^{i},t^{j}]\{t^{k},t^{l}\})$ (completeness; the $1/N_{c}$ term drops since ${\rm tr}[t^{i},t^{j}]=0$), and $\sum_{j}t^{j}[U^{\dagger}(a)U(b)]^{jk}=At^{k}A^{\dagger}$, $\sum_{i}t^{i}[U^{\dagger}(c)U(e)]^{li}=B^{\dagger}t^{l}B$, so $\mf{Q}_{d}=-2i\sum_{k,l}{\rm tr}\big([B^{\dagger}t^{l}B,At^{k}A^{\dagger}]\{t^{k},t^{l}\}\big)$. The sums over $k$ and $l$ with $\sum_{k}{\rm tr}(Pt^{k}Qt^{k})=\frac12{\rm tr}P\,{\rm tr}Q-\frac{1}{2N_{c}}{\rm tr}(PQ)$ and $\sum_{k}{\rm tr}(Pt^{k}){\rm tr}(Qt^{k})=\frac12{\rm tr}(PQ)-\frac{1}{2N_{c}}{\rm tr}P\,{\rm tr}Q$ give, after all $1/N_{c}$ terms cancel, $\sum_{k,l}(\dots)=-\frac{i}{2}\big[{\rm Im}\,Z_{1}+{\rm Im}\,Z_{2}\big]$ with $Z_{1}={\rm tr}A\,{\rm tr}B\,{\rm tr}(A^{\dagger}B^{\dagger})$, $Z_{2}={\rm tr}(ABA^{\dagger}B^{\dagger})$; hence Eq.~\eqref{eq:dabc-Qfund}. The coefficients are $N_{c}$-independent; the same result came out of an unbiased fit to all trace monomials for $N_{c}=4,\dots,7$ and was checked for $N_{c}=2,\dots,6$ (Sec.~\ref{sec:dabc-num}). These are the standard C-odd (odderon) combinations: ${\rm Im}$ of the product of three dipoles around the triangle $a,b,c$ and ${\rm Im}$ of the sextupole; both are odd under the exchange of two points.

\subsubsection{The symmetrised part of every cubic structure}\label{sec:dabc-table}
\noindent $\Delta$ for each cubic colour structure (labels and $\rho$ orderings of \texttt{fix\_singlet}; the structures are those of \texttt{Final\_terms\_fixed}, identical to \texttt{Final\_terms\_final}, \texttt{Final\_terms\_checked} and the paper appendix), with arguments ordered by the list $x,y,x',y',z,w,w'$. ``Local'' is $\Delta$ with all three $\rho$ points set equal (relevant for a local $o$, Sec.~\ref{sec:dabc-local}).
\begin{center}
\renewcommand{\arraystretch}{1.35}
\small\begin{tabular}{|l|l|l|l|}\hline
structure & $\rho$ points, type & $\Delta=\sum C^{abc}d^{abc}$ & local \\\hline
I.II, $N_{2b}$ & $x',x,y$; A23 & $2\big[\mf{G}(x,y,w')-\mf{G}(x,y,x')\big]$ & $0$\\
I.III, $C_{D}=f^{abd}\mathcal{C}_{\rm III}(z)$ & $x',x,y$; ord & $\mf{G}(x,y,x')-\mf{G}(x,y,w')+\mf{G}(x,x',z)+\mf{G}(x,z,w')$ & $\mf{G}(x,z,w')$\\
I.III, $C_{A}=f^{abd}\mathcal{C}_{\rm III}(w)$ & $x',x,y$; ord & $\mf{G}(x,y,x')-\mf{G}(x,y,w')+\mf{G}(x,x',w)+\mf{G}(x,w,w')$ & $\mf{G}(x,w,w')$\\
I.V, $N_{1}$ & $x',y,y'$; A23 & $-2\big[\mf{G}(x,x',z)+\mf{G}(x,z,w')+\mf{G}(y,x',y')+\mf{G}(y,y',w')\big]$ & $-2\mf{G}(x,y,z)-2\mf{G}(x,z,w')$\\
II.I, $N_{1a}$ & $x',y,x$; A12 & $2\big[\mf{G}(x,z,w')-\mf{G}(x,y,x')-\mf{Q}_{d}(z,y;x',w)\big]$ & $2\mf{G}(x,z,w')-2\mf{Q}_{d}(z,x;x,w)$\\
II.I, $N_{1b}$ & $x,x',y$; A23 & $2\big[\mf{G}(x,z,w)-\mf{G}(x,y,x')-\mf{Q}_{d}(z,y;x',w')\big]$ ($=N_{1a}|_{w\leftrightarrow w'}$) & $2\mf{G}(x,z,w)-2\mf{Q}_{d}(z,x;x,w')$\\
III.II & $y,x',x$; ord & $\mf{G}(x,y,x')+\mf{G}(x,x',z)+\mf{Q}_{d}(z,y;x',w)$ & $\mf{Q}_{d}(z,x;x,w)$\\
III.III$_{1}$, $\mathcal{L}(y')$ & $x',y,x$; A23 & $2\big[\mf{G}(x,y,x')-\mf{G}(x,y,y')-\mf{G}(x',z,w)+\mf{G}(y',z,w)\big]$ & $2\mf{G}(y',z,w)-2\mf{G}(x,z,w)$\\
III.III$_{1}$, $\mathcal{L}(z)$ & $x',y,x$; A23 & $2\big[\mf{G}(x,y,x')-\mf{G}(x,y,z)-\mf{G}(x',z,w)\big]$ & $-2\mf{G}(x,z,w)$\\
III.V, $\mathcal{L}_{\rm V}(y')$ & $x',x,y$; ord & $-\mf{G}(x,y,x')+\mf{G}(x,y,y')-\mf{G}(x,x',w)+\mf{G}(x,y',w)$ & $\mf{G}(x,y',w)$\\
III.V, $\mathcal{L}_{\rm V}(z)$ & $x',x,y$; ord & $-\mf{G}(x,y,x')+\mf{G}(x,y,z)-\mf{G}(x,x',w)+\mf{G}(x,z,w)$ & $\mf{G}(x,z,w)$\\
III.VI, $\mathcal{L}(y')$ & $x',y,x$; ord & $-\mf{G}(x,y,x')+\mf{G}(x,y,y')+\mf{G}(y,x',z)-\mf{G}(y,y',z)$ & $-\mf{G}(x,y',z)$\\
III.VI, $\mathcal{L}(z)$ & $x',y,x$; ord & $-\mf{G}(x,y,x')+\mf{G}(x,y,z)+\mf{G}(y,x',z)$ & $0$\\\hline
\end{tabular}
\end{center}
The $\mathcal{L}(z)$ forms are the $\mathcal{L}(y')$ forms with $y'\to z$ (using $\mf{G}(a,a,c)=0$), and $C_{A}$ is $C_{D}$ with $z\to w$. All thirteen $\Delta$ are non-zero for generic points (typical values $0.2$--$6$ for random $SU(3)$, $SU(4)$, $SU(5)$ lines, Sec.~\ref{sec:dabc-num}). The quadratic rows (I.II $N_{1}$, I.V second part, II.I second part, II.IV, III.III second part) and the DGLAP and running-coupling terms are unchanged. \emph{Pairs:} the colour identities~(colid) of the soft-remainder pairs are operator identities, so the symmetrised parts pair exactly as the contact terms do; e.g.\ $\Delta_{\rm III.V}(\mathcal{L}_{\rm V}(y'))|_{y'\to w'}=-\Delta_{\rm I.III}(C_{A})$ and $\Delta_{\rm III.VI}(\mathcal{L}(y'))|_{y'\to w',\,x\leftrightarrow y}=+\Delta_{\rm I.III}(C_{D})$ (checked numerically, $9\times10^{-16}$).

\noindent\textbf{Row factors.} From the row prefactors of \texttt{Final\_terms\_fixed} (equivalently: the contact line coefficient of \texttt{fix\_singlet} divided by $c\,i$), in units of $\mathcal{N}_{\rm f}=\frac{1}{(2\pi)^{3}}\frac{g^{4}}{16\pi^{5}}\frac{1}{k^{+}}$ ($=\mathcal{C}_{\rm f}$ of the paper), the $c$-number factor $A$ of each cubic structure is: I.II $N_{2b}$: $-\frac{1}{4\pi}\mathfrak{K}^{\rm II}_{2b}$; I.III: $2i\,\mathfrak{K}^{\rm III}_{z}$ ($C_{D}$), $2i\,\mathfrak{K}^{\rm III}_{w}$ ($C_{A}$); I.V $N_{1}$: $\frac{i}{2}\mathfrak{K}^{\rm V}_{1}$; II.I: $-i\,\mathfrak{K}_{\mathcal{T}}(w,w')$ ($N_{1a}$), $+i\,\mathfrak{K}_{\mathcal{T}}(w',w)$ ($N_{1b}$); III.II: $-2i\,\mathfrak{K}_{\mathcal{T}}(w,w')$; III.III first part: $-i[-\mathfrak{K}_{\widehat{\mathcal{E}}}+\mathfrak{K}_{\Psi}+\mathfrak{K}^{(y')}_{\delta}]$ ($\mathcal{L}(y')$), $-i\,\mathfrak{K}^{(z)}_{\delta}$ ($\mathcal{L}(z)$); III.V and III.VI: $-2i[-\mathfrak{K}_{\widehat{\mathcal{E}}}+\mathfrak{K}_{\Psi}+\mathfrak{K}^{(y')}_{\delta}]$ ($\mathcal{L}(y')$), $-2i\,\mathfrak{K}^{(z)}_{\delta}$ ($\mathcal{L}(z)$); times the phase of the row. (Spot checks against the printed prefactors: I.III $\frac{ig^{4}}{8\pi^{5}}$, II.I $\mp\frac{ig^{4}}{16\pi^{5}}$, III.II $-\frac{ig^{4}}{8\pi^{5}}$, III.III $\frac{ig^{4}}{16\pi^{5}}$, III.V $\frac{ig^{4}}{8\pi^{5}}$.)

\noindent\textbf{Singlet form of the odd part.} With the kernels of \texttt{fix\_singlet} / the paper,
\begin{align}
\frac{d^{3}N_{\rm odd}}{d^{3}k}\bigg|_{\rm singlet}&=\frac{\mathcal{N}_{\rm f}}{N_{c}^{2}-1}\,\frac{N_{c}}{N_{c}^{2}-4}\Bigg\{\int_{w,w'}e^{-ik\cdot(w-w')}\int_{x,y,x'}\frac{1}{2\pi}\,\mathfrak{K}^{\rm II}_{2b}\Big[\mf{G}(x,y,x')-\mf{G}(x,y,w')\Big]\,o(x,y,x')\notag\\
&\quad+\int_{w,w'}e^{-ik\cdot(w-w')}\int_{x,y,z,x'}2i\Big\{\mathfrak{K}^{\rm III}_{z}\Big[\mf{G}(x,y,x')-\mf{G}(x,y,w')+\mf{G}(x,x',z)+\mf{G}(x,z,w')\Big]\notag\\
&\qquad\qquad+\mathfrak{K}^{\rm III}_{w}\Big[\mf{G}(x,y,x')-\mf{G}(x,y,w')+\mf{G}(x,x',w)+\mf{G}(x,w,w')\Big]\Big\}\,o(x,y,x')\notag\\
&\quad-\int_{x,z,y,y',x',w'}e^{-ik\cdot(z-w')}\,i\,\mathfrak{K}^{\rm V}_{1}\Big[\mf{G}(x,x',z)+\mf{G}(x,z,w')+\mf{G}(y,x',y')+\mf{G}(y,y',w')\Big]\,o(y,x',y')\notag\\
&\quad-\int_{x,y,x',z,w,w'}e^{-ik\cdot(w-w')}\,2i\,\mathfrak{K}_{\mathcal{T}}(w,w')\Big[\mf{G}(x,y,x')+\mf{G}(x,x',z)+\mf{Q}_{d}(z,y;x',w)\Big]\,o(x,y,x')\notag\\
&\quad-\int_{w,x',y',z,y,x}e^{-ik\cdot(w-y')}\,2i\,\Big\{\Big[-\mathfrak{K}_{\widehat{\mathcal{E}}}+\mathfrak{K}_{\Psi}+\mathfrak{K}^{(y')}_{\delta}\Big]\Big[\mf{G}(x,y,x')-\mf{G}(x,y,y')-\mf{G}(x',z,w)+\mf{G}(y',z,w)\Big]\notag\\
&\qquad\qquad+\mathfrak{K}^{(z)}_{\delta}\Big[\mf{G}(x,y,x')-\mf{G}(x,y,z)-\mf{G}(x',z,w)\Big]\Big\}\,o(x,y,x')\notag\\
&\quad-\int_{w,y',z,y,x,x'}e^{-ik\cdot(w-y')}\,2i\,\Big\{\Big[-\mathfrak{K}_{\widehat{\mathcal{E}}}+\mathfrak{K}_{\Psi}+\mathfrak{K}^{(y')}_{\delta}\Big]\Big[-\mf{G}(x,y,x')+\mf{G}(x,y,y')-\mf{G}(x,x',w)+\mf{G}(x,y',w)\Big]\notag\\
&\qquad\qquad+\mathfrak{K}^{(z)}_{\delta}\Big[-\mf{G}(x,y,x')+\mf{G}(x,y,z)-\mf{G}(x,x',w)+\mf{G}(x,z,w)\Big]\Big\}\,o(x,y,x')\notag\\
&\quad-\int_{w,y',z,y,x,x'}e^{-ik\cdot(w-y')}\,2i\,\Big\{\Big[-\mathfrak{K}_{\widehat{\mathcal{E}}}+\mathfrak{K}_{\Psi}+\mathfrak{K}^{(y')}_{\delta}\Big]\Big[-\mf{G}(x,y,x')+\mf{G}(x,y,y')+\mf{G}(y,x',z)-\mf{G}(y,y',z)\Big]\notag\\
&\qquad\qquad+\mathfrak{K}^{(z)}_{\delta}\Big[-\mf{G}(x,y,x')+\mf{G}(x,y,z)+\mf{G}(y,x',z)\Big]\Big\}\,o(x,y,x')\Bigg\}+{\rm c.c.}\notag\\
&\quad-\frac{\mathcal{N}_{\rm f}}{N_{c}^{2}-1}\,\frac{N_{c}}{N_{c}^{2}-4}\int_{x,y,z,w,x',w'}4\sin\big(k\cdot(w-w')\big)\,\mathfrak{K}_{\mathcal{T}}(w,w')\Big[\mf{G}(x,z,w')-\mf{G}(x,y,x')-\mf{Q}_{d}(z,y;x',w)\Big]\,o(x,y,x')\,.\label{eq:dabc-odd}
\end{align}
Lines: I.II $N_{2b}$; I.III; I.V first part; III.II; III.III first part; III.V; III.VI ($+$c.c.\ for all of these); last line: II.I first part, $N_{1a}+N_{1b}$ (no c.c.), obtained from $-i\,e^{-ik\cdot(w-w')}\mathfrak{K}_{\mathcal{T}}(w,w')\Delta_{1a}+i\,e^{-ik\cdot(w-w')}\mathfrak{K}_{\mathcal{T}}(w',w)\Delta_{1b}$ by $w\leftrightarrow w'$ in the second term ($\Delta_{1b}=\Delta_{1a}|_{w\leftrightarrow w'}$). In the III.III--III.VI lines the pairs $\mathfrak{K}^{(y')}_{\delta}$, $\mathfrak{K}^{(z)}_{\delta}$ stay under one $y'$ integral, as for the contact terms. The \emph{full} singlet form of a cubic row is its \texttt{fix\_singlet} (contact) form plus its line of Eq.~\eqref{eq:dabc-odd}.

\subsubsection{Parity, reality, target average}\label{sec:dabc-parity}
\noindent $\Delta$ and $o$ are real. For Groups I and III the term plus its c.c.\ has the symmetrised part $2\,{\rm Re}[A(k)]\,\frac{N_{c}}{(N_{c}^{2}-4)(N_{c}^{2}-1)}\Delta\,o$. All kernels depend on $k$ through $e^{-ik\cdot r}$, $\mathcal{E}$, $\mathcal{J}$, $\widehat{\mathcal{E}}$, $\Psi$, $\pi\delta(s)-i{\rm P}\frac1s$, all of the form $\int dt\,(\text{real})\,e^{-ist}$, so $f(-s)=f(s)^{*}$; the cubic rows have an explicit $i$ (in I.II $N_{2b}$ the kernel $\mathfrak{K}^{\rm II}_{2b}$ is imaginary, the Fourier transform of the odd real $\mathcal{R}^{i}$). Hence $A(-k)=-A(k)^{*}$, ${\rm Re}A(-k)=-{\rm Re}A(k)$: \textbf{the symmetrised part is real and odd in $k$, pointwise in the coordinates}. For II.I the integrand of the last line of~\eqref{eq:dabc-odd} is manifestly real and odd. (Checked numerically with the closed forms of $\mathcal{E}$, $\widehat{\mathcal{E}}$, $\Psi$, $\mathcal{J}$, for I.III $C_{A}$, III.V $\mathcal{L}_{\rm V}(y')$ and II.I $N_{1a}+N_{1b}$, Sec.~\ref{sec:dabc-num}.) The contact terms remain even; hence, under (A1)+(A2$'$), \emph{the $k$-odd part of the finite terms is exactly Eq.~\eqref{eq:dabc-odd}}, and the $k$-even part is that of \texttt{fix\_singlet}.\par
\noindent\emph{When does it survive?} (i) It is C-odd in the projectile ($o$) and in the target ($\mf{G},\mf{Q}_{d}$). After an average over a C-even target ensemble (e.g.\ Gaussian MV target) it vanishes; it survives configuration by configuration (event-by-event odd harmonics) or for a target with a C-odd (odderon) component. (ii) All kernels are invariant under the simultaneous rotation by $\pi$ of all coordinates and of $k$ ($r\to-r$, $k\to-k$). Therefore, if the averaged product $\langle\Delta\rangle_{\rm target}\,o$ is invariant under the reflection of all points through a common centre, the odd part integrates to zero; a non-zero result needs a configuration (or an impact-parameter geometry) that breaks this symmetry. (iii) For a weak target it is $O(\epsilon^{3})$ in the target field (three-gluon, $d$-type exchange), and $O(o)$ in the projectile: projectile odderon $\times$ target odderon.

\subsubsection{Local three-point function}\label{sec:dabc-local}
\noindent For a local $o$ (Jeon--Venugopalan type, $o\propto\delta^{(2)}(x-y)\delta^{(2)}(x-z)$) only the ``local'' column of the table enters: I.II $N_{2b}$ and III.VI $\mathcal{L}(z)$ drop out, the other structures do not. \textbf{Not checked:} whether the kernels are integrable when all three $\rho$ points coincide (\texttt{fix\_singlet} checked only pairwise merging, for two kernels); with a local $o$ this has to be power-counted row by row.

\subsubsection{Three-$\rho$ parts of the BFKL and JIMWLK terms}\label{sec:dabc-evol}
\noindent \textbf{Statement in the files.} \texttt{Evolution\_three\_rho\_checked} (Sec.~``Cancellation of all three $\rho$'s'') and the paper (Sec.~``Terms cubic in $\rho$'') state that for each logarithm the \emph{totally symmetrised} three-$\rho$ part of the sum of all logarithmic three-$\rho$ terms (genuine ones and the three-$\rho$ remainders of the reordered four-$\rho$ terms) vanishes identically, for each of the twelve classes (kernel $\alpha,\alpha',\beta,\beta'$ $\times$ two, three, four Wilson lines). The check recorded in \texttt{Evolution\_three\_rho\_checked} was done with $SU(2)$ charges (the $SU(3)$ log of that check in \texttt{scratchpad/ev\_3rho} is empty), for which $d^{abc}=0$, so a $d$-type piece could not have been seen there; the paper states $SU(2)$ and $SU(3)$.\par
\noindent\textbf{Re-check with $SU(3)$, $SU(4)$ and the $d$-projection} [\texttt{ev\_generic.py}]. At distinct points the charges commute, so the symmetrised part of a bracket is the classical (commuting) colour tensor. For each of the twelve cancellation brackets as written, and for the collected list of all rows regrouped by kernel class, the colour tensors were evaluated with random $SU(3)$ and $SU(4)$ Wilson lines at generic points: the summed, totally symmetrised tensor (all colour components) is $\le2\times10^{-15}$, and its $d^{abc}$ projection $\le6\times10^{-15}$, while the sums of the moduli of the termwise $d$-projections are $1.6$--$32$ per bracket ($10$--$44$ per kernel class); the two-Wilson-line classes have vanishing $d$-projection term by term ($f^{abc}$ with two indices on $\rho$'s). Since the vanishing holds for the whole tensor, it holds for \emph{any} symmetrised three-point function: $d^{abc}o$ under (A2$'$), and also an $f^{abc}g_{f}$ part. Contact configurations need no separate treatment: the identity holds for arbitrary Wilson lines at the points, hence also when points coincide, and $\langle\{\rho\rho\rho\}_{\rm sym}\rangle$ is given by (A2$'$) at coincident points as well. \textbf{Result: with $\langle\rho\rho\rho\rangle_{\rm sym}\neq0$ the symmetrised three-$\rho$ parts of the logarithmic terms still cancel; there is no $\ln$-enhanced term $\propto o$}. The two-$\rho$ remainders (part of the BFKL and JIMWLK terms) and the one-$\rho$ remainders (zero average) are not affected.\par
\noindent\textbf{Why this is consistent.} The LO spectrum is quadratic in $\rho$; under (A1) it depends on the projectile only through $n$, which is C-even. The logarithmic terms at $O(g^{4})$ are its one-step evolution: in $\ln(\vee/k^{+})$ the evolution of $n$ (dilute projectile: linear BFKL, Eq.~(BFKLrhorho)), in $\ln(k^{+}/\Lambda)$ the JIMWLK evolution of the target operator at fixed $n$. A $\ln$-enhanced term $\propto o$ would mean that $\partial_{Y}n$ contains $o$ (or that the target evolution of the LO operator produces a three-$\rho$ projectile structure). Both are excluded: the evolution Hamiltonians are C-invariant, so the C-even $n$ cannot mix with the C-odd $o$; and the LO operator has no three-$\rho$ term whose evolution could appear. In addition, an $o$-term would have to come with a C-odd target structure ($\mf{G}$, $\mf{Q}_{d}$), whereas the JIMWLK term is $H_{\rm JIMWLK}$ acting on the C-even LO operator. The $o$-dependence of the spectrum starts with the non-logarithmic $g^{4}\rho^{3}$ terms of Eq.~\eqref{eq:dabc-odd} (the ``first saturation correction''); its rapidity evolution (evolution of $o$, the projectile odderon, and of the C-odd target structures) is of higher order, $\alpha_{s}\ln\times g^{4}\rho^{3}$, and is not part of this calculation. The cancellation found is therefore exactly what is required.

\subsubsection{Literature}\label{sec:dabc-lit}
\noindent Read in this session from the arXiv sources (export.arxiv.org); only what was read is quoted.
\begin{itemize}
\item McLerran, Skokov, arXiv:1611.09870 (NPA 959 (2017) 83): in the classical Yang--Mills dilute--dense calculation the leading order spectrum is symmetric configuration by configuration, $n^{(2)}(k)=n^{(2)}(-k)$; the first odd contribution is $n^{(3)}$, cubic in the projectile charge $\rho_{1}$, with the odd part $\propto{\rm Im}\{\dots f^{abc}\Omega^{a}\Omega^{b}\Omega^{c*}\}$ (their equation labelled \texttt{Eq:FinalAssymitry}); conclusions: the odd harmonics ``in single particle inclusive process \dots average out to zero for a Gaussian ensemble because they are proportional to $\rho_{1}^{3}$. Essentially they are defined by odderon exchanges.'' \emph{Agreement:} LO even, odd part cubic in $\rho$ and odd configuration by configuration, zero for a Gaussian projectile ($o=0$). Our result adds: with non-commuting charges this holds for the symmetrised part only (contact terms are even, \texttt{fix\_singlet}), and with (A2$'$) the odd part survives the projectile average, $\propto o$, in the explicit form~\eqref{eq:dabc-odd}. A term-by-term comparison with that equation (different formalism: classical fields, momentum space) was not made.
\item Kovchegov, Skokov, arXiv:1802.08166 (PRD 97 (2018) 094021): for single-gluon production ``it is impossible to have a cross section contribution which is odd under $k\to-k$ in the case of unpolarized initial and final states: \dots\ if we take the interference of a single $M_{1}$ with a single $M_{3}$ contribution \dots\ the color averaging would give zero'' (their Sec.\ on the odd-harmonic part at order $\alpha_{s}^{4}$, MV-type averaging); odd harmonics of the two-gluon correlation ``require three scatterings in the projectile and three scatterings in the target''. \emph{Consistent}: our odd part vanishes after the average over a C-even (MV) target or a $\pi$-rotation-symmetric ensemble, and over a projectile with $o=0$; it is cubic in the projectile and $O(\epsilon^{3})$ in a weak target field.
\item Kovner, Lublinsky, Skokov, arXiv:1612.07790: in double inclusive production at leading power the antisymmetric part is ``saturated by a charge conjugation odd `condensate', i.e.\ the Odderon'' of the target, odd under $S\to S^{\dagger}$, which ``averages to zero in the MV model''. \emph{Consistent} with the C-oddness of $\mf{G}$, $\mf{Q}_{d}$ in the target.
\item Jeon, Venugopalan, hep-ph/0503219 (PRD 71 (2005) 125003): weight functional with the cubic term $d_{abc}\rho^{a}\rho^{b}\rho^{c}/\kappa_{A}$ (their Eq.~(17)), which generates Odderon excitations; a concrete (local) model for $o$, Sec.~\ref{sec:dabc-ass}.
\item Li, Skokov, arXiv:2104.01879 (first saturation correction to single inclusive production; abstract only read): not compared.
\end{itemize}

\subsubsection{Numerical checks}\label{sec:dabc-num}
\noindent Scripts in \texttt{scratchpad/fix/dabc/} (\texttt{sun.py}: $SU(N)$ generators, $f$, $d$, random $V$ and $U^{ab}=2{\rm tr}(t^{a}Vt^{b}V^{\dagger})$; \texttt{invariants.py}; \texttt{dred.py}: chain reduction of the $d$-contracted structures; \texttt{dstructs.py}: the table; logs \texttt{*.log}). The colour structures are parsed from the TeX source by the parser of \texttt{fix\_singlet} (\texttt{scratchpad/fix/singlet/rows.py}); they occur verbatim in \texttt{Final\_terms\_fixed} and \texttt{Final\_terms\_checked} (\texttt{check\_sources.py}).
\begin{itemize}
\item \emph{Colour algebra} [\texttt{check\_delta.py}]: $\Delta$ of the table vs direct contraction $\sum_{\rm seq}C^{abc}_{\rm seq}d^{abc}$ of the parsed structures, $SU(3)$, $SU(4)$, $SU(5)$, three sets of random Wilson lines: largest relative deviation $2.4\times10^{-15}$; canonical (ordered-argument) forms $3\times10^{-15}$ [\texttt{canon.py}].
\item \emph{Invariants} [\texttt{t\_inv.py}, \texttt{check\_fund.py}, \texttt{fund.py}]: antisymmetry, coincident-point rules, $\mf{Q}_{d}(a,b;c,a)=-\mf{G}(a,b,c)$, $\mf{Q}_{d}(a,b;b,c)=\mf{G}|_{U\to U^{T}}$, global invariance ($\le2.4\times10^{-15}$); C-oddness under $V\to V^{*}$ (exact); fundamental forms~\eqref{eq:dabc-Qfund}, \eqref{eq:dabc-Gfund}: $\le2.3\times10^{-15}$ for $N_{c}=2,\dots,6$, 20 samples each; the unbiased least-squares fit to all 24 trace monomials gives the same coefficients for $N_{c}=4,5,6,7$ (fit residual $\le6\times10^{-15}$).
\item \emph{Operator level} [\texttt{check\_operator\_d.py}]: charges = generators on three lattice sites (explicit matrices), globally colour-invariant density matrix $e^{-H}/{\rm Tr}e^{-H}$ (ensembles of \texttt{fix\_singlet}); for every cubic structure and every assignment of its $\rho$ points to sites (random complex weights; coincident sites = contacts): exact ordered $\langle C\rho\rho\rho\rangle$ vs [contact terms of \texttt{fix\_singlet}] $+\frac{N_{c}}{(N_{c}^{2}-4)(N_{c}^{2}-1)}\Delta\,o$ with $o$ computed from the state. $SU(3)$ fundamental charges (C-odd, $\max|o|=1.11$, $f$-part $4\times10^{-18}$; $S^{abc}$ equals the (A2$'$) form to $3.5\times10^{-17}$): largest relative deviation $7.7\times10^{-16}$, the $d$-parts ranging from about $1\%$ to more than $100\%$ of the contact parts. Same with a chiral term $0.6f^{abc}\rho_{1}^{a}\rho_{2}^{b}\rho_{3}^{c}$ in $H$ ($g_{f}\neq0$, (A2$'$) violated): deviation up to $6\times10^{-2}$, removed ($\le2.7\times10^{-15}$) by adding $t\,\Phi\,g_{f}$. C-even ensemble: $o=0$, deviation $1.9\times10^{-15}$ (reduces to \texttt{fix\_singlet}).
\item \emph{Parity} [\texttt{check\_parity.py}]: $f(-s)=f(s)^{*}$ for $\mathcal{E},\widehat{\mathcal{E}},\Psi,\mathcal{J}$ (closed forms; $\mathcal{E}$ also by quadrature, $6\times10^{-17}$); for I.III $C_{A}$, III.V $\mathcal{L}_{\rm V}(y')$ (+c.c.) and II.I $N_{1a}+N_{1b}$ (integrand symmetrised under $w\leftrightarrow w'$), at 50 random configurations: symmetrised part real and odd under $k\to-k$ to rounding.
\item \emph{Weak target} [\texttt{check\_dilute.py}]: $\mf{G}$ and all thirteen $\Delta$ scale as $\epsilon^{3}$ (ratios $v(\epsilon)/v(\epsilon/2)\to8$; for II.I $N_{1a}$ and III.II with visible $O(\epsilon^{4})$ corrections).
\item \emph{Evolution terms} [\texttt{ev\_generic.py}]: Sec.~\ref{sec:dabc-evol}.
\end{itemize}

\subsubsection{Questions for the author}\label{sec:dabc-q}
\begin{itemize}
\item[(Q1)] (A2$'$) sets the $f^{abc}$ part $g_{f}$ of the symmetrised three-point function to zero. By (A1) alone $S^{abc}=d^{abc}g_{d}+f^{abc}g_{f}$, and C-evenness does not remove $g_{f}$ (\texttt{fix\_singlet}). Is $g_{f}=0$ intended (e.g.\ by an additional symmetry of the projectile)? If not, the odd part acquires $t\,\Phi\,g_{f}$ per structure, with $\Phi$ of \texttt{fix\_singlet} at separate points (verified).
\item[(Q2)] For $SU(2)$ the normalisation of (A2$'$) is $0/0$; we read (A2$'$) as $S^{abc}=0$ there. Agreed?
\item[(Q3)] Which model for $o$ is intended? For a local $o$ (Jeon--Venugopalan) two structures drop out and the integrability of the kernels at the triple coincidence of the $\rho$ points has to be checked; this was not done.
\item[(Q4)] The odd part is C-odd in the target: it vanishes after a C-even target average and survives event by event or with a target odderon. Is the result meant event by event (odd harmonics in correlations), or with a target odderon? This decides how the paper should phrase ``the spectrum is odd''.
\item[(Q5)] Presentation: should Eq.~\eqref{eq:dabc-odd} be part of the main result (Sec.~``Average over the projectile'') or an appendix? The drop-ins below put it in the main text, after the contact-term forms.
\end{itemize}

\subsubsection{Drop-in replacements}\label{sec:dabc-dropin}
\begingroup
\def\v#1{\mathbf{#1}}\providecommand{\Kv}{\mathcal{K}}\renewcommand{\eqref}[1]{(\texttt{\detokenize{#1}})}\renewcommand{\ref}[1]{\texttt{\detokenize{#1}}}\renewcommand{\cite}[1]{[\texttt{\detokenize{#1}}]}\providecommand{\pending}[1]{\textcolor{red}{[#1]}}
\noindent Labels, references and citations are printed as their keys here; in the source they are the paper's. $\mathcal{C}_{\rm f}=\mathcal{N}_{\rm f}$.

\paragraph*{(P1) Paper, \texttt{sec\_finite.tex}, Sec.~\ref{sec:fin_singlet}: item (A2) and the paragraph after the itemize}
\noindent Replace the item (A2) by:\par\smallskip
\noindent\fbox{\parbox{0.97\textwidth}{\small\begin{itemize}
\item[(A2)] \emph{Symmetrised three-point function of $d$-type:} with $S^{abc}(\v x,\v y,\v z)\equiv\langle\{\rho^a(\v x)\rho^b(\v y)\rho^c(\v z)\}_{\rm sym}\rangle$, where $\{\dots\}_{\rm sym}$ is the average over the six orderings,
\begin{align}\label{eq:A2d}
S^{abc}(\v x,\v y,\v z)=d^{abc}\,\frac{N_c}{(N_c^2-4)(N_c^2-1)}\,o(\v x,\v y,\v z)\,,\qquad o(\v x,\v y,\v z)\equiv\big\langle d^{efg}\{\rho^e(\v x)\rho^f(\v y)\rho^g(\v z)\}_{\rm sym}\big\rangle\,.
\end{align}
\end{itemize}}}\par\smallskip
\noindent and the paragraph ``(A2) holds for a Gaussian projectile \dots evolution terms.'' by:\par\smallskip
\noindent\fbox{\parbox{0.97\textwidth}{\small
The normalisation follows from $d^{abc}d^{abc}=(N_c^2-4)(N_c^2-1)/N_c$; $o$ is real and totally symmetric in its arguments. By (A1), $S^{abc}=d^{abc}g_d+f^{abc}g_f$ with $g_d$ totally symmetric and $g_f$ totally antisymmetric in the three points; (A2) states that $g_f=0$ and gives $g_d$ the normalisation of Eq.~\eqref{eq:A2d}. The function $o$ is C-odd: it is the three-point function of the cubic Casimir density, i.e.\ the odderon of the projectile. It vanishes for a Gaussian projectile (Gaussian Weyl symbol of the density matrix, as in the McLerran--Venugopalan model), whose odd moments vanish, and it is non-zero, for instance, for the weight functional with a cubic term $d_{abc}\rho^a\rho^b\rho^c$ of~\cite{Jeon:2005cf}, where it is local. For $SU(2)$, $d^{abc}=0$ and (A2) reduces to $S^{abc}=0$. Nothing else is assumed; in particular no assumption on the four-point function is needed, since the four-$\rho$ rows contribute only to the evolution terms.
}}\par\smallskip
\noindent [New bib entry needed: \texttt{Jeon:2005cf} = S.~Jeon, R.~Venugopalan, \emph{A classical Odderon in QCD at high energies}, Phys.\ Rev.\ D71 (2005) 125003, hep-ph/0503219. The key is our choice; please check against INSPIRE.]

\paragraph*{(P2) Paper, paragraph ``Quadratic and cubic terms'': the sentences after Eq.~(cubicavg)}
\noindent Replace ``Under (A2) the first term vanishes, \emph{but the contact terms do not}: \dots\ with the point $q$ replaced by $p$.'' by:\par\smallskip
\noindent\fbox{\parbox{0.97\textwidth}{\small
By (A2) the first term is $C^{abc}S^{abc}=\frac{N_c}{(N_c^2-4)(N_c^2-1)}\,C^{abc}d^{abc}\,o(p_1,p_2,p_3)$, and $2C^{abc}S^{abc}$ in the two forms with an anticommutator; we write it as $\frac{N_c}{(N_c^2-4)(N_c^2-1)}\Delta\,o(p_1,p_2,p_3)$, where $\Delta\equiv\sum C^{abc}_{\rm seq}d^{abc}$ is summed over the orderings into which the anticommutators expand. It vanishes for a Gaussian projectile, \emph{but the contact terms do not}: they are two-point terms $\propto\delta^{(2)}(p_i-p_j)\,n$, of the same order as the quadratic terms, and they are not removed by colour neutrality. \emph{The cubic terms therefore do not vanish after the average over a Gaussian projectile; only their symmetrised part does.} The contact terms are written below as quadratic terms at coincident points, $\mathfrak K|_{q\to p}$ denoting the kernel with the point $q$ replaced by $p$; the symmetrised parts are collected in Eq.~\eqref{eq:sodd}.
}}

\paragraph*{(P3) Paper, paragraph ``Parity in $\v k$'': the last three sentences}
\noindent Replace ``The symmetrised part of the cubic terms is odd under $\v k\to-\v k$, \dots\ but not for the commutator remainders of the ordered products.'' by:\par\smallskip
\noindent\fbox{\parbox{0.97\textwidth}{\small
The symmetrised part of the cubic terms is odd under $\v k\to-\v k$, their contact part is even, and the quadratic terms are even. The part of the finite terms that is odd in $\v k$ is therefore exactly the symmetrised part, Eq.~\eqref{eq:sodd}; under (A2) it is proportional to the projectile odderon $o$, and it vanishes for a Gaussian projectile. This refines the statement that the cubic, ``first saturation'' terms are odd in $\v k$~\cite{McLerran:2016snu,Kovchegov:2018jun}: it holds for their symmetrised part, configuration by configuration, but not for the commutator remainders of the ordered products; and the symmetrised part survives the average over a projectile with $o\neq0$.
}}

\paragraph*{(P4) Paper, paragraph ``Colour-singlet structures'': append after the coincident-point identities of $\mf D,\mf F,\mf Q$}
\noindent\fbox{\parbox{0.97\textwidth}{\small
The symmetrised parts of the cubic terms need, with $(D^e)_{bc}\equiv d^{ebc}$,
\begin{align}\label{eq:GQd}
\mf G(a,b,c)&\equiv f^{ijk}d^{lmn}\,U^{il}(a)\,U^{jm}(b)\,U^{kn}(c)\,,\qquad
\mf Q_d(a,b;c,e)\equiv i\,{\rm Tr}\{T^hU^\dagger(a)U(b)D^hU^\dagger(c)U(e)\}\,.
\end{align}
$\mf G$ is real and totally antisymmetric, $\mf Q_d$ is real, $\mf Q_d(a,b;c,e)=-\mf Q_d(e,c;b,a)$, $\mf Q_d(a,a;c,e)=\mf Q_d(a,b;c,c)=0$ and $\mf Q_d(a,b;c,a)=-\mf G(a,b,c)$. In terms of the fundamental Wilson lines $V$, $U^{ab}=2{\rm tr}(t^aVt^bV^\dagger)$, a Fierz rearrangement gives, exactly,
\begin{align}\label{eq:Gfund}
\mf G(a,b,c)&={\rm Im}\big[{\rm tr}(V^\dagger_aV_b)\,{\rm tr}(V^\dagger_bV_c)\,{\rm tr}(V^\dagger_cV_a)\big]+{\rm Im}\,{\rm tr}\big(V^\dagger_aV_bV^\dagger_cV_aV^\dagger_bV_c\big)\,,\notag\\
\mf Q_d(a,b;c,e)&=-{\rm Im}\big[{\rm tr}A\,{\rm tr}B\,{\rm tr}(A^\dagger B^\dagger)\big]-{\rm Im}\,{\rm tr}\big(ABA^\dagger B^\dagger\big)\,,\qquad A=V^\dagger_aV_b\,,\ B=V^\dagger_cV_e\,.
\end{align}
These are C-odd (odderon) structures of the target: under $V\to V^*$ they change sign, while $\mf D$, $\mf F$, $\mf Q$ do not; for a weak target they start at third order in the target field, and they vanish for $SU(2)$.
}}

\paragraph*{(P5) Paper: new paragraph after Eq.~(sIIIVI) and the sentence ``Rows~III.V and~III.VI are cubic \dots'' (before Table~\ref{tab:singletrows})}
\noindent\fbox{\parbox{0.97\textwidth}{\small
\paragraph{Symmetrised part of the cubic rows.} The colour factors $\Delta$ of the symmetrised parts are combinations of $\mf G$ and, in Rows~II.I and~III.II, of $\mf Q_d$. With the row prefactors and kernels of the contact forms above, the part of the finite terms odd in $\v k$ is
\begin{align}\label{eq:sodd}
\frac{d^3N_{\rm odd}}{d^3k}\bigg|_{\rm singlet}&=\frac{\mathcal C_{\rm f}}{N_c^2-1}\,\frac{N_c}{N_c^2-4}\Bigg\{\int_{w,w'}e^{-ik\cdot(w-w')}\int_{x,y,x'}\frac{1}{2\pi}\,\mathfrak K^{\rm II}_{2b}\Big[\mf G(x,y,x')-\mf G(x,y,w')\Big]\,o(x,y,x')\notag\\
&\quad+\int_{w,w'}e^{-ik\cdot(w-w')}\int_{x,y,z,x'}2i\Big\{\mathfrak K^{\rm III}_z\Big[\mf G(x,y,x')-\mf G(x,y,w')+\mf G(x,x',z)+\mf G(x,z,w')\Big]\notag\\
&\qquad\qquad+\mathfrak K^{\rm III}_w\Big[\mf G(x,y,x')-\mf G(x,y,w')+\mf G(x,x',w)+\mf G(x,w,w')\Big]\Big\}\,o(x,y,x')\notag\\
&\quad-\int_{x,z,y,y',x',w'}e^{-ik\cdot(z-w')}\,i\,\mathfrak K^{\rm V}_1\Big[\mf G(x,x',z)+\mf G(x,z,w')+\mf G(y,x',y')+\mf G(y,y',w')\Big]\,o(y,x',y')\notag\\
&\quad-\int_{x,y,x',z,w,w'}e^{-ik\cdot(w-w')}\,2i\,\mathfrak K_{\mathcal T}(w,w')\Big[\mf G(x,y,x')+\mf G(x,x',z)+\mf Q_d(z,y;x',w)\Big]\,o(x,y,x')\notag\\
&\quad-\int_{w,x',y',z,y,x}e^{-ik\cdot(w-y')}\,2i\,\Big\{\Big[-\mathfrak K_{\widehat{\mathcal E}}+\mathfrak K_\Psi+\mathfrak K^{(y')}_\delta\Big]\Big[\mf G(x,y,x')-\mf G(x,y,y')-\mf G(x',z,w)+\mf G(y',z,w)\Big]\notag\\
&\qquad\qquad+\mathfrak K^{(z)}_\delta\Big[\mf G(x,y,x')-\mf G(x,y,z)-\mf G(x',z,w)\Big]\Big\}\,o(x,y,x')\notag\\
&\quad-\int_{w,y',z,y,x,x'}e^{-ik\cdot(w-y')}\,2i\,\Big\{\Big[-\mathfrak K_{\widehat{\mathcal E}}+\mathfrak K_\Psi+\mathfrak K^{(y')}_\delta\Big]\Big[-\mf G(x,y,x')+\mf G(x,y,y')-\mf G(x,x',w)+\mf G(x,y',w)\Big]\notag\\
&\qquad\qquad+\mathfrak K^{(z)}_\delta\Big[-\mf G(x,y,x')+\mf G(x,y,z)-\mf G(x,x',w)+\mf G(x,z,w)\Big]\Big\}\,o(x,y,x')\notag\\
&\quad-\int_{w,y',z,y,x,x'}e^{-ik\cdot(w-y')}\,2i\,\Big\{\Big[-\mathfrak K_{\widehat{\mathcal E}}+\mathfrak K_\Psi+\mathfrak K^{(y')}_\delta\Big]\Big[-\mf G(x,y,x')+\mf G(x,y,y')+\mf G(y,x',z)-\mf G(y,y',z)\Big]\notag\\
&\qquad\qquad+\mathfrak K^{(z)}_\delta\Big[-\mf G(x,y,x')+\mf G(x,y,z)+\mf G(y,x',z)\Big]\Big\}\,o(x,y,x')\Bigg\}+{\rm c.c.}\notag\\
&\quad-\frac{\mathcal C_{\rm f}}{N_c^2-1}\,\frac{N_c}{N_c^2-4}\int_{x,y,z,w,x',w'}4\sin\big(\v k\cdot(\v w-\v w')\big)\,\mathfrak K_{\mathcal T}(w,w')\notag\\
&\qquad\qquad\times\Big[\mf G(x,z,w')-\mf G(x,y,x')-\mf Q_d(z,y;x',w)\Big]\,o(x,y,x')\,.
\end{align}
The lines are, in this order, Rows~I.II ($N_{2b}$), I.III, I.V (first part), III.II, III.III (first part), III.V and III.VI, each with its complex conjugate, and the first part of Row~II.I ($N_{1a}+N_{1b}$, after $\v w\leftrightarrow\v w'$ in $N_{1b}$). The pairs $\mathfrak K^{(y')}_\delta$, $\mathfrak K^{(z)}_\delta$ stay under one $\v y'$ integral. None of the thirteen colour factors vanishes identically. Each term is real and odd under $\v k\to-\v k$; it is C-odd in the projectile ($o$) and in the target ($\mf G$, $\mf Q_d$), so that it vanishes after an average over a C-even target ensemble and survives configuration by configuration or for a target with an odderon component; since all kernels are invariant under the simultaneous rotation by $\pi$ of all coordinates and of $\v k$, it also integrates to zero if the averaged product of the colour factor and $o$ is invariant under the reflection of all points through a common centre. The soft-remainder pairs of Section~\ref{sec:fin_soft} pair in Eq.~\eqref{eq:sodd} as in the contact forms, since Eqs.~\eqref{eq:colid} are operator identities. The colour factors were obtained by contracting the colour structures of Appendix~\ref{app:finite} with $d^{abc}$ and verified for $SU(3)$, $SU(4)$, $SU(5)$ (to $2\times10^{-15}$), and the complete average (contact terms and symmetrised part) at the operator level with $SU(3)$ charges in a C-odd, colour-invariant state, for which (A2) holds with $o\neq0$ (to $8\times10^{-16}$).
}}

\paragraph*{(P6) Paper, Table~\ref{tab:singletrows}: column header and caption}
\noindent Header ``after (A1), (A2)'' $\to$ ``after (A1), (A2): contact terms''; append a column ``odd part'' with entries: I.II $N_{2b}$, I.III, I.V first part, II.I first part, III.II, III.III first part, III.V, III.VI: \eqref{eq:sodd}; quadratic rows: ---. Caption:\par\smallskip
\noindent\fbox{\parbox{0.97\textwidth}{\small
Content of the rows after the average over the projectile. For every cubic row the contact terms are even in $\v k$, real after the complex conjugate is added, and do not vanish for a Gaussian, colour-neutral projectile; the symmetrised part is odd in $\v k$, proportional to the projectile odderon $o$ of Eq.~\eqref{eq:A2d}, and given in Eq.~\eqref{eq:sodd}.
}}

\paragraph*{(P7) Paper, Remarks (iii) after the table and the sentence in Sec.~\ref{sec:fin_summary}}
\noindent Remark (iii) $\to$\par\smallskip
\noindent\fbox{\parbox{0.97\textwidth}{\small
(iii) A chiral part $f^{abc}g_f$ of the symmetrised three-point function, excluded by (A2), would add to each cubic structure $t\,\Phi\,g_f(p_1,p_2,p_3)$, with $\Phi$ the factor of the contact terms evaluated at separate points and $t=1$ ($2$) for the ordered (anticommutator) forms; it is also odd in $\v k$. For a local $o$, as in~\cite{Jeon:2005cf}, the integrability of the kernels at the triple coincidence of the charge positions has not been examined.
}}\par\smallskip
\noindent In Sec.~\ref{sec:fin_summary} replace ``After the average over a Gaussian, colour-singlet projectile the finite terms are even in $\v k$; the three-$\rho$ rows contribute through the contact terms of the commutators of the ordered charges.'' by\par\smallskip
\noindent\fbox{\parbox{0.97\textwidth}{\small
After the average over a colour-singlet projectile the three-$\rho$ rows contribute through the contact terms of the commutators of the ordered charges, which are even in $\v k$, and through their symmetrised part, Eq.~\eqref{eq:sodd}, which is odd in $\v k$ and proportional to the projectile odderon $o$; for a Gaussian projectile $o=0$ and the finite terms are even in $\v k$.
}}

\paragraph*{(P8) Paper, \texttt{sec\_checks.tex}: paragraphs ``Counting of charges'' and ``Odd terms''}
\noindent ``Counting of charges'', last two sentences $\to$\par\smallskip
\noindent\fbox{\parbox{0.97\textwidth}{\small
The finite terms are of order $g^4\rho^2$ and $g^4\rho^3$. After the average over the projectile the $g^4\rho^3$ terms do not drop out: their symmetrised part is proportional to the symmetrised three-point function, i.e.\ to the projectile odderon $o$ of Eq.~\eqref{eq:A2d}, and the commutators produced by the reordering of the charges leave contact terms of order $g^4\langle\rho\rho\rangle$, of the same order as the quadratic finite terms (Section~\ref{sec:fin_singlet}). The three-$\rho$ terms that multiply a rapidity logarithm cancel in the symmetrised sense for any three-point function, so that no logarithm multiplies $o$.
}}\par\smallskip
\noindent ``Odd terms'' $\to$\par\smallskip
\noindent\fbox{\parbox{0.97\textwidth}{\small
\paragraph{Odd terms.} The three-$\rho$ terms carry an explicit factor $i$ (in Row~I.II, $N_{2b}$, it is in the Fourier transform of an odd kernel), and all kernels depend on $\v k$ through functions with $f(-s)=f(s)^*$. A three-$\rho$ term plus its complex conjugate is therefore odd under $\v k\to-\v k$ in its totally symmetrised part and even in the part produced by the commutators of the ordered charges (Section~\ref{sec:fin_singlet}). For a globally colour-invariant projectile with the symmetrised three-point function~\eqref{eq:A2d} the odd part is Eq.~\eqref{eq:sodd}: it is proportional to the projectile odderon $o$ and contains the target only through the C-odd structures $\mf G$ and $\mf Q_d$, Eq.~\eqref{eq:Gfund}, the imaginary parts of products of three dipoles and of sextupoles. It therefore requires a C-odd three-point function of the projectile, it vanishes after an average over a C-even target ensemble, and for a weak target it starts with three gluon exchanges; for a Gaussian (McLerran--Venugopalan-type) projectile, $o=0$, the spectrum is even in $\v k$. This agrees with the classical-field result that the leading-order spectrum is even configuration by configuration and that the odd part comes from the term cubic in the projectile charge, averages out for a Gaussian projectile and is ``defined by odderon exchanges''~\cite{McLerran:2016snu}, and with the counting of~\cite{Kovchegov:2018jun}, three scatterings in the projectile and three in the target. The even commutator part does \emph{not} vanish: the three-$\rho$ rows contribute to the averaged spectrum through two-point contact terms $\propto\delta^{(2)}(\v x_i-\v x_j)\,n$, which have no counterpart for classical charges. All colour factors of the averaged terms, the symmetrised parts and the parity statements were verified at the operator level, with charges represented by generators on two or three sites and globally colour-invariant density matrices, including a C-odd state with $o\neq0$ (Section~\ref{sec:fin_singlet}). A comparison of the three-$\rho$ terms with the dilute limit of the first saturation correction to single inclusive production of Li and Skokov~\cite{Li:2021ntm}, and a term-by-term comparison of Eq.~\eqref{eq:sodd} with the odd part of~\cite{McLerran:2016snu}, have not been carried out.
}}

\paragraph*{(P9) Paper, \texttt{sec\_rapidity.tex}, Sec.~\ref{sec:threerho}: append to the paragraph ``For each of the two logarithms \dots''}
\noindent\fbox{\parbox{0.97\textwidth}{\small
Since the colour tensor itself vanishes, the cancellation holds for any symmetrised three-point function of the projectile, in particular for the $d^{abc}$ (odderon) one of Eq.~\eqref{eq:A2d}, which we checked separately with $SU(3)$ and $SU(4)$ matrices: no rapidity logarithm multiplies the three-point function. This is required, since the LO spectrum depends on the projectile only through the C-even two-point function $n$, and its one-step evolution cannot produce the C-odd $o$.
}}

\paragraph*{(P10) Paper, conclusions and abstract}
\noindent \texttt{sec\_conclusions.tex}, replace ``In the averaged form the terms cubic in the charges do not vanish for a Gaussian projectile: only their symmetrised part, which is odd in $\v k$, does, and the commutators of the ordered charges leave contact terms of the same order as the quadratic terms. For such a projectile the spectrum is even in $\v k$.'' by\par\smallskip
\noindent\fbox{\parbox{0.97\textwidth}{\small
In the averaged form the terms cubic in the charges contribute in two ways. Their symmetrised part is proportional to the symmetrised three-point function of the projectile, which for a colour-singlet projectile is the $d^{abc}$ (odderon) correlator $o$; it is odd in $\v k$, and its colour factors are the C-odd target structures ${\rm Im}[{\rm tr}\,{\rm tr}\,{\rm tr}]$ and ${\rm Im}\,{\rm tr}$ of three dipoles and of a sextupole. The commutators of the ordered charges leave contact terms of the same order as the quadratic terms, which are even in $\v k$ and do not vanish for a Gaussian projectile. For a Gaussian projectile, or after an average over a C-even target, the spectrum is even in $\v k$. No rapidity logarithm multiplies the three-point function.
}}\par\smallskip
\noindent Abstract, replace the last sentence by\par\smallskip
\noindent\fbox{\parbox{0.97\textwidth}{\small
The terms cubic in $\rho$ do not vanish after the average over the projectile: their symmetrised part, odd in the transverse momentum, is proportional to the projectile odderon and to C-odd target structures, and vanishes only for a Gaussian projectile, while the commutators of the charges leave contact terms of the same order as the quadratic terms.
}}\par\smallskip
\noindent Introduction, last bullet, replace ``After the average over a Gaussian projectile the terms cubic in $\rho$ do not vanish: only their symmetrised part, which is odd in $\v k$, does, while the commutators of the ordered charges leave contact terms (Section~\ref{sec:finite}).'' by ``After the average over the projectile the terms cubic in $\rho$ do not vanish: their symmetrised part, odd in $\v k$, is proportional to the $d^{abc}$ (odderon) three-point function of the projectile, and the commutators of the ordered charges leave contact terms, even in $\v k$, which survive also for a Gaussian projectile (Section~\ref{sec:finite}).''\par
\noindent Optionally, \texttt{sec\_observable.tex} after Eq.~(singlet): ``Where the three-point function is needed we use Eq.~\eqref{eq:A2d}.''

\paragraph*{(F1) \texttt{Final\_terms\_fixed}, fragment \texttt{build/src/singlet.tex}}
\noindent (a) Item (A2) of ``Projectile average: what is assumed'' $\to$ (A2$'$) of Sec.~\ref{sec:dabc-ass} of this note (with $\langle\{\dots\}_{\rm sym}\rangle$ as defined there), and the paragraph ``(A2) holds for a Gaussian projectile \dots'' $\to$ the first paragraph of Sec.~\ref{sec:dabc-ass} (``By (A1) \dots\ $SU(2)$ \dots''), ending with ``No assumption about the four-point function is needed (four-$\rho$ rows contribute only to the evolution terms).''\par
\noindent (b) ``Cubic terms'' paragraph: ``Under (A2) the first term vanishes, \emph{but the contact terms do not}:'' $\to$ ``Under (A2$'$) the first term is $\frac{N_{c}}{(N_{c}^{2}-4)(N_{c}^{2}-1)}\Delta\,o(p_{1},p_{2},p_{3})$, $\Delta\equiv\sum_{\rm orderings}C^{abc}d^{abc}$ (Sec.~``Symmetrised part'' below); it vanishes for a Gaussian projectile, \emph{but the contact terms do not}:''.\par
\noindent (c) ``Parity in $k$'', last sentence $\to$ ``The part of the finite terms odd in $k$ is exactly the $S^{abc}$ part; under (A2$'$) it is given in Sec.~``Symmetrised part'' below; it vanishes for $o=0$.''\par
\noindent (d) New subsection before ``Which rows are cubic, which quadratic'': ``\textbf{Symmetrised part of the cubic rows ($d^{abc}$)}'' = Secs.~\ref{sec:dabc-inv}, \ref{sec:dabc-table} (with the table and Eq.~\eqref{eq:dabc-odd}) and~\ref{sec:dabc-parity} of this note, verbatim (kernel names as in the fragment).\par
\noindent (e) Table ``Which rows are cubic, which quadratic'': header ``after the average (A1), (A2)'' $\to$ ``after the average (A1), (A2$'$): contact terms''; paragraph after the table: ``For every cubic row the symmetrised part ($S^{abc}$) is zero under (A2) and would be odd in $k$ otherwise;'' $\to$ ``For every cubic row the symmetrised part is $\propto o$, odd in $k$ and given in Eq.~(dabc-odd);''.\par
\noindent (f) ``Numerical checks'': add the items of Sec.~\ref{sec:dabc-num}. ``What is not settled'' (5) $\to$ ``(5) (A2$'$) is an assumption on the projectile: $g_{f}=0$ (Q1); for a local $o$ the integrability at the triple coincidence of the $\rho$ points is not checked (Q3).''
\endgroup

\clearpage
\appendix
\section{Earlier intermediate round (before the checked copies)}\label{app:early}
\noindent Before the systematic checked copies of Table~\ref{tab:pairs} existed, a first round of corrections was made to earlier versions of the author's \texttt{Final\_terms.tex}, \texttt{Limiting\_GIII\_RI.tex}, \texttt{Group\_III\_RI.tex}, \texttt{Other\_terms.tex}, \texttt{Limiting\_Cases.tex} (folders \texttt{nlo\_review/}, \texttt{nlo\_review/rowIII\_fix/}, \texttt{stage1/}, \texttt{stage2/}, \texttt{final/}). Its list is \texttt{nlo\_review/CHANGES\_to\_propagate.md}; it is summarised here for completeness. These changes are superseded by the checked copies of Section~\ref{sec:files} (in particular the replacement of the last term of $N_{3b}$ by the recoil term and the principal-value treatment), and the files of that round were not re-extracted line by line.
\begin{itemize}
\item \texttt{Final\_terms.tex}: macros \texttt{\textbackslash chg}, \texttt{\textbackslash rnote}; a ``Conventions'' block (definitions of $K^{i}(x)$, $K(x,y,z)$, phase, the three plus prescriptions, scheme, principal value); $\int_{x,,z,w,w'}\to\int_{x,y,z,w,w'}$ and $\mathcal{K}(x,y,z)\to K(x,y,z)$ in the JIMWLK singlet formula; after $\widetilde{\mathcal P}_{gg}$: $\beta_{0}=\frac{11N_{c}}{3}$ and the statement that both plus prescriptions act on $G(\xi)=\frac{1}{\xi^{2}}\int e^{-ik\cdot(y'-y)/\xi}[\dots]$ with $G(0)=-\frac{4\pi}{k^{2}}\int_{x,x',z}K(x-z)\cdot K(x'-z)[U(z)-U(x')][U(z)-U(x)]\rho\rho$; running coupling $=-$ the $\frac{\beta_{0}}{2}\delta(1-\xi)$ part of DGLAP; ``after removing the poles in $\overline{\rm MS}$''; note on symmetrising $\rho\rho\rho$; note on $p\to0$ of $-2q^{i}\log(p^{2}/q^{2})/p^{2}$ (Group I Row II); Group I Row III: third block (prefactor $-g^{4}f^{abd}/(8\pi^{4})$, no log) deleted (principal value), second block kept with a note.
\item \texttt{Limiting\_GIII\_RI.tex}, \texttt{Group\_III\_RI.tex}: last term of $N_{3b}$: $U^{bc}(z)\to U^{bc}(w)$; for principal value every $i\pi$ is absent, the block with prefactor $g^{4}f^{abd}/8\pi^{4}$ dropped; logs of negative arguments are $\log|\cdot|$; status notes.
\item \texttt{rowIII\_fix/}: Group I Row III last term replaced by $\frac{1}{(z-w)^{2}}\log\frac{(z-w)^{2}}{(x-w)^{2}}K^{i}(x'-w')\{K^{i}(y-z)[U(z)-U(y)]-\frac{(x-w)^{2}}{(x-w)^{2}+(z-w)^{2}}K^{i}(y-w)[U(w)-U(y)]\}$ plus the contact term $\frac{\pi^{3}}{6}K(x'-w')\cdot K(y-w)[U(w)-U(y)]$; Group III Row V: the $-([\frac{k^{+}}{p^{+}(p^{+}+k^{+})}]_{+}+\frac{1}{p^{+}+k^{+}})$ term replaced by the closed form $-\frac{ig^{4}f^{c'ca}}{8\pi^{4}k^{+}}\ln^{2}\frac{2e^{-\gamma_{E}}}{|k||x-w'|}\dots$ (and its $\log(1-\frac{k^{+}}{\vee})$ term removed); \texttt{Limiting\_Cases.tex}: the $i\pi$ block of the summed Row III removed; open then: Group III Rows III and VI. Checks: \texttt{rowIII\_fix/checks/} (\texttt{v1\_rowIII.py}, \texttt{v2\_rowV.py}, \texttt{v3\_j0.py}).
\end{itemize}
\noindent These earlier forms were \emph{not} kept: the final treatment of Group I Row III is the combined kernel with the recoil term (Sec.~\ref{sec:calc-recoil}, \ref{sec:fix-soft}), and Group III Rows III, V, VI are in the author's split with the three-term BFKL list (Sec.~\ref{sec:calc-G3}).

\end{document}
